\documentclass[11pt,oneside,openany]{book}
\usepackage[a4paper,margin=1in,headheight=14pt]{geometry}
\usepackage{fontspec}
\setmainfont{Linux Libertine O}[Ligatures=TeX]
\setsansfont{Linux Biolinum O}
\usepackage{xeCJK}
\setCJKmainfont{Noto Serif CJK SC}
\setCJKsansfont{Noto Sans CJK SC}
\setCJKmonofont{Noto Sans CJK SC}
\usepackage[final]{microtype}
\usepackage{setspace}
\setstretch{1.08}
\setlength{\emergencystretch}{3em}
\usepackage{amsmath,amssymb,mathtools,mathrsfs}
\usepackage{graphicx}
\usepackage{enumitem}
\setlist{nosep,leftmargin=1.45em}
\usepackage{booktabs,tabularx,longtable,array}
\newcolumntype{Y}{>{\raggedright\arraybackslash}X}
\setlength{\tabcolsep}{5pt}
\usepackage{xcolor}
\definecolor{ink}{HTML}{000000}
\definecolor{muted}{HTML}{333333}
\definecolor{line}{HTML}{B8B8B8}
\definecolor{boxbg}{HTML}{FAFAFA}
\definecolor{accent}{HTML}{000000}
\definecolor{stopline}{HTML}{000000}
\definecolor{stopbg}{HTML}{FAFAFA}
\color{ink}
\usepackage{titlesec}
\titleformat{\part}[display]
  {\normalfont\sffamily\bfseries\raggedright\color{ink}}
  {\fontsize{10}{12}\selectfont\color{muted}\MakeUppercase{Part \thepart}}
  {8pt}{\fontsize{24}{28}\selectfont}
\titlespacing*{\part}{0pt}{0pt}{24pt}
\titleformat{\chapter}[display]
  {\normalfont\sffamily\bfseries\raggedright\hyphenpenalty=10000\exhyphenpenalty=10000\color{ink}}
  {\fontsize{10}{12}\selectfont\color{muted}\MakeUppercase{Chapter \thechapter}}
  {5pt}
  {\fontsize{22}{25}\selectfont}
\titleformat{name=\chapter,numberless}[display]
  {\normalfont\sffamily\bfseries\raggedright\hyphenpenalty=10000\exhyphenpenalty=10000\color{ink}}
  {}{0pt}{\fontsize{20}{23}\selectfont}
\titleformat{\section}
  {\normalfont\sffamily\bfseries\fontsize{15}{18}\selectfont\color{ink}}
  {\thesection}{0.65em}{}
\titleformat{\subsection}
  {\normalfont\sffamily\bfseries\fontsize{12.7}{15}\selectfont\color{ink}}
  {\thesubsection}{0.65em}{}
\titleformat{\paragraph}[runin]
  {\normalfont\sffamily\bfseries\color{ink}}
  {}{0pt}{}[.]
\titlespacing*{\chapter}{0pt}{-18pt}{22pt}
\titlespacing*{\section}{0pt}{23pt}{9pt}
\titlespacing*{\subsection}{0pt}{15pt}{6pt}
\usepackage{fancyhdr}
\pagestyle{fancy}
\fancyhf{}
\fancyfoot[C]{\small\thepage}
\renewcommand{\headrulewidth}{0pt}
\renewcommand{\footrulewidth}{0pt}
\fancypagestyle{plain}{%
  \fancyhf{}%
  \fancyfoot[C]{\small\thepage}%
  \renewcommand{\headrulewidth}{0pt}%
  \renewcommand{\footrulewidth}{0pt}%
}
\usepackage[most]{tcolorbox}
\tcbset{enhanced,breakable,colback=boxbg,colframe=line,boxrule=0.55pt,arc=1.5mm,left=3.5mm,right=3.5mm,top=2.6mm,bottom=2.8mm,before skip=10pt,after skip=12pt,fonttitle=\sffamily\bfseries\color{accent},coltitle=accent}
\newtcolorbox{sourcebox}[1][]{title={#1}}
\newtcolorbox{methodbox}[1][]{title={#1},borderline west={2pt}{0pt}{accent}}
\newtcolorbox{readerbox}[1][]{title={#1},borderline west={2pt}{0pt}{accent}}
\newtcolorbox{plainbox}[1][]{title={#1}}
\newtcolorbox{principlebox}[1][]{title={#1},borderline west={2pt}{0pt}{accent}}
\newtcolorbox{stopbox}[1][]{title={#1},colback=stopbg,colframe=stopline!55,coltitle=stopline,borderline west={2pt}{0pt}{stopline}}
\usepackage{amsthm}
\theoremstyle{definition}
\newtheorem{definition}{Definition}[chapter]
\newtheorem{example}[definition]{Example}
\newtheorem{principle}[definition]{Principle}
\newtheorem{proposition}[definition]{Proposition}
\newtheorem{theorem}[definition]{Theorem}
\newtheorem{lemma}[definition]{Lemma}
\newtheorem{corollary}[definition]{Corollary}
\newtheorem{warning}[definition]{Stop line}
\newtheorem{correspondence}[definition]{Structural correspondence}
\newtheorem{nonexample}[definition]{Non-example}
\usepackage{tikz}
\usetikzlibrary{arrows.meta,positioning,calc}
\usepackage[authoryear,round]{natbib}
\usepackage{xurl}
\usepackage{hyperref}
\hypersetup{hidelinks,bookmarksdepth=subsection,pdftitle={Structural Compatibility between Third-Order Consciousness and AF-by-Discrete: From the Organization of Experience, “Again,” and the Problem of the Observer to Groupoid Formalization},pdfauthor={Zheng Kuang},pdfsubject={Theoretical and formal cognitive science research monograph},pdfkeywords={Third-Order Consciousness, AF-by-discrete, cognitive science, structural compatibility, consciousness, self-reference, metacognition, observer position, return, event segmentation}}
\usepackage{bookmark}
\makeatletter
\renewcommand{\@pnumwidth}{3.0em}
\renewcommand{\@tocrmarg}{3.8em}
\makeatother
\renewcommand{\bibname}{References}
\newcommand{\AFBD}{AF-by-discrete}
\newcommand{\CN}[1]{{#1}}
\newcommand{\strsim}{\mathrel{\sim_{\mathrm{str}}}}
\newcommand{\ordone}{First-Order Consciousness}
\newcommand{\ordtwo}{Second-Order Consciousness}
\newcommand{\ordthree}{Third-Order Consciousness}
\newcommand{\Id}{\mathrm{Id}}

\begin{document}

\begin{titlepage}
\thispagestyle{empty}
\begin{center}
\vspace*{0.65in}
{\bfseries\LARGE Structural Compatibility between Third-Order Consciousness and AF-by-Discrete\par}
\vspace{0.6em}
{\large\itshape From the Organization of Experience, ``Again,'' and the Problem of the Observer to Groupoid Formalization\par}
\vspace{0.65in}
{\large Zheng Kuang\par}
\vspace{0.4em}
{\normalsize School of Mathematics, South China University of Technology\par}
{\normalsize 381 Wushan Rd, Tianhe District, Guangzhou 510641, China\par}
{\normalsize Email: \href{mailto:kzkzkzz@scut.edu.cn}{kzkzkzz@scut.edu.cn}\par}
{\normalsize ORCID: \href{https://orcid.org/0000-0001-6581-8207}{0000-0001-6581-8207}\par}
\vfill
{\small September 2026\par}
\vspace{0.35em}
{\small Copyright \textcopyright{} 2026 Zheng Kuang. All rights reserved.\par}
\end{center}
\end{titlepage}

\clearpage
\pagenumbering{arabic}
\setcounter{page}{1}

\chapter*{Sources, Attribution, and\\ Research Contributions}
\addcontentsline{toc}{chapter}{Sources, Attribution, and Research Contributions}

\begin{sourcebox}[Source of the Third-Order Consciousness framework]
The basic framework of \ordone, \ordtwo, and \ordthree{} in Part I---together with its accounts of how the subject is formed, why the observer is not neutral, ``the observer is the observed,'' choiceless awareness, and why the third order is not a spiritual rank---mainly consists of a review, distillation, and structured restatement of Guoxin Peng, \emph{Third-Order Consciousness: From the Biological Self and Social Self to Vacant Consciousness (The Geometry of Existence Book 4)}, edited by Pinchuan Xie and independently published in 2026 \citep{toc2026}. These ideas were not originally proposed by Zheng Kuang. The work of Part I is a reconstruction that seeks to remain as faithful as possible while permitting criticism and theoretical reduction: it preserves the problem-structure of the original work, supplies philosophical, psychological, and sociological coordinates, and examines which claims already have neighboring theories and which remain independent candidates.
\end{sourcebox}

\begin{sourcebox}[The research contribution of this project]
The research contribution of this project lies chiefly at the interface between the two native domains: it reconstructs for this book a self-contained and verifiable \AFBD{} mathematical domain from first definitions; distinguishes unlabeled background, endogenous reference positions, birooted segments, return-completion, persistence across levels, and local continuation data; establishes, item by item, reviewable relations of structural compatibility under the translation protocol and principles of the Geometry of Existence proposed by \emph{Snowfall on the Clock}; introduces the finite witness, return filtration, and residue bookkeeping used in this book, while restoring incompressibility to its strict source-expression convention; and designs counterexamples that can make the model fail, together with future empirical bridges. Here ``reconstruction'' means providing an independently readable mathematical foundation for this book. Whenever existing definitions or theorems come from the literature or from the author's earlier papers, their sources are stated explicitly in the main text.
\end{sourcebox}

\begin{sourcebox}[AI/LLM disclosure and division of labor]
The substantive content and claim structure of this English manuscript are frozen in the Bilingual Research Edition identified below. ChatGPT was used in manuscript preparation for literature retrieval, structural organization, English translation and sentence restructuring, drafting assistance, and typesetting. AI output was not treated as evidence or as an independent source of theoretical or mathematical claims. Zheng Kuang formulated the substantive research questions and mathematical materials, supplied the key corrections and cross-domain claim boundaries, reviewed the English manuscript under the frozen-source protocol, and made the final decisions. The author takes full responsibility for the claims, citations, wording, and errors retained in the final text. No structural translation made in this manuscript retroactively rewrites the original meanings of \emph{Third-Order Consciousness}, Krishnamurti's texts, the psychological literature, or the mathematical literature.
\end{sourcebox}

\begin{sourcebox}[On the Language and Genesis of the Manuscript]
The interweaving of Chinese and English in the Bilingual Research Edition should not be read as ``translation not yet completed.'' Third-Order Consciousness, relational experience, and phenomenological questions entered the project mainly in Chinese; the author's mathematical training, \AFBD{} research, verification of definitions, and most formal derivations have long unfolded in an English-language mathematical setting. Chinese therefore preserves much of the experiential texture in which the ideas first appeared, while English technical vocabulary preserves the mathematical working language actually used by the author. Their repeated encounters in the main text record how this cross-domain research in fact took place.

This fully English edition unifies the language for an international scholarly readership, but translation may not alter the mathematical content, claim levels, or provenance frozen in the Bilingual Research Edition. The mixed Chinese--English edition is therefore preserved as the Bilingual Research Edition of the project and as its canonical reference.
\end{sourcebox}

\section*{Canonical Bilingual Research Edition}
\addcontentsline{toc}{section}{Canonical Bilingual Research Edition}
The canonical reference for this English edition is the frozen Bilingual Research Edition deposited at Zenodo, DOI \href{https://doi.org/10.5281/zenodo.22202015}{10.5281/zenodo.22202015} \citep{zenodo22202015}. Whenever wording in the English edition raises a question about claim strength, mathematical notation, ordering, provenance, or the boundary between the two native domains, that deposit takes precedence over later stylistic smoothing.

Its mixed Chinese--English form has a simple historical reason. The psychological material entered the project mainly in Chinese: the discussions of consciousness, subject formation, relationships, memory, and the phenomenology of ``again'' were first thought through in Chinese. The mathematical work developed in the opposite direction. The author's training and research on groups, groupoids, Bratteli diagrams, bounded type, returns, germs, and AF-by-discrete were conducted largely in English, and much of that material entered the manuscript directly in English. In other words, the two languages did not divide the book by editorial design; they record the two working languages in which the project actually grew.

The bilingual manuscript therefore preserves more than an earlier draft. It records the history of the translation itself: psychological questions moving toward mathematics, mathematical distinctions moving back toward psychology, and the points at which one language resisted a convenient but misleading word from the other. The code-switching can look slightly eccentric---and occasionally it is funny---but it is not decorative code-switching and, for once, it is not a typesetting accident. It is part of the intellectual history of the project. The present English edition may improve cadence, syntax, and readability, but it may not silently rewrite that history or the content frozen there.

\begin{sourcebox}[On the Primary Materials from \emph{Snowfall on the Clock} and the Post-Level Index]
The archived compilation of \emph{Third-Order Consciousness} preserved by this project runs to 619 pages. Between 15 March and 10 May 2026, it contains 83 identifiable WeChat Official Account texts bearing publication dates and bylines. The compilation also includes several collateral essays written during the process and a general preface added later. The account navigation at the beginning of the compilation points to an ``Opening and General Preface'' not separately included in the archived PDF currently preserved by the project. The number 83 therefore denotes only the identifiable texts in the current archive; it does not claim to equal every related post ever published by the account.

On 30 August 2026, the project conducted post-by-post URL recovery for 96 unique permanent WeChat links preserved by the author of this book. Every link was opened, and its title and body were checked. Eighty-three links matched TOC-001--TOC-083 one-to-one, with no missing item and no duplicate assignment; the ``Opening and General Preface'' absent from the 619-page compilation was also recovered. The order in which the input links had been pasted was not used to reconstruct publication chronology. Because the live HTML exposed no reliable \texttt{publish\_time} field, the dates of the 83 records continue to follow the archive metadata preserved in the canonical bilingual deposit rather than being guessed anew from the webpages. The author subsequently supplied original permanent WeChat links for four \emph{Snowfall on the Clock} essays cited directly in the main text but absent from the preceding set of 96 crawled links. In this edition, those four are recorded bibliographically as ``URLs supplied subsequently by the author of this book'' and are not conflated with the preceding post-by-post HTML verification record.

Stable source IDs such as TOC-067 and the page numbers of the original compilation remain unchanged in the main text. The back-of-book \emph{Primary-Source Index to Third-Order Consciousness: Verified Permanent WeChat Links} gives the clickable original WeChat page associated with every source ID. The index organizes metadata, source location, URLs, and the relation of each source to the book; it does not republish the full text of those posts.
\end{sourcebox}

\chapter*{Abstract}
\addcontentsline{toc}{chapter}{Abstract}

Imagine sending someone a message and receiving no reply for a long time. Repeated checking may gradually recruit bodily tension, language, memory, identity, and judgment: ``Why isn't he replying?'', ``Why is it like this again?'', and eventually ``Why does this always happen to me?'' Taking \emph{Third-Order Consciousness: From the Biological Self and Social Self to Vacant Consciousness} by Guoxin Peng as its psychological-conceptual native domain, this book describes such shifts through three organizing foci: direct presence, subject-forming organization, and the observing position itself entering observation. ``Third order'' concerns the organization of experience, not a hierarchy of persons or a ladder toward higher consciousness.

The book asks what is meant by repetition, ``again,'' return, completion, and observer. Similarity does not by itself establish that the same pattern has returned, and a completed episode does not determine what follows. \AFBD{} mathematics is introduced as a second native domain. Its guiding picture is that a broad background may continue to unfold while relations that alter local continuation are concentrated at finite boundary support and a small number of persistent fissures. These fissures are detected where local organization can no longer continue unchanged and where apparently similar returns lead to distinguishable continuations.

The proposed bridge is structural compatibility, not identity, neural realization, or ontology. It yields a falsifiable finiteness hypothesis: subject-forming complexity may be concentrated at a small number of recurrent organizing interfaces. If so, these interfaces should remain re-identifiable as material accumulates and should predict stable differences in memory retrieval, question generation, action selection, or expected continuation. If they proliferate without limit, return reduces to content similarity, episode segmentation is not reproducible, or the formal language adds no observable distinctions or failure conditions, the bridge should be withdrawn.

This is therefore a conditional exercise in deductive theory construction. Mathematics supplies distinctions and failure tests; it does not prove \emph{Third-Order Consciousness} or claim a neural mechanism, clinical model, or ontological structure of consciousness.

\noindent\textbf{Keywords.} Third-Order Consciousness; \AFBD; structural compatibility; subject; observer; return; incompressibility; Geometry of Existence.

\setcounter{tocdepth}{0}
\tableofcontents
\clearpage

\chapter*{Reader's Guide: Two Native Domains, One Interface, Three Routes Through the Book}
\addcontentsline{toc}{chapter}{Reader's Guide: Two Native Domains, One Interface, Three Routes Through the Book}
\begingroup\small

This book keeps apart its two native domains on purpose. Part I first establishes the concepts of Third-Order Consciousness on their own terms within psychological, philosophical, and social-theoretical contexts. Part II then establishes the mathematical objects independently through their own definitions and proofs. Only in Part III do the two sides meet. This order is not a typographical convention. It is the most basic safeguard against category error in this book. The printed table of contents lists parts and chapters only, so that the main text is not crowded by more than ten pages of detailed contents; the PDF bookmarks retain section and subsection levels for close navigation.

\begin{center}
\resizebox{0.98\textwidth}{!}{%
\begin{tikzpicture}[
  node distance=7.5mm and 6mm,
  block/.style={draw=line,rounded corners=2mm,fill=boxbg,text width=0.39\textwidth,minimum height=16mm,align=center,font=\small},
  middle/.style={draw=accent,rounded corners=2mm,fill=white,text width=0.48\textwidth,minimum height=16mm,align=center,font=\small\bfseries},
  arr/.style={-{Latex[length=2mm]},draw=muted,line width=0.55pt}
]
\node[block] (n1) {\textbf{First Native Domain}\\Experience, subject, observer};
\node[block,right=of n1] (n2) {\textbf{Second Native Domain}\\Graphs, groupoids, germs, levels};
\node[middle,below=of $(n1)!0.5!(n2)$] (n3) {Structural Translation\\\normalfont Preserve relations; assert no ontological identity};
\node[block,below left=of n3] (n4) {\textbf{Formal Theory}\\witness, filtration, residue};
\node[block,below right=of n3] (n5) {\textbf{Cases and Empirical Bridges}\\observable consequences, counterexamples, exit conditions};
\draw[arr] (n1) -- (n3);
\draw[arr] (n2) -- (n3);
\draw[arr] (n3) -- (n4);
\draw[arr] (n3) -- (n5);
\end{tikzpicture}%
}
\end{center}

\noindent\textbf{Conceptual route.} Introduction $\longrightarrow$ Part I $\longrightarrow$ Part III $\longrightarrow$ Part VI. For readers chiefly concerned with Third-Order Consciousness, Krishnamurti, and methodology.

\smallskip
\noindent\textbf{Mathematical route.} Part II $\longrightarrow$ Part IV $\longrightarrow$ the corresponding chapters of Part III. For readers who prefer to verify definitions and theorems first.

\smallskip
\noindent\textbf{Empirical route.} The working examples in Part I $\longrightarrow$ the Structural Translation Cards in Part III $\longrightarrow$ the cases and experimental designs in Part V.

\begin{methodbox}[Terminological discipline]
Technical mathematical terms are permitted to perform translation work in Part III only after they have received native-domain definitions in Part II. An ordinary-language term appearing in a book title, route preview, or English-language source in Part I does not constitute a mathematical definition. In the main text, $A\strsim M$ denotes only a proposed structural compatibility; the equals sign is used only within a single mathematical native domain.
\end{methodbox}
\endgroup

\chapter*{Introduction}
\addcontentsline{toc}{chapter}{Introduction}
\markboth{Introduction}{Introduction}
\addtocontents{toc}{\protect\setcounter{tocdepth}{1}}

\begin{readerbox}[Begin with the most natural question: Why Bring Third-Order Consciousness in Front of Mathematics?]
If this is your first encounter with the book, the phrase ``\AFBD'' in the title will probably prompt an immediate question. Third-Order Consciousness is a problem about experience, the subject, and the observer. Why place it beside an unfamiliar mathematical structure? Does this merely replace psychology with a more elaborate vocabulary? Can mathematical precision genuinely yield new understanding, or will it only make a problem we can already feel and describe harder to read?

Those are exactly the questions the book is trying to answer. Mathematics is introduced neither because it is ``higher'' than psychology nor to prove that consciousness is, ontologically, some graph, groupoid, or dynamical system. It is needed because Third-Order Consciousness raises a set of highly specific questions that can readily be run together in ordinary language. When are two similar experiences merely alike, and when do they constitute ``again''? Where does an ongoing stretch of experience begin, and where does it end? Does the ``I'' already exist as the master of experience before experience occurs, or does it gradually acquire a position as experience is organized, compared, and remembered? Once a pattern has been recognized, will a newly enlarged background overturn it completely, or merely change the way we see it?

These questions sound entirely ordinary. The difficulty is that several different structures are often wrapped in the same word in ordinary language. We mix together ``repetition,'' ``cycle,'' ``remembering,'' ``it happened again,'' and ``it is always like this.'' We mix up ``I am looking from here now'' with ``I am this fixed kind of person.'' We mix up ``this story has explanatory force at present'' with ``this story is now the truth.'' Mathematics is most valuable here not because it answers life on our behalf, but because it forces apart questions that have been allowed to adhere to one another in ordinary language.
\end{readerbox}

\section{The Research Question: How Experience Becomes ``Mine,'' and How the ``I'' Enters Observation}

Let us start with a scenario that contains no mathematical symbols at all. Imagine that you sent someone a message but received no reply for a long time. At first, perhaps you would check the phone over and over again, and your attention would keep returning to the message. After a while, your body might tighten a little. Then language would intervene: ``Why isn't he replying?'' Later, earlier experiences might be brought back into the picture: ``Why is it like this again?'' Finally, a heavier judgment might be triggered: ``It seems to me that I always encounter people who are disrespectful like this.''

If we read only the final sentence, it sounds as if a fully formed ``I'' had been in charge all along: I sent the message; I waited; I became angry; I remembered; I made the judgment. If we slow the sequence down, however, several layers come into view. Your body and attention may change first. Only later do you give the event a name, connect the present experience with the past, and place both within the same pattern. That pattern may then be compressed further into a story about ``me'' and ``other people.'' Third-Order Consciousness is concerned precisely with these acts of organization.

We therefore begin not with ``What kind of entity is consciousness?'' but with a narrower and more operational question:

\begin{quote}
\textbf{How does an experience still in progress gradually become organized as ``my experience''? How do experiences at different times come to be related as ``this is happening again''? And does the observer responsible for organizing, comparing, and narrating them truly stand outside the process as a whole?}
\end{quote}

\emph{Third-Order Consciousness} organizes this problem in terms of three modes of consciousness: direct presence, subject-forming organization, and the observer itself entering observation \citep{toc2026}. ``Third order'' does not divide people into three ranks. It is closer to three focal settings. At times, life unfolds chiefly through direct coupling between body and environment. At times, language, memory, identity, and the gaze of others acquire organizational authority. And at certain moments we begin to see that the ``I'' which has been interpreting, judging, planning, and managing experience is itself shaped by the past, by language, and by relationships. Part I develops this line of thought in full.

\setcounter{section}{1}

\section{Why Formalize Structure: How Can Mathematics Give Something Back to Psychology?}

The preceding section made the aim of the book explicit: to attempt a controlled formalization, through \AFBD{}, of several organizational relations in Third-Order Consciousness. A psychological reader will therefore ask the obvious next question: \textbf{even if those relations can be expressed clearly in mathematics, why is that worth doing? What, exactly, can psychology get back from the exercise?} Psychology already studies memory, repetition, selfhood, event segmentation, and metacognition. If structural formalization merely replaces familiar psychological phenomena with a harsher symbolic vocabulary, there is no reason to undertake it.

We therefore test the psychological value of the structural bridge directly. More precisely, we ask: \textbf{once some of the organizational relations described by Third-Order Consciousness have received a clear and limited representation in the \AFBD{} native domain, can they return to the psychological scene and help us formulate questions that were previously difficult to state in a stable way?} Only if this question also has an answer does structural formalization complete the task assigned to it in this book.

\subsection{First clarify a term that will recur throughout the book: what, exactly, is an episode?}

Before continuing, we need to clarify an apparently ordinary word that can easily create confusion: \emph{episode}. Psychology, cognitive science, memory research, and narrative studies all use the word, but not with exactly the same meaning. This book does not treat it as a standard unit whose definition has already been fixed across all branches of psychology. It uses \emph{episode} as a \textbf{working term}. In the bilingual edition, the Chinese phrase ``经验片段（episode）'' was normally given at first occurrence and then shortened to \emph{episode}; the English edition can use \emph{episode} directly, while retaining the same working definition.

Why not simply say ``event''? Because \emph{event} easily suggests something already existing independently in the external world, with clear boundaries: receiving an email, having an argument, attending a meeting. Psychological research often confronts something less tidy. Many things may happen continuously in the external world, while which stretch of that flow a person organizes as ``this one thing'' can depend on current goals, attention, prediction, emotion, action, and narrative position. Conversely, a single external event may be organized into several different experiential units at different scales.

We therefore use the following plain working definition:

\begin{quote}
\textbf{At a stated scale of observation, an episode is a relatively coherent experience--action process. It can be treated as ``this segment'' because transitions before and after it alter the current mode of organization enough to distinguish this process from its neighbors.}
\end{quote}

Three phrases need to be unpacked. First, \textbf{``experience--action process''} means that an episode may include not only what happened externally, but also bodily sensation, shifts of attention, emotional change, thought, expectation, and preparation for action. Second, \textbf{``relatively coherent''} does not require internal homogeneity; it requires only that, for the present research question, the segment can still be treated as a process whose parts bear some internal relation to one another. Third, \textbf{``a stated scale of observation''} means that an episode has no universally natural duration. A few seconds can constitute an episode, and so can an hour. What determines the unit is not clock time alone, but the kind of organizational change that the researcher and participant are tracking.

Episode segmentation is therefore \textbf{scale-dependent}, but it is not arbitrary. A research-worthy segmentation should be able to point to some observable basis: a change in action goal, a substantial prediction update, a context switch, a redirection of attention, a transition in emotional organization, a change in the coordinates of self-evaluation, or a point the participant reliably experiences as ``one thing ending and another beginning.'' Classical work in Event Segmentation Theory begins from closely related questions, examining how people organize ongoing activity into relatively discrete event units and relating those boundaries to prediction updating, comprehension, and memory organization \citep{zacks2007,kurby2008}. This book uses such work as an empirical point of reference but does not identify its working notion of episode with the EST notion of event.

Two further distinctions are essential. \textbf{An episode is not a narrative.} An episode says, in effect, ``which stretch of experience is being treated as a relative unit.'' Narrative goes further: ``what did this stretch mean? how does it relate to my past? what does it show about the kind of person I am?'' The same set of episodes can be organized into different stories. Conversely, a powerful narrative can force together materials that did not originally belong to the same episode. This distinction will matter repeatedly, because the book asks how experience is first segmented and then reorganized by higher-order stories about the self.

\textbf{Nor does episode mean a diagnostic episode in the clinical sense.} Clinical psychology and psychiatry use terms such as \emph{depressive episode} and \emph{manic episode}, each with its own diagnostic criteria. The usage here is broader and more neutral. An episode is simply an experiential segment that can be identified at the current scale of analysis; no pathology is implied.

Operationally, episodes need not be drawn by the researcher ``by feel.'' For overt activity, participants can watch a video and press a key whenever ``one natural unit ends and another begins''; segmentation agreement across participants can be compared; behavioral logs, linguistic transitions, and prediction changes can also provide independent indicators. For more internal experience, researchers can use time-stamped experience reports, retrospective timeline reconstruction, or coded interviews, while explicitly checking whether later narrative has manufactured boundaries that were not present in the original process. In other words, \textbf{an episode is itself a segmentation hypothesis that can be studied, compared, and overturned.}

This bears directly on the later mathematical translation. Part III will compare some episodes with mathematical birooted traverses, but only a very limited relational skeleton will be carried across: at the current scale there is an entry, an internal process, and an exit. We will not presuppose that psychological episodes have unique natural boundaries, and still less that every experience can be neatly cut into traverses. On the contrary, if empirical work shows that some class of psychological processes cannot sustain relatively stable segmentation at all, or that different scales display no comparable nesting relation, then the corresponding mathematical mapping should fail at that point.

With this working definition in place, later questions such as ``is a larger episode composed of smaller episodes?'' and ``when are two episodes recognized as the same pattern coming back?'' can be read as questions about the segmentation and organization of experience, rather than as uses of an unexplained English label.

The first form of feedback from formalization is therefore a gain in \textbf{distinction}. In everyday psychological language, ``repetition'' can mean a stimulus recurring, a behavioral habit, rumination, situational similarity, an emotion being triggered again, or a person's very explicit sense that ``I am back in that same pattern.'' These phenomena are certainly related, but they are not identical. The mathematics developed later will force us to distinguish ``a local structure appearing again'' from ``the system organizing an occurrence as a return,'' and to distinguish ``a return has reached completion'' from ``after that completion, later continuation can still differ.'' Even if no new mechanism theory emerges, these distinctions already change the grain at which psychological questions are asked.

A second form of feedback concerns \textbf{scale}. Human experience is never organized at only one timescale. Ten seconds of argument can belong to a dinner; one dinner can belong to a relationship; a relationship can be written into the life story ``I always meet people like this.'' Psychology already studies event segmentation, autobiographical memory, and narrative identity. This book adds a further question: do the segmentations at different scales stand in compatible organizational relations? Is a larger episode genuinely composed of smaller episodes? When we compress many small events into one large story, which information has been compressed legitimately, and which differences that remain relevant to the future have been erased at the same time? The mathematical language of levels, embeddings, decomposition, and refinement turns such questions from rhetoric into structural conditions that can fail.

A third form of feedback concerns the \textbf{position of the observer}. Krishnamurti asks us to see that the observer does not stand outside what is observed. If that challenge is taken seriously, a psychological model cannot indefinitely assume a transparent, stable, omniscient ``observer inside the research participant.'' We have to ask: how is the current self-referential position formed? How does it assign some events to ``my history''? When the observing position changes, do segmentation, comparison, and prediction over the same material change with it? The value of formalization is to make ``the observer is also inside the system'' incur model-theoretic consequences rather than remain a beautiful philosophical sentence.

A fourth form of feedback concerns \textbf{falsifiability}. A good interdisciplinary model should not merely produce more explanations; it should produce exit conditions. If experience has no relatively stable episode segmentation, if segmentations at different scales bear no compatible relation, if ``again'' is only a label added retrospectively by the researcher, or if every new fact completely destroys what had counted as a return-completion, then the corresponding mathematical mapping must be weakened or removed. This matters especially in psychology: the value of formalization lies not only in making a theory more exact, but also in making it easier for reality to reject the theory.

\subsection{What can psychology actually take back after formalization?}

The language of ``higher resolution'' remains too abstract unless it can be made operational. More concretely, the book seeks to separate phenomena that are easily compressed into a single sentence of self-narration and turn them into variables that can be recorded, compared, and experimentally manipulated independently. Consider an extremely common sentence:

\begin{quote}
``Why am I doing this again?''
\end{quote}

In ordinary conversation the sentence is already perfectly natural. In research, however, it conceals at least six distinct questions. What episode has just occurred? Where did it begin, and where did it end? Is the present event merely similar in content to a past one, or has the participant organized it as ``the same pattern returning''? At what scale does ``the same'' hold---ten seconds, a day, a relationship, a whole life? If more background is added, does the original return-completion judgment remain? And even if both occurrences are described as ``again,'' are subsequent expectations, actions, and relational trajectories the same?

The function of formalization is to separate these six questions. It does not require the researcher literally to draw a participant as a groupoid. It supplies, instead, a kind of \textbf{coding checklist}: the researcher must know whether the current variable is local similarity, the boundary of an experiential segment, a participant's judgment that ``the same pattern has returned,'' the chosen scale, progressively enlarged background, or a difference in later continuation that survives even after a pattern has been recognized.

\begingroup
\small
\setlength{\tabcolsep}{3.5pt}
\begin{longtable}{@{}>{\bfseries\raggedright\arraybackslash}p{0.205\textwidth}>{\raggedright\arraybackslash}p{0.295\textwidth}>{\raggedright\arraybackslash}p{0.43\textwidth}@{}}
\toprule
Psychological question & What formalization forces us to distinguish & Candidate operationalization \\
\midrule
\endfirsthead
\toprule
Psychological question & What formalization forces us to distinguish & Candidate operationalization \\
\midrule
\endhead
``What counts as one thing?'' & The ongoing experiential flow versus segmentation into experiential units & Segmentation tasks on video or behavioral streams, time-stamped experience records, independent coding of interview material; compare with established event-segmentation paradigms \citep{zacks2007,kurby2008}.\\
``Why does it feel like `here we go again'?'' & Local reappearance versus recognition that ``the same pattern has returned'' & Control content similarity and frequency; separately record ``these look similar'' and ``I think the same pattern has returned''; then compare memory chunking and prediction of the next event.\\
``From what position am I looking now?'' & Event content versus the current self-referential/evaluative coordinates & Record current evaluative goals, ownership/agency, self-relevance, and changes in narrative subject; compare with existing metacognitive and decentering variables rather than presupposing a wholly new psychological capacity \citep{fleming2012,bernstein2015}.\\
``After more background is added, is it still the same cycle?'' & The basis for return-completion in finite background versus progressively enlarged background & Reveal material in stages and observe whether participants retain the original ``same pattern'' judgment, refine or narrow it, or withdraw it as a prior misclassification.\\
``Why can two cases both be called `again' and yet go in different directions afterward?'' & A return-completion judgment versus differences in continuation after return-completion & Match the earlier experiential segments as closely as possible, then measure differences in next-step prediction, available actions, relational expectation, or memory association.\\
``Which processes remain hard to compress into a familiar pattern?'' & Sequences that readily yield recognizable return-completion versus sequences that persistently resist it & Treat this as a candidate research problem and compare memory chunking, recall, and predictive load across sequences. The relevant mathematical terms are not introduced until Section 0.3 and Part II, and must not be translated directly into clinical labels.\\
\bottomrule
\end{longtable}
\endgroup

The most important point in this table is not that it already provides a mature psychometric instrument. It is that \textbf{after formalization, the researcher knows what to measure next, what to control, and which phenomena must not be collapsed into one another.} Part V develops these ideas into controlled experimental designs. Here the Introduction needs only to show that mathematics has not taken psychology away from the scene. It should return us to the same scene with finer observational tools.

Three direct comparisons can already be envisioned. First, while matching content similarity and frequency as closely as possible, compare segments that are ``merely repeated'' with segments that participants identify as ``another instance of the same pattern.'' Ask whether the latter are more readily compressed into a single memory chunk and whether they alter prediction of what comes next. Second, have participants make a ``same pattern'' judgment over a finite experiential segment and then progressively add background, observing whether the judgment is preserved, refined, or must be withdrawn entirely. The target here is \textbf{whether a finite judgment can be retained as the background expands}; the corresponding mathematical terminology appears only later. Third, select two experiential segments that participants both recognize as ``the same pattern,'' keep their initial structure as similar as possible, and compare whether their later predictions and action sets nevertheless differ systematically. The target here is \textbf{the possibility of distinct later continuations after an apparent return-completion}; only later will the book introduce formal language such as \emph{residue} for this difference.

Thus ``giving something back to psychology'' has at least three levels in this book: first obtain finer conceptual distinctions; then turn those distinctions into coding and experimental variables; finally let the formal structure generate empirically differential predictions. Only at the third level would we be entitled to discuss stronger mechanism links. This book mainly completes the first two and proposes a program for the third. It does not write experiments not yet conducted as established psychological facts.

This is also the minimum demand of the book on mathematics: \textbf{mathematics must make psychology draw distinctions it could not previously draw so easily, generate new questions that can be observed, compared, intervened on, or falsified, and then return to experience to be corrected there.} If a mathematical term merely renames ``anxiety,'' ``cycle,'' or ``self'' without changing any recordable variable, experimental comparison, or failure condition, it does not belong in the core theory of the book.

\section{Why the Framework Ultimately Becomes AF-by-Discrete: Ask First About Germs, Then About ``by-Discrete''}

Model selection cannot be reverse-engineered from the answer. If one first notices \AFBD{} and then redescribes each of its properties as something psychology ``must'' require, the comparison becomes a tautology. We therefore separate the selection process into three layers: the psychological native domain first states its relational requirements independently; this project then adds a further, explicitly fallible finiteness commitment; only at the end does the mathematical native domain enter the comparison.

\subsection{What the psychological native domain asks for first: six requirements stated without groupoid vocabulary}

The following six requirements are stated only in the language of Part I and ordinary psychology. They are not paraphrases of the definition of \AFBD{}.

\begin{center}
\small
\begin{tabularx}{0.95\textwidth}{>{\bfseries}p{0.08\textwidth}Y}
\toprule
No. & Psychological/phenomenological requirement \\
\midrule
$P_1$ & Local experience may occur first; the model must not preload a fully completed subject with global explanatory authority as the starting point of every process.\\
$P_2$ & Reappearance or similarity of content must remain distinguishable from the subject's judgment that ``the same pattern has returned.''\\
$P_3$ & The entry, finite process, and exit of an episode must enter the analysis; segmentation cannot remain forever a transparent presupposition of the researcher.\\
$P_4$ & The same return-completion judgment must not predetermine the same later course; subsequent prediction, action, and memory recruitment need their own bookkeeping.\\
$P_5$ & The conditions that generate the position from which experience is named, compared, evaluated, and governed must themselves be able to enter observation.\\
$P_6^{*}$ & A position does not gain unconditional authority over other positions merely by being ``more reflective''; multiple local positions do not automatically authorize a new supreme observer.\\
\bottomrule
\end{tabularx}
\end{center}

Requirements $P_1$--$P_4$ come directly from distinctions repeatedly developed in Part I. $P_5$ is definitional for Third-Order Consciousness. $P_6^{*}$ has explicit support in the original work and in the texts on ``field consciousness,'' although the chronology of its formation should remain open to a provenance audit. We avoid phrases such as ``there is no canonical super-root'' on purpose, because they would already import vocabulary from the answer-side of the comparison.

\subsection{What this project adds: two finiteness conditions that must not be given the same name}

Third-Order Consciousness itself does not imply that there are only finitely many organizing fissures. This project adds a stronger and explicitly fallible psychological modeling hypothesis:
\[
\boxed{
\begin{gathered}
H_{\mathrm{fin}}^{\mathrm{psych}}:\\[-1mm]
\text{under a stated task, scale, and re-identification rule,}\\[-1mm]
\text{persistent singular-like loci stabilize as a finite set}\\[-1mm]
\text{as the material expands.}
\end{gathered}}
\]
The sets $B_m$ in Part V place empirical pressure only on this hypothesis.

The condition in the mathematical native domain is different. For a fixed finite generating family $S$ and a fixed tile presentation, write
\[
\partial_S^\infty
:=\bigsqcup_{T\in\mathcal T_\infty}\partial_S T.
\]
The source-side finite-support condition used in this book is
\[
\boxed{
H_{\mathrm{fin}}^{\mathrm{AFBD}}(S,\mathsf B):
\quad |\partial_S^\infty|<\infty.}
\]
It counts the totality of regular and singular boundary points. The present psychological operationalization concerns only a singular-like subset. Therefore
\[
\boxed{
H_{\mathrm{fin}}^{\mathrm{psych}}
\neq
H_{\mathrm{fin}}^{\mathrm{AFBD}}.}
\]
Support for the former does not prove the full mathematical condition; failure of the former initially strikes only the singular-germ correspondence used in this book.

\subsection{The strongest competitor is really a general germ framework}

Ordinary graphs, hidden Markov models, and general state-space models can represent states and transitions, and can be enlarged to retain history. But their basic structure does not, by itself, automatically distinguish ``arriving at the same point'' from ``arriving there through different local actions,'' nor does it automatically provide germ equivalence, local isotropy, or graph-of-germs resolution. To carry $P_4$--$P_6^{*}$, such structures need additional machinery.

The strongest comparator is therefore not an obviously weaker HMM, but \textbf{a general groupoid of germs without a finite-support assumption}. A general groupoid of germs already provides a formal language in which designated positions remain internal to the local arrow structure, while retaining local action, path-sensitive sameness, and isotropy. For $P_1$--$P_6^{*}$ alone, that may already be enough.

The specific contribution of \AFBD{} begins only in the next step. It adds to general germ geometry an open AF core and a discrete non-AF complement. In the compactly generated fixed-generator tile presentation adopted here, this yields a globally finite persistent generator-boundary support. The source-side parent picture is therefore not ``a little singularity everywhere,'' but
\[
\boxed{
\begin{gathered}
\text{large AF-like background}\\[-1mm]
+
\text{globally finite persistent boundary support}\\[-1mm]
+
\text{label-preserving connector grammar}.
\end{gathered}}
\]
Within that finite boundary support, regular reversible gluing must still be separated from singular germ resolution. Part III makes cross-domain use only of the latter; the full boundary inventory remains in the mathematical native domain.

\begin{methodbox}[Compare sufficiency, not uniqueness]
This book does not prove that \AFBD{} is the unique possible formal model of Third-Order Consciousness. The more exact claim is this: general germ geometry is a strong candidate for reflexive organization; if one additionally accepts the finite-fissure commitment expressed by $H_{\mathrm{fin}}^{\mathrm{psych}}$, then \AFBD{} becomes a particularly economical, sharply constrained, and falsifiable source model. Alternative models may certainly exist, but they must state openly how they carry $P_1$--$P_6^{*}$, the finiteness commitment, and the corresponding exit conditions.
\end{methodbox}

\subsection{A necessary digression: why ``by-discrete,'' rather than immediately allowing polynomial activity?}

A more mathematical question remains, but it directly affects the psychological translation. If the book is willing to allow multiple boundary points, why insist on ``by-discrete''? Why not begin with a larger class in which boundary complexity can grow with scale, for example the classes suggested by polynomial-activity automata?

The historical prototype comes from rooted-tree automata. Let $g$ be a finite-state tree automorphism and write
\[
\theta_g(n)=\bigl|\{v\in X^n:g|_v\neq 1\}\bigr|,
\]
the number of vertices at level $n$ that still carry a nontrivial section. In Sidki's activity hierarchy, \emph{bounded automata} are precisely the degree-$0$ case: $\theta_g(n)$ is uniformly bounded in $n$. More generally, polynomial activity of degree $d$ allows $\theta_g(n)=O(n^d)$. In the Moore diagram, the distinction has a clean combinatorial meaning: the bounded case does not allow nontrivial cycles to form hierarchical chains through directed paths, whereas positive polynomial degree permits such chains. Active prefixes can then continue proliferating with level. A trajectory may, for example, wind around one cycle an arbitrary number of times, follow a connecting path into another cycle, and so on. Persistent active directions therefore need no longer collapse to a finite list. This is exactly where the finite-singular-direction picture begins to fail \citep{sidki2000,nek2022}.

The \AFBD{} framework in this book generalizes that geometry of bounded automata while replacing the regular rooted tree by a more general Bratteli/tile/groupoid setting. What must be preserved is not one particular automaton presentation, but the \textbf{global finite-boundary-support picture} behind it: relative to a fixed finite generating set, the total number of boundary points on \emph{all infinite tiles} that persistently carry non-AF continuation remains finite. Hence only finitely many infinite tiles carry boundary; the remaining infinite tiles have no persistent boundary interface. If one further assumes the bounded-type tile-inflation form used in the author's work, then finite-tile boundary cardinalities are uniformly bounded, infinite-tile boundary support is finite when infinite tiles are present, and $|V_n|,|E_n|$ are uniformly bounded. The level-total boundary burden across tile types is therefore uniformly bounded as a derived consequence, not as the general groupoid definition itself. From the rooted-tree prototype, bounded type is therefore closest to the quantitative idea of ``uniformly bounded activity,'' while \AFBD{} preserves the more basic qualitative skeleton: \textbf{the persistent boundary support of the entire limit system can still be listed finitely.}

Why does the present project stop here? The most important reason is not technical convenience, but the fact that \textbf{from the source-side finite boundary support the book extracts only a narrower finite singular-interface hypothesis for psychology}. The mathematical native domain still permits regular boundary gluing; the psychological translation asks only whether a large experiential system may be organized mainly by extensive relatively quiet local background together with a small number of repeatedly active singular-boundary-like interfaces that carry segmentation, shifts of position, and reorganization of continuation. If boundary inventory were allowed from the outset to grow with resolution at rate $n$, $n^2$, or a higher polynomial, then we would be testing a different organizational picture: the exceptional part would itself begin to form a hierarchy unfolding with scale, rather than a finite family of discrete fissures embedded in a large background.

This does not mean that polynomial-activity systems are mathematically ``too wild.'' On the contrary, they remain highly controlled and worthwhile classes of automata. Nor is the selection criterion amenability, growth, or some abstract measure of ``complexity.'' What matters to this book is the \textbf{geometry of boundary organization}: is fixed-generator non-AF information still carried by a globally finite persistent boundary support, or do active/boundary directions continue to proliferate with scale? One further distinction must be kept explicit. The finite support can contain both regular and singular points; the cross-domain translation into Third-Order Consciousness uses only the singular subset with a nontrivial local group of germs. The former source-side situation is the present by-discrete picture. The latter, proliferating situation would require a new formal language that records the boundary-growth hierarchy itself.

From the psychological side, this restriction has another important advantage: it leaves the theory with a genuine failure condition. If a researcher can say ``add one more boundary here'' every time a new episode appears, and can then allow the number of boundary loci to expand at some polynomial rate as more material arrives, the claim that ``a small number of fissures organizes a large background'' becomes almost impossible for experience to disconfirm. The model makes the stronger commitment first on purpose: \textbf{once the rules of analysis are fixed, genuinely persistent and future-relevant distinct interface loci should repeatedly point back to the same finite set.} ``Finitely many interface types'' is not strong enough, because infinitely many distinct loci could still share one type. Only when this finite-set commitment fails under the data do we have reason to upgrade the model.

A layered discipline from Part V must also be inherited here. In the mathematical native domain, unbounded growth of the boundary cardinality of individual finite AF fibres / tiles directly violates the general bounded-type estimate. Growth of a level-total aggregate by itself identifies a failure of the tile-inflation package only after the multiplicity of tile types has also been placed on the ledger; it does not tell us which clause failed. As long as the total persistent boundary set in the infinite-tile limit remains globally finite, \AFBD{} may still hold. The psychological side, by contrast, observes only candidate singular-like loci, so it cannot identify its own resolution counts with those mathematical boundary counts. If, as material and scale increase, new distinct candidate singular loci continue to appear, and no finite set can stably carry return concentration, sector passages, and continuation changes, then what exits first is the correspondence between psychological interfaces and a finite singular-boundary subset. Regular boundaries have not been operationalized, so this is not a direct empirical refutation of the complete global boundary support required by \AFBD{}.

\begin{principlebox}[Minimal singular-interface principle: complexity must be forced out by the system]
We use a conservative order of cross-domain modeling on purpose. The mathematical native domain first retains the full finite boundary support, with regular boundary gluing and singular germ resolution kept in separate accounts. The psychological side activates only a finite singular-boundary-like interface set. New singular degrees of freedom are added only when repeated observation forces the singular-interface complexity beyond the present range.

Hence
\[
\boxed{
\begin{gathered}
\text{removing the global privilege of a single observing position}\\[-1mm]
\not\Rightarrow\\[-1mm]
\text{allowing the number of boundaries to proliferate without bound}
\end{gathered}}
\]
Third-Order Consciousness can require us to abandon ``one root represents the whole.'' Multiple boundaries can likewise require us to abandon ``one interface explains every fissure.'' Both moves remain compatible with finite-interface organization. Only when the system itself repeatedly forces out a boundary hierarchy that cannot be finitized should the current by-discrete model give way.
\end{principlebox}

This is why ``by-discrete'' is not an accidental technical label in the book. On the source side the first constraint is that \textbf{regular and singular boundary points together still form a globally finite persistent boundary support}. The singular subset is therefore automatically finite, and we obtain a sufficiently narrow germ-sensitive substructure for the cross-domain discussion of ``the observer re-entering the structure,'' multi-sector passage, and nontrivial residue. Part V tests whether this singular-germ correspondence survives experience. It does not use psychological data to decide the entire \AFBD{} boundary condition directly.

\begin{readerbox}[A second coordinate not yet activated: boundaries can be counted, but they can also be pulled apart]
The preceding discussion chiefly asks a counting question: for a fixed generating family and tile presentation, how many persistent boundary points occur on all infinite tiles? The psychological side separately tracks how many organizing loci can be re-identified across materials. Tile inflation offers another, independent geometric coordinate: \textbf{once boundary blocks have been distinguished, how do they separate from one another as scale increases?}

Kuang defines \emph{expanding tile inflation} in arXiv:2409.19491 \citep[Section~3]{kuang2025tiling}. On the source side, boundary/connecting points on a finite tile whose mutual distance remains uniformly bounded are first grouped into the same block. The definition then requires that, after bounded-step telescoping, distinct boundary blocks in sufficiently high-level tiles come from distinct embeddings of tiles one level below. This condition is not about how many boundaries there are. It is about whether distinct blocks are genuinely pulled apart by large-scale geometry under inflation. Under the additional boundedness assumptions of that paper, the maximum distance between boundary points on finite tiles, and the tile diameters, both grow on an exponential scale; those estimates are then used to control the orbital geometry of infinite tiles.

There are therefore at least two independent coordinates:
\[
\boxed{
\begin{aligned}
&\text{boundary cardinality: how many persistent boundary}\\[-1mm]
&\qquad\text{points/loci are there globally?}\\[1mm]
&\text{boundary separation: how do those points/loci/blocks}\\[-1mm]
&\qquad\text{separate as scale grows?}
\end{aligned}}
\]
In the fixed-generator setup of this book, \AFBD{} constrains global persistent boundary support; general groupoid bounded type constrains each finite AF fibre / tile uniformly, while the tile-inflation form used here also controls the number of tile types per level and hence the aggregate burden. Both belong to the cardinality/control side and are not equivalent to expanding separation geometry. Expanding geometry reminds us that \emph{a globally finite boundary set may still have nontrivial large-scale relative position}. If empirical analysis later groups several loci into ``interface families,'' that is only a coarser type-level coding; the number of families cannot replace the mathematical boundary-point count.

The psychological side presently supports only a weaker question. As more history and background are recovered, do several organizing interfaces initially compressed into ``the same core issue'' gradually reveal distinct patterns of memory recruitment, return language, and continuation? This motivates an ``interface separation'' question, but does not yet provide a natural psychological metric and certainly does not justify predicting an exponential law. We therefore \textbf{do not include expanding in the structural correspondences of Part III, nor in the core empirical hypotheses of Part V}. 

One final source-side overreach must also be avoided. The definition of expanding and the distance results that follow it concern how boundary blocks separate on large scales under tile inflation. The book does not rewrite this into the additional claim that ``different types of boundary must eventually lie in different orbits.'' Only if future work discovers an independent observable of interface distance and a stable pattern of cross-scale separation would expanding tile inflation become a candidate for the next layer of formalization.
\end{readerbox}

\subsection{Why traverse, return, repetitivity, and incompressibility naturally appear together}

The attraction of \AFBD{} for this book is not exhausted by the image of ``background and boundary.'' More importantly, in the tile-inflation/bounded-type framework used here, a cluster of mathematical notions that we happen to need to distinguish can all be \textbf{defined internally within the same system}, rather than being named first in psychology and then projected backward into mathematics. Part II defines them one by one. Here we preview, in ordinary language, why they belong to one picture.

\paragraph{Repetitivity: a local structure can appear again}
A finite tile or finite pattern can reappear inside larger tiles, at larger scales, and in different positions; under suitable assumptions, the gaps between such reappearances are controlled. This tells us only that ``another copy of a local structure has appeared.'' It does not yet say that ``the system has recognized itself as having come back to the same place.'' That distinction will become the first dividing line of the book between repetition and ``again'' \citep{kuang2026a}.

\paragraph{Traverse: a finite passage from boundary to boundary}
A trajectory can enter a finite tile through one boundary point, move through the tile for a while, and leave through another boundary point, without meeting another boundary point of the same level in between. Such a finite passage is called a \emph{traverse}. It cuts a long stretch of internal motion into minimal boundary-relevant pieces with an entry, an internal process, and an exit.

\paragraph{Finite return witness: a special completion of a traverse}
If such a boundary-to-boundary passage ultimately returns to the same relevant boundary occurrence, we obtain what the book later tracks separately as a \emph{finite return witness}. The crucial point is that \textbf{return does not mean merely ``looks like something seen before.''} The internal boundary organization of the system has to say what counts as the starting point, what counts as the endpoint, and in what sense the relevant position has been reached again. Most of the trajectory may unfold in the AF-like internal region, but if the boundary organization is removed altogether, only local motion and repeated copies remain; a bare AF background does not by itself define ``coming back'' \citep{kuang2026a}.

The \emph{source return word} in the original growth paper retains one additional layer of generator context: boundary letters at the beginning and end enclose that returning internal excursion. Part II will separate these two levels rigorously. For now, the point is simply this: \emph{a traverse whose designated boundary occurrences coincide is a finite witness of return-completion; a source return word is the word-level object that retains the full boundary-entry/exit labels}. Here ``coincide'' is used to avoid any ambiguity with topological closedness; the later mathematical chapter gives the standard combinatorial meaning of \emph{closed walk} separately.

\paragraph{Incompressibility: no usable return factor occurs inside the trajectory}
In the mathematical native domain, a word is called \emph{incompressible} if it contains no return subword. This does not mean that it is ``stubborn,'' ``painful,'' or ``unable to move on.'' It says something very precise: within the entire return language currently under consideration, no contiguous factor realizes an already defined return-completion. We will later see that global incompressibility must be taken relative to all relevant returns, not merely to one kind of germ or one local perspective.

These four terms are therefore not independent labels. They form a linked chain of questions:
\[
\boxed{
\begin{gathered}
\text{Can a local structure appear again?}\\[-1mm]
\downarrow\\[-1mm]
\text{How do boundaries cut a trajectory into traverses?}\\[-1mm]
\downarrow\\[-1mm]
\text{Which traverses supply finite witnesses of return-completion?}\\[-1mm]
\downarrow\\[-1mm]
\text{Which source-word blocks belong to the return language?}\\[-1mm]
\downarrow\\[-1mm]
\text{Which words still avoid all returns?}
\end{gathered}}
\]

This is why the book treats \AFBD{} as a \textbf{structural framework for traverse--return--repetitivity--incompressibility}, not merely as an abstract class of groupoids. More strictly: \AFBD{} supplies the overall architecture of ``AF background + globally finitely many persistent continuation interfaces across all infinite tiles for a fixed generating family.'' A tile presentation further endows those interfaces with a same-label connector grammar, and general groupoid bounded type adds uniform per-fibre / per-tile boundary control; the tile-inflation form used here also adds finite boundary support on infinite tiles when present and uniform $|V_n|,|E_n|$, from which aggregate levelwise control follows. Traverse, return, repetitivity, and incompressibility obtain their exact definitions and interactions in this more concrete tile language. This qualification matters because the book must not smuggle a statement stronger than the actual theorem into the terminology itself.

\subsection{Why this happens to fit Third-Order Consciousness so well}

Only now do we reach the part of the project that is genuinely striking. Third-Order Consciousness did not predict \AFBD{} in its own conceptual native domain, and it never used the language of tiles, germs, or groupoids. Yet when the two sides are developed independently, a structural coincidence appears that deserves serious testing.

\emph{Third-Order Consciousness} first asks us not to preload a completed subject before every experience. Much experience can unfold as a direct flow of body, attention, sensation, and action. Only through certain organizing operations is experience further marked as ``mine,'' segmented into episodes, linked to the past, and described as ``this is happening again.'' One step further, the third order asks the position responsible for marking, comparing, and explaining to enter observation itself, rather than remaining permanently hidden outside the structure \citep{toc2026}.

The \AFBD{} side presents, in its own language, precisely a formal picture of \textbf{large background + globally finite boundary support}. Extensive local structure can unfold and recur inside the AF core. Finitely many boundary points determine how these local pieces connect across tiles, while singular germ data further determine which forms of return-completion retain differences of later continuation at finer resolution.

The real fit between the two sides therefore lies not in shared nouns, but in a family of relational questions:
\begin{itemize}
\item Can extensive local activity occur without presupposing a global explanatory center?
\item Can re-identifiable local organizing positions have disproportionate influence over how a large background is connected?
\item Can an episode be understood as a finite process between two specially marked positions?
\item Is ``again'' stronger than general repetitivity because it requires an organized relation of return-completion?
\item Can a system preserve an existing finite return-completion while increasing resolution and revealing distinct later continuations behind it?
\item Can the observer be understood as a position generated within these organizing relations rather than as the master outside the system?
\end{itemize}

These are the real reasons the book chooses \AFBD{}. The framework gives us a chance to separate ``flow, subject, again, completion, observer''---which are easily left as poetic expressions in Third-Order Consciousness---into several structural operations that are not equivalent to one another. More importantly, it supplies a very concrete research hypothesis: \textbf{the complexity of subject-formation in a system need not be distributed uniformly across all experience; it may be concentrated chiefly at a small number of organizing transition loci and in the relations of later continuation that they carry.} This is not an established psychological fact. It is the central hypothesis that the book hopes to turn into a testable question through structural translation.

\section{What the Book Inherits, and What It Adds: Intellectual Sources and Research Contributions}

By this point the reader knows what the book is trying to do: use the mathematical language of \AFBD{} to give a controlled formalization of a family of relations in Third-Order Consciousness concerning experience, subject, ``again,'' observer, and multiscale organization. The reader also knows why this might be useful to psychology, and why \AFBD{} was selected in particular. One question must now be made explicit: \textbf{where do these ideas come from, and at what point does the book itself begin making research claims of its own?}

This is not a ceremonial attribution paragraph. In a book that crosses psychology, philosophy, and mathematics, provenance is itself part of the theoretical architecture. If the ideas of the original author, the mathematical tools introduced here, and the cross-disciplinary mappings still awaiting tests are written as though they were all the same kind of sentence, the reader soon loses the ability to judge what sort of evidential responsibility any claim carries. A sentence may be summarizing \emph{Third-Order Consciousness}; it may be stating a mathematical fact; or it may be proposing a psychological hypothesis of the book itself. Those three sentence-types owe very different debts of evidence.

\subsection{First layer: the conceptual framework comes from Guoxin Peng\textquotesingle s \emph{Third-Order Consciousness}}

The basic problem-structure of First-, Second-, and Third-Order Consciousness comes chiefly from Guoxin Peng's \emph{Third-Order Consciousness: From the Biological Self and Social Self to Vacant Consciousness (The Geometry of Existence Book 4)}, edited by Pinchuan Xie \citep{toc2026}. Its most important lines of thought include: First-Order Consciousness as direct presence of life; Second-Order Consciousness as the formation of the subject through language, memory, identity, the other, and narrative; and the third order not as the construction of a still higher subject, but as the observer that has been explaining and managing experience itself entering observation. The original work repeatedly stresses that the three orders are not three kinds of people and not a spiritual ladder from low to high, but different modes of organization within one field of consciousness.

The task of Part I is therefore first of all \textbf{review, clarification, and critical reconstruction}. It reorganizes the central arguments of \emph{Third-Order Consciousness} into a form suitable for academic discussion and adds philosophical and psychological contexts from existentialism, Krishnamurti, event segmentation, narrative identity, metacognition, and related traditions. Those additions do not change the intellectual source. In particular, the basic ideas of ``Third-Order Consciousness,'' the loss of final authority by the observer, and the role of ``the observer is the observed'' within the third-order problem are not presented as original proposals by Zheng Kuang.

\subsection{Second layer: the Geometry of Existence and the translation protocol supply the discipline of interdisciplinary work}

The method by which the two native domains are crossed also has an explicit source. The ``Translation Protocol for a Geometry of Existence'' and ``Between Structural Blockage and Structural Overreach,'' proposed by \emph{Snowfall on the Clock}, require cross-disciplinary translation to preserve definitions within their native domains first, then compare relations rather than names, and to disclose translation strength, evidential level, failure conditions, and untranslated remainder. The Geometry of Existence is particularly alert to \emph{structural overreach}: the process by which a locally useful explanation gradually acquires the position of an ``ultimate explanation'' \citep{translation2026,geometry2026}.

This discipline matters especially here. Once \AFBD{} enters psychological language, mathematical precision can create a dangerous illusion: because the mathematical definition is exact, the psychological correspondence may seem to become exact automatically. The reverse is true. Mathematics guarantees rigor only inside the mathematical object. The bridge from the mathematical object to a psychological phenomenon must still be constructed and tested separately, and must remain removable. Thus the recurring disciplines of ``native-domain fidelity,'' ``structural mapping,'' ``stop line,'' and ``untranslated remainder'' are not ad hoc disclaimers invented for this manuscript. They implement a methodological framework that the project explicitly inherits.

\subsection{Third layer: \AFBD{} is an independent mathematical native domain}

\AFBD{} groupoids, AF groupoids, Bratteli diagrams, germs, tile inflation, bounded type, traverses, returns, and incompressibility belong to another, independent mathematical tradition. Part II builds from scratch the parts required by this book. Some specific structures and results come from the literature on groupoids, Cantor dynamics, and bounded type, and from Zheng Kuang's previous mathematical work on tile inflation, return, incompressible elements, and near-full groups \citep{kuang2026a,kuang2026b}.

This must be stated separately from the preceding layer. \textbf{\emph{Third-Order Consciousness} does not propose that consciousness is \AFBD{}, and it does not propose using traverse, return, or incompressibility to explain psychological processes.} Conversely, the existing mathematical literature on \AFBD{} did not develop those objects as a theory of Third-Order Consciousness. The two lines of thought were originally independent. That independence is exactly what makes the convergence identified in Section 0.3---``a small boundary support organizing a large background'' and the sequence repetitivity--traverse--return--incompressibility---a research problem rather than a translation of the same concepts into two vocabularies.

\begin{methodbox}[Logical autonomy across the book: shared words must not merge the native domains in advance]
Words such as ``boundary,'' ``singularity,'' ``topology,'' and ``field'' in Part I belong first to the conceptual native domain of Third-Order Consciousness and the Geometry of Existence. They do not presuppose the strict definitions of \emph{boundary point}, \emph{singular germ}, topology, or groupoid in Part II. The mathematical objects and theorems of Parts II and IV, in turn, must remain valid without any psychological interpretation. Only in Part III does the book state, relation by relation, what is allowed to cross domains. Empirical results in Part V can support, weaken, or withdraw those correspondences; they cannot retroactively rewrite either native domain.

We therefore allow the same word to have a legitimate use in both native domains, but never allow shared spelling to create automatic identification. Every genuine connection made later must retain its provenance, the relation preserved, its evidential level, and its stop line.
\end{methodbox}

\subsection{Fourth layer: the central research contribution of the book is to connect the two lines while making the bridge testable}

The original research of this book begins here: \textbf{it separates general germ geometry from the additional finite-fissure commitment and proposes that \AFBD{} can serve as a candidate formal framework for some organizational relations of Third-Order Consciousness; it then examines, item by item, what such formalization preserves, distinguishes, predicts, and where it must fail.}

That work consists of several successive steps. First, one has to select from the large mathematical language of \AFBD{} only the structures genuinely relevant to the organizational problems of Third-Order Consciousness, rather than importing any term that merely sounds similar. Second, the psychological side must stand independently first: what counts as an episode, what ``again'' means, what the current observing position is. Only after that may those notions be compared with traverse, return, root, germ resolution, and other mathematical objects. Third, resemblance must be turned into auditable relational claims: is what is preserved the structure entry--internal process--exit, or return-completion, refinement, persistence, or a difference of later continuation? Finally, formalization must be sent back to psychology and converted into questions that can be coded, compared, experimentally manipulated, or falsified.

The originality of the book therefore lies neither in giving psychological phenomena new mathematical names nor in claiming to have discovered a hidden ontology of consciousness. It is closer to \textbf{interface design}: retain the experiential and philosophical problems of Third-Order Consciousness on one side; retain the mathematical definitions of \AFBD{} on the other; then search for a set of intermediate relations precise enough to be useful and modest enough not to claim too much.

\begin{readerbox}[While reading, keep four kinds of sentence separate]
\textbf{First: conceptual claims from the original work on Third-Order Consciousness.} Examples include how the subject forms and why the observer is not the master outside the structure. Such claims should first be understood through \emph{Third-Order Consciousness} and the relevant intellectual history.

\textbf{Second: definitions or theorems in the mathematical native domain.} Examples include the definitions of AF groupoid, traverse, and return, or the persistence of a particular return under inflation. Their responsibility is mathematical definition and proof.

\textbf{Third: structural translations proposed by this book.} Examples include ``an experiential episode is structurally compatible with a birooted traverse with respect to entry--process--exit,'' or ``a small number of organizing transition loci may organize a large psychological background.'' These are the central research hypotheses of the book. They are constrained by the translation protocol developed later. Before that protocol is formally presented, the reader should understand them only as \emph{controlled structural correspondences}, not as psychological mechanisms already supported by evidence.

\textbf{Fourth: empirical predictions for future psychology.} Examples include whether candidate organizing transition loci really predict experiential segmentation, memory reorganization, or ``same pattern'' judgments disproportionately. Such questions require independent experiments. Mathematical proof cannot answer them in place of data.
\end{readerbox}

Once these four layers are kept separate, the division of responsibility across the whole book becomes clear: \textbf{\emph{Third-Order Consciousness} supplies the conceptual structure to be understood; the Geometry of Existence and the translation protocol specify how one may cross domains without overreach; \AFBD{} supplies an independent formal language; this book builds the bridge and hands it back to psychology for testing.} Every later chapter should preserve that division of labor.

\section{Three Intellectual Coordinates: Subject, Observer, and Structural Overreach}

The Introduction does not need to retell the entire history of ideas before the main text begins. Only three coordinates are indispensable to what follows: \emph{why the subject matters; why the observer also needs to be explained; and why every explanation must know where it stands}. Existentialism, Krishnamurti, and the Geometry of Existence place these three questions, respectively, in the foreground. Part I develops the first two lines; Part III turns the third into translation discipline. Here we provide only the problem map a psychological reader needs.

\subsection{Existentialism: when external essence recedes, choice and responsibility return to the subject}

One of existentialism's most important historical roles was to refuse the idea that a pre-given essence could decide an entire life on a person's behalf. The absence of an external script does not imply that ``nothing matters.'' On the contrary, choice, commitment, and responsibility become heavier because they cannot be outsourced \citep{crowell2012,flynn2006,sartre2007}.

This matters to the book because the second-order subject is not an error that ought to be abolished. Language, memory, identity, planning, ethics, and relationship all require some relatively stable organizing position. The capacity to say ``this is my commitment'' or ``this is the decision for which I am responsible'' is itself indispensable to complex human life.

Yet once the subject takes on organizational work, we must ask a new question: \emph{how was the ``I'' that is currently choosing, evaluating, and taking responsibility formed?} What past, rules, comparison classes, and social positions does it use? If those conditions remain permanently invisible, the subject can easily mistake its historical conditioning for an absolute command from outside the structure.

We therefore begin neither by abolishing the subject nor by enthroning it. We begin from a narrower proposition:
\[
\boxed{\begin{gathered}
\text{the subject must be able to work,}\\[-1mm]
\text{but need not possess final explanatory authority}
\end{gathered}}
\]

\subsection{Krishnamurti: the observer has a past and must itself enter observation}

Krishnamurti pushes the question a step further. He repeatedly argues that the psychological observer is composed of knowledge, memory, hurt, comparison, ideal, and tradition; the observer is therefore not a pane of transparent glass \citep{krishnamurti1954,krishnamurti1969,kft2026a}. When a voice says ``I will deal with my anger,'' ``I should be more mature,'' or ``I must let this go quickly,'' its calmer tone does not by itself prove that it stands outside the psychological process.

The easiest mistake is to hear ``the observer is the observed'' as a mystical doctrine of oneness. This book preserves only a narrower, auditable structural question: \textbf{the position responsible for naming, evaluating, and governing experience may itself be generated by the same past, rules, and expectations.} It cannot acquire exemption from scrutiny merely because it performs the function of observation.

This pressure test directly determines the reading of Third-Order Consciousness developed in this book. The third order is not a higher observer added on top of the old observer. It makes \emph{the conditions that generate the observing position itself} visible.

\subsection{The Geometry of Existence: an explanation must have a position, and must retain a path of revision}

The most common problem in interdisciplinary work is not a lack of explanation, but the ease with which explanation acquires global power. A judgment useful in one local setting can gradually decide in advance how every piece of evidence is to be read, and why every counterexample ``really still supports it.'' That is what the book calls \emph{structural overreach} \citep{translation2026,geometry2026}.

The Geometry of Existence contributes to the book not another psychological taxonomy, but a working discipline. It requires us repeatedly to do three things:
\begin{enumerate}
\item \textbf{Grounding:} every abstraction must be able to return to concrete events, timescales, and observable material.
\item \textbf{Positioning:} every judgment must state from what position, using which variables, and at what scale it holds.
\item \textbf{Restoring movement:} if new evidence causes the old structure to lose discriminating power, the model must allow itself to be revised, split, or withdrawn rather than protected by a larger story.
\end{enumerate}

This is also why the book repeatedly says that ``judgments must have positions.'' Making the position explicit does not entail relativism. On the contrary, only once we know where a judgment is made can we identify what material supports it, what could change it, and what it has no authority to explain.

\section{The Minimum Discipline of Interdisciplinary Translation: Move Relations First, Not Names}

The most dangerous move before mathematics enters psychology is to write an equality sign too quickly:
\[
\text{anxiety}=\text{boundary},\qquad
\text{trauma}=\text{incompressibility},\qquad
\text{self}=\text{root}.
\]
Such sentences look ``formal'' but merely replace old psychological words with mathematical ones. The strict relational content that the mathematical objects originally possessed has in fact been discarded.

The translation protocol of this book therefore permits only a slower order \citep{translation2026}:
\[
\boxed{
\begin{gathered}
\text{state the psychological problem clearly}\\[-1mm]
\downarrow\\[-1mm]
\text{verify the mathematical native meaning}\\[-1mm]
\downarrow\\[-1mm]
\text{compare only the declared relation}\\[-1mm]
\downarrow\\[-1mm]
\text{return to empirical testing}
\end{gathered}}
\]

\subsection{What, exactly, is moved by structural translation?}

What is moved first is a relation. Psychology may already have a question: \emph{why is an ongoing stretch of experience segmented as ``this one event''?} Mathematics happens to provide a birooted object that retains a finite structure of entry--internal process--exit. What deserves comparison is that relation. It does not follow that a psychological episode ``is'' a graph.

Likewise, the comparison between ``again'' and return is useful only if return preserves a testable condition of return-completion that ordinary similarity does not. If every similar experience can simply be renamed a return, the mathematical word has added no psychological resolution.

\subsection{Five evidential levels answer one question: how much proof does this sentence owe?}

The book uses Levels I--V below. A psychological reader need not memorize them on first reading. Their purpose is simply to separate claims of different strength:

\begin{center}
\small
\begin{tabularx}{0.94\textwidth}{>{\bfseries}p{0.16\textwidth}Y}
\toprule
Level I & Heuristic analogy: useful for asking questions, but no strict correspondence has yet been fixed.\\
Level II & Controlled structural mapping: both native domains have been independently defined, and the preserved relation, failure conditions, and untranslated remainder are explicit.\\
Level III & Limited formal model: actual psychological variables have been explicitly encoded and data can test the relevant formal relation.\\
Level IV & Mechanism link: independent evidence is required that the formal variables stand in a stable mediating relation to cognitive or behavioral mechanisms.\\
Level V & Ontological claim: reality is claimed to ``really be'' the mathematical structure. This book makes no such claim.\\
\bottomrule
\end{tabularx}
\end{center}

The book as a whole currently operates mainly at Level II. A theorem in Part IV can of course be fully rigorous within the mathematical native domain, but \textbf{a mathematical proof cannot pay the evidential debt of psychology}. Only if a future experimental module specifies state variables, a segmentation rule, a return test, and falsifiable predictions might some local correspondence move into Level III.

\subsection{Concept passports: what record must a mathematical term leave when it enters psychology?}

Every important correspondence in Part III receives a simplified ``Concept Passport.'' The psychological reader need not memorize the table; it is enough to check six questions:
\begin{enumerate}
\item Why does psychology need this distinction?
\item What exactly is the mathematical object in its own native domain?
\item Which relation alone is preserved across the two sides?
\item What content is explicitly not translated?
\item What observable or operationalization can test the correspondence?
\item What result would force the correspondence to exit?
\end{enumerate}

If the last question has no answer, the correspondence has already begun to acquire an authority it should not possess.

\section{Structural Compatibility: the Book Claims Comparability, Not Identity}

The notation
\[
A\sim_{\mathrm{str}}M
\]
has a limited meaning here on purpose: \textbf{on a declared set of relations, a psychological object or process $A$ and a mathematical object $M$ can be placed in a controlled comparison.}

It is not a mathematical equivalence relation. Unless explicitly stated otherwise, we assume neither reflexivity, symmetry, nor transitivity; in particular, from
\[
A\sim_{\mathrm{str}}M,
\qquad
M\sim_{\mathrm{str}}N
\]
one may not simply infer $A\sim_{\mathrm{str}}N$. Every bridge must state separately what it preserves.

\subsection{A minimal example: why event segmentation is an appropriate place to speak of structural compatibility}

An ongoing activity stream does not by itself announce where ``one event'' begins and ends. Psychology can study event boundaries through \emph{prediction change, goal transition, memory chunking}, or by asking participants to mark segmentation points directly. A mathematical birooted traverse, by contrast, strictly retains two specially designated endpoints and the finite process between them.

The materials on the two sides are entirely different, yet one narrow relation can be compared: \emph{a process becomes a comparable finite segment because its entry and exit are organizationally distinguished}. If future work shows that a certain class of psychological processes has no relatively stable endpoints, and that segmentations at different scales are wholly incompatible, then this mapping should fail. It is precisely because failure is possible that structural compatibility carries more research content than a literary analogy.

\section{Four Governing Principles That Will Recur after the Introduction}

To prevent every new term later in the book from sounding like the start of a new theory, we collect here the four most basic principles. The later notions of root, return, pNH, residue, and the case protocols are different developments of these principles.

\begin{enumerate}
\item \textbf{Principle of globally finite boundary organization.} Relative to the fixed finite generating set used in this book, the \AFBD{} picture is this: large internal structures may keep growing, yet the persistent boundary points that survive to infinite scale and govern cross-block continuation, taken over \emph{all} infinite tiles, form a finite set. Connectors at finite levels cannot be wired arbitrarily either; they must respect the original labels of the generators. Under the general groupoid bounded-type condition, the generator-boundary burden of each individual finite AF fibre / tile is uniformly bounded. In the bounded-type tile-inflation form used in the author's work, finite boundary support on infinite tiles when present and uniform bounds on $|V_n|$ and $|E_n|$ are imposed as well; consequently the aggregate boundary burden across tile types at a level is uniformly controlled.

\item \textbf{Principle of fissures forced by the system.} Candidate psychological boundaries should not be drawn in advance by the researcher. They should become visible through the repeated operation of the system itself: at certain positions the old local rule repeatedly ceases to suffice, more background has to be recruited, endpoints of many ``again'' judgments repeatedly concentrate there, and the mode of subsequent understanding or action changes in a stable way after the position is crossed. Return can help locate such a boundary-like interface, but it cannot be used to define the boundary retroactively.

\item \textbf{Principle of multiple returns.} Multiple boundary germs or organizing sectors naturally generate multiple return relations; returns along one trajectory can be nested, overlapping, or interlaced. The converse does not hold: even with only one germ type, a single return grammar may produce nested returns at different scales.

\item \textbf{Principle that an observing position has no final authority.} The current root position, the way an episode is segmented, what is recognized as a return, and how much difference the present resolution retains can all organize experience. Yet none of them gains final explanatory authority merely by virtue of performing that organizational work. The core of Third-Order Consciousness is to make these conditions of organization themselves available for scrutiny.
\end{enumerate}

\begin{plainbox}[The four principles in one sentence]
\textbf{A very large system can be organized by globally finite boundary support; the singular-like loci used in the psychological translation must be forced into visibility by the operation of the system itself; these loci can jointly generate multiple return grammars; and no current organizing position thereby acquires final authority.}
\end{plainbox}

\section{What This Book Does, and What It Explicitly Does Not Do}

By this point the Introduction has answered several distinct questions: why some organizational relations in Third-Order Consciousness should be structurally formalized, why such formalization might give something back to psychology, why \AFBD{} is the mathematical framework selected, why the Translation Protocol and the Geometry of Existence must constrain every crossing between domains, and how little ``structural compatibility'' itself commits us to. We can now gather these answers into a more direct \textbf{reader contract}: after reading the book, what do we hope to know more clearly, and which questions does the book leave open on purpose rather than pretend to have already solved?

\subsection{First, a necessary clarification: the book is still at the stage of meta-terminology}

A mature theory in psychology usually invites further expectations: stable variables, scales or task paradigms, an explicit mechanism chain, predictions comparable with existing theories, and enough data to establish which effects are actually present. This book has not yet reached that stage. It is closer to a stage of \textbf{meta-terminology and structural modeling}.

By ``meta-terminology'' we do not mean that the project is merely inventing words. The point is almost the opposite: before rushing to propose a mechanism, we first make explicit \emph{which phenomena must be distinguished, what relations may hold among them, which formulations are still asking the same question, and which formulations have already changed the question}. Many psychological disputes are difficult not only because data are lacking, but because several different processes have been packed into the same word. ``Repetition'' may mean content similarity, increased frequency, behavioral habit, rumination, or the subject's explicit experience that ``the same pattern has returned.'' ``An event'' can look like a natural unit even though where it begins and ends is itself an event-segmentation problem. ``I am observing myself'' can mix present experience, evaluative rules, self-image, and past memory.

The first task of the book, therefore, is not to announce a new psychological mechanism. It is to build a language fine-grained enough that these distinctions no longer collapse so easily into one word. Only then do we ask whether \AFBD{} can formally represent some of the relevant relations; only after that do we ask whether those formal distinctions can become operational variables and empirical predictions.

\subsection{The three kinds of output this book actually commits to}

The first is \textbf{conceptual resolution}. The book aims to separate several phenomena that everyday language easily blends together: local similarity and structural return; experiential flow and episode segmentation; current self-reference and fixed identity; finite return-completion and subsequent continuation; an already available structural witness and a retrospective narrative judgment. The value lies not in the terminology itself, but in whether these distinctions let us ask more precise questions of the same psychological material.

The second is \textbf{formal grammar}. \AFBD{} is not used to attach mathematical labels to psychological phenomena, but to express a family of relations: how a large and locally repetitive background can be organized by a small collection of interfaces; how trajectories are segmented into traverses; under what conditions a traverse gives a return; how structures across scales embed and decompose; why a finite return-completion can remain valid while finer resolution still reveals different continuations; and what it means, relative to the entire return language, for no compressible return-completion to be available. The rigorous mathematical propositions in Part IV hold only in this mathematical native domain.

The third is a \textbf{research program that can return to psychology}. If a formal distinction cannot be turned into ``what can be recorded, what should be controlled, what can be compared, and what result would refute it,'' then it has not genuinely given anything back to psychology. Part V is therefore not an appendix of examples attached to the mathematics; it is a necessary half of the book. It turns distinctions concerning episode segmentation, ``again'' judgments, contextual refinement, differences in subsequent continuation, and observer position back into candidate behavioral variables, coding schemes, and experimental questions.

The order of these three layers cannot be reversed:
\[
\boxed{
\begin{gathered}
\text{first distinguish the problem}\\[-1mm]
\downarrow\\[-1mm]
\text{then build the formal relation}\\[-1mm]
\downarrow\\[-1mm]
\text{finally return to empirical testing}
\end{gathered}}
\]

If the last step remains impossible for a long time, the preceding mathematics may still constitute an interesting structural model, but it cannot be promoted to a psychological mechanism theory.

\subsection{A psychological reader need not become a groupoid expert in order to judge whether the book is useful}

The book contains genuine mathematical definitions and proofs, but a psychological reader does not have to master groupoid theory before judging the value of the structural translation. Four ordinary questions remain sufficient throughout: Does the mathematical distinction separate two psychological phenomena that were previously mixed together? Does it reveal a new question about scale, boundaries, or differences in future continuation? Can it be translated into an observable, codable, or manipulable variable? Does it state clearly what result would make the correspondence fail?

If the answer to all four questions remains negative, mathematical elegance by itself does not amount to success in psychology. Conversely, even if only part of the \AFBD{} structure ultimately survives empirical testing, the bridge has real value if it helps psychology formulate distinctions that were previously difficult to state stably.

\subsection{What the book explicitly does not do: six boundaries}

First, \textbf{this book is not a complete theory of consciousness.} It formalizes some \emph{organizational relations} in Third-Order Consciousness. It does not explain why subjective experience exists, why qualia have a particular character, or solve the hard problem of consciousness. Phenomenological mineness, bodily feeling, affective valence, and conscious appearing itself remain important remainders.

Second, \textbf{this book is not a theory of neural mechanism.} We have no evidence that the brain contains a mathematically literal AF core, boundary germ, or traverse. Future work may certainly ask whether event boundaries, prediction updates, or distributed neural states stand in a stable relation to the formal structure. At present, however, any statement of the form ``this mathematical node is this brain region'' is structural overreach.

Third, \textbf{this book is not a system of clinical diagnosis or treatment.} Return is not relapse; incompressibility is not trauma, rumination, personality rigidity, or ``failure to heal''; and boundary is not a psychological trauma point. Clinical concepts have their own definitions, measures, ethical settings, and intervention contexts. Mathematical terms cannot replace them. A structural model can help formulate questions, but it has no authority to diagnose a person by bypassing clinical evidence.

Fourth, \textbf{this book does not turn Third-Order Consciousness into a spiritual ranking.} First-, Second-, and Third-Order Consciousness are not three kinds of people; the third order is not a higher identity. If the formalization is ultimately read as a scale for judging who is ``more third-order'' and who remains at a ``lower order,'' it has already departed from the original direction of Third-Order Consciousness, which denies the subject final authority.

Fifth, \textbf{this book does not use mathematics to prove philosophy or intellectual texts.} Krishnamurti's statement that ``the observer is the observed'' does not become an empirical scientific fact because a groupoid model is elegant; Sartre's freedom and responsibility are not proved by a formal structure either. Intellectual texts can pose structural questions, mathematics can attempt to represent some of their relations, and psychological experiments can test some empirical consequences. These remain three separate evidential accounts.

Sixth, \textbf{this book does not convert structural compatibility into ontological identity.} Even if every correspondence in Part III works extremely well, the conclusion would only be that \AFBD{} is an effective formal representation of some organizational relations in Third-Order Consciousness. It would not follow that ``consciousness is, in essence, a groupoid,'' still less that ``the universe itself is this structure.''

\subsection{What, then, would count as success?}

At the minimum level, success means that readers acquire a language more fine-grained than ``the self is a process'' or ``people repeat old patterns,'' and can unpack such large statements into finite questions. At an intermediate level, some \AFBD{} structures should organize psychological material in a stable way: local repetition should be distinguishable from the recognition that ``the same pattern has returned''; a small number of organizational transition positions may exert disproportionate influence over a large background; multi-scale episode segmentations may display stable compatibility; or different subsequent continuations may retain independent predictive value beneath the same return judgment.

A stronger success requires future work: the structural variables must not only be measurable, but generate differentiated predictions beyond existing theories; manipulating candidate organizing variables should systematically alter return judgments, memory chunking, subsequent prediction, or self-reference; and these effects should not be fully explained by existing theories of event segmentation, metacognition, decentering, memory, or emotion regulation. Only then would discussion of a Level IV mechanism link become appropriate.

Failure is equally informative. If experiments show that psychological episodes do not possess the cross-scale compatibility required here, that ``again'' is wholly determined by semantic similarity, that the proposed organizing transition positions play no special role in subsequent organization, or that graph-of-germs-style resolution only adds terminology without adding future-relevant distinctions, then the corresponding mappings should be removed. A theory that can lose part of its structure is closer to research than one that can explain everything.

This gives the rest of the book a clear standard of evaluation. The relevant question is not how complicated the later mathematics becomes, but whether at every step we still know which problem it came from, what relation it preserves, what it can return to experience, and where it must stop.

\section{How to Read This Book: Begin with Psychological Phenomena, Pass through Formalization, and Return to Psychology}

The Introduction has now completed three tasks. We know why the book attempts a controlled formalization of certain organizational relations in Third-Order Consciousness through \AFBD{}; we know that formalization is valuable only if it can return something to psychology; and we know that interdisciplinary translation must respect native-domain fidelity, separate evidential accounts, and explicit failure conditions. Before entering the body of the book, one practical question remains: \textbf{why is the book divided into six parts, and why must they be read in the present order?}

This is not a matter of layout preference. The sequence reproduces the research path of the project itself: first clarify the psychological and intellectual structure to be formalized; then build the mathematical language independently within its own domain; only after both sides can stand on their own do we build the bridge; once the bridge exists, we examine what it entails mathematically, whether it can be used in psychology, and where it must stop. The simplest route through the book is therefore
\[
\boxed{
\begin{gathered}
\text{first see the psychological phenomenon}\\
\Downarrow\\[-1mm]
\text{build an independent \AFBD{} formal language}\\
\Downarrow\\[-1mm]
\text{translate organizational relations one by one}\\
\Downarrow\\[-1mm]
\text{organize formal properties and derivable consequences}\\
\Downarrow\\[-1mm]
\text{return to psychological cases, coding, and experiments}\\
\Downarrow\\[-1mm]
\text{state clearly what the model cannot explain}
\end{gathered}}
\]

Skipping any stage produces a characteristic error. Skip Part I and the mathematics no longer knows what it is formalizing. Skip Part II and mathematical terms collapse into psychological metaphors. Skip Part III and there is no auditable bridge between the native domains. Skip Part V and the formalization can no longer genuinely give anything back to psychology. Skip Part VI and the theory becomes progressively more likely to mistake a local structure for a complete theory of consciousness.

\subsection{Part I: understand Third-Order Consciousness first as a psychological and intellectual problem}

Part I answers: \textbf{what, exactly, are we preparing to formalize?}

There is as yet no AF core, boundary, germ, traverse, or return. The reader first encounters more familiar experiential questions: how feeling and action can occur before explicit self-narration; how experience becomes organized by language, memory, role, and narrative into ``my experience''; what the word ``again'' adds in ``why am I doing this again?''; why the observer may not stand outside experience; and what Krishnamurti's statements that ``the observer is the past'' and ``the observer is the observed'' are meant to remove in terms of structural privilege.

For a psychological reader, the task of this part is not to accept a new taxonomy of consciousness, but to acquire a \emph{phenomenological problem map}: which phenomena belong primarily to immediate occurrence, which to subject-centered organization, and which to reflexive observation of that organization itself. If this layer is unclear, later mathematics may be precise while precisely formalizing an ambiguous target.

On a first reading of Part I, then, the most important question is not to memorize the names of the ``three orders,'' but repeatedly to ask: \emph{which two psychological operations that were previously mixed together has this passage separated?}

\subsection{Part II: leave psychology temporarily, and learn the mathematical language as mathematics}

Part II may be the least familiar section for psychological readers, but it is indispensable. It answers: \textbf{what does the \AFBD{} language mean when psychology is not being discussed at all?}

The book requires the mathematical native domain to stand independently first, and does so on purpose. Graphs, paths, roots, Bratteli diagrams, groupoids, germs, the AF core, tiles with boundary, traverses, returns, repetitivity, and incompressibility must all be definable, exemplifiable, and mathematically usable without psychological assistance. Only then, when Part III says that a psychological relation is structurally compatible with return, does ``return'' remain a mathematical concept rather than a metaphor invented for psychology.

A psychological reader, however, does not need to become a groupoid specialist on the first pass. What must be retained are several \textbf{structural through-lines}:
\begin{enumerate}
\item large portions of internal structure can form a relatively stable and repetitive background that is approximable through finite scales;
\item globally finite boundary support can exert disproportionate control over how these large backgrounds continue and are glued together, with the singular-germ subset carrying the nontrivial local resolution;
\item a traverse is a finite process bracketed by two distinguished positions;
\item repetitivity only says that a local form appears again, whereas return requires a stronger return-completion structure;
\item the way an already available finite witness is retained when scale increases is a question of persistence;
\item structures such as the graph of germs and residue record why different subsequent continuations may remain even after an apparent return;
\item incompressibility means that no return factor is available in the relevant language; it is not a label for psychological pathology.
\end{enumerate}

These relations form the mathematical skeleton actually used in Part III. Technical details concerning étale topology, bisections, inverse semigroups, and related constructions can be approached on a first reading with the more modest goal of understanding why they provide legitimate foundations for the later definitions, and revisited when Parts III or IV require exact reference. In other words, a psychological reader can first learn \emph{which distinctions the mathematical objects make}, and only gradually fill in \emph{how those objects are constructed rigorously}.

\subsection{Part III: this is the actual center of the book - the bridge must be built piece by piece}

Part III answers: \textbf{on which relations, exactly, are Third-Order Consciousness and \AFBD{} structurally compatible?}

This part does not place psychological words and mathematical words in two columns and join them with equality signs. Every correspondence must pass through the same checks again: What is the concept in the psychological native domain? What is the object in the mathematical native domain? Which single relation is preserved across the two? What is explicitly not translated? Which example supports the correspondence? Which non-example looks similar yet fails the conditions? Which empirical result would force this correspondence to exit?

Part III is therefore where the original formalization of the book actually occurs. The research value of comparing an episode with a birooted traverse, for example, is not that both words can refer to ``a stretch of something,'' but that we may be able to preserve the relation entry--internal process--exit and the structure of cross-scale decomposition while refusing to smuggle across emotional quality, neural mechanism, or psychological meaning. The comparison between ``again'' and return is meaningful only if return-completion adds a testable distinction beyond ordinary similarity.

While reading this part, it is useful to ask the same four questions of every correspondence:
\begin{quote}
\textbf{What did it bring from psychology? What did mathematics actually add? What was not carried across? What result would break the bridge?}
\end{quote}
As long as these questions retain clear answers, interdisciplinary translation remains auditable.

\subsection{Part IV: formal theory is not a second departure from psychology, but a check that the bridge has internal discipline}

Part IV concentrates on formal properties such as finite witnesses, decomposition, persistence, filtration, incompressibility, and residue. It answers: \textbf{once Part III has selected these structural correspondences, what mathematical consistency conditions must we accept?}

A theorem here is first and foremost a mathematical theorem. Its role is not to prove a psychological fact, but to prevent the structural translation from changing its rules whenever convenient. If, for example, a genuine return is mathematically defined through an already available finite witness of return-completion, then one cannot later enlarge the scale and declare that witness nonexistent merely to accommodate a psychological case. If actual psychological judgments revise themselves in such a nonmonotone way, the correct conclusion is that the psychological correspondence must be weakened, not that the mathematical definition should be quietly altered.

On a first reading of Part IV, a psychological reader can prioritize the ordinary-language explanations, the psychological structural notes, and the failure notes around each definition or proposition, returning to proof details as needed. What matters most is to see that \textbf{formalization has begun to impose constraints on interdisciplinary intuition}. Mathematics is no longer merely a supplier of attractive analogies; it now tells us which consequences accompany a definition once we insist on using it.

\subsection{Part V: mathematics must return to psychology, or only half the bridge has been built}

Part V answers the most practical question in the book: \textbf{can these formal distinctions become something psychology can observe, code, compare, and intervene on?}

Here the concrete materials discussed earlier reappear: nonresponse and anger in relationships, repeated checking, rumination, shifts of attentional position in research reports, and narratives that compress different experiences into ``I am always like this.'' Unlike Part I, however, we no longer merely describe experience. We ask more finely: where is the episode segmented; is what we see only local repetition, or has the participant actually made a return judgment; does that judgment survive when more background is restored; can the same return be followed by different subsequent continuations; do a small number of organizational transition positions really carry more organizing work than ordinary internal moments?

This is where the book genuinely reconnects with psychology. If a formal distinction cannot be converted into a candidate observable, a control variable, a differential prediction, or a failure condition, its interdisciplinary value must be reassessed.

\subsection{Part VI: where a theory stops is itself part of the theory}

Part VI treats phenomenological mineness, qualia, embodiment, affect, neural realization, social constitution, heterogeneous multi-singularity integration, and future formal problems. It answers: \textbf{what has this formalization not explained?}

This part is not a modest disclaimer added at the end. It has the same theoretical standing as the preceding five parts. A structural model that cannot identify its own remainder easily mistakes ``I have not represented this variable'' for ``this variable does not exist'' or ``some other structure has already explained it.'' Consciousness research makes this danger especially acute. Formally representing how a subject position is generated does not explain why experience has a first-person sense of mineness; representing episode and return does not explain affective quality, bodily realization, or neural mechanism.

Part VI therefore protects not the dignity of the theory, but reality's right to exceed it.

\subsection{On a first reading of the mathematical parts, what must be understood and what may be postponed?}

To keep the book genuinely accessible to psychological readers, we state the minimum explicitly. On a first reading of Parts II--IV, one should at least understand the ordinary-language meaning of the following claims:
\begin{itemize}
\item why ``local repetition'' is weaker than ``return'';
\item why an experience must have some meaningful entry and exit in order to count as an analyzable episode;
\item why a small number of organizational transition positions can organize a large background without thereby being identified directly with neural nodes or psychological trauma;
\item why increasing scale sometimes embeds an old structure into a larger background rather than overturning it;
\item why ``it has returned'' and ``what happens next is still different'' can both be true;
\item why incompressibility is a property in the language/factor sense, not a claim that ``this person cannot move on.''
\end{itemize}

If these relations are clear, the reader already has the mathematical intuition needed for Parts III and V. More technical proofs ensure that the author has not used these intuitions arbitrarily, but they should not become the price of entry for a psychological reader.

\subsection{Book-wide stop rules: when must we stop rather than add more explanation?}

Finally, we gather the stop lines scattered across the Introduction into one set of book-wide rules. If any later chapter violates one of them, the correct response is to retreat, not to patch the gap with more terminology.

\enlargethispage{2\baselineskip}
\begin{enumerate}
\item \textbf{Native-domain stop.} If a psychological concept has not yet been made clear within its psychological or intellectual native domain, mathematical vocabulary may not be used to give it false precision. Likewise, if a mathematical object has not been independently defined, psychological meaning may not be used to make the object stand.

\item \textbf{Scale stop.} Moving from momentary experience to personality, from individual to society, or from behavior to neural mechanism always requires an explicit mediator. Using the same vocabulary does not mean that the scale has remained the same.

\item \textbf{Evidence stop.} Mathematical proof proves only mathematical propositions; first-person description supports only claims about phenomenological structure; experimental data support only the empirical propositions that the design can test. One kind of evidence cannot pay the debt of another.

\item \textbf{Clinical stop.} Mathematical terms such as return, incompressibility, singularity, and boundary do not automatically possess clinical meanings. Unless independent clinical operationalization and validation have been established, they may not be used as diagnostic, pathological, or therapeutic labels.

\item \textbf{Observer stop.} Any purported model of the third order that ultimately still requires a highest observer who stands outside the structure, is naturally transparent, and never itself requires explanation has not solved the problem set by this book.

\item \textbf{Exit stop.} If a mathematical correspondence adds no psychological distinction, observable, prediction, or failure condition, or if empirical material persistently violates its core relation, that correspondence should exit. A theory does not fail because it loses part of its structure. Refusing to let any structure exit is what truly deprives a theory of its research character.
\end{enumerate}

\addtocontents{toc}{\protect\setcounter{tocdepth}{0}}

\part{The Phenomenological Structure of Third-Order Consciousness: Experience, Subject, and Observer}

\begin{center}
\large\itshape From direct occurrence, through subjective organization, to the observer itself entering observation
\end{center}

\begin{methodbox}[Note on Part I]
This Part remains entirely within the phenomenological, psychological, philosophical, and social-theoretical native domain of Third-Order Consciousness. It first asks how experience occurs, how the subject is formed, and how the observer loses the privilege of standing outside the structure. No later mathematical technical term is used here to explain psychological phenomena. The mathematics is defined independently only in Part II; structural translation begins only in Part III.
\end{methodbox}

\begin{sourcebox}[Intellectual source of Part I and the work of this manuscript]
The basic theoretical framework of Part I comes from Guoxin Peng's \emph{Third-Order Consciousness: From the Biological Self and Social Self to Vacant Consciousness (The Geometry of Existence Book 4)}, edited by Pinchuan Xie \citep{toc2026}; it was not originally proposed by the author of this project. The three chapters respectively review the discussions in the original work of three modes of conscious organization, the formation of the subject, and the revolution of the observer. This manuscript reorganizes the argument, distinguishes levels of claim strength, supplements the discussion with scholarly literature, and undertakes a critical reconstruction. If an existing theory can fully absorb a claim, the book permits the status of ``Third-Order Consciousness'' as an independent construct to contract accordingly. Where Krishnamurti, existentialism, psychology, or sociology is involved, the relevant ideas remain attributable to their original authors and research traditions.
\end{sourcebox}

\clearpage
\chapter{Three Modes of Conscious Organization: From Experience and Subject to the Decentering of the Observer}

\par\medskip
\noindent\textbf{Why this chapter matters: slow down the sentence that begins with ``I''}\par
A fully formed ``I'' can easily seem to stand in front of experience from beginning to end: I feel, I judge, I decide, I reflect. Yet if we slow an actual process down, we often see a different sequence. You may already have reacted physically before putting the experience into words; your attention may already have shifted somewhere; a thought may simply come to mind. Language then gives the experience a name, memory connects it with the past, and identity and principle organize it further as ``this is my experience.'' Finally, you may find yourself in a new position and say, ``I am observing myself.''

In this chapter, we do only one thing: we unfold the different operations that are ordinarily compressed into the single word ``I.'' This is necessary if Third-Order Consciousness is not to be read as three personality types, three spiritual ranks, or three mysterious states. It is first of all a distinction among different ways in which experience is organized.
\par\medskip

\par\medskip
\noindent\textbf{Sources for this chapter}\par
This chapter is a structured review of the General Preface, Afterword, and the central discussions of Parts I, II, and IV of Guoxin Peng's \emph{Third-Order Consciousness: From the Biological Self and Social Self to Vacant Consciousness (The Geometry of Existence Book 4)}, edited by Pinchuan Xie. The original work repeatedly emphasizes that the first order is the direct presence of life, the second order is the formation of the subject, and the third order deprives the already formed subject of final authority over experience; the three orders are neither three kinds of people nor a spiritual hierarchy from low to high \citep{toc2026}. This chapter preserves those conceptual relations while adding neighboring ideas from psychology and phenomenology so that the native domain used later in the book is explicit. The supplementary literature is used for positioning and comparison; it does not retroactively prove every claim made in the original work.
\par\medskip

\section{Why Begin with ``Occurrence'' Rather Than with ``I''?}

Everyday language naturally puts ``I'' at the front of the sentence. We say, ``I am thinking,'' ``I am afraid,'' ``I decided,'' ``I like this,'' or ``a proof suddenly occurred to me.'' There is nothing wrong with speaking this way. Without a relatively stable subject, responsibility, commitment, planning, and communication would all become difficult.

The question is narrower: does the grammatical subject tempt us to imagine that psychological processes are always commanded by an already completed subject, with experience occurring only afterward?

Many experiences do not unfold in this order. Imagine that you are reading a proof and suddenly see a connection: ``Wait---perhaps one could take a quotient here.'' You probably did not first issue yourself an internal command: ``Now I will produce a good idea.'' More often, knowledge accumulated through long training is already at work; the material in front of you triggers a connection, and the problem begins to take shape. Only afterward do you claim the idea as ``something I thought of.'' Likewise, after a cutting remark, your chest may tighten, your attention may narrow, and you may already feel the impulse to answer back. Only several seconds later might you put the experience into clear words: ``I was offended.''

We therefore preserve, at the outset, a very weak but crucial distinction:

\begin{quote}
\textbf{The occurrence of experience and the organization of experience as ``I am undergoing something'' need not be the same act.}
\end{quote}

Here, ``occurs first'' primarily marks an analytical priority; it need not be understood as a rigid millisecond-by-millisecond temporal order. We are not trying to prove that every sensation precedes the subject, nor are we searching for an absolutely ``selfless'' neural instant. We simply refrain, for the moment, from treating ``I'' as an unexplained starting point. Sensation, attention, tendencies toward action, memory activation, perceptual judgment, and related processes are allowed to appear first as objects of study. Only then do we ask how a stable position capable of saying ``this is mine'' acquires organizational authority within them.

This small move changes everything that follows. If the ``I'' itself must be formed, then ``my past,'' ``my principles,'' ``my professional identity,'' ``my trauma,'' and ``my choices'' cannot appear merely as finished nouns. We must also ask how body, time, language, memory, and the responses of others gradually weave them together.

\begin{example}[Unfolding a mature subject-predicate sentence]
Imagine that you take the initiative and message someone you have only recently met, but receive no reply. If we write only, ``I became angry because she did not reply,'' the whole process has already been compressed into a finished causal sentence. If we unfold it, the sequence may look more like this: you check the phone several times; your attention keeps returning to the message screen; your body becomes agitated; the thought ``what does this mean?'' comes to mind; similar past experiences are recalled; the anger grows; you begin to judge the silence as failing basic standards of courtesy; and finally you decide, ``I will not contact her again.''

None of this requires us to deny the final sentence, ``I was angry.'' It merely reminds us that a complete self-narrative is often the product of many prior acts of organization.
\end{example}

One further restriction must be in place from the beginning: \textbf{being generated does not mean being unreal}. A subject position that forms later can still genuinely alter behavior. Identity can constrain choice; commitments can operate across time; social roles can redirect attention; long-term narratives can in turn shape bodily reactions. To study the generation of the subject is not to declare the subject an illusion.

\section{The ``Three Orders'' Are Not Three Ranks but Three Organizing Focal Lengths}

The phrase ``Third-Order Consciousness'' is easily misread as an upgrade chart: first order is more primitive, second order more mature, third order highest. The reader then begins asking, ``Which order am I at now?'' or ``How do I rise to the third order?'' The Afterword of the original work explicitly resists this reading: if the third order becomes a new badge of identity, it has already been reabsorbed by the second-order subject \citep[TOC-083; pp.~613--618]{toc2026}.

A better way to understand the three orders is as three \emph{organizing focal lengths}. The original work uses a powerful image: the first order is like a ``field'' not yet excessively cut apart; the second order draws that field into paths and a relatively stable ``point''; the third order allows that point to see again that it is only a focused position within the field, not the master of the whole field \citep{toc2026}. Here ``field,'' ``path,'' and ``point'' are conceptual images used to aid understanding. Part I will not quietly convert them into mathematical objects.

In the plainest language, the three focal lengths can be understood as follows:
\begin{itemize}
\item when experience is organized mainly by bodily state, perception, feeling, immediate attention, and tendencies toward action, we call this the \textbf{first order};
\item when memory, language, identity, role, commitment, and narrative organize experience around an ``I'' extended across time, we call this the \textbf{second order};
\item when the very ``I'' that names, compares, evaluates, and narrates experience also enters the scope of observation, we call this the \textbf{third order}.
\end{itemize}

These are not stages that eliminate one another. A person can shift from one organizing focal length to another within seconds, and several can be partly active at once. The first order does not disappear when the second appears; the second does not become dispensable when the third appears.

Only after this intuition is in place do we give the three working definitions that will recur throughout the book.

\begin{definition}[First-Order Consciousness: phenomenological working definition used in this book]
\ordone{} is a mode of organization in which experience unfolds primarily through bodily state, perception, feeling, immediate attention, tendencies toward action, and coupling with the environment, while a stable subject position and its self-narrative have not yet acquired primary organizational and explanatory authority.
\end{definition}

\begin{definition}[Second-Order Consciousness: phenomenological working definition used in this book]
\ordtwo{} is a mode of organization in which memory, time, language, social role, evaluation, commitment, and narrative organize experience around a relatively stable subject position. This position can thereby say ``my past,'' ``my choices,'' ``my principles,'' and ``my relationships,'' and can use those structures in turn to interpret present experience.
\end{definition}

\begin{definition}[Third-Order Consciousness: phenomenological working definition used in this book]
\ordthree{} is a mode of organization in which the observing position responsible for naming, comparing, attributing, evaluating, and narrating experience itself enters the scope of observation. The subject is no longer presupposed as the final interpreter standing outside experience, but is treated as one component of the same process of occurrence.
\end{definition}

These three definitions belong only to the psychological and phenomenological native domain used here. They are not a standard classification already accepted across psychology, nor do they claim three corresponding neural systems. Later chapters will compare them with existing concepts such as minimal self, narrative identity, metacognition, and decentering, but proximity is not identity.

\section{First-Order Consciousness: Direct Presence Does Not Mean ``Lack of Complexity''}

``Direct presence'' is easily misheard as ``simple,'' ``primitive,'' or ``without memory.'' That would misread the first order. A trained pianist adjusting force in real time, a driver continually correcting direction on a congested road, or a mathematician scanning a proof and immediately feeling that ``something is wrong here'' may all draw on highly complex histories of learning, prediction, and pattern recognition. The complexity is already there; it simply has not been translated at every instant into explicit self-narrative.

Several neighboring concepts in psychology and phenomenology can help locate this layer.

Gallagher's notion of the \emph{minimal self} can be understood, roughly, as the idea that even without a long autobiography, experience may still have a basic first-person orientation: things happen ``from here,'' the body is ``this body,'' and an action may be experienced as initiated ``by me'' \citep{gallagher2000}. This differs from the fully developed \emph{narrative self}. Zahavi's discussion of pre-reflective self-awareness likewise reminds us that first-person givenness need not first pass through the explicit judgment ``I am now experiencing X'' \citep{zahavi2005}.

Research on the body breaks the issue down further. \emph{Interoception} concerns the sensing and integration of bodily signals such as heartbeat, breathing, hunger, pain, and visceral state; Craig's review discusses its relation to subjective feeling and representations of one's bodily condition \citep{craig2009}. The \emph{sense of ownership} and the \emph{sense of agency} must also be distinguished: the former is closer to ``this is my body/feeling,'' while the latter is closer to ``this action was initiated and controlled by me.'' Experimental and clinical work shows that the two do not always move together \citep{braun2018}.

None of these concepts is simply identical with the first order in \emph{Third-Order Consciousness}. Together, however, they support one important reading principle: the ordinary word ``I'' already compresses many distinguishable components.

\emph{Third-Order Consciousness} describes the first order as the primordial presence of life: hunger, pain, risk, desire, space, and rhythm first focus life at ``here''; the original work stresses that having a ``field'' does not yet amount to forming a stable subject node \citep[TOC-017--TOC-025; pp.~95--149]{toc2026}. This book preserves that structural distinction while keeping stronger neurobiological or cosmological explanations in a separate evidential account. Embodiment, ownership, and agency can be related to existing research; no brain region will be directly declared ``the location of First-Order Consciousness.''

\subsection{The ``I Am Here'' of the First Order}

When a foot is caught in a door, pain does not need to be written into a life story before it becomes urgent. When danger approaches from the left, the body need not convene an internal meeting before moving right. There is already direction, a local center, an organization of ``happening here,'' but this does not automatically yield ``the kind of person I am.''

We therefore keep two statements apart:

\begin{quote}
\textbf{A functional center has formed here.}\qquad
\textbf{An ultimate ontological center exists here.}
\end{quote}

The first statement is common in bodily action. The second is a much stronger philosophical claim. The first order needs the former; it does not thereby prove the latter.

\begin{example}[Skilled knowledge can work directly]
Imagine a researcher who has spent years studying a certain class of groupoids. During a talk, one keyword activates a familiar network of knowledge; examples and counterexamples connect rapidly, and a question almost seems to come to mind on its own. Only after the question has taken shape does the researcher begin to put it into words, decide whether it is worth asking, and consider how others may evaluate it. The earlier portion is not devoid of knowledge; the knowledge simply has not yet been fully converted into a subject-project about ``who I am'' and ``how well I am performing.''
\end{example}

\section{All Three Modes of Organization Can Appear within One Experience}

Understanding the orders as focal lengths has an immediate consequence: we do not need to assign every event a single order.

Imagine a researcher preparing for an important talk. Before going on stage, the researcher's heart rate rises, the palms begin to sweat, and attention keeps moving between the slides and the audience. This is one layer of direct coupling between body and environment. At the same time, an earlier failure comes back to mind, together with the thoughts ``I cannot embarrass myself again'' and ``I have to prove that I am well prepared.'' The present moment is now being organized through a self that extends across time. A little later, the researcher notices that the internal voice repeating ``I cannot fail'' is itself amplifying bodily tension and drawing on old shame and imagined evaluation by peers. At that point, the very position doing the evaluation has also entered observation.

These layers do not take turns entering the stage. The first order does not leave when the second appears, nor does the second disappear when the third appears. Actual experience is more like a shifting distribution of weight among several modes of organization: at one moment bodily reaction dominates, at another narrative takes over, and at another the observing position itself briefly becomes visible.

Accordingly, when the book later says that a person ``enters the third order,'' the default meaning is only that a particular organizational relation has become visible at that moment. It does not mean that the person has acquired an irreversible higher identity. Several minutes later, an old narrative may once again occupy center stage.

\section{The First Order Is Not a ``Pure Natural State'' Either}

Another common misunderstanding runs as follows: if the first order is more direct than narrative, perhaps it represents an original purity uncontaminated by language, culture, and history. Adult experience rarely permits such a picture.

A chess player may ``directly see'' danger in a position; a surgeon may ``directly notice'' an abnormality in tissue; a mathematician may ``directly feel'' that the rhythm of an argument is wrong. This directness is itself the sediment of long training. Language, professional habits, social experience, and even past injuries may already shape the sensitivities of attention and perception.

Thus, when this book says \emph{pre-reflective}, it means only that the present process need not first become explicit self-narrative in order to function. It does not mean that the process has no history, no learning, or access to a completely culture-free layer of ``pure experience.'' Zahavi's discussion likewise reminds us that the pre-reflective concerns a mode in which experience is given, not a primitive entity outside time and relationship \citep{zahavi2005}.

This restriction matters. Otherwise the first order is romanticized as a movement ``back to nature,'' and the third order is easily misread as bypassing the second and retreating to a wordless purity. \emph{Third-Order Consciousness} emphasizes precisely the opposite: the third order occurs after the subject has formed. It does not delete language, identity, and narrative; it prevents them from monopolizing final explanatory authority \citep[TOC-083; pp.~614--618]{toc2026}.

\section{Second-Order Consciousness: How Experience Becomes ``Mine''}

If the first order mainly asks how experience can occur directly, the second order begins from a different question: how are scattered occurrences organized into a person who can live, commit, take responsibility, and be recognized by others across time?

This step is indispensable and cannot be reduced to the negative story that ``the self begins to manufacture illusions.'' Without the second order, it would be difficult to keep yesterday's promises, take responsibility for past harm, understand a profession, an intimate relationship, or a long-term plan, or turn dispersed experience into a communicable history. The original work takes a two-sided view of the second order on purpose: it can generate foreclosure and misrecognition, but it also gives consciousness temporal persistence and entrance into language, history, relationship, and creation \citep{toc2026}.

Gallagher's narrative self, McAdams's narrative identity, and Ricoeur's discussion of personal and narrative identity all show, in different ways, that a self extended across time does not arise by simply stringing moments together. Experience is continually selected, reorganized, and interpreted \citep{gallagher2000,mcadams2013,ricoeur1992}. McAdams and McLean's notion of \emph{narrative identity} may be understood first as an evolving life story in which the past is reconstructed, the future imagined, and present action thereby given a relatively unified direction \citep{mcadams2013}. Ricoeur, in turn, draws attention to the tension between remaining the same and becoming oneself in time \citep{ricoeur1992}.

The power of the second order can be seen in a few ordinary sentences:

\begin{quote}
``I have always been like this.'' \quad ``I cannot do this again.'' \quad ``This goes against my principles.'' \quad ``I promised, so I have to go.''
\end{quote}

Each sentence inserts the present moment into a longer history of the subject. Sometimes this protects a sense of persistence across time; sometimes it overgeneralizes. One failed talk may be retained as ``I need to prepare more carefully for public speaking,'' or compressed into ``I am simply not a person who can give talks.'' Both are second-order forms of organization. The former preserves greater room for revision; the latter more readily turns a local event into identity.

The second order can therefore be understood as an enormous organizational achievement. It makes the ``I'' into an interface for responsibility across time, an interface for memory, and a social interface. What the third order will later loosen is the status of this interface as ``final master,'' not the totality of its functions.

\section{Why the Small Word ``Again'' Matters So Much to the Second Order}

A very ordinary word will recur throughout this book: \emph{again}.

``Again'' looks as though it merely marks repetition. In practice it often already performs a trans-temporal act of organization. Two sounds can occur twice and be described simply as ``two sounds.'' But when someone says, ``Why am I being ignored again?'' ``I am doubting my abilities again,'' or ``I have ended up in this kind of relationship again,'' the present experience has already been linked to the past and placed within one history of the subject.

At least three things are happening at once: certain materials from the past are retained; some similarity is recognized between present and past; and a temporally enduring ``I'' recognizes that similarity as ``the same pattern encountering me again.''

This is why \emph{again} deserves more attention than mere similarity. Similarity can be a relation between two objects. ``Again'' usually also contains a trans-temporal indexing position: who is saying ``again,'' and which two stretches of experience are being linked by that position?

The same reality can be organized at different levels. Suppose that after a date, the other person does not make further contact:

\begin{quote}
A: ``She did not reply.''\\
B: ``She did not reply, and that made me uncomfortable.''\\
C: ``Why do I keep running into this kind of disrespect again?''\\
D: ``I always attract people like this.''
\end{quote}

A mainly reports a present fact. B connects the fact with present feeling. C places the present experience into a pattern extending across time. D further compresses the pattern into a story about identity and fate. There is no simple rule that ``the farther down the list, the more wrong.'' C may genuinely identify a regularity, or may mistake rough similarity for sameness. D may help summarize long-term experience, or may harden accidental materials into identity.

Before any mathematics appears, the book therefore preserves a purely psychological--phenomenological question:

\begin{quote}
\textbf{What turns ``appearing again'' from local repetition into the subject-level experience that ``it has come back''?}
\end{quote}

This question will later become a major entry point for structural translation. In Part I, however, we only clarify the problem; we do not borrow later mathematical language to answer it in advance.

\section{Narrative Is a Tool of Compression and a Source of Explanatory Authority}

Human experience is too dense to preserve second by second, much less explain second by second. The subject has to compress. A day can be summarized as ``I spent most of today rushing to finish a paper''; ten years of professional life can become ``I have always worked at the intersection of group theory and dynamics''; a complicated relationship can be compressed into ``our values did not fit.''

This compression is not a defect of narrative. It is its function. Without compression, it would be difficult to plan, remember, explain ourselves to others, or learn from experience.

Bruner's work on narrative emphasizes the mutual shaping of life and self-understanding through story form \citep{bruner1987}. Conway and Pleydell-Pearce's self-memory system, in turn, reminds us that autobiographical memory is not the retrieval of raw footage from a warehouse. Current goals, self-knowledge, and past material jointly participate in forming a particular act of remembering \citep{conway2000}. Nelson and Fivush further stress that the development of autobiographical memory depends on language, temporal understanding, the social practice of talking with adults about the past, and developing understandings of self and others \citep{nelson2004}.

In other words, ``my past'' is itself organized. This does not make it false. A compressed map can still guide action with great accuracy. The question is what remains after compression, and whether what was omitted can still re-enter view.

The Geometry of Existence calls the latter danger \emph{structural blockage}: once a narrative becomes immune to counterexamples, it begins to move from a tool for organizing experience toward an explanatory position that exceeds its evidence \citep{geometry2026}. This will recur throughout the book, because one of the second-order subject's strongest capacities is precisely its ability to compress large amounts of complex material into an ``I'' sufficient for action.

\begin{example}[How a state is compressed into identity]
Failing to understand a talk in an unfamiliar field can be described as ``this talk lies far from my area of research.'' It can also be summarized as ``my range of knowledge is too narrow,'' and compressed further into ``perhaps I am not really a broadly educated mathematician.'' Each move adds organization and meaning, but also enlarges the inferential distance. The third order does not automatically negate the latter two sentences. It asks us to see again the variables compressed in between: distance between fields, prerequisite knowledge, attentional state that day, the comparison class being used, long-term reading history, and the situation in which the sentence was spoken.
\end{example}

\section{The Entry into the Third Order: Why Ordinary Reflection Is Not Yet Enough}

At this point the reader may ask: the second-order subject can already reflect on itself in sophisticated ways. It can keep a journal, undergo psychotherapy, learn attachment theory, and recognize defensive patterns. Why introduce a separate ``third order'' at all?

The distinction is not about whether reflection is deep enough. It is about whether reflection preserves one final position that is itself exempt from reflection.

Ordinary reflection often has a structure like this:

\begin{quote}
My emotions may be unreliable.\\
My memory may be biased.\\
My personality patterns may need revision.\\
\textbf{But the ``I'' currently analyzing all this is provisionally treated as a relatively reliable analyst.}
\end{quote}

The third order continues from the last line. It asks not only ``am I seeing myself accurately?'' but also: what is the position that is looking made from? Which memories, standards, ideals, and social evaluations does it borrow? Why does it automatically possess the authority to judge?

Someone may say, ``I must overcome my anger.'' This can help stop harmful action. Yet the position from which ``must overcome'' is spoken may itself be composed of old shame, an image of oneself, an ideal of maturity, and fear of losing control. The third order does not decide in advance that this position is good or bad. It does something more basic: it refuses to move a voice outside the psychological structure merely because the voice sounds rational.

This is the heart of what the original work calls the loss of ``final authority by the subject'': the subject can still act, take responsibility, judge, plan, and express itself, without thereby possessing natural final authority over the meaning of all experience \citep[TOC-083; pp.~613--618]{toc2026}.

The contrast between ordinary reflection and the third order can be condensed into a small table:

\begin{center}
\small
\begin{tabularx}{0.96\textwidth}{Y Y}
\toprule
Ordinary reflection more often asks & The third order continues by asking \\
\midrule
``Why am I so angry?'' & ``Where did the rule by which this anger is being evaluated come from?'' \\
``How do I get rid of this pattern?'' & ``Who recognized these events as the same pattern?'' \\
``Am I too sensitive?'' & ``At what scale, and against what comparison class, was `too sensitive' judged?'' \\
``What kind of person should I become?'' & ``How is this `should become' organizing present experience?'' \\
\bottomrule
\end{tabularx}
\end{center}

This is not an invitation to generate an infinite recursion of ``I observe myself observing myself.'' The impulse the third order is trying to undo is precisely the impulse to create yet another higher observer. The observing position is to be placed within the same process, not continue escaping upward.

\section{Why Krishnamurti Becomes a Central Stress Test for Third-Order Consciousness}

Krishnamurti becomes crucial here not because he supplies a doctrine that can be directly mathematized, but because he pushes the question ``is the observer really outside the structure?'' to a point that is difficult to evade. Chapter 13 of \emph{Third-Order Consciousness} concentrates on themes such as ``the observer is the past,'' ``the observer is the observed,'' and ``choiceless awareness,'' treating them as an important line of division between the third order and ordinary self-reflection \citep[TOC-067--TOC-071; pp.~490--512]{toc2026}.

For Krishnamurti, the psychological observer is composed of conclusions, memory, knowledge, experience, cultural conditioning, and past reactions; he also rejects the hierarchical picture in which a ``higher self'' monitors a ``lower self'' \citep{krishnamurti1954,krishnamurti1969,kft2026a,kft2026b}. For this book, that provides a stringent test:

\begin{quote}
\textbf{If the so-called third order still ultimately requires a super-observer who stands outside the structure, is naturally neutral, and is always right, then we have merely installed another layer of management software on top of the second order.}
\end{quote}

The third order is therefore closer to making organization transparent. Anger is seen, and the rule ``I should not be angry'' is also seen. Desire is seen, and the goal ``I must become a better version of myself'' is also seen. Judgment of another person is seen, and the way an old image rushes ahead to name present reality is also seen. No new position automatically becomes the final master simply because it is ``observing.''

This chapter only establishes the line of thought. Chapter 3 revisits Krishnamurti in its own right and develops the logic connecting ``the observer is the past,'' ``the observer is the observed,'' and ``choiceless awareness.''

\section{What the Third Order Is Not: Five Common Misreadings}

A new concept is most vulnerable to distortion once it becomes popular. Third-Order Consciousness is particularly vulnerable because the second-order subject can easily repackage it as ``a more advanced me.'' Before proceeding, we therefore exclude five common misreadings.

\paragraph{The third order is not emotional suppression}
A person can become highly skilled at not expressing anger while remaining organized internally by the rule ``I must control myself.'' Behavioral control is sometimes necessary, but successfully suppressing an emotion does not automatically show that the structure of the observer has become transparent.

\paragraph{The third order is not sophisticated self-analysis}
Being able to speak fluently about one's attachment style, trauma history, defense mechanisms, or personality traits may increase understanding, but it can also produce a more rigid system of self-labels. Knowing more facts \emph{about} the ``I'' and seeing \emph{how this I is organized by knowledge} are not the same thing.

\paragraph{The third order is not numbness, detachment, or loss of reality}
If a person loses contact with experience and is left only with a spectator-like feeling, that state should not be called decentering merely because it feels as though one is ``standing farther away.'' Depersonalization, dissociation, and related clinical phenomena have their own mechanisms, risks, and diagnostic contexts. This book does not exchange those concepts for Third-Order Consciousness.

\paragraph{The third order is not ``I no longer have a self''}
This sentence is especially liable to become a new badge of identity. The original work repeatedly emphasizes that the third order does not delete the functions of the subject. Language, identity, responsibility, and planning remain. What changes is that the subject no longer mistakes its current mode of organization for the final layer of all experience \citep[TOC-083; pp.~616--618]{toc2026}.

\paragraph{The third order is not nihilism}
Seeing that the subject is generated through body, language, memory, and social relationship does not imply that commitment, responsibility, and meaning are false. Generated structures still have real effects on action. The crisis of value exposed by Nietzsche cannot be reduced to ``nothing matters,'' and the intellectual history of nihilism is far more complex than the simple disappearance of value; Sartre's postwar response reconnects the absence of external guarantees with concrete responsibility \citep{nietzsche1974,gillespie1995,sartre2007}. The third order continues to ask how the one who takes responsibility is formed; it does not use that question to abolish responsibility.

\section{Closing Chapter 1: From ``Experience Occurs'' to ``The I That Organizes Experience Is Also Seen''}

This chapter has not tried to prove a complete theory of consciousness. It has performed three acts of loosening.

First, it separated ``experience occurs'' from ``I am undergoing the experience,'' so that the subject is no longer preloaded as the first cause of every psychological process.

Second, it separated direct experience from the subject extended across time, allowing us to see both the necessity of the second order and its risks: the subject makes responsibility, history, and creation possible, but can also compress local experience into excessively rigid identity.

Third, it placed ``the content being observed'' and ``the I doing the observing'' back within the same problem, so that the third order no longer means higher self-monitoring but reflexive visibility of the organizing position itself.

If the whole chapter is compressed into three sentences, they are these:

\begin{quote}
\textbf{First order:} experience unfolds through direct body--environment coupling.\\
\textbf{Second order:} experience becomes ``mine.''\\
\textbf{Third order:} how ``mine'' is formed also enters observation.
\end{quote}

What matters is not memorizing the three sentences, but preserving the six questions they raise:
\begin{enumerate}
\item How is an ongoing flow of experience segmented into units such as ``this time'' and ``that time''?
\item Under what conditions are segments from different moments organized by the subject as ``again''?
\item Why is present experience organized from a particular ``here,'' and how is that observing position formed?
\item When an established recognition of one's history receives additional background, is it canceled, revised, or placed within a larger context?
\item How does narrative compress experience while preserving necessary differences and routes for revision?
\item If the observer itself is formed by memory, language, rules, and past experience, how can it observe itself without pretending to stand outside the system?
\end{enumerate}

These questions run throughout the book. The next chapter remains within the psychological, philosophical, and social-theoretical native domain and focuses on the conditions under which the second-order subject is formed: how body, time, memory, language, others, and narrative jointly provide a temporally enduring ``I.''

\clearpage
\chapter{How the Subject Is Formed: Body, Time, Language, Others, and Narrative}

\par\medskip
\noindent\textbf{Why this chapter matters}\par
Each morning when we wake, we almost never need to prove anew that ``the person I am today'' is the same person as ``the person I was yesterday.'' The body is already here; the name is already here; unfinished business from yesterday still counts as ``my business''; other people's expectations of me remain in force; old promises continue to constrain the present; plans for the future seem to extend forward from the night before. A temporally enduring ``I'' therefore feels self-evident.

In this chapter, we take that apparent self-evidence apart again. We are not looking for a single organ hidden in the brain whose job is to ``make a self,'' nor do we dismiss the subject as an illusion. Instead, we ask how bodily ownership, temporal persistence, reconstructed memory, linguistic self-reference, the responses of others, social roles, desire, and narrative gradually become intertwined until many different processes can speak under a single grammatical subject. This is the power of Second-Order Consciousness: scattered parts of life can be organized as ``my life.'' Its risk begins at the same point. A position formed to organize experience may gradually come to mistake itself for the final master of experience.
\par\medskip

\par\medskip
\noindent\textbf{Sources for this chapter}\par
This chapter follows the line of argument in \emph{Third-Order Consciousness} concerning bodily focusing, temporal stitching, language, the social self, desire, and the ``narrative-path self.'' It especially preserves the central judgment of the original work that ``Second-Order Consciousness is the formation of the subject'' \citep{toc2026}. Literature on the minimal self, interoception, autobiographical memory, narrative identity, Vygotsky, Cooley, Mead, and Goffman supplies neighboring scholarly coordinates. The original conceptual framework, existing psychological research, and the expository reconstruction undertaken in this manuscript carry different responsibilities. The chapter remains a review and clarification of the native domain of Third-Order Consciousness; it does not import the later mathematical language in advance.
\par\medskip

\section{What the Formation of the Subject Really Has to Explain: Why Do Many Processes Speak as One ``I''?}

``The subject is constructed'' is an easy sentence to say and an easy one to empty of content. If it means only that ``the self is influenced by many factors,'' it adds almost nothing. The harder question is more concrete: bodily sensation, action, memory, language, social evaluation, and life narrative are plainly different processes. Why are they increasingly and stably booked to the same account, and why do they come to constrain one another in the name of ``I''?

Let us consider an ordinary morning. The alarm rings, and you are still tired. A few seconds later, you remember a morning meeting and the material you promised yesterday to send someone. Then you open your phone and see a message that makes you uncomfortable. In less than a minute, several times, social roles, and feelings have been gathered under the same ``I'': \emph{I} am tired, \emph{I} owe the material, \emph{I} am affected by the message, and \emph{I} will still need enough energy to work in the afternoon. Everyday consciousness does not normally treat these as the affairs of several unrelated subjects. It rapidly organizes them around one subject position.

\emph{Third-Order Consciousness} gives a concise account of the second order: body, memory, language, role, recognition, and desire are organized into a relatively stable subject; that subject gains persistence across time and, with it, the capacities for responsibility, history, relationship, and creation \citep[TOC-083; pp.~608--612]{toc2026}. Two judgments must be preserved at the same time.

First, the formation of the subject is a \textbf{real functional achievement}. If yesterday's promise cannot constrain today, if pain cannot be stably attributed to this body, if names, roles, and memories reset completely from moment to moment, it becomes difficult to sustain care, research, cooperation, legal responsibility, or long-term planning. The subject does not lose its reality merely because it is ``formed.''

Second, functional unification does not imply a constant entity independent of the processes being unified. A university is continually maintained through institutions, people, buildings, archives, and rules, yet it has real consequences. In the same way, a subject can be an ongoing organizational result while retaining powerful behavioral, social, and ethical effects. What needs explanation is how this unification is maintained, not the prior assumption of a mysterious ``true self'' behind it that performs all the unifying.

To keep later discussion from pouring everything back into the vague word ``self,'' we can begin with a functional ledger. Bodily ownership answers ``where is this happening?''; agency answers ``who is doing it?''; time and memory answer ``why do yesterday and today belong to the same history?''; language provides reusable forms such as ``I,'' ``mine,'' and ``I should''; others and institutions give identity recognizable social positions; narrative compresses these materials into a life course that can be reviewed, explained, and planned. The stability of the everyday ``I'' emerges from the long-term coupling of these layers; no additional mysterious substance standing outside them need be posited.

\section{Body: Before ``Who Am I?'' Comes ``Things Are Happening Here''}

The body is one of the easiest layers of subject formation to overlook. We usually look back at the body after we already know how to say ``I,'' and so it is tempting to imagine the order as follows: first there is a complete me, and then this me owns a body. In the order of experience, many of the earliest orientations require no autobiography.

Imagine that you strike your foot against the corner of a table. Pain immediately draws your attention to the injured foot. Or imagine that you suddenly lose your balance: your body adjusts its posture before you have time to put anything into words. As hunger increases, the tendency to look for food becomes stronger. In each case, a strong local priority is already present: this change is happening \emph{here}, and a response is required from \emph{here}. \emph{Third-Order Consciousness} describes this layer as the focusing position of life and stresses that a ``local center'' is not yet the fully formed subject with identity, story, and social status that appears later \citep[TOC-017--TOC-025; pp.~95--149]{toc2026}.

In psychology and cognitive science, \emph{interoception} generally refers to the sensing and representation of internal physiological states such as heartbeat, breathing, gastrointestinal state, temperature, and internal arousal. Craig's review relates interoception to subjective feeling and bodily-state representation \citep{craig2009}. We do not need to accept a grand explanation centered on any single brain region. We need only retain a plain fact: bodily state participates in determining what appears urgent, what captures attention, and what is readily interpreted as threat.

Two terms that are often bundled together must also be separated. The \emph{sense of ownership} can be understood roughly as ``this body, this feeling, belongs to me.'' The \emph{sense of agency} is closer to ``this action was initiated or controlled by me.'' Gallagher offered a classic distinction between the two, and later experimental and clinical research likewise shows that they do not always coincide perfectly \citep{gallagher2000,braun2018}. Even the most intimate ``I,'' then, is not an indivisible unit: \emph{belonging to me} and \emph{being done by me} are different relations.

\subsection{The Body Is Not Background Noise: It Can Change the Meaning of the Same Sentence}

Imagine receiving a perfectly ordinary work email: ``Could you revise this part again?'' When well rested, it may register as no more than a request for revision. After several nights of poor sleep, a recent failure, and sustained bodily tension, the same sentence may quickly be read as ``does this person think I am incompetent?'' The latter interpretation is not necessarily wrong, and the former is not automatically more objective. What matters is that judgment does not occur suspended outside the body.

Fatigue can shorten tolerance; sustained arousal can make threatening cues more salient; physical discomfort can change which memories and ranges of attention are currently accessible. When Third-Order Consciousness later asks that ``judgment have a position,'' the earliest layer of that demand is to write the bodily position back in. When I say, ``I am certain this is how things are,'' I can also ask, ``In what bodily and emotional condition did this certainty form?'' This does not explain all meaning away in bodily terms. It simply refuses to let interpretation pretend that it has no body.

The first thing the body gives the subject, then, is not a philosophical definition but an enduring ``here'': pain is here, hunger is here, action begins here, consequences are borne here. Name, memory, and social role will later thicken this ``here'' layer by layer until it becomes ``who I am.''

\section{Time and Memory: From ``Here'' to ``Still Here,'' and Then to ``This Is My Past''}

Bodily focusing alone is not enough to produce the subject we recognize. If every pain, hunger, comfort, and action flashed up as an unrelated event, experience might have a ``here'' but would have difficulty sustaining an ``I'' across time. \emph{Third-Order Consciousness} treats persistence across time as central to the subject's increasing thickness: what happens today can be linked to what happened yesterday, and a local center acquires the diachronic sense of ``still here'' \citep[TOC-026--TOC-030; pp.~155--180]{toc2026}.

One misunderstanding must be avoided immediately. The first order already lives in time. Breathing has rhythm, actions have sequence, danger requires prediction of the next second. What the second order adds is not ``time itself'' but time lengthened, retained, reconstructed, and attributed to the same subject. Yesterday no longer leaves only physiological traces; it can become ``the thing I did wrong yesterday.'' Tomorrow is no longer merely the next stimulus; it can become ``the kind of person I want to be next year.''

This development brings capacities indispensable to civilized life: promise, planning, responsibility, waiting, delayed gratification. Something promised yesterday remains in force today; a skill learned ten years ago still counts as ``my ability''; a goal that has not yet occurred can constrain present action. At the same time, anxiety, regret, and identity pressure acquire longer temporal reach. The more successfully the subject organizes itself across time, the more persistently it can also be surrounded by its own history and future.

\subsection{Autobiographical Memory Is Not an Internal Video Archive}

``My past'' can feel like a warehouse already there, ready to be retrieved. Actual autobiographical memory is much more complicated. Conway and Pleydell-Pearce's self-memory system model emphasizes the interaction between retrieval of particular autobiographical memories, the goals of the current \emph{working self}, retrieval cues, and structures of autobiographical knowledge \citep{conway2000}. We do not simply pull a complete recording from the brain; present concerns participate in deciding which fragments become accessible, how they are arranged, and what receives greater weight.

This does not mean ``all memories are made up.'' Real events, other people's testimony, photographs, written records, and repeated experience constrain memory. A more precise distinction is between what happened in the past and the way a usable history is organized from the past in the present.

When a young researcher receives several rejections in succession, for example, earlier episodes of ``not being recognized'' may suddenly become densely available and form the historical feeling ``I have never been recognized.'' Several years later, if some projects mature, the same past can be reorganized as ``those failures forced me to narrow the problem and build my own direction.'' Both narratives may cite real events. They differ in selection, linkage, and evaluation.

Nelson and Fivush's work on the development of autobiographical memory adds another reminder: the mature category ``my past'' does not arise automatically inside an isolated brain. Language ability, temporal understanding, parent--child conversations about the past, culturally patterned modes of narration, and understandings of self and other all participate in its formation \citep{nelson2004}. When an adult says ``this is my own memory,'' the word \emph{own} therefore does not mean ``without social history.''

\begin{example}[One research history, two different life-lines]
A researcher looking back over the years before and after a postdoctoral period can say, ``I spent years failing as I chased problems that were too hard.'' The same person can also say, ``It took me several years to learn how to break large problems into structures I could actually work on.'' Neither sentence need falsify a single concrete event, yet each chooses a different beginning, causal center, and outcome. Third-order observation does not select the more encouraging story; it first makes the act of organizing the story visible.
\end{example}

The subject is therefore both shaped by the past and continually reorganizing it in the present. This already raises an important third-order question: when I say ``my history proves that I am this kind of person,'' which parts of my history have I brought into the present judgment? Which parts have I left out? How has my present goal changed what I can see as evidence?

\section{Language: Experience Acquires Names, and a Stable Subject as Well}

Body and memory can give experience local temporal persistence, but the second-order subject also requires a stronger tool of compression: language. Through language, the past can be brought back again and again, feelings can be classified, social rules can continue to operate when no external commander is present, and the pronoun ``I'' can occupy a reusable position.

\emph{Third-Order Consciousness} makes a strong claim about language: language does not merely translate a completed inner experience into sound. Naming, grammar, and symbol participate in segmenting experience and continually reinforce the position of ``I'' as the center to which experience is attributed \citep[TOC-035--TOC-040; pp.~212--255]{toc2026}. This book preserves the central relation that language participates in organizing experience, while remaining restrained about stronger forms of linguistic determinism. Pain does not cease to exist when it lacks a name, nor can every form of subjectivity be caused by subject--verb--object grammar alone. The weaker proposition we actually need is this: whenever experience is preserved across time, communicated publicly, used in self-regulation, or attributed to identity, reusable symbolic systems participate deeply in that organization.

The Vygotskian tradition concerning the development of language and thought offers a useful bridge. Children initially hear externally spoken directives such as ``wait,'' ``do not touch,'' and ``finish this first.'' Such social speech can gradually become part of internal self-regulation \citep{vygotsky1986}. In adulthood, even when no one else is in the room, we still say to ourselves, ``finish this first,'' ``do not look so emotional,'' or ``you should be more mature.'' External rules have become internal positions from which speech occurs.

\subsection{Naming Can Give Access to Experience, and It Can Also Seal Experience Off}

The power of language works in two directions. Learning words such as ``panic attack,'' ``burnout,'' or ``grief'' can allow a person, perhaps for the first time, to place vague and isolated suffering within public knowledge, find similar experiences, seek professional help, and gain a communicable interface. Naming here enlarges the space for action.

The same word can continue farther. ``I have been noticeably anxious lately'' can become ``I am an anxious person,'' and then ``I will ruin any important occasion.'' The distance increases as language moves from describing a state to defining identity. The question worth observing is not simply whether ``labels are good or bad,'' but how much work a given word is now being asked to do. Does it still point back to concrete experience, or has it begun issuing verdicts about the future?

This is also why a phrase such as ``I should'' deserves attention. ``I should be more generous,'' ``I must act professionally,'' and ``I cannot contact them again'' may sometimes supply necessary rules; at other times they import social evaluation directly into internal life. The third order does not simply delete such rules. It gives the rules a history as well: who taught me to judge this way? In what relationships did the rule once protect me? How broadly does it still apply now?

Language thus performs two operations. It makes experience readable, and it makes experience attributable. A flowing process is repeatedly phrased as ``I feel,'' ``I remember,'' and ``I decide.'' The subject position gradually becomes like a fixed chair, making it easier and easier to forget that what occupies the chair is continually changing.

\section{Others and Society: I Am Not Only Living; I Am Also Being Seen, Called, and Expected}

Even with body, memory, and language, the socially meaningful ``I'' cannot be completed in solitude. A name becomes an identity because other people call it. Positions such as ``student,'' ``teacher,'' ``expert,'' ``partner,'' and ``good child'' carry weight because interaction attaches expectations, evaluations, rights, and responsibilities to them.

\emph{Third-Order Consciousness} treats ``the gaze of others'' as an important force that thickens the second-order subject: how others see, expect, recognize, and reject me no longer remains only an external event, but can be written into the subject and become part of self-evaluation \citep[TOC-032--TOC-034; pp.~195--211]{toc2026}. The other person need not literally be present for this inscription to operate. Someone alone in a room can still worry, ``what will people think?'' because social evaluation has already been internalized as material used in present judgment.

Cooley's tradition of the \emph{looking-glass self} emphasizes that we often form self-feelings through imagined images of ourselves in other people's eyes and through imagined evaluations of those images \citep{cooley1902}. Mead places the self within social interaction and role-taking, using distinctions such as ``I'' and ``me'' and the notion of the generalized other to discuss how social rules enter the relation of the individual to the self \citep{mead1934}. Goffman, beginning from concrete interaction scenes, self-presentation, and impression management, further shows that ``the kind of person I am'' always appears within a situational order that makes a certain self intelligible to others \citep{goffman1959}.

The common lesson is not that ``other people determine everything.'' It is that other people's responses do more than evaluate an already completed subject. They help supply coordinates through which the subject understands itself.

\subsection{Academic Identity: Why a Familiar Field and an Unfamiliar Field Can Feel Like Two Different ``I''s}

``I am a researcher in this field'' is never only a sentence said privately in the mind. It is maintained through training, advisers, papers, conferences, peers, journal classifications, employment structures, and repeated engagement with problems. During a talk in a familiar area, keywords can immediately activate dense existing networks and questions may arise naturally. In an unfamiliar area, the same intelligent person may suddenly fall silent and even think, ``Why do I know nothing?''

A real difference in knowledge must be distinguished here from an evaluation of the subject. An unfamiliar talk may indeed expose a limit in one's knowledge. Yet ``I cannot immediately generate a question in this setting'' does not automatically mean ``my knowledge as a whole is impoverished.'' The jump from a local state of access to a global identity judgment passes through comparison classes, distance between fields, professional self-esteem, and social evaluation.

Such examples show that stability and situatedness coexist in the subject. A person does not become someone else simply by entering a different seminar room, but different settings activate different identity coordinates and make certain memories, capacities, and evaluations more salient. The entrance to the third order is here: not only ``how capable am I?'' but ``which social position is evaluating my capacity at this moment?''

\section{Roles and Institutions: Society Does Not Enter the Subject Only through ``What Other People Think''}

It is still too narrow to understand all social influence as ``the gaze of others.'' Real roles carry real consequences. Students, teachers, physicians, parents, editors, referees, employees, and managers possess different permissions, responsibilities, and behavioral rules. Institutions repeatedly confirm these positions and give them durability through salary, evaluation, reward and punishment, qualification, contract, and procedure.

Berger and Luckmann's analysis of institutionalization, knowledge, and socialization shows that shared reality is not invented from zero by each individual every morning. Institutions give repeatability to a set of roles and interpretive forms \citep{berger1966}. Goffman's interaction analysis similarly shows how a scene determines which forms of self-presentation are feasible and which expressions carry higher cost \citep{goffman1959}.

Thus, after someone says psychologically, ``I no longer want work to define my identity,'' work identity may still enter life through evaluations, income, colleagues, students, and deadlines. The subject is not maintained by an internal story alone. External structures continue calling it into action.

This matters greatly for psychology. Many conflicts that appear to be private ``internal friction'' are in fact simultaneous demands from several roles placed upon one person. As a researcher one may need to maintain standards; as a collaborator one may need to compromise; as a child one may need to provide care; as a partner one may need to be present. If all these tensions are classified as ``personality problems,'' social structure disappears from the explanation.

The second-order subject is therefore not a lonely sphere. It is more like a public position repeatedly summoned by many relationships: a name calls you out; a role requires a response; an institution stores your history; other people continually send expectations back toward you. Over time, these external relations become the internal voice of ``how I should be.''

\section{Narrative: How the Subject Compresses Fragments into ``One Life of Mine''}

By this point, the body has supplied ``here''; time and memory have supplied ``still the same person''; language has supplied a reusable grammatical subject; others and institutions have supplied that subject with social position. Yet life remains too heterogeneous. Success and failure conflict with one another, relationships change, memory is incomplete, desires are inconsistent. The second order therefore needs one final powerful mode of organization: it turns heterogeneous material into a coherent course one can live through.

\emph{Third-Order Consciousness} calls the result the ``narrative-path self'': the subject continually books memory, shame, recognition, desire, failure, and moments of achievement to a common account, until life increasingly resembles a line that can be reviewed, explained, and made coherent \citep[TOC-052--TOC-056; pp.~351--391]{toc2026}. The original work criticizes this process sharply but also acknowledges the need for persistence across time. Without some kind of story, it is difficult to explain why today's cost is worth bearing for a goal ten years away, or to communicate dispersed experiences to another person.

Bruner stresses the reciprocal shaping of life and narrative understanding \citep{bruner1987}. McAdams's work on narrative identity treats an internalized and evolving life story as an important level of personality organization, joining reconstructed past and imagined future so as to provide life with unity and purpose \citep{mcadams2013}. Ricoeur's discussion of narrative identity shows how personal identity can be maintained through ongoing change without requiring an eternally fixed set of attributes \citep{ricoeur1992}.

Together these approaches remind us that narrative is not decorative material added after the facts have already been written. It participates in deciding what counts as a beginning, what counts as a turning point, what is treated as a cause, what is worth remembering, and where the future ought to continue.

\subsection{Why Narrative Is Necessary: Life Cannot Be Stored Frame by Frame}

No one can preserve every minute of a year and then consult all of it before making a decision. The subject must compress. A week of research can become ``I spent most of the week revising a proof.'' Several professional years can become ``I have always worked where group theory and dynamics meet.'' A complicated relationship may eventually become ``we wanted different things from the relationship.'' Such compression increases the efficiency of action and makes experience communicable.

Danger begins when the inferential span of compression becomes too large. Failing to understand a talk can be summarized as ``this talk is far from my field,'' then as ``my knowledge is too narrow,'' and then as ``I will always be confined to a very narrow circle.'' Each step may preserve some information, but the scope of the conclusion expands faster than the evidence necessarily does.

The progression can be written as an ordinary psychological chain:
\[
\begin{gathered}
\text{one state}\longrightarrow\text{repeated description}\longrightarrow\text{stable trait}\\
\longrightarrow\text{identity}\longrightarrow\text{prediction of fate}.
\end{gathered}
\]

This is not a formula for error. Long-term evidence can genuinely support stable traits. What the third order asks is whether every expansion preserves its sources, its scale, and the conditions under which it can be revised. When a sentence such as ``this is simply what I am'' has become incapable of being altered by any new material, narrative begins to move from a tool of compression toward an explanation that claims final authority.

\section{Desire and ``Becoming'': Why the Subject Keeps Sending the Present toward a Future Version of Itself}

Once the subject possesses a story, it rarely remains at ``what am I now?'' It continues asking: what am I missing? Do others recognize me? Who am I to become next? Desire and comparison thus become important fuel for Second-Order Consciousness.

In its discussion of desire, \emph{Third-Order Consciousness} repeatedly emphasizes that an object often does more than simply exist as an object. It can also carry status, recognition, imagined completeness, and an image of a future self \citep[TOC-046--TOC-051; pp.~305--350]{toc2026}. Position, relationship, wealth, knowledge, and achievement can all have practical value and subject-level meaning at the same time. Losing an object can therefore mean not only ``the thing is gone'' but ``the person I thought I would become through this has failed to appear.''

It would be inappropriate to compress Schopenhauer, Nietzsche, Lacan, Foucault, Girard, and their very different theories of desire into a single mechanism \citep{schopenhauer1969,nietzsche1997,lacan2006,foucault1978,girard1965}. This chapter retains only one ordinary relation directly relevant to Third-Order Consciousness: the second-order subject often lives between ``who I am now'' and ``who I should or hope to become.''

This distance is not inherently harmful. Learning, training, ethical change, and professional planning all depend on orientation toward the future. The problem appears in the next chapter: if even ``observing myself'' is taken over by a future ideal--I will learn to observe so that I become more mature, calmer, more advanced--has observation been assigned a direction by the subject-project before it even begins? This is where Krishnamurti's criticism of \emph{psychological becoming} enters.

\section{How Eight Functions Are Stitched into One Subject: Unity Is Ongoing Work}

We can now come back to the question with which the chapter began. An everyday subject performs at least the following kinds of work:

\begin{center}
\small
\begin{tabularx}{0.96\textwidth}{>{\bfseries}p{0.28\textwidth}Y}
\toprule
Functional layer & What it makes possible \\
\midrule
Bodily ownership & ``This is happening to this body,'' giving local states priority. \\
Sense of agency & ``This action was initiated/controlled by me,'' allowing action to be attributed. \\
Temporal persistence & Organizes today, yesterday, and tomorrow into the diachronic process of one person. \\
Memory organization & Retrieves from a vast past the historical material relevant to current goals. \\
Linguistic self-reference & Stabilizes ``I,'' ``mine,'' and ``I should,'' making reusable forms of self-regulation possible. \\
Social role & Gives the subject a position within other people and institutions that is recognizable, expectable, and accountable. \\
Desire and value-direction & Gives objects, recognition, and future selves different weights, allowing ``what I want'' and ``what I want to become'' to organize present choice. \\
Narrative identity & Compresses heterogeneous experience into a life course of ``who I am, where I came from, and where I am going.'' \\
\bottomrule
\end{tabularx}
\end{center}

These eight layers are not installed once and for all in a fixed order. They shape one another from early development onward and continue to change throughout adult life. Bodily fatigue alters memory retrieval; social evaluation alters bodily tension; linguistic labels alter attention; narrative goals affect which parts of the past are brought back into view; changes of role can force an entire life story to be rewritten.

There is therefore nothing surprising about the subjective feeling ``I am one whole person.'' When many systems have coordinated over long periods, the sense of unity itself becomes very stable. The task of Third-Order Consciousness is not to smash that sense of unity, but to make visible that stable unity is ongoing work: it has components, a history, and possibilities of reorganization.

\begin{definition}[Second-order subject: functional working definition]
This book understands the second-order subject as a functional organizing position continually maintained through time. It can attribute experience as ``mine,'' organize past and future into a traceable history, recruit language, social roles, and evaluative rules, and preserve action, commitment, and identity across time through narrative compression. This subject has real behavioral and social efficacy, but the book does not presuppose it to be an unchanging entity independent of the processes that sustain it.
\end{definition}

This definition is crucial for the remainder of Part I. It explains why the subject can neither be deified nor casually abolished: a position formed through many processes can still decide, take responsibility, hurt another person, care for another person, write a paper, and keep a promise.

\section{Two Complete Cases: How the Subject Turns Local Experience into ``Something about Me'' within Minutes}

\subsection{No Reply: From a Fact to a Principle, a History, and a Future Script}

Imagine that you sent a message and received no reply. If we record only the order in which things happen, the beginning may be simple: you check the phone, your attention keeps returning to it, irritation grows in your body, and different guesses come to mind. As second-order organization thickens, more layers are brought in: an interaction rule---``normally one should reply''; a subject position---``I already took the initiative''; a link to the past---``this has happened to me before''; an identity principle---``I am not a person who repeatedly begs for attention''; and a future script---``if there is contact again, I will make my principles clear.''

These elements cannot simply be dismissed as ``all psychological projection.'' Norms of courtesy may be real, and self-respect may protect necessary personal limits. What is worth observing is how a local nonresponse becomes, within a very short time, evidence about the past, a judgment about identity, and a script for the future. The second-order subject is doing precisely this kind of expansion.

The action may already be complete--one has decided not to make further contact--while the imagined conversation continues. This raises an important question: once action has ended, why does the subject still need, in imagination, to complete an explanation, correction, or restoration of position? The next chapter will not rush to answer. It will place the internal observer who says ``I need to get this straight'' within the scope of observation as well.

\subsection{An Unfamiliar Academic Talk: From Local Accessibility to Global Self-Evaluation}

During a talk in one's own area, years of accumulated knowledge allow keywords to activate examples, counterexamples, and questions rapidly. During a talk in an unfamiliar area, the network is sparse and many connections must be constructed on the spot. A global judgment can then arise with remarkable ease: ``my knowledge is too narrow.''

This sentence may contain a real warning: the researcher may genuinely benefit from reading more broadly. Yet it also merges several different quantities--distance between the talk and one's own field, prerequisite knowledge, immediate attentional load, the people used for comparison, long-term breadth of knowledge, and a professional standard concerning ``how broad a good mathematician should be.''

The second-order subject favors this compression because it needs a rapid answer to ``what does this mean for me?'' The entrance to the third order is not the claim that ``this judgment must be wrong.'' It is the visibility of the compressing process itself: from which local facts did I jump to a conclusion about my whole self? If the scene, scale, or comparison class changed, would the conclusion retain the same strength?

Both cases show that the subject is not formed only slowly over a developmental history. In each concrete moment, the subject position can also be taken up again, reinforced, and revised. Body, past, rule, role, and future script can reconverge on the position ``I'' within seconds.

\section{Closing the Chapter: The Second Order Is Not an Error to Be Destroyed but a Condition the Third Order Must See}

At this point, the central tension in the account of the second order given in the original work can be stated more clearly. The second order stitches direct experience into history, places the body within relationship, and organizes language and role into a sustainable subject. Because of it, human beings can sustain long-term commitments, take responsibility, preserve knowledge, enter communities, create work, and say ``this is my life.'' The original text is explicit on this point: the second order carries risks of foreclosure, but it is also a necessary condition of temporal persistence, responsibility, language, history, relationship, and creation \citep[TOC-083; pp.~608--612]{toc2026}.

The problem begins when the subject gradually forgets its own history of formation. Bodily position, memory selection, social evaluation, linguistic rules, and narrative compression recede into the background. What remains is a sentence such as ``this is simply who I am,'' ``I can see myself objectively,'' or ``I know what this means.'' A functional position then begins to acquire final explanatory authority.

The third order therefore does not add a new capacity outside the subject. It first makes the subject itself an observable question again:

\begin{quote}
\textbf{What materials formed this ``I'' that is remembering, judging, explaining, and planning? Where did it obtain its present standards of evaluation? Why, at this moment, does it connect certain facts while allowing others to recede into the background?}
\end{quote}

As long as this question is still answered by a naturally neutral ``higher self,'' we have not really left the second order. The next chapter therefore enters the crisis of the observer directly: what ordinary reflection, metacognition, and decentering each accomplish; why Krishnamurti insists that ``the observer is the past''; what internal relation of authority is undone by ``the observer is the observed''; and why even the project ``I must learn to observe myself better'' can become another way in which the subject perpetuates itself.

\clearpage
\chapter{The Crisis of the Observer: From Reflection to Krishnamurti's Revolution in Observation}

\par\medskip
\noindent\textbf{Why this chapter matters}\par
The first two chapters have brought the problem to this point: experience can occur before it is fully organized, and the subject can gradually form through body, memory, language, others, and narrative. Yet once the subject learns to reflect, it is easy to preserve one last safe position: \emph{I may be influenced by emotion, desire, and story, but at least there is still a relatively lucid ``I'' that can step back and watch all of this.}

The real difficulty of Third-Order Consciousness lies here. It is not satisfied with helping the subject ``reflect better.'' It asks further: where does the observer that reflects on, evaluates, and manages experience come from? If this observer too is composed of past experience, comparison, ideals, language, and relational history, then ``I observe myself'' is no longer a pure subject facing a problematic object. We now have different positions within the same psychological process. This chapter makes that shift explicit step by step.
\par\medskip

\par\medskip
\noindent\textbf{Sources for this chapter}\par
This chapter mainly reviews Chapter 13, ``The Revolution of Observation,'' and related sections of Guoxin Peng's \emph{Third-Order Consciousness: From the Biological Self and Social Self to Vacant Consciousness (The Geometry of Existence Book 4)}, edited by Pinchuan Xie, especially its lines of thought concerning observation not as self-cultivation, ``the observer is the past,'' ``the observer is the observed,'' choiceless awareness, and relationship as the testing ground \citep{toc2026}. These ideas were not originally proposed by the author of this project. The chapter also revisits Krishnamurti's texts and foundation materials to verify the central propositions \citep{krishnamurti1954,krishnamurti1969,kft2026a}, and uses literature on metacognition, decentering, and related topics as neighboring psychological coordinates. Every later working definition and structural reconstruction of the ``third order'' is an academic restatement built on those sources; it does not retroactively rewrite the original authors.
\par\medskip

\section{From ``I Can Reflect'' to ``Who Is Reflecting?'': How the Crisis of the Observer Appears}

Modern people do not lack the capacity to reflect on themselves. We replay arguments, analyze why we procrastinate, ask after an emotional episode ``what just happened?'' and use psychological language to say ``that was my psychological projection,'' ``I am repeating a relationship pattern,'' or ``an old experience has been triggered.'' Such abilities usually offer more resolution than simply being carried away by a reaction, and they can directly improve behavior.

We therefore do not begin from the claim that reflection is useless. The question is more precise: when more and more of experience becomes an object of analysis, does the position doing the analysis also need an account of where it comes from?

Imagine that someone becomes angry because an expected reply never arrived. Later, the person reflects: ``I know this is my issue. I have to train myself not to be affected by things like this.'' We can now see a clear two-layer structure. One layer consists of anger, disappointment, and impulse. The other is a calmer manager that names and evaluates the first layer and designs a plan for improvement.

Such management is not inherently wrong. It can help prevent impulsive action, help establish workable limits, and form an important part of psychotherapy and self-regulation. Third-Order Consciousness adds only one question:

\begin{quote}
\textbf{Why does this calmer manager assume that ``not being angry'' is the correct goal? Whose standard is it using? What past does it carry? Why does it believe itself to be closer to the facts than the emotion it is managing?}
\end{quote}

As long as this question is not asked, reflection easily preserves an implicit picture: experience can be confused, but the observer is relatively transparent; emotion requires explanation, but the explainer does not.

This is what the chapter calls the \emph{crisis of the observer}. The ``observer'' is not a mysterious entity and not a homunculus hidden in the brain. It is first of all a \textbf{functional position}: at a given moment, experience is organized from that position, and the person says, ``I know what happened,'' ``I will judge what this means,'' or ``I will decide how this should be handled.'' Third-Order Consciousness asks that this position itself enter the problem.

\section{How Far Psychology Has Already Gone: Metacognition and Decentering as Neighboring Coordinates}

Before turning to Krishnamurti, we need to respect work that psychology has already done. Otherwise Third-Order Consciousness can too easily be written as though psychology had never studied self-observation, which is plainly false.

\subsection{Metacognition: Cognition Can Turn Its Own Reliability into an Object of Study}

In the plainest sense, \emph{metacognition} is the monitoring and evaluation of one's own cognition. A person can not only answer a question but also judge, ``How confident am I in this answer?''; not only make a perceptual choice but also report how reliable that choice feels. Research shows that task performance and sensitivity to the reliability of one's own performance can be measured separately \citep{fleming2012,fleming2012a}.

This matters greatly here. It shows that ``cognition evaluating its own cognition at a second level'' can enter experiment and need not remain philosophical language. Yet Third-Order Consciousness is not the same construct as metacognitive accuracy. Someone can know with great precision that ``I am only sixty percent confident in this judgment'' without asking: why am I using this particular standard of evaluation here? Why does this error threaten my professional identity while another error does not?

Metacognition therefore supplies a mature experimental neighborhood. Third-Order Consciousness places further attention on how the evaluative position itself is formed. This does not mean that the third order is ``higher'' than metacognition. Their focal questions differ.

\subsection{Decentering: Not Immediately Treating a Thought as ``Me'' or as Fact}

\emph{Decentering} is often used for a process in which one loosens identification with internal experience. Thoughts still occur, but need not immediately be treated as facts, commands, or complete definitions of the self. Bernstein and colleagues discuss meta-awareness, disidentification from internal experience, and reduced reactivity within a related process model \citep{bernstein2015}.

This is clearly close to Third-Order Consciousness. ``The thought `he does not respect me' has appeared in my mind'' is not the same sentence as ``he does not respect me.'' The first already allows the thought to appear as a psychological event rather than automatically acquiring the status of fact.

The third order, however, continues with a point that can easily be missed: can the observing position that says ``I have learned to treat thoughts as thoughts'' become a new stable identity? A person can readily turn disidentification into ``I am someone who observes well, who is mature, who knows how to decenter.'' The old center has not necessarily disappeared; it may simply have acquired a more refined description of itself.

Metacognition and decentering are therefore psychological neighborhoods that this book must engage seriously. They cannot be swallowed by Third-Order Consciousness, nor can they simply be treated as synonyms. If the third order is ever operationalized, it must show what difference it adds that these existing constructs cannot fully explain.

\section{Critical Reconstruction: What Have Neighboring Theories Already Explained, and What Is Left for Third-Order Consciousness?}

Third-Order Consciousness is still developing. This book has no reason to protect it as an irreducible proper name. A stricter strategy is to let neighboring theories first take away the parts they already explain, and then ask whether a stable relational increment remains.

\begin{methodbox}[Keep the terminology separate: the third order of this book is not Higher-Order Theory]
In philosophy of consciousness, higher-order theories usually ask under what conditions a first-order mental state becomes conscious by virtue of a higher-order representation; Rosenthal's higher-order-thought approach is a representative example \citep{rosenthal2005}. The ``third order'' in this book does not mean a representation nested to a third level, nor is it defined by ``another representation about the current state.'' It describes a change in organizational relation: do the conditions that form the position responsible for naming, evaluating, and governing experience themselves enter the same scope of observation?
\end{methodbox}

\paragraph{Self-model theory has already taken away a large part of the claim that ``the self is not a substance''}
Metzinger understands the phenomenal self as the content of a transparent self-model rather than as an additional entity \citep{metzinger2003}. Therefore, ``the self is a constructed model/interface'' cannot serve as the exclusive novelty of Third-Order Consciousness. If the third order is to add anything, it must explain further how a position currently using a self-model to make judgments can itself be identified as a historically formed, scale-dependent, revisable organizing position.

\paragraph{Predictive processing has already taken away a large part of the claim that ``perception is not a transparent window''}
Predictive processing and interoceptive inference understand perception as inference constrained by models, error, and bodily regulation \citep{seth2015}. Accordingly, phrases such as ``the senses are filters'' or ``experience is compressed'' can function only as sourced scientific metaphors or weaker conceptual claims; rhetorical repetition cannot promote them into new cognitive mechanisms. The additional third-order question is not to announce once more that perception is constructive, but to ask whether the rules currently used to explain, correct, and evaluate perception remain transparent themselves.

\paragraph{Metacognition, decentering, and self-distancing are the most direct empirical competitors}
Metacognition can measure monitoring of the reliability of one's own judgment; decentering can loosen the automatic equation of thought with fact; self-distancing has an experimental tradition comparing immersed and distanced reflection \citep{fleming2012,bernstein2015,kross2011}. If every observable effect attributed to the third order can be fully absorbed by these constructs, the book should contract its claim that Third-Order Consciousness is an independent psychological capacity, while retaining its value as an integrative problem-framework.

\paragraph{Second-order cybernetics already comes very close to ``the observer enters the system''}
Second-order cybernetics places the observer within the observed system and asks how the describer participates in description \citep{vonfoerster2003}. This may absorb much of the conceptual novelty of $P_5$, and may even supply an earlier theoretical ancestor for $P_6^*$. A possible distinctive feature of Third-Order Consciousness is to bring such reflexivity back to the formation of the subject, narrative compression, corrective purpose, and concrete relational practice. But any advantage of this combination must be argued for; naming it does not establish it.

\paragraph{Buddhist traditions of non-self remove any claim to originality for ``the subject is not a permanent substance''}
Indian Buddhist traditions have long developed multiple systematic arguments against a permanent, independent, autonomous self \citep{siderits2007}. This book cannot present ``there is no substantial self'' as a new discovery. A narrower candidate contribution of Third-Order Consciousness is that it retains a functional subject, responsibility, and workable limits while studying how an already formed subject loses final authority. It neither requires the subject to be erased nor turns non-self into a rank of personhood.

\begin{center}
\small
\begin{tabularx}{0.96\textwidth}{>{\bfseries}p{0.31\textwidth}Y}
\toprule
Result of critical reconstruction & How the book should adjust \\
\midrule
Neighboring theories fully explain all relevant distinctions & Demote ``Third-Order Consciousness'' to a clear natural-language entry point; retain the mathematical formalization of the family of established organizational relations. \\
Neighboring theories explain only some distinctions & State the remaining relations explicitly; do not use a single umbrella term to conceal which parts have already been absorbed. \\
Stable incremental predictions appear & Only then advance the corresponding local module to Level III; even then, do not declare that a complete theory of consciousness has been established. \\
\bottomrule
\end{tabularx}
\end{center}

\begin{methodbox}[Faithfulness to a source theory does not mean protecting it from reduction]
This book is faithful to \emph{Third-Order Consciousness} by preserving the problems it actually raises, recording their intellectual provenance, and allowing strict comparison to change their epistemic status. If critical reconstruction eventually narrows the proper name, the formalization can remain: what is being formalized is a set of relations, not a label whose originality must be preserved forever.
\end{methodbox}

\section{\texorpdfstring{How Observation Can Become a New Project of\\Self-Cultivation}{How Observation Can Become a New Project of Self-Cultivation}}

Before introducing ``the observer is the past,'' \emph{Third-Order Consciousness} addresses an important trap: observation itself can be reabsorbed by the second-order subject \citep[TOC-066; pp.~480--484]{toc2026}.

It is easy to hear the third order as a new capacity: before, I did not know how to be aware; now I will practice awareness. Before, I was carried away by emotion; later I will become calmer. Before, I was not mature enough; now I will cultivate a more advanced power of observation. Such wishes are not necessarily harmful. Attention training, emotion regulation, and behavioral control all have practical value in their own domains.

The problem is that once ``observation'' is written from the beginning into a program of self-upgrading, the original subject structure can remain intact. There is a ``not-good-enough me'' in the present, a ``more lucid me'' in the future, and between them a system of training, comparison, accumulation, and scoring. The subject has not lost its central position. It has acquired a new form of capital: ``my awareness,'' ``my growth,'' ``my serenity.''

Following Krishnamurti, \emph{Third-Order Consciousness} makes the point sharply: observation can count as revolutionary not because it is ``deeper'' than ordinary reflection, but because it no longer presupposes an engineer standing over experience and treating the self as a project to be optimized \citep{toc2026}.

\begin{example}[When ``awareness'' itself becomes performance]
A person records every day how long they meditated, rates their emotions, and judges whether they successfully ``maintained observation.'' Such records may be useful. But if anger one day is followed immediately by the judgment, ``My awareness was terrible today; I have regressed again,'' then the previous performance structure has simply changed indicators. The secular version was ``I was not successful enough today.'' The spiritualized version becomes ``I was not lucid enough today.''
\end{example}

This does not require abandoning training. It requires training to remain available for examination as well. What does it optimize? What does it cost? Has it quietly re-centered the entire process on ``becoming a better me''? Third-Order Consciousness is not opposed to method. It is opposed to allowing method to become a new identity through which the observer perpetuates itself without scrutiny.

\section{Krishnamurti's First Cut: The Observer Is the Past}

Only here does the deeper break begin. Krishnamurti repeatedly argues that the psychological observer is not a transparent pane of glass. It is composed of knowledge, memory, experience, conclusions, hurt, tradition, and expectation; in this important sense, ``the observer is the past'' \citep{krishnamurti1954,krishnamurti1969,kft2026a}. \emph{Third-Order Consciousness} treats this as the threshold of the revolution in observation: the past is no longer merely content being observed; it is itself functioning as the eye through which the present is seen \citep[TOC-067; pp.~485--490]{toc2026}.

Take an ordinary case. Someone sees that a partner has arrived later than agreed and almost immediately thinks, ``He does not take me seriously--again.'' The judgment has not come from nowhere. It may carry several earlier delays, an old argument, memories of being disregarded, a rule about respect, fatigue that day, and an identity narrative saying ``I cannot let people treat me like this anymore.''

Not all of this historical material is wrong. In fact, reasonable judgment depends on the past. Physicians rely on training to recognize symptoms; researchers rely on accumulated knowledge to judge whether a proof is plausible; someone in a relationship is entirely justified in remembering repeated broken promises. The issue is not ``having a past.'' It is this:

\begin{quote}
\textbf{When does the past cease to appear as ``information supplied by the past'' and begin to appear as ``I am simply seeing the facts objectively right now''?}
\end{quote}

Once that shift occurs, the observer acquires unusual power. It need not disclose where it came from before deciding what deserves attention, what is merely an exception, which evidence should be amplified, and which counterexamples can be ignored.

This is why ``the observer is the past'' must not be crudely translated as ``memory is bad.'' Krishnamurti is not asking people to become amnesic and is not opposed to professional knowledge. The more precise structural question is whether memory has usurped sovereignty over observation.

The distinction can be written very simply:
\[
\begin{gathered}
\text{what the present judgment says}\\[-1mm]
\neq\\[-1mm]
\text{the historical position from which the judgment occurs}.
\end{gathered}
\]

Third-Order Consciousness requires the latter to have a chance of becoming visible too.

\section{Words, Images, and Facts: How the Past Enters the Present before We Speak}

``The observer is the past'' does not operate only through conscious recollection. The past often enters the present through words, images of persons, and already formed categories.

Imagine that a relationship has gradually been compressed into one sentence: ``He is unreliable.'' The next time the person is five minutes late, the new fact can be fitted into the old model almost immediately. If the person arrives early, that may be interpreted as ``only an exception.'' The model is not worthless--historical stability should rationally affect prediction. The question is whether the model still permits new material to change it.

Krishnamurti repeatedly warns that words and images can stand between a person and fact; \emph{Third-Order Consciousness} interprets this as the past organizing present perception in advance through naming and image \citep[TOC-067; pp.~486--490]{toc2026}. This connects directly to Chapter 2. Language allows experience to be preserved and compared, but when a label arrives too quickly, a new experience may be taken over by its old name before it has fully shown itself.

A psychological reader can understand this as a question about how prior models affect the encoding of new evidence. No mystical vocabulary is required. The third order does not demand that we ``see the world without models''; that is probably impossible. It asks that model and present material become visible together.

A more operational formulation is therefore not ``never look at a person through the past,'' but something like this:

\begin{quote}
\textbf{On the basis of the last ten interactions, my present level of trust in this relationship is low. At the same time, today's concrete behavior still enters as new material and is, in principle, capable of altering my judgment about the overall pattern.}
\end{quote}

This does not erase history, and it does not give history unconditional final authority.

\section{Krishnamurti's Second Cut: Why ``The Observer Is the Observed'' Removes the Privilege of Internal War}

If the observer is itself composed of the past, the next step is sharper. One of Krishnamurti's most famous and most easily mystified statements is \emph{the observer is the observed}. \emph{Third-Order Consciousness} stresses that this is not, first of all, a lyrical statement about cosmic unity. It is a challenge to a split relation within the subject \citep[TOC-068; pp.~491--495]{toc2026}.

In ordinary language we naturally say:

\begin{quote}
``I am dealing with my anger.''
\end{quote}

This sentence is useful and often functionally appropriate. Yet it implies a picture: on one side there is an emotion that must be handled; on the other side, a relatively independent and rational subject that manages the emotion. Experience has been divided into case and judge.

But if the judge is formed by the same histories of comparison, idealization, shame, injury, and fear, then it cannot declare itself outside the battlefield merely because it speaks in a calmer tone. The voice saying ``I must overcome fear'' may itself be organized by the fear ``I cannot appear weak.'' The voice saying ``I should rise above jealousy'' may still be comparing ``who I am now'' with ``a more advanced me.''

The point of ``the observer is the observed'' is therefore not to abolish discrimination or practical action, but to remove the unexamined transcendence of the internal manager.

A dangerous misunderstanding must be prevented. Krishnamurti's formulation cannot be used to dismiss psychotherapy, cognitive reappraisal, behavioral training, emotion regulation, or risk control. Those practices have their own evidence and domains of application. The book preserves a more basic structural reminder: \textbf{effective regulatory function does not automatically grant the regulator an ontologically ahistorical and unconditional highest position.}

\begin{sourcebox}[The original intellectual claim is stronger than the empirical proposition retained by this manuscript]
In \emph{Third-Order Consciousness}'s interpretation of Krishnamurti, ``the observer is the observed'' is pushed to a stronger judgment: when the psychological division between observer and observed truly ends, the internal conflict previously maintained by that division loses the basis on which it continued \citep[TOC-068--TOC-069; pp.~491--500]{toc2026}. This book must record the strength of that intellectual claim faithfully, but may not write it as though a universal causal law had already been confirmed experimentally.

Accordingly, Part I preserves the original ``division--conflict'' insight at the level of intellectual history. When the idea enters the research language of this project, it is retained for now only as a structural stress test: do not presuppose an ahistorical internal judge, and then investigate empirically under what conditions second-layer self-governance weakens, prolongs, or reactivates first-layer processes. Stronger mechanistic claims remain for independent empirical research.
\end{sourcebox}

\section{Choiceless Awareness: Why Observation Cannot Be Pre-Assigned to ``Who I Want to Become''}

Once the observer loses its natural privilege, a reader almost inevitably asks, ``Then how should I observe?'' This is precisely the moment Krishnamurti treats with greatest suspicion. ``Tell me what to do'' can quickly rebuild a method, a goal, and a future achievement: I cannot observe now, but after training I will; I am not free now, but later I will become freer.

Following Krishnamurti, \emph{Third-Order Consciousness} develops the question as \emph{choiceless awareness}: observation does not begin by serving the result ``I want to become a better version of myself'' \citep[TOC-069; pp.~496--500]{toc2026}. This does not mean that ordinary life can contain no choices, nor that people should not set goals. Two kinds of time must be distinguished clearly.

\begin{quote}
\textbf{Technical time:} ``The paper is due next week, so I need to read two articles today.''\\[1mm]
\textbf{Identity time:} ``Only when I finally become disciplined enough, calm enough, and free enough of jealousy will I count as a truly mature person.''
\end{quote}

The first organizes tasks and is indispensable to life. The second can turn an ideal future self into a center that continually judges the present. Krishnamurti's criticism of \emph{becoming} is directed primarily at this psychological movement: the past shapes an ideal future self, the future self comes back to measure the present self, and the subject perpetuates itself along the timeline ``I am not good enough yet'' \citep{krishnamurti1954,krishnamurti1969}.

``Choiceless'' therefore cannot be understood as numbness, passivity, or cynicism. Danger must still be avoided, harm must still be stopped, commitments can still be kept. What is provisionally withdrawn is a subtler act: the moment experience arises, it is not immediately requisitioned as raw material for a project of self-improvement.

It can be said in very ordinary language:

\begin{quote}
\textbf{First see what is happening. Then see who is in such a hurry to make it become some result.}
\end{quote}

This is also why Krishnamurti repeatedly remains wary of fixed method and authority. The problem is not learning itself. It is that once a road to freedom is standardized in advance, the subject easily reoccupies the center of the road: I follow the steps, I accumulate results, I compare my progress, I finally arrive. Even ``selflessness'' can become the most refined form of ``my achievement.''

Here again, Krishnamurti's own position must remain distinct from the evidential discipline of the book. As an intellectual claim about psychological freedom, his suspicion of ``method as the road to freedom'' is far more radical than the ordinary statement ``methods should also be evaluated''; \emph{Third-Order Consciousness} preserves that sharpness on purpose \citep[TOC-069; pp.~496--500]{toc2026}. This book will not domesticate him into a new training technique, but neither will it infer from his philosophical position that every skill-training, therapy, or behavioral exercise is ineffective. The former belongs to the source idea; the latter requires its own empirical evidence.

\section{Relationship Is the Testing Ground: Why Lucidity in Solitude Is Not Enough}

If the changes above occur only in private introspection, they can easily become another private self-image. \emph{Third-Order Consciousness} therefore places relationship at the center of the revolution of observation. The real test is not ``how peaceful I feel when alone,'' but what happens to old structures when another person refuses, arrives late, does not reply, misunderstands, challenges, or simply stops cooperating with my narrative \citep[TOC-070; pp.~501--506]{toc2026}.

Relationship matters because it activates many layers at once: expectation, dependence, comparison, hurt, possession, the need for confirmation, identity position, and images of other people formed from the past. We think we are facing ``what this person has done now,'' when often we are also facing ``the accumulated version of this person already carried by me.''

A sentence such as ``you always do this'' is rarely a simple description of a present action. It has already compressed multiple past segments into a person-model. Such a model may have predictive value; if someone repeatedly behaves in the same way, ignoring history would be irrational. The third-order question is still not ``forget the past.'' It is to know how much of the present response comes from the fact in front of us and how much from an already formed image and expectation.

This makes relationship a very concrete stress test:

\begin{quote}
\textbf{When reality does not continue according to my expectation, can I see at once the external fact, my first-layer reaction, and the observing position that rapidly begins to explain, assign responsibility, defend identity, and arrange the future?}
\end{quote}

If not, then ``I understand that the observer is the past'' may remain knowledge without altering actual organization.

\subsection{Relationship Is Not a Private Vacuum: The Observer Also Carries Institutional Position}

A sociological restriction must be added. The observer's ``past'' is not only childhood memory. It also includes education, profession, gender, class, family role, and institutional expectation. A professor judging a student as ``lacking independence,'' a manager judging an employee as ``not proactive enough,'' or a family member judging another as ``immature'' may all be carrying standards of the ``ideal member'' supplied by the institution in which the judgment occurs.

Mead, Goffman, Berger and Luckmann, and related social-theoretical traditions remind us that the self is not completed outside society and then inserted into relationship. Roles, recognition, evaluation, and institutional classifications participate in forming the subject itself \citep{mead1934,goffman1959,berger1966}. If the third order is genuinely to make the observer visible, then it cannot trace only private psychological history; it must also see the social position behind the observing position.

This also prevents third-order language from being misused to erase real power differences. Adviser and student, physician and patient, manager and employee are not merely two symmetrical psychological individuals. If one side is constrained by genuine institutional power and is then simply told ``that is only your psychological projection,'' decentering becomes a way of making the less powerful party pay the structural cost alone. This book explicitly rejects that use.

\section{The Value and Limit of the Bohm Dialogues: Resonance Is Not Physics Proving Psychology}

Krishnamurti's long dialogue with David Bohm is attractive to many readers and especially vulnerable to interdisciplinary overreach. Popular accounts sometimes turn it into ``quantum physics proves Eastern wisdom.'' This book explicitly rejects that move. Bohm's achievements in physics do not automatically constitute experimental support for a psychological proposition; that would violate the evidential separation established in the Introduction \citep{translation2026}.

What deserves to be preserved in the dialogues is that both participants are highly sensitive to fragmentation and to the assumption that the observer somehow stands outside the total process \citep{krishnamurtiBohm1985}. For Krishnamurti, the issue is how the psychological observer is composed of the past. In Bohm's broader thought, the concern involves fragmented thinking that mistakes its a partial view for the whole.

At most, we can say that there is a structural resonance in the posture of the problem: both sides are suspicious of an absolutely uninvolved vantage. We cannot go on to say that the two share one mechanism, still less that physics provides a natural-scientific proof of Third-Order Consciousness.

This distinction matters especially here. However elegant the later mathematics becomes, we will still have no right to say ``mathematics has finally proved what Krishnamurti really meant.'' Mathematics can provide a new interface for comparison. It cannot seize interpretive authority over the source text.

\section{Third Order and Freedom: From ``I Have More Choices'' to ``The Chooser Also Has Conditions of Formation''}

The Introduction has already discussed a historical contribution of existentialism: when external essence no longer decides for the person, choice and responsibility become heavier. Krishnamurti places ``choice is freedom'' under pressure again. If the chooser is strongly conditioned by fear, tradition, comparison, past hurt, and social standards, a larger number of options does not automatically amount to psychological freedom \citep{krishnamurti1969}.

Two levels must not be confused. Choice in social and political life is genuinely important; the third order does not use ``choicelessness'' to diminish real autonomy. It raises an additional question: \emph{when I choose, what conditions are organizing the choice?}

Someone may have several jobs available and nevertheless experience one of them automatically as ``not respectable'' because of family expectations. Someone may be free to decide whether to continue a relationship, while automatically equating ``contacting the other person again'' with ``losing dignity.'' In such cases, recording only ``what did the person choose?'' is insufficient to describe the organization of the subject.

We therefore understand freedom in the third-order sense as an increase in the \emph{transparency of conditions}. Action still occurs and responsibility remains, but the agent is less likely to mistake historical conditions for absolute commands coming from outside the structure.

\section{Making the Observer Visible Does Not Mean the Subject Must Disappear}

The argument can now be misread in the opposite direction: if the observer is only part of the structure, is the ideal state one with no observer, no subject, and no judgment? No.

The second-order subject performs many indispensable functions. Writing a paper requires sustained goals; signing a contract requires identity to persist across time; caring for another person requires commitment; refusing harm requires clear judgment. The third order does not zero these functions out. It alters their authority.

A subject can say, ``On the basis of these repeatedly observed behaviors, I judge that this relationship should not continue.'' The third order does not demand that the judgment be withdrawn. It asks that we can also say which facts were used, which parts of the past were brought into the judgment, at what temporal scale, according to which values, and what new material could alter the conclusion.

The subject therefore moves from ``final master'' back to ``a working interface with a position.'' Making the position explicit does not imply relativism. On the contrary, only when a position is marked can a judgment take responsibility more clearly. This is consistent with the Geometry of Existence principle that judgments must have positions \citep{geometry2026}.

\section{Two Complete Cases: Writing the Observer Back into the Same Process}

The preceding concepts can easily become attractive self-knowledge unless they come back to concrete events. We therefore revisit two cases that have appeared repeatedly and ask what changes when ``the observer enters observation.''

\subsection{No Reply, Anger, and Corrective Imagination}

Let us return to the unanswered-message scenario. Imagine that someone takes the initiative and contacts a person they have only recently met, but receives no reply. Before adding any interpretation, let us unfold the sequence in time. The person checks the phone repeatedly; their attention keeps returning to it; they become increasingly agitated; the thought ``this does not meet basic standards of courtesy'' begins to take shape; anger grows; and eventually they decide not to make further contact.

At this point, the action in the outside world may already be over. The person has decided not to make further contact. Yet the thinking continues. In imagination, a future conversation is rehearsed: if she appears again, I will tell her why this behavior is disrespectful; I will make my principle clear; I will force the interaction to reach an explicit conclusion. Each rehearsal brings the anger back.

A little later, another layer may enter: ``Why am I still thinking about this? I should have let it go by now.'' The original anger is now accompanied by self-evaluation and may eventually become another question: ``Am I too easily affected?''

If the whole process is described only as ``me versus my anger,'' these layers are flattened. Once the observer is written back into the process, we can distinguish at least: the external fact, bodily reaction, rule-based judgment, decision about action, corrective imagination, reactivated emotion, and the later layer of self-management.

The third order does not announce that ``the true cause is a need for completion,'' or any other single mechanism. Corrective imagination may involve fairness norms, prediction error, self-respect, attachment, anger habits, and other variables. For now we obtain only a cleaner description: \textbf{the manager that wants to end the cycle can itself become part of the cycle.}

\subsection{An Unfamiliar Talk and ``My Knowledge Is Too Narrow''}

Now let us return to the academic example. In a familiar area, concepts have many accessible connections and questions may come naturally. In a completely unfamiliar area, those connections are sparse, and the immediate experience may simply be ``my mind is blank.'' A few seconds later, however, this local difficulty may trigger a much heavier judgment: ``my knowledge is too narrow.''

The judgment may contain real information. A person's range of reading, repertoire of methods, and background knowledge can certainly be compared objectively. The third order does not comfort itself with the slogan ``that is only an illusion generated by structure.''

It simply unfolds a transition that had been omitted:

\begin{quote}
\textbf{Why did reduced accessibility within this particular talk become so quickly a long-term identity judgment about ``the kind of mathematician I am''?}
\end{quote}

Once asked, more variables reappear: how far is the talk from my area? Is prerequisite knowledge missing? Is the exposition itself clear? Am I tired today? Who am I using as the comparison class? Does my professional identity depend strongly on the standard that ``a good mathematician should quickly generate questions''?

The third order does not add a more positive conclusion. It makes the mechanism by which evaluation is generated part of experience as well.

\section{How Psychology Could Test These Distinctions: Do Not Begin by Manufacturing a ``Third-Order Score''}

If this chapter were reduced to ``Third-Order Consciousness means having no observer,'' psychology could hardly use it. A more useful approach is to decompose the intellectual proposal into processes that can be studied separately, rather than immediately designing a ``Third-Order Consciousness scale.''

First, distinguish present experiential content from the rule by which the experience is being evaluated. Participants might report not only ``I am angry now'' but also ``I am currently evaluating myself by the rule that `a mature person should not be angry.' '' The intensity and time course of these two variables can change independently.

Second, distinguish first-layer reaction from second-layer self-governance. Bodily arousal, emotion, and action tendency can be recorded first; later shame, self-blame, comparison, suppression, or attempts at improvement can be recorded separately. The research question then becomes: does the second layer reduce, prolong, or reactivate the first-layer response?

Third, manipulate self-referential frames. Does the same fact produce different memory selection, event segmentation, attribution, and prediction under prompts such as ``as a researcher,'' ``as a person being evaluated,'' ``as a friend,'' or ``as someone who must preserve dignity''? If so, the observing position begins to have measurable consequences.

Fourth, any new variable must compete with existing constructs for explanatory work. Metacognition, decentering, self-distancing, emotion regulation, and autobiographical memory are already mature research areas. If future work introduces ``observer endogeneity'' as a variable, it must show whether that variable produces incremental prediction beyond what the existing variables explain. Otherwise there is no need to create another construct.

Accordingly, this book proposes empirical bridges but does not write unfinished operationalization as though it were an established psychological mechanism. If future experiments show that the distinctions are not stable, the theory should contract accordingly.

\section{A Refined Definition of Third-Order Consciousness: The Subject Retains Function but Loses Final Authority}

After the first three chapters, the working definition from Chapter 1 can be refined. It remains a research formulation used by this book after reviewing \emph{Third-Order Consciousness}; it is not an established diagnostic standard or a unified construct already accepted by psychology.

\begin{definition}[Third-Order Consciousness: refined Part I definition]
\ordthree{} is a mode of consciousness in which the conditions that organize experience themselves enter observation. Feelings, thoughts, memories, and narratives can still be observed. At the same time, the subject position responsible for naming, comparing, evaluating, attributing, and changing those experiences is understood as a process-position jointly formed by bodily state, memory, language, social rules, relational history, values, and current goals.

The third order does not require the subject to stop acting and does not require the subject to disappear. It removes one default privilege: a position cannot be treated as the final master, standing outside the same process and needing no explanation, merely because that position is currently ``observing.''
\end{definition}

The relation among the three orders can again be compressed into three sentences:

\begin{quote}
\textbf{First order:} experience unfolds through direct body--environment coupling.\\
\textbf{Second order:} experience is organized as ``my experience.''\\
\textbf{Third order:} even ``how my experience is being organized'' enters observation.
\end{quote}

These are not three spiritual ranks but three distinct focal questions.

\section{The Last Gate before Mathematics: What Kind of Formal Language Do We Now Need?}

Part I can now end. Without borrowing later mathematical terminology, we have gradually clarified the problems that require formalization.

We need a language capable of describing: how ongoing occurrence forms comparable segments; how two different moments move from ``somewhat similar'' to ``again''; how an observing position acquires special organizational force within the process itself; how an established historical recognition is preserved, refined, or reinterpreted as more background enters; how a compressed story can conceal differences that still matter for the future; and, above all, how the observing position belongs to the same dynamic process rather than standing outside it.

At the same time, the formal language must obey three disciplines already established in Part I.

\begin{enumerate}
\item \textbf{Krishnamurti's stress test:} the third order cannot be defined by appealing to an observer outside the structure who is never itself explained.
\item \textbf{The discipline of the Geometry of Existence:} every abstraction must be able to lead back to concrete events, an explicit scale of observation, and conditions under which it can be revised.
\item \textbf{The discipline of the translation protocol:} mathematical objects must first stand independently in the mathematical native domain; mathematical definitions cannot be reverse-engineered from psychological intuition.
\end{enumerate}

For this reason, the next Part begins with a move that may look abrupt but is in fact necessary: we temporarily stop talking altogether about ``I,'' anger, subject, and Third-Order Consciousness, and begin again from the most basic mathematical objects. Only if the mathematical language can stand on its own without psychological assistance can the later correspondences become genuine structural comparisons rather than psychological words exchanged for mathematical symbols.

\section{Conclusion of Part I: What Have We Gained?}

Part I has not produced a mechanistic theory of consciousness. It has produced a sufficiently clear map of the problem.

First, we distinguished direct occurrence of experience, subjective organization of experience, and the observer itself entering observation as three different focal questions. Second, we decomposed the second-order subject back into conditions such as body, time, memory, language, others, institutions, and narrative, rather than treating the ``I'' as a single substance. Third, through Krishnamurti's revolution in observation, we made explicit that the third order cannot be rewritten as a higher, purer, more competent super-subject whose task is to manage itself.

What later Parts actually need to formalize is therefore not ``what consciousness is,'' but a family of relations:
\begin{enumerate}
\item how ongoing experience acquires segmental structure;
\item how moments separated in time are linked as ``happening again'';
\item how a local observing position is generated and acquires organizational force;
\item how a recognition grounded in history is preserved, refined, or revealed as a misrecognition as background expands;
\item how narrative compression preserves useful history without erasing differences that still matter for the future;
\item how an observing position is placed back within the same process instead of pretending to stand outside it.
\end{enumerate}

The next Part translates these questions into mathematics for the first time. Before that translation begins, the book again executes its most important methodological discipline: \textbf{let the two native domains stand independently before proposing structural translation between them.}

\part{The Mathematical Language of AF-by-Discrete: From Finite Graphs to Groupoids and Germs}

\begin{center}
\large\itshape First let the mathematical objects stand on their own in their native domain; only then discuss any correspondence with consciousness.
\end{center}

\begin{methodbox}[Note on Part II]
This Part temporarily leaves the contexts of psychology, philosophy, and social theory and establishes only the mathematical objects needed later. The dependency order of the terminology is: graphs, roots, and biroots; Cantor spaces and local homeomorphisms; Bratteli diagrams and path spaces; groupoids, local actions, and germs; the tail groupoid and AF groupoid; graphs with boundary and tile inflation; the strict definition of \AFBD; orbital graphs, Cayley graphs, and graphs of germs; and only then trajectories, traverses, returns, and incompressibility. This Part gives only definitions, basic examples, and necessary background. The formal theorems on persistence, filtration, decomposition coherence, and residue are concentrated in Part IV, so that the same theory is not proved twice.
\end{methodbox}

\begin{plainbox}[Terminology rule for Part II]
The main text retains standard mathematical terms because they will later need to be cited precisely. In ordinary explanations, however, \emph{root} can first be understood as ``a designated reference position''; \emph{germ} as ``local transformation information retaining only how an action behaves near a point''; \emph{witness} as ``a piece of structural evidence that can be checked finitely''; and \emph{resolution} as ``the degree of discrimination currently retained.'' First understand what distinctions these terms make; then learn the full formal constructions.
\end{plainbox}

\clearpage

\chapter{Graphs, Paths, Roots, and Birooted Structures: A Minimal Language of Finite Structure}
\label{chap:graphs-en}

\par\medskip
\noindent\textbf{Why this chapter matters}\par
Mathematics begins here. For the moment, set aside all the psychology from the preceding part. Here we study only the simplest building blocks: vertices, edges, walking along edges, and, within a graph, marking ``start here'' and ``end here.'' The groupoids, tiles, traverses, and returns that may look complicated later will ultimately reduce to these finite objects.
\par\medskip

\section{Why Begin with Finite Graphs?}

All later hierarchies, boundaries, completion relations, and decompositions ultimately have to reduce to vertices, edges, and finite trajectories along edges. Here we establish only the discrete language.

\begin{definition}[Directed graph]
A directed graph $\Gamma$ consists of a vertex set $V(\Gamma)$, an edge set $E(\Gamma)$, and source and range maps
\[
s,r:E(\Gamma)\longrightarrow V(\Gamma).
\]
For an edge $e$ with $s(e)=u$ and $r(e)=v$, write $u\xrightarrow{e}v$. A directed graph equipped in addition with a labeling map $\ell:E(\Gamma)\to S$ is called an $S$-labeled directed graph.
\end{definition}

\begin{definition}[walk, closed walk, and simple path]
A walk of length $m$ is a composable sequence of edges $e_1\cdots e_m$, where $r(e_i)=s(e_{i+1})$. We call such a walk a \emph{closed walk} if its initial and terminal vertices coincide. We call it a \emph{simple path} if no vertex is repeated except possibly the initial/terminal vertex. When this book formally treats trajectories of group actions, it preferentially uses \emph{trajectory} or \emph{walk}.
\end{definition}

\begin{plainbox}
Here ``closed'' means only that the initial vertex equals the terminal vertex. It is not yet the return defined later in this book, because a return additionally requires a finite block with boundary and the condition that the boundary is not met again in between.
\end{plainbox}
\vspace{-8pt}
\section{Roots and Birooted Structures}

\begin{definition}[rooted graph]
A rooted graph is a pair $(\Gamma,o)$, where $o\in V(\Gamma)$ is the root. A rooted isomorphism must send the root to the root.
\end{definition}

\begin{definition}[birooted graph]
A birooted graph is a triple $(\Gamma;o_-,o_+)$, where $o_-$ and $o_+$ are respectively the initial root and the terminal root. We allow $o_-=o_+$; in this case it is called \emph{closed-birooted}. A birooted isomorphism must preserve both roots and their order.
\end{definition}

\begin{center}
\begin{tikzpicture}[>=Latex,node distance=14mm,
  v/.style={circle,draw,inner sep=1.8pt},
  rr/.style={circle,draw,very thick,inner sep=2.2pt}]
\node[rr] (a) {$o_-$};
\node[v,right=of a] (b) {};
\node[v,right=of b] (c) {};
\node[rr,right=of c] (d) {$o_+$};
\draw[->] (a)--(b); \draw[->] (b)--(c); \draw[->] (c)--(d);
\node[below=8mm of b,xshift=7mm] {A birooted graph records an entry and an exit};
\end{tikzpicture}
\end{center}
\vspace{-7pt}
\begin{definition}[occurrence and type]
We call $(Q;x,y)\subseteq\Gamma$ a \emph{birooted occurrence} when a larger graph $\Gamma$ contains a subgraph $Q$ with designated vertices $x,y\in V(Q)$. Forgetting its specific position in $\Gamma$ and retaining only its birooted isomorphism class gives a \emph{birooted type}.
\end{definition}

\begin{principle}[occurrences first]
Hierarchical embeddings, persistence, and decomposition must first be defined on concrete occurrences. A type forgets position, adjacent connectors, and ambient context; it is suitable for classification and counting, but cannot carry all structural information.
\end{principle}
\vspace{-9pt}
\section{Embeddings and Inductive Systems}

\begin{definition}[graph embedding]
A graph embedding $\iota:\Gamma\hookrightarrow\Gamma'$ is an injective graph map preserving source, range, and all declared labels. For rooted objects, a graph embedding must also preserve the roots.
\end{definition}

A compatible sequence of embeddings
{\setlength{\abovedisplayskip}{3pt}\setlength{\belowdisplayskip}{3pt}
\[
\Gamma_1\hookrightarrow\Gamma_2\hookrightarrow\cdots,
\qquad
\iota_{n,k}=\iota_{m,k}\circ\iota_{n,m},
\]
}
allows us to discuss an inductive limit. The infinite tiles introduced later arise from compatible embeddings of finite occurrences of this kind.

\chapter{Topological Preliminaries: Cantor Spaces, Clopen Sets, and Local Homeomorphisms}
\label{chap:topology-en}

\par\medskip
\noindent\textbf{Why this chapter matters}\par
The preceding chapter dealt only with finite discrete objects: whether two vertices are the same vertex, whether edges can be composed, and whether a small graph is genuinely embedded in a larger one. We now need a new language for expressing ``what happens near a point.'' Topology provides exactly this level of resolution. It does not require a numerical distance to be assigned to every pair of points; it specifies instead which sets may serve as neighborhoods, what continuity means, and what it means for a map to be locally invertible. Every later comparison of local actions is built on this language.
\par\medskip

\section{Why Move from Discrete Graphs to Topology?}

Sameness in a finite graph is very rigid: two vertices are either the same or different; two edges are either the very same edge or they are not. In an infinite space, however, we often need more refined questions. Although two points are different, do they share enough finite information? Although a transformation cannot be inverted on the whole space, is it completely invertible on some small region? Although two local actions differ far away, do they agree strictly near the point currently under consideration?

Topology gives mathematical form to these notions of ``nearby,'' ``local,'' and ``continuous change.'' It does not begin by asking how far apart two points are. It asks instead: \emph{which sets around a point count as its neighborhoods, and whether a map sends sufficiently small neighborhoods to other neighborhoods in a stable way.}

\begin{plainbox}[This chapter does only one thing]
Here we make no psychological correspondence, and we do not yet introduce the algebraic objects defined later. We answer only a sequence of basic questions: how to say ``nearby'' precisely; how to tell whether a map preserves local structure; how to partition a space into finite pieces that are both open and topologically well separated; and how to describe a transformation that is defined only on part of the space but is completely invertible there. The next chapter then uses these concepts to equip an infinite path space with a topology.
\end{plainbox}

\section{Open Sets, Neighborhoods, and Topology: First Make ``Nearby'' Precise}

\begin{definition}[topological space]
Let $X$ be a set. A \emph{topology} on $X$ is a family of subsets $\tau\subseteq\mathcal P(X)$ satisfying:
\begin{enumerate}
\item $\varnothing,X\in\tau$;
\item arbitrary unions of sets in $\tau$ remain in $\tau$;
\item finite intersections of sets in $\tau$ remain in $\tau$.
\end{enumerate}
The sets in $\tau$ are called \emph{open sets}, and the pair $(X,\tau)$ is called a topological space.
\end{definition}

\begin{definition}[neighborhood and closed set]
Let $x\in X$. A set $N\subseteq X$ is called a \emph{neighborhood} of $x$ if there exists an open set $U$ such that
\[
x\in U\subseteq N.
\]
A set $F\subseteq X$ is called \emph{closed} if $X\setminus F$ is open.
\end{definition}

A neighborhood need not itself be open; it only has to contain a genuinely open neighborhood. This distinction will be convenient later, because a statement such as ``two maps agree on some neighborhood of $x$'' allows that neighborhood to be shrunk to an open set without changing the local information.

\begin{definition}[basis for a topology]
A family of open sets $\mathcal B$ is called a \emph{basis} for the topology of $X$ if every open set can be written as a union of members of $\mathcal B$. Equivalently, for every open set $U$ and every $x\in U$, there exists $B\in\mathcal B$ such that
\[
x\in B\subseteq U.
\]
\end{definition}

\begin{plainbox}[Why a ``basis'' matters]
A space may have very many open sets, but we usually do not want to list them one by one. If we can find a collection of simple small open sets from which every open set can be assembled, it is enough to understand those basic pieces. The infinite sequence space below has an especially natural kind of basic piece: fix a finite prefix and leave everything afterward unrestricted.
\end{plainbox}

\section{Continuity and Homeomorphism: When Does a Map Preserve Topological Structure?}

\begin{definition}[continuous map]
A map $f:X\to Y$ is called \emph{continuous} if, for every open set $V$ in $Y$, the preimage
\[
f^{-1}(V)=\{x\in X:f(x)\in V\}
\]
is open in $X$.
\end{definition}

\begin{plainbox}[Why define continuity by open preimages?]
Intuitively, continuity means that sufficiently small local changes in the input do not cause the output to jump abruptly to a completely uncontrolled location. Topology has no numerical distance, so an $\varepsilon$--$\delta$ condition is not available in this generality. ``The preimage of every open set is open'' is precisely the coordinate-free version of this local stability.
\end{plainbox}

\begin{definition}[homeomorphism]
We call a bijection $f:X\to Y$ a \emph{homeomorphism} if both $f$ and $f^{-1}$ are continuous, and we write
\[
X\cong Y.
\]
In this case, $X$ and $Y$ have the same topological structure.
\end{definition}

A homeomorphism is much stronger than the mere existence of a bijection. A bijection says only that two sets have the same number of points; a homeomorphism also requires their local neighborhood structures to be preserved in both directions. When later chapters say that two local regions are ``completely the same,'' the intended meaning is usually homeomorphism in this sense, not visual resemblance.

\section{Clopen Sets and Zero-Dimensional Structure: Why Cantor Spaces Are Especially Suited to Discretization}

\begin{definition}[clopen]
A set $C\subseteq X$ is called \emph{clopen} if it is both open and closed.
\end{definition}

\begin{plainbox}[Why ``clopen'' is not a contradiction]
``Open'' and ``closed'' are not mutually exclusive everyday states here. Openness describes the local neighborhood structure; closedness means that the complement is open. The sets $\varnothing$ and $X$ are always both open and closed, and highly disconnected spaces contain many nontrivial clopen sets. Such sets are especially useful as pieces of finite partitions: each piece is an open region in its own right, while remaining topologically well separated from the other pieces.
\end{plainbox}

\begin{definition}[clopen partition]
We call a finite family $\mathcal P=\{C_1,\ldots,C_m\}$ a \emph{finite clopen partition} of $X$ if its members are nonempty clopen sets, are pairwise disjoint, and satisfy
\[
X=C_1\sqcup\cdots\sqcup C_m.
\]
\end{definition}

A finite clopen partition is the model of ``finite resolution'' used repeatedly later. At this resolution, we temporarily distinguish only which $C_i$ contains a point, ignoring finer information inside each block. If each $C_i$ is then subdivided into smaller clopen pieces, the resolution becomes finer.

\begin{definition}[zero-dimensional space]
A topological space is called \emph{zero-dimensional} if it has a basis consisting of clopen sets.
\end{definition}

Zero-dimensional does not mean that the space ``contains nothing.'' It means that arbitrary local information can be approximated by small clopen pieces. For this book, the property is especially important because it connects continuous topology with finite combinatorial structure: local neighborhoods can be cut into finitely manageable clopen regions.

\section{Compactness, No Isolated Points, and Total Disconnectedness}

\begin{definition}[compactness]
A topological space $X$ is called \emph{compact} if every open cover
\[
X=\bigcup_{i\in I}U_i
\]
has a finite subcover
\[
X=U_{i_1}\cup\cdots\cup U_{i_m}.
\]
\end{definition}

\begin{plainbox}[What compactness means here]
Compact does not mean ``small.'' A space can contain uncountably many points and still be compact. The intuition useful later is this: if a family of local open sets already covers the whole space, we never need infinitely many mutually independent local rules merely to describe this particular cover; finitely many of those open sets will still cover the space.
\end{plainbox}

\begin{definition}[isolated point and no isolated points]
A point $x\in X$ is called \emph{isolated} if the singleton $\{x\}$ is open. We say that $X$ has \emph{no isolated points} if no point of $X$ is isolated (it is \emph{perfect} in the present compact metrizable setting).
\end{definition}

If $x$ is isolated, a sufficiently small neighborhood can pick it out by itself. Cantor-type spaces have exactly the opposite character: no matter how small a neighborhood becomes, it still contains other points. Every point can be localized with increasing precision without suddenly becoming an isolated atom at any finite resolution.

\begin{definition}[connected and totally disconnected]
A space is called \emph{connected} if it cannot be written as the union of two nonempty disjoint open sets. A space $X$ is called \emph{totally disconnected} if every connected subset of $X$ contains at most one point.
\end{definition}

In the compact metrizable setting relevant to this book, Cantor spaces are not only totally disconnected but also zero-dimensional, and therefore possess rich clopen bases. Intuitively, they are not ``blocks'' glued together from continuous curves; they are more like limit spaces left behind by infinitely repeated discrete branching.

\section{Cantor Spaces: The Standard Stage for Infinite Discrete Choice}

\begin{definition}[metric and metrizable space]
A \emph{metric} on a set $X$ is a function
\[
d:X\times X\longrightarrow[0,\infty)
\]
satisfying $d(x,y)=0$ if and only if $x=y$, symmetry $d(x,y)=d(y,x)$, and the triangle inequality
\[
d(x,z)\le d(x,y)+d(y,z).
\]
A topological space is called \emph{metrizable} if its topology can be generated by the open balls of some metric,
\[
B_d(x,\varepsilon)=\{y\in X:d(x,y)<\varepsilon\}.
\]
\end{definition}

\begin{plainbox}[Why we need only know that ``some distance exists'']
Metrizability does not require a space to come equipped with one uniquely correct numerical distance. It requires only that at least one metric generate the same open sets and neighborhoods. What later chapters actually use is the topology; the choice of a compatible metric usually serves only to turn ``sharing a longer stretch of finite information'' into a number.
\end{plainbox}

\begin{definition}[Cantor space]
A nonempty, compact, metrizable, totally disconnected space with no isolated points is called a \emph{Cantor space}.
\end{definition}

The classical uniqueness theorem for Cantor spaces says that all such spaces are homeomorphic. Thus, when only the topology matters, the standard binary sequence space
\[
X_2=\{0,1\}^{\mathbb N}
\]
can be used as the most intuitive model.

\subsection{The Binary Sequence Space}

A point
\[
x=x_1x_2x_3\cdots\in X_2
\]
is an infinite sequence of $0$--$1$ choices. If $u=u_1\cdots u_n\in\{0,1\}^n$ is a finite word, define the corresponding \emph{cylinder set}
\[
[u]=\{x\in X_2:x_1=u_1,\ldots,x_n=u_n\}.
\]
In other words, $[u]$ fixes only the first $n$ coordinates and leaves the entire infinite tail free.

\begin{example}[the first three coordinates fixed]
The set
\[
[010]=\{010x_4x_5x_6\cdots:x_i\in\{0,1\}\}
\]
contains all infinite binary sequences beginning with $010$. It is not a single point, but an entire local block in which infinitely many later choices remain available.
\end{example}

These cylinders are all clopen, and all cylinders together form a basis for the topology of $X_2$. Thus the real meaning of ``two sequences are close'' becomes very concrete: they share a long initial prefix.

To express this intuition as a distance, define
\[
d(x,y)=
\begin{cases}
0,&x=y,\\
2^{-N(x,y)},&x\neq y,
\end{cases}
\]
where $N(x,y)$ is the first coordinate at which the two sequences differ. The longer the shared prefix, the smaller the distance. The topology generated by this metric is exactly the cylinder topology.

\begin{plainbox}[``Close'' does not mean numerically close in value]
For example, the sequences $000000\cdots$ and $000001\cdots$ agree in their first five coordinates and are therefore close in this topology; by contrast, $000000\cdots$ and $100000\cdots$ differ at the first coordinate and immediately fall into different basic neighborhoods. Local structure is determined by how many levels of finite prefix agree, not by turning an entire infinite sequence into an ordinary real number and comparing numerical magnitudes.
\end{plainbox}

\subsection{Convergence Means That Longer and Longer Finite Information Stabilizes}

In $X_2$, a sequence of points $x^{(k)}$ converges to $x$ if and only if, for every fixed prefix length $n$, once $k$ is sufficiently large the first $n$ coordinates of $x^{(k)}$ agree with those of $x$.

This gives an important intuition for the hierarchical structures used later: \emph{the stabilization of finite levels does not require that the entire infinite object be known at some finite stage; it requires only that every fixed finite resolution eventually stop changing.}

\section{Local Homeomorphisms: Globally Many-to-One, Yet Locally Completely Invertible}

\begin{definition}[local homeomorphism]
A map $f:X\to Y$ is called a \emph{local homeomorphism} if, for every $x\in X$, there exists an open neighborhood $U$ of $x$ such that
\[
f|_U:U\longrightarrow f(U)
\]
is a homeomorphism and $f(U)$ is open in $Y$.
\end{definition}

A local homeomorphism need not be injective on the whole space. It requires only that every point have a surrounding region on which the map is completely invertible.

\begin{example}[the left shift]
Define on $X_2$
\[
\sigma(x_1x_2x_3\cdots)=x_2x_3x_4\cdots.
\]
The map $\sigma$ is not globally injective because
\[
0x_2x_3\cdots\quad\text{and}\quad 1x_2x_3\cdots
\]
have the same image. But after restricting $\sigma$ to $[0]$ or to $[1]$, it is a homeomorphism onto all of $X_2$. Hence $\sigma$ is a local homeomorphism.
\end{example}

\begin{plainbox}[The point most easily confused]
``Locally invertible'' does not mean ``approximately invertible.'' On every chosen neighborhood, the restricted map must genuinely be an exact homeomorphism. What is allowed to fail is \emph{global injectivity}, not local exactness.
\end{plainbox}

\section{Partial Homeomorphisms: Invertible Actions Defined Only on Part of the Space}

\begin{definition}[partial homeomorphism]
Let $X$ be a topological space. We call a homeomorphism
\[
h:D\longrightarrow R
\]
a \emph{partial homeomorphism} of $X$ if $D,R\subseteq X$ are open. The set $D$ is its domain and $R$ its range.
\end{definition}

The only difference between a partial homeomorphism and an ordinary homeomorphism is that the former has no obligation to act on every point of $X$. It need only be completely invertible on its own domain.

\begin{example}[finite-prefix replacement]
In $X_2$, define
\[
h:[010]\longrightarrow[111],
\qquad
h(010z)=111z,
\]
where $z\in\{0,1\}^{\mathbb N}$ denotes an arbitrary infinite tail. The map $h$ changes only the first three coordinates and preserves the rest of the tail coordinate by coordinate. Its inverse is
\[
h^{-1}(111z)=010z.
\]
Because $[010]$ and $[111]$ are both clopen, $h$ is an especially regular partial homeomorphism.
\end{example}

\subsection{Partial Homeomorphisms Compose Naturally}

Let
\[
h:D_h\to R_h,
\qquad
k:D_k\to R_k
\]
be partial homeomorphisms. Then $k\circ h$ is defined only on those points where $h$ can first be applied and where $h(x)$ then lies in $D_k$:
\[
D_{k\circ h}=\{x\in D_h:h(x)\in D_k\}.
\]
On this domain, $k\circ h$ is again a partial homeomorphism. Every $h$ also comes with an exact inverse map
\[
h^{-1}:R_h\longrightarrow D_h.
\]

This is precisely the basic fact needed later to organize ``locally invertible actions'': an action may be executable only on part of the space, and whether two actions can be composed depends on whether the image of the first falls within the domain of the second.

\section{Local Agreement: Difference Far Away Does Not Prevent Exact Agreement Near the Current Point}

We can now define a relation that will be used repeatedly later, though this chapter does not yet give it any further name.

\begin{definition}[local agreement near a point]
Let $h:D_h\to R_h$ and $k:D_k\to R_k$ be partial homeomorphisms of $X$, and let $x\in D_h\cap D_k$. We say that $h$ and $k$ \emph{agree locally near $x$} if there exists an open neighborhood $U\subseteq D_h\cap D_k$ of $x$ such that
\[
h|_U=k|_U.
\]
\end{definition}

\begin{example}[the same local action, different behavior far away]
Suppose $h$ replaces $010z$ by $111z$ on $[010]$, while $k$ has a larger domain, performs exactly the same replacement on $[010]$, and also has an additional action on another clopen region disjoint from $[010]$. Then for every $x\in[010]$, the maps $h$ and $k$ agree locally near $x$; whether they are the same far away is irrelevant to this local judgment.
\end{example}

This is the central point to retain from the chapter: local information can be independent of global expression. Two global objects can be different and yet induce exactly the same local behavior near a point. The mathematics developed later will treat this local agreement itself as data worth preserving.

\section{Four Terms That Must Be Kept Separate: close, local, clopen, closed}

In Chinese, these words can easily be blurred together through expressions for ``near,'' ``local,'' and ``closed/completed,'' but their mathematical meanings are entirely different.

\begin{center}
\small
\begin{tabularx}{0.96\textwidth}{>{\bfseries}p{2.15cm}Y Y}
\toprule
Term & Meaning in this chapter & Must not be misread as \\
\midrule
close / ``near'' & sharing sufficiently fine neighborhood information in the given topology or metric & two objects being ``approximately the same'' \\
local & a strict property on an actual neighborhood & approximately true, or true with a small error \\
clopen & both open and closed & the everyday contradiction of being ``both open and shut'' \\
closed walk & a finite walk from the preceding chapter whose initial and terminal vertices coincide & a closed subset in the topological sense \\
\bottomrule
\end{tabularx}
\end{center}

\begin{stopbox}[Do not let one source-language form of ``closed/completion'' switch domains unnoticed]
A topological ``closed set,'' a graph-theoretic ``closed walk,'' and the various completion relations that appear later in the book belong to different definitions. Similar word forms cannot replace definitions. Every use must be returned to the formal condition in its current mathematical native domain.
\end{stopbox}

\section{Chapter Summary: From ``Adjacency'' to ``Local Structure''}

The preceding chapter established a finite discrete language; this chapter adds a topological level of resolution. At this point, seven things should be kept firmly in mind:
\begin{enumerate}
\item topology specifies what ``nearby'' means through open sets and neighborhoods;
\item a basis lets us generate all open sets from a small collection of standard local blocks;
\item a homeomorphism preserves local topological structure, not merely the number of points;
\item clopen sets are useful both as local neighborhoods and as pieces of finite partitions;
\item a Cantor space can be observed through increasingly fine clopen blocks and has no isolated points;
\item a local homeomorphism may be globally many-to-one but must be strictly invertible near every point;
\item partial homeomorphisms can be composed, and two such actions can agree exactly near a point even when they differ completely far away.
\end{enumerate}

The next chapter constructs, for the first time, the infinite path space actually used in this book. There, ``fixing a finite prefix'' becomes a basic clopen neighborhood, and the compactness, Cantor property, and language of local structure developed here acquire a concrete hierarchical model.

\chapter{Bratteli Diagrams and Path Spaces: Organizing Finite Levels into an Infinite Object}
\label{chap:bratteli-en}

\par\medskip
\noindent\textbf{Why this chapter matters}\par
The preceding chapter showed how, in a Cantor space, a finite prefix can determine a small neighborhood. We now explain where such a structure comes from. A Bratteli diagram breaks an infinite object into level after level of finite choices: each level has only finitely many vertices and finitely many edges, yet proceeding indefinitely yields a complete infinite path. More importantly, different finite histories are allowed to merge at the same vertex, making a Bratteli diagram better suited than an ordinary rooted tree to situations in which finite histories are numerous while their subsequent organization can be shared.
\par\medskip

\section{Why Do We Need a Graph That Is Layered but Not Necessarily a Tree?}

The first chapter of this Part already introduced ordinary finite graphs. Why introduce another kind of graph? Because later we need to retain two features at once: at every fixed scale there is only finite data, while the scale can continue indefinitely, and different finite histories can merge again at a higher level.

A rooted tree is very good at representing continual branching: once two paths separate, they normally do not merge again. A Bratteli diagram permits something different. Two paths with different prefixes may arrive at the same vertex at a higher level; from that vertex onward, they again have the same possible continuations. This ability to ``merge again after branching'' is the basis of the tail equivalence and finite-level approximation used later.

Bratteli diagrams were originally introduced by Ola Bratteli for inductive limits of finite-dimensional $C^*$-algebras, and later became a standard tool in Cantor dynamics and the theory of equivalence relations \citep{bratteli1972,gps1995,putnam2010}. This chapter uses only the combinatorial and topological parts needed later in the book.

\section{Bratteli Diagrams: Every Level Is Finite, but There Is No Last Level}

\begin{definition}[Bratteli diagram]
A Bratteli diagram
\[
\mathsf B=(V,E,\mathbf s,\mathbf r)
\]
consists of two sequences of finite nonempty sets
\[
V=\bigsqcup_{n\ge 0}V_n,
\qquad
E=\bigsqcup_{n\ge 1}E_n,
\]
together with maps
\[
\mathbf s:E_n\longrightarrow V_{n-1},
\qquad
\mathbf r:E_n\longrightarrow V_n.
\]
We take $V_0=\{v_0\}$ to be a singleton and require every restriction $\mathbf s|_{E_n}$ and $\mathbf r|_{E_n}$ to be surjective. Thus every edge goes only from level $n-1$ to level $n$, and every vertex has at least one incoming edge and at least one edge continuing to the next level.
\end{definition}

\begin{plainbox}[In ordinary language]
Draw the diagram one horizontal level at a time. The set $V_n$ consists of the finitely many ``state positions'' allowed at level $n$, while $E_n$ records the permitted connections from the preceding level to this one. Each level is finite, so the diagram can be drawn completely to any fixed depth; there are infinitely many levels, so the whole object can still carry infinite information.
\end{plainbox}

A ``vertex'' here does not record the complete history of a path. Two different finite paths may end at the same vertex. A vertex is therefore closer to saying ``which kind of position you occupy upon reaching this level,'' while the path records ``how, specifically, you arrived here.''

\begin{center}
\begin{tikzpicture}[x=1.55cm,y=1.15cm,>=Latex,
 v/.style={circle,draw,inner sep=1.8pt}]
\node[v] (v0) at (0,0) {$v_0$};
\node[v] (a1) at (1,0.65) {}; \node[v] (b1) at (1,-0.65) {};
\node[v] (a2) at (2,0.75) {}; \node[v] (b2) at (2,-0.75) {};
\node[v] (a3) at (3,0.55) {}; \node[v] (b3) at (3,-0.55) {};
\foreach \u/\v in {v0/a1,v0/b1,a1/a2,a1/b2,b1/a2,b1/b2,a2/a3,a2/b3,b2/a3,b2/b3}\draw[->] (\u)--(\v);
\node[above] at (0,0.35) {$V_0$};
\node[above] at (1,1.00) {$V_1$};
\node[above] at (2,1.10) {$V_2$};
\node[above] at (3,0.90) {$V_3$};
\end{tikzpicture}
\end{center}

In the diagram, the two different vertices in $V_1$ can both reach the same vertex in $V_2$. This is one of the most important distinctions between a Bratteli diagram and an ordinary rooted tree: \textbf{finite histories may differ while positions at a higher level may merge again.}

\section{Finite Paths: What Is Actually Preserved Between Levels}

\begin{definition}[finite Bratteli path]
Let $0\le n<m$. A finite path from $V_n$ to $V_m$ is an edge sequence
\[
\gamma=e_{n+1}e_{n+2}\cdots e_m
\]
satisfying
\[
\mathbf r(e_k)=\mathbf s(e_{k+1})
\qquad(n+1\le k<m).
\]
Its source and range are denoted by
\[
\mathbf s(\gamma)=\mathbf s(e_{n+1}),
\qquad
\mathbf r(\gamma)=\mathbf r(e_m).
\]
The set of finite paths of length $m$ beginning at the root vertex $v_0$ is denoted by $\Omega_m(\mathsf B)$.
\end{definition}

If $\gamma\in\Omega_n(\mathsf B)$ and $\eta$ is a finite path from $\mathbf r(\gamma)$ to a higher level, then they can be concatenated as $\gamma\eta$. This is exactly the same concatenation of walks as in the first chapter of this Part, except that the level structure now forces the order of composition, so there is no issue of ``walking backward'' or ``skipping a level.''

\begin{example}[binary Bratteli diagram]
Suppose every level has a single vertex and every $E_n$ has two edges, labeled $0$ and $1$. Then a path of length $n$ is simply a binary word
\[
a_1a_2\cdots a_n,
\qquad a_i\in\{0,1\}.
\]
Although all different finite histories end at the same vertex, their edge sequences still distinguish them. This example brings the binary sequence space of the preceding chapter directly into the present one.
\end{example}

\section{The Infinite Path Space: An Infinite Object Determined Level by Level by Finite Prefixes}

\begin{definition}[infinite path space]
The infinite path space $\Omega(\mathsf B)$ of a Bratteli diagram $\mathsf B$ is the set of all infinite composable edge sequences
\[
x=e_1e_2e_3\cdots,
\qquad
\mathbf r(e_n)=\mathbf s(e_{n+1}).
\]
Its prefix of length $n$ is denoted by
\[
x|_n=e_1e_2\cdots e_n\in\Omega_n(\mathsf B).
\]
\end{definition}

No infinite path can ever be written down completely at a fixed finite level, but each finite level supplies a progressively longer prefix. This is the most basic form of approximating an infinite object by finite data.

\begin{definition}[cylinder set]
For $\gamma\in\Omega_n(\mathsf B)$, define the cylinder
\[
Z(\gamma)=\{x\in\Omega(\mathsf B):x|_n=\gamma\}.
\]
Thus $Z(\gamma)$ consists of all infinite paths beginning with $\gamma$.
\end{definition}

For fixed $n$, all cylinders of length $n$ form a finite clopen partition:
\[
\Omega(\mathsf B)=\bigsqcup_{\gamma\in\Omega_n(\mathsf B)}Z(\gamma).
\]
The partition at level $n+1$ refines the one at level $n$, because each $Z(\gamma)$ decomposes into all of its one-edge extensions:
\[
Z(\gamma)=
\bigsqcup_{\substack{e\in E_{n+1}\\ \mathbf s(e)=\mathbf r(\gamma)}}
Z(\gamma e).
\]

\begin{principle}[finite prefixes provide nested resolutions]
A Bratteli path space naturally carries a sequence of increasingly fine finite clopen partitions. Changing the level $n$ does not replace the object under study; it changes the finite resolution at which the same infinite path space is being viewed.
\end{principle}

\section{The Topology of the Path Space: Sharing a Long Prefix Means Being Close}

Taking all cylinders as a basis defines a topology on $\Omega(\mathsf B)$. The preceding chapter defined a topological basis abstractly; here we finally encounter its most important concrete example.

If $x\neq y$, let
\[
N(x,y)=\min\{n\ge1:e_n(x)\neq e_n(y)\},
\]
and define
\[
d(x,y)=2^{-N(x,y)},
\qquad d(x,x)=0.
\]
This metric is compatible with the cylinder topology: the more initial edges two paths share, the smaller their distance.

\begin{proposition}[basic topological properties of the path space]
Under the finiteness and surjectivity assumptions above, $\Omega(\mathsf B)$ is a nonempty, compact, metrizable, totally disconnected space, and the cylinders form a clopen basis. If the path space has no isolated points, then it is a Cantor space.
\end{proposition}

\begin{plainbox}[Why ``no isolated points'' must be stated separately]
A Bratteli diagram does not automatically produce a Cantor space. If, beyond some level, a finite path has a unique continuation, then the corresponding infinite path may eventually be completely locked in by finite information and become an isolated point. Only when genuine branching remains visible at higher levels after every finite prefix does the path space have no isolated points.
\end{plainbox}

In the binary one-vertex example, $\Omega(\mathsf B)$ is exactly the space from the preceding chapter,
\[
\{0,1\}^{\mathbb N},
\]
and $Z(a_1\cdots a_n)$ is the ordinary binary cylinder. The abstract Cantor topology of the preceding chapter is therefore connected directly to the hierarchical diagram of this chapter.

\section{Tail Equivalence: Not ``Looking Similar,'' but Eventually Following Exactly the Same Tail}

Infinite paths may differ at the beginning and then coincide completely thereafter. This relation will be important later, so we define it separately now.

\begin{definition}[cofinal / tail equivalence]
Two infinite paths
\[
x=e_1e_2\cdots,
\qquad
y=f_1f_2\cdots
\]
are called \emph{cofinal}, or \emph{tail equivalent}, written
\[
x\sim_{\mathrm{tail}}y,
\]
if there exists $N$ such that
\[
e_n=f_n
\qquad\text{for all }n>N.
\]
The tail-equivalence class of $x$ is denoted by
\[
\operatorname{Cof}(x)=\{y\in\Omega(\mathsf B):y\sim_{\mathrm{tail}}x\}.
\]
\end{definition}

In the binary one-vertex example, tail equivalence means that two binary sequences differ in only finitely many coordinates. For example,
\[
0010110101\cdots
\qquad\text{and}\qquad
1110110101\cdots
\]
are tail equivalent if they agree completely from the fourth coordinate onward.

\begin{proposition}
Tail equivalence is an equivalence relation on $\Omega(\mathsf B)$.
\end{proposition}

\begin{plainbox}[Why this does not conflict with the cylinder topology]
The cylinder topology asks, ``Looking from the beginning, how close are these two paths now?'' Tail equivalence asks, ``After ignoring a finite beginning, do these two paths eventually coincide exactly?'' One is topological nearness; the other is an equivalence relation. They answer different questions. Two tail-equivalent paths may even differ completely in their early prefixes.
\end{plainbox}

\section{Simplicity: Every Current Region Eventually Reaches Every Region at a Higher Level}

\begin{definition}[simple Bratteli diagram]
A Bratteli diagram $\mathsf B$ is called \emph{simple} if for every $n$ there exists $m>n$ such that, for every
\[
v\in V_n,
\qquad w\in V_m,
\]
there is at least one finite path from $v$ to $w$.
\end{definition}

This condition can be read directly from the tail-equivalence classes. For the purposes of this chapter, we record only the following standard fact; proofs can be found in the classical literature on Bratteli diagrams and Cantor dynamics \citep{gps1995,putnam2010}.

\begin{proposition}[simplicity and dense tail classes]
The diagram $\mathsf B$ is simple if and only if, for every $x\in\Omega(\mathsf B)$, the tail-equivalence class $\operatorname{Cof}(x)$ is dense in $\Omega(\mathsf B)$; that is, every nonempty cylinder contains a path tail equivalent to $x$.
\end{proposition}

\begin{plainbox}[Why this matters]
Given any finite prefix $\gamma$, simplicity guarantees that one can first reach the position specified by $\gamma$ and then, at a higher level, reconnect to some tail of $x$. Thus a tail-equivalence class cannot remain trapped in only one local region of the path space. This background condition will recur whenever we later need finite patterns to be encountered repeatedly throughout the space.
\end{plainbox}

\section{Telescoping: Skipping Intermediate Levels without Changing Infinite Paths}

The numerical labels of the levels are not sacred. Sometimes many consecutive levels carry only repetitive intermediate bookkeeping, and we want to combine them into a larger scale. This is telescoping.

\begin{definition}[telescoping]
Choose a strictly increasing cofinal sequence of levels
\[
0=n_0<n_1<n_2<\cdots.
\]
The vertices at the new level $k$ are $V_{n_k}$. A new edge at level $k$ is no longer a single edge of the original diagram, but an entire finite path in the original diagram from $V_{n_{k-1}}$ to $V_{n_k}$. The resulting Bratteli diagram is denoted by $\mathsf B'$ and is called the telescoping of $\mathsf B$ along $(n_k)$.
\end{definition}

\begin{plainbox}[In ordinary language]
A stretch that originally required passing through five levels can now be packaged as one ``large edge.'' The intermediate levels have not been declared nonexistent; they are simply omitted from the current presentation. Whenever necessary, the large edge can still be expanded back into the original finite path.
\end{plainbox}

\begin{proposition}[telescoping preserves the path space and tail equivalence]
Every infinite path in $\mathsf B$ can be divided into blocks along the intervals
\[
[n_{k-1}+1,n_k]
\]
to obtain an infinite path in $\mathsf B'$. Conversely, expanding each ``large edge'' of $\mathsf B'$ uniquely recovers an infinite path in $\mathsf B$. This correspondence gives a natural homeomorphism
\[
\Theta:\Omega(\mathsf B)\longrightarrow\Omega(\mathsf B'),
\]
and
\[
x\sim_{\mathrm{tail}}y
\quad\Longleftrightarrow\quad
\Theta(x)\sim_{\mathrm{tail}}\Theta(y).
\]
\end{proposition}

\begin{warning}[telescoping is not an arbitrary quotient]
Telescoping merely regroups levels: every new edge still retains a complete old path, so two originally distinct infinite paths are not forcibly glued into a single point. An arbitrary quotient, by contrast, may genuinely identify different paths and discard future information. This distinction must be preserved whenever the book later discusses ``legitimate compression.''
\end{warning}

Telescoping also supplies a discipline that will recur later: \textbf{a specific integer level is a presentation coordinate, not an absolute structure.} If intermediate levels are removed, the same finite phenomenon may change from ``appearing at level 12'' to ``appearing at level 5 in the new diagram.'' Therefore, any conclusion genuinely intended to be independent of presentation cannot depend only on a particular numerical level itself.

\section{Chapter Conclusion: What Has the Bratteli Diagram Added?}

From the ordinary graphs in the first chapter of this Part to the present point, we have added four kinds of structure.

\begin{enumerate}
\item \textbf{Levels.} Every finite scale has its own finite vertex and edge data.
\item \textbf{Infinite paths.} An infinite object is determined jointly by all compatible finite prefixes.
\item \textbf{Nested resolutions.} Cylinders of length $n$ give a finite clopen partition, which becomes progressively finer as $n$ increases.
\item \textbf{Scale compression that can change the presentation.} Telescoping can skip intermediate levels while preserving the infinite path space and tail equivalence.
\end{enumerate}

At this stage we still have only layered graphs, paths, and topology. The next chapter adds a new mathematical language: instead of asking only ``what relation holds between two points,'' we will write the relation itself as a system of arrows that can be composed, inverted, and varied over local neighborhoods.

\chapter{Groupoids, Local Actions, and Germs: From ``Where Can It Move?'' to ``How Does It Act Locally?''}
\label{chap:groupoids-en}

\par\medskip
\noindent\textbf{Why this chapter matters}\par
The preceding chapters described only positions, paths, neighborhoods, and levels. We now add a new kind of object: arrows themselves become mathematical data. An arrow does not merely say that two points are related; it records ``how one moves locally from this point to that point.'' Different arrows can be composed when the matching conditions allow it, and every arrow can also be followed in reverse. The second half of this chapter will define the germ in the title precisely. A germ is not a form of ``vague similarity'': it forgets everything a local transformation does far away and retains only how that transformation actually acts near a specified point.
\par\medskip

\par\medskip
\noindent\textbf{Mathematical sources and notation for this chapter}\par
Here we establish only the mathematical native domain. The basic language of groupoids, étale groupoids, and bisections follows the standard framework of Renault, with the terminology of ample groupoids used commonly in Matui's work; partial homeomorphisms, inverse semigroup actions, and the germ construction follow the relevant treatments of Exel and Nekrashevych \citep{renault1980,matui2012,exel1998,exel2008,nek2022}. Different sources sometimes adopt opposite conventions for multiplication order and action notation. Accordingly, whenever composition is involved, this chapter also states in words which step is performed first and which is performed second.
\par\medskip

\section{Why ``One Global Transformation'' Is Not Enough}

The most familiar picture of a group action is that every group element gives an invertible transformation of the entire space. The tiles, boundaries, and local continuations used later will not always be so tidy. An action may be defined only on part of the space; two local actions can be continued one after the other only when the endpoint of the first lies in the domain of the second; and there may be different modes of local action between the same source and range points.

It is therefore not enough to record only that ``$x$ is related to $y$.'' We also want to retain an arrow from $x$ to $y$ and be able to answer:
\begin{itemize}
\item where the arrow starts and where it ends;
\item which two arrows may be composed in succession;
\item which arrow results from that composition;
\item how to reverse direction along the same local relation;
\item and, in a topological space, how a compatible family of arrows can be read as a locally invertible map.
\end{itemize}
This is exactly the information that groupoid language is designed to preserve.

\section{Groupoids: A System of Invertible Arrows That Can Be Composed Locally}

\begin{definition}[groupoid]
A groupoid $\mathcal G$ consists of an arrow set $\mathcal G$, a unit set $\mathcal G^{(0)}$, source and range maps
\[
\mathbf s,\mathbf r:\mathcal G\longrightarrow\mathcal G^{(0)},
\]
and partially defined multiplication and inversion. The set of composable pairs is
\[
\mathcal G^{(2)}
=
\{(g,h)\in\mathcal G\times\mathcal G:\mathbf s(g)=\mathbf r(h)\}.
\]
For $(g,h)\in\mathcal G^{(2)}$, their product is written $gh$ and means: first follow $h$, then follow $g$. These operations satisfy:
\begin{enumerate}
\item multiplication is associative whenever the products involved are defined;
\item every $x\in\mathcal G^{(0)}$ has a unit arrow $1_x:x\to x$;
\item every $g\in\mathcal G$ has an inverse arrow $g^{-1}$ satisfying
\[
g^{-1}g=1_{\mathbf s(g)},
\qquad
gg^{-1}=1_{\mathbf r(g)}.
\]
\end{enumerate}
\end{definition}

\begin{plainbox}[In ordinary language]
Write an arrow $g$ as
\[
\mathbf s(g)\xrightarrow{\ g\ }\mathbf r(g).
\]
The central point of a groupoid is not that ``every action can be multiplied.'' Only arrows whose endpoints match can be composed. If
\[
x\xrightarrow{\ h\ }y\xrightarrow{\ g\ }z,
\]
then and only then can we form
\[
x\xrightarrow{\ gh\ }z.
\]
Every step can also be traversed in reverse. A groupoid therefore retains both ``positions'' and ``locally invertible transformations.''
\end{plainbox}

\begin{example}[a group is a groupoid with one unit]
If $G$ is a group, take a single unit point $*$ and regard every $g\in G$ as an arrow $*\to *$. Then all arrows are composable, and groupoid multiplication is exactly the original group multiplication. A group is therefore a special case of a groupoid.
\end{example}

\begin{example}[equivalence-relation groupoid]
Let $R\subseteq X\times X$ be an equivalence relation. Regard $(x,y)\in R$ as an arrow from $y$ to $x$:
\[
\mathbf s(x,y)=y,
\qquad
\mathbf r(x,y)=x.
\]
If $(x,y),(y,z)\in R$, then
\[
(x,y)(y,z)=(x,z),
\qquad
(x,y)^{-1}=(y,x).
\]
This example records only which points belong to the same equivalence class; there is at most one arrow between each ordered pair of related points.
\end{example}

\begin{example}[transformation groupoid of a discrete group action]
Let a discrete group $G$ act on a topological space $X$ by homeomorphisms. Define
\[
G\ltimes X=\{(g,x):g\in G,\ x\in X\},
\]
and regard $(g,x)$ as the arrow
\[
x\longrightarrow gx.
\]
Thus
\[
\mathbf s(g,x)=x,
\qquad
\mathbf r(g,x)=gx,
\]
and two successive steps
\[
x\xrightarrow{(g,x)}gx
\xrightarrow{(h,gx)}hgx
\]
compose to $(hg,x)$. Here the arrow records not only its starting and ending points, but also which group element performs the step.
\end{example}

\section{Orbits, Isotropy, Principality, and Minimality}

Once a groupoid has arrows, it can distinguish ``where can one reach?'' from ``which nontrivial loops remain at the same point?'' These two questions will recur throughout the book.

\begin{definition}[orbit and isotropy group]
For $x\in\mathcal G^{(0)}$, define its orbit by
\[
\operatorname{Orb}_{\mathcal G}(x)
=
\mathbf r\bigl(\mathbf s^{-1}(x)\bigr),
\]
the set of all unit points reachable by arrows starting at $x$. The isotropy group at $x$ is
\[
\mathcal G_x^x
=
\{g\in\mathcal G:\mathbf s(g)=\mathbf r(g)=x\}.
\]
\end{definition}

\begin{definition}[principality and minimality]
A groupoid $\mathcal G$ is called \emph{principal} if, for every $x\in\mathcal G^{(0)}$, the isotropy group $\mathcal G_x^x$ contains only the unit arrow. A topological groupoid $\mathcal G$ is called \emph{minimal} if every orbit $\operatorname{Orb}_{\mathcal G}(x)$ is dense in the unit space.
\end{definition}

\begin{plainbox}[In ordinary language]
Principality excludes nontrivial local symmetry that stays at the same point. Minimality says that no nonempty open region can remain permanently isolated from the rest of an orbit world. These are completely different properties. A groupoid can be minimal without being principal, or principal without being minimal.
\end{plainbox}

For a transformation groupoid $G\ltimes X$, the isotropy at $x$ is the usual stabilizer
\[
G_x=\{g\in G:gx=x\}.
\]
For an equivalence-relation groupoid, every isotropy group is automatically trivial, so such a groupoid is automatically principal.

\section{Topological Groupoids and the Étale Condition}

So far a groupoid has only been a set-theoretic object. If its unit space is a Cantor space, later we also need to know how an arrow varies over ``nearby points.'' We therefore place a topology on the arrow set itself.

\begin{definition}[topological groupoid]
A groupoid $\mathcal G$ is called a \emph{topological groupoid} if $\mathcal G$ is a topological space and the following maps are continuous:
\begin{itemize}
\item multiplication $\mathcal G^{(2)}\to\mathcal G$;
\item inversion $g\mapsto g^{-1}$;
\item the source and range maps $\mathbf s,\mathbf r$.
\end{itemize}
The unit space $\mathcal G^{(0)}$ is given its natural subspace topology.
\end{definition}

\begin{definition}[étale groupoid]
A topological groupoid $\mathcal G$ is called \emph{étale} if the source map
\[
\mathbf s:\mathcal G\longrightarrow\mathcal G^{(0)}
\]
is a local homeomorphism. Since inversion exchanges source and range, the range map $\mathbf r$ is then also a local homeomorphism.
\end{definition}

\section{Bisections: Reading a Family of Arrows as a Local Function}

\begin{definition}[bisection]
A subset $B\subseteq\mathcal G$ is called a \emph{bisection} if both restrictions
\[
\mathbf s|_B,
\qquad
\mathbf r|_B
\]
are injective. We call a bisection $B$ an \emph{open bisection} if $B$ is open.
\end{definition}

\begin{lemma}
Let $\mathcal G$ be étale and let $B\subseteq\mathcal G$ be an open bisection. Then
\[
\mathbf s|_B:B\longrightarrow\mathbf s(B),
\qquad
\mathbf r|_B:B\longrightarrow\mathbf r(B)
\]
are homeomorphisms. Consequently $B$ induces a partial homeomorphism
\[
\theta_B
=
\mathbf r\circ(\mathbf s|_B)^{-1}
:
\mathbf s(B)\longrightarrow\mathbf r(B).
\]
\end{lemma}

\begin{proof}
Because $\mathbf s$ and $\mathbf r$ are local homeomorphisms, they are open maps. After restriction to a bisection they are also injective, and hence become continuous open bijections onto their images. They are therefore homeomorphisms. The displayed map $\theta_B$ is consequently a homeomorphism between two open subsets of the unit space.
\end{proof}

If $B,C$ are open bisections, then the inverse set
\[
B^{-1}=\{g^{-1}:g\in B\}
\]
and the composable product
\[
BC
=
\{gh:g\in B,\ h\in C,\ \mathbf s(g)=\mathbf r(h)\}
\]
are again open bisections. The corresponding local maps are precisely the corresponding inverse and composition.

\begin{definition}[ample groupoid]
An étale groupoid is called an \emph{ample groupoid} if its topology has a basis consisting of compact-open bisections \citep{renault1980,matui2012}.
\end{definition}

\begin{plainbox}[In ordinary language]
For a Cantor-type unit space, a compact-open bisection is especially convenient finite local data: its domain and range can be cut out exactly by clopen pieces, while on each piece there is a completely reversible local rule. The AF groupoids introduced later work precisely in this discrete and topologically stable language.
\end{plainbox}

\section{Why Partial Homeomorphisms Form an Inverse Semigroup}

The preceding chapter already defined a partial homeomorphism. We now collect all such locally invertible maps into a single algebraic object.

\begin{definition}[partial-homeomorphism inverse semigroup]
Let $\mathcal I(X)$ denote the set of all partial homeomorphisms between open subsets of $X$. For
\[
\psi:D_\psi\longrightarrow R_\psi,
\qquad
\phi:D_\phi\longrightarrow R_\phi,
\]
the composition $\phi\circ\psi$ is defined on
\[
D_\psi\cap\psi^{-1}(D_\phi).
\]
The inverse of every $\phi$ is the usual homeomorphism
\[
\phi^{-1}:R_\phi\longrightarrow D_\phi.
\]
Under this composition, $\mathcal I(X)$ is an inverse semigroup.
\end{definition}

\begin{definition}[inverse semigroup]
A semigroup $S$ is called an \emph{inverse semigroup} if every $s\in S$ has a unique element $s^*$ satisfying
\[
ss^*s=s,
\qquad
s^*ss^*=s^*.
\]
\end{definition}

\begin{definition}[idempotent and inverse semigroup action]
An element $e\in S$ is called \emph{idempotent} if $e^2=e$. In $\mathcal I(X)$, the idempotents are precisely the partial identity maps $\operatorname{id}_U$ on open subsets $U$. An action of an inverse semigroup $S$ on $X$ is a map
\[
S\longrightarrow\mathcal I(X)
\]
that preserves multiplication and inverses, with the partial domains covering the unit space under consideration.
\end{definition}

\begin{principle}[two equivalent viewpoints on local actions]
On the étale-groupoid side, open bisections form an inverse semigroup. Conversely, an inverse semigroup of partial homeomorphisms can generate a local arrow system. The first viewpoint starts from ``arrows'' and obtains ``local functions''; the second starts from ``local functions'' and reconstructs ``arrows.'' The germ construction is the key step in this second direction.
\end{principle}

\section{Germ: What Actually Happens Near the Same Point}

Before defining a germ, three different statements must be kept separate:
\begin{enumerate}
\item $\phi(x)=\psi(x)$: the two transformations merely take the same value at this one point;
\item $\phi$ and $\psi$ are ``very close'' at many nearby points: this is a metric or approximation question;
\item $\phi$ and $\psi$ are \emph{exactly equal} on a genuine open neighborhood of $x$: this is the condition required for a germ.
\end{enumerate}

\begin{example}[the same value at one point does not imply the same germ]
On $\mathbb R$, let
\[
\phi(t)=t,
\qquad
\psi(t)=t+t^3.
\]
Both are homeomorphisms and $\phi(0)=\psi(0)=0$. But $\phi$ and $\psi$ are not equal on any open neighborhood of $0$. Thus ``sending the same point to the same point'' is far weaker than ``having the same local action.''
\end{example}

\begin{definition}[germ equivalence]
Let $\mathscr S\subseteq\mathcal I(X)$ be closed under composition and inversion, and let $x$ lie in the domains of both $\phi$ and $\psi$. We say that $(\phi,x)$ and $(\psi,x)$ define the same germ if there exists an open neighborhood
\[
x\in U\subseteq D_\phi\cap D_\psi
\]
such that
\[
\phi|_U=\psi|_U.
\]
The corresponding equivalence class is denoted by $[\phi,x]$.
\end{definition}

\begin{plainbox}[In ordinary language]
A germ discards differences far away. Two transformations may be very different across the whole space, but if the local routes they prescribe agree exactly near the present point, then they represent the same germ there. A germ is therefore an \emph{exact local equivalence class}, not a similarity relation.
\end{plainbox}

For an abstract inverse semigroup, the same relation can be written algebraically: on an appropriate idempotent neighborhood $e$ one requires $se=te$. When the partial homeomorphisms are closed under restriction to open subsets, this is equivalent to exact agreement on some open neighborhood. This book mainly uses the latter description because it keeps the local meaning transparent.

\section{Groupoid of Germs: Reassembling Local Equivalence Classes as Arrows}

\begin{definition}[groupoid of germs]
The set of all germs $[\phi,x]$, with $x\in D_\phi$, forms the groupoid of germs
\[
\operatorname{Germ}(\mathscr S\curvearrowright X).
\]
Its source and range maps are
\[
\mathbf s([\phi,x])=x,
\qquad
\mathbf r([\phi,x])=\phi(x).
\]
Whenever $\phi(x)\in D_\psi$,
\[
[\psi,\phi(x)]\,[\phi,x]
=
[\psi\circ\phi,x],
\]
meaning that one first applies $\phi$ and then $\psi$. The inverse is
\[
[\phi,x]^{-1}
=
[\phi^{-1},\phi(x)].
\]
\end{definition}

\begin{definition}[germ topology]
For $\phi\in\mathscr S$ and an open subset $U\subseteq D_\phi$, set
\[
[\phi,U]
=
\{[\phi,x]:x\in U\}.
\]
These sets form a basis for the topology on the groupoid of germs.
\end{definition}

\begin{lemma}
The groupoid of germs $\operatorname{Germ}(\mathscr S\curvearrowright X)$ is étale.
\end{lemma}

\begin{proof}
Every basic set $[\phi,U]$ is an open bisection: the source map sends $[\phi,x]$ to $x$, hence restricts to a bijection
\[
\mathbf s|_{[\phi,U]}:[\phi,U]\longrightarrow U,
\]
and by the definition of the germ topology this map is a homeomorphism. Thus the source map is a local homeomorphism at every germ, which is precisely the étale condition.
\end{proof}

\section{Transformation Groupoid versus Groupoid of Germs: Two Levels of Bookkeeping}

If a discrete group $G$ acts on $X$, there is a natural map
\[
G\ltimes X
\longrightarrow
\operatorname{Germ}(G\curvearrowright X),
\qquad
(g,x)\longmapsto[g,x].
\]
The groupoid of germs on the right identifies group elements that act identically on some neighborhood of $x$.

\begin{example}[a globally nontrivial element can have a trivial germ]
Let $g\neq e$ be a homeomorphism of $X$ whose support lies away from $x$. Then there is an open neighborhood $U$ of $x$ such that
\[
g|_U=\operatorname{id}_U.
\]
In the transformation groupoid, $(g,x)$ and $(e,x)$ remain arrows carrying different group-element labels; but in the groupoid of germs,
\[
[g,x]=[e,x].
\]
The germ quotient forgets the global name here and retains only the local action.
\end{example}

\begin{warning}[``forgetting what happens far away'' does not mean ``forgetting everything'']
The germ quotient removes global differences that have no effect near the present point. If two transformations both fix $x$ but act differently on every neighborhood of $x$, their germs are still different. The singular germs introduced later depend precisely on local differences of this kind that cannot be removed on any neighborhood.
\end{warning}

\section{Group of Germs and Singular Points}

\begin{definition}[germ isotropy and group of germs]
In a groupoid of germs, the isotropy group at $x$,
\[
H_x
=
\operatorname{Germ}(\mathscr S\curvearrowright X)_x^x
=
\{[\phi,x]:\phi(x)=x\},
\]
is called the \emph{group of germs} at $x$.
\end{definition}

If a group $G$ acts on $X$ by homeomorphisms, define the stabilizer
\[
G_x=\{g\in G:gx=x\}
\]
and the neighborhood stabilizer
\[
G_{(x)}
=
\left\{
 g\in G_x:
 \begin{gathered}
 \text{there exists an open neighborhood }U\text{ of }x\\
 \text{such that }g|_U=\operatorname{id}_U
 \end{gathered}
\right\}.
\]
The subgroup $G_{(x)}$ is normal in $G_x$, and the kernel of the germ map $G_x\to H_x$ is exactly $G_{(x)}$. Hence
\[
\boxed{
H_x\cong G_x/G_{(x)}.
}
\]

\begin{plainbox}[In ordinary language]
The stabilizer asks only whether an action ``still lands at $x$.'' The group of germs further declares trivial every stabilizer element that in fact does nothing near $x$. Thus $H_x$ measures how much genuine local symmetry remains after fixing $x$ and removing everything that can be made locally invisible on a neighborhood.
\end{plainbox}

\begin{definition}[regular and singular point]
We call $x$ a \emph{regular point} if $H_x$ is trivial, and a \emph{singular point} if $H_x$ is nontrivial.
\end{definition}

\begin{plainbox}[A fixed point is not automatically singular]
A nonidentity transformation may fix $x$ and nevertheless have the identity germ there, provided it is the identity on some neighborhood of $x$. Singularity requires a nontrivial local action that fixes $x$ and cannot be eliminated by passing to any sufficiently small neighborhood.
\end{plainbox}

\begin{stopbox}[Terminology stop line]
In this book, \emph{singular} is a strict groupoid/germ term: it means nontrivial germ isotropy. It carries no judgment that the point is ``bad,'' ``defective,'' or ``pathological.'' Pure non-Hausdorffness, introduced below, adds a further topological condition and must not be conflated with singularity itself.
\end{stopbox}

\begin{plainbox}[Boundary and singularity are two different axes]
The tile language introduced later will define a \emph{boundary point}. A boundary point can still be regular: it may participate in a generator arrow joining different tile pieces while its group of germs $H_x$ remains trivial. A \emph{singular boundary germ} means specifically that a persistent boundary point also carries nontrivial germ isotropy. Boundary records that ``cross-piece joining is required''; singularity records that ``at the same orbit position, germ information remains that is not removed by passing to the orbital quotient.'' The two notions are never interchangeable.
\end{plainbox}

\section{After Singularity, Ask Again: Can This Difference Be Separated in a Hausdorff Way?}

The condition $H_x\neq\{1\}$ says only that the same orbit position retains nontrivial germ information. A separate question remains: once this germ difference has appeared, can it occupy a local position near $x$ that is cleanly separated from the identity germ? The following terminology follows the singularity framework used by Nekrashevych \citep{nek2018}.

For $g\in G$, write
\[
\operatorname{Fix}^{\circ}(g)
=
\operatorname{Int}\{y\in X:gy=y\}.
\]

\begin{definition}[Hausdorff, non-Hausdorff, and purely non-Hausdorff singularities]
Let $x$ be a singular point.
\begin{enumerate}
\item The point $x$ is a \emph{Hausdorff singularity} if, for every
\[
g\in G_x\setminus G_{(x)},
\]
one has
\[
x\notin\overline{\operatorname{Fix}^{\circ}(g)}.
\]
\item The point $x$ is \emph{non-Hausdorff} if the preceding condition fails; equivalently, there exists a nontrivial germ represented by some $g\in G_x\setminus G_{(x)}$ such that
\[
x\in\overline{\operatorname{Fix}^{\circ}(g)}.
\]
\item The point $x$ is \emph{purely non-Hausdorff} if every $g\in G_x$ satisfies
\[
x\in\overline{\operatorname{Fix}^{\circ}(g)}.
\]
For the identity element this condition is automatic; the substantive requirement is that every nontrivial germ direction can be approximated arbitrarily closely to $x$ by open regions on which the representing action is locally the identity.
\end{enumerate}
\end{definition}

\begin{plainbox}[What the three cases distinguish]
\begin{itemize}
\item \textbf{Hausdorff singularity:} a genuine germ difference exists, but it remains stably distinguishable from the identity in sufficiently small neighborhoods; locally, the difference can be separated off.
\item \textbf{Non-Hausdorff singularity:} at least one nontrivial germ differs from the identity at $x$, yet arbitrarily near $x$ there are open regions on which the action becomes locally identical to the identity; the genuine difference has no completely separate local territory.
\item \textbf{Purely non-Hausdorff singularity:} this inseparability occurs in every germ direction, rather than arising accidentally for only one of them.
\end{itemize}
\end{plainbox}

Thus
\[
\boxed{
\begin{gathered}
\text{singular}
=
\text{Hausdorff singular}
\;\sqcup\;
\text{non-Hausdorff singular},\\[1mm]
\text{purely non-Hausdorff}
\subseteq
\text{non-Hausdorff singular}.
\end{gathered}
}
\]

The new question is therefore on another axis: \emph{to what extent can an already existing local germ difference be isolated again as a locally separable object?} This question is distinct from asking whether nontrivial germ isotropy exists in the first place.

\begin{warning}[an explanatory axis is not a set-theoretic inclusion chain]
Later chapters may place ``boundary-free,'' ``regular boundary,'' ``Hausdorff singular,'' ``non-Hausdorff singular,'' and ``purely non-Hausdorff'' along a single explanatory axis. That is not a chain of set inclusions. Regular boundary belongs to a different boundary/singularity axis, while purely non-Hausdorff singularities form a strong subclass of non-Hausdorff singularities. These categories must not be collapsed into one linear hierarchy.
\end{warning}

\section{Chapter Conclusion: From Global Names to Local Action}

We have now added eight distinctions.
\begin{enumerate}
\item A \textbf{groupoid} upgrades a relation to arrows carrying source, range, composition, and inversion.
\item An \textbf{orbit} records which units are reachable; \textbf{isotropy} records local symmetries based at individual units.
\item The \textbf{étale} condition guarantees that arrows vary locally and continuously with their source points.
\item An \textbf{open bisection} can be read as a partial homeomorphism.
\item Partial homeomorphisms naturally form an \textbf{inverse semigroup}.
\item \textbf{Germ equivalence} retains only the actual local action near the designated point.
\item The \textbf{groupoid of germs} reassembles these local equivalence classes into an étale groupoid.
\item The \textbf{group of germs} $H_x$ measures local symmetries fixing $x$ that cannot be removed on a neighborhood.
\end{enumerate}

The two opposite channels can be summarized as
\[
\boxed{
\begin{gathered}
\text{open bisections of an étale groupoid}
\longrightarrow
\text{partial homeomorphisms},\\[1mm]
\text{local partial homeomorphisms}
\longrightarrow
\text{germs}
\longrightarrow
\text{groupoid of germs}.
\end{gathered}
}
\]

The next chapter applies this language to a Bratteli path space. Tail equivalence will become a concrete étale groupoid, after which we can state precisely what ``approximately finite'' means in an AF groupoid.

\chapter{Tail Groupoid and AF Groupoid: An Infinite Union of Finite Approximations}
\label{chap:tail-af-en}

\par\medskip
\noindent\textbf{Why this chapter matters}\par
The preceding chapter already gave us Bratteli diagrams, infinite paths, and tail equivalence, but so far ``tail equivalence'' is only a set-theoretic relation. We now upgrade it to a genuine groupoid: we retain not only which paths eventually have the same tail, but also the arrows from one path to another, the local topology on those arrows, and the successive finite stages inside a single object. Here, for the first time, \emph{approximately finite} in AF acquires its strict meaning. It does not mean ``more or less finite.'' It means that an infinite groupoid is \emph{exactly} the union of an increasing sequence of elementary subgroupoids.
\par\medskip

\par\medskip
\noindent\textbf{Mathematical sources and scope of this chapter}\par
The standard background on Bratteli tail equivalence, AF equivalence relations, and Cantor dynamics follows Giordano--Putnam--Skau and Putnam; the terminology of elementary and AF groupoids is used in Matui's Cantor-groupoid setting \citep{gps1995,putnam2010,matui2012}. This chapter remains entirely within the mathematical native domain: the tail groupoid, its compact-open bisection basis, finite-stage subgroupoids, the AF filtration, minimality, and presentation dependence. Only the next chapter introduces interface data describing how finite pieces may continue outward; no terminology defined there is used in advance here.
\par\medskip

\section{From Tail Equivalence to the Tail Groupoid: Why Retain the Arrows?}

The preceding chapter defined
\[
x\sim_{\mathrm{tail}} y
\]
to mean that two infinite paths agree edge by edge from some level onward. If we treat this only as an equivalence relation, we know which points belong to the same cofinality class, but we have not yet retained ``from $y$ to $x$'' itself as an object. The first thing a groupoid does is make this directional information explicit.

\begin{definition}[Bratteli tail groupoid]
Let $\mathsf B$ be a Bratteli diagram with path space $\Omega(\mathsf B)$. Define
\[
\mathfrak T_{\mathsf B}
=
\{(x,y)\in\Omega(\mathsf B)\times\Omega(\mathsf B):x\sim_{\mathrm{tail}}y\}.
\]
Regard $(x,y)$ as an arrow from $y$ to $x$, and set
\[
\mathbf s(x,y)=y,
\qquad
\mathbf r(x,y)=x,
\]
\[
(x,y)^{-1}=(y,x),
\qquad
(x,y)(y,z)=(x,z).
\]
The unit $x\in\Omega(\mathsf B)$ is identified with $(x,x)$.
\end{definition}

This is simply the equivalence-relation groupoid of the tail-equivalence relation from the preceding chapter. Two properties are therefore immediate.

\begin{proposition}[orbits and isotropy]
For every $x\in\Omega(\mathsf B)$, the orbit of $x$ under $\mathfrak T_{\mathsf B}$ is exactly its cofinality class:
\[
\operatorname{Orb}_{\mathfrak T_{\mathsf B}}(x)
=
\operatorname{Cof}(x).
\]
Moreover, $\mathfrak T_{\mathsf B}$ is principal: if an arrow has both source and range equal to $x$, then it can only be the unit arrow $(x,x)$.
\end{proposition}

\begin{plainbox}[In ordinary language]
There is no additional isotropy hidden at a point. The tail groupoid records only which distinct paths eventually become cofinal. Once both source and range are the same $x$, an equivalence-relation groupoid has no second, distinct arrow $x\to x$. If nontrivial isotropy appears later in a larger groupoid, it must come from local arrow structure added outside the AF part.
\end{plainbox}

\section{Basic Bisections: Replacing a Finite Prefix on a Cylinder}

A set-theoretic groupoid is not yet enough. To obtain an \(\acute{e}\)tale groupoid, we must put a topology on the arrow space. The most natural local operation is this: preserve the common infinite tail and replace only a finite prefix.

For $v\in V_n$, write
\[
\Omega_n(v)
=
\{\alpha\in\Omega_n(\mathsf B):\mathbf r(\alpha)=v\}.
\]
If $\alpha,\beta\in\Omega_n(v)$, define
\[
Z(\alpha,\beta)
=
\{(\alpha z,\beta z): z \text{ is a compatible infinite tail beginning at }v\}.
\]

\begin{plainbox}[What does this set actually do?]
Fix two finite paths $\alpha$ and $\beta$ of the same length and with the same terminal vertex. No matter which compatible infinite tail $z$ is chosen, the arrow
\[
(\alpha z,\beta z)
\]
does exactly the same thing: it replaces the initial prefix $\beta$ by $\alpha$ and leaves the tail $z$ completely unchanged. Thus an entire set $Z(\alpha,\beta)$ is not a scattering of unrelated arrows; it is one uniform local rule acting on a cylinder.
\end{plainbox}

Its source and range are precisely the two cylinders
\[
\mathbf s\bigl(Z(\alpha,\beta)\bigr)=Z(\beta),
\qquad
\mathbf r\bigl(Z(\alpha,\beta)\bigr)=Z(\alpha),
\]
and the restrictions of both source and range to $Z(\alpha,\beta)$ are homeomorphisms. Hence $Z(\alpha,\beta)$ is a bisection.

The inverse and the most basic composition rule are equally transparent:
\[
Z(\alpha,\beta)^{-1}=Z(\beta,\alpha),
\]
and, at the same level when the middle prefix matches exactly,
\[
Z(\alpha,\beta)Z(\beta,\gamma)=Z(\alpha,\gamma).
\]

\begin{definition}[tail groupoid topology]
Equip $\mathfrak T_{\mathsf B}$ with the topology having all sets $Z(\alpha,\beta)$ as a basis, where $\alpha$ and $\beta$ have the same length and terminate at the same vertex. This is the standard \(\acute{e}\)tale topology on the Bratteli tail groupoid.

Each $Z(\alpha,\beta)$ is homeomorphic to the corresponding infinite-tail space and is therefore a compact-open bisection.
\end{definition}

\begin{proposition}[the tail groupoid is ample]
If $\Omega(\mathsf B)$ is a Cantor space, then $\mathfrak T_{\mathsf B}$ is a principal ample groupoid whose unit space is naturally identified with $\Omega(\mathsf B)$ \citep{gps1995,putnam2010}.
\end{proposition}

\begin{warning}[do not silently use the subspace topology from the product space]
Although, as a set,
\[
\mathfrak T_{\mathsf B}
\subseteq
\Omega(\mathsf B)\times\Omega(\mathsf B),
\]
the groupoid topology in this chapter is not obtained merely by saying ``take the subspace topology from the product.'' The \(\acute{e}\)tale structure is determined by the prefix-replacement bisections above. In particular, the finite-stage subgroupoids introduced next are open in the groupoid topology; intuition from the product topology must not be substituted for this definition.
\end{warning}

\section[The Finite Stages Tn: The Unit Space Remains Infinite, but Every Orbit Is Finite]{The Finite Stages $\mathfrak T_n$: The Unit Space Remains Infinite, but Every Orbit Is Finite}

Before ``approximately finite'' appears in the formal definition, consider the most natural sequence of finite stages inside the tail groupoid. For $n\ge0$, define
\[
\mathfrak T_n
=
\{(x,y)\in\mathfrak T_{\mathsf B}:x_k=y_k\text{ for every }k>n\}.
\]
Thus, inside $\mathfrak T_n$, two paths may differ only through their first $n$ levels; from the $(n+1)$st edge onward they must agree exactly.

If $\Omega_n(v)$ denotes the length-$n$ paths ending at $v\in V_n$, then
\[
\boxed{
\mathfrak T_n
=
\bigsqcup_{v\in V_n}
\ \bigsqcup_{\alpha,\beta\in\Omega_n(v)}
Z(\alpha,\beta).}
\]
There are only finitely many bisections on the right because $V_n$ and all the sets $\Omega_n(v)$ are finite.

\begin{plainbox}[Exactly what is finite here?]
The unit space of $\mathfrak T_n$ is still the entire Cantor space, so $\mathfrak T_n$ itself usually has uncountably many arrows. Its ``finiteness'' does \emph{not} mean that the subgroupoid has finitely many elements. Rather, after an infinite tail has been fixed, only finitely many length-$n$ prefixes are available for prefix replacement, so every orbit is finite; at the same time, the whole subgroupoid is assembled from finitely many compact-open bisections.
\end{plainbox}

\begin{definition}[elementary groupoid]
In the Cantor-unit setting of this book, let $\mathcal G$ be an ample groupoid. A subgroupoid $\mathcal K\subseteq\mathcal G$ is called an \emph{elementary subgroupoid} if $\mathcal K$ is compact-open and principal and
\[
\mathcal K^{(0)}=\mathcal G^{(0)}
\]
\citep{matui2012}.
\end{definition}

\begin{proposition}[every $\mathfrak T_n$ is elementary]
For every $n$, $\mathfrak T_n$ is an elementary subgroupoid of $\mathfrak T_{\mathsf B}$, and
\[
\mathfrak T_0
\subseteq
\mathfrak T_1
\subseteq
\mathfrak T_2
\subseteq\cdots.
\]
\end{proposition}

\begin{proof}
The finite bisection decomposition above shows that $\mathfrak T_n$ is compact-open. It contains every unit because $(x,x)\in\mathfrak T_n$, and it remains an equivalence-relation groupoid, hence is principal. If two paths already agree after level $n$, then they certainly agree after level $n+1$, so $\mathfrak T_n\subseteq\mathfrak T_{n+1}$.
\end{proof}

\begin{example}[one vertex and two edges at every level]
Consider the Bratteli diagram with one vertex at every level and two edges, labeled $0$ and $1$, from each level to the next. Its path space is
\[
\{0,1\}^{\mathbb N}.
\]
Inside $\mathfrak T_n$, once the infinite tail beginning with coordinate $n+1$ is fixed, the first $n$ coordinates may be chosen arbitrarily. Hence every $\mathfrak T_n$-orbit has exactly $2^n$ points. As $n$ increases, increasingly long finite prefixes may be modified; taking all stages together yields the full tail-equivalence relation, in which two sequences may differ in only finitely many coordinates.
\end{example}

\section{AF Groupoids: Here the ``Approximation'' Is an Exact Union}

Only now do we define AF.

\begin{definition}[AF groupoid]
We call an ample groupoid $\mathcal H$ an \emph{AF (approximately finite) groupoid} if there are increasing elementary subgroupoids
\[
\mathcal H_0\subseteq\mathcal H_1\subseteq\mathcal H_2\subseteq\cdots,
\qquad
\mathcal H_n^{(0)}=\mathcal H^{(0)},
\]
such that
\[
\boxed{
\mathcal H=\bigcup_{n\ge0}\mathcal H_n.}
\]
This is the standard AF condition \citep{matui2012}.
\end{definition}

\begin{plainbox}[The easiest misunderstanding of ``approximately finite'']
There is no numerical approximation here with an error $\varepsilon_n\to0$, nor is $\mathcal H$ obtained only after taking the closure of the finite-stage subgroupoids. The definition requires an exact equality of sets: \textbf{every concrete arrow already appears at some finite stage.} The infinitude comes from there being no single final stage, not from each stage merely making an approximate guess at the complete arrow set.
\end{plainbox}

For the Bratteli tail groupoid, tail equivalence itself means agreement from some level onward. Consequently every arrow $(x,y)$ belongs to some $\mathfrak T_n$, and therefore
\[
\boxed{
\mathfrak T_{\mathsf B}
=
\bigcup_{n\ge0}\mathfrak T_n.}
\]
Hence:

\begin{theorem}[the Bratteli tail groupoid is AF]
If $\Omega(\mathsf B)$ is a Cantor space, then its tail groupoid $\mathfrak T_{\mathsf B}$ is an AF groupoid.
\end{theorem}

An AF groupoid is automatically principal as well. Every isotropy arrow lies in some elementary stage, and every elementary stage is principal.

\begin{plainbox}[Why this matters later]
This point will be important throughout the book: the AF part itself has no nontrivial germ isotropy. If a larger system later develops nontrivial isotropy or singular germs, that additional information cannot be said to have been ``already inside the AF core.'' An AF groupoid provides a principal background relation exhausted by elementary finite stages; new local isotropy belongs to the structure added later.
\end{plainbox}

\section{Simple Bratteli Diagrams and Minimal Tail Groupoids}

The preceding chapter recorded the standard fact that $\mathsf B$ is simple if and only if every cofinality class is dense in the path space. Since an orbit is now exactly a cofinality class, the same fact can be translated directly into groupoid language.

\begin{proposition}[simplicity $\Longleftrightarrow$ minimality]
Assume that $\Omega(\mathsf B)$ is a Cantor space. Then
\[
\mathsf B\text{ is simple}
\quad\Longleftrightarrow\quad
\mathfrak T_{\mathsf B}\text{ is minimal}.
\]
\end{proposition}

\begin{proof}
The $\mathfrak T_{\mathsf B}$-orbit of $x$ is exactly $\operatorname{Cof}(x)$. A groupoid is minimal when every orbit is dense in the unit space, while the preceding chapter already established the equivalence between simplicity of $\mathsf B$ and density of every cofinality class \citep{gps1995,putnam2010}.
\end{proof}

\begin{plainbox}[Why ``simple'' is not a decorative hypothesis]
If the Bratteli diagram is not simple, the path space may split into regions that can never be reached from one another by tail arrows. The AF background is then no longer a single minimal dynamical whole.
\end{plainbox}

\section{What Is an AF Core? A Role Name, Not a New Kind of Object}

Later chapters study groupoids larger than an AF groupoid. If a larger groupoid $\mathcal G$ contains a distinguished AF subgroupoid
\[
\mathcal H\subseteq\mathcal G,
\qquad
\mathcal H^{(0)}=\mathcal G^{(0)},
\]
this book will usually call $\mathcal H$ the \emph{AF core}. The phrase describes only the role it plays as the ``AF background'' inside the larger structure; it does not alter the definition of an AF groupoid.

At this point three layers of information must be kept permanently separate.

\begin{center}
\small
\begin{tabularx}{\textwidth}{>{\bfseries\raggedright\arraybackslash}p{0.25\textwidth}Y Y}
\toprule
Layer & Concrete data & Status \\
\midrule
AF groupoid intrinsic data
& unit space, arrows, source/range, composition, inverse, topology
& belongs to the groupoid itself and should persist under a change of presentation \\
\addlinespace
AF presentation data
& a particular Bratteli diagram, a chosen sequence of elementary stages $\mathcal H_n$, literal level numbers, cylinder names
& displays the AF structure, but is usually not a canonical invariant \\
\addlinespace
chosen generator data
& generators selected in a Cayley/orbital graph, edge labels, colors, and letters
& a further layer of presentation; it must not masquerade as intrinsic labeling of the AF core \\
\bottomrule
\end{tabularx}
\end{center}

\begin{principle}[the AF core has no intrinsic generator labels]
The intrinsic data of an AF core do not include the names of a chosen generating family, nor do they include edge labels drawn from those generators. Two entirely different generating or labeling schemes may describe the same AF groupoid.
\end{principle}

\begin{plainbox}[``No labels'' emphatically does not mean ``no structure'']
The AF core still has a rich system of arrows, orbits, topology, and finite-stage exhaustion. Once a Bratteli presentation is chosen, one can also see a whole sequence of finite-stage subgroupoids. What is absent is only a set of externally chosen generator names treated as intrinsic attributes. Thus, if a picture uses three colors labeled $a,b,c$, one cannot reverse the logic and say that a point of the AF core ``naturally carries the label $a$.''
\end{plainbox}

\section{What Does Telescoping Do on the Groupoid Side? Stage Numbers Are Not Intrinsic Invariants}

The preceding chapter showed that Bratteli telescoping packages several consecutive levels into a single step without changing the path space or tail equivalence. On the groupoid side, the same operation amounts to passing from
\[
\mathfrak T_0
\subseteq
\mathfrak T_1
\subseteq
\mathfrak T_2
\subseteq\cdots
\]
to a cofinal subsequence
\[
\mathfrak T_{n_0}
\subseteq
\mathfrak T_{n_1}
\subseteq
\mathfrak T_{n_2}
\subseteq\cdots,
\qquad n_k\to\infty.
\]
Because the subsequence is cofinal, we still have
\[
\mathfrak T_{\mathsf B}
=
\bigcup_k\mathfrak T_{n_k}.
\]

\begin{warning}[a literal stage number is only a presentation coordinate]
An arrow may first enter $\mathfrak T_{17}$ in the current presentation; after telescoping, it may already appear at stage $4$ of the new presentation. The integer stage number itself is therefore generally not an intrinsic invariant of the AF groupoid. What remains stable is the fact that the arrow belongs to some finite stage, together with the cofinal finite-stage exhaustion as a whole.

This discipline will recur later. Any number that depends on ``the level at which something is born'' must first be checked: does it describe the groupoid itself, or only the current Bratteli presentation?
\end{warning}

\section{What Has This Chapter Given Us, and What Has It Still Not Given Us?}

The mathematical role of the AF background is now precise. It gives us:
\begin{enumerate}
\item a principal \(\acute{e}\)tale/ample arrow system on a Cantor unit space;
\item a genuinely increasing sequence of elementary subgroupoids, with every arrow eventually entering some finite stage;
\item in the Bratteli setting, the standard local rule that a finite prefix may change while the common tail remains untouched;
\item minimality under the simplicity hypothesis;
\item a clear separation between finite-stage presentation data and intrinsic groupoid data;
\item no canonical generator labels, and no additional interface semantics produced automatically by such labels.
\end{enumerate}

But this chapter has not yet supplied interfaces recording \emph{where} a finite piece may continue outward, nor has it defined how old finite pieces are copied, embedded, and joined through such interfaces to form new pieces. Those data are not secretly contained in the definition of AF. The next chapter first defines these finite-piece interfaces independently, and only then studies their level-by-level assembly and growth.

\begin{readerbox}[Shortest summary of the chapter]
The Bratteli tail groupoid turns ``eventual cofinality'' into a system of locally invertible arrows. At stage $n$, only the first $n$ levels may be changed, so the stage is elementary. The exact union of all stages is the full tail groupoid, hence it is AF. The strength of the AF core lies in an arrow structure that is exhausted by finite stages, not in generator labels written in advance on its points.
\end{readerbox}

\chapter{Graph with Boundary and Tile Inflation: How Finite Pieces Continue to Grow}
\label{chap:tiles-en}

\par\medskip
\noindent\textbf{Why this chapter matters}\par
Up to the preceding chapter, we have learned to use an AF groupoid to describe how increasingly large finite relations combine exactly into an infinite object. One question was left unanswered on purpose: if a finite piece does not yet contain every possible continuation, how can the mathematics record that ``continuation is still possible here''?

A graph with boundary answers exactly this question. It allows a finite graph to retain a small number of edges whose second endpoint is still undefined; tile inflation then repeatedly copies old tiles and inserts connectors between such unresolved interfaces. The central point is not that there are ``many boundaries.'' It is that \textbf{large bodies of already determined internal structure can keep being copied, while the extra data that decide how distinct copies are joined are concentrated at a small number of boundary interfaces.}
\par\medskip

\par\medskip
\noindent\textbf{Mathematical sources and indexing convention for this chapter}\par
The definitions of graph with boundary, compatible boundary edges, connectors, tile inflation, infinite tile, and boundary point in this chapter come from Kuang's tile-inflation framework \citep{kuang2026a}, and use the same geometric language as the near-full work \citep{kuang2026b}. The source paper numbers Bratteli levels beginning with $V_1$; this book has already fixed $V_0$ as the top level, so all indices here are rewritten uniformly as
\[
E_n:V_{n-1}\longrightarrow V_n.
\]
Apart from this shift, the relational structure of the original construction is preserved.
\par\medskip

\section{Why Ordinary Finite Graphs Are Not Enough: We Need to Retain ``Unfinished Endpoints''}

Every edge in an ordinary graph has two endpoints that already exist. If our only object of study were a finite graph whose incidences had all been settled, that would be enough. Tile inflation, however, is designed precisely so that a finite piece can keep growing. At the current scale, some edges therefore display only one endpoint: they tell us that a connection has already extended out from the internal part of the piece, while its other endpoint will be determined only in a larger tile.

A graph with boundary is consequently not a ``badly drawn graph.'' It is a finite graph that \emph{keeps some connection data unresolved on purpose: the second endpoint has not yet been determined}. Internal edges whose incidences are already fixed continue to behave exactly like ordinary graph edges; only a small number of edges are allowed to hang temporarily.

\begin{definition}[graph with boundary]
A \emph{graph with boundary} $\Gamma$ consists of a vertex set $V(\Gamma)$, an edge set $E(\Gamma)$, two partially defined maps
\[
\mathsf s,\mathsf r:E(\Gamma)\dashrightarrow V(\Gamma),
\]
and a fixed-point-free involution
\[
e\longmapsto e^{-1},
\]
satisfying the following conditions:
\begin{enumerate}
\item every edge $e$ belongs to the domain of at least one of $\mathsf s$ and $\mathsf r$; in other words, an edge may not have both endpoints completely unknown;
\item $(e^{-1})^{-1}=e$ and $e^{-1}\neq e$;
\item inversion exchanges whichever source and range values have already been defined: whenever the corresponding side is defined,
\[
\mathsf s(e)=\mathsf r(e^{-1}),
\qquad\text{or equivalently}\qquad
\mathsf r(e)=\mathsf s(e^{-1}).
\]
\end{enumerate}
\end{definition}

\begin{plainbox}[In ordinary language]
An ordinary internal edge is a pair of oppositely directed arrows whose two endpoints both lie in the current finite piece. A boundary edge is more like a short line protruding through a wall: we know either which internal vertex it leaves or which internal vertex it enters, but its other endpoint has not yet appeared at the current scale. Taking the inverse merely records the same potential connection from the opposite direction.
\end{plainbox}

\section{Labels, Boundary Edges, and Boundary Vertices}

Fix a finite symmetric label set $\mathcal S$ with involution
\[
F\longmapsto F^{-1}.
\]
Edge labels must be compatible with inversion: if an edge $e$ is labeled $F$, then $e^{-1}$ is labeled $F^{-1}$. These labels still belong to the chosen presentation. They indicate a type of local action and must not be conflated with intrinsic data of the AF core from the preceding chapter.

For a label $F$, the ordinary internal situation is a pair of vertices $u,v$ with
\[
u\xrightarrow{F}v,
\qquad
v\xrightarrow{F^{-1}}u.
\]
There is also the case in which only half of the arrow has acquired endpoints. For example, an edge $e$ labeled $F$ may satisfy $\mathsf s(e)=u$ while $\mathsf r(e)$ is undefined; its inverse edge $e^{-1}$ is labeled $F^{-1}$, satisfies $\mathsf r(e^{-1})=u$, and has $\mathsf s(e^{-1})$ undefined.

\begin{definition}[boundary edge and boundary vertex]
We call an edge $e$ a \emph{boundary edge} if exactly one of its endpoints is defined. The defined vertex incident to this boundary edge is called the corresponding \emph{boundary vertex}. An edge for which both endpoints are defined is called an \emph{internal edge}.
\end{definition}

\begin{plainbox}[In ordinary language]
``Boundary'' describes a \emph{relational position}. A vertex is a boundary vertex not because it has a special internal shape, but because the current finite piece still carries at that vertex a connection whose other endpoint has not yet been supplied. A larger tile may later join that hanging connection to another copy.
\end{plainbox}

\begin{center}
\begin{tikzpicture}[>=Latex,node distance=13mm,
 v/.style={circle,draw,inner sep=1.8pt}, b/.style={circle,draw,very thick,inner sep=2.2pt}]
\node[v] (a) {};
\node[v,right=of a] (b) {};
\node[b,right=of b] (c) {};
\draw[<->] (a)--node[above]{internal}(b);
\draw[->] (b)--(c);
\draw[->,dashed] (c)--++(1.4,0) node[right]{undefined endpoint};
\node[below=4mm of c]{boundary vertex};
\end{tikzpicture}
\end{center}

Chapter 1 of this book already defined a well-labeled graph. Finite tiles here are likewise required to be well-labeled: for every vertex and every label $F$, there is at most one outgoing edge labeled $F$ and at most one incoming edge labeled $F$. This uniqueness is crucial. When a label $F$ later defines a local transformation on path space, a point cannot simultaneously have two competing $F$-continuations nearby.

\section{Compatible Boundary Edges: Which Hanging Endpoints May Be Joined?}

Not every pair of boundary edges may be connected. Their labels must match, and the two missing endpoints must face complementary directions.

\begin{definition}[compatible boundary edges and admissible boundary vertices]
Let $e_1,e_2$ be boundary edges in two level-$n$ tiles with the same label $F\in\mathcal S$. We call $e_1$ and $e_2$ \emph{compatible} if either
\begin{itemize}
\item the source of $e_1$ is defined and its range is undefined, while the range of $e_2$ is defined and its source is undefined; or
\item conversely, $e_1$ is missing its source while $e_2$ is missing its range.
\end{itemize}
Their respective defined boundary vertices are called a pair of \emph{admissible boundary vertices}.
\end{definition}

\begin{plainbox}[Why the same label, rather than inverse labels?]
This is an easy point to write incorrectly. The two hanging edges being paired both carry the same label $F$ in the original definition: one is missing a range and the other is missing a source. Once they are paired, together they form a genuine $F$-arrow. The reverse arrow is then supplied automatically by their inverse edges and carries label $F^{-1}$. Thus the input to a connector is a pair of boundary edges with \textbf{the same label and complementary missing-endpoint directions}, not a pair whose labels are mutual inverses.
\end{plainbox}

If $F=F^{-1}$, then in some situations a boundary vertex can even be admissible with itself. This is only a special case of the label involution and requires no additional rule.

\section{Finite Tiles: Finite History Blocks Beneath the Same Higher-Level Vertex}

Now place graphs with boundary on the Bratteli diagram already fixed in the preceding chapters:
\[
\mathsf B=(V,E,\mathbf s,\mathbf r),
\qquad
E_n:V_{n-1}\longrightarrow V_n.
\]

For $v\in V_n$, the vertex set of the level-$n$ tile $\mathcal T_{v,n}$ consists of all finite Bratteli paths of length $n$ that start at $V_0$ and terminate at $v$. In other words, one tile places into the same finite graph with boundary those finite paths that have different early histories but merge at the same vertex $v$ on level $n$.

\begin{plainbox}[A Bratteli path and a tile vertex play two different roles]
A finite Bratteli path $\eta=e_1\cdots e_n$ is, on the one hand, a path in the Bratteli diagram and, on the other, a vertex of the tile $\mathcal T_{v,n}$. In the first role it records a hierarchical history; in the second it marks the concrete position occupied by that history in the current finite tile. Later, writing $\eta e$ extends the old history by one more Bratteli edge and simultaneously places the old tile vertex into a copy inside the higher-level tile.
\end{plainbox}

\section{The First Step of Inflation: Copy the Old Tiles First}

Take $v\in V_{n+1}$. Every Bratteli edge
\[
e:u\longrightarrow v,
\qquad \mathbf r(e)=v,
\]
tells us that the level-$(n+1)$ tile $\mathcal T_{v,n+1}$ must first contain a copy of the level-$n$ tile $\mathcal T_{u,n}$.

More precisely, for every such $e$ there is an embedding
\[
\phi_{u,e}:\mathcal T_{u,n}\hookrightarrow\mathcal T_{v,n+1},
\qquad
\eta\longmapsto\eta e.
\]
Before any connectors are inserted, these copies may be regarded as mutually separate finite pieces.

\begin{center}
\begin{tikzpicture}[>=Latex,node distance=11mm,
 tile/.style={draw,rounded corners,minimum width=3.0cm,minimum height=1.35cm,fill=boxbg}]
\node[tile] (t1) {$\mathcal T_{u_1,n}$ copy};
\node[tile,right=13mm of t1] (t2) {$\mathcal T_{u_2,n}$ copy};
\node[tile,right=13mm of t2] (t3) {$\mathcal T_{u_3,n}$ copy};
\node[below=8mm of t2,align=center] {Copy first; for now, do not decide\\how the copies are connected};
\end{tikzpicture}
\end{center}

This step uses only the Bratteli incidence data. It tells us which lower-level copies appear inside the higher-level tile, but it does not yet tell us how those copies are to be joined at their boundaries. That is exactly the task of the connector data.

\section{Connectors: The Data That Actually Join Distinct Occurrences}

\begin{definition}[level-$(n+1)$ connector]
A level-$(n+1)$ \emph{connector} is a triple
\[
(\xi_1e_1,\xi_2e_2,F),
\]
where $e_1,e_2$ both terminate at the same vertex $v\in V_{n+1}$; each $\xi_i$ is an admissible boundary vertex in the corresponding level-$n$ tile; and at these two boundary vertices there is a compatible pair of boundary edges, both labeled by $F$. The points $\xi_1e_1$ and $\xi_2e_2$ are called the \emph{connecting points} of this connector.
\end{definition}

Once a connector has been chosen, the two complementary hanging $F$-interfaces are turned into a pair of genuine internal arrows. Their direction is determined by which hanging edge is missing a source and which is missing a range. Thus one possibility is
\[
\xi_1e_1\xrightarrow{F}\xi_2e_2,
\qquad
\xi_2e_2\xrightarrow{F^{-1}}\xi_1e_1,
\]
while the whole orientation may instead be reversed, in which case the first displayed arrow carries $F^{-1}$ and the reverse arrow carries $F$. The key point is not to choose an orientation in advance. A connector must \emph{supply the unique internal connection required by the two compatible hanging edges}. Performing this operation for all level-$(n+1)$ connectors produces the new tile $\mathcal T_{v,n+1}$. Boundary interfaces not consumed by a connector remain available for treatment at higher levels.

\begin{plainbox}[A connector joins concrete occurrences, not abstract types]
The same level-$n$ tile type may appear many times inside a higher-level tile. A connector joins the particular points $\xi_1e_1$ and $\xi_2e_2$; it does not merely assert that ``type A is connected to type B.'' This is exactly why Chapter 1 insisted on occurrence-first bookkeeping. If an occurrence is forgotten into its type too early, one can simultaneously lose which copy it belongs to, which connector lies next to it, and how its later continuation unfolds.
\end{plainbox}

\begin{center}
\begin{tikzpicture}[>=Latex,node distance=11mm,
 tile/.style={draw,rounded corners,minimum width=2.85cm,minimum height=1.45cm,fill=boxbg},
 b/.style={circle,draw,very thick,inner sep=2pt}]
\node[tile] (t1) {$\mathcal T_{u_1,n}$ copy};
\node[tile,right=25mm of t1] (t2) {$\mathcal T_{u_2,n}$ copy};
\node[b] (b1) at ($(t1.east)+(-0.15,0)$) {};
\node[b] (b2) at ($(t2.west)+(0.15,0)$) {};
\draw[->,very thick] (b1) to[bend left=16] node[above]{$F$} (b2);
\draw[->,very thick] (b2) to[bend left=16] node[below]{$F^{-1}$} (b1);
\node[below=9mm of $(t1)!0.5!(t2)$,align=center] {A connector turns two concrete boundary\\occurrences into an internal connection};
\end{tikzpicture}
\end{center}

\section{How Boundary and Connecting Points Change Under Inflation}

There is an important one-way rule here, and it will be used repeatedly later. If a level-$n$ vertex is not a boundary vertex, then extending it along a Bratteli edge to level $n+1$ does not suddenly make it a boundary point or a connecting point. Every higher-level interface must have lower-level boundary ancestry.

Conversely, when a level-$n$ boundary vertex is extended to the next level, it has two basic possible fates:
\begin{enumerate}
\item it retains some hanging boundary interface and therefore remains a level-$(n+1)$ boundary vertex;
\item it is used by a connector, thereby becoming a connecting point, while the original hanging interface has become internal in the higher-level tile.
\end{enumerate}
Both may happen simultaneously because a single vertex can carry more than one boundary label.

\begin{principle}[boundary ancestry]
Under tile inflation, new boundary/connector organization can propagate only along boundary data that already exist. A non-boundary vertex cannot spontaneously acquire boundary or connector status at a higher level.
\end{principle}

\begin{plainbox}[Why this matters]
Boundary data are not an arbitrary new list of special points chosen afresh at every level. A higher-level interface has its own finite history: it must continue from a hanging connection already present at a lower level, or be absorbed by a connector at some stage. Later, when the book asks where a ``crack'' came from, what can actually be traced is this finite-level ancestry.
\end{plainbox}

\section{Infinite Tiles: Taking an Inductive Limit Along a Bratteli Path}

Take an infinite path
\[
x=e_1e_2e_3\cdots\in\Omega(\mathsf B),
\]
and write $v_n=\mathbf r(e_n)\in V_n$. Along the successive edges of $x$, the append-one-edge embeddings above give a compatible system
\[
\mathcal T_{v_1,1}
\xhookrightarrow{\phi_{v_1,e_2}}
\mathcal T_{v_2,2}
\xhookrightarrow{\phi_{v_2,e_3}}
\mathcal T_{v_3,3}
\hookrightarrow\cdots.
\]

\begin{definition}[infinite tile]
The graph inductive limit of the finite tiles above is called the \emph{infinite tile} of $x$ and is denoted by $\mathcal T_x$:
\[
\mathcal T_x
=
\varinjlim_n \mathcal T_{v_n,n}.
\]
\end{definition}

In this inductive limit, different finite paths that eventually enter the same tail can occur inside the same infinite tile, in keeping with the Bratteli tail-equivalence language from the preceding chapters. At this point we still need only regard an infinite tile as a compatible union of finite occurrences; we do not introduce the \AFBD{} groupoid definition from the next chapter in advance.

\section{Persistent Boundary Points: Which Interfaces Survive to Infinite Scale?}

A finite boundary need not remain a boundary forever. It may be absorbed by a connector at the next level and thereby become an ordinary internal connection of a larger tile. What enters the boundary of an infinite tile is the compatible boundary history that continues to retain boundary status at every deeper level.

\begin{definition}[boundary point of an infinite tile]
Let $\xi$ be a vertex of the infinite tile $\mathcal T_x$, represented by compatible finite truncations. We call $\xi$ a \emph{boundary point} of $\mathcal T_x$ if, for every level $n$, the truncation $\xi|_n$ is a boundary vertex of the corresponding finite tile. After telescoping or reindexing the Bratteli presentation from a new base level, ``every level'' means every level in that new presentation; the definition still requires boundary status to persist along the chosen system of levels.
\end{definition}

\begin{plainbox}[Finite boundary and infinite boundary must be kept separate]
A boundary at level $n$ says only that ``an interface is still hanging at this scale.'' An infinite boundary requires that the interface never become fully internal as the scale increases. The former may appear and later disappear; the latter records a boundary history that genuinely persists to the inductive limit.
\end{plainbox}

The following three cases must therefore be distinguished:
\begin{center}
\small
\begin{tabularx}{0.97\textwidth}{>{\bfseries}p{2.7cm}Y Y}
\toprule
Finite-level status & What may happen at the next level & Limit meaning \\
\midrule
non-boundary vertex & continues as an ordinary internal vertex & cannot spontaneously become boundary/connector \\
\addlinespace
boundary vertex & is absorbed by a connector & a finite boundary witness exists, but no persistent boundary results \\
\addlinespace
boundary vertex & retains a boundary interface & if this persists cofinally, it yields an infinite boundary point \\
\bottomrule
\end{tabularx}
\end{center}

\section{How Labels Produce Partial Homeomorphisms on the Path Space}

At this point the abstract language of partial homeomorphisms and inverse semigroups from the preceding chapters finally acquires a concrete source. Let $F\in\mathcal S$. If a vertex $x\in\Omega(\mathsf B)$ in an infinite tile has a genuine internal outgoing edge
\[
x\xrightarrow{F}y,
\]
then well-labeledness guarantees that there is at most one such $y$, so we may define
\[
F(x)=y.
\]

More importantly, this does not hold only accidentally at a single point. If the $F$-edge has already appeared at some finite level, then every path sharing a sufficiently long common Bratteli prefix sees the same finite local rule. Consequently the domain of $F$ is open, and $F^{-1}$ supplies the inverse local rule. The labels in a tile inflation therefore naturally produce partial homeomorphisms of the path space.

\begin{plainbox}[Here boundary first appears as a location where the definition has not yet been extended]
At a boundary point, the finite tiles may continue to display a hanging edge labeled $F$ without ever supplying its other endpoint inside the current infinite tile. The internal tile action may therefore still be undefined there. When a larger groupoid is constructed later, one must decide whether, and in what manner, these boundary interfaces can be extended by additional discrete germs. This is precisely the interface through which the book passes from tile inflation to \AFBD{}.
\end{plainbox}

In some especially regular situations, these local homeomorphisms can be extended to the \emph{closure of their domains}, with labeled arrows added between the corresponding boundary points; one then obtains a partial action with clopen domains. This additional extension is not guaranteed by the definition of graph with boundary itself. We therefore record it here only as extra structure that must be checked in the next chapter, rather than building it silently into tile inflation.

\section{A Minimal Two-Copy Inflation Picture}

Compress the preceding definitions into a minimal model. Suppose a higher-level vertex $v\in V_{n+1}$ has two incoming Bratteli edges
\[
e_1:u_1\longrightarrow v,
\qquad
e_2:u_2\longrightarrow v.
\]
Then $\mathcal T_{v,n+1}$ first contains the two concrete copies
\[
\phi_{u_1,e_1}(\mathcal T_{u_1,n})
\quad\text{and}\quad
\phi_{u_2,e_2}(\mathcal T_{u_2,n}).
\]
If $\xi_1e_1$ in the first copy and $\xi_2e_2$ in the second carry compatible $F$-boundary edges, we may choose the connector
\[
(\xi_1e_1,\xi_2e_2,F)
\]
and insert the internal $F/F^{-1}$ edge pair.

Thus one inflation step separates into three strictly different operations:
\[
\boxed{
\begin{gathered}
\text{copy the old tile}
\;\longrightarrow\;
\text{identify admissible boundary occurrences}\\[1mm]
\longrightarrow\;
\text{use connectors to insert new internal connections}
\end{gathered}}
\]

The first operation preserves old structure; the second merely reads boundary data that are already present; only the third genuinely creates new adjacency between distinct copies. When later chapters say that ``a small amount of boundary data organizes a large body of internal structure,'' this asymmetric division of labor is what the mathematics actually means.

\section{Chapter Summary: What Does Tile Inflation Add?}

We have now established the first genuinely geometric layer of Part II that goes beyond the AF core. What must be retained is not the loose statement that ``tiles get larger,'' but the following seven distinctions:
\begin{enumerate}
\item a graph with boundary allows a small number of edges to have only one endpoint defined temporarily;
\item a boundary vertex is an interface position in the current finite piece, not a special vertex type;
\item compatible boundary edges must carry the same label and have complementary missing-endpoint directions;
\item a higher-level tile is first built from copies of lower-level tile occurrences;
\item a connector acts on concrete occurrences and turns hanging interfaces into internal edges;
\item non-boundary data do not spontaneously acquire boundary/connector identity at higher levels, while a finite boundary can either be absorbed by a connector or continue to persist;
\item taking the inductive limit of compatible embeddings along an infinite Bratteli path produces an infinite tile, and only a boundary history that persists cofinally becomes an infinite boundary point.
\end{enumerate}

More compactly: Bratteli incidence data tell us \emph{which old pieces are copied into the same higher-level piece}; boundary data tell us \emph{which positions still permit additional continuation}; connector data tell us \emph{how those copies are actually joined at the higher level}.

The next chapter will give the strict definition of \AFBD{} for the first time on this concrete model. It must continue to respect the discipline established here: the AF core, the Bratteli presentation, boundary interfaces, connector data, and additional germs are distinct layers of mathematical data. They must not be blended into a single object merely because they eventually occur inside the same groupoid.

\chapter{AF-by-Discrete Groupoids: Strict Definition and Tile Presentation}
\label{chap:afbd-en}

\par\medskip
\noindent\textbf{Why this chapter matters}\par
The preceding two chapters established the AF core and tile inflation separately. We can now place them inside a single strict object: an \AFBD{} groupoid. The name invites two misunderstandings. First, ``discrete'' does not mean ``there are only finitely many extra arrows.'' Second, the appearance of a boundary edge in a finite tile does not automatically mean that a non-AF germ has already appeared there. The purpose of this chapter is to separate these layers once and for all: what belongs to the definition itself, what belongs only to a tile presentation, and what lands in $\mathfrak G\setminus\mathfrak T_{\mathsf B}$ only after it persists through the infinite hierarchy.
\par\medskip

\par\medskip
\noindent\textbf{Mathematical sources for this chapter}\par
The definition of \AFBD{} follows Nekrashevych's AF-by-discrete framework and is fixed here over the Bratteli tail groupoid in the form used in Kuang's work on bounded-type tile inflations \citep{nek2022,kuang2026b}. The propositions that the tail groupoid is an open subgroupoid, that the internal part of a finite tile is an elementary AF Cayley graph, and that a boundary edge may become internal at a deeper level are organized from Kuang's tile-presentation results \citep{kuang2026b}. The aim here is not to reinvent these objects, but to make explicit the logical order on which the later chapters depend.
\par\medskip

\section{First Clarify What ``Discrete'' Means}

\begin{definition}[discrete subspace]
Let $D\subseteq Y$ be a subset of a topological space. We call $D$ \emph{discrete} in the subspace topology if for every $d\in D$ there exists an open set $U_d\subseteq Y$ such that
\[
U_d\cap D=\{d\}.
\]
\end{definition}

\begin{warning}[Discrete does not mean finite]
Whenever the later text says that the non-AF part is discrete, this must never be rewritten automatically as ``there are only finitely many non-AF arrows.'' What turns discreteness into finiteness is the additional hypothesis of compactness. That implication will be proved separately below.
\end{warning}

\section{The Strict Definition of AF-by-Discrete}

\begin{definition}[AF-by-discrete over a Bratteli diagram]
Let $\mathsf B$ be a Bratteli diagram, let $\Omega(\mathsf B)$ be its infinite path space, and let $\mathfrak T_{\mathsf B}$ be the tail groupoid. Let $\mathfrak G$ be an \emph{\'etale} groupoid with the same unit space $\Omega(\mathsf B)$. We call $\mathfrak G$ over $\mathsf B$ \emph{AF-by-discrete} if
\[
\mathfrak T_{\mathsf B}\subseteq\mathfrak G
\]
is an open subgroupoid and
\[
\mathfrak G\setminus\mathfrak T_{\mathsf B}
\]
is discrete in the subspace topology. In this given presentation, $\mathfrak T_{\mathsf B}$ is called the \emph{AF core} of $\mathfrak G$ \citep{nek2022,kuang2026b}.
\end{definition}

\begin{plainbox}[In ordinary language]
The definition has only two substantive requirements. First, a large AF relation must sit inside the ambient groupoid as an open subgroupoid. Second, all arrows outside that AF core must be locally separable from one another in the subspace topology. The object therefore has an open background that is exhausted by finite stages, together with additional arrows that form a discrete subspace.

The simplest example is the purely AF case. If $\mathfrak G=\mathfrak T_{\mathsf B}$, then the complement is empty, and the empty set is of course discrete. Thus an AF groupoid itself is a degenerate but legitimate \AFBD{} groupoid. The phrase ``by-discrete'' does not require the existence of extra arrows.
\end{plainbox}

\begin{stopbox}[What the name does not secretly promise]
The phrase ``by-discrete'' in the definition does not assert the existence of a short exact sequence, nor does it say that $\mathfrak G/\mathfrak T_{\mathsf B}$ is a discrete group. In fact, $\mathfrak G\setminus\mathfrak T_{\mathsf B}$ is generally not a subgroupoid: the product of two non-AF arrows may lie back in the AF core. The name records only a topological division of labor between an open AF core and its discrete complement.
\end{stopbox}

There is also a presentation warning. Throughout this book, the phrase ``the AF core'' is relative to the currently chosen Bratteli presentation and the embedding
\[
\mathfrak T_{\mathsf B}\hookrightarrow\mathfrak G.
\]
The definition itself does not prove that a given $\mathfrak G$ admits only one AF presentation. Whenever later chapters refer to a level, tile, boundary, or connector, those objects must therefore be remembered as presentation-dependent.

\section{From Abstract Discreteness to the Fixed-Generator Finite-Boundary Consequence}

The abstract definition says ``open AF core plus discrete complement.'' Tile language requires one more layer: fix compact-open generators and specify how non-AF generator germs are read as persistent boundary incidences. The dictionary between these two layers may not be suppressed by writing ``equivalently.''

\begin{proposition}[A compact region sees only finitely many non-AF arrows]
Suppose $\mathfrak G$ over $\mathsf B$ is \AFBD{} and $K\subseteq\mathfrak G$ is compact. Then
\[
K\cap\bigl(\mathfrak G\setminus\mathfrak T_{\mathsf B}\bigr)
\]
is finite.
\end{proposition}

\begin{proof}
Because $\mathfrak T_{\mathsf B}$ is open, its complement
\[
D:=\mathfrak G\setminus\mathfrak T_{\mathsf B}
\]
is \emph{closed}. Hence $K\cap D$ is compact. On the other hand, $D$ is discrete in the subspace topology, so $K\cap D$ is a compact discrete space and must therefore be finite.
\end{proof}

Now fix a finite compact-open generating family invariant under inversion
\[
S=\{F_1,\ldots,F_r\}.
\]
Set
\[
E_S:=\bigcup_{i=1}^{r}\bigl(F_i\cap D\bigr).
\]
Each $F_i$ is compact, so the preceding proposition gives $|E_S|<\infty$.

\begin{definition}[the infinite-tile indexing convention used in this book]
Let $\mathcal T_\infty$ denote the actual \emph{unrooted} AF-core infinite-tile pieces that occur in the current fixed-generator tile presentation. The same piece is not listed again merely because a different root is chosen. When a rooted object is needed, it will be written $(T,x)$.
\end{definition}

\begin{proposition}[fixed-generator finite-boundary consequence]
Assume that the tile presentation satisfies the standard generator--boundary dictionary: every persistent generator-boundary point is the source or range of some $\gamma\in E_S$, while a generator germ that lies in the AF core is not counted as a persistent boundary incidence. Define
\[
\partial_S^\infty
:=
\bigsqcup_{T\in\mathcal T_\infty}\partial_S T.
\]
Then
\[
\partial_S^\infty\subseteq s(E_S)\cup r(E_S),
\qquad
|\partial_S^\infty|<\infty.
\]
In particular, the set of boundary-bearing infinite tiles
\[
\mathcal T_\infty^\partial
:=
\{T\in\mathcal T_\infty:\partial_S T\neq\varnothing\}
\]
is finite.
\end{proposition}

\begin{proof}
The set $E_S$ is finite, and both source and range are single-valued maps, so $s(E_S)\cup r(E_S)$ is finite. The generator--boundary dictionary places every persistent boundary point in this finite endpoint set. Since $\mathcal T_\infty$ does not list the same unrooted piece repeatedly, the disjoint union does not create infinitely many artificial copies by changing the root.
\end{proof}

These finitely many boundary points must still be divided according to germ isotropy:
\[
\partial_{S,\mathrm{reg}}^\infty
:=
\{b\in\partial_S^\infty:H_b=\{1\}\},
\qquad
\partial_{S,\mathrm{sing}}^\infty
:=
\{\xi\in\partial_S^\infty:H_\xi\neq\{1\}\}.
\]
Thus
\[
\partial_S^\infty
=
\partial_{S,\mathrm{reg}}^\infty
\sqcup
\partial_{S,\mathrm{sing}}^\infty,
\qquad
|\partial_{S,\mathrm{sing}}^\infty|<\infty.
\]
The fixed-generator finite-support consequence of \AFBD{} counts \emph{all} boundary points. Finiteness of the singular subset is only a corollary.

\begin{warning}[This chapter uses only the forward implication]
Under the generator--boundary dictionary above, the statement proved here is
\[
\boxed{
\begin{gathered}
\text{open AF core + discrete complement + compact generators}\\[-1mm]
\Longrightarrow\\[-1mm]
\text{global finite persistent generator-boundary support}.
\end{gathered}}
\]
Finite boundary support by itself does not automatically imply that the entire non-AF complement is discrete in the groupoid topology. If a particular tile theorem supplies a converse, its continuity, generation, and presentation hypotheses must be stated separately. Later chapters will not turn this forward consequence into an unconditional equivalence.
\end{warning}

\section{How Tile Inflation Produces the Ambient Groupoid}

Now consider the tile inflation of the preceding chapter. Let the finite symmetric label set satisfy $S=S^{-1}$. After the required continuous extensions have been made at the relevant infinite boundary points, each label $F\in S$ defines a partial homeomorphism between a clopen domain and a clopen range. These partial homeomorphisms generate an inverse semigroup
\[
\mathcal G_0=\langle S\rangle,
\]
and hence a groupoid of germs
\[
\mathfrak G
=
\operatorname{Germ}\bigl(\mathcal G_0\curvearrowright\Omega(\mathsf B)\bigr).
\]
The same symbol $F$ is often used for the label, the partial homeomorphism, and the compact-open bisection of germs, but the three roles must be distinguished conceptually: the label belongs to the presentation, the partial homeomorphism is a local action, and the bisection is a set of arrows in the groupoid.

Tile inflation produces an ambient groupoid of germs; it does not automatically produce an \AFBD{} groupoid. Two additional facts must be verified: the tail groupoid must indeed sit inside it openly, and the remaining germs must form a discrete complement. Connected finite tiles automatically provide the first fact.

\begin{proposition}[the tail groupoid is an open subgroupoid]
Assume that all finite tiles in the preceding tile inflation are connected, that the labels have been extended to partial homeomorphisms at the required boundary points, and let
\[
\mathfrak G
=
\operatorname{Germ}\bigl(\mathcal G_0\curvearrowright\Omega(\mathsf B)\bigr).
\]
Then
\[
\mathfrak T_{\mathsf B}\subseteq\mathfrak G
\]
is an open subgroupoid \citep{kuang2026b}.
\end{proposition}

\begin{proof}
Take finite paths $\alpha$ and $\beta$ of the same length $n$ that terminate at the same Bratteli vertex. They belong to the same level-$n$ tile. Since the tile is connected, there is a finite path along internal labeled edges from $\beta$ to $\alpha$. Let $F$ be the composition of the labels along that path. An internal tile edge changes only a finite prefix and leaves the compatible infinite tail unchanged, so for every compatible tail $z$,
\[
F(\beta z)=\alpha z.
\]
Hence the basic tail bisection
\[
Z(\alpha,\beta)
=
\{(\alpha z,\beta z):z\text{ compatible}\}
\]
is contained in $\mathfrak G$. The preceding chapter showed that these sets $Z(\alpha,\beta)$ form a topology basis for $\mathfrak T_{\mathsf B}$, so the entire tail groupoid is an open subgroupoid of $\mathfrak G$.
\end{proof}

Thus, for tile inflation to give an \AFBD{} presentation, the genuine ``by-discrete'' condition still remains:
\[
\boxed{
\mathfrak G\setminus\mathfrak T_{\mathsf B}
\quad\text{must be discrete}.}
\]
This is an additional assumption or a property that must be proved; it cannot be deduced merely from the existence of finite tiles.

\section{What Exactly Is the Internal Part of a Finite Tile? An Elementary AF Cayley Piece}

The finite tile $T_{v,n}$ from the preceding chapter contains both internal edges and boundary edges. We can now answer exactly which part belongs to the level-$n$ AF groupoid.

Write
\[
T_{v,n}^{\circ}
\]
for the graph obtained by deleting all boundary edges. For every internal labeled arrow
\[
\beta\xrightarrow{F}\alpha,
\]
assign the label $F$ to the basic tail bisection $Z(\alpha,\beta)$. Let $\mathcal S_n$ be the finite set of all such bisections. Since the finite tiles are connected, $\mathcal S_n$ generates the elementary subgroupoid $\mathfrak T_{\mathsf B,n}$ from the preceding chapter.

\begin{proposition}[finite tile and elementary Cayley graph]
Fix $v\in V_{n+1}$ and $\gamma_0\in T_{v,n}$, and then fix a compatible infinite tail $z$ beginning at $v$. Let
\[
x=\gamma_0z.
\]
Then
\[
\Phi_x:T_{v,n}^{\circ}\longrightarrow(\mathfrak T_{\mathsf B,n})_x,
\qquad
\Phi_x(\gamma)=(\gamma z,\gamma_0z),
\]
is a rooted labeled graph isomorphism. The graph on the left is rooted at $\gamma_0$, and the graph on the right is the Cayley graph of the elementary groupoid $\mathfrak T_{\mathsf B,n}$ in the source fibre $x$ with respect to $\mathcal S_n$ \citep{kuang2026b}.
\end{proposition}

\begin{proof}
The arrows in the source fibre $(\mathfrak T_{\mathsf B,n})_x$ are precisely
\[
(\gamma z,\gamma_0z),
\qquad
\gamma\in T_{v,n}.
\]
Hence $\Phi_x$ is bijective on vertices and sends $\gamma_0$ to the unit arrow. If the tile contains an internal edge $\beta\xrightarrow{F}\alpha$, the unique arrow in $Z(\alpha,\beta)$ whose source is $\beta z$ is $(\alpha z,\beta z)$, and
\[
(\alpha z,\beta z)(\beta z,\gamma_0z)
=
(\alpha z,\gamma_0z).
\]
Thus the label and Cayley adjacency match exactly. The converse is the same argument in reverse.
\end{proof}

\begin{plainbox}[What this proposition actually fixes]
The internal part of a finite tile is not a loose metaphor for the AF core. After the boundary edges are deleted, it really is, for every compatible tail, the rooted labeled Cayley graph of the finite-stage elementary groupoid. Boundary edges must be interpreted separately; they do not belong to this finite-stage Cayley graph.
\end{plainbox}

\section{The Crucial Distinction: A Finite Boundary Edge Is Not a Non-AF Germ}

This is the easiest point in the chapter to state incorrectly, and the most important one to get right. If an edge is still a boundary edge at level $n$, that says only this: on the current cylinder it cannot yet be represented uniformly by a level-$n$ finite prefix replacement. A connector may internalize it at a deeper level. A genuinely non-AF germ arises only when the edge never becomes internal along a concrete infinite path.

\begin{proposition}[internalization criterion for a boundary edge]
Suppose $\beta\in T_{v,n}$ has an outgoing boundary edge labeled $F$. Choose a compatible tail $z$ and set
\[
y=\beta z\in\operatorname{Dom}(F).
\]
The germ $[F,y]$ gives an $F$-edge in the ambient Cayley graph, but at level $n$ it does not belong to the image of $T_{v,n}^{\circ}$. The following three conditions are equivalent \citep{kuang2026b}:
\begin{enumerate}
\item $[F,y]\in\mathfrak T_{\mathsf B}$;
\item there exist $m\ge n$, a length-$m$ extension $\beta'$ of $\beta$, and another length-$m$ path $\alpha'$ such that
\[
Z(\alpha',\beta')\subseteq F;
\]
\item along the concrete infinite path $y$, after passing to deeper tile inflations, this $F$-boundary edge becomes an internal tile edge at some level $m$.
\end{enumerate}
\end{proposition}

\begin{proof}
At level $n$, the definition of a boundary edge says exactly that $F$ has not yet become a uniform level-$n$ prefix replacement on the entire cylinder $[\beta]$, so it is not represented inside $T_{v,n}^{\circ}$.

If the second condition holds, then on the deeper cylinder $[\beta']$ the map $F$ agrees with the basic tail bisection $Z(\alpha',\beta')$. Hence $[F,y]$ is a tail germ; in tile language, the edge has become internal.

Conversely, suppose $[F,y]\in\mathfrak T_{\mathsf B}$. Since both $F$ and $\mathfrak T_{\mathsf B}$ are open around that germ, the neighborhood can be shrunk to a basic tail bisection $Z(\alpha',\beta')\subseteq F$. Deepen the level if necessary so that $\beta'$ extends $\beta$; this gives the third condition.
\end{proof}

The distinction to remember is therefore
\[
\boxed{
\begin{gathered}
\text{finite-level boundary}\neq\text{non-AF germ},\\[1mm]
\text{persistent generator boundary}
\Longleftrightarrow
\text{germ outside AF core}.
\end{gathered}}
\]

\begin{warning}[Do not rewrite ``not yet internalized'' as ``can never be internalized'']
A finite boundary is a finite-scale state of the presentation; non-AF membership is an infinite-level fact about whether the actual germ belongs to $\mathfrak T_{\mathsf B}$. Every later discussion of cracks, persistent boundary points, singular germs, and graphs of germs depends on this distinction.
\end{warning}

\section{Which Part of the Orbital Graph Is an Infinite Tile?}

Fix $\gamma\in\Omega(\mathsf B)$. Taking the inductive limit of the finite tiles along the truncations of $\gamma$ produces an infinite tile. If we retain only those generator edges that eventually belong to the AF core, write the resulting internal infinite tile as $T_\gamma^{\circ}$. It can be identified directly inside the ambient orbital graph.

\begin{principle}[infinite tile = AF-core orbital piece]
The vertices of $T_\gamma^{\circ}$ are precisely the points in the $\mathfrak T_{\mathsf B}$-orbit of $\gamma$. A generator edge is retained if and only if its germ lies in the AF core. The remaining generator edges occur at persistent boundary points and connect distinct AF infinite-tile pieces \citep{kuang2026b}.
\end{principle}

This gives the most accurate picture of the \AFBD{} tile geometry:
\[
\boxed{
\begin{gathered}
\text{large AF orbital pieces}\\[-1mm]
+\ \text{finitely many persistent-boundary generator interfaces}.
\end{gathered}}
\]
The word ``finitely'' is always relative to the current fixed finite/compact generating family. In abstract groupoid language the non-AF complement is discrete; compact generators turn this into finitely many exceptional generator germs; tile language displays the same finiteness as a globally finite set of persistent boundary points across \emph{all} infinite tiles. Hence only finitely many infinite tiles carry boundary, while the remaining infinite tiles are boundary-free.

This does not say that the whole groupoid contains only finitely many non-AF arrows. Nor does it say that an aggregate over all finite-tile types at every level is already uniformly bounded. The general bounded-type condition introduced below instead controls each finite AF fibre / tile uniformly; in the finite-rank tile-inflation form used here, the aggregate bound follows only after the number of tile types per level is controlled as well.

Moreover, every AF infinite tile lifts isomorphically to a subgraph of the corresponding groupoid Cayley graph. Passing to germ resolution does not destroy the AF piece itself; it only unfolds local continuation data that the orbital quotient had identified. The next chapter studies this resolution explicitly.

\section{The Four Layers of Data Must Remain Permanently Separate}

By this point many different objects have appeared in the tile presentation. The easiest mistake is to call all of them ``boundary data.'' The following table fixes the usage for the remainder of the book.

\begingroup
\footnotesize
\setlength{\tabcolsep}{3.2pt}
\renewcommand{\arraystretch}{1.16}
\begin{longtable}{>{\raggedright\arraybackslash\bfseries}p{0.20\textwidth}>{\raggedright\arraybackslash}p{0.38\textwidth}>{\raggedright\arraybackslash}p{0.32\textwidth}}
\toprule
Object & Mathematical status & Must not be substituted by \\
\midrule
\endfirsthead
\toprule
Object & Mathematical status & Must not be substituted by \\
\midrule
\endhead
$\mathfrak T_{\mathsf B}$ & the designated open tail subgroupoid in the given \AFBD{} presentation & a chosen system of generator labels, or a particular finite tile \\
$\mathfrak T_{\mathsf B,n}$ & a finite-stage elementary subgroupoid of the AF core & an intrinsic ``absolute level $n$'' \\
$T_{v,n}^{\circ}$ & a rooted labeled Cayley piece of $\mathfrak T_{\mathsf B,n}$ & the entire ambient groupoid \\
finite boundary edge & a presentation datum not yet internalized by a uniform prefix replacement at the current level & automatically a non-AF germ \\
connector & finite combinatorial data describing how higher-level inflation joins concrete tile copies & an independently existing class of arrows in the groupoid \\
persistent boundary germ & a local-action class that never enters the AF core along a concrete infinite path & an arbitrary finite boundary occurrence \\
chosen labels & a presentation generating partial homeomorphisms / bisections & intrinsic labels of the AF core \\
\bottomrule
\end{longtable}
\endgroup

The table also explains why the preceding chapter insisted on occurrence-first bookkeeping. An abstract boundary type by itself cannot determine whether a concrete germ is ultimately AF; one must follow a concrete occurrence and its concrete infinite continuation to see whether it becomes internal at a deeper level.

\section{Bounded Type: General Groupoid Form and the Tile-Inflation Form Used Here}

The preceding sections fixed the \AFBD{} picture used in this book: an AF core remains open, the non-AF complement is discrete, and in the compactly generated fixed-generator tile presentation the persistent generator-boundary support across the infinite-tile limit is finite. The phrase \emph{bounded type} will be used at two levels that should remain explicitly distinguished.

\begin{definition}[groupoid of bounded type: general formulation]
Let $\mathcal G$ be an étale groupoid with Cantor unit space. Following \citet[Definition~5.2.16]{nek2022}, we say that $\mathcal G$ is a \emph{groupoid of bounded type} if there exists an increasing sequence
\[
\mathfrak F_0\subseteq \mathfrak F_1\subseteq\cdots
\]
of elementary subgroupoids such that
\[
\mathcal G\setminus\bigcup_{n\ge0}\mathfrak F_n
\]
is discrete and, for every compact subset $C\subseteq\mathcal G$, there exists $K_C<\infty$ satisfying
\[
\boxed{
\bigl|C\mathfrak F_n x\setminus \mathfrak F_n x\bigr|
\le K_C
\qquad
\text{for every }x\in\mathcal G^{(0)}\text{ and every }n.}
\]
For a compactly generated AF-by-discrete groupoid, it is enough to verify the uniform estimate for one compact symmetric generating set $S$ \citep[Proposition~5.2.17]{nek2022}.
\end{definition}

In the Bratteli/tile presentation used in this book, the elementary fibres $\mathfrak F_nx$ are represented by finite AF tiles. Under the generator--boundary dictionary already established, the fixed-generator estimate is therefore the \emph{per-fibre / per-finite-tile} statement
\[
\boxed{
\sup_n\;
\sup_{T\in\mathcal T_n}
|\partial_S T|
<\infty,}
\]
where $\mathcal T_n$ denotes the finite tile types at level $n$. This is the finite-level uniformity supplied by the general groupoid definition. It is not, by itself, the sum of the boundary cardinalities of every tile type at one level.

\begin{definition}[bounded-type tile-inflation form used in the author's work]
In the tile-inflation formulation used in the author's work \citep[Definition~2.13]{kuang2025tiling}\citep[Definition~3.1]{kuang2026b}, bounded type is recorded by the following package:
\begin{enumerate}
\item the boundary cardinalities of all finite tiles are uniformly bounded;
\item the infinite tiles, when present, have only finitely many boundary points;
\item the cardinalities $|V_n|$ and $|E_n|$ are uniformly bounded in $n$.
\end{enumerate}
When the book works directly with tile inflations, this is the presentation-level form that will be used.
\end{definition}

The third clause is additional combinatorial control on the presentation. It has an important consequence for the aggregate burden across tile types. Define
\[
L_n(S)
:=
\sum_{T\in\mathcal T_n}|\partial_S T|.
\]
If
\[
B:=\sup_n\sup_{T\in\mathcal T_n}|\partial_S T|<\infty
\qquad\text{and}\qquad
M:=\sup_n|\mathcal T_n|<\infty,
\]
then
\[
\boxed{L_n(S)\le MB\quad\text{for every }n.}
\]
Thus the \emph{level-total} intuition used in earlier drafts is valid in this finite-rank tile-inflation setting as a \textbf{derived consequence}: the boundary burden of each tile is uniformly bounded and the number of tile types per level is uniformly controlled. The additional $|E_n|$ bound belongs to the tile-inflation package as well, although it is not needed for this particular inequality.

\begin{plainbox}[The distinction in one display]
\[
\boxed{
\begin{array}{ll}
\text{general groupoid bounded type:}
&
\text{discrete remainder + uniform per-fibre boundary burden;}\\[1mm]
\text{tile-inflation form used here:}
&
\text{uniform finite-tile boundaries}\\
&
+\text{ finite infinite-tile boundary support when present}\\
&
+\text{ uniformly bounded }|V_n|,|E_n|.
\end{array}}
\]
A level-total boundary bound is not taken as the general definition. In the finite-rank tile-inflation form, it follows from the first and third clauses.
\end{plainbox}

\begin{principle}[label-preserving connector grammar]
Numerical control is not the whole story. For every generator label $F\in S$, a finite connector may pair only two compatible boundary interfaces with the \textbf{same label $F$ and complementary missing-endpoint directions}. The pairing produces a genuine $F$-arrow, and the reverse arrow automatically carries $F^{-1}$. Thus the phrase ``a large system organized by finitely many boundary interfaces'' always contains two pieces of information:
\[
\boxed{
\text{finite interface data}
+\text{ label-preserving gluing rules}.}
\]
Persistent boundary germs are the local-action classes of those boundary histories that remain outside the AF core in the infinite limit; a connector is finite-level presentation data. They share generator-label grammar, but they are not literally the same kind of object.
\end{principle}

\begin{warning}[Keep the general groupoid statement and the tile-inflation package distinct]
For the general groupoid definition, use the per-fibre estimate above. For tile-inflation arguments in this book, also retain finite boundary support on infinite tiles when they occur and the uniform $|V_n|,|E_n|$ controls. A level-total boundary bound may be invoked only as a derived consequence after the number of tile types per level has been controlled.
\end{warning}

\section{Chapter Summary: What Is the Real Division of Labor in ``AF + Discrete''?}

The chapter can be compressed into seven mathematical facts that must be carried forward:
\begin{enumerate}
\item The abstract \AFBD{} definition is an open AF core plus a discrete complement. Under a fixed finite/compact generating family and the generator--boundary dictionary of Proposition~10.6, it implies that the total number of persistent generator-boundary points across all infinite tiles is finite, and hence that only finitely many infinite tiles carry boundary.
\item ``Finitely many interfaces'' does not mean that the entire non-AF part is globally finite. It is a statement about orbital/tile organization relative to the chosen generators.
\item Connected tile inflation automatically embeds the tail groupoid openly into the ambient groupoid of germs, but the ``by-discrete'' condition still has to be verified.
\item After boundary edges are deleted, a finite tile is not an analogy: it is an actual Cayley piece of the finite-stage elementary AF groupoid.
\item A finite boundary edge means only that the edge has not yet been internalized at the current scale. It becomes a persistent non-AF interface in the fixed-generator picture only if it remains a boundary edge forever along the concrete infinite path.
\item A connector preserves the generator-label grammar: compatible boundary edges must carry the same label and have complementary missing-endpoint directions.
\item General groupoid bounded type adds uniform boundary control for each finite AF fibre / tile. In the tile-inflation form used here, uniform $|V_n|,|E_n|$ supplies the additional finite-rank control from which a whole-level aggregate bound follows.
\end{enumerate}

\enlargethispage{2\baselineskip}
The next chapter increases the resolution. An orbital graph identifies local actions that have the same endpoint, whereas a groupoid Cayley graph and a graph of germs also record \emph{how} that endpoint was reached. Thus a finite pattern that has already reached completion downstairs may still retain nontrivial displacement at a finer germ resolution.

\chapter{Orbital Graph, Cayley Graph, and Graph of Germs}
\label{chap:orbital-cayley-germs-en}

\par\medskip
\noindent\textbf{Why this chapter matters}\par
The preceding chapter wrote an \AFBD{} groupoid as ``an AF core plus discrete non-AF germs.'' Looking only at the groupoid itself, however, does not yet make one crucial distinction visually obvious: two local actions may send the same starting point to the same endpoint and still be different arrows.

Here we draw that distinction with three related graphs that retain different amounts of information. The orbital graph records only ``where one finally arrived.'' The groupoid Cayley graph records ``through which local arrow one arrived.'' When the groupoid comes from germs of a group action, this Cayley graph is precisely what is usually called the graph of germs. All later phenomena in which a pattern has already formed a closed walk in the base graph while a finer lift still lands on another sheet rest on the covering structure established in this chapter.
\par\medskip

\par\medskip
\noindent\textbf{Mathematical sources for this chapter}\par
This chapter uses the Cayley-graph language of compactly generated \(\acute{e}\)tale groupoids and follows Nekrashevych's groupoid geometry together with the notation of the author's Near-full Part I \citep{nek2022,kuang2026b}. In particular, the groupoid Cayley graph, the orbital graph, and the covering property of the range map come from that standard framework. In the special case of a group action, the vertices of the groupoid Cayley graph are the germs of the action, so the Cayley graph becomes the graph of germs; its fibres over the orbital graph are controlled by the local group of germs. For concrete background on bounded-type group actions and orbital-graph methods, see also Cantu \citep{cantu2019}.
\par\medskip

\section{Why an Orbital Graph Is Not Yet Enough}

Let \(\mathfrak G\) be a compactly generated \(\acute{e}\)tale groupoid with unit space \(X\). Fix a finite family of compact-open bisections
\[
\mathcal S=\mathcal S^{-1}
\]
that generates \(\mathfrak G\). Here \(\mathcal S=\mathcal S^{-1}\) means that the family is invariant under inversion.

For \(F\in\mathcal S\), the symbol \(F\) is not merely an abstract letter. It is a compatible family of local arrows. If \(y\in\mathbf s(F)\), then because \(F\) is a bisection there is a unique arrow
\[
\eta_F(y)\in F,
\qquad
\mathbf s(\eta_F(y))=y.
\]

Now fix \(x\in X\). There are two different questions one may ask.

First, which points can be reached from \(x\)? This gives the orbit
\[
\mathcal O(x)
=
\{\mathbf r(\gamma):\gamma\in\mathfrak G,\ \mathbf s(\gamma)=x\}.
\]

Second, which concrete arrows start at \(x\)? This gives the source fibre
\[
\mathfrak G_x:=\mathbf s^{-1}(x).
\]

There is a natural forgetful map
\[
\mathbf r:\mathfrak G_x\longrightarrow\mathcal O(x),
\qquad
\gamma\longmapsto\mathbf r(\gamma).
\]
It forgets \emph{how} one arrived and remembers only \emph{where} one arrived. If two different arrows have the same range, this map genuinely discards information.

\begin{plainbox}[In ordinary language]
Suppose two different local actions both carry \(x\) to \(y\). The orbital graph identifies them at the single vertex \(y\). The Cayley graph introduced next retains two different vertices, because the arrows themselves have not yet been forgotten.
\end{plainbox}

\section{The Groupoid Cayley Graph: Its Vertices Are Arrows from the Root, Not Points of the Space}

\begin{definition}[groupoid Cayley graph]
Fix \(x\in X\). The rooted labeled Cayley graph
\[
\mathcal C(x,\mathcal S)
\]
has vertex set \(\mathfrak G_x\) and root \(1_x\). For \(\gamma\in\mathfrak G_x\) and \(F\in\mathcal S\) satisfying \(\mathbf r(\gamma)\in\mathbf s(F)\), let \(\eta\in F\) be the unique arrow with
\[
\mathbf s(\eta)=\mathbf r(\gamma).
\]
Place a directed edge labeled \(F\)
\[
\gamma\xrightarrow{F}\eta\gamma.
\]
\end{definition}

\begin{plainbox}[How to read this definition]
Starting from the root, every step composes on the left with the local arrow selected from a generating bisection. A vertex
\[
\gamma:x\longrightarrow y
\]
therefore remembers not only the endpoint \(y\), but also the genuine groupoid arrow by which that endpoint has been reached.
\end{plainbox}

Because each \(F\) is a bisection, at every vertex there is at most one outgoing and at most one incoming edge carrying the label \(F\). Since \(\mathcal S\) generates \(\mathfrak G\), the Cayley graph is connected.

There is also a useful change-of-root observation. If
\[
\gamma:x\longrightarrow y
\]
is an arrow of the groupoid, then right multiplication by \(\gamma\) gives an isomorphism of unrooted labeled graphs
\[
\mathcal C(y,\mathcal S)
\longrightarrow
\mathcal C(x,\mathcal S),
\qquad
\delta\longmapsto\delta\gamma,
\]
sending \(1_y\) to \(\gamma\). Thus the unrooted Cayley graph depends only on the orbit; choosing a different unit in the same orbit merely chooses a different root.

\section{The Orbital Graph: Retaining Only the Range}

\begin{definition}[orbital graph]
Fix \(x\in X\). The rooted labeled orbital graph
\[
\Gamma(x,\mathcal S)
\]
has vertex set \(\mathcal O(x)\) and root \(x\). For \(F\in\mathcal S\) and
\[
y\in\mathbf s(F)\cap\mathcal O(x),
\]
place a directed edge
\[
y\xrightarrow{F}\mathbf r(\eta_F(y)).
\]
\end{definition}

Three objects must therefore be kept separate:
\begin{itemize}
\item for a finitely generated group \(G\), the ordinary Cayley graph has vertices \(g\in G\);
\item for a groupoid, the Cayley graph rooted at \(x\) has vertices given by arrows with source \(x\);
\item the orbital graph has vertices given by points in the orbit of \(x\).
\end{itemize}
Only in special situations do these graphs coincide.

\section{The Range Map Is a Covering: Locally the Same, Globally on Different Sheets}

The most important standard fact is that the forgetful map above is not arbitrary. It is a covering of labeled directed graphs \citep{nek2022,kuang2026b}.

\begin{proposition}[the Cayley graph covers the orbital graph]
The range map
\[
\mathbf r:\mathcal C(x,\mathcal S)
\longrightarrow
\Gamma(x,\mathcal S)
\]
is a covering of labeled directed graphs. At every vertex it induces bijections, label by label, on both outgoing and incoming edges.
\end{proposition}

\begin{proof}
Let \(\gamma\in\mathfrak G_x\) and put \(y=\mathbf r(\gamma)\). If \(F\in\mathcal S\) is defined at \(y\), the bisection property gives a unique \(\eta\in F\) with \(\mathbf s(\eta)=y\). Hence there is a unique \(F\)-edge in the Cayley graph
\[
\gamma\xrightarrow{F}\eta\gamma.
\]
Its image under the range map is precisely the \(F\)-edge
\[
y\xrightarrow{F}\mathbf r(\eta)
\]
in the orbital graph. The same argument, using the inverse bisection, applies to incoming edges. Thus the labeled incoming and outgoing stars are carried bijectively to those downstairs.
\end{proof}

This covering explains a point that is very easy to state incorrectly: \textbf{the Cayley graph does not add ``more locally available directions.''} In the covering sense, the local labeled neighborhood on every sheet is exactly the same as the one downstairs. What the Cayley graph adds is \emph{sheet position}: the same subsequent label sequence, started from different lifts, can remain represented by different arrows throughout.

\section{The Fibre Is Controlled by Isotropy}

Fix the root \(x\). Its isotropy group is
\[
\mathfrak G_x^x
=
\{h\in\mathfrak G:\mathbf s(h)=\mathbf r(h)=x\}.
\]
For \(h\in\mathfrak G_x^x\), right multiplication
\[
R_h:\mathfrak G_x\longrightarrow\mathfrak G_x,
\qquad
\gamma\longmapsto\gamma h
\]
preserves labels and range, and therefore defines a covering automorphism of the Cayley graph.

\begin{proposition}[isotropy fibres]
The right-multiplication action of \(\mathfrak G_x^x\) on \(\mathcal C(x,\mathcal S)\) is free, and its orbits are exactly the fibres of the range map. Hence
\[
\mathcal C(x,\mathcal S)/\mathfrak G_x^x
\cong
\Gamma(x,\mathcal S).
\]
\end{proposition}

\begin{proof}
If \(\gamma h=\gamma\), cancellation of \(\gamma\) gives \(h=1_x\), so the action is free. If
\[
\mathbf r(\gamma_1)=\mathbf r(\gamma_2),
\]
then
\[
h=\gamma_1^{-1}\gamma_2\in\mathfrak G_x^x
\]
and \(\gamma_1h=\gamma_2\). Thus any two points in the same fibre are related by a unique isotropy element.
\end{proof}

\begin{plainbox}[In ordinary language]
A vertex \(y\) of the orbital graph may have several lifts upstairs. Every one of them represents ``an arrow from \(x\) to \(y\).'' They differ only in which isotropy class has been accumulated by the time \(y\) is reached. Quotienting these lifts once more gives the orbital graph.
\end{plainbox}

\section{The Special Case of a Group Action: The Cayley Graph Is the Graph of Germs}

Now let a finitely generated group \(G=\langle S\rangle\) act on \(X\), and let
\[
\mathfrak G=\operatorname{Germ}(G\curvearrowright X).
\]
Fix \(x\in X\). Recall
\[
G_x=\{g\in G:g(x)=x\},
\]
and
\[
G_{(x)}
=
\left\{g\in G:
\begin{gathered}
\text{there exists a neighborhood }U\text{ of }x\text{ such that}\\
g|_U=\operatorname{id}
\end{gathered}
\right\}.
\]
The local group of germs is therefore
\[
H_x=G_x/G_{(x)}.
\]

In this situation the vertices of the groupoid Cayley graph may be written as germs
\[
[g,x].
\]
The graph is therefore usually denoted
\[
\widetilde\Gamma_x
\]
and called the \emph{graph of germs at \(x\)}. The orbital graph is still denoted \(\Gamma_x\).

With right-action notation, these two graphs are the Schreier graphs
\[
\widetilde\Gamma_x
\cong
\operatorname{Sch}(G_{(x)}\backslash G,S),
\qquad
\Gamma_x
\cong
\operatorname{Sch}(G_x\backslash G,S).
\]
Thus the fibre of the natural quotient
\[
G_{(x)}\backslash G
\longrightarrow
G_x\backslash G
\]
is indexed precisely by
\[
H_x=G_x/G_{(x)}.
\]

This is also why the later chapters connect singularity with graph-of-germs resolution. A singular point is not singular because the orbital graph acquires an additional ``strange vertex.'' It is singular because the range projection at that root has a nontrivial germ fibre.

\section{Downstairs Loop and Upstairs Displacement}

Consider a finite labeled walk in the orbital graph that begins and ends at \(x\):
\[
x=y_0\xrightarrow{F_1}y_1
\xrightarrow{F_2}\cdots
\xrightarrow{F_m}y_m=x.
\]
Because
\[
\mathbf r:\mathcal C(x,\mathcal S)\to\Gamma(x,\mathcal S)
\]
is a covering, there is a unique lift starting at the root \(1_x\):
\[
1_x=\gamma_0\xrightarrow{F_1}\gamma_1
\xrightarrow{F_2}\cdots
\xrightarrow{F_m}\gamma_m.
\]
Since the downstairs walk ends at \(x\),
\[
\gamma_m\in\mathfrak G_x^x.
\]

\begin{principle}[A closed walk and a trivial lift are two different things]
The fact that a walk in the orbital graph is closed implies only that the endpoint of its lift lies in the isotropy fibre:
\[
\text{closed downstairs}
\quad\Longrightarrow\quad
\gamma_m\in\mathfrak G_x^x.
\]
Only when
\[
\gamma_m=1_x
\]
does the lift also end at its original root in the Cayley graph. Therefore
\[
\text{closed downstairs}
\quad\not\Longrightarrow\quad
\text{closed at the same upstairs root}.
\]
\end{principle}

\begin{plainbox}[In ordinary language]
Downstairs asks only whether the spatial position has come back to \(x\). Upstairs asks one further question: when the path reaches \(x\) again, is the local arrow class it carries the identity class? These two questions can have different answers.
\end{plainbox}

For a group action, \(\gamma_m\) is a germ class
\[
[h,x]\in H_x.
\]
Thus an orbital loop can close completely in the base graph while carrying the root of the graph of germs from the identity sheet to another germ sheet. When later chapters keep finer continuation data, they are continuing the bookkeeping built on this basic covering fact.

\section{Where Does an Infinite Tile Sit in an AF-by-Discrete Tile Presentation?}

The preceding chapter distinguished a finite boundary edge from a persistent non-AF germ. The three graphs now make that distinction more precise.

Fix an infinite path \(x\). Let \(T_x^{\circ}\) denote the infinite tile obtained from tile inflation after deleting the boundary edges that still point outside. Under the hypotheses of the preceding chapter, \(T_x^{\circ}\) records exactly the portion of the orbital geometry reachable through AF-core generator germs: its vertices belong to the tail orbit, and the germs of its internal generator edges lie in \(\mathfrak T_{\mathsf B}\). Near-full Part I further shows that each such infinite tile lifts isomorphically to a subgraph of the ambient groupoid Cayley graph \citep{kuang2026b}.

Thus the following three layers must remain permanently distinct:
\[
\boxed{
\begin{gathered}
\text{finite tile / infinite tile in the orbital presentation}\\
\downarrow\ \text{lift}\\
\text{AF-core Cayley piece}
\subseteq
\text{full groupoid Cayley graph}.
\end{gathered}}
\]

Persistent boundary generator edges join distinct infinite AF pieces in the ambient orbital graph. A finite-level boundary edge that becomes internal at a deeper level still belongs to AF geometry. Only a germ that is never internalized along the concrete infinite path lies genuinely outside the AF core.

\section{Boundary Is Not the Same as Singular: Three Infinite-Tile / Graph-of-Germs Cases}

The global finiteness condition of the preceding chapter says that the total set of persistent boundary points across all infinite tiles is finite. The next question is local: is the group of germs at such a point trivial or nontrivial?

Write
\[
\mathcal T_\infty^\partial
:=
\{T\in\mathcal T_\infty:\partial_S T\neq\varnothing\}.
\]
This set is finite. But a boundary-bearing tile is not automatically a ``singular tile.'' Three cases must be separated.

\paragraph{First case: a boundary-free infinite tile.}
Suppose \(T\) has no boundary. For every root \(x\in T\), there is no persistent generator boundary and \(H_x\) is trivial. Hence, as rooted labeled graphs,
\[
\widetilde\Gamma_x
\cong
\Gamma_x
\cong
(T,x).
\]
The tile, the orbital graph, and the graph of germs coincide.

\paragraph{Second case: a regular boundary point.}
Let
\[
b\in\partial_{S,\mathrm{reg}}^\infty,
\qquad
H_b=\{1\}.
\]
The point \(b\) participates in a persistent cross-tile generator connection, so the rooted orbital graph may differ from the single tile that produced \(b\). But there is no branching in the isotropy fibre:
\[
\widetilde\Gamma_b\cong\Gamma_b.
\]
In the simplest case, a boundary pair carrying the label \(F\) connects \(b\) to a boundary point \(b'\) in another tile, and the orbital graph---hence also the graph of germs---glues the two infinite tiles by reversible arrows
\[
b\xrightarrow{F}b',
\qquad
b'\xrightarrow{F^{-1}}b.
\]
This is tile gluing, not nontrivial germ-sheet resolution.

\begin{example}[the regular boundary of the binary odometer]
View the binary odometer as the adding-one action on the binary path space. The generator sends
\[
1^\omega\longmapsto0^\omega.
\]
This is a genuine persistent boundary transition. But the odometer action is free, so the stabilizer at the relevant point, and therefore its group of germs, is trivial. Thus \(1^\omega\) (and the corresponding point at the other end) is a regular boundary point: the graph of germs has no more sheets than the orbital graph. It merely joins adjacent infinite-tile pieces through this pair of reversible labeled boundary arrows.
\end{example}

\paragraph{Third case: a singular boundary point.}
If
\[
\xi\in\partial_S^\infty,
\qquad
H_\xi\neq\{1\},
\]
then \(\xi\) is a singular boundary point in the strict sense used in this book. Now
\[
\widetilde\Gamma_\xi\longrightarrow\Gamma_\xi
\]
has a nontrivial germ fibre. The orbital quotient compresses distinct local germ classes into the same orbital position, while the graph of germs unfolds those sheets again. From the tile viewpoint, this is the \emph{singular-germ resolution} used repeatedly later in the book: rooted copies of finitely many source infinite tiles may be glued according to germ continuation into one resolved graph, and the number of copies appearing inside one resolved graph may itself be finite or infinite.

Because
\[
\partial_{S,\mathrm{sing}}^\infty
\subseteq
\partial_S^\infty
\]
and the latter set is finite, only finitely many singular roots produce nontrivial germ sheets. But one warning is essential: \textbf{``the graph of germs differs from a single producing tile'' is not equivalent to singularity.} Regular boundary gluing can already make an orbital graph, and therefore a graph of germs, cross between two tiles. The extra content of singularity is that \(\widetilde\Gamma_\xi\) has germ sheets beyond those visible in \(\Gamma_\xi\).

\begin{plainbox}[The three cases in one display]
\[
\boxed{
\begin{array}{lll}
\partial_S T=\varnothing
&\Rightarrow&
\widetilde\Gamma_x\cong\Gamma_x\cong(T,x);\\[1mm]
b\in\partial_{S,\mathrm{reg}}^\infty
&\Rightarrow&
\begin{gathered}\widetilde\Gamma_b\cong\Gamma_b,\text{ but }\Gamma_b\text{ may be glued from tiles}\\\text{along reversible boundary arrows}\end{gathered};\\[1mm]
\xi\in\partial_{S,\mathrm{sing}}^\infty
&\Rightarrow&
\widetilde\Gamma_\xi\to\Gamma_\xi\text{ has nontrivial germ fibre / sheets.}
\end{array}}
\]
Thus the global finite condition in \AFBD{} is \emph{finite boundary support}. The central discussions in this book of Third-Order observation, nontrivial residue, the pNH benchmark, and heterogeneous germ sheets use only the \emph{singular boundary germs} within that finite support.
\end{plainbox}

\section{The Analytic-Continuation / Riemann-Surface Analogy: Retain Only the Covering Relation}

Only now, in the third case of a singular boundary point, does a natural analogy appear. On a Riemann surface, a closed loop in the base space may lift to another sheet. Here, a closed labeled walk in the orbital graph may likewise lift to the graph of germs from the identity sheet to a nontrivial isotropy sheet. Regular boundary gluing has no such nontrivial sheet displacement.

The analogy retains only the following formal relations:
\[
\text{base quotient}
\quad\longleftarrow\quad
\text{resolved covering space},
\]
and
\[
\text{closed base walk}
\quad\leadsto\quad
\text{possibly non-closed lift at the chosen sheet}.
\]

\begin{stopbox}[Do not upgrade the analogy into a mechanism]
A graph of germs is not a Riemann surface, and germ continuation is not complex analytic continuation. The former arises from equivalence classes and covering geometry in an \(\acute{e}\)tale groupoid / local action; the latter arises from the theory of analytic continuation. Later chapters borrow only the structural intuition of ``base and sheets'' and ``closed loop and lift,'' not the analytic mechanism.
\end{stopbox}

\section{The Orbit Principle for Singular Germs: Same Orbit Forces Isomorphism; Different Orbits May Be Heterogeneous}

A fact that is easy to state either too strongly or too weakly must now be fixed precisely. For any groupoid, the isomorphism type of the isotropy group is constant along an orbit. If \(x,y\in X\) lie in the same \(\mathfrak G\)-orbit, then by definition there is an arrow
\[
\gamma:x\longrightarrow y.
\]
Conjugation gives a group isomorphism
\[
\operatorname{Ad}_\gamma:\mathfrak G_x^x
\longrightarrow
\mathfrak G_y^y,
\qquad
h\longmapsto\gamma h\gamma^{-1}.
\]
Thus, for a groupoid of germs,
\[
\boxed{
x\text{ and }y\text{ lie in the same orbit}
\quad\Longrightarrow\quad
H_x\cong H_y.}
\]
In particular, two singular points in the same orbit necessarily have isomorphic groups of germs.

Here ``isomorphic'' must not be heard as ``canonically identified without any choice.'' If another arrow \(\gamma':x\to y\) is chosen, the resulting conjugation map generally differs from \(\operatorname{Ad}_\gamma\) by an inner automorphism of the target isotropy group. The choice-independent conclusion is that the \emph{isomorphism type of the group of germs is an orbit invariant}.

Conversely, if two singular points belong to different orbits, there is no groupoid arrow joining them, so the groupoid structure itself does not force their local groups of germs to be isomorphic; they may indeed be non-isomorphic. The correct heterogeneity statement is therefore
\[
\boxed{
\begin{gathered}
\text{different-orbit singular sectors may carry}\\
\text{non-isomorphic groups of germs}
\end{gathered}}
\]
not the unrestricted claim that ``any two different singular points may have non-isomorphic groups of germs.''

\begin{warning}[finite boundary geometry cannot by itself prove ``same orbit'']
A finite path inside one finite tile may connect two boundary occurrences. This does not automatically imply that there is a groupoid arrow between the corresponding limiting singular points. Only a genuine orbit relation between the infinite points triggers isotropy conjugacy. Finite-level incidence, same-orbit relation, and isomorphism of groups of germs must therefore be recorded separately.
\end{warning}

This orbit principle becomes important when later chapters handle multiple singularities. Part II retains the graph of germs and local isotropy of each singular point separately, while recording that local groups of germs within the same orbit have the same isomorphism type. A genuinely heterogeneous germ profile can occur across different orbits. Further integration across orbits is left to the future problems in Part VI.

\section{Chapter Summary: What Is Lost at Each of the Three Forgetful Steps?}

For a groupoid of germs arising from a group action, the chapter can be compressed into three layers of data:
\[
\boxed{
\begin{array}{ccc}
\text{ordinary group word / presentation}
&\longrightarrow&
\text{groupoid arrow or germ}
\\[1mm]
&&\downarrow\ \mathbf r
\\[-1mm]
&&
\text{orbit point}.
\end{array}}
\]
For a group action, more concretely,
\[
G
\longrightarrow
G_{(x)}\backslash G
\longrightarrow
G_x\backslash G.
\]

The first step forgets distinctions among group words that are already locally equivalent. The second forgets distinct germs that arrive at the same orbit point. The graph of germs sits precisely between them: it no longer retains irrelevant global names, but it still retains the local arrow class that is invisible to the orbital graph.

\enlargethispage{4\baselineskip}
Before entering the next chapter, four facts must therefore remain fixed:
\begin{enumerate}
\item vertices of the orbital graph are orbit points;
\item vertices of the groupoid Cayley graph are arrows starting from the root;
\item for the groupoid of germs of a group action, this Cayley graph is the graph of germs;
\item a base loop may be closed while its lift from a specified sheet still ends with nontrivial isotropy displacement.
\end{enumerate}

The next chapter will define special pieces of concrete labeled walks on these graphs and will introduce, in order, trajectory, traverse, return, and incompressibility. At that point, ``closing'' will acquire a more restricted meaning than an arbitrary graph loop.

\chapter{Trajectory, Traverse, Return, and Incompressibility: Basic Definitions}
\label{chap:trajectory-traverse-return-en}

\par\medskip
\noindent\textbf{Why this chapter matters}\par
Only here, near the end of Part II, do we finally reach four terms that will recur throughout the rest of the book: \emph{trajectory}, \emph{traverse}, \emph{return}, and \emph{incompressibility}. They may all look like ways of talking about ``walking for a while,'' but they perform completely different mathematical tasks.

The easiest mistakes are to call the reappearance of the same local graph a return, to call an arbitrary closed walk a return, or to infer incompressibility directly from the fact that one particular lift fails to come back to the identity sheet. This chapter rebuilds the relations in dependency order. First, a concrete word actually traces a trajectory from a concrete starting point. A finite-tile boundary then cuts that trajectory into traverses. A return is the special case in which the two ends land at the same boundary occurrence. Only after those notions are in place does incompressibility become a factor-avoidance condition on the whole word.
\par\medskip

\par\medskip
\noindent\textbf{Mathematical sources and the bookkeeping introduced in this manuscript}\par
The definitions of trajectory, the original return word, and incompressibility come from Kuang, \emph{Growth of groups with incompressible elements, I}, Definitions~3.9--3.11; the finite-tile traverse comes from Definition~4.1 of the same paper \citep{kuang2026a}. This growth-theoretic line of work also inherits the traverse-counting method developed by Bartholdi--Nekrashevych--Zheng for linear Schreier graphs \citep{bnz2022}. The return/incompressibility definitions and the source-expression convention used here, however, follow Kuang's bounded-type formulation.

In the original paper, a return word includes the generator labels at the beginning and end that enter and leave the boundary. For the multiscale witness theory developed later in this book, we additionally give a separate name, \emph{finite return witness}, to the finite segment inside the tile that begins at a boundary position and comes back to that same position. This is occurrence-first bookkeeping added around the original definition; it must not be presented as terminology used verbatim in the source paper.
\par\medskip

\section{Why Word, Trajectory, and Occurrence Must Be Separated First}

From this chapter onward, we use right-action notation to write words in the generators; that is, the generators are read from left to right. Let
\[
S=S^{-1}
\]
be a finite symmetric labeling set. A formal word
\[
w=F_1F_2\cdots F_m\in S^*
\]
is, by itself, only a string of labels. Since the labels usually come from partial homeomorphisms, the same word need not be executable from every vertex.

\begin{definition}[admissible word at a starting vertex]
Let \(v\) be a vertex in an orbital/tile graph. We call the word \(w\) \emph{admissible at \(v\)} if the successive actions
\[
v_i:=v\cdot F_1F_2\cdots F_i,
\qquad 1\le i\le m,
\]
are defined at every step; equivalently, these labels are composable at that occurrence.
\end{definition}

This occurrence-first discipline will remain in force from now on. Return, persistence, and residue belong first to trajectory occurrences that actually take place. Retaining only an abstract word or type forgets where it occurred, which boundary context it passed through, and whether it was genuinely admissible there.

\section{Trajectory: The Finite Route Actually Traced by a Generator Word}

\begin{definition}[trajectory / walk]\label{def:trajectory-basic-en}
Let \(w=F_1\cdots F_m\) be admissible at a vertex \(v\). We call the sequence
\[
\operatorname{Traj}_v(w)
=(v_0,v_1,\ldots,v_m),
\qquad
v_0=v,
\qquad
v_i=v\cdot F_1\cdots F_i
\]
the \emph{trajectory} or \emph{walk} of \(w\) from \(v\). The vertices \(v_0\) and \(v_m\) are its initial vertex and final vertex, respectively \citep{kuang2026a}.
\end{definition}

Definition~3.9 of the source paper writes the successive vertices as
\[
v\cdot F_1,
\quad
v\cdot F_1F_2,
\quad\ldots\quad,
 v\cdot F_1\cdots F_m,
\]
while calling \(v\) the initial vertex. Writing \(v=v_0\) explicitly inside the sequence here only makes later cutting indices easier to read; it does not change the content of the definition.

\begin{example}[the same word at different occurrences]
Let \(w=F_1F_2F_3\). At one vertex \(v\), all three steps may be executable, giving
\[
v=v_0\longrightarrow v_1\longrightarrow v_2\longrightarrow v_3.
\]
At another vertex \(u\), \(F_1\) may be defined while \(F_2\) is undefined at \(u\cdot F_1\). Then \(w\) still exists as a formal word, but \(\operatorname{Traj}_u(w)\) does not exist.
\end{example}

Consequently, whenever a later chapter says that a factor forms a traverse or a return, it must also be able to answer three questions: at which starting occurrence is the factor executed? In which finite tile does it lie? At which level is the asserted boundary set being taken?

\section{Traverse: A Boundary-to-Boundary Excursion in the Same Finite Tile and at the Same Level}

Let \(T\) be a finite tile at a fixed Bratteli level, with boundary vertex set
\[
\partial T\subseteq V(T).
\]
The word \emph{boundary} has already been defined independently in the preceding chapters; this chapter does not assign it a new meaning.

\begin{definition}[traverse occurrence]\label{def:traverse-basic-en}
Let \(u\) be a contiguous factor of a word whose trajectory from a concrete starting vertex lies entirely in the finite tile \(T\). We call the trajectory
\[
\gamma=(b_-=z_0,z_1,\ldots,z_\ell=b_+)
\]
a \emph{traverse} if
\[
b_-,b_+\in\partial T,
\qquad
z_i\notin\partial T\quad(0<i<\ell).
\]
The decorated datum
\[
\mathsf T=(T;b_-,\gamma,b_+)
\]
is called a \emph{traverse occurrence}. The two endpoints belong to the boundary set of the \emph{same finite tile and the same Bratteli level} \citep{kuang2026a}.
\end{definition}

\begin{plainbox}[In ordinary language]
A traverse starts from a boundary vertex, travels for a while through the internal part of the current tile, meets no other boundary of the same level in between, and then lands for the first time on a boundary vertex again. It is therefore naturally a birooted object: \(b_-\) is the entry and \(b_+\) is the exit.
\end{plainbox}

\begin{center}
\begin{tikzpicture}[>=Latex,node distance=11mm,
 b/.style={circle,draw,very thick,inner sep=2.3pt},
 v/.style={circle,draw,inner sep=1.6pt}]
\node[b] (a) {$b_-$};
\node[v,right=of a] (x) {};
\node[v,right=of x] (y) {};
\node[v,right=of y] (z) {};
\node[b,right=of z] (d) {$b_+$};
\draw[->] (a)--(x); \draw[->] (x)--(y); \draw[->] (y)--(z); \draw[->] (z)--(d);
\node[below=7mm of y]{no further contact with $\partial T$ in between};
\end{tikzpicture}
\end{center}

The definition of a traverse is stronger than ``an arbitrary boundary-to-boundary path,'' because it forbids another boundary visit in the middle. If a long trajectory successively visits the boundary occurrences
\[
b_0,b_1,b_2,\ldots,
\]
then the minimal segments between successive visits are the elementary traverses. This minimality later supplies definite cutting points for lower-level decomposition.

\begin{principle}[occurrence and type must again be separated]
The object \((T;b_-,\gamma,b_+)\) is first a concrete occurrence. Only after forgetting its position in the ambient tile/orbital graph and retaining its birooted isomorphism class do we obtain a traverse type. Later witness persistence, connector context, and germ-sensitive data cannot be recovered from the bare type alone.
\end{principle}

\section{Finite-Tile Return: A Return Is the Special Same-Boundary Case of a Traverse}

Only now are we in a position to define the finite witness that the rest of the book will repeatedly use.

\begin{definition}[finite return witness / closed traverse]\label{def:finite-return-witness-en}
A \emph{finite return witness} is a closed traverse occurrence
\[
W=(T;b,\gamma,b),
\]
that is, the special case of Definition~\ref{def:traverse-basic-en} in which
\[
b_-=b_+=b.
\]
The two occurrences of \(b\) here must be the \emph{same boundary vertex occurrence}, not merely two points of the same boundary type or two different points carrying the same label.
\end{definition}

Here and throughout this chapter, \emph{closed} in \emph{closed traverse} is a combinatorial same-endpoint condition. It does not mean that a subset is closed in the topological sense.

Thus a return in a finite tile is precisely the special same-boundary case of a birooted traverse:
\[
\boxed{
\text{boundary }b
\longrightarrow
\text{internal tile excursion}
\longrightarrow
\text{the same boundary }b.}
\]

\begin{center}
\begin{tikzpicture}[>=Latex,node distance=11mm,
 b/.style={circle,draw,very thick,inner sep=2.4pt},
 v/.style={circle,draw,inner sep=1.6pt}]
\node[b] (a) {$b$};
\node[v,right=of a] (x) {};
\node[v,right=of x] (y) {};
\draw[->] (a)--(x); \draw[->] (x)--(y);
\draw[->,bend left=45] (y) to node[above]{same boundary occurrence} (a);
\end{tikzpicture}
\end{center}

This also shows why three common confusions fail:
\[
\boxed{
\text{repetition}
\not\Rightarrow
\text{return},
\qquad
\text{closed walk}
\not\Rightarrow
\text{elementary return}.}
\]
The reappearance of an isomorphic tile copy shows only a recurrence of local type. It specifies neither an actual trajectory nor the same boundary coordinate. Conversely, even an ordinary closed walk that begins and ends at the same vertex may pass through other boundary vertices several times in between. In that case it should usually be cut into several traverses rather than treated as a single elementary return witness.

\section{Exact Relation to the Return Word in the Source Paper: Boundary Letters Enclose a Same-Boundary Internal Excursion}

The return word in the original growth paper is defined at a boundary point of an infinite tile. To prevent the finite witness introduced here from being collapsed into the source terminology, we write out the original structure in full.

Let \(\xi\in\Omega(\mathsf B)\) be a boundary point of the infinite tile \(T_\xi\), and suppose that
\[
g=F_1F_2\cdots F_m
\]
is composable at the relevant positions. Definition~3.10 of the source paper calls \(g\) a \emph{return word} if and only if:
\begin{enumerate}
\item \(F_1\) and \(F_m\) are boundary-edge labels at \(\xi\);
\item the trajectory begins at \(\xi\cdot F_1^{-1}\), and the first step through \(F_1\) enters \(\xi\);
\item the middle labels \(F_2,\ldots,F_{m-1}\) correspond to edges in the internal part of \(T_\xi\), and none of those edges is a boundary edge;
\item the final step leaves from the same boundary point \(\xi\) through \(F_m\), so the final vertex is \(\xi\cdot F_m\) \citep{kuang2026a}.
\end{enumerate}

The trajectory therefore has the form
\[
\xi\cdot F_1^{-1}
\xrightarrow{F_1}
\xi
\xrightarrow{F_2\cdots F_{m-1}}
\xi
\xrightarrow{F_m}
\xi\cdot F_m.
\]
In other words, the original return word packages a complete generator block of the form
\[
\begin{gathered}
\text{enter the boundary}
\longrightarrow
\text{same-boundary internal excursion}
\\[-1mm]
\longrightarrow
\text{leave from the same boundary}
\end{gathered}
\]
The middle segment is precisely the closed-birooted witness that the later chapters of this book will track separately.

Because the internal excursion has finite length while the infinite tile is an inductive limit of finite tiles, for every such return word there is a sufficiently deep level \(n\) at which the entire internal excursion already lies in a finite truncation tile. Since \(\xi\) is a persistent boundary point, its truncation \(\xi_n\) is still a boundary vertex. Hence every return word in the original sense produces a finite return witness at a sufficiently deep level.

The converse requires care. A finite closed-traverse occurrence is, by itself, only finite evidence. Promoting it to a return word in the sense of the source paper additionally requires an actually extendable infinite boundary context together with the corresponding entering and exiting boundary labels. A finite witness and a global source return word are closely related, but they are not interchangeable without further conditions.

\section{Source-Side Prototype with Multiple Singular Boundary Points: How a Bridge-Type Word Can Organize a Return Jointly}

\begin{sourcebox}[This structure already appears in the early preprint; it is not invented by this book]
Version~1 of Kuang's preprint, arXiv:2402.16238v1, Subsection~5.2, treats the case of ``multiple types of bridges'' \citep[Subsection~5.2]{kuang2024v1}. That version describes a tile-inflation picture with multiple germ-defining singularities / infinite boundary points. The published version later reorganized Section~5. The v1 citation here therefore points specifically to this multi-boundary prototype; it does not replace the main theorem of the published paper.
\end{sourcebox}

This prototype matters because the designated infinite boundary points are precisely germ-defining singularities. It shows that the germ-sensitive organization of returns studied in this book need not, even at the source, be centered on a unique singular boundary point. Regular boundary gluing carries no return-sector decoration in this subsection.

Let \(S\) be a finite symmetric generator set. In Subsection~5.2 of v1, every bridge type \(x\in S\) corresponds to a designated infinite boundary point \(\xi_x\). The level-\(n+1\) tile is assembled from the various \(x\)-bridge pieces \(T_{n,x}\). A single higher-level tile may therefore contain several adjacent bridge types.

For fixed \(x\), a rank-\(n\) traverse of type \(x\) is a boundary-to-boundary walk through \(T_{n,x}\) that meets no boundary in between. Let
\[
\Theta_{n,x}
\]
denote the set of such traverses, and set
\[
\Theta_n=\bigcup_{x\in S}\Theta_{n,x}.
\]
If a rank-\((n+1)\) traverse \(\tau\) induces rank-\(n\) traverses
\[
\Phi(\tau)=(\tau_1,\ldots,\tau_k),
\]
we may simultaneously record the bridge types through which they pass:
\[
\operatorname{typ}\Phi(\tau)=x_1x_2\cdots x_k.
\]
This type word is a finite encoding of the multi-singular-boundary geometry developed in this book. It records how a single higher-level trajectory successively crosses different germ-defining boundary/bridge sectors instead of projecting every crossing back to one distinguished site.

The source proof then searches this induced type word for a return subword. In particular, a return factor can be jointly determined by several successive bridge types; there is no reason for it to belong to a language controlled by one boundary point alone. The source side therefore already contains the pattern
\[
\boxed{
\begin{gathered}
\text{multiple singular boundary sectors}
\longrightarrow
\text{bridge-type word}
\\[-1mm]
\longrightarrow
\text{return factor}
\end{gathered}}
\]

\begin{warning}[multi-singular-boundary geometry does not by itself force a return]
Proposition~5.7 in Subsection~5.2 \emph{assumes} that there exists \(d>1\) such that every reduced word of length greater than \(d\) contains a return subword, and then derives a growth estimate from that hypothesis. In other words, the germ-defining singular boundary points / bridge types provide a structural prototype for a return grammar. They do not by themselves prove finite incompressibility or uniform return latency.
\end{warning}

Example~5.8 of that subsection gives a concrete three-sector picture. The generators \(a,b,c\) correspond respectively to the infinite boundary points
\[
1^\omega,\qquad 5^\omega,\qquad 6^\omega.
\]
By tracing paths in a finite orbital graph while ignoring boundary edges, one can check that every reduced word of length greater than \(12\) contains a return subword \citep[Example~5.8]{kuang2024v1}. This example will later serve as the most direct mathematical reference in the book for the idea that multiple organizing positions can jointly determine a single return grammar.

\begin{plainbox}[Keep the two meanings of ``multiple'' separate]
\textbf{Multiple germs at one singular site} means fixing \(\xi\) and comparing different germ sheets / resolutions. \textbf{Multiple germ-defining singular boundary points / bridge types} means recording which sectors a single higher-level trajectory visits in succession. The former is the direction of the pNH benchmark in Chapter~28; the latter is the multi-singular-boundary prototype of this section. Regular boundary points provide only ordinary reversible gluing and belong to neither of these two nontrivial germ structures. They must not later be merged into a single phrase such as ``multiple germs.''
\end{plainbox}

\section{Return Occurrence and Return Language: A Word Is Only a Label String; the Witness Says Where It Came Back}

Let \(u\) be a contiguous generator factor. Even when its spelling is fixed, it can behave differently at different tile occurrences. Word-level return status and the finite evidence that realizes it should therefore be recorded separately. If \(u\) is a source return word and
\[
W_R=(T;b,\gamma,b)
\]
is the finite return witness supplied by the internal same-boundary excursion in one admissible occurrence, then the later chapters record the complete witnessed datum as
\[
R=(u;W_R).
\]
Here \(u\) retains the complete spelling of the source word, including the initial and final boundary entry/exit labels, while \(W_R\) records only the finite internal witness responsible for the same-boundary return.

\begin{definition}[return language]\label{def:return-language-en}
In a system with a declared boundary/tile convention, let
\[
\mathcal R
\]
denote the set of admissible return words being used. When discussing Kuang's growth theorem, \(\mathcal R\) is the set of return words from Definition~3.10 of the source paper. In the finite-witness calculus developed later in this book, every accepted return factor must additionally carry a concrete closed-birooted witness.
\end{definition}

\begin{plainbox}[In ordinary language]
A literal word cannot replace occurrence bookkeeping. If we later ask how a return persists at a higher level, which germ sheet a lift reaches, or what residue remains, we must go back to the concrete witness rather than look only at the character string.
\end{plainbox}

This is also why the later theory always keeps
\[
\text{return status}
\qquad\text{and}\qquad
\text{finer germ continuation data}
\]
as different coordinates. A downstairs return that already has a finite witness is not revoked merely because a finer lift ends on a nonidentity sheet.

\section{Incompressibility: Factor Avoidance Relative to the Entire Return Language}

For a word
\[
w=F_1F_2\cdots F_n,
\]
let
\[
\operatorname{Fac}(w)
\]
denote the set of all its contiguous factors.

\begin{definition}[incompressible word]\label{def:incompressible-basic-en}
Relative to a declared return language \(\mathcal R\), a word \(w\) is called \emph{incompressible} if
\[
\operatorname{Fac}(w)\cap\mathcal R=\varnothing.
\]
Equivalently, no contiguous factor of \(w\) is an admissible return word \citep{kuang2026a}. Denote the language of all such words by
\[
\mathcal{IC}.
\]
\end{definition}

This is the language-theoretic core of Definition~3.11 in the source paper. For the structural theory developed later in this book, we state it first as a property of words on purpose, rather than immediately upgrading it to an ``intrinsic property of an abstract group element.'' A group element can have different generator representatives. Unless one additionally fixes a reduced or canonical representation, or proves independence from the chosen representation, factor structure belongs first to the concrete word.

\begin{warning}[one lift failing to come back is nowhere near enough to imply incompressibility]
Incompressibility is \emph{global return-language avoidance}. Even if the lift of one particular occurrence fails to return to the identity sheet in the graph of germs, this does not make the whole word incompressible. If another contiguous factor forms a return in some admissible boundary context, Definition~\ref{def:incompressible-basic-en} has already failed.
\end{warning}

The following statements must therefore remain permanently separate:
\[
\begin{array}{c}
\text{this local graph appeared again};\\
\text{this ordinary walk begins and ends at the same vertex};\\
\text{this factor has a finite closed-birooted witness};\\
\text{this source return word holds in an infinite boundary context};\\
\text{this complete word contains no return factor}.
\end{array}
\]
They are, respectively, recurrence, the closed-walk condition, finite witness, source return, and incompressibility. They are not interchangeable.

\section{Scale Discipline: Elementary Traverse Sets Do Not Simply Grow Monotonically}

Boundary is level-relative. Suppose that the same actual trajectory is placed simultaneously inside a level-\(n\) tile and a higher level-\(m\) tile. An endpoint of a level-\(n\) traverse may already have become a connector/internal position at level \(m\). One therefore cannot write
\[
\Theta_n\subseteq\Theta_m
\]
to mean that all elementary traverses automatically survive unchanged at higher levels.

What is stable is the occurrence witness. A finite same-boundary subtrajectory that already exists at a lower level can be embedded as a subconfiguration in a larger tile, but it need not remain an elementary traverse of the higher-level tile. Conversely, when a high-level traverse is viewed at lower levels, it can decompose into several lower-level traverses, many of which have distinct boundary endpoints.

This scale discipline also explains why we define only the basic objects here. It does not prematurely claim monotonicity for traverse sets, return sets, or birth levels. The later formal theory will define witness embedding and traverse decomposition separately; these are two different longitudinal operations pointing in opposite directions.

\section{Six Concepts in One Comparison Table}

\begin{center}
\small
\begin{tabularx}{0.96\textwidth}{>{\raggedright\arraybackslash\bfseries}p{3.35cm}Y}
\toprule
Object & What the mathematics actually requires \\
\midrule
local repetition / recurrent copy
& A finite rooted/birooted type appears again; no closed trajectory connecting the two occurrences is required.\\
closed walk
& The initial and final vertices coincide in the ordinary graph-theoretic sense; the walk may pass through boundary vertices many times in between.\\
traverse
& In the same finite tile and at the same level, a path goes from one boundary occurrence to another without meeting any other boundary occurrence in between.\\
finite return witness
& The special traverse with \(b_-=b_+\), hence with the same boundary occurrence at both ends.\\
source return word
& The full generator block from the original growth paper: boundary entry + same-boundary internal excursion + exit from that same boundary.\\
incompressible word
& No contiguous factor of the complete word belongs to the declared return language.\\
\bottomrule
\end{tabularx}
\end{center}

\section{Chapter Summary: The Logical Order of Return Cannot Be Reversed}

The chapter establishes only the following dependency chain:
\[
\boxed{
\begin{gathered}
\text{admissible word}
\longrightarrow
\text{trajectory occurrence}
\\[-1mm]
\Downarrow
\\[-1mm]
\text{boundary-bracketed traverse}
\longrightarrow
\text{closed traverse / finite return witness}
\\[-1mm]
\Downarrow
\\[-1mm]
\text{return language}
\longrightarrow
\text{incompressibility}
\end{gathered}}
\]

The three most important disciplines are:
\begin{enumerate}
\item \textbf{Repetition is not return.} Repetition supplies reappearance; the boundary-organized same-boundary condition supplies ``coming back.''
\item \textbf{Return occurs through an internal tile excursion, but it cannot be defined from the bare internal part / AF core alone.} It depends on boundary positions.
\item \textbf{Incompressibility is not a property of one germ or of one failed return.} It is a factor-avoidance property relative to the entire return language.
\end{enumerate}

\part{Structural Translation: Compatibility between Third-Order Consciousness and AF-by-Discrete}

\begin{center}
\large\itshape Starting from Two Independently Defined Native Domains, Building Auditable Interfaces One by One
\end{center}

\begin{methodbox}[Core statement of Part III]
The first two Parts kept two territories separate on purpose. Part I began from psychological, social, and philosophical reality: how experience becomes ``my experience,'' how a subject takes shape through memory, language, and relationships, and why the observer does not in fact stand outside experience. Part II then built the mathematical native domain of \AFBD{} on its own terms. Only in Part III do the two lines genuinely meet for the first time.

The aim here is neither to make psychology obey mathematics nor to dress familiar psychological phenomena in technical language. For a psychology reader, mathematics is worth introducing only when it makes psychological relations that were previously collapsed together more distinguishable. It should help us ask more clearly: when is a similarity merely a repetition, and when has it been organized as ``again''? Why is an unbroken stretch of experience segmented as ``this episode''? From where does a self-judgment acquire organizing authority? Once the background is enlarged, is the old structure overturned, or merely placed back inside a larger relation?

Part III therefore treats \AFBD{} as a \emph{controlled formal language}. The present cross-domain claims are mainly Level II: controlled structural mappings. Only when explicit psychological or experiential variables are encoded in a restricted model and the relevant relations can be tested there does a local module qualify for Level III: a limited formal model of experience. The rigor of a mathematical theorem belongs to the mathematical native domain; it does not automatically upgrade a psychological claim.

The notation
\[
A\strsim M
\]
means only that, on a publicly declared set of relations, the psychological object or process $A$ and the mathematical object $M$ exhibit candidate structural compatibility. It never means that they are the same thing, and still less that the brain ``is'' a groupoid.
\end{methodbox}

\begin{readerbox}[How psychology readers should read this Part]
On a first pass through each interface card, read only three entries: \textbf{Psychological motivation}, \textbf{Preserved relation}, and \textbf{Failure condition}. First check that the interface addresses a genuine psychological problem; only then ask what extra distinction mathematics contributes. The mathematical native domain, candidate observables, and technical stop lines can be read more closely on a second pass.
\end{readerbox}

\chapter{The Translation Interface: How Two Native Domains Can Meet on Legitimate Terms}
\label{chap:interface-en}

\par\medskip
\noindent\textbf{Why this chapter matters}\par
Psychology has never lacked explanations. When someone says, ``I've been rejected again,'' we can talk about attachment, schemas, memory bias, social comparison, self-esteem, defenses, current emotion, or relationship history. The real difficulty is often not that we lack words, but that the words arrive too quickly: once one explanation fills the scene, it becomes harder to see which stretch of experience was cut out, what made two experiences count as ``the same kind of thing,'' and who is completing that judgment.

Part III introduces mathematics precisely to leave some space inside explanation. It does not add a still larger story. It tries instead to provide a relational grammar, so that before saying ``this is the same pattern,'' ``this happened again,'' or ``this is just who I am,'' we ask several additional questions: \textbf{where does it begin, where does it end, what is preserved, what has changed, and where does the current observation position stand?}
\par\medskip

\section{Why Psychology Needs a Relational Grammar}

The processes studied by psychology often unfold without premarked breaks, overlap one another, and remain subject to reinterpretation. A day of experience has no natural tick mark announcing, ``this is where one event ends.'' Memory is not a videotape that retrieves the past unchanged. And when someone says today, ``I am back in the same old problem again,'' that does not imply that today and the past are the same in every respect. Research on event segmentation has shown that people actively organize ongoing activity into relatively distinct episodes, and that this segmentation is related to prediction, goals, and comprehension \citep{zacks2007,kurby2008}. Research on autobiographical memory likewise reminds us that current goals participate in the retrieval and organization of past material \citep{conway2000}.

Several ordinary sentences already contain relational operations:
\begin{itemize}
\item ``This happened again'' places two moments inside the same comparison relation.
\item ``This began that day'' marks an entrance within an ongoing stream of experience.
\item ``I have always been like this'' compresses local events into a history of the subject.
\item ``Now I see it differently'' or ``I have let it go'' assumes that the current observation position can re-evaluate a previous one.
\end{itemize}

Existing psychology can study all of these. This book does not replace event segmentation, autobiographical memory, metacognition, self-reference, or narrative identity. It asks for one additional resolution: can we isolate recurrent \emph{relational operations} that ordinary explanatory language tends to merge? Is repetition the same as return? Is a relation that has reached return-completion thereby sealed against all future continuation? Is the observer a fixed point of view left unaffected by the process it observes?

\begin{plainbox}[If mathematics only renames psychology, Part III has failed]
If ``self'' is merely renamed a root, ``cycle'' is merely renamed a return, or familiar psychological descriptions are simply replaced by mathematical nouns, this Part adds nothing. Mathematics becomes useful only if its formal distinctions generate questions that were previously difficult to pose cleanly and that can later be recorded, compared, and potentially refuted.
\end{plainbox}

\section{The Most Dangerous Kind of Translation Is Mistaking Renaming for Understanding}

Three sentences should immediately trigger suspicion:
\[
\begin{gathered}
\text{``anxiety is a boundary germ,''}\\
\text{``trauma is an incompressible element,''}\\
\text{``the self is a root.''}
\end{gathered}
\]

Their problem is not merely that they are mathematically imprecise. A mathematical noun has been granted an evidential position it has not earned. Anxiety may involve bodily arousal, attachment expectations, social evaluation, memory, actual risk, and present interaction. Calling it a boundary does not say which relation is preserved; worse, it can erase precisely the distinctions psychology needs to keep.

The \emph{Geometry of Existence} calls this danger \emph{structural overreach}: a local explanatory tool gradually acquires authority over relations it was never entitled to organize \citep{translation2026,geometry2026}. The translation protocol used in this book therefore insists on the slower order
\[
\boxed{
\begin{gathered}
\text{first state where psychology lacks a needed distinction}\\[-1mm]
\Downarrow\\[-1mm]
\text{then ask whether mathematics can supply one}
\end{gathered}}
\]

What crosses the interface is a \textbf{relation}, not a noun. A psychological episode and a mathematical traverse need not contain the same material. The candidate comparison concerns a process bracketed by an entrance and an exit. Likewise, the psychological recognition of ``again'' does not become the mathematical object called return. The relevant comparison is whether ``again'' adds a recognized return-completion relation beyond the weaker fact that two contents look alike.

\section{The Two Domains Meet at Interfaces Defined by Psychological Problems}

The following list is a \emph{problem map} on purpose, not a correspondence table. The psychological problem on the left must exist first. Only then may the mathematical vocabulary on the right be asked whether it increases resolution.

\begin{center}
\small
\begin{tabularx}{0.96\textwidth}{>{\raggedright\arraybackslash\bfseries}p{3.4cm}Y}
\toprule
Psychological problem & Candidate mathematical interface / relation to inspect \\
\midrule
Experience can unfold before it is stably written as ``about me.'' & AF-like local flow: distinguish local occurrence from the later subject-label that acquires organizing authority.\\
Why does a current position suddenly acquire special significance? & Root / boundary: ask where local continuation becomes insufficient and an additional reference position becomes necessary.\\
Why is an ongoing stream of experience cut into ``this episode''? & Birooted traverse: make entrance, exit, and the scale of the finite segment explicit.\\
When does ``appearing again'' become ``it has come back again''? & Return: distinguish content similarity from a structured relation of return-completion.\\
How does a pattern enter a longer self-history? & Filtration / inflation: distinguish later reinterpretation from the continued existence of an earlier finite witness.\\
Why can the ``same pattern'' have a different consequence this time? & Graph of germs / residue: distinguish whether return-completion has occurred from how continuation differs afterward.\\
Is the observer really outside the process? & Reflexive access to rooting / resolution: treat the current observation position as a variable that can itself be seen, compared, and moved.\\
\bottomrule
\end{tabularx}
\end{center}

For concise bookkeeping, write
\[
\mathcal V_{3O}
=\{\text{event flow, ego-position, episode, again, history, observer}\},
\]
\[
\mathcal V_{AFD}
=\{\text{AF core, root, traverse, return, filtration, germ resolution}\},
\]
and use the dashed arrow
\[
\Phi:\mathcal V_{3O}\dashrightarrow\mathcal V_{AFD}.
\]
Here $\mathcal V_{3O}$ is not a completed mathematical category; it is the phenomenological and relational vocabulary independently developed in Part I. $\mathcal V_{AFD}$ is the mathematical vocabulary of Part II. The symbol $\Phi$ is \textbf{only an interface index}: it records which objects are worth placing in the same table. The substantive claim always lies in the sentence beneath the arrow---exactly which relation is preserved and exactly what is discarded. No functorial structure or total correspondence is being asserted.

\section{Structural Compatibility: How Strong a Claim Do We Actually Permit?}

For psychology readers, we keep the commitment limited on purpose. If two very different materials both exhibit an organization of the form ``entrance--process--exit,'' they may be comparable on that relation without sharing a mechanism, ontology, or material realization.

\begin{definition}[structural compatibility]
Let $\mathfrak A$ and $\mathfrak M$ be objects and relations drawn from different native domains. We say that $\mathfrak A$ and $\mathfrak M$ are \emph{structurally compatible relative to a declared relation family} if we declare in advance a finite or otherwise controlled family of relations and can show a comparable organization on those specified relations, while continuing to distinguish the materials, mechanisms, dimensions, and unmatched parts of the two sides. We write
\[
\mathfrak A\strsim\mathfrak M.
\]
\end{definition}

\begin{plainbox}[Do not read too much into this symbol]
The notation $\strsim$ means only: ``on the relations just declared, these two objects may be placed side by side in a controlled comparison.'' It is \textbf{not} a mathematical equivalence relation. In particular, no transitivity is assumed: from $A\strsim B$ and $B\strsim C$ one cannot simply infer $A\strsim C$. Every bridge must state separately what it preserves.
\end{plainbox}

For example,
\[
\text{episode}\strsim\text{birooted traverse}
\]
can mean that, on the selected relation, both involve an internal process bracketed by two functionally meaningful positions. The emotional texture of the episode, its social meaning, memory mechanisms, and neural realization are not thereby translated.

\begin{principle}[compatibility does not imply sameness]
\[
\mathfrak A\strsim\mathfrak M
\quad\not\Rightarrow\quad
\mathfrak A=\mathfrak M.
\]
Structural compatibility cannot be used to infer the same material realization, the same causal mechanism, or the same ontological nature.
\end{principle}

\section{\texorpdfstring{Where Formalization Begins:\\ When Metaphor Starts to Lose Its Freedom}{Where Formalization Begins: When Metaphor Starts to Lose Its Freedom}}

To say that a relation is being formalized does not mean forcing every experience into symbols. Formalization begins at the point where, after borrowing a mathematical relation, we are no longer free to change its meaning whenever a case becomes inconvenient.

Several consequences follow immediately.
\begin{enumerate}
\item If a psychological episode is proposed as a candidate return, the analysis must specify its starting position and what satisfies the return-completion condition. ``It feels like before'' is not enough.
\item If a finite return-completion has a clear witness, later context may reinterpret what the witness means, but a new story cannot simply pretend that the earlier finite witness never occurred.
\item If one kind of content repeatedly fails to complete a relevant return relation, that does not make the whole history globally incompressible. Another local factor may already contain a return.
\item If increasing the resolution reveals no difference in any continuation that can later be distinguished, then the additional vocabulary may have added no structure.
\end{enumerate}

This is the role of \AFBD{} here. It does not tell psychology what the true answer must be. It disciplines the use of a relation: if we choose to borrow a formal distinction, we also accept the distinctions and restrictions that make that formal notion what it is, and we abandon shortcuts that would erase them.

This book therefore goes beyond Level I rhetorical analogy but remains primarily at Level II controlled structural mapping. A local module reaches Level III only when explicit state variables, an episode segmentation rule, an endpoint rule or return test, and formally testable relations are supplied. A theorem proved later in Part IV remains a theorem of the mathematical native domain; it does not pay the evidential cost owed by psychology.

\section{Every Translation Must First Answer: Why Does Psychology Need This?}

Every important interface in this Part will be built in the same order. The narrative center of gravity remains the psychological problem throughout; mathematics is used to help psychology see a relation more clearly, not to lead the reader on a tour of a mathematical building.

\begin{enumerate}
\item \textbf{Psychological motivation.} What has been collapsed in ordinary language or the current narrative? Why does it need to be separated?
\item \textbf{Existing psychological coordinates.} How does mature psychology already describe the issue? The book must not pretend that an established literature is absent.
\item \textbf{Relation not yet separated.} Is the missing distinction entrance versus exit, repetition versus return-completion, coarse versus fine resolution, observing position versus observed content, or something else?
\item \textbf{Mathematical interface.} Only here may terms such as root, traverse, return, filtration, or germ resolution enter, and they must retain the definitions established in Part II.
\item \textbf{Preserved and discarded material.} Which relation crosses the interface, and which psychological material is explicitly not translated?
\item \textbf{Real example and non-example.} What case actually satisfies the proposed relation, and what merely looks similar while failing it?
\item \textbf{Observable and failure condition.} What could be recorded, compared, manipulated, or used to make the interface exit?
\item \textbf{Translation remainder.} Which aspects---for example qualia, phenomenological mineness, bodily mechanism, social realization, or some other variable---still have no legitimate interface?
\end{enumerate}

\section{Concept Passport: Leaving a Travel Record When a Concept Crosses Domains}

The translation protocol therefore gives every important cross-domain concept a \emph{Concept Passport} \citep{translation2026}. Without such a record, a term may travel for dozens of pages and we may forget that it originally carried three pieces of luggage but acquired ten more along the way.

\begin{center}
\small
\begin{longtable}{>{\raggedright\arraybackslash\bfseries}p{0.27\textwidth}>{\raggedright\arraybackslash}p{0.64\textwidth}}
Psychological motivation & Why does psychology need this distinction in the first place?\\
Native domain & In which domain was the concept originally defined before any cross-domain use?\\
Original definition & What does the term mean there, independently of the translation?\\
Preserved relation & Which exact relation is allowed to cross the interface?\\
Discarded content & What material, mechanism, or meaning is explicitly not transported?\\
Translation level & What evidential level does the present claim actually occupy?\\
Real-world example & What concrete case makes the intended relation visible?\\
Candidate observables & What could later be recorded, compared, or manipulated?\\
Non-example & What looks superficially similar but does not satisfy the relation?\\
Failure condition & What result would force this interface to be weakened or removed?\\
Reverse-translation check & Which established psychological constructs must remain able to compete with the new vocabulary?\\
Translation remainder & What still has no legitimate cross-domain translation?\\
\end{longtable}
\end{center}

This is not an administrative form. It is a device for preserving complexity. A mathematical word that enters psychology must remember that it has not acquired unlimited explanatory authority merely by crossing the border.

\section{Stable, Provisional, and Remainder: Three Status Labels for Cross-Domain Work}

Part III also marks statements by status.

\begin{description}[leftmargin=2.4cm,style=nextline]
\item[Stable.] The native-domain definition is clear, or the book has frozen the statement as the current working interface.
\item[Provisional.] The proposed relation appears to increase resolution, but still requires independent psychological evidence.
\item[Remainder.] No legitimate translation is presently available; leaving the space blank is more accurate than inventing an interface.
\end{description}

Remainder is not a defect. Interdisciplinary theory often fears blank spaces, but a blank can be more honest than a sentence that cannot be tested. If we do not know how phenomenological mineness arises, it remains a remainder. If we do not know which psychological or neural variable could correspond to a return residue, it remains a remainder. A theory that knows what it cannot see has at least a chance to keep seeing.

``Stable'' must also not be misread as ``psychologically proven.'' The mathematical statement that a finite return witness is a closed birooted traverse can be Stable in its native domain. The proposed structural compatibility between the ordinary recognition of ``again'' and this closed-birooted completion relation is, at most, a currently frozen Level II interface unless independent psychological evidence is supplied. The assertion that a return residue \emph{is} a neural memory trace remains Remainder.

\section{Where the Original Work of This Project Actually Lies}

The framework of Third-Order Consciousness in Part I comes from Guoxin Peng\textquotesingle s \emph{Third-Order Consciousness: From the Biological Self and Social Self to Vacant Consciousness (The Geometry of Existence Book 4)}, edited by Pinchuan Xie; this manuscript reviews, compares, reorganizes, and critically reconstructs it. Krishnamurti, Sartre, Mead, Goffman, Ricoeur, and the psychological literatures cited in Part I remain in their own historical and disciplinary contexts.

The original work of this project occurs between the two territories: selecting \AFBD{} as a controlled formal language; constructing the interfaces involving root, birooted traverse, return, filtration, and germ resolution; and subjecting those interfaces to the translation protocol and the checks against structural overreach supplied by the Geometry of Existence \citep{translation2026,geometry2026}.

The originality does not come from translating more psychological nouns into mathematical nouns. It lies, just as importantly, in \textbf{deleting what cannot legitimately be translated}, and in making the few bridges that remain clear enough about their preserved relations, evidential status, and failure conditions that they can actually fail.

\section{Chapter Summary: Once Mathematics Enters, Psychological Reality Should Become Clearer, Not More Distant}

Part III will not prove that ``consciousness is essentially a groupoid.'' Its aim is more restrained and more practical: to let a psychology reader see relations hidden by familiar words such as ``I,'' ``again,'' ``pattern,'' ``past,'' and ``observer.''

If the language of this Part is working, the reader should leave a chapter remembering a sharpened psychological question rather than a mathematical symbol:
\begin{quote}
What did I just treat as the same thing?\\
How was this episode cut out?\\
Who is recognizing this return?\\
After the scale changes, what remains of the earlier judgment?
\end{quote}

Chapter 14 begins with the quietest layer: what can already happen before ``I'' acquires primary organizing authority. Experience can occur, repeat, learn, and form local order. That is where First-Order Consciousness first meets AF-like flow.

\chapter{First-Order Consciousness and AF-like Flow: How Experience Operates Before Self-Narrative}
\label{chap:firstorder-aflike-en}

\par\medskip
\noindent\textbf{Why this chapter matters}\par
It is easy to imagine consciousness as a process in which the ``I'' appears first and only then begins to feel, judge, and act. Psychological reality often runs in precisely the opposite order. You may already have pulled your hand away from the hot rim of a cup, moved aside from a vehicle that suddenly approached, or noticed that something is wrong with a sentence before you can explain any of it to yourself. Only afterward does the ``I'' begin to put the experience into words: ``That just frightened me,'' ``Why am I so tense again?'' ``I really am no good at this.''

First-Order Consciousness directs our attention to precisely this span of time, so easily covered over by narrative: life is already sensing, filtering, predicting, approaching and avoiding, learning and adjusting, while a stable self-story has not yet acquired primary organizing authority. This chapter introduces AF-like flow not to turn First-Order Consciousness into a mathematical object, but to give this psychological fact a clearer structural language: complex local order can operate before ``an explanation about me'' takes charge.
\par\medskip

\begin{methodbox}[Interface card: First-Order Consciousness and AF-like flow]
\textbf{Psychological motivation:} Much bodily, perceptual, attentional, and skilled processing in experience does not wait for self-narrative to approve it before beginning. We need to separate ``the process is already organized'' from ``the process has already been stably written as a story about me.''

\textbf{Native domain of Third-Order Consciousness:} direct body--environment coupling, affective and action orientation, minimal bodily orientation, and persistence of survival organization; a stable identity narrative has not yet acquired primary organizing authority.

\textbf{Mathematical interface:} the AF core provides a controlled background that is locally approximable, repeatable, and embeddable; the definition of these local relations does not require additional non-AF boundary/germ data.

\textbf{Preserved relation:} local order can be rich and stable without taking a higher-level self-indexed history as a condition of its definition.

\textbf{Level:} Level II.

\textbf{Not claimed:} First-Order Consciousness ``is'' an AF groupoid; nor is the first order claimed to lack memory, prediction, value orientation, bodily centering, or a learning history.
\end{methodbox}

\section{Psychological Motivation: Experience Does Not Wait for the ``I'' to Finish Speaking Before It Begins}

An utterly ordinary fact of psychological life can easily be flattened by language: many processes happen first; explanation arrives afterward.

When someone walks down a familiar staircase, the feet automatically adjust their height. When listening to one's native language, one does not calculate the syntax word by word before understanding a sentence. In skilled driving, directional correction, speed control, and visual search occur extensively outside explicit self-statement. Pain, hunger, fright, and bodily tension certainly do not first ask, ``What does this have to do with my life story?'' before deciding whether to enter attention.

This does not mean that such processes contain ``no cognition.'' Quite the contrary: they may contain long-term learning, prediction, action preparation, value weighting, and highly refined bodily regulation. Part I has already emphasized that First-Order Consciousness is not a blank substrate, but the primordial presence of life: the body is continually touched by the world and continually responds to it selectively \citep{toc2026}.

Research on the minimal self, interoception, the sense of ownership, and the sense of agency likewise reminds us, from different directions, that first-person orientation and bodily ownership do not need to be written into a long autobiography before they can exist \citep{gallagher2000,craig2009,braun2018}.

The distinction this chapter needs, then, is neither ``conscious/unconscious'' nor ``simple/complex.'' It is a distinction between two kinds of organizing authority:
\begin{center}
\fbox{\parbox{0.90\linewidth}{\centering
\textit{the process is being organized by body, environment, and learning history}\\[0.35em]
$\neq$\\[0.35em]
\textit{the process is being organized by a stable self-narrative.}
}}
\end{center}
This is the psychological opening through which AF-like flow can enter.

\section{The First Order Does Not Mean ``There Is No I'': It Already Has Bodily Direction, but Is Not Yet a Narrative Subject}

A dangerous misreading has to be corrected first. If the first order is described as ``there is no I,'' a psychology reader can easily imagine a neutral stream with no ownership, value, or center. That is not the meaning of \emph{Third-Order Consciousness}.

At the first order, pain can pull the system sharply toward ``this cannot continue here''; hunger can raise the priority of certain cues; threat can alter attention and action preparation. Bodily boundaries mean that a stimulus is no longer merely a change occurring somewhere in the world, but gradually acquires the direction of ``coming toward here.'' Guoxin Peng describes this process as the stream of experience tightening around bodily position under the pressure of survival: a complete subject does not first exist and then have sensations attached to it. Rather, sensations, alarms, and actions repeatedly converge on the same living interface, gradually forming the earliest ``I am here'' \citep{toc2026}.

Psychology permits a similar distinction. Gallagher's minimal self primarily concerns the first-person givenness of experience, ownership, and agency; it is not the same as the later narrative self maintained through language, roles, and autobiography \citep{gallagher2000}. Thus the first order, as this book uses the term, permits:
\begin{itemize}
\item a bodily center and spatial orientation;
\item tendencies to approach, avoid, or freeze;
\item skilled pathways produced by learning;
\item rhythm, expectation, and local prediction;
\item the most basic direction of ownership: ``this is what this body is undergoing.''
\end{itemize}

What it does not yet require is something else: drawing these changes into a stable, self-referential history and continually answering, ``What does this show about what kind of person I am?'' ``Why do I always do this?'' ``What does this have to do with my past and future?''

In other words, the first order is not centerless. Rather, \textbf{the center has not yet expanded into narrative sovereignty over the whole world of experience}.

\section{Why AF-like Is Needed Here, Rather Than Only the Term ``Automatic Processing''}

At this point a psychology reader may ask: if research on automatic processing, habit, embodied cognition, and related topics already exists, why introduce AF-like flow at all?

Because this book is trying to separate one further relation. The question is not merely whether a process is automatic, but whether a very large experiential background can be composed of many locally reproducible, embeddable, relatively stable pieces of organization, while only a small number of positions genuinely require the system to change how continuation is organized.

The AF core from Part II provides exactly such a formal picture. Mathematically, it is approximated stage by stage by increasingly large finite structures. Local patterns can recur and enter larger tiles without every occurrence requiring an additional boundary/germ label. What matters most here is not the technical definition of AF, but the picture it gives us:
\begin{center}
\fbox{\parbox{0.90\linewidth}{\centering
\textit{a great deal of local order can sustain itself,}\\[0.35em]
\textit{while special interfaces carry additional organization only at a few positions.}
}}
\end{center}

Bringing this picture back to psychology does not yield the claim ``the brain is an AF groupoid.'' It yields a candidate question: can much of a person's perception, action, linguistic skill, and routine regulation operate as a relatively stable background, while subject-centered explanation becomes conspicuous only at a smaller number of moments---when prediction fails, relational rules conflict, or identity has to be repositioned?

If the answer is positive even in part, AF-like flow is not merely another name for ``automatic processing.'' It begins to help us study how background and interface divide their work.

\section{The First Structural Translation: Being Organized Does Not Mean Having Already Been Named by ``Me''}

On the psychological side, the relation we preserve is:
\begin{center}
\parbox{0.92\linewidth}{\centering
body, perception, attention, and skilled processing can form stable local order,\\[0.25em]
without a stable self-narrative acquiring explanatory authority at every step.
}
\end{center}

On the mathematical side, the relation we preserve is:
\begin{center}
\parbox{0.92\linewidth}{\centering
AF-local moves, copies, and embeddings can be defined and composed,\\[0.25em]
without additionally invoking non-AF boundary/germ data.
}
\end{center}

Hence:
\begin{correspondence}[First-Order Consciousness and AF-like flow]\leavevmode\par
\noindent\hfill
\fbox{\parbox{0.68\textwidth}{\centering\small
\textit{an experiential background in which stable self-narrative has not acquired primary organizing authority}\\[0.35em]
$\strsim$\\[0.35em]
\textit{an AF-like local background sustainable without additional boundary data}
}}\hfill\mbox{}\par
The preserved relation is this: \textbf{local order can first be established, repeated, and maintained, while incorporating those occurrences into ``my history'' belongs to another layer of organization.}
\end{correspondence}

Here ``first'' primarily marks a relation of organizational dependence. It does not place the AF core and boundary data into a mathematical time sequence. In a complete \AFBD{} system the two can coexist; in psychological life, pre-reflective processing and self-narrative also often operate at the same time. What we are comparing is whether a local process must borrow the latter in order to count as a process at all.

This is a Level II structural mapping. AF-like refers to an organizational relation, not to a brain region, neural-network topology, or the ontology of consciousness.

\section{A More Important Distinction: Repetition Can Occur Before ``Again'' Appears}

One of the most useful contributions of the first order to later chapters is that it separates repetition from recognized return.

Life is full of repetition. Breathing is rhythmic; walking is rhythmic; skilled actions repeatedly draw on familiar pathways. A person may repeatedly tense, avoid, or approach under similar cues. The very possibility of learning already means that the past remains in the present system in some form. Yet such repetitions do not automatically say:
\begin{quote}
``Look, I am back in that same old problem again.''
\end{quote}

The latter sentence has done something extra. It places a present occurrence and a past occurrence within the same self-indexed history and declares that, at a relevant position, they count as a return.

Thus:
\[
\boxed{\text{repetition can be lived before it is recognized as return.}}
\]

This is much more accurate than saying ``the first order has no memory.'' \emph{Third-Order Consciousness} itself describes first-order time in terms of rhythm, repetition, reinforcement, and persistence of survival organization, while reserving ``my past'' and ``my future'' for a later and stronger organization of the subject \citep{toc2026}. The first order can have habits without yet writing habit into character; it can contain repetition without yet writing repetition into destiny.

The AF core offers only a restrained formal reminder here: the reappearance of copies does not mean that a rooted return-completion relation has already been established. The later chapters on return make this distinction formal.

\section{Psychological Reality I: When Danger Arrives, the Body Does Not Wait for Narrative to Convene a Meeting}

When a car suddenly approaches from the side, a person may first turn away, draw in the shoulders, raise muscular tension, or even take a step before verbal judgment catches up: ``That almost hit me.'' A few seconds later, a second layer of story may appear: ``Why am I always so inattentive?'' ``People here drive with no regard for anyone.'' ``I must never walk this way again.''

These moments help separate three kinds of organization:
\begin{enumerate}
\item immediate body--environment coupling has already changed action;
\item the event is subsequently owned as ``something that happened to me'';
\item a still later narrative begins to write the event into a story about competence, other people, or rules for the future.
\end{enumerate}

The AF-like mapping picks out only the structural difference between the first layer and the latter two. It does not deny that the first layer already has a bodily center, and it does not claim that real neural responses possess an approximately finite mathematical structure.

\section{Psychological Reality II: Skill Is Itself a Historically Shaped Form of ``No Need for Self-Narration''}

Let us take a scene with neither danger nor strong emotion. When an experienced typist writes their own name, they do not need to issue a key-by-key command to their fingers--``Which letter am I pressing now?'' Likewise, someone who has spent years reading material in a field may look at a familiar structure and immediately feel that ``something is wrong here'' or ``this argument is missing a step,'' sometimes several seconds before being able to say exactly why.

Such fluency is not ``primitive.'' It is precisely the result of historical training. The past has already entered present action and perception; it simply does not appear every time in autobiographical form. \textbf{History can settle into skill without having to rise into story every time.}

This matters greatly for the first-order mapping. An AF-like background is not a ``pure present with no past.'' It is closer to a background repeatedly shaped by a great deal of past history and now capable of operating locally with fluency. Only when the process jams, rules conflict, or the present result cannot continue along a familiar route does the explanatory ``I'' often become suddenly loud.

\begin{nonexample}[``First order means no memory, no learning, no meaning'']
If the first order is understood as an instantaneous stimulus--response machine with no learning history, this translation should stop immediately. The first order permits habit, prediction, bodily value, skilled processing, and structural traces left by the past. It simply does not give a stable self-narrative primary organizing authority over every local process.
\end{nonexample}

\section[``Not Too Large'': Why We Need a Background--Interface Hypothesis]{``Not Too Large'': Why We Need a\\Background--Interface Hypothesis}

One of the deepest inspirations that \AFBD{} offers this project can be expressed in an ordinary sentence: \textbf{a system can be extremely rich without every place in it being equally complex.}

In the fixed-generator \AFBD{} tile picture, a large amount of structure belongs to an AF background that can be approached through finite scales. The persistent boundary points that genuinely survive on infinite tiles are finite in total across \emph{all} infinite tiles; consequently only finitely many infinite tiles carry boundary. Regular boundaries preserve reversible gluing across tiles, while singular boundaries additionally carry nontrivial germ-resolution data. If the tile-inflation form of bounded type used here is also assumed, then finite-tile boundary cardinalities and the number of tile types per level are uniformly controlled, so the aggregate boundary/connector burden across tile types at each finite Bratteli level has a uniform bound as a derived consequence. Roughly speaking, this is a ``not too large'' mode of organization: complexity is not spread uniformly across the whole space, but organized by finitely many interfaces; bounded type further ensures that the boundary burden of each finite tile does not explode with scale, while the finite-rank presentation controls the aggregate across tile types.

What makes this picture attractive for psychology is not that we have proved the brain ``almost finite.'' It is that ordinary experience suggests a possibility worth testing: much processing may run routinely, while the moments that genuinely force the system to redraw ``what is happening, where I stand, and what to do next'' may be much sparser than background processing.

If such turning positions can eventually be measured independently, and if they prove more predictive than ordinary moments of segmentation, memory updating, or self-referential reorganization, then ``background + sparse interfaces'' will have acquired empirical content. Conversely, if psychological complexity is always dispersed across the entire space of psychological conditions, if no relatively stable background can be found, and if no small number of positions carries a disproportionate share of the organization of continuation, then this AF-like mapping should exit.

\begin{plainbox}[``Not too large'' does not mean psychological life is simple]
The intuitions of tame, amenable, or almost-finite cannot be turned directly into personality vocabulary. A person's experience may be extraordinarily rich, full of memory and social meaning. The question is only \emph{how its organizational complexity is distributed}. ``Not too large'' is closer to saying that most local relations can be handled repeatedly in finite ways, while the interfaces that genuinely change the overall manner of continuation are relatively concentrated.
\end{plainbox}

\section{Concept Passport}

\begin{longtable}{>{\raggedright\arraybackslash\bfseries}p{0.24\textwidth}>{\raggedright\arraybackslash}p{0.68\textwidth}}
Psychological motivation & Distinguish ``experience is already highly organized'' from ``experience has already been taken over by a stable self-narrative.''\\
Original psychological concepts & First-Order Consciousness, pre-reflective experience, minimal self, interoceptive / embodied organization.\\
Original mathematical concepts & AF core and AF-like local background.\\
Preserved relation & Local processes can repeat, learn, embed, and maintain order without additional self-indexed / boundary organization as a prerequisite of their definition.\\
Discarded content & Bodily material, affective value, neural mechanisms, qualia, and the literal mathematical meaning of finite orbits.\\
Translation level & Level II.\\
Real-world examples & Immediate bodily adjustment under danger; skilled action and perception operating before explicit self-explanation.\\
Candidate observables & The temporal relation between automatic/skilled processing and explicit self-reference; the distribution of routine periods and highly organized turning positions; whether the latter more strongly predict segmentation, memory updating, or subsequent prediction.\\
Non-example & Describing the first order as lacking memory, learning, or a bodily center; or directly calling all automatic processing AF.\\
Failure condition & No relatively stable local background can be found; organizational complexity is evenly dispersed throughout the system; the proposed interfaces generate no additional observable distinction.\\
Reverse-translation check & Do not use the AF property of having ``no intrinsic labels'' to erase valence, ownership, agency, or bodily orientation in psychology.\\
Translation remainder & Phenomenological mineness, qualia, actual neural realization, and how the bodily center of the first order develops into a full subject.\\
\end{longtable}

\section{Chapter Conclusion: Life Lives First; Only Later Does It Explain Itself}

The most important compatibility to preserve between the first order and AF-like flow is not ``lower,'' ``simple,'' or ``without a self.'' Quite the contrary, it lets us see something more powerful: \textbf{before self-narrative arrives, life already has an order of its own.}

The body is already assigning priorities; perception is already filtering; learning history has already settled into skilled pathways; rhythm has already organized time. These forms of order simply have not yet been brought under a single historical language of ``how I have always been,'' ``what this says about who I am,'' or ``which old problem of mine this is again.''

An AF-like background gives this layer of experience a restrained geometric picture: broad local flow can recur, while additional subject-centered organization becomes necessary only at certain positions. The next chapter turns to precisely those positions---moments when familiar local order is suddenly no longer enough and the system must invoke new relational information to determine how to continue. At that point the background first reveals a fissure, and boundary and root acquire their genuine psychological motivation.

\chapter{Fissure, Boundary, and Endogenous Root: When the Old Route Is Suddenly No Longer Enough}
\label{chap:boundary-root-en}

\par\medskip
\noindent\textbf{Why this chapter matters}\par
Most of the time, psychological life does not need to reinvent itself. We move forward through familiar actions, relational scripts, networks of knowledge, and bodily rhythms; many judgments never need to be put into complete sentences before we have already moved on.

What deserves attention is another kind of moment: what used to carry us forward almost automatically is suddenly no longer enough.

An unanswered message, a counterexample in a paper, a conflict of roles, or a key definition in a talk that never quite becomes intelligible can make experience hesitate at a particular point. The world has not collapsed, and the emotion need not be intense. It is simply that, from here onward, the old local rules can no longer quietly determine the next step. More relational history, identity, or normative information has to be brought in. This book provisionally calls such a position a \emph{continuation crack}. Mathematical boundary enters here not because it sounds like a ``psychological boundary,'' but because it too distinguishes a structure of the following kind: \textbf{a local background may ordinarily be sufficient to operate, while at a few positions continuation requires additional data.}
\par\medskip

\begin{methodbox}[Interface card: continuation crack, boundary, and ego-position]
\textbf{Psychological motivation:} Why can some very slight events suddenly alter subsequent interpretation, action, memory recruitment, and self-reference, while other intense events continue to run along familiar scripts?

\textbf{Psychological native domain:} the constraints that an existing routine organization places on the next step suddenly weaken; several continuations become viable at once; the system begins recruiting additional relational history, roles, rules, or self-evaluation.

\textbf{Mathematical native domain:} in the tile-inflation picture, an internal AF-like region supplies a large locally manageable background. For cross-domain translation, this book selects only \emph{singular persistent boundary germs}---that is, boundary points with $H_\xi\neq\{1\}$---as the candidate interface prototype. Regular boundary points perform only reversible tile gluing and do not enter the third-order correspondence of this chapter.

\textbf{Preserved relation:} a decline in local sufficiency; the entry of additional relational data; and the resulting salience of an internal reference position.

\textbf{Level:} Level II.\qquad
\textbf{Not claimed:} trauma, anxiety, conflict, intense emotion, or what ordinary language calls an ``interpersonal boundary'' does not thereby become a mathematical boundary.
\end{methodbox}

\begin{warning}[From this point onward, the boundary translation used in this book refers only to singular boundary germs]
Mathematically, $\partial_S^\infty$ contains both regular and singular boundary points. A regular boundary point has trivial group of germs. It can connect two tile pieces through a pair of invertible labeled boundary arrows, yet it does not produce a nontrivial germ fibre.

The psychological translations later in this book concerning the continuation crack, the third-order observing position, return residue, and multi-sector germ structure are directed by default only at $\partial_{S,\mathrm{sing}}^\infty$. Thus, in the cross-domain chapters, \emph{boundary-like interface} is shorthand for \emph{singular-boundary-like interface}. It is not a psychologization of all mathematical boundary points.
\end{warning}

\section{The Real Psychological Question: Why Do Some Moments Suddenly Become ``Heavier''?}

Imagine an ordinary morning in which you are handling dozens of emails. Most messages can be classified almost as soon as they are read: answer this one, archive that one, forward another, leave one for later. Then a short message arrives: ``Could we talk for a moment?'' Its literal content is very small. The bodily reaction need not be dramatic. Yet attention is suddenly redistributed. What is the problem? Which matter is this about? From what role am I being addressed? Should I explain something, prepare for something, or simply wait?

Something similar happens in research. While reading a paper in a familiar area, a trained reader may absorb most technical details through an already established network. A counterexample then appears, and the old route of proof loses direction. What changes is not merely the quantity of information. The local organization that had been determining the next step is no longer sufficient. New definitions, comparison objects, and sometimes even one's relation to the research problem begin to enter.

Psychology already has mature ways to study neighboring phenomena: event segmentation, prediction updating, uncertainty, schema violation, and self-reference. The purpose here is not to cover all of them with a new word. Event Segmentation Theory, for example, treats ongoing experience as organized through event models and relates prediction error to model updating and segmentation \citep{zacks2007,kurby2008}. Here we ask a narrower relational question:

\begin{center}
\fbox{\parbox{0.90\linewidth}{\centering
\textbf{When the old local organization is no longer enough, what additional relations begin to acquire organizing authority?}
}}
\end{center}

The continuation crack is meant to mark not the place where experience hurts most, but the place where continuation begins to require a different relational setting.

\section{Continuation Crack: The Moment an Old Script Loses Its Automaticity}

We can now make the working term more precise.

\begin{definition}[Candidate continuation crack]
Relative to an explicitly declared scale of analysis, we call a position a \emph{candidate continuation crack} if the currently activated routine response, local feeling rule, task rule, or existing narrative is no longer sufficient to stably constrain how what comes next is to be understood, acted upon, or placed into history, several continuations thereby become viable at once, and the system is prompted to recruit additional relational history, identity, norms, comparison frames, or self-evaluation.
\end{definition}

\begin{plainbox}[A familiar road reaches a fork]
This does not mean that the psyche has ``cracked.'' The simpler image is that a familiar road reaches a fork. The road behind has not disappeared, and the feet can still move; what has failed is automatic navigation. New questions begin to matter: What am I actually facing? From which position am I to continue? Which parts of the past have now become relevant?

A continuation crack is therefore a candidate site of \emph{relational reorganization}, not a clinical symptom.
\end{plainbox}

The content can vary greatly. The other person suddenly goes silent, and the old script no longer decides whether this means busyness, rejection, or the end of an interaction. A skilled action fails, and automatic execution gives way to explicit monitoring. A counterexample appears, and the current proof framework no longer decides the next move. A friend suddenly speaks as a manager, and the relation must be reread through another role. An ordinary evaluation can even shift a task from ``finish the work'' to ``prove that I am capable.''

These situations are not grouped together because all of them feel unpleasant. They are grouped together because they share an organizational action: \textbf{the authority of a local rule declines, and additional relations begin to enter.}

\section{The Boundary Is Not Drawn by the Analyst: How the System Forces the Crack into View}

If a continuation crack depends only on the researcher saying afterward, ``this point was important,'' then the hypothesis can almost never fail. A boundary-like hypothesis with genuine empirical content should have a direction of \textbf{endogenous detection}: first observe where the system repeatedly loses local sufficiency, and only then decide whether that position deserves to be marked as a candidate interface.

In practice, four kinds of evidence can be sought. They are independent in principle, although they can reinforce one another:
\begin{enumerate}
\item \textbf{recurrent local-sufficiency failure:} across several comparable episodes, an existing policy, prediction, role script, or narrative rule repeatedly fails at similar positions;
\item \textbf{extra-data recruitment:} whenever the process reaches such a position, relational history, norms, identity, or comparison frames that normally need not be explicitly recruited become relevant;
\item \textbf{return concentration:} across different episodes, candidate return endpoints and judgments of ``again''---that is, judgments that the process has re-entered the same organizing position---repeatedly cluster near this kind of transition;
\item \textbf{continuation consequence:} once the position is crossed, next-event prediction, the available action set, memory retrieval, or the self-reference frame changes in a relatively stable way.
\end{enumerate}

The methodological role of return is especially important here. Suppose a person repeatedly undergoes a pattern of the form
\[
\begin{gathered}
\text{local operation}\ \longrightarrow\ \text{loss of sufficiency at a certain position}\\
\longrightarrow\ \text{recruitment of additional relations}\\
\longrightarrow\ \text{re-entry into the same organizing position}.
\end{gathered}
\]
Then repeated returns can function like probes, helping locate within a long trajectory a crack that the system keeps encountering. If several different return grammars share some endpoints, they may even provide a kind of ``triangulation'' of a candidate boundary sector.

But this is not a reversal of definition. A mathematical boundary germ is not created by return, and on the psychological side the sentence ``I am doing this again'' is not enough to announce that a boundary has been found. The more accurate relation is:
\begin{center}
\fbox{\parbox{0.90\linewidth}{\centering
\textit{A boundary-like crack is a candidate interface repeatedly forced into view by the system;}\\[0.25em]
\textit{return is one source of evidence for detecting it.}
}}
\end{center}
If a supposed boundary appears only after the researcher tells a story about it, disappears entirely under another equally reasonable recording scheme, and has no stable consequence for any continuation variable, it should not be retained.

\section{Intensity Is Not Boundary; Mild Events Can Still Change the Entire Continuation}

Here we have to resist a natural psychologizing impulse: reading boundary as ``the point of strongest emotion.''

A person who suddenly sees a snake may experience an extremely intense fear response while still following a highly familiar approach--avoidance program: attention narrows, the body withdraws, and distance increases. The intensity is high, yet continuation need not be ambiguous. Conversely, a softly spoken ``we can talk about it later'' may produce little physiological arousal while causing roles, expectations, and the future of a relationship to be recalculated.

The distinction preserved by this chapter is therefore
\[
\boxed{\text{structural importance}\neq\text{subjective intensity}.}
\]

Psychologically, this distinction matters because we often treat the most painful point as the most important one. Yet the position with the greatest organizing force for later prediction, event segmentation, and self-evaluation need not coincide with the emotional peak. A boundary-like candidate interface is better tested by asking whether it changes what counts as relevant, which parts of the past are brought back into view, who is regarded as responsible, and which futures suddenly become live options.

\section{Why Mathematical Boundary Is Useful Here}

Part II has already carried the burden of defining boundary mathematically and specifying its technical details. Part III keeps only the structural function that is most explanatory for the present comparison.

In the tile-inflation picture, many positions inside a tile can be handled through the existing AF-like local organization. At certain boundary positions, however, the information available within the current tile is no longer sufficient to determine the full continuation; connector data, a germ sheet, or a larger surrounding structure must also be known. In the simplest possible language:

\begin{center}
\fbox{\parbox{0.88\linewidth}{\centering
\textit{Most of the time, local information is enough; at a few positions, it suddenly is not.}
}}
\end{center}

This is what makes mathematical boundary genuinely attractive to the psychological problem. It does not tell us what emotion a person has ``at a boundary,'' nor which brain region is ``responsible for boundary.'' It supplies a formal picture of nonuniform organization: complexity need not be spread evenly across an entire system; a small number of interface positions can have disproportionate influence over how continuation is organized.

This motivates the following controlled correspondence.

\begin{correspondence}[boundary-like position and continuation crack]
\[
\boxed{
\begin{gathered}
\text{a position where existing experiential organization no longer suffices}\\[-1mm]
\text{to stably constrain the next step}\\[1mm]
\strsim\\[1mm]
\text{a boundary position where AF-local information is insufficient}\\[-1mm]
\text{and additional germ/connector data are required.}
\end{gathered}}
\]
This remains Level II. A psychological crack does not ``become'' a mathematical boundary. What is compared is the same kind of relational action: once local sufficiency declines, additional relational data begin to determine continuation.
\end{correspondence}

\section{Both Sides Speak of ``Additional Data,'' but Do Not String Them into the Same Causal Chain}

To keep a logical boundary from being obscured by elegant language, a simple discipline has to remain in force.

On the psychological side, the sequence may be written as
\[
\begin{gathered}
\text{familiar script is operating}\\
\downarrow\\
\text{local organization is no longer sufficient}\\
\downarrow\\
\text{more relations, history, roles, or rules are recruited}\\
\downarrow\\
\text{a reference position becomes more salient.}
\end{gathered}
\]

The mathematical side is a different and independent structure:
\[
\begin{gathered}
\text{AF-like local background}\\
\downarrow\\
\text{persistent boundary position}\\
\downarrow\\
\text{non-AF germ / connector data}\\
\downarrow\\
\text{distinguished root in a chosen presentation.}
\end{gathered}
\]

The comparison concerns the organizational shape of the two sequences. It does not identify their causal mechanisms. Otherwise ordinary psychological contents such as rejection, shame, control, or responsibility would simply be pasted onto mathematical boundary data, and the translation would have lost precisely the distinction it was introduced to preserve.

\section{A Crack Does Not Create an ``I''; It Makes the Current Reference Position Visible}

This is where root acquires its psychological motivation.

First-Order Consciousness did not mean the absence of bodily centering, ownership, or first-person direction. Likewise, a crack or conflict does not create a subject ex nihilo. The wrong formula would be
\[
\text{crack}\longrightarrow\text{the ego is created}.
\]
Many self-indexing processes are already at work. What uncertainty can do is make one reference position acquire much greater organizing authority than it had a moment before.

Take the unanswered message. The decisive shift may be from a diffuse question---``what is happening?''---to a more specific position: ``I am the person who has already taken the initiative.'' Once this position becomes salient, different pasts, rules, and fairness judgments become relevant, and the space of reasonable next actions changes.

The same thing can happen in an unfamiliar academic talk. A momentary blank can be organized around
\[
\text{``My knowledge is too narrow.''}
\]
or around
\[
\text{``I am a researcher who temporarily lacks coordinates in this field.''}
\]
The local facts need not change. What changes are the pasts recruited, the futures that remain available, and the way evaluation is organized.

This is the psychological motivation for introducing root: not to invent another self, but to mark \emph{where the current organization is being indexed from}.

\section{Ego-Position: Where Is Experience Being Reordered Around Right Now?}

\begin{definition}[ego-position]
An \emph{ego-position} is an internal reference position that acquires a prominent indexing role in the current organization of experience. It helps determine which materials are relevant to the present problem, which parts of the past are brought into the present organization, who is regarded as needing to act or bear responsibility, and which futures are treated as still viable. It is a local and movable coordinate of organization, not the whole subject and still less an observer standing outside experience.
\end{definition}

The question is therefore narrower than ``Who am I really?'' It is closer to: \emph{around which ``I-here'' is this material being rearranged right now?}

An ego-position may be organized through bodily location, social role, responsibility, relationship history, or an evaluative goal. It may dominate for seconds or remain stable for much longer. The term is meant to mark current indexing authority, not to duplicate every construct already contained in self-concept.

A rooted graph $(\Gamma,o)$ provides the controlled mathematical image. The root does not create the graph. It declares a basepoint from which the current structure is viewed and organized. With the restrictions already stated, we therefore write:

\begin{correspondence}
\textbf{boundary/germ-distinguished root and ego-position.}\par
\begin{center}
\fbox{\parbox{0.62\textwidth}{\centering\small
\textit{an internal reference position that becomes salient in a continuation problem}\\[0.35em]
$\strsim$\\[0.35em]
\textit{a root distinguished by boundary/germ organization}
}}
\end{center}
\end{correspondence}

The most important stopping sentence remains:
\[
\boxed{\text{root models an ego-position, not the whole subject.}}
\]

\section{Root Is Narrower Than Self-Concept, Identity, and the Narrative Self}

Psychology readers especially need protection from conceptual inflation here. Root is not a new name for self-concept, and it is not a mathematical version of narrative identity. A person's self-concept can include relatively stable knowledge such as ``I am a researcher,'' ``I am a friend,'' ``I am a child of my parents,'' or ``I am a cautious person.'' Narrative identity organizes past and future into a longer life story \citep{mcadams2013,ricoeur1992}. Ego-position is much narrower. It asks only: \textbf{which internal position is acquiring indexing authority right now?}

The distinction can be summarized roughly as follows:

\begin{center}
\small
\begin{tabularx}{0.94\textwidth}{>{\bfseries}p{0.30\textwidth}X}
\toprule
Concept & Distinction relevant to this chapter \\
\midrule
minimal self & Experience has first-person givenness and ownership/agency; a complete narrative is not required \citep{gallagher2000}.\\
self-concept / role & Relatively stable knowledge about what kind of person one is and what social position one occupies.\\
narrative identity & A longer-term life story organizing past, present, and future \citep{mcadams2013}.\\
ego-position / root proxy & Which internal position is currently determining relevance, allocation of responsibility, and continuation.\\
\bottomrule
\end{tabularx}
\end{center}

The model therefore has not discovered a mysterious psychological ``root faculty'' unknown to existing psychology. It proposes a candidate \textbf{relational variable}: with content held comparable, does a change in the current self-reference position systematically alter episode segmentation, memory recruitment, return judgment, and future prediction? If not, the psychological value of the root mapping must shrink.

\section{The Same Fact, Viewed from Different Roots, Grows into Different Futures}

Root is most useful when it helps separate a change in the facts from a change in the reference position.

Take the fact ``the other person has not replied.'' Different current ego-positions tend to recruit different continuations:

\begin{center}
\small
\begin{tabularx}{0.94\textwidth}{>{\bfseries}p{0.39\textwidth}X}
\toprule
Current ego-position & Continuation it tends to recruit \\
\midrule
``I am the person who has already taken the initiative.'' & Whether to initiate again; whether reciprocity needs to be protected as a principle.\\
``I am the person who may be rejected.'' & Earlier rejection experiences, self-worth, and defensive responses to risk.\\
``I am the person gathering information.'' & Which external variables remain unknown and when it is worth judging again.\\
``I am the person who must educate or correct the other person.'' & Corrective dialogue, norms, and responsibility in what follows.\\
\bottomrule
\end{tabularx}
\end{center}

Each position may contain true information, and each can also overreach. The task of the third order is not to select a single ``most correct root.'' It is first to see which position has quietly acquired the power to write local material as though it were the whole.

This is also where the methodological meaning of the \emph{Geometry of Existence} genuinely enters. Structural overreach often occurs not because a judgment is entirely false, but because a local position forgets its scale and writes ``what can be seen from here'' as ``all the facts'' \citep{geometry2026}.

\section{Example One: No Reply---How a Mild Fact Becomes a Subject-Relevant Event}

At the level of external fact, there is only one sentence: the message has not been answered.

For a time, this may lead only to repeated checking, while the existing script is still enough to sustain ordinary activity. At some point, however, local interpretation is no longer enough: Is the other person busy, uninterested, impolite, or has the interaction already ended? Relational history and familiar rules are then brought in: ``I initiate at most twice.'' ``A normal person should respond.'' ``I do not want to spend time on someone who leaves things uncertain.''

This is the candidate continuation crack. What is worth recording is not primarily whether the person has become angry, but three things:
\begin{enumerate}
\item which old script loses its sufficiency here;
\item which additional relations and rules begin to enter;
\item which ego-position consequently acquires primary organizing authority.
\end{enumerate}

If the finally salient position becomes ``I am the party who must uphold the principle,'' then the event has already been transformed from an external silence into the entrance to a subject-relevant episode. The principle may be reasonable. The model is not evaluating whether the principle is right; it is asking how the principle acquired its organizing position.

\section{Example Two: An Unfamiliar Talk---How a Knowledge Gap Gets Written as ``About Me''}

While listening to a talk in a familiar field, a network built through long training can perform a great deal of work automatically: terminology is classified, counterexamples are retrieved, and questions arise naturally. On entering an unfamiliar field, that network may suddenly fail. The first crack can be very small: \emph{I do not know which conceptual system to use in order to continue.}

The psychological system can nevertheless search very quickly for a more stable root. It may write:
\[
\text{``My knowledge is too narrow.''}
\]
A local accessibility failure has then been raised into a global self-evaluation. Or it may write:
\[
\text{``At present I lack an interface to this field.''}
\]
From that position, a very different future continuation becomes reasonable: record the terminology, look for a survey, and ask whether one's own research problem can be reformulated through this new language.

Neither narrative is guaranteed to be correct. What matters is that Part III can now ask a checkable question: \textbf{under different ego-positions, which memories are brought back by the same comprehension failure, which questions remain visible, and which actions are treated as possible?}

\begin{nonexample}[Calling every form of discomfort a boundary]
If a researcher can say only ``this makes me anxious, therefore this is a boundary,'' yet cannot identify where a local organization loses sufficiency, which alternative continuations become available, or what additional relational data are recruited, then the label has added no psychological resolution. It is merely a technicalized adjective for discomfort.
\end{nonexample}

\section{Why Root Usually Becomes Visible Sparsely, Rather Than Shouting ``I'' Every Second}

During skilled action, a person does not need to reinforce the sentence ``I am here'' every second. Driving, reading, walking, typing, and talking with familiar people can continue through rich bodily and cognitive organization while explicit self-reference remains in the background. It is more often when a route is interrupted, responsibility becomes ambiguous, roles conflict, or persistence across time is called into question that an ego-position is suddenly illuminated.

This is structurally compatible with the bounded-type / AF-by-discrete intuition that a small number of interfaces can organize a large background: most of the time the system flows through existing local order; a few transition positions have disproportionate influence on continuation.

One limitation must remain explicit: \textbf{what is sparse is the appearance of some salient re-rooting events, not selfhood itself.} First-person direction, bodily ownership, and longer-term identity can persist throughout. Here we study only what happens when a reference position moves from background to foreground and begins reorganizing the relations that follow.

\section{If This Correspondence Has Psychological Value, What New Questions Should It Generate?}

The point of the boundary/root mapping is not to give psychology two elegant new terms. It is to force research design to distinguish variables that are often allowed to collapse into one another:
\begin{itemize}
\item Can affective intensity and continuation reorganization be separated?
\item Which transition positions disproportionately predict event segmentation or memory updating?
\item With the external facts held fixed, does changing the self-reference frame alter subsequent prediction and return judgment?
\item Can a change in ego-position be completely explained by existing constructs of self-reference, appraisal, metacognition, or role?
\item If root language adds no predictive power at all, should it be removed?
\end{itemize}

This is where the writing of Part III differs most sharply from that of Part II: once a mathematical concept enters, it has to be sent back into psychological reality and asked what, exactly, it helps us \emph{see that we could not previously distinguish}.

\section{Concept Passport}

\begin{longtable}{>{\bfseries}p{0.25\textwidth}p{0.66\textwidth}}
Psychological motivation & Why some transitions change interpretation, action, memory recruitment, and self-reference, while emotional intensity alone is insufficient to explain the change.\\
Psychological concepts & continuation crack; ego-position.\\
Mathematical concepts & persistent boundary point; boundary germ; rooted presentation.\\
Preserved relations & local-sufficiency failure; additional continuation data; salience of an internal reference position.\\
Discarded content & mathematical topology and group-action mechanisms; concrete psychological content; affective valence; neural realization.\\
Translation level & Level II.\\
Candidate observables & changes in prediction/update, event segmentation, memory retrieval, self-reference frame, and responsibility/action choice before and after a transition; across episodes, whether local-sufficiency failures and return endpoints / sector passages repeatedly cluster near similar transitions.\\
Non-example & Strong emotion, trauma, conflict, or interpersonal ``boundary-setting'' automatically becoming mathematical boundary; declaring a boundary on the basis of one surprise or one return.\\
Failure conditions & Inability to identify where local organization becomes insufficient; a candidate interface that can only be assigned post hoc by the researcher and cannot be independently localized through repeated continuation failure / return concentration; root shift adds no explanatory power for any continuation variable.\\
Reverse-translation check & Do not use boundary/root to overwrite existing constructs such as event segmentation, self-reference, appraisal, role, or ownership/agency.\\
Translation remainder & phenomenological mineness; integration of the whole subject; coordination among multiple ego-positions; actual neural and social mechanisms.\\
\end{longtable}

\section{Chapter Conclusion: The Subject Position Often Becomes Visible Where ``the Road Is No Longer Enough''}

What this chapter is meant to leave behind is not two new terms, boundary and root, but a very simple psychological picture.

Most of the time, life flows along existing paths. The body knows how to adjust; skilled knowledge knows how to continue; a relational script already gives us some sense of what the next line is likely to be. At some point, however, the old road may still be there but no longer be enough to determine the next step. More of the past, more rules, more roles, and more possible futures begin to enter. At precisely this point, a reference position that had remained in the background is illuminated: \emph{From where am I seeing this? Whom am I now treating as the person who must take responsibility? Which pasts suddenly become relevant because of this position?}

Mathematical boundary helps mark a position at which ``the local is no longer enough.'' Root helps mark ``from where the relations are being reorganized.'' Neither tells psychology what the content must be, and neither tells a person how to act. They simply slow down a process that often happens too quickly to be seen, allowing us for the first time to notice:
\begin{center}
\fbox{\parbox{0.92\linewidth}{\centering
\textit{Experience does not always begin with a fixed ``I'' waiting for the world to be interpreted;}\\[0.25em]
\textit{often, difficulty in continuation brings a reference position from background to foreground,}\\[0.25em]
\textit{making the current ``I see from here'' suddenly visible.}
}}
\end{center}

The next chapter will continue from this psychological reality. Once an entrance position and an exit position are available, we can finally ask: \textbf{why is ongoing experience cut into ``this segment''?} Only then does the birooted traverse have a reason to enter.

\chapter{Birooted Traverse and Episode: How Experience Becomes ``This Segment''}
\label{chap:birooted-episode-en}

\par\medskip
\noindent\textbf{Why this chapter matters}\par
Life does not happen one square at a time. Conversation, waiting, bodily sensation, thought, and action often seep into one another; actual experience is more like an ongoing river. Yet as soon as we remember, judge, or tell a story, we quite naturally say: ``that argument,'' ``that talk,'' ``those days when I kept waiting for a reply.'' How does ongoing happening suddenly acquire a beginning and an end?

Here we focus on precisely this ordinary-looking yet very deep act. A mathematical \emph{birooted traverse} does not tell psychology what an event ultimately is. It supplies only a restrained pair of brackets: where the process enters, where it exits, and what is provisionally treated as one stretch of process between those two positions.
\par\medskip

\begin{methodbox}[Interface card: episode and birooted traverse]
\textbf{Psychological motivation:} Ongoing experience comes with no natural cuts, yet perception, memory, and narrative must continually organize it into ``this one thing.'' Changing the entrance or exit can also change the psychological meaning of the event itself.

\textbf{Psychological native domain:} event segmentation; changes of goal or prediction; subject-relevant episode.

\textbf{Mathematical native domain:} a traverse is a finite excursion bracketed by two relevant boundary positions; at the current scale, its internal segment does not hit another boundary of the same level.

\textbf{Preserved relation:} entrance, internal process, exit, and the possibility of comparing segmentations as scale changes.

\textbf{Level:} Level II.

\textbf{Not claimed:} the brain naturally cuts experience into mathematical traverses, or a psychological episode has one uniquely correct segmentation.
\end{methodbox}

\section{Psychological Entry Point: The World Keeps Happening, Yet Experience Is Remembered as ``Events''}

Imagine an argument. Where does it begin?

Does it begin when the other person says the hurtful sentence, or when you notice that they have been distracted all evening? Does it begin with last night's silence, or with an older conflict that was never resolved? And where does it end? When the two of you stop talking, when one person leaves the room, when you make contact again the next day, or only when you stop rehearsing the conversation in your mind?

``That argument'' is easily treated in ordinary language as an object that was already there. Psychology has to be more careful. Research on event segmentation has long shown that people organize ongoing activity into relatively discrete event units; changes in goals, contextual updating, and whether the current event model still predicts what comes next can all influence where we experience one event as ending and another as beginning \citep{zacks2007,kurby2008}.

The cut that marks an episode, then, is not a seam lying on the surface of the world waiting to be discovered. It is part of perception and organization.

We do not merely undergo events. We are also continually deciding what should count as the same event.

This is where the birooted structure first acquires a genuine psychological motivation.

\section{An Episode Is Not a Piece Cut from a Timeline}

The easiest mistake is to understand an episode as a time interval: nine to nine-thirty is one segment, nine-thirty to ten another. Such a partition is of course useful for recording, but it need not capture psychological organization.

Psychology is more interested in another question: what change makes the system begin to organize what follows through a different prediction, goal, role, or self-reference? The entrance and exit of an episode are therefore often functional rather than merely clock-based.

\begin{definition}[Subject-relevant episode]
Relative to an explicitly specified scale of observation, we call a stretch of ongoing experience a \emph{subject-relevant episode} if it is bracketed by two turning positions with an organizing role and no same-level repositioning occurs again in between. At the entrance, a current question, goal, role, or predictive structure acquires organizing authority; at the exit, that organization reaches completion for now, loses its force, changes direction, or is taken over by another organization.
\end{definition}

\begin{plainbox}[The central issue is not duration]
The central issue is not how many minutes an episode lasts. It is why the current model can legitimately say, ``this begins to count here,'' and why it can provisionally stop counting the same process here. With a different question, the same ten minutes may form a complete episode or only a small part of a larger one.
\end{plainbox}

Scale is therefore decisive. A two-hour meeting can contain hundreds of small shifts of attention. If the research question concerns only entering and leaving a condition of ``being evaluated,'' the whole meeting may still count as one coarse episode. If the question concerns prediction failures or shifts of attention, many finer episodes may be required.

An event does not become more ``real'' merely because it is cut more finely. The quality of a segmentation depends on whether it preserves the organizational change that matters to the current question.

\section{Why One Boundary Position Is Not Enough: An Event Needs Two Endpoints}

The previous chapter needed only one transition position to ask whether the old route had become insufficient. But to call something ``this one thing,'' more is required. We need to know:
\begin{itemize}
\item from which position the process enters the new organization;
\item at which position that organization ends, reaches completion for now, or is replaced;
\item and why what lies between can provisionally count as the same process.
\end{itemize}

This is the psychological reason for moving from a rooted to a birooted viewpoint.

One root answers: \emph{from where is organization occurring now?}

Two roots make a different question possible: \emph{from here to there, what does this middle count as?}

A sense of event therefore acquires duration for the first time. This is not because time suddenly appears. It is because one organizing position is taken up by another.

\section{Why Mathematics Enters Here: A Traverse Is a Disciplined Pair of Brackets}

Only at this point do we need to go back to the mathematical native domain of Part II.

In the tile picture, a traverse can be written
\[
\mathsf T=(P;b_-,\gamma,b_+),
\]
where $b_-$ is the relevant boundary position through which the current tile excursion is entered, $b_+$ is the exit position, and $\gamma$ is the finite trajectory between them. At the current scale, the internal part does not hit another relevant boundary of the same level.

A psychology reader need not imagine that a tile graph literally exists in the brain. The only feature worth carrying across here is a punctuation role:
\[
\text{entrance}\quad [\ \text{an internal process}\ ]\quad \text{exit}.
\]

A traverse does not tell us whether the brackets should contain ``anger,'' ``working memory,'' or ``social role.'' It imposes a different discipline: once we claim that this is one segment, we must explain why these two positions deserve to be endpoints and why, at the current scale, the middle does not contain another segmentation of the same level.

\begin{correspondence}[birooted traverse and episode]
\[
\boxed{
\begin{gathered}
\text{a stretch of experience bracketed by two functional turning positions}\\
\strsim\\
\text{a finite traverse bracketed by two boundary positions}
\end{gathered}}
\]
What is preserved is the relational skeleton ``entrance--internal process--exit,'' together with the constraint that the internal segment does not retrigger segmentation at the same level at the current scale.
\end{correspondence}

This is a Level II mapping. It does not identify the neural mechanisms of event segmentation, the computation of prediction error, or the subjective sense of an event with tile geometry.

\section{The Same Material Can Become a Different Event When the Endpoints Change}

The greatest psychological value of the birooted viewpoint is that it forces us to acknowledge that \textbf{what counts as the same episode is sensitive to endpoints}.

Take the same fact that a message has not received a reply. We can segment it in at least several different ways:

\begin{center}
\small
\begin{tabularx}{0.98\textwidth}{>{\raggedright\arraybackslash\bfseries}p{0.20\textwidth}>{\raggedright\arraybackslash}p{0.30\textwidth}Y}
\toprule
Analytic question & Candidate entrance / exit & Resulting episode \\
\midrule
Waiting behavior & send the message / stop checking & a process of ``waiting--checking--exiting'' \\
Relationship judgment & notice uncertainty in the interaction / decide not to continue & an evaluation of ``whether the relationship continues'' \\
Self-evaluation & ``maybe something is wrong with me'' appears / the evaluation temporarily loosens & a self-evaluative episode \\
Rumination & begin rehearsing the conversation internally / attention is genuinely taken over by another task & an internal-simulation episode \\
\bottomrule
\end{tabularx}
\end{center}

These segmentations are not four mutually exclusive ``truths.'' They answer different questions and may be nested within one another. The real danger begins when one segmentation quietly acquires exclusive explanatory authority and then says, ``this is where the thing really began, and this is where it really ended.''

This is also the kind of structural overreach against which the Geometry of Existence warns us: \textbf{a bracket that is locally useful should not automatically become the only legitimate line of demarcation in reality} \citep{geometry2026}.

\section{The Inside Can Be Complex Without Becoming ``My Story'' Every Second}

A single episode can contain a great deal of complex processing: bodily tension, visual search, automatic semantic activation, skilled action, movement of attention, and rapid judgment. None of this has to be narrated every second as ``I, as this kind of person, am now doing this particular thing.'' Yet once the entrance and exit have been distinguished by subject-relevant organization, these internal activities can be gathered as a whole into one sentence:
\begin{center}
\emph{``That was something I went through.''}
\end{center}

The structure here is therefore not
\[
\text{First Order}\longrightarrow\text{Second Order}
\]
as though a switch were being flipped. It is closer to this: rich pre-reflective processing continues to flow, while certain endpoints begin to acquire a subjectifying punctuation role.

This also makes a claim from Part I clearer. The subject does not need to occupy the entire experiential scene. It may acquire organizing authority only at a small number of positions, yet gather a long stretch of intervening flow retrospectively into ``my episode.''

\section{From ``Here'' to ``There,'' and Then to ``Back Here Again''}

Place three minimal structures side by side and the psychological intuition becomes very clear:
\[
\begin{aligned}
(\Gamma,b) &:\quad \text{this is the current reference position;}\\
(\Gamma;b_-,b_+) &:\quad \text{from here, to there;}\\
(\Gamma;b,b) &:\quad \text{starting here, and coming back here again.}
\end{aligned}
\]

The first line gives only a position. The second gives, for the first time, a bracketed stretch of experience. The third adds a return-completion condition to the episode: the exit is recognized as the same relevant position as the entrance.

Here we deal only with the second line. The ``again'' of the next chapter must not be smuggled in early. Having an episode does not mean that the episode already constitutes a return.

\section{Example One: Is an Argument One Event or Several?}

Two people begin arguing over one sentence. At a coarse scale, we may treat ``the argument begins--they stop speaking'' as one episode. But suppose the material shows several clear re-rootings:
\begin{enumerate}
\item at first the issue is ``why were you late?'';
\item in the middle it becomes ``you never respect my time'';
\item later it becomes ``are we fundamentally incompatible?''
\end{enumerate}

The external interaction is still one argument, but psychologically its organization may already have been taken over three times.

If the question concerns how long the emotional arousal lasted, one episode may be sufficient. If the question concerns how self-reference escalated, several episodes may be needed.

The mathematical condition that the internal segment has no same-level relevant boundary gives a useful discipline here. If a repositioning of the same analytic level genuinely occurred in the middle, we should not preserve an elegant story by pretending that the entire stretch remained one undivided event.

\section{Example Two: How a Vague Research Discomfort Becomes a Workable Episode}

Imagine a researcher reading a paper and initially feeling only that ``something is wrong here.'' After checking the argument several times, the researcher locates the problem: a principal assumption of the argument fails in the nonprincipal setting of the researcher's own work. That localization can serve as an entrance.

The following days may contain examples, subcases, failed approaches, and literature searches. Eventually a new organization is reached: ``finite singular germs and multiple singularities should be separated into two research routes.'' That decision can serve as an exit.

What matters for the coarse episode is not each calculation in isolation, but the sequence
\[
\begin{gathered}
\text{vague discomfort}\\[-0.2em]
\downarrow\\[-0.2em]
\text{a localized problem}\\[-0.2em]
\downarrow\\[-0.2em]
\text{a new research decomposition}.
\end{gathered}
\]

The proof attempts and technical subproblems inside the episode can of course be refined into smaller episodes when a finer research question requires it.

\section{An Episode Is Not Ten Minutes Arbitrarily Cut Out by an External Observer}

\begin{nonexample}[Mechanical ten-minute bins]
Suppose a researcher mechanically divides one hour of psychological activity into six ten-minute windows. These windows may be perfectly legitimate measurement bins, but they do not thereby become six birooted episodes. The traverse mapping requires the endpoints to carry an organizing role for the question being studied. External timing supplies only sampling cut points; it does not explain why experience is reorganized there.
\end{nonexample}

This distinction matters for experiments. Fixed time windows are useful statistically. But if the research question concerns event segmentation, memory updating, or judgments of ``again,'' an additional measure is needed for where participants experience a change in goal, prediction, or self-reference strong enough to reorganize the event.

\section{How Adjacent Episodes Hand Off: First Study Process Succession, Without Prematurely Defining a Whole Subject}

Once an episode has two endpoints, a smaller and safer question follows naturally: how does the exit of one process become the entrance to the next?

Mathematically, a sequence of composable traverses can be written
\[
b_0\xrightarrow{\mathsf T_1}b_1
\xrightarrow{\mathsf T_2}b_2
\xrightarrow{\mathsf T_3}\cdots,
\]
where the exit of one traverse is composable, in the declared structure, with the entrance of the next.

On the psychological side we borrow only a weak episode-handoff reading. A commitment made today can enter tomorrow's action. One relationship judgment can become the starting point of the next interaction. The exit of one research problem can become the entrance to the next problem. What is being described is a traceable succession among processes, not a deduction from composability that a complete subject has thereby been defined.

For this reason, this chapter does not call the chain a ``formal model of the persistence of the subject across time.'' The full persistence of subjecthood across time also involves autobiographical memory, bodily persistence, social roles, phenomenological mineness, and the integration of multiple self-positions. All of these go far beyond what composable traverses themselves can carry.

\begin{stopbox}[Limit]
This section preserves only the restricted relation of episode handoff / root-to-root composability. On the mathematical side, the orbit and germ structure of different singular sectors likewise cannot be erased at will by a psychological narrative.
\end{stopbox}

\section{Why Large Events Can Be Composed of Smaller Events}

Psychological life is inherently multiscale. ``Four years at university'' can be an episode; ``a thesis defense'' can be an episode; ``one question during the defense'' can also be an episode. The issue is not which level is the most real, but whether a traceable relation can be maintained among the levels.

Part IV will formally treat the decomposition of a high-scale traverse. Part III borrows only the psychological question: at a finer scale, can a larger episode be unfolded into several smaller episodes and transitions without completely rewriting the original occurrence order?

For example, ``a failed talk'' may have only an entrance and an exit at a coarse scale. A finer view may reveal: insufficient preparation; a smooth opening; loss of the thread at a key definition; the onset of self-evaluation; regaining a foothold through an example; and a decision afterward to repair the missing background. A coarse narrative can compress these materials. But if it directly rewrites ``I failed to understand something for one minute'' as ``the entire talk was a failure,'' the operation is no longer merely scale compression; it may have become narrative overreach.

Multiscale structure therefore functions first as a psychological discipline: a macroscopic story may omit details, but it should not arbitrarily rewrite the basic order and differences among those details.

\section{Coherence Does Not Require Memory Never to Change}

The mathematical native domain can impose strict coherence on decompositions. Human memory is different. Remembering is reconstructive, and later information can genuinely change where an event now seems to have begun or ended.

Part III therefore keeps only a weaker test:

\begin{plainbox}[A weak cross-scale coherence test]
If a coarse segmentation and a fine segmentation both claim to describe the same underlying occurrences, we should be able to indicate how they correspond. If no shared backbone can be identified at all, we should admit that we may already be dealing with two different narrative constructions, rather than simply ``the same event at different resolutions.''
\end{plainbox}

This does not force the psyche to obey a mathematical decomposition theorem. It only prevents ``higher resolution'' from becoming a license to rewrite the past arbitrarily.

\section{If This Interface Works, What Should Psychology Be Able to Observe?}

The birooted interface should generate questions that can leave the page:
\begin{itemize}
\item Do event-segmentation judgments shift systematically with goals, predictions, and self-reference frames?
\item If candidate entrances and exits are changed, do memory organization, causal judgment, and ``same event'' judgments also change?
\item Do subjectively indicated segmentation points predict later memory updating or behavioral change better than equal-length time windows?
\item Do coarse and fine segmentations show repeatable nesting or refinement, or are they almost entirely dependent on post hoc narrative?
\item If endpoint structure adds no explanatory power beyond existing psychological variables, should the birooted mapping be removed?
\end{itemize}

Mathematics is thus handed back to psychology. It does not answer for us what an event is; it forces us to measure how events are bracketed.

\section{Concept Passport}

\begin{longtable}{>{\raggedright\arraybackslash\bfseries}p{0.25\textwidth}>{\raggedright\arraybackslash}p{0.66\textwidth}}
Psychological motivation & Why ongoing experience is divided into ``this one thing''; why changing entrance or exit can change what counts as the same event and later memory, causal, and narrative organization.\\
Psychological concepts & event segmentation; subject-relevant episode; functionally defined entrance and exit.\\
Mathematical concepts & birooted traverse; boundary-to-boundary finite excursion.\\
Preserved relations & Two endpoints bracket an internal process; the internal segment has no same-level boundary hit at the current scale; episodes can be composed and compared across scales.\\
Discarded content & tile labels; the concrete group action; psychological content; affective valence; duration units; neural mechanisms of event segmentation.\\
Translation level & Level II.\\
Candidate observables & subjective event-segmentation judgments; changes in goal or prediction; memory updating; causal grouping; same-event judgments; multiscale segmentation.\\
Non-example & mechanical time windows; segmenting merely because affect is intense; the researcher assigning endpoints arbitrarily after the fact in order to fit a model.\\
Failure conditions & Endpoints have no additional organizing effect on any psychological variable; segmentations at different scales show no repeatable relation; all episode cut points can only be imposed arbitrarily by an external observer.\\
Reverse-translation check & Do not use traverse to replace event segmentation, event model, goal structure, prediction updating, or other native psychological concepts. Mathematics adds only the structural question of two functionally defined endpoints and their relation.\\
Translation remainder & neural formation of episodes; the subjective sense of an event; phenomenological mineness; and how multiple people in social interaction negotiate where the same event begins and ends.\\
\end{longtable}

\section{Chapter Conclusion: Events Are Brackets Placed into Flow}

What this chapter is meant to leave behind is not a terminological correspondence from psychological episode to mathematical traverse, but a way of looking at experience.

Life keeps happening. Bodily states do not come with neat periods; thoughts do not automatically end when the clock reaches a certain minute; a relationship leaves no clean cutting line in reality. Yet for action, memory, and narration, we have to keep placing brackets into the flow: from here onward, these things provisionally count together; up to here, another organization takes over.

The value of the birooted traverse lies only in making those brackets auditable. It forces us to ask: why these two endpoints? Is there really no same-level repositioning in between? If we change the question, the scale, or the root, does the original ``this one thing'' still hold?

An episode therefore ceases to be merely a box for content and becomes a relational act:
\begin{center}
\emph{Ongoing experience itself has no natural period;}\\[0.3em]
\emph{an entrance and an exit provisionally organize a stretch of flow into ``this one thing.''}
\end{center}

The next chapter adds only one small but decisive condition: what happens if the exit is not merely another position, but is recognized as the same relevant position as the entrance? Only then does ``this one thing'' acquire, for the first time, the structure of ``again.''

\chapter{Return and ``Again'': When Repetition Becomes Coming Back}
\label{chap:return-again-en}

\par\medskip
\noindent\textbf{Why this chapter matters}\par
``She did not reply.'' And ``She did not reply \textbf{again}.'' Only one small word has been added, yet the psychological structure has already changed.

The first sentence remains mainly in the present: a fact occurred. The second opens a temporal tunnel in an instant. Earlier episodes are brought back into view; the current event is placed beside them; one kind of similarity is selected; and finally these previously separate moments are compressed into the same subject-history: \emph{this is not the first time; I have come back here.}

Here we study that tiny yet powerful operation. Mathematical \emph{return} does not decide for psychology which patterns in a person's life are real. It adds only one discipline to ``again'': \textbf{similarity is not enough; to say that something has come back, one must explain how the current episode returns to an already distinguished relevant position.}
\par\medskip

\begin{methodbox}[Interface card: ``again'' and return]
\textbf{Psychological motivation:} People classify present and past experiences as instances of the same pattern. Such cross-temporal classification supports learning, prediction, and persistence of self across time, but it can also overcompress a local similarity into ``I am always like this.''

\textbf{Psychological native domain:} autobiographical memory; pattern recognition; self-indexing; ``again'' judgments; narrative compression.

\textbf{Mathematical native domain:} this chapter mainly borrows the finite return witness: a traverse starts at a relevant root and ends at that same root. The generator/context conditions of a source return word remain in the mathematical native domain.

\textbf{Preserved relation:} recognized return-completion, rather than general similarity of content; ``the same'' is always relative to a declared coordinate and scale.

\textbf{Level:} Level II.

\textbf{Not claimed:} every experience that ``feels like before'' constitutes a return, or the strict mathematical same-endpoint condition can determine a person's true psychological pattern.
\end{methodbox}

\section{How One Small Word, ``Again,'' Pulls the Present into the Past}

Consider four very ordinary sentences:
\begin{quote}
A: She did not reply.\\
B: She did not reply, and that made me uncomfortable.\\
C: Why am I \textbf{again} encountering this kind of disrespect?\\
D: I \textbf{always} encounter people like this.
\end{quote}

A mainly reports a current fact. B links that fact to a current feeling. With C, the temporal structure changes for the first time in an obvious way: the subject is no longer processing only today's silence, but places it in the same category as certain past events. D moves one step farther: what was still a local return is compressed into a more stable identity--fate narrative.

``Again,'' then, is not a trivial adverb. It is one of the smallest hinges of Second-Order Consciousness:
\[
\boxed{
\text{current event}
\longrightarrow
\text{cross-temporal pattern}
\longrightarrow
\text{possible identity narrative}.}
\]

\emph{Third-Order Consciousness} repeatedly emphasizes that First-Order Consciousness can already contain rhythm, habit, and repetition, without yet fully appropriating temporal persistence as ``my past'' and ``my future'' \citep{toc2026}. At the second order, memory, language, and subject-position begin sewing dispersed moments into one history. ``Again'' is one of the most ordinary grammatical forms of this stitching.

It can be extremely useful. Someone who says, ``I procrastinated before a deadline again,'' may have recognized a genuinely stable behavioral regularity. Someone who says, ``I lost my footing in this kind of relationship again,'' may retrieve material that had previously been ignored. Yet the same operation can also foreclose too quickly: one similarity is written as repetition; several repetitions are written as ``always''; and ``always'' is finally written as ``this is simply the kind of person I am.''

The real question of this chapter is therefore not \emph{why do people repeat?} It is:
\begin{quote}
\textbf{Under what conditions does an episode now in progress become recognized by the subject as ``I have come back to the same place''?}
\end{quote}

\section{``Again'' Contains at Least Four Psychological Operations}

Unpack the sentence ``I am doing this again,'' and several pieces of work have already been compressed into it.

\begin{enumerate}
\item \textbf{There must first be a current episode.} If the present material has not even been segmented as ``this stretch,'' there is not yet a determinate question of where it has come back to.

\item \textbf{Some part of the past must be brought back.} Autobiographical memory is not a video warehouse from which the past is neutrally replayed. What becomes accessible is affected by the current working self, goals, and self-knowledge \citep{conway2000}.

\item \textbf{A relevant similarity must be selected.} Present and past are never identical in every respect. To say ``again'' is already to decide which differences can be ignored and which relation is important enough to preserve.

\item \textbf{That similarity must be assigned to the same subject-history.} ``These two things resemble one another'' is not yet ``I have come back here again.'' The latter compresses moments at different times into a self-indexed history.
\end{enumerate}

These operations often arrive in everyday experience as a single, holistic ``again.'' Third-Order Consciousness slows them down and separates them.

This is adjacent to established work on autobiographical memory and narrative identity. Past experience is not merely stored; it is selected and reorganized under current goals, self-knowledge, and life-story organization \citep{conway2000,bruner1987,mcadams2013}. This book does not replace those theories. It isolates a narrower relational question: \textbf{when does a current episode come to be judged as completing a return to a previously meaningful position?}

\section{Similarity Is Not Coming Back: What the Psychological System Performs Is Cross-Temporal Classification}

Two heartbeats can resemble one another. Two subway trains can arrive at the same platform. A person can take the same route to work every morning. None of these recurrences automatically means, ``I am back in the same old problem.''

At minimum,
\[
\boxed{
\text{recurrence / similarity}
\neq
\text{recognized return}.}
\]

The decisive question is what counts as the \emph{same relevant position}.

Suppose two relational episodes both contain ``the other person did not reply promptly.'' If response latency is the variable under examination, the episodes may be treated as comparable. But if commitment, asymmetry of power, whether the message was seen, the history of the relationship, or the practical stakes are relevant, the two episodes may no longer occupy the same coordinate.

``The same pattern'' is therefore always a categorization relative to variables, scale, and the subject's current problem. ``Again'' is not literal sameness, and it is not sameness of emotion. It is closer to the following operation: although much of the content has changed, the system recognizes the endpoint of the current episode as coming back to a previously meaningful organizing position.

This is the point at which mathematical return begins to help.

\section{Why Mathematics Enters Here: Return Gives ``Coming Back'' a Minimal Skeleton}

In the mathematical native domain, a finite return witness is a \emph{closed traverse}: it starts from a boundary root $b$, passes through a finite internal process, and ends at the same relevant root $b$. The important point is not that the route ``looks like something seen before.'' It is that, in the declared structure, its initial and terminal positions satisfy the same-endpoint condition.

For the psychology reader, the distinction can be compressed into one sentence:
\begin{center}
\emph{Similarity asks, ``Does it look alike?'' Return asks, ``Has it come back to the same relevant position?''}
\end{center}

\begin{correspondence}[return and ``again'']
\[
\boxed{
\begin{gathered}
\text{an episode is recognized as ``I have come back here again''}\\[0.25em]
\strsim\\[0.25em]
\text{a finite traverse completes a return}\\[-0.05em]
\text{at the same relevant root}
\end{gathered}}
\]
\end{correspondence}

This is a Level II mapping. Mathematical endpoint equality is a strict relation. On the psychological side, by contrast, ``the same relevant position'' must be operationalized through segmentation, self-reference, memory, and task scale; it cannot be obtained directly from the mathematical definition.

\begin{plainbox}[Mathematical terminology boundary]
The \emph{return word} in the source growth paper is more specific than the finite witness borrowed in this chapter: it also involves a generator word and an admissible boundary context. Part II has already separated these notions. For psychological translation, this chapter borrows only the relational skeleton of a \emph{closed birooted excursion}. When later chapters enter return language, incompressibility, or Kuang's growth theorem, the source-domain return words must be restored in full.

Here ``closed'' is strictly combinatorial: the designated endpoints coincide. It is not topological closedness and it carries no psychological completion sense.
\end{plainbox}

\section{The Route Is Traveled Internally; ``Coming Back'' Is Determined by a Coordinate}

The internal process of a return can be utterly ordinary: waiting, checking, bodily tension, imagined dialogue, continuing to work, checking again. What gives it the meaning of ``coming back'' is not those internal contents themselves, but the entrance and exit being recognized by the subject's history as the same relevant position.

The mathematical picture is simple:
\[
b\longrightarrow a_1\longrightarrow\cdots\longrightarrow a_k\longrightarrow b.
\]

Most of the path can occur within routine internal activity; the return-completion condition is determined by boundary/root organization.

The psychological side has the same relational form. In the middle of a relational cycle, a person may have very different thoughts and perform very different actions. Yet it is the final sentence
\begin{center}
\emph{``In the end I still came back to the position where I had to prove that I was not being slighted.''}
\end{center}
that organizes the long stretch of material as a return.

A sentence more suited to Part III is therefore:
\begin{center}
\emph{The route is traveled in experience; history gives the route its coordinate of coming back.}
\end{center}

This also explains why bare AF-like repetitivity is never enough. Local patterns may recur, but only boundary/root organization tells us which recurrence the subject has admitted into the grammar of ``again.''

\section{``The Same Place'' Never Means That Everything Is the Same}

A mistake is easily made in everyday narrative here: once I say ``again,'' it can sound as though today must be identical to the past. The opposite is true. A recognized return requires only that a relevant coordinate be identified as completing a return; many other variables may change.

The same structure ``rapid self-doubt after criticism'' may recur, while this time you still complete the response. The same structure ``the other person did not reply'' may recur, while this time you do not continue checking. The same structure ``I cannot follow an unfamiliar talk'' may recur, while this time you record two connecting points that later enter your own research. At a coarse scale these episodes may still count as returns, even though their finer continuations have already diverged.

Mathematically, this is exactly why finite return-completion and higher germ resolution can be separated: downstairs one can come back to the same $b$, while upstairs the lift can end on a different sheet. Chapter 20 will formally discuss \emph{return residue}.

This chapter retains only one psychological discipline:
\begin{quote}
\textbf{When you say ``again,'' ask at the same time: what exactly has been recognized as the same, and what has in fact become different?}
\end{quote}

If ``again'' preserves only sameness while systematically erasing difference, pattern recognition begins to slide into narrative overcompression.

\section{``Again'' Is Both a Capacity for Learning and a Narrative Risk}

Cross-temporal classification is indispensable. It lets people discover regularities, predict risk, maintain commitments, and sustain an autobiographical self that persists across time. The second order cannot function if every event remains a completely isolated occurrence.

But compression has a cost. Three very different academic setbacks can quickly become ``I failed again at the critical moment.'' This sentence may identify a genuine preparation problem. It may also collapse a rejected paper, an unfamiliar talk, and one stalled proof attempt into a single identity pattern.

The escalation often has a recognizable form:
\[
\begin{gathered}
\text{this happened}\\[-0.15em]
\downarrow\\[-0.15em]
\text{I did this again}\\[-0.15em]
\downarrow\\[-0.15em]
\text{I always do this}\\[-0.15em]
\downarrow\\[-0.15em]
\text{this is who I am}.
\end{gathered}
\]

This chain is not inevitable, and the sentence farther to the right is not automatically more false than the sentence to its left. It only reminds us that \textbf{each step to the right carries a greater burden of generalization.}

The Geometry of Existence is particularly important here. Structural overreach can begin at the moment when ``again'' first acquires excessive explanatory authority: a finite return-completion originally says only that a relational skeleton has reappeared, yet it is rapidly promoted into a judgment about character, fate, or the other person as a whole \citep{geometry2026}.

Third-Order Consciousness does not forbid pattern recognition. It makes the radius of generalization visible again.

\section{A Return Hypothesis Supported by Evidence, and an ``Again'' Sustained Only by a Label}

Psychology cannot directly observe an absolute endpoint equality in the same way as mathematics. Rather than declaring a particular ``again'' genuine or false, it is more useful to distinguish what evidential structure currently supports it.

\begin{definition}[witness-supported psychological return hypothesis]
Relative to a declared segmentation of experiential material, comparison scale, and endpoint-identification rule, we call an ``again'' judgment a \emph{witness-supported return hypothesis} if it can identify the current and past episodes, the relevant position treated as the same, and the behavioral, mnemonic, or relational material that makes this return-completion checkable.
\end{definition}

\begin{definition}[label-only return hypothesis]
We call an account a \emph{label-only return hypothesis} if it compresses several episodes directly into the same pattern mainly because of an already available broad label---for example, ``rejection,'' ``need for control,'' or ``I am simply not good enough''---but cannot specify which endpoints, process relations, or observable differences support that classification.
\end{definition}

These names are not another packaging of ``true'' and ``false.'' A witness-supported hypothesis can be weakened by new material; a label-only hypothesis can later acquire genuine evidence. The distinction serves only to turn ``I am doing this again'' from an opaque complete judgment back into a hypothesis that can continue to be tested.

\section{Example One: Does No Reply Really Mean ``I Have Been Rejected Again''?}

The current fact remains only this: the message has not received a reply.

If several past episodes share a similar relational structure---active investment; interaction becomes uncertain; the subject rapidly shifts into checking; silence is interpreted as a question of status; a corrective dialogue forms in the mind; and finally the subject withdraws---then ``again'' does have a structural candidate worth testing. The key is not that the past also ``felt bad,'' but that several entrance, process, and exit relations show a repeatable organization.

Conversely, if the supposed history has been reduced to the umbrella label ``other people do not respect me,'' while the past material includes actual rejection in one case, illness in another, different communication habits in another, and perhaps no established commitment at all in yet another, then ``again'' may already have flattened the differences too early.

Third-Order Consciousness does not decide for the subject whether contact should continue. It adds only one question:
\begin{quote}
\textbf{When you say ``again,'' which stretch of history have you brought into the judgment, and which other parts of the past have you left out?}
\end{quote}

This directly echoes the self-memory system: the current working self participates in determining which autobiographical memories are easier to access \citep{conway2000}. ``Again'' therefore does not merely connect the past; it also \emph{selects the past}.

\section{Example Two: Does Failing to Follow a Talk Mean ``I Am Trapped by My Limitations Again''?}

Imagine that a person is listening to a talk in an unfamiliar area and temporarily cannot recall several key definitions. At the level of the current event, this is first a local failure of accessibility.

If the person's past shows a repeatable organization
\[
\begin{gathered}
\text{unfamiliar terminology}\\[-0.15em]
\downarrow\\[-0.15em]
\text{attention loses its foothold}\\[-0.15em]
\downarrow\\[-0.15em]
\text{global self-devaluation}\\[-0.15em]
\downarrow\\[-0.15em]
\text{withdrawal from asking questions}\\[-0.15em]
\downarrow\\[-0.15em]
\text{retreat to one's own narrow research area},
\end{gathered}
\]
then ``I have come back to the position of being limited by my knowledge again'' can become a return hypothesis.

But suppose this occasion was simply marked by lack of sleep, or the talk depended heavily on unstated background, while the subject still retained two new questions and continued following the literature afterward. To place this episode directly inside the old pattern ``my knowledge is too narrow'' may instead erase the most important new continuation.

Return mapping helps psychology ask a finer question here: \textbf{which variables govern same-event similarity, and which govern the judgment that the same self-pattern has recurred?}

\section{Why Return Creates a Compressible History}

Once a return category has been established, the subject no longer has to preserve every detail anew each time. A long process can be compressed into a macro-symbol:
\begin{center}
\emph{``This is that uncertainty--checking--self-evaluation--withdrawal loop.''}
\end{center}

This is part of the economy of pattern recognition, and one of the conditions under which narrative identity can sustain persistence across time. A life in which every decision required rereading the whole past from scratch could scarcely maintain commitments, learning, or long-term plans.

But a macro-symbol also selects the future. As soon as a current event is recognized as ``this again,'' attention begins preferentially seeking the evidence required by the old pattern; memories consistent with the pattern become easier to retrieve; and certain actions can begin to look justified earlier. Return is therefore not only a summary of the past. It can also become an entrance into future continuation.

This is why ``again'' is so small and yet so heavy. It turns an episode into part of history, while allowing history to re-enter the current episode.

\section{New Context Can Rewrite an Interpretation, but It Does Not Remove the Burden of Evidence}

Psychological memory forgets and reconstructs; segmentation of past episodes can genuinely change in light of new information. For this reason, we do not claim that a psychological return has the same absolute empirical persistence as a mathematical witness.

It retains only a weaker research discipline: if a return hypothesis was once supported by explicit material and a larger context later leads us to withdraw it, we should state which relation has been overturned.

For example, later learning that the other person was hospitalized may invalidate the earlier endpoint identification ``silence = rejection.'' That is genuinely new evidence. Conversely, merely preferring a new self-story today does not justify declaring that the past checking, withdrawal, or self-devaluation never occurred. That too would be narrative overreach.

The value of mathematical persistence in Part III is not to tell a person that ``the past can never be rewritten.'' It is to propose an ethic of evidence:
\begin{center}
\emph{Interpretations may be updated;}\\[0.2em]
\emph{when a structural judgment is withdrawn, disclose what new material disqualified the old witness.}
\end{center}

\section{If This Interface Works, What Should Psychology Be Able to Observe?}

If the return/``again'' mapping is more than elegant language, it should motivate concrete research questions:
\begin{itemize}
\item After controlling content similarity, does manipulating episode endpoints or the self-reference frame systematically change ``again / same pattern'' judgments?
\item Before participants say ``again,'' are autobiographical memories consistent with the current pattern more readily retrieved, while structurally inconsistent memories are more likely to be neglected?
\item Does an ``again'' judgment predict subsequent behavior, expectation, or responsibility judgments better than a simple similarity rating?
\item Does placing the same event in a different historical context change whether it is recognized as a return?
\item Can we distinguish an ``again'' judgment supported by an explicit relational witness from one driven only by a ready-made identity label? Do the two have different consequences for memory, prediction, and behavior?
\item If endpoint return-completion has no explanatory value for any additional psychological variable, should the return mapping be reduced to ordinary similarity / categorization language?
\end{itemize}

Only in this way does mathematical return genuinely give something back to psychology: it forces the sentence ``I am doing this again'' to be decomposed into relations that can be manipulated and measured separately, rather than adding mystery to it.

\section{Concept Passport}

\begingroup\small
\begin{longtable}{>{\raggedright\arraybackslash\bfseries}p{0.25\textwidth}>{\raggedright\arraybackslash}p{0.66\textwidth}}
Psychological motivation & How a current episode is pulled into the past and assigned to the same self-history; how that classification can both support learning and produce narrative overcompression.\\
Psychological concepts & ``again'' judgment; autobiographical retrieval; self-indexed pattern recognition; narrative compression.\\
Mathematical concepts & finite return witness / closed birooted traverse; a source return word retains its generator/context conditions.\\
Preserved relations & The episode has already been segmented; relative to a declared coordinate, the endpoint satisfies a same-endpoint return-completion condition; recognized return is stronger than general similarity.\\
Discarded content & generator spelling; tile geometry; specific mechanisms of memory; emotional content; clinical meaning; neural implementation.\\
Translation level & Level II; the strictness of the mathematical same-endpoint condition does not automatically strengthen the psychological correspondence.\\
Candidate observables & ``again'' / same-pattern judgments; autobiographical retrieval; endpoint manipulation; self-reference frame; similarity ratings; subsequent prediction; behavioral choice.\\
Non-example & two episodes having the same emotion; literal events being similar; purely periodic behavior; or a strong identity label, and therefore automatically constituting a return.\\
Failure conditions & ``Again'' is fully explained by ordinary content similarity; endpoint/self-history structure has no additional predictive value; a return hypothesis yields no checkable relation among episodes.\\
Reverse-translation check & Do not use return to replace autobiographical memory, categorization, schema, narrative identity, or other established psychological constructs. The mapping adds only the relational question: when is similarity organized into self-indexed return-completion?\\
Translation remainder & The mechanism of the subjective sense of sameness; the actual neural processes of memory reconstruction; how pattern categories develop; and how social and cultural contexts determine which repetitions are worth naming ``again.''\\
\end{longtable}
\endgroup

\section{Chapter Conclusion: ``Again'' Is a Hinge in Time}

What this chapter is meant to leave behind is not the claim that psychological ``again'' equals mathematical return. It is something finer.

The present could simply have remained the present. A silence, a mistake, or a familiar wave of bodily tension could occur and then pass. But once ``again'' appears, the present folds back into history: earlier episodes are recalled, differences are filtered again, a relevant position is recognized as reappearing, and the current experience is no longer merely this moment. It becomes a new member of ``my pattern.''

This is the power of ``again,'' and also its risk.

It allows learning: ``there really is a regularity here.'' It can also foreclose too quickly: ``see, I really am always like this.'' Third-Order Consciousness does not abolish the word. It turns it from an automatic verdict back into a relation that can be inspected:
\begin{center}
\emph{similarity does not automatically mean coming back;}\\[0.25em]
\emph{``again'' means that the subject compresses present and past into the same relevant coordinate;}\\[0.25em]
\emph{and every such compression should preserve both what is the same and what is different.}
\end{center}

``Again'' is therefore a hinge in time: it folds the present toward the past, while allowing the past to regain the power to enter the present.

The next chapter asks a more difficult question. A person's history never contains only one kind of ``again.'' ``Being ignored,'' ``losing control,'' ``having to prove myself,'' and ``I am not good enough'' can all compete for explanatory authority within the same trajectory. At that point a single return is no longer enough. We need to study an entire return grammar, and what it means for a stretch of experience to avoid, for a time, every already available way of coming back.

\chapter{Multiple ``Agains'' and Temporarily Incompressible Experience}
\label{chap:multiple-agains-en}

\par\medskip
\noindent\textbf{Why this chapter matters}\par
A person rarely has only one ``again.'' The same unanswered message can be organized as ``I am being ignored again,'' ``I am losing control again,'' ``I cannot stop checking again,'' or ``I have again become the person who needs to prove that I am fine.'' The facts may be the same, yet different parts of the past are brought back, different endpoints matter, and different continuations become available.

Here we ask what happens when several return grammars overlap, nest, or compete for explanatory authority. The mathematical notion of incompressibility enters only after these distinctions are fixed. Here it never means that a person is stubborn, pathological, or impossible to understand. The psychological interface is much more cautious: relative to a declared grammar, the current experience may remain \emph{temporarily unchunkable by return}.
\par\medskip

\begin{methodbox}[Interface card: multiple return grammars and temporary unchunkability]
\textbf{Psychological motivation:} Rumination, checking, relational cycles, and self-evaluation often contain more than one repetitive structure. Looking only for one ``root pattern'' can compress overlapping return-completions into a single story.

\textbf{Psychological native domain:} repetitive thought; overlapping self-patterns; ``again'' judgments; model-relative temporary unchunkability.

\textbf{Mathematical native domain:} the source-domain return language $\mathcal R$ collects admissible return factors; incompressibility requires the entire factor set to avoid $\mathcal R$.

\textbf{Preserved relation:} whether a long trajectory contains a return must be assessed across different, overlapping, and nested local factors, rather than by inspecting only one germ or one narrative.

\textbf{Level:} Level II. Only a future module with an explicit psychological coding alphabet, episode factors, and formal return rules could move locally to Level III.

\textbf{Not claimed:} rumination, trauma, suffering, stubbornness, or ``being unable to stop'' is itself mathematical incompressibility.
\end{methodbox}

\section{Psychological Entry Point: Several ``Agains'' May Coexist within One Stretch of Repetitive Thought}

Complicated repetition is easily compressed into one sentence in everyday language: ``I am spiraling again.'' ``I am trapped in this relationship pattern again.'' ``This old problem of mine is acting up again.'' Such sentences are efficient because they immediately give disorderly experience a name. Structural differences can also be erased by the same efficiency.

Research on repetitive thought and rumination has long warned that ``thinking something over and over'' is not a single process. Watkins emphasizes that the consequences of repetitive thought vary with content, level of abstraction, context, and mode of processing; repetition can be constructive in some settings and rigid or depleting in others \citep{watkins2008}. Nolen-Hoeksema, Wisco, and Lyubomirsky likewise review the complex relations among rumination, negative affect, problem solving, and social functioning \citep{nolen2008}.

This chapter does not redefine rumination. It extracts a narrower question alongside those mature literatures:
\begin{quote}
\textbf{When a person says, ``I keep going around the same circle,'' is there really only one circle?}
\end{quote}

Imagine an evening spent waiting for a reply. It may contain all of the following:
\begin{itemize}
\item \textbf{a checking cycle:} uncertainty $\to$ open the phone $\to$ no new message $\to$ put it down briefly $\to$ uncertainty rises again;
\item \textbf{a relationship-status loop:} silence $\to$ ``am I unimportant?'' $\to$ search for earlier episodes of being ignored $\to$ greater confidence that the present silence has status meaning;
\item \textbf{a rule loop:} ``a normal person should reply'' $\to$ imagine correcting the other person $\to$ the rule has still not been acknowledged $\to$ return to the corrective impulse;
\item \textbf{a meta-evaluation loop:} notice that one is still thinking $\to$ ``why have I not let go yet?'' $\to$ self-management $\to$ one's attention returns to the matter again.
\end{itemize}

They share one evening, but they are not the same return. One action may even belong to several grammars at once. Third-Order Consciousness therefore does not begin by hunting for ``the deepest one.'' Its first task is to \textbf{prevent the loudest cycle from swallowing other structures that are operating at the same time}.

\section{Why One ``Big Cycle'' Is Often a Premature Compression}

Suppose we draw one enormous circle:
\[
\text{being ignored}
\longrightarrow
\text{pain}
\longrightarrow
\text{rumination}
\longrightarrow
\text{being ignored again}.
\]
A great deal can now appear to have been explained. Yet almost no resolution has been gained. Why does checking stop on some occasions but not others? Why does the same silence produce anger in one episode and self-devaluation in another? Why can an internal corrective dialogue continue after the decision not to contact the person again has already been made?

The phrase ``a long-term pattern'' may therefore hide several overlapping local return-completions rather than one totalizing root circle. A checking cycle may have one exit condition; a status cycle another; a rule-enforcement cycle still another. Compressing them all into a single ``old problem'' makes the story smoother while reducing the number of places at which the process can actually be distinguished.

The methodological value of the \emph{Geometry of Existence} is especially clear here. The search for a root cause can itself become a form of structural overreach: a locally useful return acquires more explanatory authority than its evidence supports. The more cautious strategy is to keep several cycles visible at once, while also preserving whatever material does not currently belong to any recognized loop \citep{geometry2026}.

\section{Psychological Return Grammar: Which ``Agains'' the Current Model Recognizes}

To keep these different ``agains'' separate, the psychological side needs a working bookkeeping device.

\begin{definition}[Psychological return grammar]
Relative to a declared segmentation of experiential episodes, an observation scale, and endpoint-identification rules, let $\widehat{\mathcal R}$ denote all psychological return patterns that the current model can support with finite witnesses. We call $\widehat{\mathcal R}$ the \emph{current psychological return grammar} of the model.
\end{definition}

The words \emph{current model} matter. $\widehat{\mathcal R}$ is not a final dictionary secretly stored in the brain. It is a candidate grammar that may change when data are added or removed, episodes are resegmented, or endpoint-identification rules are revised.

The same episode can meet several patterns at once. A shorter return may sit inside a longer one; two returns may share part of their material; one return may begin before another and end inside it. Experience therefore need not look like a set of boxes in which each item receives one label. It may be better represented, at this stage, as several translucent relational nets laid over the same material.

\begin{plainbox}[Do not mistake a grammar for a diagnostic classification]
Labels such as ``ignored return,'' ``control return,'' ``self-evaluation return,'' or ``checking return'' are working names for hypotheses about relations. They are not diagnoses, personality types, or established psychological modules. The data that matter are which episode is being tracked, which endpoints are being identified, what is being recognized as the same, and how the next step changes.
\end{plainbox}

\section{Why Mathematics Enters Here: It Forces Us to Inspect All Factors, Not Only the Most Salient Loop}

Part II fixed a source-domain return language $\mathcal R\subseteq S^*$. Every accepted return factor satisfies the relevant admissibility conditions and, when an occurrence is being tracked, is supported by a concrete witness. For a long word $w$, the question is whether \emph{any} contiguous factor belongs to $\mathcal R$:
\[
\operatorname{Fac}(w)\cap\mathcal R
\stackrel{?}{\neq}
\varnothing.
\]

Translated back to the psychological side, the discipline is simple. If one claims that an entire stretch of experience contains no recognizable return-completion, it is not enough to show that the most familiar loop is absent. Other possible returns must also be checked.

This chapter does not redefine the mathematical $\mathcal R$. The psychological side always uses the hatted notation $\widehat{\mathcal R}$; generator words, germ types, and admissibility remain in their mathematical native domain. What crosses the interface is only the relational discipline that \textbf{a global judgment must audit all relevant local factors}.

\section{A More Direct Mathematical Prototype: A Return Grammar Can Already Be Jointly Organized by Multiple Boundary Sectors}

The psychological examples above can make ``multiple agains'' look like an interpretive layer added only because human experience is complicated. There is, however, a more direct source-side prototype. In Kuang's preprint v1, Subsection~5.2, return structures already appear with multiple boundary points and multiple bridge types \citep{kuang2024v1}.

There, a high-level traverse can induce a sequence of lower-level traverses
\[
\Phi(\tau)=(\tau_1,\ldots,\tau_k),
\]
together with a bridge-type word
\[
x_1x_2\cdots x_k,
\]
where different $x_i$ may correspond to different infinite boundary points. A return can therefore be witnessed by a type word passing across several boundary sectors. One need not first choose a single singular root and assign every return to it.

This permits one important structural move here. If a psychological model has predeclared a family $\Sigma$ of organizing positions or sectors, one may provisionally keep candidate return grammars on separate ledgers, for example
\[
\widehat{\mathcal R}
=
\bigcup_{\sigma\in\Sigma}
\widehat{\mathcal R}_{\sigma},
\]
or, more generally, record finite type sequences that pass across sectors. The latter is often closer to the source-side multi-boundary picture: the relevant information may be the ordered sequence of organizing sectors through which an episode chain passes, rather than a simple union of independent grammars.

\begin{warning}[Do Not Directly Rename Psychological Sectors as Singularities]
The psychological sectors in this chapter are still Level II organizing positions. A checking sector, relationship-status sector, rule-evaluation sector, or meta-evaluation sector is not literally a mathematical boundary point, and it does not automatically define a different orbit or a different singularity in a groupoid of germs.

The preserved relation is narrower: a complete return grammar can be jointly organized by several local organizing positions, and there is no structural reason to grant one of them unconditional global authority.
\end{warning}

\begin{principle}[Multi-Boundary Nesting Principle: Why Different ``Agains'' Can Be Nested]
A long trajectory may carry several families of returns. Suppose two return intervals are
\[
I_1=[i,j],
\qquad
I_2=[k,\ell],
\]
organized by different boundary/germ types or by different sector-passage words. They can be nested,
\[
i<k\le \ell<j,
\]
overlap,
\[
i<k\le j<\ell,
\]
or appear in more complicated interleavings. Multi-boundary and multi-germ geometry therefore provides a source-domain prototype for heterogeneous nested return structure.
\end{principle}

\begin{warning}[The Converse Does Not Hold]
Do not infer
\[
\text{nested returns}
\quad\Longleftrightarrow\quad
\text{multiple singular boundary germs}.
\]
Multiple singular boundary germs provide a natural source of heterogeneous return families, but they are not necessary for nesting. A single germ type or a single return grammar can produce nested returns at different scales. Every return must still be located relative to a declared boundary/germ/sector and scale.
\end{warning}

Some returns therefore occur within one sector, while others become visible only after the passage from one organizing sector to another is recorded.

\begin{plainbox}[Repeated returns can also help detect interfaces]
Multiple return families are not useful only after a boundary has already been identified. They can also act as probes. If independent return hypotheses repeatedly place endpoints, position switches, or sector passages near the same transition, that transition becomes a stronger candidate for a boundary-like interface.

Likewise, if a larger return and a nested smaller return repeatedly share an entrance or exit, one should inspect whether local organization loses sufficiency there again and again. The concentration of returns supplies detection evidence; it does \textbf{not} define the boundary. The same germ or type can generate nesting at several scales, and a single return-completion can occur accidentally.
\end{plainbox}

\section{The Most Counterintuitive Point: A Great Deal of Repetition Often Means Precisely That the Process Is Compressible}

Everyday psychological language easily aligns three ideas: something happens repeatedly, the person ``cannot stop,'' and the process is therefore ``incompressible.'' Mathematical incompressibility points almost in the opposite direction.

Suppose a checking sequence clearly repeats
\[
\text{uncertainty}
\longrightarrow
\text{checking}
\longrightarrow
\text{brief relief}
\longrightarrow
\text{uncertainty returns}.
\]
The sequence already contains a conspicuous return. It may be exhausting and difficult to exit, yet structurally it is easy to compress: a long string of actions can be chunked as ``the checking cycle.''

Hence
\[
\boxed{
\text{high repetition}
\neq
\text{incompressibility}.}
\]
A highly repetitive process may in fact be highly compressible.

This distinction matters psychologically because it prevents a mathematical word from quietly becoming a synonym for ``severe,'' ``stubborn,'' ``distressing,'' or ``emotionally impossible to escape.'' Intensity of suffering, difficulty of exit, and structural compressibility are three different ledgers.

\section{What, Then, Is Psychological Incompressibility Trying to Capture?}

In the mathematical native domain, for a generator word $w$,
\[
\operatorname{Fac}(w)\cap\mathcal R=\varnothing
\]
means that none of its contiguous factors is an admissible return word. This is global incompressibility.

The psychological side cannot simply copy that definition. There is no generally accepted psychological alphabet here, and no God's-eye complete return grammar. At the present Level II, the interface therefore uses a more cautious working concept.

\begin{definition}[Model-Relative Temporary Unchunkability]
Relative to a declared segmentation of experiential episodes, rooting convention, observable variables, and psychological return grammar $\widehat{\mathcal R}$, we say that a trajectory is \emph{temporarily unchunkable by return under the current model} if all subepisodes that the current model permits us to inspect have not yet acquired a witness-supported return.
\end{definition}

The word \emph{temporary} is deliberate. Failure of an old grammar to compress an experience has at least two very different possible explanations:
\begin{enumerate}
\item \textbf{The experience may contain genuine structural novelty.} Old loops do not cover the current change, and the new continuation deserves to be preserved rather than forced back into an earlier pattern.
\item \textbf{Our observational language may be too coarse.} The episode may have been cut incorrectly; relevant variables may have been omitted; endpoints may not have been measured; or the researcher may simply recognize only familiar returns.
\end{enumerate}

Psychological incompressibility therefore first functions as an alarm about the current explanatory power of the model, not as a judgment about the person.

\begin{correspondence}[Global Factor Avoidance and Model-Relative Temporary Unchunkability]
At the present Level II, the preserved relation is
\[
\begin{gathered}
\boxed{
\begin{gathered}
\text{relative to a declared grammar, all checkable local segments}\\[-1mm]
\text{still lack a witness-supported return}
\end{gathered}}
\\[1mm]
\strsim
\\[1mm]
\boxed{
\begin{gathered}
\text{all contiguous factors of a word}\\[-1mm]
\text{avoid the return language}
\end{gathered}}
\end{gathered}.
\]
This correspondence does not rename trauma, rumination, personality rigidity, or ``being unable to let go'' as incompressibility.
\end{correspondence}

\section{A Small but Important Formal Discipline: If the Whole Has No Return, a Subwindow Cannot Secretly Contain One}

Part IV will state the corresponding mathematical factoriality formally. Here we borrow only a consistency requirement for model-relative judgments.

If a whole trajectory is declared to contain no return under a given grammar, one cannot later admit that a subwindow satisfies that same grammar and still preserve the original global no-return claim. A global statement must remain accountable to its local factors.

Many sweeping self-descriptions fail precisely here. ``I made no progress all year.'' ``This relationship never reached a genuine return-completion.'' ``I have always been unable to get out of any pattern.'' Once the timeline is enlarged and inspected locally, smaller episodes may reveal return-completions that had been erased by the larger narrative.

This use of factoriality does not force psychological reality to obey an algebraic property. It imposes a discipline on the model: \textbf{a global narrative must remain open to audit by local material}.

\section{Overlapping Returns: Why the Same Action Can Belong to Several Cycles at Once}

Imagine opening the phone for the third time during the same evening. That single action may belong simultaneously to several structures:
\begin{itemize}
\item \textbf{checking cycle:} uncertainty again drives checking, and the person returns to a familiar behavioral position;
\item \textbf{relationship-status loop:} the action rechecks whether the other person treats the subject as important;
\item \textbf{rule loop:} the action gathers further evidence about whether the expected norm of replying has been violated;
\item \textbf{meta-evaluation loop:} the action confirms ``I am still checking,'' which in turn activates self-management.
\end{itemize}

The same behavior is therefore not the same return. The relevant endpoints, the parts of the past brought into the episode, the next-event predictions, and the exit conditions can differ. Counting phone openings gives behavioral frequency. It does not by itself reveal the self-indexed grammar under which each opening is being organized.

Repetition count is one-dimensional; overlapping return-completion structures may be multidimensional.

\section{Example from a Research Project: The Same Stretch of Work Can Be Cut by Two ``Again'' Grammars at Once}

A research project can carry more than one return grammar at the same time. One sequence may be
\[
\begin{gathered}
\text{excitement about a new question}
\longrightarrow
\text{technical details expand}\\[-1mm]
\longrightarrow
\text{doubt its value}
\longrightarrow
\text{return to excitement about a new question}.
\end{gathered}
\]
This is one return.

At the same time, another sequence may run through the same period:
\[
\begin{gathered}
\text{see other people's work}
\longrightarrow
\text{social comparison}\\[-1mm]
\longrightarrow
\text{doubt one's own range}
\longrightarrow
\text{retreat back into familiar territory}.
\end{gathered}
\]
The second sequence can pass through part of the first while answering a different question.

If both are compressed into ``I always doubt myself,'' the narrative becomes smooth but two different intervention points disappear. One process may call for control of project scope; the other may call for changing the comparison frame. This is the practical value of multiple return grammars for psychology: they do not guarantee a deeper explanation, but they may locate action more precisely.

\section{Persistent Incompressibility and Psychological Robustness Audits: Do Not Mix the Two Axes}

One failure to detect a return proves almost nothing. The current grammar may be too poor, the segmentation may be unsuitable, or the participant may not yet have formed a stable ``again'' judgment.

Part IV will define return filtration and persistent incompressibility inside a \emph{fixed} formal coding and presentation. As the permitted finite support level increases, source return words with admissible witnesses accumulate into $\mathcal R_{\le n}$. A coded word that continues to avoid this expanding return dictionary belongs to the corresponding persistent incompressible remainder.

The psychological side must keep a separate ledger. Changing episode granularity, adding observable variables, switching the self-reference frame, or even revising the return grammar is \textbf{not} the same as moving up one level in that mathematical filtration. These changes are better treated as a family of predeclared \emph{robustness audits}: the researcher actively changes modeling choices to test whether the original non-detection was merely produced by one segmentation or one variable set.

Third Order therefore borrows only a weaker epistemic discipline: if a judgment of the form ``there really is no old return-completion available here'' continues to hold under several predeclared and mutually comparable modeling schemes, it deserves more attention than a single non-detection under one encoding.

For example, one can separately audit:
\begin{itemize}
\item different but predeclared episode granularities;
\item additional observable variables relevant to the current question;
\item different self-reference frames;
\item overlapping return grammars rather than a search for one root loop.
\end{itemize}

If all these robustness audits still fail to produce a witness-supported return-completion, structural novelty deserves to be protected more seriously. This remains a robustness judgment about a psychological research protocol. It is not the formal persistent incompressibility of Part IV.

One stop line is especially important:
\begin{quote}
\textbf{Failure to be recognized by an old pattern does not mean that the experience is pathological. Sometimes the old pattern is precisely what should give way.}
\end{quote}

This is also what separates Third-Order Consciousness from a simple technique of pattern recognition. Third Order can expose hidden returns, but it must equally allow reality to exceed an existing grammar.

\section{When the Model Cannot Compress the Experience, Add Variables Before Adding Grand Labels}

If a trajectory remains unchunkable under the current grammar, the most dangerous response is to invent larger abstractions immediately: ``this is deep trauma,'' ``this is my core personality,'' ``this is a fate pattern.'' Such concepts may eventually be useful, but if they add no new observable variables, they make the explanation fuller without making the structure clearer.

A response more consistent with the method of this book is
\[
\boxed{
\begin{gathered}
\text{refine observables}
\longrightarrow
\text{re-segment / re-root}\\[-1mm]
\longrightarrow
\text{re-test returns}
\longrightarrow
\text{hidden return-completion or protected novelty}.
\end{gathered}}
\]

Adding variables means something concrete. ``Distress'' can be separated into action choice, attention shift, memory retrieval, and next-event prediction. ``The same pattern'' can be separated into endpoint judgment, self-reference frame, and next-step expectation. Refinement earns its place only when the new variable changes continuation.

The \emph{Geometry of Existence} again supplies a plain methodological direction: \textbf{when an explanation is insufficient, what is needed is not necessarily a larger meaning; often reality first needs a passage to reopen} \citep{geometry2026}.

\section{If This Interface Holds, What Should Psychology Be Able to Observe?}

If the return-grammar / incompressibility interface adds value, it should at least generate empirical questions of the following kinds:
\begin{itemize}
\item Within the same stretch of repetitive thought, can participants stably distinguish two or more overlapping ``again'' patterns rather than reporting only one overall rumination intensity?
\item After behavioral frequency and negative affect are controlled, do different return grammars predict different next-event expectations, memory retrieval, or action choices?
\item Are subepisodes that participants identify as returns easier to chunk, compress in memory, or use to predict subsequent structure?
\item When segmentation, self-reference frame, or observable variables change, can a trajectory previously judged unchunkable acquire a new witness-supported return?
\item If all putative return grammars are completely explained by existing variables from rumination, habit, appraisal, schema, or event segmentation, should the psychological claims made in AF-by-discrete language be reduced?
\end{itemize}

Watkins's emphasis on different modes of repetitive thought is a useful reminder here: repetition by itself does not determine structural function \citep{watkins2008}. If the present project cannot add measurable relations beyond distinctions already available in those literatures, it should not preserve extra terminology merely because the mathematics is elegant.

\section{Concept Passport}

\begingroup\footnotesize
\begin{longtable}{>{\raggedright\arraybackslash\bfseries}p{0.25\textwidth}>{\raggedright\arraybackslash}p{0.66\textwidth}}
\toprule
Psychological motivation & A stretch of repetitive experience may contain several overlapping or nested ``agains''; one ``big cycle'' can obscure different predictions, memory recruitment, and action exits.\\
Psychological concept & Repetitive thought; overlapping return grammars; ``again'' judgments; model-relative temporary unchunkability.\\
Mathematical concept & Return language; factor avoidance; relative/global incompressibility; factoriality.\\
Preserved relation & A global judgment must audit all relevant factors; different returns can overlap, nest, or interleave; absence of return and high-frequency repetition are different properties.\\
Discarded content & Group-growth conclusions; psychopathology classification; intensity of suffering; clinical diagnosis; literal realization of mathematical generators or germs.\\
Translation level & Level II; a local module could move to Level III only after an explicit psychological coding alphabet and formal return rules are specified.\\
Candidate observables & Overlapping ``again'' judgments; episode segmentation; chunking; memory retrieval; next-event prediction; action choice; stability of return judgments under refinement.\\
Non-example & Calling rumination, checking, trauma, stubbornness, or ``being unable to stop'' incompressible; inspecting only one salient loop and then declaring global avoidance.\\
Failure conditions & Return grammars are entirely retrospective labels; different grammars make no differential prediction for any observable; alleged unchunkability disappears arbitrarily under slight resegmentation.\\
Reverse-translation check & Do not replace mature constructs such as rumination, repetitive thought, habit, schema, or appraisal with return language. The added question is only how several self-indexed return-completions overlap and when no usable return-completion remains.\\
Translation remainder & The actual coding alphabet of the psychological system; how return categories develop; how social and cultural settings select which cycles are worth naming; neural and bodily mechanisms.\\
\bottomrule
\end{longtable}
\endgroup

\section{Chapter Conclusion: Do Not Rush to Compress Complex Experience into One ``Old Problem''}

What this chapter tries to preserve is one more degree of resolution in the study of repetition.

An experience can repeat conspicuously and still be easy to compress into a familiar loop. It can look confused because several grammars overlap. And some stretches, even after scale is changed, roots are reconsidered, or additional variables are introduced, may still fail to return to any old position recognized by the current model.

The last case must be handled cautiously. It may reveal the poverty of the model. It may also protect something genuinely new in the experience. In Part III, incompressibility must therefore never become a heavy personality label. The psychological interface is closer to a region that, for now, refuses to be chunked by an old story.

Third-Order Consciousness asks three questions at once:
\begin{quote}
\textbf{Are there other ``agains'' that we have not yet seen?}\
\textbf{Did we segment the episode incorrectly?}\
\textbf{Or, this time, did reality genuinely fail to return?}
\end{quote}

All three possibilities must remain open.

Chapter 19 adds scale to the problem. At what resolution is a return first recognized? As context expands, how should an earlier return-completion be preserved, revised, or reinterpreted? Only then does filtration acquire its psychological meaning.

\chapter{Scale, Memory, and Narrative Rewriting: How the Past Is Seen Anew Later}
\label{chap:scale-memory-narrative-en}

\par\medskip
\noindent\textbf{Why this chapter matters}\par
Years later, a person may come to understand an old event differently. An experience once described as ``I was rejected'' may later be understood as ``we had never established shared expectations in the first place.'' A talk once remembered as ``I failed completely'' may look different when old notes reveal that two key questions were, in fact, understood. This raises a difficult question: when later understanding becomes deeper, does that mean we now see the past more clearly, or that the past itself has been reorganized?

Psychology has long known that memory is not a recording stored in a warehouse waiting to be retrieved. Present goals, self-knowledge, and narrative frames all participate in the construction of a particular act of remembering \citep{conway2000,bruner1987}.

For this reason, we do not translate mathematical \emph{persistence} as ``memory never changes.'' It borrows inflation, filtration, and telescoping only to separate several operations that are often blended together: adding context, discovering new structure, changing interpretation, resegmenting events, and actually changing the material that is remembered.
\par\medskip

\begin{methodbox}[Interface card: context expansion, filtration, and narrative rewriting]
\textbf{Psychological motivation:} When more background becomes available later, or when a person remembers or retells an event again, which elements have merely been placed in a larger context, which interpretations have changed, which episode boundaries have been redrawn, and which remembered details themselves have been reconstructed?

\textbf{Psychological native domain:} autobiographical-memory reconstruction; narrative revision; event resegmentation; context-dependent memory retrieval.

\textbf{Mathematical native domain:} a finite witness remains traceable along specified nested tile occurrences / support-preserving embeddings; source-return-word filtration accumulates as the permitted support level increases; telescoping changes the presentation while preserving the declared structural relations.

\textbf{Preserved relation:} ``adding context'' and ``rewriting existing material'' must be kept on separate ledgers; if a new coarse narrative claims to arise from the same underlying occurrences, it should state what the compression preserved.

\textbf{Level:} Level II.

\textbf{Not claimed:} human memory accumulates monotonically according to a filtration, or a newer narrative is necessarily closer to the truth than an older one.
\end{methodbox}

\section{Psychological Entry Point: ``I Finally Understand It Now'' Can Contain Four Entirely Different Changes}

When looking back on the past, it is easy to say:
\begin{quote}
``Now I finally understand what really happened then.''
\end{quote}

The sentence sounds as though the same past has simply been viewed at higher resolution. Psychological reality is much less orderly. At least four different changes may be compressed into that one sentence:

\begin{center}
\begin{tabularx}{0.94\textwidth}{>{\bfseries\raggedright\arraybackslash}p{0.235\textwidth}Y}
\toprule
Change & What has actually changed psychologically \\
\midrule
Added context & The original event material remains largely the same, while new antecedents, consequences, information about other people, or longer-term background enter. \\
Reinterpretation & The same main occurrences acquire new causal, evaluative, or identity-related meaning. \\
Resegmentation & We begin to think the event actually started earlier, ended later, or should be divided internally into several different episodes. \\
Memory reconstruction & The specific content, sequence, details, or even confidence that is recalled changes. \\
\bottomrule
\end{tabularx}
\end{center}

All four can produce the feeling that ``the past has changed,'' but they cannot be represented by the same formal model.

What \AFBD\ contributes most usefully here is an exceptionally clean contrast. If an already existing finite configuration is merely placed inside a larger context, the original finite configuration does not automatically disappear because the background has grown.

This is not a law of the brain. It is an analytic ruler. It lets us ask whether a given ``reinterpretation'' really adds context, or whether it has already altered the event structure that was previously taken to be present.

\section{Memory Is Not a Recording: The Present Self Also Participates in How the Past Comes Back}

Research on autobiographical memory is especially clear that ``remembering the past'' is not the opening of a static archive. Conway and Pleydell-Pearce's self-memory system emphasizes that the current goals of the working self participate in determining which autobiographical knowledge becomes more accessible and how a specific memory is constructed \citep{conway2000}. Bruner's discussion of narrative likewise stresses that people do not first possess a perfectly neutral record of a life and then add a story from outside; stories themselves participate in how a life is organized and understood \citep{bruner1987}.

Third Order therefore has to be very cautious about one temptation:
\[
\begin{gathered}
\text{``today's interpretation is more mature''}\\[-1mm]
\Longrightarrow\\[-1mm]
\text{``today I finally see the one true form of the past.''}
\end{gathered}
\]
That implication does not follow automatically.

The person you are today has new language, new relational experience, a new ordering of values, and a newly operative root. These may help reveal variables that were genuinely overlooked in the past; they may also make some materials especially easy to retrieve. Reinterpreting the past is itself an organizing act occurring in the present.

This does not push everything into relativism. On the contrary, it forces us to separate the evidence: which elements are traceable old materials, which are memory contents reappearing now, and which are interpretations added today.

\section{Why Mathematics Enters Here: Tile Inflation Provides an Ideal Contrast for ``Adding Context Only''}

On the formal side, a finite return witness can be carried into a larger tile context along a specified lower-level tile occurrence. Part IV will treat this copy-relative / occurrence-supported embedding rigorously. Here we do not assume a global canonical inclusion from an abstract level-$n$ tile into an arbitrary higher tile without specifying the copy. Part III borrows only one minimal idea:

\begin{plainbox}
\begin{center}
\emph{A larger context can contain more structure,\\
without rewriting the finite subconfiguration already embedded.}
\end{center}
\end{plainbox}

This supplies psychology not with a memory mechanism, but with a counterfactual standard. If we say that we have ``only added background,'' then as far as possible we should be able to identify which earlier materials have been preserved and what has been added.

Suppose, for example, that after a date the other person did not stay in contact. Years later, you learn that they had been going through a family crisis at the time. The new context may completely change the explanation of why they did not reply. But if there is a traceable sequence---you sent a message, there was no reply for several days, and contact later stopped---the new explanation does not automatically erase that occurrence sequence.

We therefore map mathematical persistence to the \textbf{traceability of an evidence ledger}, rather than to the claim that ``memory is permanently preserved.''

\begin{correspondence}[Finite-Witness Persistence and Traceable Material]
Suppose a psychological return hypothesis was supported by an explicit behavioral sequence, contemporaneous records, verifiable events, or a predeclared episode witness. Later expansion of context may change its interpretation, comparison scale, or continuation meaning. But if the original witness is to be withdrawn, the present model requires a separate account of which earlier material was found to be erroneous, omitted, missegmented, or incomparable.
\end{correspondence}

\textbf{Structural note.} This is a research-accounting principle, not a natural law of autobiographical memory. If the witness depends purely on later subjective recollection, then the witness itself may change through reconstruction, and mathematical persistence cannot be applied directly.

\section{Discovering It Later Does Not Mean the Pattern Started Later}

Mathematical filtration offers a distinction that is especially useful for psychology to borrow---and especially easy to misuse: \textbf{the level at which a structure is first seen is not the same as the time at which that structure originated in reality}.

Part IV will separate two quantities inside a fixed presentation: the support level $\ell(R)$ of a concrete witnessed datum, and the earliest witness level $b_{\mathrm{word}}(u)$ of a source return word. Part III takes only the research lesson they share. If a person first says on the third similar interaction, ``So this really does repeat,'' that does not mean the pattern first came into existence on the third occasion. Likewise, if a researcher first detects a return-completion only after adding a new variable, that does not mean the return-completion appeared only after the experimenter detected it.

The psychological side should therefore distinguish:
\begin{itemize}
\item \textbf{occurrence time:} when the relevant events occurred;
\item \textbf{recognition time:} when the subject first recognized them as a particular pattern;
\item \textbf{model-detection level:} the scale at which the current research protocol first acquires sufficient witness support.
\end{itemize}

The three may be far apart.

This distinction matters especially in psychology because many experiences of ``sudden insight'' create a strong illusion: either the pattern seems not to have existed until it was named, or, once named, it suddenly appears to have been everywhere since childhood. Both impressions require independent evidence.

\section{Filtration Corresponds Best Not to ``More and More Memory,'' but to the Discovery Sequence of a Research Procedure}

Mathematically, as more levels are permitted, more return witnesses may appear; correspondingly, the remainder that has not met any return shrinks. This direction is a monotonicity property of the formal system.

Psychological reality is different. A pattern recognized today may be withdrawn tomorrow; a stretch of the past may be resegmented; an old label may become invalid because of new evidence. It is therefore illegitimate to translate
\[
\mathcal R_{\le 1}\subseteq \mathcal R_{\le 2}\subseteq\cdots
\]
directly as ``a person's structural memory becomes progressively more complete.''

A better psychological reading places filtration in the research procedure rather than in an ``internal memory store.'' Predeclare a sequence of increasingly rich observables, context windows, or segmentation scales, and record at each stage which return hypotheses gain witness support. New scales can then add candidate structures. If an old candidate is withdrawn, the record must also state whether the withdrawal occurred at the experiential level, the coding level, or the model level.

In other words:
\begin{plainbox}
\begin{center}
\emph{Filtration models cumulative access conditions,\\
not cumulative human memory.}
\end{center}
\end{plainbox}

This is what makes mathematical monotonicity genuinely useful: it governs how we keep research records, rather than inventing for the brain a monotone memory mechanism that does not exist.

\section{``I Used to Understand It Wrongly'' Has at Least Three Entirely Different Meanings}

When a person withdraws an old narrative, Third Order should not immediately say that ``the old return was false.'' It should first ask what, exactly, has been withdrawn.

\begin{enumerate}
\item \textbf{The interpretation is withdrawn.} The event sequence and episode witness remain largely intact, but ``this was disrespect'' becomes ``there was insufficient information.''
\item \textbf{The segmentation is withdrawn.} Later material shows that what had been treated as one episode actually contained an important re-rooting, so the old return-completion test itself must be redone.
\item \textbf{The material is withdrawn.} New evidence shows that a remembered detail, sequence, or contemporaneous judgment is unreliable; the original witness packet has changed.
\end{enumerate}

The theoretical consequences of these three cases are completely different.

The first is closest to ``the same structure receives a new interpretation.'' The second has already changed the segmentation. The third changes even the data we regard as available. Only the first admits the cleanest analogy with mathematical inflation/persistence. The second and third force the mapping to weaken, or even to exit.

\section{Narrative Compression: Why Stories Are Necessary, and Why They Are Dangerous}

Without compression, a life is almost impossible to narrate. Years of experience cannot be preserved second by second in autobiography; enormous numbers of micro-events have to be compressed into macro units such as ``that relationship,'' ``my PhD years,'' or ``those years when I was always trying to prove myself.'' The power of narrative identity comes precisely from this capacity to make enormous temporal scales available for use \citep{mcadams2013,bruner1987}.

Mathematical telescoping provides a beautiful formal analogy here, but one that must be used with restraint: it skips intermediate levels and makes the presentation coarser, without implying that the intermediate structure never existed.

We therefore propose a Level II candidate:
\[
\text{structurally faithful narrative compression}
\strsim
\text{telescoped presentation}.
\]

Psychologically, ``faithful'' cannot be judged by how compelling the story sounds. The relevant question is whether the differences compressed away would change our judgments about future continuation, responsibility, behavioral choice, or counterfactuals.

If not, the compression may simply be omission. If they would, yet the narrative still glues them into a single ``this is just who I am,'' the operation looks more like structural overreach.

A sentence such as ``I always end up with people who do not respect me'' may telescope ten very different interactions into a single life theme. The theme may capture a genuine regularity. It may instead glue failures of reciprocity, scheduling constraints, communication habits, value conflicts, and one accidental silence into one identity story. Third Order does not solve the problem by opposing stories. It asks: \textbf{which relations that distinguish possible futures did this compression preserve?}

\section{Untelescoping: When a Large Word Becomes Too Heavy, Put the Intermediate Levels Back In}

Third Order sometimes needs no deeper explanation. It needs the explanation to be taken apart.

\begin{quote}
``I am a failure.''
\end{quote}

The sentence is heavy because it jumps in one step from a small number of episodes to an entire identity. At that point, a very ordinary operation becomes useful: put the omitted intermediate levels back in.

\begin{itemize}
\item Which specific setbacks were included in this judgment?
\item Were the task demands the same each time?
\item What resources, attentional conditions, and social conditions were present then?
\item Which contrary episodes did not enter the current retrieval?
\item Between ``I did poorly this time'' and ``I am a failure,'' which unspoken inferences were made?
\end{itemize}

This is what the book actually means when it borrows \emph{untelescoping}: restoring the intermediate relations flattened by a macroscopic narrative.

Here it meets the ``dimensional switching'' of the \emph{Geometry of Existence}. When one coordinate loses resolution, the first move is not to manufacture a still grander meaning, but to return the obscured relations to the scene \citep{geometry2026}.

\section{Two Examples: The Same ``Reinterpretation'' Can Be Structurally Completely Different}

\textbf{Example one: an unfamiliar talk and ``my range of knowledge is too narrow.''} Imagine that you attended a talk but could not formulate a question. Years later, you might still conclude that you genuinely lacked the relevant background knowledge. Third Order does not require this conclusion to be softened. It requires the scale to be opened: what prerequisites did that talk require? What was your attentional condition that day? How did you perform when asking questions in your own area? Did you later acquire the relevant background?

If the larger context still supports ``my range of knowledge needs to expand,'' then the judgment survives untelescoping. If it supports only ``I had no coordinates in that subfield,'' then ``my range of knowledge is too narrow'' is a scale jump from local accessibility to global identity.

\textbf{Example two: was a relationship really marked by ``long-term value incompatibility''?} ``We had long-term value incompatibilities'' can be a legitimate macro narrative if episodes across different times consistently leave similar future-relevant continuation differences in responsibility, communication, future planning, and behavioral choice. But if the entire judgment is retrospectively colored mainly by the last unanswered message, causing a large body of contrary earlier material to disappear from memory, then what occurred is not only telescoping; retrieval bias and narrative reconstruction may also be involved.

In both cases one can say, ``Only later did I finally see it clearly.'' The first looks more like a stable compression formed after context was expanded. The second reminds us that an ending can also reach backward and rewrite how we remember the beginning.

\section{If This Interface Holds, How Should Psychology Test It?}

If the scale/filtration mapping is genuinely to return something useful to psychology, it should at least motivate research questions of the following kinds:
\begin{itemize}
\item After different contexts are added to the same autobiographical material, which event judgments remain stable, and which narrative meanings change?
\item In participants' later retellings, which occurrence backbones remain traceable across time, and which details and episode boundaries move?
\item When a researcher untelescopes a global self-judgment, does restoring intermediate episodes change subsequent prediction, causal attribution, or behavioral choice?
\item Can recognition time, event occurrence time, and model-detection level be separated empirically?
\item If a putative ``persistent witness'' cannot be supported by a contemporaneous record, behavioral trace, or independent report and depends entirely on later memory, should the persistence mapping actively exit?
\end{itemize}

These questions are more discriminating than simply asking whether memory is reliable. They ask which kind of change occurs in the data, which in the segmentation, and which in the interpretation.

\section{Concept Passport}

\begingroup\footnotesize
\begin{longtable}{>{\raggedright\arraybackslash\bfseries}p{0.25\textwidth}>{\raggedright\arraybackslash}p{0.66\textwidth}}
\toprule
Psychological motivation & When a person later ``reinterprets the past,'' how can we distinguish added context, changed interpretation, redrawn event boundaries, and memory reconstruction?\\
Psychological concept & Autobiographical recollection; narrative revision; event resegmentation; context-dependent memory retrieval.\\
Mathematical concept & Occurrence-supported witness embedding; finite-witness persistence; return filtration; telescoping / untelescoping.\\
Preserved relation & Keep larger context and data rewriting on separate ledgers; preserve traceability of a finite witness; a coarse presentation should state how it comes from finer occurrences.\\
Discarded content & Human memory mechanisms; forgetting; reconsolidation; false-memory mechanisms; specific neural realization.\\
Translation level & Level II.\\
Candidate observables & Longitudinal retelling; stability of event boundaries; context manipulation; memory retrieval; causal attribution; subsequent prediction; contemporaneous records.\\
Non-example & Writing ``I now have a new interpretation'' directly as ``this is what the past really was''; treating filtration as monotonic accumulation inside human memory.\\
Failure conditions & Remembered content itself is highly unstable and lacks independent traces; narratives at different scales share no common occurrence backbone; differences removed by telescoping systematically change future-relevant judgments.\\
Reverse-translation check & Do not use persistence to deny memory reconstruction; do not use telescoping to replace theories of autobiographical memory or narrative identity.\\
Translation remainder & The actual mechanisms of memory reconstruction; forgetting and reconsolidation; false memories; the ways socially shared narratives alter an individual's past.\\
\bottomrule
\end{longtable}
\endgroup

\section{Stop Line for This Chapter: The Past Is Not a Stone, and the Model Must Not Pretend It Is}

The human past is not a stone buried underground, waiting only for a better theory to dig it up. We preserve traces, and we forget. We call the past back, and in calling it back we also reorganize it. New relationships, new language, and new self-positions all change which materials come back most readily and how they are understood.

Mathematical persistence may therefore enter only in a highly controlled sense. When a finite body of evidence has been explicitly declared to exist and the operation is only an expansion of context, the model requires that evidence to remain traceable. Once remembered content, an episode boundary, or the evidential status itself changes, the ledger must change as well; ``persistence'' can no longer be invoked to underwrite it.

\section{Chapter Conclusion: A Larger Context Should Not Automatically Become a Larger Story}

Knowing more later is a real and important change. But ``knowing more'' does not always mean ``finally finding the uniquely correct past.'' Sometimes we have added context. Sometimes we have changed the segmentation. Sometimes we have changed the story. Sometimes memory itself has changed.

Third Order does something modest here on purpose: every time we say ``I understand it now,'' it asks us to specify \emph{which level actually changed}.

\begin{quote}
\textbf{Adding context is not the same as erasing old material;}\\
\textbf{changing interpretation is not the same as changing an occurrence;}\\
\textbf{and genuine memory reconstruction cannot be disguised as mere refinement.}
\end{quote}

Only when these three ledgers can be kept separate do we earn the right to speak of filtration, telescoping, and so-called ``structural memory.''

The next chapter moves to a completely different kind of ``seeing more finely.'' Instead of enlarging the temporal background of the past, it asks whether finer continuations may still differ even after a return has already been established. Only then do the graph of germs and residue acquire their genuine psychological motivation.

\chapter{The Same ``Again,'' Different Afters: What Remains after Return-Completion}
\label{chap:resolution-translation-en}

\par\medskip
\noindent\textbf{Why this chapter matters}\par
Psychological change is easily trapped between two extreme narratives. One says, ``If this pattern has appeared again, then I have not changed at all.'' The other says, ``If I handled it differently this time, then it is no longer the same pattern.'' Reality is usually finer-grained than either sentence. An old pattern can be triggered again while the actions, predictions, relational positions, and self-interpretations that follow it have changed.

Here we preserve precisely the possibility that both can be true at once: \textbf{return-completion can be the same while continuation can differ}. The mathematical graph of germs and return residue are not names for ``psychological residue.'' They only supply a second coordinate, so that acknowledging an ``again'' does not force everything that happens afterward to be compressed into sameness.
\par\medskip

\begin{methodbox}[Interface card: why the same return can still have a different afterward]
\textbf{Psychological motivation:} The reappearance of an old pattern does not automatically mean that learning, action, and future choice have all returned to their starting point. Psychology needs to record separately whether one has entered the same organization again and how the process continues after entering it.

\textbf{Psychological native domain:} pattern recurrence; context-sensitive responding; subsequent prediction; action choice; self-reference; relationship updating.

\textbf{Mathematical native domain:} a finite witness $W_R$ records downstairs return-completion. When it is paired with a source return word $u$ to form the witnessed datum $R=(u;W_R)$, and a singular sector $\xi$ together with a germ-decorated lift / connector context is further fixed, the corresponding displacement is denoted $\rho_\xi(R)$ and records finer continuation. The bare return shape is generally insufficient to determine residue.

\textbf{Preserved relation:} return-completion and future-relevant continuation are two coordinates that may vary together or independently.

\textbf{Level:} Level II.

\textbf{Not claimed:} $\rho_\xi(R)$ is an emotional residue, a memory engram, a neural trace, or some directly measurable psychological entity; nor is $\rho_\xi$ miswritten as a function of the bare birooted return type alone.
\end{methodbox}

\section{Psychological Entry Point: An Old Pattern Coming Back Does Not Mean Everything Has Returned to the Starting Point}

Imagine two very similar events in a relationship.

The first time, you initiate contact and the other person does not respond for a long time. You check repeatedly, begin to wonder whether you are being slighted, rehearse a future corrective conversation in your head, and eventually withdraw from contact.

Several months later, the other person remains silent in a similar situation. Your body still tightens a little, and the old thought---``am I being ignored again?''---comes back quickly. In other words, a return may genuinely have occurred. This time, however, you do not keep sending messages and do not rehearse how to educate the other person. You still dislike the uncertainty, but you treat it as new information about the relationship and continue with your life.

If we have only one coordinate, we are forced into a choice between two statements:
\begin{quote}
``If the old feeling has appeared again, I have not changed at all.''

or:

``If I acted differently this time, then this is no longer the same pattern.''
\end{quote}

Part III rejects this forced choice. A more precise description may be:
\begin{plainbox}
\begin{center}
\emph{The same subject-relevant return-completion has appeared again;\\
but the continuations available after it have changed.}
\end{center}
\end{plainbox}

This is the psychological reality that all of the mathematical structure in this chapter is meant to serve.

\section{Two Questions Must Be Separated: Did It Come Back, and How Does It Continue after Coming Back?}

Chapter 17 used return to address the first question:
\begin{quote}
``Did this episode come back to the same relevant position?''
\end{quote}

This chapter adds a second:
\begin{quote}
``Even if it came back, how can things continue from this position?''
\end{quote}

These two questions are easily glued together in self-narrative. A person may say, ``I became anxious again, so I am still the same person I used to be.'' Or: ``I have learned not to reply anymore, so the old pattern has disappeared completely.'' The first expands return-completion into an unchanged identity; the second may deny a recurrence that genuinely occurred in order to prove growth.

Third Order prefers to keep two ledgers:

\begin{center}
\begin{tabularx}{0.94\textwidth}{>{\bfseries\raggedright\arraybackslash}p{0.31\textwidth}Y}
\toprule
Question & What psychology needs to record \\
\midrule
Was a return formed? & How the episode was segmented; which endpoint was recognized as the same relevant position; what historical material supports the ``again'' judgment. \\
How does the process continue after the return? & Whether the next action, prediction, allocation of attention, memory retrieval, relational choice, self-evaluation, and related variables still follow the old path. \\
\bottomrule
\end{tabularx}
\end{center}

The first row concerns return-completion; the second concerns continuation. Only by separating them can we describe an ordinary form of change often hidden by language: \textbf{the old pattern was triggered, but it no longer possesses the same future}.

\section{What Counts as a Difference That Affects Continuation: Not Every Difference Deserves Promotion}

Restraint is needed here as well. Every recurrence differs in countless details: the weather is different, the clothes are different, the exact wording differs, and the heart-rate curve cannot be exactly identical. If every microscopic difference is called residue, the model quickly loses resolution.

Here we care only about a narrower kind of difference: \textbf{does it change which paths remain feasible, predictable, or more likely to occur next?}

For example, after the same broad ``initiate contact--wait--no response'' episode:
\begin{itemize}
\item one occurrence may make continued checking more likely, while another allows checking to stop quickly;
\item one may retrieve an autobiographical cluster organized around ``I am always rejected,'' while another mainly retrieves the judgment ``there is not enough information yet'';
\item one may push future relationship choices toward defensiveness and withdrawal, while another merely adjusts expectations about the present person;
\item one may give organizing authority to ``I need to prove that I do not care,'' while another generates no new self-evaluative episode.
\end{itemize}

These differences matter not because they are somehow ``deeper,'' but because they alter continuation.

The psychological side can therefore begin with an entirely ordinary working question:
\begin{quote}
\textbf{Once we have already acknowledged ``this is the same again,'' which remaining differences systematically change the next step?}
\end{quote}

Only at this point does the graph of germs have reason to enter.

\section{Why Mathematics Enters Here: The Same Downstairs Return-Completion Can Have Different Resolved Continuations}

Part II has already established the mathematical native domain. Part III borrows only one relational skeleton.

Let $R=(u;W_R)$ be a witnessed return datum, with finite witness
\[
W_R=(Q;b,\gamma,b)
\]
satisfying, in the coarser tile/orbital picture,
\[
b\xrightarrow{\gamma}b.
\]
In other words, the downstairs endpoint has completed a return.

But after lifting the same path to the graph of germs, one may have
\[
(\mathrm{Id},\widetilde b)
\xrightarrow{\widetilde\gamma}
(h,\widetilde b'),
\qquad h\neq e.
\]
From the coarser coordinate it has indeed ``come back.'' On the finer germ sheet, however, it has not returned to exactly the same continuation position.

Once the concrete lift and connector / germ context for this return occurrence have been fixed, this book abbreviates the two pieces of data as
\[
(R,\rho_\xi(R)).
\]
Here the finite witness $W_R$ records downstairs return-completion, while $\rho_\xi(R)$ records the finer germ displacement of that decorated witnessed return datum. Strictly speaking, $\rho_\xi(R)$ is a context-dependent shorthand: knowing only the bare birooted shape or abstract return type is generally insufficient to recover the residue.

Psychology readers do not need to interpret this formula as saying that a hidden group element exists in the brain. The one sentence worth carrying across is:
\begin{plainbox}
\begin{center}
\emph{Returning to the same coarse position does not imply that all finer differences relevant to the future have been erased.}
\end{center}
\end{plainbox}

The graph of germs is not introduced here to deny the downstairs return. It prevents a coarse return-completion from declaring all continuations identical in advance.

\section{A Mathematical Image for Intuition: Returning to the Same Place while Standing on a Different Sheet}

Analytic continuation offers a beautiful image here, but it must remain tightly limited.

Take the complex logarithm and go once around the origin in $\mathbb C^\times$. In the base space the path returns to its original point. Lift it to the logarithm's Riemann surface, however, and the starting point $w$ moves to
\[
w+2\pi i.
\]
Thus both are true:
\[
\text{return-completion on the base}
\qquad + \qquad
\text{displacement on the sheet}.
\]

The image is compelling not because psychological life is ``really like complex analysis,'' but because it lets us imagine a cycle that genuinely completes without returning the system to an entirely identical continuation.

The stop line must follow immediately. A Riemann surface arises from the multivaluedness of an analytic function; a graph of germs arises from a local group action; psychological experience is neither. This chapter borrows only the relational intuition of \textbf{coarse-level return-completion together with fine-level displacement}.

\section{Return Residue: Do Not Read It as ``How Much Emotion Is Left''}

We can now state the central structural correspondence of the chapter.

\begin{correspondence}[Return Residue and the Continuation Distinction]
For a germ-decorated return occurrence in which a concrete lift / connector context has already been fixed,
\[
\begin{gathered}
\begin{gathered}
\text{differences that, after an episode has been recognized as the same kind of}\\[-1mm]
\text{return-completion, still systematically alter future continuation}
\end{gathered}
\\[1mm]
\strsim
\\[1mm]
\rho_\xi(R)\in H_\xi.
\end{gathered}
\]
Here $\rho_\xi(R)$ must not be understood as a property of the bare episode shape.
\end{correspondence}

The most important danger here comes from the ordinary-language sense of ``residue.'' A person may still feel very distressed after an event, and that may of course matter. But the mere persistence of distress does not tell us how continuation has changed.

Sometimes emotion remains intense while the next steps remain exactly the same: the person still runs the same checking cycle, avoidance pattern, and explanatory route. At other times subjective discomfort is already low, yet a new relationship judgment permanently changes what the person will approach, refuse, or predict in the future. The latter is closer to the continuation distinction at issue here.

Therefore,
\[
\boxed{\text{emotional persistence}\neq\text{return residue}.}
\]

In this book, residue is only a structural ledger: once return-completion has been established, which differences remain useful for the future?

\begin{nonexample}[``It still hurts, therefore the residue is nontrivial'']
Subjective suffering, bodily activation, or vividness of memory cannot by themselves establish the continuation distinction corresponding to $\rho_\xi(R)$. They may serve as candidate observables, but they acquire the structural role needed in this chapter only when they systematically change subsequent action, prediction, retrieval, or relational trajectory.
\end{nonexample}

\section{Why Some Repetitions Look More Like a Spiral than a Circle}

Everyday language often says, ``I have gone around and come back again.'' Yet some repetitions are more like passing through a similar direction again while climbing a spiral staircase. In a planar projection one appears to return to the same angle; in a finer vertical coordinate the position is different.

This is the most valuable role of the ``spiral'' metaphor in this book. It lets us reject two extremes at once:
\begin{itemize}
\item ``It appeared again, so all progress has been erased.''
\item ``I changed a little, so the old pattern no longer exists.''
\end{itemize}

The mathematical spiral prototype, however, is much stricter than this psychological metaphor. Only when a germ-decorated witnessed return datum $R$ is fixed, its repeated composition is meaningful inside the same group of germs, and its residue continues to be recorded by the same
\[
\rho_\xi(R)=h,
\]
can we rigorously distinguish:
\begin{enumerate}
\item $h=e$: the resolved continuation also completes;
\item $h^k=e$ for some $k>1$: after several iterations, the finer position completes;
\item $h^k\neq e$ for all $k>0$: the downstairs return keeps appearing while the upstairs continuation does not complete.
\end{enumerate}

Only the third case supplies a genuine mathematical prototype for a spiral. The psychological side currently borrows it only to pose a question. Everyday expressions such as ``repetition while growing'' or ``trauma spiral'' cannot be identified directly with $h^k$.

\section{The Same Appearance Can Still Carry a Different Continuation}

Even when two episodes have the same shape at a coarse scale, context may change their future meaning.

Both can be summarized as ``an argument.'' One occurs in a long-term secure relationship in which repair is available; the other occurs in a relationship with pronounced power asymmetry and high exit costs. The surface episode shape may be similar while the predictions, action spaces, and responsibility structures that follow are entirely different.

The same is true in research. Two unfamiliar talks may both contain
\[
\text{cannot follow}
\longrightarrow
\text{brief blankness}
\longrightarrow
\text{stop asking questions}.
\]
After the first, the researcher returns to ``I can only work inside my own small area.'' After the second, despite again asking no question, the researcher records one genuinely new connector and later begins reading another field. The coarse return-completion may be similar; future continuation is not.

Mathematically, a bare birooted type is generally insufficient to determine residue; the connector / germ context must be retained. The psychological side does not borrow that theorem itself, but a research discipline:
\begin{quote}
\textbf{If a difference changes the future, do not quotient it out merely because the episode looks similar.}
\end{quote}

\section{Higher Resolution Does Not Revoke a Return: Growth Does Not Require Pretending ``It Never Repeated''}

This is the most important discipline in the chapter for psychology readers.

If a return already has an explicit witness, later increasing the resolution and discovering that action, prediction, or relational position differs does not require us to say retroactively:
\begin{quote}
``So there was never an old pattern after all.''
\end{quote}

A more precise description is:
\begin{plainbox}
\begin{center}
\emph{past return-completion retained;\\
future-relevant continuation refined.}
\end{center}
\end{plainbox}

Psychological change does not require history to be whitewashed. A person can say honestly:
\begin{quote}
``Yes, the same kind of silence triggered me again; but this time I did not continue writing it into the same story.''
\end{quote}

This has more structural resolution than ``I am completely beyond being triggered.'' It allows recurrence and change to coexist, and it prevents growth itself from becoming a new identity performance.

\section{What Third Order Adds Is Not More Detail, but the Question: What Did That ``Again'' Just Compress Away?}

Second-order subjectivity is highly skilled at using macro-symbols: ``rejected again,'' ``failed again,'' ``another person like this.'' These compressions are extremely efficient, but they can fold together differences that still matter for continuation.

Third-order resolution does not require infinite analysis of every detail. It continues with only one question:
\begin{quote}
\textbf{Which differences that would have led to different futures were just compressed into sameness by this ``again''?}
\end{quote}

If the answer is ``none---the added detail changes no prediction, action, or relational choice,'' then increasing resolution further is merely semantic inflation. If the answer is ``yes---the return-completion is the same, but one variable now systematically changes the next step,'' then the coarse return and finer continuation are worth recording separately.

The psychological value of the graph-of-germs mapping is therefore not that ``finer is more advanced.'' It supplies a resolution principle with an exit condition: \textbf{only differences that affect continuation deserve to remain visible}.

\section{If This Interface Holds, What Else Should Psychology Measure?}

If residue is more than an elegant metaphor, it has to return to observable differences. Once participants judge two episodes to be ``the same pattern / again,'' future experiments can continue by comparing:
\begin{itemize}
\item whether the next action choice differs;
\item whether prediction of the future outcome differs;
\item whether autobiographical retrieval draws on different material;
\item whether approach / avoidance, checking, or repair behavior changes;
\item whether the self-reference frame shifts;
\item whether the same event shape produces a different continuation in different relational contexts;
\item whether adding a candidate variable genuinely improves prediction of subsequent behavior, rather than merely adding narrative complexity.
\end{itemize}

If a proposed residue has no additional discriminating power for any future variable, it has not completed the psychological translation and should exit rather than survive as a mysterious ``deep trace.''

\section{Concept Passport}

\begingroup\scriptsize
\begin{longtable}{>{\raggedright\arraybackslash\bfseries}p{0.25\textwidth}>{\raggedright\arraybackslash}p{0.66\textwidth}}
\toprule
Psychological motivation & Why ``the same pattern appeared again'' and ``how this occurrence continues'' must be kept on separate ledgers; how recurrence and change can be described as simultaneously true.\\
Psychological concept & Pattern recurrence; context-sensitive continuation; subsequent prediction; action choice; self-reference / relationship updating.\\
Mathematical concept & Graph of germs; covering sheets; finite return witness; return residue / germ displacement.\\
Preserved relation & Coarse return-completion and finer continuation can remain distinct; similar episode shape need not determine the same future.\\
Discarded content & Analytic-function mechanism; literal realization of groups of germs; the everyday meaning of emotional ``residue''; memory engrams; neural mechanisms.\\
Translation level & Level II.\\
Candidate observables & Subsequent action; subsequent prediction; retrieval pattern; approach / avoidance; checking / repair behavior; self-reference frame; relational context.\\
Non-example & Treating ``it still hurts'' as automatically implying nontrivial residue; declaring continuation different merely because details differ; directly rewriting psychological ``spiral growth'' as $h^k$.\\
Failure conditions & Adding the proposed residue variable changes no subsequent prediction; continuation difference is entirely free post-hoc naming; same-return judgments have no reproducible relation to finer variables.\\
Reverse-translation check & Do not use residue to replace established psychological constructs of learning, context effects, memory updating, or behavioral change. It adds only the structural question of keeping return-completion and future continuation on separate ledgers.\\
Translation remainder & The actual neural, bodily, and social realization of continuation differences; which psychological variables are stable future-relevant coordinates.\\
\bottomrule
\end{longtable}
\endgroup

\section{Structural Conclusion: Return Is a Loop, Not a Foreclosure of Fate}

``Again'' is powerful because it suddenly gives the present a history. For that very reason, it can also lead the subject to assume that once an old pattern appears again, the past has already determined the future.

This chapter preserves a more open geometry. An episode can come back to the same relevant coordinate. That return-completion can be real, traceable, and worth acknowledging. At the same time, finer continuation can still move.

\begin{plainbox}
\begin{center}
\emph{The same ``again'' need not lead to the same afterward;\\
return says that the path has formed a loop,\\
it has not written the ending of the future.}
\end{center}
\end{plainbox}

This is also Third Order's most important loosening of repetition: it does not prove growth by denying an old pattern, and it does not let an old pattern steal the future merely because it has appeared again. The next chapter will use a genuine mathematical theorem as a pressure test: under what strict conditions can alternative contexts internal to the system force a previously open traverse to acquire a return witness?

\chapter{No Perspective Has Final Authority: Why ``Looking from Another Angle'' Must Also Be Constrained by Evidence}
\label{chap:constrained-perspective-en}

\par\medskip
\noindent\textbf{Why this chapter matters}\par
In psychology, ``look at it from another angle'' is a very common suggestion. A silence can be understood as ``disrespect,'' or as ``the other person has not replied yet''; failing to understand a talk can be understood as ``I am not capable enough,'' or as ``I lack background in this area.'' Sometimes a change of frame genuinely reveals something that had been hidden. Sometimes it merely replaces a painful story with a more pleasant one.

We therefore ask only one question: if the current interpretation has no final authority, what entitles a new interpretation to enter? We will borrow a genuine mathematical theorem to constrain the question. Psychology readers do not first need to understand a purely non-Hausdorff singularity or keep track of a group of germs. The most important point of the chapter can be stated in ordinary language: another perspective cannot appear out of nowhere. It must come from a predeclared, rule-governed mode of organization, and after the perspective changes it must return to the same material and show exactly what additional structure it explains.
\par\medskip

\begin{methodbox}[Interface card: constrained ``change of perspective'']
\textbf{Psychological motivation:} The same event can support different interpretations. We need to know when another frame genuinely reveals more structure and when it merely substitutes one story for another.

\textbf{Mathematical native domain:} The two endpoints of the same finite process may be different on a ``fine map'' yet be identified as the same position on another, coarser map permitted by the system itself. A strict theorem supplies a special instance of this phenomenon.

\textbf{Part III retains only:} a new frame must have a source, change an explicit relation, and return to the same material for testing.

\textbf{Level:} strict in the mathematical native domain; cross-domain use remains Level II.
\end{methodbox}

\section{Set Mathematics Aside First: This Chapter Really Asks Only One Question}

Imagine that you sent someone a message and received no reply for two days. A thought might quickly come to mind:
\begin{quote}
``She does not respect me.''
\end{quote}
A friend might say:
\begin{quote}
``Maybe she is just busy.''
\end{quote}
The second sentence is gentler, but gentleness does not automatically make it closer to the facts. It may be true, or it may simply be an interpretation temporarily added to make you feel better.

Likewise, imagine a researcher who cannot follow an unfamiliar talk and thinks:
\begin{quote}
``My knowledge is too narrow.''
\end{quote}
The researcher might then replace it with:
\begin{quote}
``I simply lack background in this field.''
\end{quote}
That sounds more flexible, but the same questions remain. Does understanding actually improve once the relevant background is supplied? Can the researcher then formulate questions? If nothing changes, the second sentence may simply be a kinder way of saying the same thing.

Third Order therefore cannot stop at:
\begin{quote}
``The current interpretation may be wrong, so think of another one.''
\end{quote}
If any new interpretation is admissible, the theory loses its constraints. The real question is:
\begin{quote}
\textbf{Once the current interpretation loses final authority, what constrains the alternative?}
\end{quote}
Mathematics has only this methodological job here.

\section{Look at a Map First: Why the Same Route Sometimes Counts as Coming Back and Sometimes Does Not}

Let us first take an example unrelated to psychology. Imagine that you entered a building through its south door and eventually left through its north door. If the map resolves individual doors, the south and north doors are obviously different positions:
\begin{quote}
You started at the south door and ended at the north door. You did not come back to the same door.
\end{quote}
But if the current question is only whether you left the building, the map can legitimately be coarser. Both doors belong to the same building, so at the scale of the building one can say:
\begin{quote}
You started from this building and ended at this building. At this resolution, a return has formed.
\end{quote}
Both statements can be true because they answer questions at different resolutions.

\begin{plainbox}[The key point]
This does not mean that a coarse view may arbitrarily call different things the same. If we simply call the whole city ``one place,'' then every route can be said to ``come back,'' and the claim has no informational value. A legitimate coarsening must say why, for the present question, this particular difference may be forgotten while other differences must still be retained.
\end{plainbox}

This is the most basic mathematical intuition of the chapter. A system may permit more than one resolution. At a finer resolution two endpoints are different; at another, rule-governed coarser resolution they may be identified. For the psychological meaning of the chapter, understanding this much already gets us halfway there.

\section{What Does This Have to Do with Psychology? A Difference That Looks ``Crucial'' Now Need Not Remain Crucial in Every Frame}

Whenever people interpret experience, they are also deciding which differences deserve to be preserved. Imagine that the other person has not replied for two days. If the current question is ``does this relationship respect me?'', promptness may be decisive and may directly alter what happens next. If the current question is instead ``do I have enough information to infer the other person's relational intention?'', the very same delay may not yet justify a conclusion.

The unfamiliar-talk example has the same structure. Under ``a competent researcher should be able to formulate questions quickly,'' drawing a blank may be magnified into information about identity. Under ``do I currently possess the required background?'', it is first treated as a local accessibility problem.

No frame is automatically correct because it is kinder, harsher, or more ``psychological.'' Third Order initially acknowledges only this:
\begin{quote}
\textbf{A difference treated as crucial by the current frame need not be equally crucial in every evidence-based frame.}
\end{quote}
But a second sentence must immediately follow:
\begin{quote}
\textbf{The new frame cannot therefore be invented at will.}
\end{quote}
The work of Chapter 21 is to make that second sentence precise.

There is also a structural correspondence here that matters to the whole book. The mathematical audit later in the chapter will show that a purely non-Hausdorff singularity is not merely a ``special point.'' At the translation layer of Part III, it provides an especially useful benchmark for a third-order observation site. But the claim must be stated correctly from the beginning: it is not the singular point $\xi$ alone that equals a third-order observation site. The formal prototype is $\xi$ together with the admissible nearby resolutions supplied by the system itself. A distinguished root by itself still corresponds only to a current ego-position. The third-order feature appears when the distinction structure generated at that position can itself be re-examined in another legitimate coordinate system.

\section{The Real Dialectic: Once a Singularity Becomes Visible, Can It Be Reconstituted as a New Final Object?}

After separating regular from singular boundary points, one more distinction is essential: singular does not automatically mean non-Hausdorff, and it certainly does not automatically mean purely non-Hausdorff. This matters for Third-Order Consciousness not because the mathematical vocabulary sounds ``deeper,'' but because it distinguishes two different reflexive structures: the observer enters the picture, and the observer that has entered the picture may or may not be able to reconstitute itself as the new master of the picture.

\begin{sourcebox}[Two independent native domains are being connected here, not synonyms]
In the mathematical native domain, Hausdorff, non-Hausdorff, and purely non-Hausdorff singularities are strict topological conditions governing how interiors of fixed-point sets accumulate near a singular point \citep{nek2018}. Snowfall on the Clock uses ``singularity'' in the language of the Geometry of Existence: the coding, boundaries, and feedback of an existing structure can no longer process a disturbance in the same way; the subject is a central singularity forced by living, linguistic, and social organization; the task of Third-Order Consciousness is not to delete this center but to open a passage at the singularity \citep{snowfall2026e,snowfall2026f,snowfall2026g}. A later discussion of ``node awareness'' emphasizes that stable boundaries and open interfaces can coexist: a position can be real without possessing the whole \citep{snowfall2026h}.

This section proposes only a controlled structural correspondence. Mathematical separability / non-separability supplies a formal prototype for asking whether an observation position, once exposed, can again acquire final authority. It does not claim that these intellectual texts are already discussing groupoid topology.
\end{sourcebox}

The four mathematical cases can first be placed side by side:

\begin{center}
\small
\begin{tabularx}{0.98\textwidth}{>{\bfseries\raggedright\arraybackslash}p{0.25\textwidth}p{0.15\textwidth}Y}
\toprule
Position & Germ fibre & Key local property \\
\midrule
Regular boundary & $H_b=1$ & Cross-tile gluing may be required, but there is no germ multiplicity. \\
Hausdorff singularity & $H_\xi\neq 1$ & A germ difference exists and remains locally separable from the identity on a sufficiently small neighborhood. \\
Non-Hausdorff singularity & $H_\xi\neq 1$ & At least some nontrivial germ differences are approached by arbitrarily nearby identity regions. \\
Purely non-Hausdorff & $H_\xi\neq 1$ & Every germ direction has such nearby trivialization. \\
\bottomrule
\end{tabularx}
\end{center}

This axis yields three increasingly fine structural sentences for the psychological translation, but it must never be turned into a hierarchy of persons.

\textbf{First: at the translation layer, singularity supplies a prototype in which the organizing position can no longer remain hidden from the record.} A generic singularity already means that the orbital endpoint does not exhaust the available information. At the same position one must also record germ information about ``how one arrived here.'' In the Level-II language of this book, the organizing operation that says ``this is me'' or ``this is another instance'' can no longer remain hidden; it too has to enter the record.

\textbf{Second: a Hausdorff singularity leaves the newly visible germ difference locally separable.} If the singularity is Hausdorff, a nontrivial germ difference remains separable from the identity in a sufficiently small neighborhood. At the translation layer, this leaves open the possibility that an exposed observation position can be reified as a new local object. Psychologically, this suggests a common danger: genuine reflexive insight may occur, yet ``my observer pattern,'' ``my controller,'' or ``my higher perspective'' is reified as a new stable managerial center. The claim is not that a Hausdorff singularity is a meta-ego. It is only that the source-side structure does not prevent such re-stabilization.

\textbf{Third: non-Hausdorffness removes clean local separability from the identity for at least one germ direction.} Suppose
\[
[g,\xi]\neq[1,\xi],
\]
yet arbitrarily near $\xi$ there are open regions on which $g$ agrees locally with the identity. Then the difference is real, while at the same time it cannot acquire a local territory completely separated from the identity. At the translation layer, this supplies a prototype in which a real distinction cannot be assigned a completely separate local territory.

This step is not claimed as an idea unique to Third-Order Consciousness. Traditions such as second-order cybernetics already articulate observer endogeneity. Non-Hausdorff geometry is used here to refine a shared problem: after an observation position enters the system, can it once again acquire local final authority? \citep{vonfoerster2003} The formal distinction is therefore:
\begin{center}
\resizebox{0.96\linewidth}{!}{$\displaystyle
\begin{gathered}
\text{source side: singularity retains nontrivial germ data beyond the orbital endpoint;}\\
\text{source side: non-Hausdorffness makes at least one such germ difference}\\
\text{approachable by arbitrarily nearby identity regions.}
\end{gathered}
$}
\end{center}

Purely non-Hausdorffness extends this property to every germ direction. For every attempt to stabilize the singular position as some final local difference, the system supplies endogenous regular contexts arbitrarily close by in which that difference locally trivializes. It therefore becomes one of the strongest source-side prototypes of an ``empty position'' used in this book: the position remains, and the difference remains real, but no germ direction gains permanent possession of the whole neighborhood.

\begin{principle}[Difference Exists, but Difference Cannot Possess the Whole Neighborhood]
The deepest relation borrowed here from purely non-Hausdorffness is not that ``everything ultimately becomes the same.'' Quite the opposite: nontrivial germs remain genuinely nontrivial at $\xi$. What is removed is their local final authority. The conditions
\[
[g,\xi]\neq[1,\xi]
\qquad\text{and}\qquad
\xi\in\overline{\operatorname{Fix}^{\circ}(g)}
\]
can hold at the same time. \emph{Difference survives; monopoly does not.}
\end{principle}

This also leaves a useful future intuition for non-Hausdorff but not purely non-Hausdorff singularities. A system may contain some organizing distinctions that can already be relativized by nearby admissible contexts, while other distinctions remain Hausdorff-stable. If psychology eventually produces corresponding evidence, third-order access may look more like sector-specific or relation-specific openness than a one-time jump from ``0'' to ``1.'' The present manuscript does not promote this intuition to an empirical claim; it is retained only as a future hypothesis.

\begin{warning}[Never Turn This into a Spiritual Hierarchy]
This book does not adopt a correspondence such as
\begin{center}
\resizebox{0.96\linewidth}{!}{$\displaystyle
\text{Hausdorff singular}=\text{Second Order},
\qquad
\text{purely non-Hausdorff}=\text{Third Order}.
$}
\end{center}
The Hausdorff, non-Hausdorff, and purely non-Hausdorff cases describe local separation structure among germ differences. First-, Second-, and Third-Order Consciousness describe modes of conscious organization. At most, the two axes are Level-II compatible with respect to the relation ``does this organizing difference still possess final authority?''
\end{warning}

\section{What the Mathematical Theorem Gives Us Is Not ``Another Answer,'' but Three Rules}

The preceding chapters have repeatedly said that the current root is not the whole system, that the present segmentation of experience is not the only possible segmentation, and that the current return grammar may omit structure. The danger is that once ``other modes of organization'' are allowed, the theory can become an explanatory machine that never fails.

We therefore use a special theorem of Kuang as a benchmark \citep{kuang2026a}. Its technical content will be examined later. Psychology borrows only three disciplines here.

\textbf{First: another perspective must have a source.} A new frame does not qualify merely because it makes sense after the fact. It should arise from variables that can be stated in advance, for example:
\begin{itemize}
\item whether the current position is that of ``the person being evaluated'' or ``the person gathering information'';
\item whether the scale is ten minutes, one interaction, or six months of relational history;
\item whether the current goal is to judge safety, judge reciprocity, or decide the next action;
\item which evidence threshold or appraisal rule is being used.
\end{itemize}
If a new ``deeper perspective'' can be invented whenever a counterexample appears, the model has no failure condition.

\textbf{Second: another perspective must change an explicit relation.} ``This way of thinking feels better'' is not structural evidence. A new frame should alter at least one traceable relation:
\begin{itemize}
\item which two moments count as the same episode;
\item whether this occurrence is judged to be ``again'';
\item which parts of the past are retrieved;
\item what is predicted next;
\item who is taken to need to act;
\item whether checking, approach, avoidance, or repair choice changes.
\end{itemize}
If no relation changes and only the wording changes, the operation may be mere renaming.

\textbf{Third: another perspective must return to the same material.} This is the most important rule. A new interpretation may not abandon the original scene and begin a grand story elsewhere. It must face again the same conversation records, the same behavioral order, the same episodes, and the same comparable repetitions. The most important ``formula'' in the chapter is therefore not a group-theoretic one:
\begin{center}
\resizebox{0.96\linewidth}{!}{$\displaystyle
\boxed{
\text{has a source}
\longrightarrow
\text{changes an explicit relation}
\longrightarrow
\text{returns to the same material}.}
$}
\end{center}
Third-order openness does not mean ``all perspectives are right.'' It means that the current perspective has no monopoly, and the alternative perspective has no exemption from evidence.

\section{Example One: No Reply---``She Does Not Respect Me'' and ``There Is Not Enough Information Yet''}

Return to the unanswered-message example. The current frame is:
\begin{quote}
``She does not respect me.''
\end{quote}
This frame may recruit several earlier similar interactions, place the current silence inside a return grammar of ``being slighted,'' and quickly generate withdrawal or corrective fantasy.

Now use another frame:
\begin{quote}
``At present there is not enough information to judge her relational intention.''
\end{quote}
The second frame is not automatically better because it sounds calmer. It has to be tested. Do the two frames make different predictions about the following?
\begin{itemize}
\item If the other person later explains the delay, which frame updates more readily?
\item If similar silences occur consistently across multiple interactions, does the second frame also begin to shift toward ``nonreciprocity''?
\item Do the two frames retrieve different material from the past?
\item Do they lead to different checking frequencies, probabilities of contacting again, or times of withdrawal?
\end{itemize}
If all of the answers are the same, ``changing perspective'' has added little structural information. If systematic differences appear, we have reason to say that the current root does not exhaust the organizational possibilities of the experience.

\section{Example Two: An Unfamiliar Talk---A Kinder Interpretation Must Also Be Able to Fail}

The current frame is:
\begin{quote}
``My knowledge is too narrow, so I cannot understand this.''
\end{quote}
The alternative frame is:
\begin{quote}
``What I currently lack is interface knowledge for this field.''
\end{quote}
Third Order does not adopt the second sentence merely because it is kinder. It must make predictions that can fail:
\begin{itemize}
\item Does understanding improve substantially after one page of terminology is supplied?
\item Does performance improve in another equally unfamiliar talk whose background is explained more clearly?
\item After a few key concepts are filled in, can questions actually be generated?
\item If none of these changes help, should the ``interface-knowledge deficit'' frame contract?
\end{itemize}
At the same time, ``my knowledge is too narrow'' does not become more realistic merely because it sounds harsher. It too must account for differences among talks, changes after background is supplied, and longitudinal reading data.

What Third Order does is remove both stories from the judge's seat. They are candidate organizations, and each must state its source, predictions, and failure conditions.

\section{Mathematical Audit I: For Purely Non-Hausdorffness, First Understand Only One Local Phenomenon}

Only now do we enter the technical language. At a special position $\xi$, a nontrivial germ can still be distinguished from the identity. Roughly, this means that some difference in continuation genuinely exists at $\xi$.

The special feature of purely non-Hausdorffness is that arbitrarily close to $\xi$ there are regular regions in which this distinction locally disappears: the corresponding action agrees there with the identity. It therefore gives not ``perspectives may arbitrarily change truth,'' but a highly constrained local structure:
\[
\begin{gathered}
\text{at the singular position, the distinction is preserved;}\\
\text{in some nearby regular contexts supplied by the system itself,}\\
\text{the distinction can locally disappear.}
\end{gathered}
\]
Which contexts are admissible and which germs can trivialize there are determined by the mathematical structure. The observer is not free to invent the identifications.

\subsection*{Why this supplies a formal prototype of a third-order observation site}

We can now connect the root of Chapter 15 to the present purely non-Hausdorff structure more precisely. If we retain only the singular root $\xi$, it is first merely a distinguished position: relations are currently organized from here, and some continuation differences are visible from here. This remains at the same level as the earlier Part-III relation
\[
\text{root}\strsim\text{ego-position}.
\]
A special root does not automatically become Third Order merely because it is special.

What purely non-Hausdorffness adds is a family of endogenous nearby regular contexts. A germ distinction that remains nontrivial at $\xi$ may locally trivialize in those admissible contexts. For the first time, we obtain a structure consisting of ``a position + the possibility of relativizing the resolution generated at that position.'' For Part-III bookkeeping we may temporarily write
\[
\mathcal O_\xi=(\xi;\mathcal C_\xi),
\]
where $\mathcal C_\xi$ denotes only the admissible nearby regular contexts / limit resolutions associated with $\xi$. This notation is not a new standard mathematical object; it is translation-layer bookkeeping. In the special setting of Theorem 5.1, the relevant contexts are further controlled by a finite family and generate the regular-limit graphs $\Lambda_i$ used below.

\begin{correspondence}[pNH Observation Site and the Third-Order Observation-Site Benchmark]
At Level II, the controlled relation is
\begin{center}
\resizebox{0.96\linewidth}{!}{$\displaystyle
\begin{array}{rcl}
\text{distinguished root at }\xi &\strsim& \text{current ego-position},\\[1mm]
\text{root }\xi+\text{admissible nearby resolutions}
&\strsim&
\text{formal prototype of a third-order observation site}.
\end{array}
$}
\end{center}
What is preserved is this: the current observation position remains operative, but it no longer monopolizes the rule determining which distinctions must be retained. The system contains constrained alternative resolutions in which distinctions generated by the present position can themselves be made visible differently.
\end{correspondence}

This is why the package is closer to Third Order than a bare root. Third Order does not replace the current position; it makes visible the way the current position itself generates differences.

Mathematical objects do not literally ``observe'' anything. ``Observation site'' here is structural translation. Nor do we claim that every third-order observation site must possess purely non-Hausdorff structure, or that an actual psychological system contains a literal singular point. The correspondence remains a Level-II benchmark.

\section{Mathematical Audit II: What Exactly Happens in ``One's Traverse Is Another's Return''?}

Fix a type-$h$ traverse. At the fine resolution of the graph of germs, it runs from the identity branch to the $h$-branch. The endpoints are therefore different, so it is an open traverse.

The third hypothesis of the later Theorem 5.1 can now be stated directly. It requires a pre-fixed finite partition $\mathcal P$ of $\Omega(B)\setminus\{\xi\}$ such that for every $h\in H$ there is a piece $P_i\in\mathcal P$ accumulating at the singular point $\xi$ on which the restriction of $h$ agrees with the identity:
\[
h|_{P_i}=\Id|_{P_i}.
\]
Thus $h$ may give a nontrivial germ at $\xi$, while along one preauthorized class of nearby regular contexts approaching $\xi$ the germ distinction locally disappears. The context is not an interpretation chosen ad hoc by an observer; it belongs to the finite family fixed in the theorem's hypotheses.

Only then does a quotient enter. Take regular points $\zeta_n\in P_i$ with $\zeta_n\to\xi$. Limits of the corresponding rooted orbital graphs give a regular-limit graph $\Lambda_i$. In that limit graph, germs that become trivial along $P_i$ are identified. If the current $h$ lies in the corresponding trivialized subgroup, the identity branch and the $h$-branch, distinct in the graph of germs, collapse in $\Lambda_i$. Thus the same finite subwalk is
\[
\Id\longrightarrow h
\]
at the finer resolution but becomes
\[
\Id\longrightarrow\Id
\]
in $\Lambda_i$, thereby satisfying the relevant same-endpoint condition.

This still does not mean that the actual generator word already contains a return. Strong repetitivity supplies the final step: the required finite configuration occurs repeatedly in the actual orbital / tile structure, so a genuine closed subwalk---a concrete return witness---can be found. In the fewest possible mathematical words:
\[
\boxed{
\begin{gathered}
\text{open at finer germ resolution}\\
\Downarrow\\
\text{same-endpoint in an admissible regular-limit quotient}\\
\Downarrow\\
\text{realized again as an actual finite return witness.}
\end{gathered}}
\]
This is the source-domain prototype of the three psychological disciplines above. Alternative organization is not arbitrary, and in the end it must land back in actual finite structure.

\section{The Formal Theorem Is Deferred to Part IV: Here We Retain Only Its Dependency Structure}

Kuang's Theorem 5.1 will be stated in Chapter 28 with its original five hypotheses. Part III does not reproduce it here. It retains only the dependency diagram useful for translation:
\[
\begin{gathered}
\text{pNH nearby trivialization}\\
+\ \text{finite germ and context control}\\
+\ \text{thin / one-ended orbital geometry}\\
+\ \text{strong repetitivity}\\
\Downarrow\\
\text{uniform return forcing in the source model.}
\end{gathered}
\]
The conclusion is strong because the conditions work together. Part III has no authority to abbreviate this as ``pNH automatically produces psychological return,'' nor may it translate unique singularity, a finite abelian group of germs, or thinness into personality properties.

\section{Which Conditions May Be Borrowed, and Which Should Remain in Mathematics?}

To keep technical assumptions from filling Part III again, we divide them into three groups:

\begin{center}
\small
\begin{tabularx}{0.98\textwidth}{>{\bfseries\raggedright\arraybackslash}p{0.28\textwidth}Y}
\toprule
How Part III treats it & Condition \\
\midrule
Borrow one structural intuition & Purely non-Hausdorffness: a distinction retained in the current coordinate need not remain visible in every admissible nearby context. \\
Borrow one methodological discipline & Finite context family and strong repetitivity: alternative organization must be constrained and must return to the same finite material. \\
Do not translate at present & Unique singularity; finite abelian group of germs; invariant domains; thin / one-ended orbital geometry. They matter to the theorem but currently have no reliable psychological counterpart. \\
\bottomrule
\end{tabularx}
\end{center}

The table is itself an act of self-restraint in the Geometry of Existence. A condition can be important in a mathematical proof without being entitled to a psychological counterpart.

\section{What Additional Questions Does This Benchmark Actually Let Psychology Ask?}

If this chapter eventually enters empirical work, it should not try to verify whether ``the brain is a groupoid.'' It should ask smaller questions:
\begin{itemize}
\item For the same event sequence, does changing a predeclared self-reference / appraisal frame systematically alter whether participants judge that they have ``returned to the same pattern again''?
\item Does the change also alter memory retrieval, subsequent prediction, or action choice, rather than verbal report alone?
\item Can the effect of an alternative frame recur across multiple comparable episodes?
\item If all such effects are fully explained by existing constructs such as decentering, metacognition, or emotion regulation, should this project stop adding new terminology? \citep{fleming2012,bernstein2015}
\end{itemize}
The actual small hypothesis is only this:
\begin{quote}
\textbf{Does the judgment ``does this count as returning to the same position again?'' change with an operationalizable organizing coordinate rather than being determined by surface similarity alone?}
\end{quote}
If not, this part of the AF-by-discrete mapping should contract.

\section{Stop Line: This Theorem Does Not Mean That ``All Long-Term Suffering Can Eventually Be Explained''}

Real experience may contain genuine novelty. It may contain several irreducible organizing centers. There may be no finite family of alternative frames, and no comparable structure resembling strong repetitivity. Some problems may simply be inappropriate for return language.

\begin{warning}[Benchmark Stop Line]
This chapter does not claim that the brain has a purely non-Hausdorff singularity; that personality states or psychological states form a finite group of germs; that a life contains only one singular center; or that sufficiently long experience must contain a psychological return. The mathematical theorem exerts only one methodological pressure on Third Order: changing coordinates is permitted, but the new coordinate must be constrained and must bring the material with it.
\end{warning}

\section{Concept Passport}

\begingroup\scriptsize
\begin{longtable}{>{\raggedright\arraybackslash\bfseries}p{0.25\textwidth}>{\raggedright\arraybackslash}p{0.66\textwidth}}
\toprule
Psychological motivation & When does ``looking from another angle'' genuinely increase resolution, and when does it merely replace the current story with another story?\\
Psychological concept & Perspective switching; neighboring questions in reappraisal / decentering; alternative self-reference / evaluative frames.\\
Mathematical concept & Unique purely non-Hausdorff singularity; regular-limit quotient; finite context family; strong repetitivity; return forcing.\\
Preserved relation & The current coordinate does not monopolize the judgment of ``whether it came back''; alternative organization must have a source, change an explicit relation, and return to the same finite material.\\
Discarded content & Literal psychological realization of the group of germs; counting personality states; reification of orbital geometry as a psychological entity; mathematical growth conclusions.\\
Translation level & The mathematical theorem is strict in its native domain; the cross-domain lesson is Level II.\\
Candidate observables & Frame manipulation; return / endpoint judgment; memory retrieval; subsequent prediction; action choice; replication across episodes.\\
Non-example & ``Just look at it from another angle''; changing wording while changing no relational variable; inventing after the fact a new story that can never fail.\\
Failure conditions & Alternative frames cannot be predeclared; frame changes alter no structural judgment or future variable; all effects are fully explained by established constructs.\\
Reverse-translation check & Compete with mature constructs such as decentering, metacognition, appraisal, and emotion regulation; do not overwrite them with the vocabulary of the present model.\\
Translation remainder & How real cognitive systems generate alternative frames; how multiple frames compete; the corresponding neural, bodily, and social mechanisms.\\
\bottomrule
\end{longtable}
\endgroup

\section{Chapter Conclusion: A Genuine Change of Perspective Must Bring the Material Back}

What this chapter is trying to leave behind is not the phrase ``purely non-Hausdorff.'' It leaves a much more ordinary and stringent principle. A current interpretation can be coherent, can have a great deal of history behind it, and can produce the powerful feeling that ``this is simply how things are.'' Coherence does not confer final authority. Another organization may enter, but the words ``look from another angle'' do not give it immunity from testing.

Third Order therefore ultimately requires:
\begin{plainbox}
\begin{center}
\emph{The current perspective has no final authority;\
the alternative perspective has no exemption from evidence;\
a genuine shift of perspective must have a source, alter relations, and bring the material back.}
\end{center}
\end{plainbox}

In still more ordinary language:
\begin{quote}
\textbf{Of course you may look from another angle. But after you do, look again at the same thing and tell me exactly what more you can now see.}
\end{quote}

The next chapter applies this discipline to the observer itself. If every perspective has a source, then the observation position that keeps choosing perspectives, segmenting events, recognizing ``again,'' and deciding what counts as evidence must enter the same picture as well.

\chapter{When the Observer Enters the Picture: Reflexive Traditions and the Question of Final Authority in Third-Order Consciousness}
\label{chap:observer-enters-picture-en}

\par\medskip
\noindent\textbf{Why this chapter matters}\par
The preceding chapters have made many distinctions: where experience is simply flowing and where a root begins to acquire organizing authority; how a stretch of experience is bracketed as an episode; when ``similar'' becomes ``again''; how several ``agains'' can overlap; whether later context lets us see the past more finely or actually rewrites it; and why the same return can still lead to a different afterward.

All of these questions eventually return to a deeper one: who is doing the segmenting, comparing, naming, and judging? When I say, ``This is rejection again,'' ``This relationship is simply like this,'' ``I should let go now,'' or ``I finally understand my own pattern,'' does the observation position that is drawing the conclusion really stand outside the entire process?

Within its own theoretical line, Third-Order Consciousness pushes the reflexive problem to this point. It does not manufacture a higher and clearer observer. It lets the observer itself enter the picture: where does it look from, what past does it use for comparison, at what scale does it segment, why does it select this particular grammar of ``again,'' and toward what result is it urgently trying to move experience?
\par\medskip

\par\medskip
\noindent\textbf{Where this chapter stands in the history of ideas: reflexivity is not the property of one tradition}\par
The critical reconstruction in Part I already acknowledged that including the observer in the observed system is not a discovery unique to Third-Order Consciousness. Second-order cybernetics systematically discussed observer participation, while Krishnamurti repeatedly asked from the side of psychological observation whether the observer is really outside experience \citep{vonfoerster2003,krishnamurti1954,krishnamurti1969}. We therefore make no claim to priority for observer endogeneity. It tests a narrower question: once the observer is put back inside the system, can the newly exposed position again acquire unconditional judicial authority? A candidate distinctive contribution of Third-Order Consciousness is to place that question back inside subject formation, narrative compression, corrective aims, and concrete relational practice. The germ / non-Hausdorff geometry that follows adds further formal resolution to this shared reflexive problem.
\par\medskip

\begin{methodbox}[Interface card: Third-Order Consciousness and the reflexive visibility of organizing conditions]
\textbf{Psychological motivation:} Reflection, metacognition, and self-regulation may all retain a manager that is treated as transparently given. Third Order asks further what histories, rules, roles, aims, and present framings organize that very position.

\textbf{Third-Order native domain:} The subject loses final authority. It retains functions of action, responsibility, and narrative, but no longer possesses unconditional authority over experience.

\textbf{Formal interface:} Rooting, episode segmentation, return grammar, scale / context, and continuation resolution cease to be merely background parameters and become explicit organizing conditions open to audit.

\textbf{Preserved relation:} the observer is endogenous; the current root has a position; organizing rules can be repositioned, untelescoped, or resolved; alternative perspectives must still return to the same material.

\textbf{Level:} Level-II structural model.

\textbf{Not claimed:} The brain literally stores a set of groupoid parameters, or Third Order is a permanent state, a personality rank, or an established neural mechanism.
\end{methodbox}

\section{From Here On, No New Mathematical Objects Are Needed}

By this point in Part III, there is no need to keep adding mathematical terminology. What mattered in the preceding chapters was not how many terms we learned, but that each term forced psychology to ask one more question that an ordinary sentence could otherwise swallow.

\begin{center}
\scriptsize
\begin{tabularx}{0.98\textwidth}{>{\bfseries\raggedright\arraybackslash}p{0.30\textwidth}Y}
\toprule
Earlier interface & What additional question it forces psychology to ask \\
\midrule
AF-like background & Which complex processes already operate before a stable self-narrative takes organizing control? \\
Boundary / root & Which transitions make the existing script insufficient? Around what position is experience now being reorganized? \\
Birooted episode & Why is ongoing experience taken to begin here, end there, and count as ``this one thing''? \\
Return / ``again'' & What makes two occurrences not merely similar, but incorporated into one subject-relevant history? \\
Return grammar & How many different ``agains'' are competing for explanatory authority within the same stretch of experience? \\
Multi-boundary bridge-type prototype & Can an ``again'' be jointly organized by passage through several organizing sectors rather than one unique root? \\
Incompressible remainder & Which material cannot yet be absorbed by an old grammar: new structure, or a model that is too coarse? \\
Scale / filtration & When the past is ``understood later,'' has background been added, interpretation changed, segmentation changed, or has the remembered material itself changed? \\
Residue / resolution & Even when the same pattern returns, has its future continuation changed this time? \\
pNH observation-site benchmark & Can the distinction structure retained by the current root be compared under other admissible resolutions generated by the system itself? Can an alternative perspective return to the same material and change a testable relation? \\
\bottomrule
\end{tabularx}
\end{center}

All of these questions now converge on the same place: the observation position performing these acts of organization must also have a position of its own.

This is both an internal convergence point of Third-Order Consciousness and a problem space shared with reflexive traditions such as second-order cybernetics. From here on, this book does not treat observer endogeneity itself as a novel discovery of Third Order. What remains to be tested is whether the framework adds further distinctions concerning subject formation, narrative, corrective aims, and whether the newly visible observer re-acquires final authority.

\begin{plainbox}[A particularly clean mathematical example: root versus pNH observation site]
The root of Chapter 15 is already enough to express: ``I am currently looking from here.'' That can still belong to second-order subject formation.

The pNH benchmark of Chapter 21 adds one further step. The current position $\xi$ not only preserves a distinction; the same system supplies a family of admissible nearby resolutions through which we can see why that distinction matters here and why it may weaken or collapse under another legitimate organization.

Thus this book keeps
\[
\text{root at }\xi
\qquad\text{and}\qquad
\mathcal O_\xi=(\xi;\mathcal C_\xi)
\]
separate. The former is the formal prototype of an ego-position; the latter is the benchmark for a third-order observation site. The third-order feature does not come from the word ``singular.'' It comes from the current position being relativizable by its own admissible alternative resolutions.
\end{plainbox}

This gives ``the observer also enters the picture'' a very concrete formal reading. The observer does not disappear, and the current root is not cancelled. What changes is that the coordinates through which it distinguishes, completes returns, and evaluates experience become objects that can themselves be compared.

\begin{plainbox}[An equally important mathematical example: multiple boundary sites do not form one larger root]
The pNH benchmark of Chapter 21 concerns resolution plurality \emph{within} one observation site: fix $\xi$ and compare how the germ distinctions it retains behave under different admissible nearby resolutions.

The multi-bridge prototype in Kuang's preprint v1 concerns something else \citep{kuang2024v1}: one high-level traverse can pass successively through distinct boundary / bridge sectors and produce a type word
\[
x_1x_2\cdots x_k.
\]
A return grammar can therefore be jointly organized by the ordered passage through several sites.

The two prototypes must remain side by side:
\begin{center}
\resizebox{0.96\linewidth}{!}{$\displaystyle
\begin{array}{ll}
\text{single-site pNH:} & \text{resolution plurality within one position;}\\
\text{multi-boundary:} & \text{passage plurality among several organizing sectors.}
\end{array}
$}
\end{center}
\end{plainbox}

Translated back to Third-Order Consciousness, ``the subject loses final authority'' therefore has a second meaning. Not only does no single root possess final authority over resolution; no single organizing sector automatically monopolizes the entire return grammar. Third Order can ask: which position supplies the entrance to the present ``again,'' through which positions does it pass, and at which position does it acquire return-completion?

This still does not define a ``complete structural subject.'' How several organizing sectors integrate, whether different sectors carry different local germ profiles, and how transport among them works belong to the later integration problem. The important distinction is that multi-site return coding already has a source-side mathematical prototype; complete subject integration remains future work.

\section{The Real Turn: from ``I Observe Experience'' to ``The Position of Observation Also Belongs to Experience''}

Ordinary reflection can already say:
\begin{quote}
``I am angry.''
\end{quote}
Metacognition can add:
\begin{quote}
``I know that I am angry now, and I am uncertain about the reliability of my present judgment.''
\end{quote}
Decentering can add another degree of distance:
\begin{quote}
``The thought `he does not respect me' is appearing. That thought need not automatically be the fact.''
\end{quote}

All of these capacities have genuine psychological value. Third Order does not need to disparage them in order to establish itself. It adds only one further question:
\begin{quote}
\textbf{At this moment, who or what is deciding what counts as fact, what counts as maturity, what counts as the same pattern, and what ought to be changed?}
\end{quote}

For example, imagine that a person becomes angry after receiving no reply. A little later, a calmer thought may come to mind:
\begin{quote}
``I should not care this much. I should be more mature and let it go sooner.''
\end{quote}
That sentence may be helpful and may reduce impulsive action. Third Order does not place it outside the structure merely because its tone is calmer. It continues to ask: why is ``not caring'' the goal? From which value standard does ``mature'' come? Does this manager also carry histories of fear of losing control, maintenance of self-image, and striving for self-upgrade?

The direction of observation therefore changes:
\begin{center}
\resizebox{0.96\linewidth}{!}{$\displaystyle
\boxed{
\text{experience enters observation}
\longrightarrow
\text{the conditions organizing experience also enter observation}.}
$}
\end{center}

One of the central claims of Third-Order Consciousness in its own source tradition is precisely that the subject loses final authority: the subject remains a focal point, observation site, and working interface, but is no longer presumed to be the final owner of existence \citep{toc2026}. This meets Krishnamurti's pressure test that ``the observer is the past'': the past no longer appears only as something observed; it may also be the eye through which the present is being viewed \citep{krishnamurti1954,krishnamurti1969,kft2026a}.

\section{Third Order Is Not One More Step Back: Otherwise We Manufacture an Infinite Tower of ``Higher Selves''}

A natural misunderstanding is: if the present observer is biased, step back once more and find a higher observer that can observe it. The structure then becomes
\begin{center}
\resizebox{0.96\linewidth}{!}{$\displaystyle
\text{I observe emotion}
\longrightarrow
\text{a higher I observes the first I}
\longrightarrow
\text{an even higher I observes the first two}
\longrightarrow\cdots
$}
\end{center}
Nothing has been solved; the judge's seat has merely been moved higher and higher.

Third Order removes the assumption that the highest currently visible level is naturally exempt from explanation. It does not ask how to find the purest observer. It asks how any position that acquires explanatory authority is generated from the same psychological scene.

The following four sentences make the distinction concrete:

\begin{center}
\small
\begin{tabularx}{0.98\textwidth}{>{\raggedright\arraybackslash}p{0.35\textwidth}Y}
\toprule
Sentence & What it adds \\
\midrule
``I am angry.'' & Current experiential content enters language. \\
``I notice that I am angry.'' & Meta-awareness is added to the current state. \\
``I should not be angry; I need to control myself.'' & A managerial position begins to evaluate and govern first-layer experience. \\
``Even this position that says `I should not be angry' has a source.'' & The manager itself falls back into the same process and loses automatic structure-external privilege. \\
\bottomrule
\end{tabularx}
\end{center}

The fourth sentence is the structural change that matters most here. It does not deny the first three. It keeps one question from being prematurely foreclosed: \emph{what entitles this position that is explaining everything to stand where it stands?}

\section{``The Observer Is the Observed'': What Is Cancelled Is the Internal Judge's Immunity from Examination}

In ordinary psychological language, experience is easily divided in two: on one side are the emotions, desires, and impulses that need to be managed; on the other is the ``I'' that governs them.
\begin{quote}
``I am dealing with my anger.''
\end{quote}
The sentence has obvious practical uses. The problem is that it can quietly create an unexamined courtroom: anger sits in the defendant's chair, while the rational ``I'' sits on the bench.

If the judge itself is made of comparison, shame, ideals, self-image, earlier hurt, and social norms, then a calmer tone alone cannot establish that it is outside the conflict. The interpretation of ``the observer is the observed'' in Third-Order Consciousness brings this split to the foreground: the observer ceases to be an external judge and returns as a role within the same scene \citep{toc2026}.

This does not mean that ``anger is me,'' that emotion should not be regulated, or that all psychotherapy is a form of division. Effective regulation can remain effective. Third Order retains a more basic discipline:
\begin{quote}
\textbf{A regulatory function can be effective without thereby granting the regulator a highest identity without history.}
\end{quote}
Third Order therefore does not cancel the capacity for judgment; it cancels the judge's immunity from examination.

\section{The First Fire and the Second Courtroom: Why Secondary Conflict Becomes Visible}

Psychological suffering often has at least two layers. The first may be the direct event: anger, fear, jealousy, hurt, the urge to ask, or the urge to flee. The second is the courtroom built immediately afterward:
\begin{quote}
``Why am I still like this?''\quad
``A mature person should not care so much.''\quad
``I have to prove that I have moved on.''
\end{quote}

The second layer can sometimes aid behavioral control. It can also add shame, suppression, comparison, and prolonged rumination to the first-layer response. The earliest change under Third Order need not be to extinguish the first fire; it may be to make the second courtroom visible. The manager rushing to repair, correct, and upgrade the self is also a psychological event occurring now.

This connects directly with the empirical question already raised in Part I. First-layer response and second-layer self-governance can be recorded separately, and one can ask whether the latter shortens, prolongs, or reactivates the former, rather than treating the presence of a ``higher evaluation'' as automatic evidence of successful regulation \citep{bernstein2015,fleming2012}.

\section{Third Order Does Not Delete the Root: the Point Remains, but No Longer Obscures the Whole Scene}

One of the most important original intuitions of Third-Order Consciousness must remain intact here: Third Order is not the absence of a point. It is that once the point is opened again, the wider scene becomes visible again \citep{toc2026}.

A subject position is still necessary. Writing a paper requires a persistent aim; signing a contract requires identity persistence; caring for someone requires commitment; confronting genuine harm requires clear judgment. A system with no root at all could neither act nor assume responsibility.

What Third Order changes is not the existence of the root, but its authority. A subject can perfectly well say:
\begin{quote}
``Given these repeatedly observed behaviors, I have decided not to continue this relationship.''
\end{quote}
Third Order does not require the decision to be withdrawn. It requires us to be able to state at the same time which facts were used, which parts of the past were brought into the judgment, at what scale different episodes were treated as the same pattern, which value rule was applied, and what new material could change the conclusion.

The root therefore steps down from ``final owner'' to a located working interface:
\begin{quote}
\textbf{At this moment I must act from some position; but no current position exhausts the entire structure.}
\end{quote}
This is not relativism. Precisely because the position is explicit, judgment can bear responsibility more clearly.

\section{Not the Absence of Narrative: Narrative Begins to Know That It Has a Scale}

The second-order subject has to compress history. Without narrative compression we could not promise, plan, learn, or organize decades of life into units from which action is possible.

Third Order does not demand that stories be destroyed. It only asks a story to remember again that it is a presentation. Sentences such as ``I have always been disrespected,'' ``I have never been good at forming relationships,'' and ``my knowledge is too narrow'' may contain real information, but each crosses scales. Third Order asks:
\begin{itemize}
\item Which episodes does this sentence compress?
\item Which counterexamples are skipped in the middle?
\item Which omitted details would leave the conclusion unchanged?
\item Which omitted differences would actually change future continuation?
\end{itemize}
Third Order is therefore not the absence of stories. It is the loss of ontological status by the story. Compression is not prohibited; compression must be capable of being untelescoped.

In one sentence:
\begin{quote}
\textbf{Narrative may still guide the way, but it may not quietly declare the map to be the world itself.}
\end{quote}

\section{``Again'' Also Enters Observation: Who Decides That This Still Counts as the Same Pattern?}

Chapters 17 and 18 showed that ``again'' is one of the smallest yet most powerful temporal acts of Second-Order Consciousness. It pulls the current episode into the past and suppresses enough differences to permit a return-completion.

Third Order now turns this operation itself into an object of observation. When the sentence
\begin{quote}
``I have been rejected again''
\end{quote}
appears, we ask not only whether it is true or false, but also:
\begin{itemize}
\item How was the current episode segmented?
\item Which parts of the past were retrieved, and which did not appear?
\item Which common relations are preserved by the category ``rejected''?
\item Which differences that affect continuation were suppressed by the old label this time?
\end{itemize}
``Again'' thus becomes more than a call to the past; it becomes an organizing operation open to audit. The subject may retain a genuine pattern while allowing residue to say: yes, this has happened again; but what follows need not be the same afterward.

This is an important form of third-order openness. It does not obtain freedom by denying an old pattern, and it does not let the old pattern write the future in advance.

\section{Four Formal Operations Can Now Be Translated into Four Ordinary Questions}

The formal side has introduced Inflate, Telescope, Untelescope, and Resolve. At this convergence chapter there is no need for psychology readers to memorize them as technical operations. They correspond to four distinct questions:

\begin{center}
\scriptsize
\begin{tabularx}{0.98\textwidth}{>{\bfseries\raggedright\arraybackslash}p{0.18\textwidth}p{0.31\textwidth}Y}
\toprule
Formal name & Ordinary question & Third-order discipline \\
\midrule
Inflate & If this event is placed inside a larger context, what else becomes visible? & Adding context does not automatically make the old interpretation more correct, and does not automatically erase existing material. \\
Telescope & Which details can be compressed into one sentence? & Compression can serve action, but must not erase differences that change continuation. \\
Untelescope & What was omitted between a local fact and a global identity judgment? & Restore episodes, variables, and inferential links flattened by a large word. \\
Resolve & Even if we acknowledge ``this happened again,'' what is different in what follows this time? & Keep return-completion and continuation on separate ledgers; do not use repetition to deny change, or change to deny repetition. \\
\bottomrule
\end{tabularx}
\end{center}

These four questions do not automatically form a therapeutic procedure and need not be executed in order during real experience. They are simply four ways of giving the current interpretation a position again.

\section{Putting the Whole Structure Back into One Evening: Convergence of the Unanswered-Message Case}

We can now put all of the preceding interfaces back into one simple scene instead of leaving each term suspended in isolation.

Imagine that a person contacts someone they have recently met and receives no reply. At first, much of the process remains in the background. Their attention keeps returning to the phone; they become increasingly irritated and briefly lose focus on their work. After a while, the current interaction script no longer tells them what to do next. At that point an ego-position becomes salient: ``I am the one who already initiated contact and must now maintain my principles.''

The experience can then be segmented into episodes: sending the message, waiting, a checking cycle, deciding to stop contact, and continuing to rehearse a future conversation in imagination. If several imagined rounds repeatedly come back to ``the rule has not been acknowledged, and I need to make myself clear,'' we have a candidate return. At the same time, other return grammars may coexist: ``I am being ignored again,'' ``I am proving again that I do not care,'' ``why am I still thinking about this again?''

Later, more context becomes available. The subject may discover that the earlier narrative ``every silence means disrespect'' compressed several different kinds of events; or the subject may discover that a particular pattern really is stable. Even if the return remains, the present occurrence can have a different residue: there is no renewed contact, no prolonged checking, and no conversion of the other person's behavior into proof of one's own value.

Only here does the specifically third-order move appear. It no longer asks only:
\begin{quote}
``Why is she not replying?''
\end{quote}
or only:
\begin{quote}
``Why am I so angry?''
\end{quote}
It begins to see at once:
\begin{quote}
\textbf{Who or what segmented the silence as this episode? Who selected the grammar ``ignored again''? Who stipulated that ``a mature person should let go sooner''? Who is in a hurry to turn this event into a conclusion about me?}
\end{quote}
The word ``who'' here is not a search for another little person in the head. It tracks which organizing position, which history, which rule, and which scale are acquiring explanatory authority.

Third Order does not decide whether the person should contact the other again. The action rule may remain exactly the same. What changes is that external fact, first-layer response, return grammar, value rule, and the later manager are no longer compressed into one indivisible ``I.''

\section{The Academic Case Is the Same: from ``I Cannot Follow This'' to ``What Kind of Researcher Am I?''}

During an unfamiliar talk, the initial facts may be small: insufficient prerequisites, a missing terminological interface, several minutes of lost tracking. The subject can rapidly cross scales and write:
\begin{quote}
``My knowledge is too narrow.''
\end{quote}
That judgment is not necessarily wrong. Third Order does not immediately replace it with a more positive story.

It restores the omitted organizing conditions to the picture: the distance between the talk and one's own field, the day's attentional state, the comparison target, professional-identity expectations that ``I should be able to formulate questions quickly,'' and which earlier successes or failures the current root has retrieved.

The question therefore expands from
\begin{quote}
``Is this judgment correct?''
\end{quote}
to
\begin{quote}
\textbf{How did this judgment move from a local accessibility failure to a global self-evaluation?}
\end{quote}

If comprehension improves substantially after a finite terminological interface is supplied, the original frame should be revised. If multiple comparable contexts show the same knowledge gap, the harsher judgment may gain support. Third Order is not responsible for making conclusions kinder. It is responsible for making the structure that generates evaluation visible, comparable, and revisable.

\section{Third Order and Freedom: Not More Options, but Conditions No Longer Masquerading as Commands}

Third Order is easily paraphrased as ``having more choices.'' That is insufficient. A person can have many options while a hidden root has already decided which options are ``like me,'' which are humiliating, which are mature, and which are unacceptable. There may be many options while the conditions organizing them remain hidden.

We therefore describe third-order freedom as an increase in the transparency of conditions. Action still occurs and responsibility remains, but the actor is less likely to hear historical conditions, role demands, and narrative rules as absolute commands arriving from outside the structure.

This does not abolish real power. There are genuine institutional asymmetries between advisor and student, manager and employee, doctor and patient. Third Order cannot psychologize structural constraints into ``that is merely your projection.'' Giving the observer a position means that private history and social position both enter the picture, rather than analyzing only how the less powerful party perceives the problem.

\section{Third Order Can Be Practiced, but It Cannot Become a Medal}

The analyses above can of course be trained. A person can become increasingly skillful at distinguishing fact from interpretation, recognizing ego-positions, resegmenting episodes, checking the witness behind an ``again,'' restoring compressed intermediate layers, and testing alternative frames. These abilities can improve.

But this is also where Second Order can most easily reabsorb Third Order:
\begin{quote}
``I am very good at observing myself.''\quad
``I am much more aware than I used to be.''\quad
``I have entered Third Order.''
\end{quote}
Third Order itself then becomes identity capital; observational ability becomes a new performance metric; the subject returns to the throne in more sophisticated clothing.

This is why Krishnamurti's suspicion of method, psychological achievement, and becoming remains relevant \citep{krishnamurti1954,krishnamurti1969}. Third Order may have trainable prerequisite capacities, but it cannot be completed by the identity claim ``I have become a third-order person.''

\section{If This Enters Psychological Research, Do Not Begin by Manufacturing a ``Third-Order Total Score''}

If the only empirical product of this chapter were a single Third-Order Consciousness scale, a complex structure would probably have been compressed too early into one score.

A better research route would operationalize several processes separately:
\begin{itemize}
\item \textbf{Awareness of the source of the observation position:} can individuals identify which histories, roles, or aims generate the current evaluative rule, rather than merely report the evaluation itself?
\item \textbf{Root sensitivity:} when the self-reference frame changes, do event segmentation, memory retrieval, return judgments, and subsequent predictions move systematically?
\item \textbf{Second-order governance:} after a first-layer emotion, do shame, self-blame, suppression, or improvement attempts shorten, prolong, or reactivate the first-layer response?
\item \textbf{Scale transparency:} can individuals distinguish a local episode, a long-term pattern, and an identity judgment, and restore compressed intermediate layers after manipulation?
\item \textbf{Return-completion--continuation separation:} after acknowledging an ``again,'' can participants also identify genuine differences in the present continuation?
\item \textbf{Alternative-frame discipline:} can alternative interpretations be predeclared, return to the same material, and generate differential predictions?
\end{itemize}

These variables must compete for explanatory power with established constructs such as metacognition, decentering, self-distancing, emotion regulation, and autobiographical memory \citep{fleming2012,bernstein2015}. If ``Third Order'' adds no stable distinction, there is no reason to create an additional global construct.

\section{Concept Passport}

\begingroup\scriptsize
\begin{longtable}{>{\raggedright\arraybackslash\bfseries}p{0.25\textwidth}>{\raggedright\arraybackslash}p{0.66\textwidth}}
\toprule
Psychological motivation & As more content becomes reflectively accessible, must the observation position that names, compares, evaluates, and governs experience also disclose its own source?\\
Third-Order native domain & The subject loses final authority; the observer is formed by past, language, value, relationships, and social position; the observer cannot be presumed to stand outside the observed.\\
Formal interface & Rooting; episode segmentation; return grammar; scale / context; filtration; continuation resolution become explicit, auditable parameters.\\
Preserved relation & Observer endogeneity; organizing conditions enter observation; the current root can be repositioned; narrative can be untelescoped; return-completion and continuation can be kept on separate ledgers.\\
Discarded content & Qualia; neural mechanism; mathematical identification of the complete subject; ontological claims of permanent selflessness; a hierarchy of third-order personalities.\\
Translation level & Level-II operational structural model.\\
Candidate observables & Reports of the source of an observation position; self-reference manipulations; event boundaries; return judgments; second-order self-governance; scale switching; subsequent prediction and action choice.\\
Non-example & ``I am already very good at awareness''; acquiring immunity by speaking calmly; rewriting all real power problems as personal projection; treating Third Order as lower affect or a higher personality level.\\
Failure conditions & Proposed observer parameters cannot be operationalized reliably; they add no explanatory value for behavior, memory, or prediction; all effects are more simply explained by existing constructs.\\
Reverse-translation check & Do not use Third Order to overwrite metacognition, decentering, self-distancing, or emotion regulation; do not declare the mathematical model a scientific proof of Krishnamurti's thought.\\
Translation remainder & Phenomenological mineness; qualia; integration among multiple subject positions; neural and social mechanisms of observer visibility.\\
\bottomrule
\end{longtable}
\endgroup

\section{A Formal Interim Definition of Third-Order Consciousness}

At this point we can state the formal definition used in Part III. It adds a structural interface to the intellectual definition of Part I, but it is still not a mechanistic definition.

\begin{definition}[Third-Order Consciousness: Part-III Structural-Model Definition]
Third-Order Consciousness is modeled as the reflexive visibility of organizing conditions: the system not only processes current experience, but can also make available for access and comparison the conditions participating in that organization, including the current ego-position, episode segmentation, return grammar, history / scale presentation, and continuation resolution. The ego-position carrying out observation, evaluation, and governance remains part of the same architecture and therefore does not automatically acquire structure-external privilege.
\end{definition}

The Chapter-21 package
\[
\mathcal O_\xi=(\xi;\mathcal C_\xi)
\]
is not a necessary condition of this definition. It is an especially clear mathematical benchmark: a distinguished root remains, while the distinction structure retained there can be compared again through admissible resolutions generated by the system itself. In other words, the pNH observation-site package displays, in a controlled mathematical setting, how ``the current position remains valid while the resolution rules of that position can themselves be relativized.'' It is not a general mathematical definition of Third-Order Consciousness.

Likewise, the multi-boundary bridge-type prototype is not a defining condition. It supplies an orthogonal formal reminder: a return grammar may be jointly encoded by passage through several boundary sectors. The operational definition of Third Order therefore does not require a unique observation site or a unique root. It requires only that the positions, segmentations, and return grammars currently organizing experience be able to enter the same auditable process.

This definition is the present Level-II operational structural definition of the project. It formalizes the relation ``how the conditions of observation enter observation.'' It does not explain qualia or phenomenological mineness, and it does not claim that the brain literally implements an AF-by-discrete groupoid.

\section{Chapter Conclusion: the Point Has Not Disappeared; the Scene Reappears}

The core of Third Order is, in the end, simple. Experience still happens. The body still hurts. Relationships can still wound. A researcher still needs to judge whether a proof is correct, and a person still needs to decide whether to continue a relationship. Language does not disappear, identity does not disappear, and narrative does not disappear.

What changes is their authority.

A root can still organize action, but it knows that it has a position. A story can still compress history, but it knows that it has a scale. An ``again'' can still identify a real pattern, but it must make room for difference. An observer can still regulate experience, but observing no longer automatically places it outside experience.

Third Order therefore does not delete the ``I.'' It places the ``I'' back into the scene from which it came:
\begin{plainbox}
\begin{center}
\emph{not no ``I,'' but an ``I'' without final authority;\
not no story, but a story that knows it is one scale;\
not no judgment, but judgment that also reveals where it arises from.}
\end{center}
\end{plainbox}

Third-Order Consciousness describes this as the wider scene appearing again once the point is opened back into relation \citep{toc2026}. In the more restrained structural language of this project: the observation position continues to exist, but it reconnects with the background, history, segmentation, and continuation relations from which it emerged.

This also explains why Third Order is not an escape from experience, but a more thorough return to experience. In the end, even the ``I'' most eager to step forward and say ``let me handle all of this'' falls back into everything that is happening.

The next chapter adds no new correspondence. Part III ends by auditing the bridge: which mappings actually stand, which remain provisional hypotheses, and which things must explicitly remain outside the bridge.

\chapter{The Bridge Reaches This Far: Structural Audit and Stop Lines for Part III}
\label{chap:part3-audit-en}

Part III stops adding correspondences here. Its task has not been to ``put psychological life inside mathematics,'' but to test whether relations already defined independently on both sides can preserve the same constraints across the interface. If a correspondence can survive only by repeatedly rewriting the mathematical definition, adding hidden variables, or changing psychological segmentation after the fact, it should exit.

\section{What Part III Has Actually Connected}

Only the following families of relations remain in the core:

\begingroup\scriptsize
\begin{longtable}{>{\raggedright\arraybackslash\bfseries}p{0.22\textwidth}>{\raggedright\arraybackslash}p{0.34\textwidth}>{\raggedright\arraybackslash}p{0.24\textwidth}}
\toprule
Relation & Current claim & Main exit condition \\
\midrule
Episode / birooted traverse & Entrance--finite process--exit can serve as a Level-II structural skeleton. & Psychological segmentation fails to yield reproducible endpoint relations over time.\\
Repetition / recognized return & ``Appearing again'' and ``being organized as coming back'' must be kept on separate ledgers. & After controlling similarity and frequency, return-completion adds no predictive information.\\
Organizing position / root & Where comparison and retrieval of history currently occur is not identical with the whole system. & The proposed position cannot be encoded independently, or is only a post-hoc label.\\
Return-completion / continuation & The same return-completion can preserve different future profiles. & Future differences are fully explained by surface cues or a simpler model.\\
Finite singular-like loci & The additional empirical hypothesis $H^{\mathrm{psych}}_{\mathrm{fin}}$ undertaken by this project. & Expanding material continually forces new loci and the old finite set loses predictive value.\\
Non-Hausdorff / pNH benchmark & A resolution generated on the mathematical side: becoming visible does not imply acquiring final authority again. & It generates no question beyond existing reflexive / decentering language.\\
\bottomrule
\end{longtable}
\endgroup

The last row also enforces a provenance discipline. Observer endogeneity is not unique to Third-Order Consciousness; traditions such as second-order cybernetics have discussed observer participation systematically \citep{vonfoerster2003}. If the mathematics adds anything, the gain should appear in a finer distinction---for example, ``the observation position has entered the system'' and ``that position has or has not reacquired local final authority'' are not the same statement.

\section{Three Statuses: Standing, Awaiting Test, and Left outside the Bridge}

What already stands are definitions internal to the two native domains and several restricted relations. On the mathematical side, root, traverse, return, germ, and residue each have their own object level. On the psychological side, episode, return judgment, self-reference frame, and continuation profile must first be established in ordinary language and by independent coding.

What remains to be tested is whether the cross-domain interface genuinely increases empirical resolution. A Level-II correspondence can move locally toward Level III only if it supplies incremental information under held-out material, baseline comparison, and reproducible coding.

Explicitly left outside the bridge are phenomenological mineness, qualia, full embodiment, affective valence, neural implementation, complete subject integration, and any claim that ``the ontology of consciousness is a groupoid.'' If later chapters touch these questions, they can appear only as remainders or future problems.

\section{How to Retreat When Something Fails}

Retreat follows a minimum-modification rule. If a single-root hypothesis fails, only the single-root hypothesis is withdrawn. If a particular singular-interface detector fails, only that empirical translation is withdrawn. Failure of $H^{\mathrm{psych}}_{\mathrm{fin}}$ does not rewrite source-side AF-by-discrete mathematics. Failure of a psychological case likewise cannot modify a purely mathematical theorem. Conversely, a mathematical theorem cannot pay the evidential cost of psychology merely by being correct.

Part III therefore retains only one global claim:
\begin{plainbox}
\textbf{AF-by-discrete / germ geometry supplies a controlled formal language for a family of psychological questions about position, return-completion, scale, and continuation. The present cross-domain evidence remains mainly at Level II, and only the singular-germ substructure of the source-side boundary geometry is activated for this interface.}
\end{plainbox}

The bridge stands because neither bank has disappeared. From Part IV onward, psychological language temporarily recedes. Mathematics assumes responsibility only for the formal consequences it has already promised.

\part{Formal Theory: Concrete Occurrences, Finite Witnesses, Return Language, and Germ Resolution}

\begin{sourcebox}[Part IV note]
Part III explained why root, episode, return, scale, and resolution are needed. Part IV does only one thing: it checks whether the mathematical relations being borrowed actually stand. Every definition and theorem is stated in a fixed mathematical presentation. If generators, telescoping, tile presentation, or germ decoration are changed, one must state again which objects may still be identified as ``the same.'' Psychological examples from Part III appear here only as reading aids; they are not mathematical conclusions.
\end{sourcebox}

\begin{plainbox}[Standing setup for Part IV: keep four object layers separate]
Throughout Part IV, fix a Bratteli / tile-inflation presentation, a generator / local-action alphabet $S$, and the corresponding finite tiles and boundary convention. Bookkeeping always proceeds through four layers:
\[
\begin{gathered}
\text{a source word / factor}\\
\Downarrow\\
\text{its concrete occurrence and finite witness at a specified position}\\
\Downarrow\\
\text{its type after concrete position is forgotten}\\
\Downarrow\\
\text{if finer continuation is required, germ decoration and a lift are added.}
\end{gathered}
\]
The notation $\mathcal R\subseteq S^*$ always denotes source return words / factors. $W_R$ provides concrete evidence for the closed internal excursion. Only $R=(u;W_R)$ is a witnessed datum. A type is a classification obtained by forgetting occurrence data. $\rho_\xi$ is defined only after a singular sector, anchor, and germ decoration have all been fixed. Many apparently minor notational restrictions below prevent these four layers from impersonating one another.
\end{plainbox}

\begin{plainbox}[Part IV also inherits a lower-level finite-interface principle]
The calculus is not taking place on an arbitrary huge graph. Relative to a fixed finite generating set, AF-by-discrete already compresses global organization into finitely many persistent boundary sectors; finite connectors continue to use the same label grammar as generator germs. General groupoid bounded type adds uniform boundary control for each finite AF fibre / tile. In the finite-rank tile-inflation form used here, the bounded number of tile types also controls the aggregate boundary / connector burden across a level. Thus the witness, decomposition, multi-boundary return, and residue calculus of Part IV lives on the layered structure
\[
\begin{gathered}
\text{large AF-like internal regions}\\
+\ \text{finite persistent interfaces}\\
+\ \text{uniform per-tile control, with aggregate level control in the finite-rank tile form.}
\end{gathered}
\]
Without the bounded-type finite-tile control, only the first two items remain available. And without separate finite-rank control on the number of tile types, an aggregate levelwise bound may not be used silently.
\end{plainbox}

\begin{warning}[Mathematical Definitions and Empirical Detection Must Remain Separate]
In the mathematical native domain of Part IV, boundary points / germs are already given by tile presentation, boundary ancestry, and infinite continuation. Return witnesses do not define them. The idea of ``detecting boundaries through repeated returns'' belongs only to the reverse-translation / system-identification protocols of Parts III and V: when an interface is unknown in a target domain, repeated returns, traverses, and continuation failures may help localize relations worth testing as candidate boundary analogues. That detection protocol must not be written back into the source-domain definition.
\end{warning}

\chapter{Finite-Scale Calculus: Concrete Occurrences, Finite Witnesses, and Two Scale Operations}
\label{chap:finite-scale-calculus-en}

\par\medskip
\noindent\textbf{Why this chapter matters}\par
Part IV begins strict mathematics here, but psychology readers only need to grasp two very ordinary questions first.

First, if the same shape appears twice, how do we know which occurrence we are tracking? An argument, a period of waiting, or a talk can be ``very similar'' to an earlier one, but saying only that the shape is the same forgets where it occurred, which context followed it, and what it later led to. Mathematically, this is the distinction between type and occurrence.

Second, a change of scale actually contains two very different operations. One places an already identified small piece of structure inside a larger context and keeps tracking it. The other takes a larger process apart and asks which smaller pieces compose it. The first is embedding; the second is decomposition. They point in opposite directions, but they are not inverses.

The task of this chapter is to make both operations type-correct mathematical objects. Whether persistence, filtration, and residue stand later depends on not switching objects silently here.
\par\medskip

\begin{plainbox}[Four sentences for psychology readers to remember]
\begin{center}
\resizebox{0.96\linewidth}{!}{$\displaystyle
\begin{array}{rcl}
\text{occurrence} &=& \text{``this very instance, at this very position''},\\
\text{type} &=& \text{``after forgetting position, which shape does it belong to?''},\\
E &=& \text{``enlarge the context while retaining the same finite witness''},\\
D &=& \text{``fix the same high-level trajectory and cut it downward.''}
\end{array}
$}
\end{center}
All later notation in the chapter only makes these four sentences precise. The psychological language here is navigation only; it is not part of the mathematical propositions.
\end{plainbox}

\section{Why ``the Same Shape'' Is Not Enough: Fix This Concrete Occurrence First}

Let $P_m$ be a level-$m$ finite tile. For $n<m$, write
\[
\operatorname{Occ}_n(P_m)
\]
for the set of all concrete level-$n$ subtile occurrences produced inside $P_m$ by tile inflation. If
\[
Q\in\operatorname{Occ}_n(P_m),
\]
then $Q$ is already an actual copy inside $P_m$, so there is a natural inclusion
\[
j_Q:Q\hookrightarrow P_m.
\]

Why can this information not be omitted? The same abstract tile type may occur many times inside $P_m$. If one merely writes
\[
P_n\hookrightarrow P_m
\]
without specifying which $Q\in\operatorname{Occ}_n(P_m)$ is selected, the endpoint, connector, and germ context to be tracked later may all become ambiguous.

\begin{definition}[Nested Occurrence Chain]
We call a chain of concrete subtile occurrences a \emph{nested occurrence chain} if $n<m<\ell$ and
\[
Q_n\subseteq Q_m\subseteq P_\ell,
\qquad
Q_n\in\operatorname{Occ}_n(Q_m),
\qquad
Q_m\in\operatorname{Occ}_m(P_\ell).
\]
Its inclusion maps compose as ordinary subgraph inclusions.
\end{definition}

\begin{plainbox}[In ordinary language]
A tile type is like a template; an occurrence is this particular instance of the template in the scene. Two copies may be completely isomorphic while meeting different connectors or sitting in different larger contexts. Whenever Part IV discusses persistence, connector data, or residue, it first retains ``where exactly is this occurrence?'' Position is forgotten only later, for classification or counting.
\end{plainbox}

This convention is compatible with the original bounded-type counting picture. A high-level tile is assembled from copies of lower-level tiles using boundary / connector data, and a deeper-level traverse induces earlier-level traverses along those actual copies \citep{kuang2026a}. Here we extract the occurrence-level bookkeeping and keep it separate from the later last-moment selection.

\section{Traverse Is Already Defined: What We Add Is a Witness That Can Be Carried into a Larger Context}

Part II, Definition 12.4, already defined a concrete traverse occurrence
\[
\tau=(Q;b_-,\gamma,b_+),
\]
where $b_-,b_+\in\partial Q$ and the internal part of $\gamma$ meets no further point of $\partial Q$. Part IV retains that source object; it does not invent a second definition of traverse.

What is needed here is a slightly broader bookkeeping object. Once a traverse is placed inside a larger tile, its original endpoints need no longer be boundary points of the larger tile, but we still need to recognize that the original finite relation remains present.

\begin{definition}[Finite Birooted Witness]
A \emph{finite birooted witness} is a quadruple
\[
W=(Q;b_-,\gamma,b_+),
\]
where $Q$ is a concrete finite-tile occurrence, $\gamma$ is a finite trajectory in $Q$ joining $b_-$ to $b_+$, and the two endpoints are the distinguished positions inherited by the witness at its scale of origin. We call $W$ a \emph{closed witness} if $b_-=b_+$.
\end{definition}

\begin{plainbox}[A witness as a finite evidence card]
A witness card retains at least four things: the concrete support in which it lies, where it starts, which actual trajectory it follows, and where it ends. After the witness is placed in a larger context, its identity is tracked by these data. It need not remain an elementary traverse of the larger tile.
\end{plainbox}

Thus persistence will preserve finite occurrence data, not the false set inclusion
\[
\mathcal T_n\subseteq\mathcal T_{n+1}.
\]

\begin{example}[Part-III pointer, not a mathematical conclusion: contextualization and cancellation]
Part III said that an earlier episode can be placed inside a larger historical context and reinterpreted without thereby declaring that the earlier finite relation never occurred. The only formal support supplied here is this: if one wants to say that ``the old structure is still traceable,'' the preserved object should be a concrete witness, not a lower-level elementary traverse incorrectly reclassified as a higher-level elementary traverse. This example adds no psychological mechanism claim.
\end{example}

\section{Embedding: Enlarge the Context without Rewriting the Finite Relation Already Identified}

Fix a concrete occurrence relation
\[
Q_n\subseteq Q_m,
\]
where $Q_n$ is a level-$n$ occurrence and $Q_m$ a level-$m$ occurrence. Denote the inclusion by
\[
j^{Q_n\subset Q_m}_{n,m}:Q_n\hookrightarrow Q_m.
\]
If
\[
W=(Q_n;b_-,\gamma,b_+),
\]
define
\[
E_{Q_n,Q_m}(W)
:=
\bigl(Q_n\subseteq Q_m;
 j(b_-),j(\gamma),j(b_+)\bigr),
\]
where $j=j^{Q_n\subset Q_m}_{n,m}$. The first component explicitly keeps the support occurrence $Q_n\subseteq Q_m$; the support has not been silently replaced by the whole of $Q_m$.

\begin{definition}[Witness Embedding]
The map $E_{Q_n,Q_m}$ above is called the \emph{witness embedding} along the concrete occurrence relation $Q_n\subseteq Q_m$. It preserves the original witness's support, endpoints, and actual trajectory. It does not require the image to become an elementary traverse of $Q_m$.
\end{definition}

\begin{plainbox}[What the operation actually does]
Embedding is modest on purpose: the context grows, while the old witness is left unchanged. It neither reinterprets the trajectory nor searches at a higher level for a ``similar'' traverse. It only says that this already specified finite substructure is now regarded as part of a larger context.
\end{plainbox}

Only when the surrounding discussion has fixed a unique occurrence chain do we abbreviate this as
\[
E_{n,m}(W).
\]
This abbreviation never asserts a unique map from an abstract level-$n$ tile to an abstract level-$m$ tile without a choice of copy.

\begin{lemma}[Embedding Coherence]
Let
\[
Q_n\subseteq Q_m\subseteq Q_\ell
\]
be a nested occurrence chain. For every finite witness $W$ supported in $Q_n$,
\[
E_{Q_m,Q_\ell}\bigl(E_{Q_n,Q_m}(W)\bigr)
=E_{Q_n,Q_\ell}(W),
\]
where equality means two bookkeeping descriptions of the same concrete subgraph occurrence, the same endpoints, and the same trajectory.
\end{lemma}

\begin{proof}
This is associativity of inclusion. The composite $Q_n\hookrightarrow Q_m\hookrightarrow Q_\ell$ equals the direct inclusion $Q_n\hookrightarrow Q_\ell$. Hence the images of the support occurrence, both endpoints, and the trajectory are identical.
\end{proof}

\begin{plainbox}[Do not read more into the lemma]
Embedding coherence does not prove that every pair of isomorphic tile types comes with a canonical map. It says only that once a concrete nested chain has been chosen, enlarging along that chain in two steps and in one step gives the same occurrence.
\end{plainbox}

\section{Decomposition: Not Moving a Small Object Upward, but Cutting a High-Level Trajectory Downward}

The easiest mistake at this point is to treat decomposition as the reverse of embedding. They answer different questions.

Embedding asks:
\begin{quote}
``Where is this already existing finite witness after it is placed in a larger context?''
\end{quote}
Decomposition asks:
\begin{quote}
``Which lower-level traverses actually compose this high-level traverse when it is viewed downward?''
\end{quote}

Let
\[
\tau=(P_m;b_-,\gamma,b_+)
\]
be a concrete traverse of a level-$m$ tile $P_m$, and fix $n<m$. Consider all level-$n$ subtile occurrences inside $P_m$,
\[
\operatorname{Occ}_n(P_m).
\]
Along the temporal order of $\gamma$, take all maximal subtrajectories with the following property: for some
\[
Q\in\operatorname{Occ}_n(P_m),
\]
the subtrajectory lies entirely in $Q$, its two endpoints belong to $\partial Q$, and its internal part meets no other point of $\partial Q$. These maximal pieces are exactly the level-$n$ traverse occurrences actually induced by $\gamma$. List them in the order in which they occur:
\[
\widetilde\tau_1,\ldots,\widetilde\tau_k.
\]

\begin{definition}[Induced Traverse Sequence of a Concrete Occurrence]
Define
\[
D^{\mathrm{occ}}_{m,n}(\tau)
=(\widetilde\tau_1,\ldots,\widetilde\tau_k).
\]
It records the level-$n$ traverse occurrences actually traversed and completed by the high-level trajectory, retaining the concrete subtile copy in which each occurrence lies.
\end{definition}

\begin{plainbox}[A psychology reader may think ``resegmentation,'' but only as navigation]
This is not interpretive freedom to divide a large event into any desired number of pieces. Mathematically, the cutting is determined by the same actual trajectory, a fixed lower-level subtile system, and a fixed boundary condition. Precisely because the cutting rule is fixed in advance, the coherence theorem below has content.
\end{plainbox}

This is the same source-side picture as the mechanism in Kuang's growth proof by which a deeper-level traverse induces a finite sequence of earlier-level traverses \citep{kuang2026a}. But it must be kept distinct from the last-moment map. The last-moment map is a selection map applied \emph{after} an induced sequence exists: under additional boundary-truncation conditions it selects one distinguished last-induced traverse. It is not the decomposition itself.

\section{Why the Induced Traverses Alone Are Still Not Enough: How One Passes between Them May Matter}

The sequence
\[
D^{\mathrm{occ}}_{m,n}(\tau)
=(\widetilde\tau_1,\ldots,\widetilde\tau_k)
\]
is enough for many traverse-counting problems, but it omits one thing: how does the actual trajectory pass from one subtile occurrence to the next between successive induced traverses?

If only coarse counting is needed, these intermediate pieces can be forgotten. If germ type, residue, or exact composition is to be studied, the connector / boundary data in the intermediate pieces may determine the result and cannot be quotiented out prematurely.

Therefore retain the complementary subtrajectories of $\gamma$. In temporal order, write uniquely
\[
\gamma
=c_0*\gamma_1*c_1*\gamma_2*\cdots*\gamma_k*c_k,
\]
where $\gamma_i$ is the trajectory of $\widetilde\tau_i$, while $c_i$ is the actual complementary connector / boundary block between neighboring induced traverses. The endpoint blocks $c_0$ and $c_k$ are allowed to be empty.

\begin{definition}[Full Structural Decomposition]
Define
\[
D^{\mathrm{str}}_{m,n}(\tau)
=(c_0,\widetilde\tau_1,c_1,\ldots,\widetilde\tau_k,c_k).
\]
It records the complete ordered cutting of the actual trajectory $\gamma$ under the level-$n$ subtile decomposition.
\end{definition}

After forgetting the complementary blocks and then forgetting occurrence locations, one obtains the traverse-type word projection
\[
D^{\mathrm{tr}}_{m,n}(\tau)
:=
\operatorname{Type}\circ\operatorname{Forget}_c
\bigl(D^{\mathrm{str}}_{m,n}(\tau)\bigr).
\]

\begin{plainbox}[Three information levels]
\[
\text{actual structural cutting}
\Longrightarrow
\text{occurrence sequence}
\Longrightarrow
\text{type word}.
\]
Every step to the right forgets data on purpose. A counting problem may need only the rightmost object; connector-sensitive residue requires returning to the left. Formalization does not mean retaining the maximum possible information forever. It means knowing exactly at which step a piece of information was forgotten.
\end{plainbox}

\begin{nonexample}[Transport from Strong Repetitivity Is Not Decomposition]
Strong repetitivity may let us find an isomorphic finite subwalk in another tile copy, for counting or for constructing a return witness. Moving a trajectory to another copy is transport. Cutting the original actual trajectory into the subtile pieces it already traverses is decomposition. The two operations must not be written as the same map.
\end{nonexample}

\section{Decomposition Coherence: Why ``Cut Coarsely First, Then Refine'' Should Agree with ``Cut Finely at Once''}

This is the most technical section of the chapter, but the question itself is simple. If coarse and fine resolutions both claim to cut the same trajectory, passing through an intermediate resolution should not manufacture a different order of occurrences.

Before stating the theorem, perform a type check. An expression such as
\[
D^{\mathrm{str}}_{m,n}\bigl(D^{\mathrm{str}}_{\ell,m}(\tau)\bigr)
\]
is not type-correct, because the output of $D^{\mathrm{str}}_{\ell,m}(\tau)$ is already an alternating structural word, not a single level-$m$ traverse. A segmentation list cannot pretend to be one object and simply be fed back into the same decomposition map.

The correct operation is refinement of a structural cutting. Let $n<m<\ell$. First perform the structural cutting of $\tau$ at level $m$. Then keep the same actual trajectory fixed, further subdivide every level-$m$ traverse block and complementary block according to level-$n$ subtile occurrences, and merge adjacent empty blocks. Denote this refinement operation by
\[
\operatorname{Ref}_{m\to n}.
\]

\begin{theorem}[Decomposition Refinement Coherence]
Fix the tile-inflation presentation and a concrete high-level traverse occurrence $\tau$. If $n<m<\ell$, then
\[
D^{\mathrm{str}}_{\ell,n}(\tau)
=
\operatorname{Ref}_{m\to n}
\bigl(D^{\mathrm{str}}_{\ell,m}(\tau)\bigr).
\]
That is, directly cutting the same actual trajectory at level $n$ and first cutting it at level $m$ and then refining that cutting along the same trajectory to level $n$ produce the same ordered structural decomposition.
\end{theorem}

\begin{proof}
Level-$n$ subtile occurrences give a nested refinement of level-$m$ subtile occurrences. Neither procedure moves the trajectory or replaces an occurrence. Both procedures merely mark, on the same finite edge sequence $\gamma$, the cutting times determined by level-$n$ boundary sets and subtile supports.

Marking level-$m$ cuts first and then adding all level-$n$ cuts inside the resulting blocks gives the same temporally ordered cut set as marking all level-$n$ cuts at once. The complementary connector / boundary blocks are the same remaining subtrajectories of the same $\gamma$, so the structural words agree.
\end{proof}

\begin{corollary}[Coherence of Projected Traverse Words]
Under the hypotheses above, if each level uses the same occurrence-to-type convention, then
\[
D^{\mathrm{tr}}_{\ell,n}(\tau)
=
\operatorname{Proj}^{\mathrm{tr}}
\operatorname{Ref}_{m\to n}
\bigl(D^{\mathrm{str}}_{\ell,m}(\tau)\bigr).
\]
Thus the intermediate level used on the way down to level $n$ does not alter the final order of lower-level induced traverses. Differences arise only from occurrence / connector decorations that we later choose to forget.
\end{corollary}

\begin{example}[Part-III pointer, not a mathematical conclusion: a formal support for legitimate untelescoping]
Part III required that if coarse and fine narratives both claim to arise from the same underlying occurrences, they should be able to state how they correspond. The theorem supplies only a mathematical benchmark: for a fixed actual trajectory in a nested occurrence presentation, legitimate refinement preserves the coherence of cutting order. Whether real human memory obeys anything analogous is a separate empirical ledger.
\end{example}

\section{Why Embedding and Decomposition Are Never Inverses in General}

The two operations can now be placed side by side:
\begin{center}
\resizebox{0.96\linewidth}{!}{$\displaystyle
\begin{array}{rcl}
E:&\text{supported finite witness}&\longrightarrow
\text{the same witness inside a larger chosen occurrence context},\\[1mm]
D:&\text{high-level traverse occurrence}&\longrightarrow
\text{ordered lower-level cutting of the same actual trajectory}.
\end{array}
$}
\end{center}

They not only point in opposite directions; their input and output objects are different:
\begin{itemize}
\item $E$ does not resegment the trajectory; it preserves an already existing finite witness;
\item $D$ does not move a lower-level witness to a higher level; it analyzes the internal composition of a high-level elementary traverse;
\item $E$ requires a specified nested occurrence relation;
\item $D$ requires a specified high-level occurrence and lower-level subtile resolution.
\end{itemize}
Thus in general there is no identity
\[
D\circ E=\Id
\qquad\text{or}\qquad
E\circ D=\Id.
\]

\begin{plainbox}[Why psychology readers should remember the distinction]
``Put the same old event inside a larger context'' and ``split one large event into finer episodes'' are both easily called ``understanding it again'' in ordinary speech. Formal language forces them apart. The first changes context; the second changes segmentation resolution. Part IV proves only that the mathematical operations are different. Whether psychology needs the same separate ledgers was already raised as a Level-II question in Part III.
\end{plainbox}

\section{Formal Disciplines Frozen by This Chapter}

All later chapters follow four rules:
\begin{enumerate}
\item \textbf{Occurrence first.} For embedding, persistence, connector data, or residue, fix the concrete occurrence before speaking about type.
\item \textbf{Embedding is copy-relative.} Strict witness embedding exists only along a specified nested occurrence chain. There is no unconditionally unique abstract map $P_n\hookrightarrow P_m$.
\item \textbf{Decomposition follows the actual trajectory.} Induced traverses arise from real cutting of the same high-level trajectory through lower-level subtile occurrences; they are not produced by arbitrary transport using strong repetitivity.
\item \textbf{Coherence is refinement coherence.} What composes legitimately is refinement of nested cuttings, not the ill-typed operation of treating a structural word as one traverse and feeding it back into $D_{m,n}$.
\end{enumerate}

\begin{plainbox}[What enters the next chapter]
If only one distinction is retained, retain this one: persistence tracks whether \emph{the same finite witness is still there}; decomposition analyzes \emph{what composes a high-level traverse}. The next chapter uses only the first operation to study persistence of return witnesses. The chapter after that turns to return language and factor avoidance. The two lines are separated from this point onward.
\end{plainbox}

\chapter{Return Witnesses, Birth Levels, and Persistence: How a Return-Completion Is Tracked}
\label{chap:return-persistence-en}

\par\medskip
\noindent\textbf{Why this chapter matters}\par
The preceding chapter separated two vertical operations. Embedding places the \emph{same finite structure} inside a larger context; decomposition cuts a larger trajectory downward into smaller pieces. Here we follow only the first line and ask one concrete question: once a return has an explicit finite witness, how do we continue to recognize at higher levels that ``this is still the same return-completion''?

Three misunderstandings are especially easy here. First, a source word and the occurrence witnessing its return-completion are not the same object. Second, ``first visible at level $n$'' is a resolution coordinate in a presentation, not ``first happened at that time'' in reality. Third, mathematical persistence says only that a finite witness remains traceable along a specified embedding. It does not say that human memory cannot forget or reconstruct. The task of the chapter is to keep all three ledgers separate.
\par\medskip

\section{Begin with Three Cards: Word, Occurrence, and Return-Completion Witness}

Psychology readers can understand this section first as an ordinary discipline of evidence. If we say that ``this process forms a return,'' at least three different questions are involved:
\begin{center}
\resizebox{0.96\linewidth}{!}{$\displaystyle
\begin{array}{rcl}
\text{word data} &:& \text{what are the complete source-word labels?}\\
\text{occurrence data} &:& \text{in which boundary context is it realized?}\\
\text{return-completion witness} &:& \text{which internal excursion proves the same-endpoint condition?}
\end{array}
$}
\end{center}
If these three are compressed into one symbol too early, birth, persistence, and residue will later begin to switch object levels silently.

Retain the source-side notation frozen in Part II. Let $S$ be the fixed generator / local-action alphabet and let
\[
\mathcal R\subseteq S^*
\]
be the set of admissible return words / factors. Part IV does not redefine $\mathcal R$.

\begin{definition}[Witnessed Return Datum]
A \emph{witnessed return datum} supported at level $n$ is written
\[
R=(u;W_R),
\qquad
u\in\mathcal R,
\qquad
W_R=(Q;b,\gamma,b),
\]
where $u$ retains the complete spelling of the Part-II source return word, including the boundary entry and exit labels; $Q$ is a concrete level-$n$ tile occurrence; and $\gamma$ is the finite trajectory induced by the internal closed excursion of $u$ in that admissible occurrence.

Thus
\[
W_R=(Q;b,\gamma,b)
\]
records only the finite internal witness actually responsible for return-completion: both distinguished endpoint occurrences are the same occurrence $b$. The complete word $u$ and the internal witness $W_R$ belong to different object layers.
\end{definition}

Accordingly, $u$ answers ``what is the complete source word?''; $Q$ and $\gamma$ answer ``where is its internal closed excursion finitely witnessed on this occurrence?''; and $b=b$ answers ``why is this internal witness closed?'' The same literal factor $u$ may occur in different tile copies, connector contexts, or singular sectors and thereby produce different witnessed data.

\begin{plainbox}[One sentence to remember]
The word is the name; the witness is the evidence. One name can have many concrete pieces of evidence. If a concrete piece of evidence loses its occurrence support, later embedding and residue no longer have enough information to track.
\end{plainbox}

\begin{example}[Part-III pointer, not a mathematical conclusion: why ``again'' cannot rely on a similarity label alone]
Part III distinguished ``I am anxious again'' from ``I have returned to the same subject-relevant position.'' The mathematics here supplies only an evidential structure: literal content can be similar, but a return datum must still state how one concrete finite occurrence completes a return. The example helps explain why formal objects are layered; it does not identify a psychological episode with a tile.
\end{example}

\section{Support Level: First State ``Which Level Supports This Witness,'' Then Speak about Birth}

Earlier notation can easily mix ``the level at which this particular witness is currently recorded'' with ``the earliest level at which this return can be witnessed.'' We define the first quantity before the second.

\begin{definition}[Support Level]
Let
\[
R=(u;W_R),
\qquad
W_R=(Q;b,\gamma,b),
\]
where $Q$ is a concrete level-$n$ tile occurrence. We call
\[
\ell(R)=n
\]
the \emph{support level} of the witnessed datum.
\end{definition}

The number $\ell(R)$ says only at which finite scale this particular concrete witness is presently recorded. It has no ``earliest'' meaning. If the same source word $u$ has concrete witnesses at levels $3$ and $7$, then those data have support levels $3$ and $7$. The latter is not a ``new return'' merely because it was recorded at a later level.

\begin{plainbox}[Especially important for psychology readers]
A level here is neither clock time nor a developmental stage of life. It is closer to a finite resolution coordinate in the current mathematical presentation. Therefore
\begin{center}
\resizebox{0.96\linewidth}{!}{$\displaystyle
\text{support level}
\neq
\text{time of real occurrence}
\neq
\text{time at which the subject first became aware of it}.
$}
\end{center}
A person who says only on the third similar interaction, ``so this is a pattern I repeat,'' has not thereby shown that the pattern was born on the third interaction. Likewise, first seeing a mathematical witness clearly at one truncation does not mean that its trajectory began to exist only at that level.
\end{plainbox}

\section{Birth Level: Earliest Visibility Is Not Absolute Origin}

Within a fixed presentation we may ask a further question. For an admissible occurrence of a source return word, at which finite level can its internal closed excursion first obtain a witness? This is the only sense in which the chapter uses ``birth.''

\begin{definition}[Word Birth Level]
For $u\in\mathcal R$, whenever a finite witnessed realization exists, define
\begin{center}
\resizebox{0.96\linewidth}{!}{$\displaystyle
b_{\mathrm{word}}(u)
:=
\min\bigl\{
\ell(R):R=(u;W_R)\text{ is an admissible witnessed return datum}
\bigr\}.
$}
\end{center}
This is the \emph{word birth level} of $u$ in the current presentation. The level locates the finite tile scale supporting its internal closed witness; it does not compress the complete source word into a tile-internal path.
\end{definition}

In the finite-witness setup of Part IV, this notation is used only for return words possessing such a finite realization. The integer depends on the presentation. Bratteli telescoping, level renumbering, generator choice, or the convention for finite realization may change it. Hence
\[
\boxed{b_{\mathrm{word}}(u)\text{ is a presentation coordinate, not an intrinsic invariant}.}
\]

If descendants of one particular witness later need to be tracked, the more accurate phrase is ``the level at which this supported witness lineage begins to be recorded,'' not a relabeling of every descendant's support level as a new birth.

\section{Persistence: Enlarging the Context Does Not Turn the Same Closed Witness into an Open One}

We can now state the main structural theorem of the chapter. Suppose a closed witness has already been fixed, and we have explicitly specified which lower-level tile occurrence it follows into a higher-level context. The question is simply: after this finite evidence card is placed inside a larger folder, do its equal endpoints remain equal?

\begin{theorem}[Persistence of a Closed Witness under Supported Inflation]
Let
\[
R=(u;W_R),
\qquad
W_R=(Q_n;b,\gamma,b)
\]
be a witnessed return datum with support level $n$, and fix a concrete nested occurrence relation
\[
Q_n\subseteq Q_m.
\]
Then the occurrence-supported witness embedding of Chapter 24 gives
\[
E_{Q_n,Q_m}(W_R),
\]
which is still a closed finite witness supported on the same embedded copy $Q_n\subseteq Q_m$. In other words, merely enlarging the ambient tile context does not turn this already existing return-completion into an open witness.
\end{theorem}

\begin{proof}
Embedding does not re-cut the trajectory or reselect the endpoints. It places the same concrete support, the same $\gamma$, and the same endpoint occurrence $b$ inside a larger tile. Both endpoints of the original witness are $b$, and after inclusion they still land on the same vertex. Hence the same-endpoint property is preserved.

One point requires emphasis: $E_{Q_n,Q_m}(W_R)$ need not be an elementary return traverse of $Q_m$. A position that was a boundary point of $Q_n$ can become merely an internal distinguished point or a connector-adjacent position in $Q_m$. The theorem preserves a finite closed subwitness, not elementary-traverse status.
\end{proof}

\begin{warning}[The Domain of Persistence Is a Supported Occurrence]
Without first specifying the concrete occurrence relation
\[
Q_n\subseteq Q_m,
\]
one cannot infer from equality of tile type an unconditional abstract map $P_n\hookrightarrow P_m$. A higher-level tile may contain several isomorphic lower-level copies. Persistence must say \emph{which copy} is being preserved.
\end{warning}

\begin{plainbox}[What ``a return witness has birth but no death under inflation'' actually means]
The precise claim is only this: once a closed finite witness and a supported descendant chain have been fixed, that subwitness itself is not erased by later enlargement of context.

It does \emph{not} mean:
\begin{itemize}
\item the witness remains an elementary return in every higher tile;
\item all presentations assign it the same birth level;
\item once human memory forms, it can never forget;
\item later psychological interpretation cannot reject a false memory or resegment an episode.
\end{itemize}
What mathematics preserves is an explicit finite object, not a mechanism of subjective memory.
\end{plainbox}

\section{Return Filtration: Accumulating ``What Can Already Be Witnessed up to This Level''}

Once word birth level is available, we can define a genuinely type-correct cumulative language.

\begin{definition}[Return Language Visible up to Level $n$]
Define
\[
\mathcal R_{\le n}
=
\{u\in\mathcal R:b_{\mathrm{word}}(u)\le n\}.
\]
Equivalently, $u\in\mathcal R_{\le n}$ if and only if there exists an admissible witnessed datum $R=(u;W_R)$ satisfying $\ell(R)\le n$.
\end{definition}

Thus, in the present finite-realizable setup,
\[
\mathcal R_{\le1}\subseteq\mathcal R_{\le2}\subseteq\cdots,
\qquad
\mathcal R=\bigcup_{n\ge1}\mathcal R_{\le n}.
\]
Two different kinds of monotonicity must be kept separate:
\begin{center}
\resizebox{0.96\linewidth}{!}{$\displaystyle
\begin{array}{ll}
\mathcal R_{\le n}\text{ grows:} & \text{this follows from the cumulative definition ``up to level }n\text{'';}\\
W_R\text{ remains traceable in larger context:} & \text{this follows from the persistence theorem.}
\end{array}
$}
\end{center}
The first is bookkeeping. The second is structural preservation. Packaging the first as a deep theorem would exaggerate the mathematical content of the filtration itself.

\begin{example}[Part-III pointer, not a mathematical conclusion: discovering more does not mean memory can only increase]
In psychological research, new context or new variables may lead a person to recognize a previously unnoticed pattern, but can also lead to withdrawal of an earlier false positive. That learning trajectory is generally not monotone. The present $\mathcal R_{\le n}$ describes only the mathematical axis ``the permitted finite witness level increases'' inside a fixed formal presentation. It cannot be used directly as an autobiographical learning curve.
\end{example}

\section{A Large Return Can Be Composed of Many Open Pieces}

Persistence asks how ``the same small witness is carried upward.'' This section records another fact that is easy to confuse with it: when a higher-level return is decomposed downward, not every lower-level piece need itself be closed.

Let $R^{(m)}$ be a level-$m$ elementary return occurrence. Using the actual-trajectory cutting of Chapter 24, at a level $n<m$ obtain
\[
D^{\mathrm{occ}}_{m,n}(R^{(m)})
=(\widetilde\tau_1,\ldots,\widetilde\tau_k).
\]
Every $\widetilde\tau_i$ may be an open traverse. Return-completion occurs at the composite scale. When the induced pieces and the complementary connector blocks are reassembled in the original trajectory order, only the total start and total endpoint return to the same higher-level endpoint.

\begin{proposition}[Composite Return-Completion under Compatible Cutting]
If $R^{(m)}$ is a higher-level return, then its occurrence-level lower-scale cutting can be concatenated in actual trajectory order. Reassembling all induced traverse pieces and complementary connector blocks recovers the original closed trajectory. No individual $\widetilde\tau_i$ is required to be a return.
\end{proposition}

\begin{proof}
This is a direct consequence of the full structural decomposition of Chapter 24. Cutting merely inserts lower-level boundary cut points into the original trajectory. It neither replaces the trajectory nor changes its total start and endpoint. Hence the full concatenation remains closed.
\end{proof}

\begin{plainbox}[Why this matters to psychology readers]
Part III discussed how a ``large event'' can be made from smaller episodes. The only mathematical discipline supplied here is that global return-completion does not require every local piece to complete a return independently. Therefore, when a macroscopic return is observed, one must not mechanically search every microscopic segment for a smaller return of the same shape. The scale can change, and the return-completion relation can change with it.
\end{plainbox}

\section{Three Scale Operations Kept Separate Again}

The scale language of this and the preceding chapter can now be compressed into three operations:

\begin{center}
\small
\begin{tabularx}{0.96\textwidth}{>{\bfseries\raggedright\arraybackslash}p{0.22\textwidth}Y}
\toprule
Operation & What it actually does \\
\midrule
Inflate & Along a specified occurrence inclusion, place the same finite witness inside a larger context.\\
Telescope & Legitimately omit intermediate Bratteli levels and change the scale coordinates of the presentation.\\
Untelescope & Restore intermediate refinement skipped inside the same nested inflation system.\\
\bottomrule
\end{tabularx}
\end{center}

Resolve does not appear in this table because it belongs to the graph-of-germs resolution axis, not the Bratteli / tile scale axis.

\begin{plainbox}[Four sentences to carry forward]
\begin{enumerate}
\item $u$ is the source word; $W_R$ is the concrete return-completion witness.
\item $\ell(R)$ is the support level of this datum; $b_{\mathrm{word}}(u)$ records the earliest finite level at which $u$ is visible in the current presentation.
\item Persistence preserves the same closed subwitness along a specified occurrence chain. It does not say that ``memory never changes.''
\item The increase of $\mathcal R_{\le n}$ is cumulative bookkeeping; the real structural content lies in witness preservation.
\end{enumerate}
The next chapter stops following the fate of one witness at higher levels. It gathers all admissible return words into a forbidden-factor language and uses that language to define incompressibility strictly.
\end{plainbox}

\chapter{Return Language and Incompressibility: How a Word Avoids Every Return}
\label{chap:return-language-en}

\par\medskip
\noindent\textbf{Why this chapter matters}\par
The preceding chapter kept following one concrete return witness as it remained present at larger scales. We now ask a different question. If we collect all finite words permitted to count as returns, when does a longer word avoid every return completely?

Here ``language'' does not mean natural language or the way a person tells a story about the self. It has the elementary formal-language meaning: a set of finite words. Once return words are collected into a set $\mathcal R$, incompressibility becomes a precise forbidden-factor condition: no contiguous subword of the given word belongs to $\mathcal R$.

For psychology readers, retain only one sentence for now. The question is not ``this person is difficult to understand.'' It is: ``inside the declared return dictionary, this coded trajectory contains no return factor.''
\par\medskip

\section{Specify the Universe First: Admissible Words, Factors, and the Return Language}

Part IV has fixed the generator / local-action alphabet $S$. Not every formal string $s_1\cdots s_N\in S^*$ need be realizable in the present local-action system. Denote the actually realizable finite words by
\[
\mathcal L_{\mathrm{adm}}\subseteq S^*.
\]
Thus $\mathcal L_{\mathrm{adm}}$ is the admissible source language of the current presentation: every $w\in\mathcal L_{\mathrm{adm}}$ has at least one legitimate finite realization. Since a contiguous segment of an actual trajectory is again an actual trajectory, we assume $\mathcal L_{\mathrm{adm}}$ is closed under contiguous factors.

For
\[
w=s_1\cdots s_N\in\mathcal L_{\mathrm{adm}},
\]
define its nonempty contiguous factors by
\[
\operatorname{Fac}(w)
=
\{s_i\cdots s_j:1\le i\le j\le N\}.
\]
A factor is therefore a contiguous block cut from the original trajectory, not an arbitrary selection of letters.

As in the preceding chapter,
\[
\mathcal R\subseteq\mathcal L_{\mathrm{adm}}
\]
always denotes the current source-side admissible return language. Thus $u\in\mathcal R$ only when it can acquire a concrete closed witness in the sense of Chapter 25.

\begin{plainbox}[For psychology readers: think of a finite pattern dictionary]
The set $\mathcal R$ is like a declared ``return dictionary.'' Given a longer word $w$, the chapter no longer asks whether the whole word vaguely resembles a pattern. It inspects the word block by block:
\[
\text{does some }u\in\operatorname{Fac}(w)\text{ already belong to }\mathcal R?
\]
If yes, a finite return factor has been found. Only if the answer is no does incompressibility arise. This is much stricter than saying that the whole process ``feels repetitive.''
\end{plainbox}

\section{Return Classes and Intervals: One Long Word Can Hit Several Kinds of Return at Once}

Sometimes the return language must be classified according to a predeclared witness / decorated criterion. Use a class index
\[
\theta\in\Theta,
\qquad
\mathcal R_\theta\subseteq\mathcal R.
\]
The symbol $\theta$ serves only as a class index; $\rho$ is reserved in Part IV for germ residue in the next chapter.

If the family $\{\mathcal R_\theta\}_{\theta\in\Theta}$ is declared complete, then
\[
\mathcal R=\bigcup_{\theta\in\Theta}\mathcal R_\theta.
\]
Without such a completeness assumption, the classes are only local sublanguages and cannot replace the full $\mathcal R$.

\section{Multi-Boundary Bridge-Type Coding: A Return Class Need Not Come from One Site}

A source-side structure already present in the literature must now be incorporated formally into Part IV. Subsection 5.2 of Kuang's preprint v1 treats multiple types of bridges \citep{kuang2024v1}. For every bridge type $x$, there is a traverse family $\Theta_{n,x}$, with
\[
\Theta_n=\bigcup_x\Theta_{n,x}.
\]
A higher-level traverse $\tau$ induces
\[
\Phi(\tau)=(\tau_1,\ldots,\tau_k),
\]
and the type map produces a finite word
\[
\operatorname{typ}\Phi(\tau)=x_1\cdots x_k.
\]
The $x_i$ may correspond to different infinite boundary points. Thus the type word itself records how one finite trajectory passes successively through several boundary / bridge sectors.

This gives a more accurate picture than simply writing the return language as a union of separate $\mathcal R_x$. A return factor may depend on a contiguous block of a multi-sector type word, for example
\[
x_i x_{i+1}\cdots x_j,
\]
rather than being witnessed independently by one $x$. Sector data are therefore more naturally treated as internal coding of return words, not merely as an external class label.

\begin{definition}[Sector-Coded Return Datum: Bookkeeping Form]
In a fixed multi-bridge presentation, let $R=(u;W_R)$ be a witnessed return datum whose relevant induced-traverse decomposition carries the bridge-type word
\[
x(R)=x_i\cdots x_j.
\]
We call
\[
(R;x(R))
\]
a \emph{sector-coded return datum}. The word $x(R)$ records only source-side boundary / bridge-type passage. It is not the germ residue of the next chapter, and distinct sectors are not required to belong to distinct groupoid orbits.
\end{definition}

\begin{plainbox}[Why a union of single-site dictionaries can lose information]
If a return forms by passing successively through bridge types $x,y,z$, forcing it into ``a return belonging to $x$'' or ``a return belonging to $y$'' loses the order of passage. The more faithful object is the type word $xyz$ together with its witnessed factor. Multiple boundary points first produce a path grammar, not merely many local dictionaries.
\end{plainbox}

Proposition 5.7 of the cited Subsection 5.2 assumes a uniform $d$ such that every sufficiently long reduced word contains a return and then derives the growth estimate by multi-type last-moment counting. It does not derive that $d$ from multi-boundary structure alone. Example 5.8 obtains $d=12$ in a concrete example with three boundary sectors $1^\omega,5^\omega,6^\omega$ corresponding to $a,b,c$, by finite orbital-graph tracing \citep{kuang2024v1}.

Part IV therefore contains two distinct questions:
\begin{center}
\resizebox{0.96\linewidth}{!}{$\displaystyle
\begin{array}{ll}
\text{multi-boundary coding:} & \text{how is a return encoded across several sectors?}\\
\text{pNH forcing:} & \text{in a special single-site setting, when can a uniform return bound be proved?}
\end{array}
$}
\end{center}
They must remain separate.

\begin{definition}[Return Intervals Relative to a Declared Classification]
For $w=s_1\cdots s_N\in\mathcal L_{\mathrm{adm}}$, define
\[
\operatorname{Ret}_\Theta(w)
=
\{(i,j,\theta):1\le i\le j\le N,\ w[i,j]\in\mathcal R_\theta\},
\]
where $w[i,j]=s_i\cdots s_j$. This set records all return intervals in $w$ recognized by the current class family.
\end{definition}

There is no reason for different return intervals to be disjoint. They may overlap or nest. The same factor may satisfy several classification criteria and hence correspond to several values of $\theta$. Thus ``this word has a return'' usually does not mean that it contains one unique cycle.

\begin{principle}[How Boundary Sectors Can Generate Nested Returns]
Multi-boundary / multi-germ coding gives a concrete source of nesting. Suppose
\[
(i,j,\theta_1),(k,\ell,\theta_2)\in\operatorname{Ret}_\Theta(w).
\]
Then
\[
i<k\le\ell<j
\]
means that one return interval is strictly nested inside the other, while
\[
i<k\le j<\ell
\]
means that the two returns partially overlap. Distinct singular boundary germs / bridge sectors often naturally generate heterogeneous returns with $\theta_1\neq\theta_2$, and a sector-passage word can allow a larger return to cross several local return families.

There is no converse. Even when $\theta_1=\theta_2$, the same germ / type may produce nested intervals at different scales or different occurrence positions. Multi-boundary geometry is therefore an important source prototype for nested-return structure, not a necessary condition for it.
\end{principle}

\begin{plainbox}[Why overlap must be allowed]
Part III already showed, in the discussion of rumination, that the same trajectory can be hit by several different ``again'' structures. In source-side terms, when a system has several boundary / germ sectors one should not expect different return families automatically to partition a word into disjoint blocks. A large return can contain a smaller return witnessed by another sector, and two returns can share only a middle portion. Global incompressibility must therefore inspect the complete return language rather than only the most salient class.
\end{plainbox}

\section{Relative and Global Incompressibility: How Much of the Return Language Is Being Avoided?}

\begin{definition}[Relative Incompressibility]
Relative to a declared sublanguage $\mathcal R_\theta$, a word $w$ is \emph{relatively incompressible with respect to class $\theta$} if
\[
\operatorname{Fac}(w)\cap\mathcal R_\theta=\varnothing.
\]
\end{definition}

This is only a local conclusion. It says that $w$ hits no return of this class; it may still contain a return factor belonging to another class.

\begin{definition}[Global Incompressibility]
We say that $w\in\mathcal L_{\mathrm{adm}}$ is \emph{globally incompressible in the current presentation} if
\[
\operatorname{Fac}(w)\cap\mathcal R=\varnothing.
\]
We write
\[
w\in\mathrm{IC}.
\]
\end{definition}

This is exactly the language form of the source-side definition: an admissible expression is incompressible if it contains no return subword \citep{kuang2026a}. Here ``global'' means only relative to the currently declared complete $\mathcal R$. It does not mean an absolute invariant across generators, tile presentations, or normal-form conventions.

\begin{warning}[This Chapter Is Primarily a Word-Level Statement]
In this chapter, $\mathrm{IC}$ is a language of admissible words / expressions. If incompressibility is later to be descended to abstract group elements, one must separately specify a representative convention or normal form, or prove insensitivity to the relevant choice of representative. No such presentation-independence is silently assumed here.
\end{warning}

\begin{nonexample}[Non-Return of a Single Germ Does Not Imply Global Incompressibility]
Even if the lifted continuation of one fixed germ type never completes a return, that excludes only one family of return-completions. It does not prevent another contiguous factor of $w$ from being witnessed as a return by another boundary / germ configuration. Therefore
\[
\text{germwise non-return}
\not\Rightarrow
\operatorname{Fac}(w)\cap\mathcal R=\varnothing.
\]
\end{nonexample}

\section{Compressibility Has a Finite Certificate}

Global incompressibility looks like a negative condition requiring inspection of all factors. Its negation, however, has a finite certificate.

\begin{proposition}[Finite Certificate of Compressibility]
For $w\in\mathcal L_{\mathrm{adm}}$, the following are equivalent:
\begin{enumerate}
\item $w\notin\mathrm{IC}$;
\item there exists a contiguous factor $u\in\operatorname{Fac}(w)$ together with a witnessed return datum
\[
R=(u;W_R)
\]
in the sense of Chapter 25, with $u\in\mathcal R$.
\end{enumerate}
\end{proposition}

\begin{proof}
By definition, $w\notin\mathrm{IC}$ is equivalent to
\[
\operatorname{Fac}(w)\cap\mathcal R\neq\varnothing.
\]
Choose $u$ from the intersection. By the standing convention on $\mathcal R$, $u$ has an admissible concrete closed witness, giving $R=(u;W_R)$. The converse is immediate.
\end{proof}

\begin{plainbox}[Why state this separately?]
``Incompressible'' is a global avoidance statement. ``Not incompressible'' does not require an explanation of the whole word. One finite return witness suffices. This is the point of the witness-first discipline of Part IV: compressibility leaves concrete, finite, checkable evidence.
\end{plainbox}

\section{Finite-Level Filtration: Returns Accumulate while Words That Have Not Hit Any Return Become Fewer}

The preceding chapter defined the cumulative return language
\[
\mathcal R_{\le n}
=
\{u\in\mathcal R:b_{\mathrm{word}}(u)\le n\}.
\]
Now define
\[
\mathrm{IC}_{\le n}
=
\{w\in\mathcal L_{\mathrm{adm}}:
\operatorname{Fac}(w)\cap\mathcal R_{\le n}=\varnothing\}.
\]

\begin{proposition}[Incompressibility Filtration]
In the present finite-realizable setup,
\[
\mathrm{IC}_{\le1}\supseteq\mathrm{IC}_{\le2}\supseteq\cdots,
\qquad
\mathrm{IC}=\bigcap_{n\ge1}\mathrm{IC}_{\le n}.
\]
\end{proposition}

\begin{proof}
Since
\[
\mathcal R_{\le n}\subseteq\mathcal R_{\le n+1},
\]
avoiding the larger forbidden family is the stronger condition. Hence
\[
\mathrm{IC}_{\le n+1}\subseteq\mathrm{IC}_{\le n}.
\]
Also, from Chapter 25,
\[
\mathcal R=\bigcup_{n\ge1}\mathcal R_{\le n}.
\]
A word avoids every $\mathcal R_{\le n}$ if and only if it avoids the complete $\mathcal R$, which gives the intersection formula.
\end{proof}

The two filtrations therefore point in opposite directions:
\[
\mathcal R_{\le n}\nearrow
\qquad\Longleftrightarrow\qquad
\mathrm{IC}_{\le n}\searrow.
\]
As the permitted witness level grows, the returns we cumulatively recognize can only increase, so the words that have ``still hit no return up to this level'' can only decrease. These are dual bookkeeping directions inside one fixed formal presentation.

\begin{warning}[Changing the Model Is Not ``Moving Forward along the Filtration'']
Changing the generator alphabet, tile presentation, episode coding, psychological observables, or the researcher's segmentation rule changes the admissible language or return dictionary. That is a change of formal model, not an increase of $n$ along one fixed $\mathrm{IC}_{\le n}$ filtration.
\end{warning}

\section{Factoriality: If the Whole Segment Avoids Every Return, Every Contiguous Subsegment Does Too}

\begin{proposition}[Factoriality of the Incompressible Language]
For every $n$, the language $\mathrm{IC}_{\le n}$ is closed under contiguous factors, and so is global $\mathrm{IC}$. Thus if
\[
w\in\mathrm{IC}_{\le n},
\qquad
v\in\operatorname{Fac}(w),
\]
then $v\in\mathrm{IC}_{\le n}$. The same statement holds globally.
\end{proposition}

\begin{proof}
If $v$ contained some $u\in\mathcal R_{\le n}$, then $u$ would also be a contiguous factor of $w$, contradicting $w\in\mathrm{IC}_{\le n}$. Replace $\mathcal R_{\le n}$ by $\mathcal R$ for the global case.
\end{proof}

This property is a direct formal consequence of a forbidden-factor definition. It requires no additional bounded-type geometry. It is worth recording because it shows that $\mathrm{IC}$ is genuinely a factorial language rather than a collection of unrelated exceptional words.

\begin{plainbox}[Do not psychologize factoriality]
Mathematically, ``a long word is incompressible $\Rightarrow$ every factor is incompressible'' is a hereditary property forced by the definition. It does not mean that every local part of a complicated real-life problem is equally complicated, or that cutting psychological material into shorter pieces must preserve the same explanatory difficulty.
\end{plainbox}

\section{Finite Incompressible Language and a Uniform Upper Bound on Waiting for a Return}

The pNH benchmark of the next chapter will contain a stronger conclusion: there exists a uniform $M$ such that every sufficiently long admissible word contains a return. Before entering that special theorem, isolate the pure language-theoretic content of the statement.

\begin{proposition}[Bounded Avoidance Is Equivalent to Uniform Return]
Assume the alphabet $S$ is finite. The following are equivalent:
\begin{enumerate}
\item $\mathrm{IC}$ is finite as a word language;
\item incompressible word lengths are uniformly bounded:
\[
\sup\{|w|:w\in\mathrm{IC}\}<\infty;
\]
\item there exists $M$ such that every $w\in\mathcal L_{\mathrm{adm}}$ with $|w|>M$ satisfies
\[
\operatorname{Fac}(w)\cap\mathcal R\neq\varnothing.
\]
Equivalently, every sufficiently long admissible word contains a return factor.
\end{enumerate}
\end{proposition}

\begin{proof}
$(1)\Rightarrow(2)$: the set of lengths of words in a finite set has a maximum.

$(2)\Rightarrow(3)$: choose $M$ strictly larger than every incompressible length. If $|w|>M$ and $w$ contained no return factor, then $w\in\mathrm{IC}$, a contradiction.

$(3)\Rightarrow(2)$ is immediate. Finally, because $S$ is finite, there are only finitely many words of length at most a fixed $M$, so $(2)\Rightarrow(1)$.
\end{proof}

\begin{plainbox}[Where the next chapter adds real content]
The proposition above is only language logic over a finite alphabet. It does not tell us \emph{why} such an $M$ should exist. The special unique-pNH theorem in the next chapters supplies a nontrivial geometric mechanism: under strong hypotheses, finitely controlled regular-limit quotients together with strong repetitivity force a uniform upper bound on waiting for a return.
\end{plainbox}

\section{Part-III Pointer: Unchunkability Borrows a Relation, Not the Mathematical Definition}

Part III borrowed factor avoidance to discuss ``model-relative temporary unchunkability.'' Part IV freezes the mathematical content as
\[
\mathrm{IC}
=
\{w\in\mathcal L_{\mathrm{adm}}:
\operatorname{Fac}(w)\cap\mathcal R=\varnothing\}.
\]
Intensity of suffering, difficulty exiting, checking frequency, duration of rumination, or ``temporarily not being able to explain something'' are not part of this definition.

Three reminders should be kept explicit:
\[
\begin{gathered}
\text{high repetition}\not\Rightarrow\text{incompressibility},\\
\text{hard to exit}\not\Rightarrow\text{incompressibility},\\
\text{absence of one return class}\not\Rightarrow\text{global incompressibility}.
\end{gathered}
\]
A highly repetitive trajectory often hits return factors very easily and may therefore be highly compressible.

\begin{plainbox}[Five things to carry forward]
\begin{enumerate}
\item ``Language'' here means a set of finite words; $\mathcal L_{\mathrm{adm}}$ consists of source words that can actually be realized.
\item $\mathcal R$ is the complete return dictionary; $\mathcal R_\theta$ are only declared subclasses.
\item Global incompressibility is complete forbidden-factor avoidance: $\operatorname{Fac}(w)\cap\mathcal R=\varnothing$.
\item As $\mathcal R_{\le n}$ increases, $\mathrm{IC}_{\le n}$ decreases, and $\mathrm{IC}$ is a factorial language.
\item Over a finite alphabet, finite $\mathrm{IC}$ is equivalent to a uniform $M$ such that every sufficiently long admissible word contains a return.
\end{enumerate}
The next question is therefore precise: what special AF-by-discrete / pNH geometry can actually produce such a uniform $M$?
\end{plainbox}

\chapter{Germ-Resolved Return Residue: After Return-Completion, Does the Lift Also Come Back?}
\label{chap:germ-resolved-residue-en}

\par\medskip
\noindent\textbf{Why this chapter matters}\par
The preceding chapter made return precise: if a finite trajectory comes back to the same relevant endpoint at a declared scale, it has a concrete return-completion witness. But ``having come back'' is not always the end of the structure.

At the finer resolution of the graph of germs, the lift of a trajectory that is already closed in the projected graph may travel from one fibre sheet to another. Two questions must therefore be separated: first, whether the projected trajectory is closed; second, whether that return-completion leaves a fibre displacement behind. This chapter calls the second quantity the \emph{return residue}. It is not a label automatically carried by a bare return. It is defined only after the singular sector, base lift, and germ/connector decoration have all been fixed.
\par\medskip

\section{Why Return-Completion Needs a Second Coordinate}

Ignore the terminology for a moment and look only at the schematic picture
\[
\begin{array}{ccc}
\widetilde b_e & \xrightarrow{\widetilde\gamma} & \widetilde b_h\\
\downarrow && \downarrow\\
b & \xrightarrow{\gamma} & b.
\end{array}
\]
The downstairs trajectory $\gamma$ has already returned to the same vertex $b$, so its return-completion is not in doubt. Its lift $\widetilde\gamma$, however, may travel from one sheet to another. Thus
\begin{center}
\resizebox{0.96\linewidth}{!}{$\displaystyle
\boxed{\text{return-completion in the projected graph}}
\qquad\text{and}\qquad
\boxed{\text{return-completion in the covering fibre}}
$}
\end{center}
are two different facts.

This is precisely the information that a graph of germs retains beyond an orbital graph. An orbital graph records only orbit position; the graph of germs also retains the different local germ data that may lie over the same orbit position. A return may therefore already be completed at orbital resolution while leaving a nontrivial displacement at germ resolution.

\begin{plainbox}[What the psychology reader needs from this distinction]
The mathematical role of this chapter is not to rename an emotion a ``germ.'' It trains one distinction only: ``I have come back to the same position again'' and ``after coming back, every distinction relevant to later continuation is also the same'' are not the same statement. Part III, Chapter 20 borrowed exactly this relational skeleton.
\end{plainbox}

\section{Fix the Fibre First: Residue Has a Narrower Domain than Return}

Fix a singular sector $\xi$. Write its stabilizer, neighborhood stabilizer, and group of germs as
\[
G_\xi,\qquad G_{(\xi)},\qquad H_\xi=G_\xi/G_{(\xi)}.
\]
In the graph-of-germs setup used here, there is a covering
\[
\pi_\xi:\widetilde\Gamma_\xi\longrightarrow\Gamma_\xi,
\]
whose fibre is controlled by the germ classes in $H_\xi$ \citep{nek2022,kuang2026a}. For a fixed fibre, the natural formulation is that it is an $H_\xi$-torsor. Intuitively, the sheets in the fibre have only relative positions; no sheet is intrinsically marked ``identity.'' Only after a starting lift $\widetilde b_e$ is chosen as an origin can the remaining sheets be labeled relative to it by elements of $H_\xi$.

\begin{definition}[$\xi$-Anchored Germ-Decorated Return Datum]
A $\xi$-anchored germ-decorated return datum consists of the following data:
\begin{enumerate}
\item a witnessed return datum
\[
R=(u;W_R),\qquad W_R=(Q;b,\gamma,b);
\]
\item a realization compatible with the $\xi$-sector, under which $\gamma$ projects to a based closed trajectory in $\Gamma_\xi$ at a vertex $b$;
\item a chosen starting lift $\widetilde b_e\in\pi_\xi^{-1}(b)$;
\item label/bisection, connector, and local germ decorations sufficient to determine each step of the lift.
\end{enumerate}
Without these additional data, this chapter recognizes $R$ as a return but leaves $\rho_\xi(R)$ undefined.
\end{definition}

\begin{plainbox}[Why these data cannot be omitted]
A bare return tells us only that a finite trajectory has completed a return. Residue asks more: from which fibre sheet did it start, along which local-action lift did it travel, and on which sheet did it end? If even the starting sheet and connector context have not been fixed, the question ``what is the residue?'' has no unique answer.
\end{plainbox}

\section{Residue: Record the Lifted Endpoint Displacement}

Lift the fixed decorated trajectory uniquely from $\widetilde b_e$. We now fix a convention for the rest of the chapter. The fibre is a right $H_\xi$-torsor, with deck action written $\widetilde b\cdot k$. Cayley edges arise from left multiplication by local arrows and therefore commute with the right deck action. If the lifted endpoint is
\[
\widetilde b_e\cdot h^{-1},
\]
we define the residue to be $h$. We insert the inverse on purpose so that, under the fixed convention below, residues multiply in the order in which the corresponding based returns are traversed.

\begin{definition}[Return Residue]
We define the \emph{return residue} by
\[
\rho_\xi(R)=h\in H_\xi
\]
when the lifted endpoint satisfies
\[
\operatorname{end}(\widetilde\gamma)=\widetilde b_e\cdot h^{-1}.
\]
The complete germ-resolved return datum is written
\[
(R,\rho_\xi(R)).
\]
We call the anchored return \emph{fibre-trivial} if $\rho_\xi(R)=e$, and \emph{fibre-nontrivial} if $\rho_\xi(R)\neq e$.
\end{definition}

Thus a return now has two independent ledger entries:
\[
\begin{aligned}
\text{return-completion coordinate:}&\quad b\xrightarrow{\gamma}b,\\
\text{resolution coordinate:}&\quad \rho_\xi(R)\in H_\xi.
\end{aligned}
\]
The first asks whether the trajectory comes back. The second asks what net displacement remains after that return is resolved at the finer germ level.

\begin{proposition}[Well-Definedness after Fixing the Anchor]
In a well-labeled, uniquely liftable local-action setup, once the concrete decorated realization above and the starting lift $\widetilde b_e$ are fixed, the lifted endpoint is unique. Hence $\rho_\xi(R)$ is well-defined.
\end{proposition}

\begin{proof}
At each step of the trajectory, the current lifted germ vertex together with the chosen label/bisection uniquely determines the next lifted vertex. Iterating finitely many times uniquely determines the entire lifted walk and its endpoint fibre position. Relative to the fixed origin $\widetilde b_e$, the corresponding $H_\xi$-label is therefore unique.
\end{proof}

\section{Changing the Starting Lift: The Exact Value Is Generally Defined Only up to Conjugacy}

A simple fact can easily be hidden by notation. We chose $\widetilde b_e$ by hand as the reference sheet. That choice is part of the anchor data. If the fibre origin is changed, the residue of the same projected loop need not remain literally the same group element.

\begin{proposition}[Change of Fibre Anchor]
Fix the same decorated projected return and put
\[
\widetilde b_k:=\widetilde b_e\cdot k.
\]
If the old residue is $\rho_\xi(R)$, then with $\widetilde b_k$ as the starting lift,
\[
\rho_{\xi,\widetilde b_k}(R)
=
 k^{-1}\rho_{\xi,\widetilde b_e}(R)k.
\]
\end{proposition}

\begin{proof}
If the old lift ends at $\widetilde b_e\cdot\rho^{-1}$, right deck-equivariance shows that the new lift beginning at $\widetilde b_e\cdot k$ ends at $\widetilde b_e\cdot\rho^{-1}k$. In coordinates based at the new origin,
\[
\widetilde b_e\cdot\rho^{-1}k
=
(\widetilde b_e\cdot k)\cdot(k^{-1}\rho^{-1}k),
\]
so the new residue is $k^{-1}\rho k$.
\end{proof}

\begin{corollary}[Anchor-Independent Information]
If no fibre origin is fixed, the conjugacy class of the residue is more stable than a specific group element. In particular, when $H_\xi$ is abelian, conjugation is trivial and the exact residue is unchanged by a change of fibre origin.
\end{corollary}

\begin{plainbox}[Why this will matter in the next chapter]
The Theorem 5.1 benchmark in Chapter 28 assumes that the group of germs $H$ is finite abelian. Abelianity has its own role in the source theorem. Here it also brings one extra convenience: once the singular sector is fixed, changing the fibre origin no longer changes the exact residue. This convenience is still not translated into a psychological mechanism.
\end{plainbox}

\section{Enlarging the Ambient Context Does Not Automatically Change a Fixed Residue}

Chapter 24 emphasized that witness embedding is occurrence-supported. Suppose
\[
Q_n\subseteq Q_m
\]
is a chosen occurrence relation and the support of $R=(u;W_R)$ lies in $Q_n$. Define
\[
R^{(m)}=\bigl(u;E_{Q_n,Q_m}(W_R)\bigr),
\]
carrying the original $\xi$-anchor, starting lift, labels, and germ/connector decorations along the same finite trajectory.

\begin{proposition}[Stability of Residue under Compatible Context Enlargement]
Under the compatibility assumptions above,
\[
\rho_\xi(R^{(m)})=\rho_\xi(R).
\]
\end{proposition}

\begin{proof}
The ambient tile is larger, but the finite labeled trajectory used to calculate the residue and its anchored lift have not been rewritten. The lifted path is therefore identical step by step, and so is its endpoint fibre displacement.
\end{proof}

The proposition is weak on purpose. It does not say that arbitrary re-embedding, arbitrary transport between isomorphic copies, or an arbitrary change of singular sector preserves residue. It says only this: if we truly add ambient context without changing the finite decorated lift, the residue does not change merely because the surroundings became larger.

\section{Why an Undecorated Birooted Type Generally Does Not Determine Residue}

Two return occurrences may have the same abstract closed-birooted shape while lifting to different germ sheets in different connector contexts. Therefore
\[
[W_{R_1}]_{\mathrm{br}}=[W_{R_2}]_{\mathrm{br}}
\quad\not\Rightarrow\quad
\rho_\xi(R_1)=\rho_\xi(R_2).
\]

\begin{proposition}[Residue Does Not Descend to Bare Type without Additional Decoration]
If residue is to become type-level data, the equivalence relation used must preserve every datum that affects the anchored lift endpoint, including at least the relevant connector/boundary context, the singular-sector anchor, and the germ-lift decoration. Bare birooted isomorphism is generally insufficient.
\end{proposition}

\begin{definition}[Germ-Decorated Return Type]
Retain the bare birooted return shape together with decorations sufficient to determine the admissible connector context, the $\xi$-anchor, and the anchored lift endpoint. The equivalence class under decoration-preserving equivalence is called a \emph{germ-decorated return type}.
\end{definition}

\begin{plainbox}[A formal reminder that matters for psychology]
``These episodes look like the same kind'' corresponds only to similarity at the level of bare shape. If a difference changes future continuation, the formal model cannot quotient that difference away merely because the outward forms are similar. The point of residue calculus is precisely to preserve cases in which return-completion is already the same while finer continuation data remain different.
\end{plainbox}

\section{Composition: Only Admissible Based Returns Can Be Multiplied}

Residues cannot be multiplied merely because two ``agains'' have been noticed. One must first fix the same $\xi$-sector, compatible fibre anchor, and compatible decorations. The two return data must also admit based concatenation along an actual trajectory.

Write $R_1*R_2$ for the operation of first traversing the underlying based germ-decorated closed trajectory of $R_1$ and then continuing along that of $R_2$. This notation records based trajectory concatenation. It is not mechanical concatenation of two complete source return words as bare strings. Under the fixed right-torsor/inverse-label convention above, the multiplication order follows the traversal order.

\begin{proposition}[Multiplicativity of Residue]
If $R_1$ and $R_2$ are based return data in the same $\xi$-sector with compatible anchor and decorations, and if $R_1*R_2$ is admissible, then
\[
\rho_\xi(R_1*R_2)=\rho_\xi(R_1)\rho_\xi(R_2).
\]
The remainder of this manuscript uses only the fixed right-torsor/inverse-label convention of this section.
\end{proposition}

\begin{proof}
Write $\rho_i=\rho_\xi(R_i)$. By the fixed convention, the lift of $R_i$ from the chosen origin ends at the sheet labeled by $\rho_i^{-1}$. Right deck-equivariance, together with the stipulated based-concatenation order, gives the endpoint of the composite lift as
\[
\widetilde b_e\cdot(\rho_1\rho_2)^{-1}.
\]
The definition of residue therefore gives $\rho_\xi(R_1*R_2)=\rho_1\rho_2$.
\end{proof}

\begin{warning}[No Universal Cocycle Has Been Constructed Here]
This chapter records only a concrete monodromy-like multiplication on fixed singular-sector, fixed based/anchored data. It does not thereby produce a canonical cocycle on the entire path category or full groupoid, still less a Morita-invariant cohomology class. If such objects are needed later, they must be defined separately and their well-definedness proved.
\end{warning}

\section{Repeated Return: When Does \texorpdfstring{$h^k$}{h to the k-th power} Actually Make Sense?}

Let $R$ be a germ-decorated based return and put $h=\rho_\xi(R)$. Only when
\[
R^k=\underbrace{R*\cdots*R}_{k\text{ times}}
\]
is genuinely admissible in the same sector and anchor, with compatible decoration at every repetition, do we have
\[
\rho_\xi(R^k)=h^k.
\]
There are then three purely mathematical cases:
\begin{itemize}
\item if $h=e$, each projected return-completion is also fibre-trivial;
\item if $1<\operatorname{ord}(h)=q<\infty$, a single return leaves a residue, but the accumulated fibre displacement vanishes after the $q$th repetition;
\item if $\operatorname{ord}(h)=\infty$, each projected trajectory is closed while the resolved fibre never returns to its original sheet.
\end{itemize}

The last picture is striking, but it is emphatically not incompressibility. Quite the opposite: every datum $R=(u;W_R)$ already contains a return factor $u\in\mathcal R$ with a finite witness. What is nontrivial is the finer displacement behind the return-completion, not the absence of return-completion.

\section{Part-III Pointer: The Same ``Again'' Can Have a Different Afterwards}

Chapter 20 borrowed the relational skeleton of this chapter by writing one return on two ledgers:
\[
\text{return-completion}
\quad+
\text{continuation-resolution datum}.
\]
On the psychological side, this permits questions such as whether, after the same recognized return, subsequent prediction, action choice, or self-reference has changed. The stop lines must remain explicit:
\begin{itemize}
\item ``I still feel bad'' is not, by itself, residue;
\item ``I still remember it'' is not, by itself, residue;
\item ``my body still reacts'' is not, by itself, residue;
\item only when a predeclared finer resolution identifies a difference that stably changes continuation can that difference become a candidate residue-like observable.
\end{itemize}
Part IV establishes no psychological mechanism here. It proves only that, mathematically, return-completion and finer fibre displacement can be recorded separately.

\section{Chapter Conclusion: Return Is Not the End of Resolution}

The return calculus has acquired an additional layer, but the boundaries between its objects are now sharper:
\[
R=(u;W_R)
\quad\text{records concrete return-completion},
\]
while
\[
(R,\rho_\xi(R))
\quad\text{also records fibre displacement at anchored germ resolution}.
\]
Four rules are now frozen. First, a bare return does not automatically possess a residue. Second, exact residue becomes well-defined only after the anchor is fixed; changing the fibre origin generally preserves only its conjugacy class. Third, an undecorated birooted type generally does not determine residue. Fourth, multiplication and powers are valid only for actually admissible based compositions.

The next chapter changes the question. Residue tells us what an already existing return may still leave behind. The unique purely non-Hausdorff benchmark asks why, under a strong package of hypotheses, sufficiently long admissible words cannot avoid return forever.

\chapter{A Controlled Benchmark: Why pNH Is Not Enough, and When It Really Forces Return}
\label{chap:pnh-benchmark-en}

\par\medskip
\noindent\textbf{Why this chapter matters}\par
The calculus developed in the preceding chapters comes mainly from definitions, finite witnesses, and nested occurrence structure; it does not require a purely non-Hausdorff singularity. Here, for the first time, we enter a result that genuinely depends on special dynamical geometry.

For a psychology reader, the easiest mistake is to think that because pNH allows a distinction to disappear in regular contexts arbitrarily close to the singular point, a sufficiently long trajectory must somehow ``come back.'' Mathematically this is false. pNH supplies only a local possibility. Theorem 5.1 needs those possibilities to be organized into a finite, repeatable, uniformly controlled structure before uniform return forcing follows. Here we do nothing more, and nothing less, than unpack that logical chain.
\par\medskip

\begin{warning}[pNH Is a Strong Branch, Not a Synonym for Singularity]
Nekrashevych's Definition 2.2 first separates singularities into Hausdorff and non-Hausdorff cases, and then isolates purely non-Hausdorff singularities as a stronger case within the latter \citep{nek2018}. The regular-limit trivialization mechanism below therefore does not follow from bare singularity. A Hausdorff singularity specifically lacks the condition that identity regions accumulate arbitrarily close to the nontrivial germ. pNH is selected here because the theorem genuinely needs system-endogenous nearby trivializations, not because all singular germs carry the same psychological meaning.
\end{warning}

\section{First Gate: Purely Non-Hausdorff Means Only That the Difference Can Become Locally Invisible Arbitrarily Nearby}

Set the theorem aside for a moment and consider pNH itself. Let a group $G$ act on a Cantor space $X$, and let $\xi\in X$ be a germ-defining singular point. Write $G_\xi$ for the stabilizer and $G_{(\xi)}$ for the neighborhood stabilizer. If $g\in G_\xi$ is not the identity on any neighborhood of $\xi$, then it defines a nontrivial germ at $\xi$.

The purely non-Hausdorff condition can be written
\[
\xi\in\overline{\operatorname{Int}\operatorname{Fix}(g)}
\qquad\text{for every }g\in G_\xi.
\]
Here $\operatorname{Int}\operatorname{Fix}(g)$ is the topological interior of the region fixed pointwise by $g$. Thus, even when $g$ retains a nontrivial germ distinction at the singular focus $\xi$, every neighborhood of $\xi$ contains regular open regions on which $g$ is locally indistinguishable from the identity \citep{kuang2026a}.

\begin{plainbox}[Keep only this picture]
At $\xi$,
\[
(g,\xi)\neq(\operatorname{Id},\xi)
\qquad\text{as germs}.
\]
But however closely one zooms in, one can find a nearby regular region $U$ with
\[
g|_U=\operatorname{Id}|_U.
\]
Thus pNH supplies a kind of local resolution instability: the distinction retained at the current singular point need not remain visible in every nearby regular context.
\end{plainbox}

This is already enough to explain why Part III treats the pNH package as a benchmark for a third-order observation site: a distinction retained at the current root is not automatically a distinction that every admissible resolution must preserve. But there is still no uniform return theorem.

The reason is simple. pNH tells us that for each germ some trivializing regions can be found arbitrarily close to $\xi$. It does not automatically tell us whether all germ types can be controlled by one finite family of contexts, whether those local contexts recur in actual orbital/tile geometry in a controlled way, whether the waiting scales for different germ types admit a common bound, or whether sufficiently long admissible words must therefore hit genuine return witnesses. More conditions are needed.

\section{Second Gate: Condition (3) Turns ``Trivializable Somewhere'' into a Finite Menu}

The third hypothesis of Theorem 5.1 introduces a finite partition
\[
\mathcal P=\{P_1,\ldots,P_d\}
\qquad\text{of }\Omega(B)\setminus\{\xi\}.
\]
It requires that for every
\[
h\in H=G_\xi/G_{(\xi)},
\]
there be at least one piece $P_i$ accumulating at $\xi$ on which
\[
h|_{P_i}=\operatorname{Id}|_{P_i},
\]
and conversely that each $P_i$ trivialize at least one germ type.

The strengthening beyond pNH can be compressed into one contrast:
\begin{center}
\resizebox{0.96\linewidth}{!}{$\displaystyle
\begin{array}{ll}
\text{pNH:}&\text{each germ can become locally trivial somewhere arbitrarily nearby;}\\[1mm]
\text{Condition (3):}&\text{one predeclared finite menu covers all relevant germ types.}
\end{array}
$}
\end{center}
Finiteness is crucial. It compresses local phenomena scattered among infinitely many nearby regular points into finitely many manageable context classes. Without this step, even if each $h$ can separately disappear somewhere nearby, all such behavior need not admit uniform control at one common scale.

For each $P_i$, choose regular points
\[
\zeta_n\in P_i,
\qquad
\zeta_n\longrightarrow\xi.
\]
The corresponding rooted orbital graphs converge to a regular-limit graph $\Lambda_i$. Let $H_i\le H$ be generated by those elements that become trivial germs along this regular context. Then $\Lambda_i$ may be viewed as a controlled quotient of the graph of germs: sheet differences belonging to $H_i$ are identified in the quotient. Hence if $h\in H_i$, the distinct
\[
\operatorname{Id}\text{-sheet}
\qquad\text{and}\qquad
h\text{-sheet}
\]
of the fully resolved germ fibre are identified in $\Lambda_i$.

\begin{plainbox}[Why this still gives only a possible return-completion]
At this stage we know only that if a finite subwalk is projected to the appropriate regular-limit quotient $\Lambda_i$, its endpoints can be identified with the same branch. This makes return-completion possible at a theorem-controlled resolution. It does not yet say that the required configuration actually occurs in the orbital graph under study. Turning quotient-level possibility into a concrete return witness requires the next gate.
\end{plainbox}

\section{Third Gate: Strong Repetitivity Brings Quotient-Level Possibility Back to the Actual Orbit}

The role of strong repetitivity can be summarized in one sentence: a finite central configuration visible in a regular-limit graph can recur in the actual orbital/tile geometry in a controlled way.

The original proof takes finite central parts of the $\Lambda_i$ and organizes them into bridges. Strong repetitivity ensures that these finite patterns do not live only inside a limit object; they can be realized again in actual orbital graphs or truncation tiles. Only then can
\[
\text{``the endpoints are identified in }\Lambda_i\text{''}
\]
be converted into
\[
\text{``an actual word contains a finite closed subwalk.''}
\]
Schematically,
\begin{center}
\resizebox{0.96\linewidth}{!}{$\displaystyle
\boxed{\text{return-completion possible in a regular-limit quotient}}
\quad\xRightarrow{\text{strong repetitivity}}\quad
\boxed{\text{actual finite return witness}}.
$}
\end{center}

This step is also useful for disciplining the language of Part III. An ``alternative admissible perspective'' that works only at the level of interpretation, but never returns to the same actual material with a repeatable finite witness, has not acquired structural force in the sense of this benchmark.

\section{Fourth Gate: Why a Single Uniform Length Bound \texorpdfstring{$M$}{M} Can Finally Be Chosen}

Even if each germ type eventually produces a return somewhere, one final question remains: why should there be one common $M$ such that every sufficiently long admissible word contains a return?

The finiteness of the group of germs, the finite partition, and the controlled singular orbital/tile geometry in Theorem 5.1 leave only finitely many relevant germ/context configurations. Strong repetitivity provides uniform recurrence control for these finite configurations. The original proof can therefore choose a common scale covering every relevant case and obtain
\[
\exists M>0
\quad\text{such that}\quad
|w|>M
\Longrightarrow
w\text{ contains a return subword}.
\]
Together with the language-theoretic equivalence in the preceding chapter, this shows that under the composable $S$-expression convention of the source, the incompressible word language has bounded length. No presentation-independence across arbitrary representatives is asserted here.

The uniformity comes from the entire package of theorem hypotheses working together. It cannot be credited to pNH alone.

\section{Kuang Theorem 5.1: Reassemble the Four Gates into the Formal Statement}

The formal theorem can now be read without treating its five assumptions as an arbitrary checklist.

\begin{theorem}[Finite Incompressibility under a Unique pNH Singularity]
Let $G$ be a bounded-type group and assume:
\begin{enumerate}
\item $G$ has a unique purely non-Hausdorff singularity $\xi\in\Omega(B)$;
\item $H=G_\xi/G_{(\xi)}$ is finite abelian and has invariant domains;
\item there exists a finite partition $\mathcal P$ of $\Omega(B)\setminus\{\xi\}$ such that, for every $h\in H$, some $P\in\mathcal P$ accumulating at $\xi$ satisfies $h|_P=\operatorname{Id}|_P$, and conversely every $P\in\mathcal P$ admits at least one $h\in H$ with $h|_P=\operatorname{Id}|_P$;
\item the orbital graph $\Gamma_\xi$ is thin and one-ended;
\item the truncation tiles of $\xi$ are strongly repetitive.
\end{enumerate}
Then $G$ has only finitely many incompressible elements. Under the expression-level convention of Definition 3.11 in the source paper, an element is considered together with a given composable $S$-expression $g=F_1\cdots F_n$; the proof directly yields a uniform bound $M$ such that every admissible $S$-word of length greater than $M$ contains a return subword \citep{kuang2026a}. This book does not automatically upgrade that terminology to a group-element invariant independent of the chosen representative.
\end{theorem}

\begin{warning}[The Most Important Logical Stop Line]
The theorem has the form ``many hypotheses jointly imply the conclusion.'' Its logic is not an automatic escalation from one property:
\[
\begin{gathered}
pNH
+
\text{finite germ/algebraic control}
+
\text{finite partition cover}\\
+
\text{orbital geometry}
+
\text{strong repetitivity}
\Longrightarrow
\text{uniform return bound }M.
\end{gathered}
\]
In particular, pNH alone does not imply $M$, and the remaining hypotheses may not be discarded as technical decoration.
\end{warning}

\section{``One's Traverse Is Another's Return'': The Phrase Can Now Be Made Precise}

Fix a type-$h$ traverse and the germ-resolved anchor/decorated lift required by Chapter 27. In the fully resolved graph of germs, it may travel from the identity sheet to the $h$-sheet and hence fail to close at sheet resolution:
\[
\widetilde b_e\longrightarrow\widetilde b_h.
\]
If the finite context family in Condition (3) contains $P_i$ with $h\in H_i$, then the regular-limit quotient
\[
\widetilde\Gamma_\xi\longrightarrow\Lambda_i
\]
collapses this $h$-distinction, so the endpoints of the same finite subwalk land on the same branch in $\Lambda_i$. It is closed at that quotient resolution. Strong repetitivity then returns the required finite configuration to the actual orbital/tile structure and produces a genuine closed subwalk, hence a concrete return witness.

Schematically,
\[
\begin{gathered}
\text{open in the fully resolved germ fibre}\\
\Downarrow\ \text{theorem-controlled regular-limit quotient}\\
\text{closed after a specified sheet distinction is collapsed}\\
\Downarrow\ \text{strong repetitivity}\\
\text{actual finite return witness}.
\end{gathered}
\]
Thus ``another'' never means ``pick any viewpoint you like.'' It refers to a finite family of theorem-authorized, system-endogenous quotient contexts traceable to regular contexts.

\section{Assumption Audit: What Each Hypothesis Does, and What We Leave Untranslated on Purpose}

\begingroup\small
\begin{longtable}{>{\raggedright\arraybackslash\bfseries}p{0.23\textwidth}>{\raggedright\arraybackslash}p{0.35\textwidth}>{\raggedright\arraybackslash}p{0.32\textwidth}}
\toprule
Mathematical hypothesis & What it supplies in the benchmark & Treatment in Part III\\
\midrule
unique singularity & restricts the theorem to one singular sector and avoids multi-sector bookkeeping & scope simplification only; not ``a person has only one center''\\
purely non-Hausdorff & arbitrarily near regular trivializations for every stabilizer germ & the current distinction need not remain visible in every admissible nearby resolution\\
finite $H$ & finitely many germ types, supporting uniform control & not translated into a finite number of personality states\\
abelian & removes part of the noncommutative bookkeeping; in Chapter 27 also removes conjugacy ambiguity under fibre-origin change & no independent psychological correspondence\\
invariant domains & algebraic/action regularity assumed by the source theorem & currently untranslated\\
finite partition condition & compresses scattered trivializations into a predeclared finite context family covering relevant germs & alternative contexts must be constrained in advance, not invented without limit after the fact\\
thin + one-ended & controls singular orbital/tile geometry in the theorem & untranslated technical remainder\\
strong repetitivity & re-realizes finite regular-limit possibilities in actual orbital/tile structure and supports uniform recurrence control & alternative organization must return to actual comparable material\\
\bottomrule
\end{longtable}
\endgroup

The discipline established in Part III still applies: mathematical importance in a proof does not itself grant a hypothesis permission to enter psychology.

\section{Structural Return-Forcing Schema: Extract the Proof Architecture without Stealing the Theorem}

Inside this benchmark, the proof architecture may be compressed as
\[
\begin{gathered}
\text{arbitrarily near local trivializations}
+
\text{finite controlled context cover}\\
+
\text{recurrent realization in actual geometry}
+
\text{finite/uniform control}\\
\Longrightarrow
\text{a uniform upper bound on waiting for a return}.
\end{gathered}
\]
The value of this schema is to show where the theorem obtains its force. It is not promoted into a new general theorem.

\begin{warning}[Not a General AF-by-Discrete Theorem]
This manuscript has not proved any of the following:
\begin{itemize}
\item every pNH singularity has a uniform upper bound on waiting for a return;
\item every AF-by-discrete groupoid has finite incompressibility;
\item multi-singularity systems automatically admit one common uniform $M$;
\item a psychological process that has a ``third-order observation site'' must eventually produce a return.
\end{itemize}
Outside the specific hypotheses of Theorem 5.1, these conclusions must be proved anew or abandoned.
\end{warning}

\section{Relation to the Multi-Boundary Prototype: Two Benchmarks Solve Different Problems}

The present benchmark must now be placed beside the multi-boundary bridge-type prototype added in Chapter 26. Both come from the bounded-type/tile-inflation world, but they address different kinds of complexity.

\begingroup\small
\begin{longtable}{>{\raggedright\arraybackslash\bfseries}p{0.21\textwidth}>{\raggedright\arraybackslash}p{0.35\textwidth}>{\raggedright\arraybackslash}p{0.34\textwidth}}
\toprule
 & Multi-boundary / multiple bridge types & Unique-pNH benchmark\\
\midrule
Distinguished data & several infinite boundary points / bridge types $\xi_x$ & several germ sheets over one singular point $\xi$\\
Finite coding & $\operatorname{typ}\Phi(\tau)=x_1\cdots x_k$ records cross-sector passage & $\widetilde\Gamma_\xi\to\Lambda_i$ records different admissible resolutions of one site\\
Core problem & how a return grammar is jointly organized by several sectors & how distinctions retained at one site trivialize in nearby regular resolutions\\
Uniform return bound & Proposition 5.7 assumes a uniform $d$; the general setup does not automatically produce one & Theorem 5.1 genuinely derives finite incompressibility / a uniform $M$ under extra hypotheses\\
Weakest Part-III translation & no organizing sector must monopolize the whole return grammar & the resolution of the current root has no final authority\\
\bottomrule
\end{longtable}
\endgroup

Thus the single-pNH case is not the only ``third-order mathematical example'' in the book. More accurately, the book now has two complementary prototypes:
\[
\boxed{\text{multi-site passage plurality}}
\qquad\text{and}\qquad
\boxed{\text{single-site resolution plurality}}.
\]
The former already has a source-domain prototype in Subsection 5.2 of the v1 preprint \citep{kuang2024v1}; the latter is supplied especially clearly by the pNH theorem family. A complete structural subject combining both would require additional cross-sector compatibility data, which remain outside the present formal core.

\section{Chapter Conclusion: pNH Gives Possibility; the Theorem Gives Forcing}

The central distinction of this chapter can be frozen in one line:
\begin{center}
\resizebox{0.96\linewidth}{!}{$\displaystyle
\boxed{\text{pNH gives local trivialization possibilities; Theorem 5.1 gives uniform return forcing.}}
$}
\end{center}
Between pNH and a uniform $M$ lies an entire package of additional controls: finite germ complexity, a finite partition cover, controlled orbital geometry, and strong repetitivity. Precisely because these costs remain visible, Chapter 21 can borrow ``one's traverse is another's return'' as a bounded benchmark rather than as permission to reinterpret anything from any angle until it looks like a return.

\chapter{Formal Audit: Assemble the Mathematical Map of Part IV}
\label{chap:formal-audit-en}

Part IV adds no further theorem here. It closes by assembling the formal relations already established into one dependency map: which conclusions follow from definitions, which depend on additional geometry, and which hold only inside a fixed presentation.

\section{Five Independent Formal Branches}

\begingroup\small
\begin{longtable}{>{\raggedright\arraybackslash\bfseries}p{0.25\textwidth}>{\raggedright\arraybackslash}p{0.65\textwidth}}
\toprule
Branch & Frozen content\\
\midrule
occurrence / type & When persistence, connectors, or residue are involved, fix the concrete occurrence first; type is the classification obtained after position has been forgotten.\\
embedding / decomposition & Embedding preserves the same finite witness inside a larger context; decomposition cuts the same actual trajectory at a lower level. The operations are not inverses.\\
return / witness / language & Source word, concrete return-completion witness, and the full return language are separate ledgers; filtration is cumulative bookkeeping inside a fixed presentation.\\
incompressibility & Contiguous-factor avoidance relative to the complete return language, not the psychological condition of being ``unable to get out.''\\
germ residue / pNH benchmark & Residue is defined only with a fixed singular sector, anchor, and decoration; passing from pNH to uniform return forcing still requires the complete hypotheses of Theorem 5.1.\\
\bottomrule
\end{longtable}
\endgroup

\section{Presentation Dependence Must Remain Visible}

Generators, Bratteli telescoping, tile occurrences, boundary conventions, birth levels, and concrete residue coordinates may all change with the chosen presentation. These chapters use them only within the presentation already declared. More stable content includes the types of relations defined, the isomorphism type of isotropy along one orbit, and the fact that under fixed decoration return-completion and finer displacement can be recorded separately.

In particular, AF-by-discrete and bounded type remain distinct, and the general groupoid statement must remain separate from the tile-inflation package used in the author's work. In the general groupoid formulation, bounded type gives uniform per-fibre / per-finite-tile boundary control. In the tile-inflation form used here, finite boundary support on infinite tiles when present and uniform $|V_n|,|E_n|$ are imposed as well, so a level-total boundary burden across tile types is uniformly bounded as a consequence. No theorem may silently replace one formulation by a stronger one without carrying the additional hypotheses. Chapter 10 used only the forward implication from abstract AF-by-discrete structure to the finite-support picture under the stated hypotheses; it did not smuggle the converse into terminology.

Part IV has likewise not proved that human memory obeys formal persistence, that psychological episodes possess Bratteli levels, that the brain contains germ fibres, or that Third-Order Consciousness is non-Hausdorffness. It establishes only this: once these formal relations are chosen, which objects must be kept on separate ledgers, which compositions are legitimate, and which strong conclusions require extra hypotheses. If empirical material fails to satisfy those relations, the correspondence must change; the mathematical definitions frozen here do not.

\part{Controlled Cases and Empirical Bridges: From Structural Explanation to Testable Questions}

\begin{plainbox}[Cases are not proofs; counterexamples determine boundaries; mechanism claims require independent measurement]
Part III established structural correspondences and Part IV froze the finite-scale formal core. Part V introduces no new mathematical object and does not treat cases as ``validation'' of the theory. It does three things only: tests whether the existing language actually provides more resolution than ordinary narrative; actively seeks situations in which the model does not apply; and rewrites Level-II structural mappings as questions that future psychological or behavioral studies could test.
\end{plainbox}

\begin{principlebox}[Describe first, segment second; segment first, correspond third]
Whenever a case enters this model, preserve the raw event order first, then disclose the episode segmentation, the current question, and the observation position. Only after those analytic choices are visible should one discuss a psychological return hypothesis, a continuation profile, or a structural comparison with Part IV. It is prohibited to start from a desired mathematical conclusion and select events backward until a case appears to fit the model.
\end{principlebox}

\begin{plainbox}[Case notation and mathematical notation remain on separate ledgers]
Part IV froze the mathematical-native-domain notation
\[
R=(u;W_R),\qquad \rho_\xi(R).
\]
Part V never writes a real psychological episode directly as one of those mathematical objects. It uses weaker case notation on purpose: $D$ for the raw event record, $\Pi$ for a candidate segmentation, $Q$ for the currently explicit question or return-completion criterion, $\Sigma$ for one or more candidate organizing positions/sectors, $\mathcal R_{\mathrm{psych}}$ for a psychological return hypothesis, and $C^+(E)$ for the case-level continuation profile after an experiential segment $E$. From the outset $\Sigma$ may contain several positions: a case should test whether a return is organized mainly within one position or whether the order of passage across several positions must be recorded. The symbols $R,W_R,\rho_\xi$ reappear only when the text explicitly says that it is referring back to Part IV.
\end{plainbox}

\chapter{Case Protocol: How to Use the Model without Fitting a Person into a Formula}
\label{chap:case-protocol-en}

\par\medskip
\noindent\textbf{Why this chapter matters}\par
From this part onward, mathematics no longer stands in front of experience. The preceding parts have already supplied formal distinctions such as boundary, return, germ, and residue. The harder question now runs in the opposite direction: when we face a person's behavior, narrative, memory, and relational experience again, how do we know that we are seeing structure in the material rather than merely seeing the mathematical vocabulary we have just learned?

For this reason, we proceed more slowly on purpose. We first record the facts, and only then disclose how the analyst segmented them. We first look for transition points repeatedly forced by the material itself, and only then ask whether any of them deserve comparison with mathematical boundary structure. If a case fits the model only after repeated changes of segmentation, question, or explanation, the model should be weakened before the facts are revised.
\par\medskip

\section{Why a Case Must Move More Slowly than a Theory}

A reader already fluent in the vocabulary of this book can easily begin to ``see'' boundary, return, residue, and incompressibility everywhere. That fluency is itself a risk. Once a concept becomes useful enough, a researcher begins to recognize it faster and faster: repeatedly opening a phone looks like return; a relational turning point looks like boundary; ``why am I doing this again?'' looks like some kind of return-completion. But looking similar is not the same as the material already supporting the structural judgment.

Mathematical concepts, once they become too convenient, can absorb unanalyzed material just as psychological labels can. This is precisely the kind of symbolic overreach criticized by the \emph{Geometry of Existence}: a concept begins as a tool for seeing relations and ends by dictating what reality must look like \citep{geometry2026}.

A working case therefore begins as no more than a sequence of events. Someone sends a message; no reply arrives; after some time the chat is opened again; an interpretation forms; affect changes; another action follows. At this point there is no return and no boundary. Even the judgment that these events belong to one episode is still an analytic choice.

Call the unstructured material $D$. The researcher must then disclose each layer of judgment: which events are provisionally grouped into one segment, what question is actually being answered, from which organizing position the comparison is made, which repetitions are provisionally understood as ``again,'' and which entirely different psychological accounts could also explain the same material. The order of case analysis is therefore
\[
\begin{gathered}
\text{what happened}\\
\Downarrow\\
\text{how we segmented it, what we are asking, and from where we judge}\\
\Downarrow\\
\text{whether there is a candidate return}\\
\Downarrow\\
\text{alternative explanations and alternative resolutions}\\
\Downarrow\\
\text{later predictions and independent observables}\\
\Downarrow\\
\text{what observation would make this explanation fail.}
\end{gathered}
\]
Only when precise bookkeeping is needed does this part abbreviate these layers with \(D\), \(\Pi\), \(Q\), \(\Sigma\), \(\widehat{\mathcal R}_{\mathrm{psych}}\), and \(C^+\). These symbols are not new psychological entities. Their first purpose is to prevent the researcher from changing the question after seeing the outcome.

\section{Ten Entries in a Controlled Case Record}

\begin{principlebox}[Case protocol: ten basic records]
A controlled case should disclose at least the following ten entries.
\begin{enumerate}
\item \textbf{Raw events.} Describe observable behavior, temporal order, and the participant's own words as far as possible before applying theoretical labels.
\item \textbf{Candidate segmentation.} State which events are provisionally treated as one episode and why the cuts are placed where they are. Segmentation is an analytic choice, not a label already attached to the material.
\item \textbf{Current question $Q$.} Declare in advance what is being asked: whether action has ended, whether relational intention is known, whether the same pattern has appeared again, or something else. Different $Q$ may legitimately require different distinctions to be preserved, but the question may not be changed ad hoc after the result is known.
\item \textbf{Candidate organizing positions.} Who is comparing what, by which standard, from which position, and with what prediction about the next step? Does the episode contain a change of position? Only after these psychological descriptions are explicit do we abbreviate them as $\Sigma=\{\sigma_1,\ldots,\sigma_r\}$. ``Position'' here does not presuppose a mathematical root or singular point.
\item \textbf{Evidence for ``again.''} If two finite episodes are claimed to stand in a return-like relation, state which finite facts support the judgment: the same action, the same unresolved question, the same evaluative position, the same expectation, or a more complex return-completion relation. The psychological material is not renamed a formal witness $W_R$ from Part IV.
\item \textbf{Alternative resolution and alternative model.} Supply at least one alternative question/frame for the same material and one baseline psychological model that can explain the material without this framework.
\item \textbf{Scale and robustness check.} Record what happens if the episode segmentation is made finer or coarser or the self-reference frame is changed. These are case audits, not the formal filtration of Part IV.
\item \textbf{Continuation profile.} Record predictions, available actions, expectations, or next transitions after the episode, written $C^+(E)$. This is a case-level proxy, not germ residue $\rho_\xi(R)$.
\item \textbf{Independent observable consequences and exit conditions.} If the structural analysis adds information, what independent judgment, memory, behavior, or prediction should change? What outcome would force the present correspondence to exit?
\item \textbf{Translation remainder.} Record experiential qualities, motives, social meanings, and mechanisms that the model still does not explain.
\end{enumerate}
\end{principlebox}

\begin{plainbox}[In ordinary language]
The purpose of this format is to make explanation slower than fact, and to make the question itself public. First record what happened; then say how the material was cut, what is being asked, and from where it is judged. Only then introduce ``again,'' continuation, or a mathematical interface. At every step a reader should be able to say ``I disagree with this segmentation'' or ``you changed the question'' without thereby rejecting the entire model.
\end{plainbox}

\section{A Boundary Must Be Forced out by the System, Not Drawn First by the Theory}

The easiest place for a formal model to cheat is at the boundary. If every change in a case can be followed by the retrospective announcement ``there is a new boundary here,'' then any sufficiently complicated experience can eventually be segmented until it fits the model. Such a boundary carries little empirical content; it is first of all the analyst's cut.

This book requires the opposite order. Observe how the system repeatedly runs, then ask whether some positions are forced into fissures again and again. A stronger evidence pattern should arise across several comparable episodes, not from one segment alone. For example, we may repeatedly see a familiar local policy or prediction fail at similar positions; the system must recruit more distant context; candidate return endpoints or passages between organizing positions concentrate near these transitions; and after crossing them the available actions, expectations, or continuation profile $C^+$ change in a stable way.

Only after such patterns appear do we have reason to propose a candidate singular-boundary-like interface. The working order is therefore
\[
\begin{gathered}
\text{repeated continuation failure / candidate-return concentration}\\
\Downarrow\\
\text{candidate singular-boundary-like interface}\\
\Downarrow\\
\text{independent continuation test}.
\end{gathered}
\]
A boundary is not a line drawn in advance by the researcher. It is a position gradually exposed as the system repeatedly runs, returns, and fails. If one must already know the theoretical answer in order to locate the boundary, the interface hypothesis depends too heavily on the researcher.

\begin{plainbox}[Keep AF-by-discrete, the singular subset, and bounded type on separate case ledgers]
The mathematical native domain begins from two strict facts. Under a fixed generating set and tile presentation, AF-by-discrete yields a global finite-support picture for all persistent boundary points; its singular subset
\[
\partial^{\infty}_{S,\mathrm{sing}}
=
\{\xi\in\partial_S^\infty:H_\xi\neq\{1\}\}
\]
is automatically finite. The general groupoid bounded-type condition further requires the generator-boundary burden of each individual finite tile / AF fibre to be uniformly bounded across levels; the tile-inflation form used here also controls finite boundary support on infinite tiles when present and the Bratteli level sizes $|V_n|,|E_n|$.

The psychological translation in this book targets only singular boundary germs. What a case can directly pressure-test is therefore a narrower singular-interface shadow. Let $B_m$ denote the number of candidate singular-interface loci that can be reidentified across material. If those loci proliferate without bound as material accumulates, the present singular-germ translation loses support. Conversely, even if $B_m$ appears stable, this does not empirically establish the complete AF-by-discrete global boundary count, because regular boundary points are not measured by the present psychological operationalization.

Likewise, the $N_r$ used later is defined as a per-unit maximum over predeclared comparable episode/block units at each resolution. It is at most a singular-sector proxy for the per-fibre / per-tile direction of general groupoid bounded type. It omits regular boundary data and does not operationalize the separate $|V_n|,|E_n|$ controls of the tile-inflation package. The empirical procedure therefore exerts only structural pressure; it does not turn psychological samples into a mathematical definition.
\end{plainbox}

\begin{warning}[A Case Is Not a Diagnosis]
The words root, return, incompressibility, and residue in this model do not correspond to clinical diagnoses. Repetitive thought does not make a person ``mathematically incompressible,'' and repeated relational conflict does not acquire a return witness merely because the language of cycles seems apt.
\end{warning}

\section{A Twenty-Second Micro-Example}

Consider the following material:
\begin{quote}
The phone vibrates. A person sees that it is not the message they are waiting for and turns the phone face down. Ten seconds later they pick it up and refresh. Then they say, ``Forget it, I won't look anymore.'' Half a minute later they open the same interface again and think, ``Why am I waiting again?''
\end{quote}

If the researcher writes ``checking cycle'' immediately, much of the structure has already been hidden. A slower record asks whether the first look and the second refresh belong to one episode; whether ``I won't look anymore'' supplies an endpoint; whether the ``again'' linking the third opening to the first is organized by the same action, the same expectation, or the same self-evaluation; and whether the third opening has shifted in purpose from ``waiting for a reply'' to ``checking whether I am still waiting.'' The same visible action may occupy a different organizing position.

The value of the model is therefore not to rename repetition, but to force the distinctions
\[
\text{same action}
\neq
\text{same episode}
\neq
\text{same return}
\neq
\text{same continuation}.
\]
More importantly, the fissure worth following may not be the third phone opening itself. It may become visible only across several similar episodes: whenever ``waiting for the other person's reply'' turns into ``observing why I am still waiting,'' the evaluative standard, action options, and later predictions may change together.

\begin{plainbox}[One episode can suggest an interface; repeated organization begins to probe it]
A single case can propose a candidate interface. Only continuation change that recurs across multiple comparable episodes gradually adds empirical weight to that candidate position. The later cases in Part V follow this rule.
\end{plainbox}

\section{The Minimum Requirement of Third-Order Observation: Keep the Position of Interpretation in the Record}

Third-order observation does not give the researcher a higher or more objective position. It first removes that privilege.

Suppose the researcher or participant says, ``Obviously this is because he is afraid of rejection.'' That sentence already changes which later facts appear important: later silence may be read as avoidance, another phone check as anxiety, and an ordinary explanation may be absorbed into the same story. The interpretation must therefore enter the case record as well. Which earlier material supplied it? Which segmentation does it use? Which alternatives does it exclude? Does it pre-organize every later event as evidence of ``fear of rejection''?

A third-order record therefore retains at least one additional layer: who is interpreting now, what past material is being used, how the material has been segmented, and which future possibilities the interpretation thereby excludes in advance. This does not require the absence of judgment. It requires judgment to retain its coordinates.

First-order material tells us what happened. Second-order organization tells us how those events became ``my experience.'' Third-order recording asks one more question: how is this ``my experience'' now being observed, named, and reorganized? Its minimum requirement can be compressed into one line:
\begin{center}
\resizebox{0.96\linewidth}{!}{$\displaystyle
\boxed{\text{the observed material enters the record; the observation position enters the record as well.}}
$}
\end{center}
This directly inherits the methodology of the Translation Protocol and the Geometry of Existence: a judgment must have a position, and a model must preserve alternative projections and failure conditions \citep{translation2026,geometry2026}.

\chapter{Nonresponse, Anger, and Corrective Simulation: A Controlled Case of ``Needing Completion''}
\label{chap:nonresponse-case-en}

\section{Raw Events: Let the Facts Remain Sparse First}

Let us consider an anonymized working case. Two people have recently begun interacting. One of them initiates contact and receives no reply. Over the next several hours the following sequence occurs:
\begin{center}
\resizebox{0.96\linewidth}{!}{$\displaystyle
\begin{array}{ll}
e_1:&\text{send a message;}\\
e_2:&\text{wait and check the chat interface several times;}\\
e_3:&\text{confirm that there is still no reply;}\\
e_4:&\text{anger or irritation rises;}\\
e_5:&\text{form the judgment ``there should at least be a response'';}\\
e_6:&\text{decide not to initiate again for now;}\\
e_7:&\text{begin rehearsing an imagined conversation about ``explaining the principle'';}\\
e_8:&\text{after the simulation ends, notice that anger has been reactivated.}
\end{array}
$}
\end{center}
Write $D=(e_1,\ldots,e_8)$. At this stage we possess only an event order. The material has not established a mechanism in which nonresponse causes anger, nor has it shown that the person possesses one unified ``need for completion.''

``Corrective simulation'' is also only a case-level descriptive phrase here. It means repeatedly imagining, without continuing the real interaction, how to make the other person understand a principle, acknowledge a fact, or accept an evaluation. The phrase is not a clinical diagnosis and does not presuppose why the simulation occurs.

\section{Declare the Question First: There Are Actually Three Different Completion Questions}

The ordinary psychological sentence ``I need closure'' easily compresses different questions into one. Following the Chapter 30 protocol, separate at least three questions:
\begin{center}
\resizebox{0.96\linewidth}{!}{$\displaystyle
\begin{aligned}
Q_A:&\quad\text{Do I still need to take another action?}\\
Q_E:&\quad\text{Do I know the other person's intention, reason, or relational stance?}\\
Q_N:&\quad\text{Have I received the response or normative acknowledgment I believe was due?}
\end{aligned}
$}
\end{center}
There is also a fourth descriptive question:
\[
Q_S:\quad\text{Is the internal corrective simulation still continuing?}
\]
Unlike the first three, $Q_S$ first records an ongoing process; it does not prescribe when that process ought to end.

At one moment it is perfectly possible that
\begin{center}
\resizebox{0.96\linewidth}{!}{$\displaystyle
Q_A=\text{settled},\qquad
Q_E=\text{unresolved},\qquad
Q_N=\text{unresolved},\qquad
Q_S=\text{ongoing}.
$}
\end{center}
A decision not to contact the person again does not reveal why they were silent. Lacking an explanation does not automatically mean there remains a real-world action that must be performed. The first information gain of the case is simply to prevent nonresponse, non-understanding, non-acknowledgment, and ``I am still thinking about it'' from being swallowed by one label.

\begin{plainbox}[Why $Q$ must be fixed first]
If the researcher first says ``this person lacks psychological closure,'' then after seeing that action has stopped silently changes the intended criterion to ``emotional resolution,'' and after affect subsides changes it again to ``the other person never acknowledged it,'' the model can never fail. $Q$ has a modest function: say in advance what would count as completion, then observe whether it actually occurred.
\end{plainbox}

\section{Candidate Episode Segmentation: One Event Sequence Allows More than One Legitimate Cut}

One coarse segmentation is
\[
E_1=(e_1,e_2,e_3,e_4,e_5,e_6),
\qquad
E_2=(e_7,e_8).
\]
This $\Pi_1$ places waiting, evaluation, and the action decision in the first segment and the later internal simulation in the second. It is suitable for asking whether simulation continues after real-world action has ended.

Another segmentation $\Pi_2$ begins a new segment at $e_5$:
\[
E_1'=(e_1,e_2,e_3,e_4),
\qquad
E_2'=(e_5,e_6,e_7,e_8).
\]
Here the working hypothesis is that once the normative judgment ``there should at least be a response'' enters, the withdrawal decision and corrective simulation may belong to one longer organizing process.

Both cuts are compatible with the material, but they answer different questions. For this reason, we do not treat one episode segmentation as the uniquely correct answer hidden inside the facts. The researcher must record which $\Pi$ is being used and test whether the conclusion survives reasonable alternatives.

\section{Candidate Organizing Positions: After Anger Appears, Who Is Requiring What?}

At least three candidate organizing positions can be distinguished:
\begin{enumerate}
\item $\sigma_A$, the \textbf{action position}: ``Should I keep contacting them?'' It selects the next action.
\item $\sigma_N$, the \textbf{normative-evaluation position}: ``Basic reciprocity should include a response.'' It judges whether what happened satisfies a rule.
\item $\sigma_C$, the \textbf{corrective position}: ``The other person should understand why this is inappropriate.'' It keeps attention directed toward an interaction that has temporarily stopped and tries, in imagination, to change the other person's understanding.
\end{enumerate}
We may provisionally write
\[
\Sigma=\{\sigma_A,\sigma_N,\sigma_C\}.
\]
These positions may switch within seconds or overlap. They are not three personalities or three brain regions. The inventory forces only one question: what is currently deciding that ``this is not over''---the next action, a normative judgment, or the demand that the other person understand something?

Third-order observation adds little here, but what it adds matters. It does not abolish the judgment ``the other person should respond.'' It keeps the position issuing that judgment in the record. The rule may still be reasonable without automatically acquiring explanatory authority over everything that follows.

\section{Where Is the Interface? Nonresponse Itself Cannot Be Renamed a Boundary}

Chapter 30 requires an interface to be exposed by repeated system behavior. This case therefore cannot declare $e_3$ a boundary-like moment merely because it is the discovery that no reply has arrived.

From this single episode we can propose at most two candidate interfaces:
\begin{center}
\resizebox{0.96\linewidth}{!}{$\displaystyle
\begin{aligned}
b_1:&\quad\text{from waiting for a response to ending real-world action;}\\
b_2:&\quad\text{from ``I will not act again'' to ``I still need to correct the other person in imagination.''}
\end{aligned}
$}
\end{center}
Actual boundary detection requires more comparable episodes. If attention, prediction, and action repeatedly reorganize near $b_1$ across later interactions, and corrective simulation repeatedly begins near $b_2$, then those positions gradually acquire empirical weight as singular-boundary-like candidates. Conversely, if every case requires inventing a new turning point that cannot be reidentified with earlier loci, the finite singular-interface/persistent-locus picture receives no support.

The detection order is
\begin{center}
\resizebox{0.96\linewidth}{!}{$\displaystyle
\text{repeated organizing shift}
\longrightarrow
\text{candidate interface}
\longrightarrow
\text{independent continuation difference}.
$}
\end{center}
One silence can create a problem. Only repeated organizing pressure begins to probe a structural fissure.

\section{The Same Facts, Two Legitimate Resolutions of Completion}

Return now to the observation-site benchmark of Chapter 21. We do not reproduce the purely non-Hausdorff theorem here. We borrow only one methodological discipline: the same material can preserve different distinctions under different questions that were declared in advance.

For $Q_E$---``Do I know why the other person was silent and how they understood the interaction?''---the material remains unresolved. There is no answer and therefore no motive or relationship-intention data.

For $Q_A$---``Do I still have another action to perform?''---suppose the participant's predeclared action rule is ``send one message; if there is no reply within a specified time window, stop initiating.'' The same material can then complete the action episode:
\[
\text{send message}
\longrightarrow
\text{wait}
\longrightarrow
\text{no reply}
\longrightarrow
\text{stop further action}.
\]
Both of the following can therefore be true:
\[
\boxed{\text{the other person's relational intention remains unresolved}}
\]
and
\[
\boxed{\text{my present action episode has reached completion}.}
\]
This is not ``changing the story to feel better.'' Both $Q$ were declared before the analysis and both can fail: a later reply changes $Q_E$; a decision to keep contacting the person means $Q_A$ cannot be recorded as completed.

\begin{plainbox}[Reference back to Chapter 21: borrow the benchmark, do not turn the case into pNH]
In the mathematical benchmark, a fine germ resolution can preserve a distinction that a system-authorized regular-limit resolution collapses. The psychological side retains only this relation: a distinction that must be preserved under one question can be provisionally irrelevant to the completion criterion of another legitimate question. The case does not claim that nonresponse forms a purely non-Hausdorff singularity, and changing the question is not literally a mathematical quotient.
\end{plainbox}

This bookkeeping also has a practical consequence. If the participant accepts only ``she must tell me why'' as a completion criterion, action, explanation, and normative acknowledgment become tied together. If the researcher instead forces the material into ``you have already decided not to contact her, so it is over,'' real uncertainty, anger, and value judgment are erased. Third-order recording allows all of them to remain present and requires each to state which question it answers.

\section{Return Hypothesis: One Corrective Simulation Is Not Yet an ``Again''}

So far the chapter has described only one episode. It cannot establish a psychological return from this material alone. Return audit begins only when the participant or researcher proposes an ``again'' across episodes.

Suppose the participant says, ``Why am I teaching people a lesson in my head again?'' This proposes a candidate $\widehat{\mathcal R}_{\mathrm{psych}}$, but it still does not provide a witness. One must locate an independently recorded finite past episode $E_{\mathrm{past}}$ and compare it with $E_{\mathrm{now}}$. A stronger return hypothesis requires several relations to recur together:
\begin{itemize}
\item a similar entrance condition: real interaction has become blocked or stopped;
\item a similar organizing shift from action/evaluation into the corrective position $\sigma_C$;
\item a similar finite process: rehearsing an imagined conversation in which the other person is made to understand a principle;
\item a describable exit condition: fatigue, attentional redirection, renewed action, or another event ends the simulation.
\end{itemize}
If both episodes are merely called ``anger,'' or both include one phone check, the evidence supports only content repetition. Structural return requires stronger preservation of relations.

Thus
\[
\boxed{\text{repeated affect label}\not\Rightarrow\text{witness-supported psychological return}.}
\]
If no inspectable past episode exists, ``I am always like this'' remains at the narrative level; the chapter may not supply a return witness on its behalf.

\section{After the Same Return, the Afterwards Can Still Differ}

Even if two episodes eventually support the same psychological return hypothesis, their later continuations need not be the same. After one corrective simulation the participant may have
\begin{center}
\resizebox{0.96\linewidth}{!}{$\displaystyle
C^+(E_{\mathrm{past}})
=
\left\{\begin{gathered}
\text{stronger withdrawal tendency, lower probability of renewed contact,}\\
\text{continued attitude checking}
\end{gathered}\right\},
$}
\end{center}
whereas after a similar current episode,
\begin{center}
\resizebox{0.96\linewidth}{!}{$\displaystyle
C^+(E_{\mathrm{now}})
=
\left\{\begin{gathered}
\text{stop the simulation, return attention to one's own action rule,}\\
\text{allow the unknown to remain unknown}
\end{gathered}\right\}.
$}
\end{center}
These sets are only schematic case-level continuation profiles; an actual study can replace them with probabilities, ratings, or behavioral choices. The key relation is
\[
C^+(E_{\mathrm{past}})\neq C^+(E_{\mathrm{now}}).
\]
``It happened again'' and ``what happens after it this time'' are thereby separated. The mathematical residue $\rho_\xi(R)$ of Part IV supplies a formal prototype for this bookkeeping. The case has not measured mathematical residue and is not entitled to rename continuation difference a trauma remnant, learning trace, or neural mechanism.

\section{Alternative Models, Independent Observables, and Exit Conditions}

Many more mature psychological accounts are available for this case. Norm violation can evoke anger; uncertainty can sustain attention; unfinished goals can remain cognitively accessible; rejection-related appraisal, counterfactual thinking, and emotion regulation may contribute to later simulation. AF-by-discrete language has no priority over them.

The incremental questions of the present analysis can therefore be made falsifiable:

\begingroup\small
\begin{longtable}{>{\raggedright\arraybackslash\bfseries}p{0.22\textwidth}>{\raggedright\arraybackslash}p{0.36\textwidth}>{\raggedright\arraybackslash}p{0.32\textwidth}}
\toprule
Layer & Candidate observable & Result that weakens this analysis\\
\midrule
$Q_A$ action completion & willingness to contact again, actual follow-up, next action choice & recorded action completion persistently disagrees with actual action choice\\
$Q_E/Q_N$ explanation and acknowledgment & confidence about the other's intention; judgment about what information or response is still missing & these judgments show no stable relation to the separate completion ledgers\\
$\sigma_C$ corrective position & imagined-dialogue frequency/duration; whether the participant most wants to change their own action or the other's understanding & simulation does not vary at all with organizing position or completion question\\
return candidate & cross-episode similarity of entrance, position shift, finite process, and exit & only affect names or surface actions are similar; no relational sequence is preserved\\
continuation $C^+$ & later predictions, checking, approach/avoidance, repair choice & the return judgment adds no resolution to later prediction\\
\bottomrule
\end{longtable}
\endgroup

If existing models explain these variables more simply, while the bookkeeping in \(Q\), \(\Pi\), \(\Sigma\), \(\widehat{\mathcal R}_{\mathrm{psych}}\), and \(C^+\) adds no reliable prediction, the AF-by-discrete correspondence is unnecessary in this case.

There are also direct exit conditions: if internal simulation occurs only once and has no identifiable episode endpoints; if a small reasonable change in segmentation makes the return judgment disappear completely; if every candidate interface can be located only by retrospective interpretation; or if the researcher cannot specify in advance what would count as completion, this chapter should stop using return/boundary language.

\section{Translation Remainder: Structural Bookkeeping Does Not Explain Why Anger Feels This Way}

Even if every structural record above holds, many central questions remain unanswered. Why does silence produce this particular bodily feeling? Why does this interaction rule carry this amount of value weight for the participant? How does relational history enter the current appraisal? How do gender, cultural etiquette, and platform norms change the meaning of ``one should reply''? How are anger and bodily response implemented? Why does corrective simulation persist more readily for some people and at some times?

These belong to the translation remainder of the case. The structural model tracks which distinctions are preserved, which questions remain unresolved, and which observation position continues to organize what follows. It does not thereby acquire a complete theory of emotion mechanisms.

\section{Chapter Conclusion: Which Question Actually Needs Completion?}

The most important distinction left by this small case is
\[
\begin{gathered}
\text{action may already have ended;}\\
\text{relational intention may still be unknown;}\\
\text{normative acknowledgment may still be absent;}\\
\text{internal simulation may still continue.}
\end{gathered}
\]
Third-order analysis adds one more question: who is elevating which of these ledgers into the criterion that ``the matter is not over''?

If this bookkeeping holds, a vague ordinary-language ``need for closure'' no longer has to explain the entire material. We can track $Q$, episode segmentation, organizing shifts, candidate interfaces, return evidence, and continuation difference separately. Every step can later be overturned by new facts.

\chapter{An Unfamiliar Talk and the Sudden Blankness of One's Own Field: Attentional Center Is Not Total Knowledge}
\label{chap:unfamiliar-talk-en}

\section{Separate the Three Things Most Easily Confused}

Imagine that a researcher attends a talk in a neighboring but unfamiliar area. In their own work, many examples, counterexamples, and open problems are readily available, but the speaker uses a different vocabulary, a different canon of problems, and different proof habits. After ten or fifteen minutes the researcher is using most of their effort merely to keep track of definitions and is generating almost no questions. Then comes the judgment:
\begin{quote}
``Is my knowledge too narrow? Why can I suddenly remember almost nothing from my own field?''
\end{quote}

The judgment may contain genuine information. The talk may indeed expose a knowledge gap, and the researcher may indeed need broader reading. The case protocol first stops a cross-scale compression. At least three quantities must be separated:
\[
\begin{aligned}
K:&\quad\text{long-term knowledge and skill stock;}\\
A_t:&\quad\text{knowledge and questions currently accessible for the task at time }t;\\
J_t:&\quad\text{the evaluation at time }t\text{ of ``what kind of researcher am I?''}
\end{aligned}
\]
Classic cognitive psychology supplies important background for the first two. Working-memory research emphasizes that complex comprehension requires temporarily maintaining and manipulating limited information \citep{baddeley1992}. Tulving and Pearlstone's experiments on availability and accessibility show that whether stored information can be retrieved at a given moment depends on retrieval cues and organizational conditions \citep{tulving1966}. Those studies did not study the academic identity in this chapter and do not establish a root shift. They supply only a necessary psychological discipline:
\[
\boxed{\text{cannot retrieve it now}\not\Rightarrow\text{did not know it in the longer term}.}
\]
The reverse overreach is also forbidden. Recovering many ideas when familiar cues return does not prove that one's knowledge is sufficiently broad. Accessibility and breadth of knowledge must be measured separately.

\section{Raw Events: Record How the Blankness Develops}

A slower case record is
\[
D=(e_1,\ldots,e_9):
\]
\begin{itemize}
\item $e_1$: enter the unfamiliar vocabulary and problem background of the talk;
\item $e_2$: continuously maintain new definitions and their dependencies;
\item $e_3$: spare attention available for generating one's own questions decreases;
\item $e_4$: lose one local derivation and have to reconnect;
\item $e_5$: notice ``I cannot come up with a question right now'';
\item $e_6$: begin monitoring whether other people follow quickly or ask questions;
\item $e_7$: the question shifts from ``how does this argument go?'' to ``am I broad enough?'';
\item $e_8$: one's own projects, papers, and familiar questions recede further from access;
\item $e_9$: form the global judgment ``my knowledge is narrow.''
\end{itemize}
The record does not write $e_3$ as an already established ``working-memory overload'' mechanism or $e_7$ as a mathematical root. It leaves that identification open on purpose. It first preserves the observable or reportable order: difficulty following, self-monitoring, comparison, self-evaluation, and temporary withdrawal of familiar material.

\begin{plainbox}[What matters first for the psychology reader]
The phenomenon to explain is not merely ``not understanding.'' The more important event is an expansion of scope: how a local access difficulty can, within a short period, acquire interpretive authority over overall competence and professional identity.
\end{plainbox}

\section{Fix the Questions First: Knowledge, Access, and Evaluation Cannot Be Settled with One Answer}

Following the Chapter 30 protocol, declare three questions:
\begin{center}
\resizebox{0.96\linewidth}{!}{$\displaystyle
\begin{aligned}
Q_K:&\quad\text{How much transferable knowledge does the researcher actually possess in the relevant area?}\\
Q_A:&\quad\text{How much material related to their own research can they access in the present context?}\\
Q_J:&\quad\text{How is the current experience being written into a judgment about research identity?}
\end{aligned}
$}
\end{center}
The three can have very different states, for example
\[
Q_K=\text{unknown},
\qquad
Q_A=\text{temporarily low},
\qquad
Q_J=\text{strongly negative}.
\]
There is no contradiction. The feeling of a blank mind first supports a local statement about $Q_A$; its evidential force for $Q_K$ depends on independent knowledge measures; $Q_J$ introduces comparison targets, professional standards, and self-narrative.

Segmentation must also be disclosed. One $\Pi_1$ treats $e_1$--$e_5$ as the task episode of following the talk and $e_6$--$e_9$ as a self-evaluation episode. Another $\Pi_2$ begins at the local loss of the derivation in $e_4$ and treats $e_4$--$e_9$ as an ongoing mismatch-monitoring-evaluation process. These cuts answer different questions. If later root-shift or return judgments hold only under one segmentation, that sensitivity must be reported.

The distinction is important. Third-order observation preserves the evidential path. If independent measures repeatedly show that the researcher genuinely lacks required background, $Q_K$ may receive a severe answer. The severe conclusion must still state which measurements, comparison scales, and task conditions support it. Severity does not exempt a judgment from an evidential path.

\section{Candidate Organizing Shift: From ``Understanding This Talk'' to ``Judging Who I Am''}

The dominant organizing position in the first part of the talk can be written provisionally as
\begin{center}
\resizebox{0.96\linewidth}{!}{$\displaystyle
\sigma_{\mathrm{task}}=\text{``How do I understand the definitions, examples, and argument in front of me?''}
$}
\end{center}
Later, another position may acquire organizing authority:
\begin{center}
\resizebox{0.96\linewidth}{!}{$\displaystyle
\sigma_{\mathrm{self}}=\text{``Where do I stand relative to this field, these listeners, and the standard of a `good researcher'?''}
$}
\end{center}
This is the case-level root-shift candidate. The root proxy is narrow on purpose: it records which reference position currently determines what material matters, which parts of the past become relevant, and which next action appears reasonable. It is not identical with attention and not identical with the whole self-concept.

The two positions generate different continuations. If $\sigma_{\mathrm{task}}$ remains dominant, the next action may be to note a definition, seek an example, or mark a lemma that was lost. If $\sigma_{\mathrm{self}}$ takes over, the next action may instead be to compare oneself with the audience, retrieve evidence of one's knowledge gaps, search for proof that ``I am too narrow,'' or stop following the talk.

The empirical emphasis is therefore on
\begin{center}
\resizebox{0.96\linewidth}{!}{$\displaystyle
\boxed{
\begin{gathered}
\text{after the same local comprehension difficulty passes through different organizing positions,}\\
\text{does }C^+\text{ change systematically?}
\end{gathered}}
$}
\end{center}
If the putative root shift is fully redundant with ordinary attentional shift, task difficulty, or self-evaluation and adds no later prediction, the term need not be retained.

\section{Where Is the Fissure? One Moment of Not Understanding Is Not Yet a Boundary}

A single lost derivation at $e_4$ cannot be renamed a boundary. Such failures are ordinary in unfamiliar talks. Most require only one definition, a glance back at a slide, or temporary acceptance of a claim; the existing task organization can continue.

The more informative candidate is repeated local-sufficiency failure: the original policy of ``keep following the talk'' no longer organizes the next step, so another set of materials becomes relevant---comparison targets, professional standards, earlier failures, and the long-term identity narrative---and behavior changes accordingly. From one episode we can propose at most
\[
b^*:\quad
\text{repair current understanding}
\longrightarrow
\text{evaluate one's overall position}.
\]
Actual boundary detection requires comparable settings. If across several unfamiliar talks the researcher repeatedly shows $\sigma_{\mathrm{task}}\to\sigma_{\mathrm{self}}$ after similar local access failures, and question generation, memory retrieval, and action choice change reliably after that transition, then $b^*$ gradually acquires empirical weight as a singular-boundary-like candidate.

\begin{plainbox}[The fissure must be forced out repeatedly]
First observe where continuation pressure repeatedly gathers, then propose a candidate interface. One episode of confusion supplies only a clue. Organizing shifts that cluster stably across multiple comparable episodes begin to probe a fissure.
\end{plainbox}

A finite singular-interface audit can also be attempted. Fix a finite case coding such as ``follow definition,'' ``request local repair,'' ``recruit own field,'' ``enter self-comparison,'' and ``leave task,'' and ask whether key transitions across talks fall into a stable descriptive vocabulary. This checks type-level coding only; it cannot support a finite-locus claim. Finitely many recurrent types can cover an indefinitely growing number of distinct loci.

The pressure corresponding to the finite singular-interface hypothesis of the book is stronger. After a cross-episode locus-identity rule is frozen, do the concrete persistent singular-like loci keep being reidentified as one finite set? If each new body of material forces the analyst to invent a new type, or if old types still contain an ever-growing set of nonmergeable loci, the empirical translation of singular germs receives no corresponding support. A bounded-type-inspired audit may further ask whether the number of candidate singular-like interface occurrences remains uniformly controlled as episode resolution becomes finer. This remains only a singular-sector proxy; it neither turns psychological data into literal Bratteli boundary points nor counts regular boundary occurrences.

\section{``I Am Like This at Every Unfamiliar Talk'': Return Needs More than a Self-Narrative}

The researcher may then say, ``Whenever I listen to other people's talks, I feel that I know too little.'' This proposes a psychological return hypothesis. It does not yet provide a return witness.

The easiest spurious return is a merger by identity label:
\begin{center}
\resizebox{0.96\linewidth}{!}{$\displaystyle
\text{temporarily losing fluency for different reasons}
\longrightarrow
\text{``I have exposed my narrowness again.''}
$}
\end{center}
One talk may fail because of missing prerequisites, another because of fatigue, one because the pace is extreme, another because the researcher is under private stress. If all are retrospectively renamed ``I am too narrow,'' the narrative label performs the merger first and then uses the merged result as evidence for ``again.''

A stronger return hypothesis should compare finite processes:
\begin{itemize}
\item do they enter through a localizable access difficulty?
\item do they show a similar $\sigma_{\mathrm{task}}\to\sigma_{\mathrm{self}}$ shift?
\item do they recruit similar comparison standards or past material?
\item do they show similar exits, such as abandoning the talk, searching for one's own field, recovering afterward, or continuing self-evaluation?
\end{itemize}
Hence
\[
\boxed{\text{same self-conclusion}\not\Rightarrow\text{same return structure}.}
\]

\section{Restore Access Conditions: A Minimal Counterfactual Test}

This case allows a small counterfactual test. If the blankness of one's own field involves accessibility and organizing-position change, one can alter access conditions before trying to alter the researcher's self-evaluation.

During or after the talk, provide one limited cue related to the researcher's own work: a page from a recent paper, three familiar keywords, or a preprepared question from their field. Then measure again:
\begin{itemize}
\item speed and number of questions generated from the researcher's own field;
\item number of connections made to the unfamiliar talk;
\item intensity of the subjective blankness;
\item whether the global self-evaluation $Q_J$ changes;
\item whether the next action is to keep understanding, seek bridging material, or withdraw into self-comparison.
\end{itemize}
If a small familiar cue restores $A_t$ and question generation substantially while independent knowledge measures remain stable, the earlier inability to retrieve material cannot directly serve as evidence that total knowledge has fallen. If the cue has little effect and stable gaps appear across tasks, the knowledge-deficit explanation gains weight.

The structural comparison is between two paths:
\begin{center}
\resizebox{0.96\linewidth}{!}{$\displaystyle
\text{local access failure}\to\text{interface repair}\to\text{continued understanding or new connection},
$}
\end{center}
and
\begin{center}
\resizebox{0.96\linewidth}{!}{$\displaystyle
\text{local access failure}\to\text{global self-evaluation}\to\text{more comparison and task withdrawal}.
$}
\end{center}
If the structural model adds information, these paths should produce stable differences in later attention, retrieval, and action.

\section{Alternative Models, Independent Observables, and Exit Conditions}

This case lies first within mature psychological domains: working-memory load, retrieval cues, task difficulty, expertise gaps, social comparison, performance anxiety, fatigue, and general self-efficacy can all explain sudden blankness. AF-by-discrete has no explanatory priority.

Its incremental question is narrower: after those variables are controlled, does recording ``local-sufficiency failure $\to$ organizing shift $\to$ continuation change'' still add resolution?

\begingroup\small
\begin{longtable}{>{\raggedright\arraybackslash\bfseries}p{0.22\textwidth}>{\raggedright\arraybackslash}p{0.36\textwidth}>{\raggedright\arraybackslash}p{0.32\textwidth}}
\toprule
Layer & Candidate observable & Result that weakens this analysis\\
\midrule
$Q_K$ long-term knowledge & independent knowledge items, transfer items, delayed tests & blankness tracks stable knowledge deficit almost completely\\
$Q_A$ current accessibility & retrieval after familiar cues, question generation, counterexample generation & cue/context reinstatement has no systematic effect\\
$\sigma$ organizing position & task/self focus, comparison target, later retrieved content and action choice & the shift is fully explained by task difficulty or attention\\
candidate interface & clustering of organizing shifts and continuation change across talks & turning positions cannot be repeatedly localized and must be assigned post hoc case by case\\
return candidate & cross-episode relations among entrance, shift, recruited material, and exit & only the label ``I am too narrow'' repeats; finite process does not\\
\bottomrule
\end{longtable}
\endgroup

A direct exit condition remains: if across reasonable segmentations, familiar cues, and comparison frames the alleged root shift changes neither retrieval, question generation, action choice, nor later prediction, the chapter should stop using AF-by-discrete language. ``Performance decline on an unfamiliar task plus negative self-evaluation'' is then sufficient.

\section{Translation Remainder: Structural Analysis Does Not Explain Why ``Breadth'' Matters So Much to a Researcher}

Even if the structural analysis is correct, it does not explain why a particular person puts breadth of knowledge at the center of professional value; how advisor culture, disciplinary institutions, and conference norms define what a ``good researcher'' should understand; how competition, age, position, and peer comparison change the weight of blankness; how fatigue and anxiety shape the experience; or why some fields genuinely lack transferable interfaces.

These remain in the translation remainder. More precise recording of organizing position does not automatically yield a social-psychological mechanism, educational mechanism, or neural implementation.

\section{Chapter Conclusion: Blankness Has a Position, and Judgment Has a Position}

The case leaves four statements:
\begin{center}
\resizebox{0.96\linewidth}{!}{$\displaystyle
\begin{gathered}
\text{one episode of blankness does not prove that long-term knowledge has been globally rewritten;}\\
\text{currently accessible knowledge can change sharply with task, cue, and attention conditions;}\\
\text{local access difficulty can be compressed into a judgment about the whole self;}\\
\text{Third Order requires the position and path of that compression to remain in the record.}
\end{gathered}
$}
\end{center}
``I cannot remember it right now'' is first of all a local fact. Do not cancel it with ``actually you know everything,'' and do not immediately upgrade it to ``I objectively discovered that I am narrow.'' The stricter task is to measure $K,A_t,J_t$ separately, track whether $\sigma_{\mathrm{task}}\to\sigma_{\mathrm{self}}$ recurs, and observe which fissures are genuinely forced out by repeated tasks.

That is the contribution of this case to Part V: a judgment about ability regains an occurrence position, finite evidence, and exit conditions.

\chapter{Repeated Checking, Rumination, and Overlapping Returns: Repetitive Thought Is Not One Return}
\label{chap:overlapping-returns-cases-en}

\section{Return Psychological Terms to Psychology First}

Research on repetitive thought and rumination has long warned that ``thinking repeatedly'' is not one homogeneous process. Watkins emphasizes that the consequences of repetitive thought vary with content, level of abstraction, context, and processing mode; the same recurrence of thought can be constructive or become rigid and depleting \citep{watkins2008}. Nolen-Hoeksema, Wisco, and Lyubomirsky likewise review the complex relations among rumination, negative affect, problem solving, and social functioning \citep{nolen2008}.

This book does not redefine those psychological constructs as mathematical return. The first ledger separation is
\begin{center}
\resizebox{0.96\linewidth}{!}{$\displaystyle
\boxed{\text{repetitive thought is a class of experiential phenomena; return is a finite completion relation in this model}.}
$}
\end{center}
The two may be structurally compatible in some cases, but they are not synonyms. A thought occurring ten times does not by itself establish a return; a clear return does not require the same sentence to repeat mechanically ten times.

We therefore do not count only how many times a thought occurs. Instead, we track three narrower questions: when the sequence is re-entered, is the current question the same? Is the organizing position the same? After return-completion, are the available next steps the same?

\section{A Checking Sequence: The Action Repeats, but the Entrance Has Already Changed}

Imagine that a person is waiting for an important email:
\[
\begin{array}{ll}
C_1:&\text{open email--no new message--close it;}\\
C_2:&\text{open messaging app--no new message--close it;}\\
C_3:&\text{open email again--no new message--close it;}\\
C_4:&\text{begin checking whether the previous email was written incorrectly;}\\
C_5:&\text{open email again.}
\end{array}
\]
If classified only by the action ``open email,'' $C_1,C_3,C_5$ look nearly identical. The case protocol requires the outward action to be separated from the organizing question.

The entrance to $C_1$ may be simply
\[
Q_1=\text{``Is there a new message?''}
\]
By $C_3$, the question may already include temporal evaluation:
\[
Q_3=\text{``It has been this long; why is there still nothing?''}
\]
After the self-check in $C_4$, $C_5$ may begin from
\[
Q_5=\text{``Did I write something wrong, and is that why there is no reply?''}
\]
The same action can therefore be entered from different positions and lead to different afterwards: continuing work, continuing to wait, self-blame, rereading old mail, rehearsing a repair message, or withdrawing for a while.

The minimum discipline is
\begin{center}
\resizebox{0.96\linewidth}{!}{$\displaystyle
\text{action repetition}
\neq
\text{question repetition}
\neq
\text{return repetition}
\neq
\text{continuation repetition}.
$}
\end{center}
If the researcher draws one ``big cycle'' merely because the behavior looks the same, the structural changes that need explanation have already been erased by the diagram.

\section{Several ``Agains'' Can Overlap: Multiple Return Hypotheses May Run through One Thought Sequence}

A single thought sequence may support several case-level return hypotheses simultaneously. For bookkeeping, write provisionally
\[
\widehat R_{\mathrm{check}},
\qquad
\widehat R_{\mathrm{self\text{-}blame}},
\qquad
\widehat R_{\mathrm{future\text{-}simulation}}.
\]
The hats matter: these remain psychological case-level candidate classifications, not the formal return language defined in Part IV.

A longer sequence may simultaneously contain:
\begin{itemize}
\item a checking return: ``open--nothing--close--open again'';
\item a self-evaluation return: ``no response--suspect myself--search for an error'';
\item a future-simulation return: ``imagine how the other person might reply--rehearse my next line.''
\end{itemize}
The hypotheses can share events. Opening the inbox again may complete a local checking return while serving only as an intermediate step in the self-blame sequence.

Returns can also nest. A longer $\widehat R_{\mathrm{check}}$ can contain a shorter $\widehat R_{\mathrm{self\text{-}blame}}$ that ends while the former continues; two return intervals may instead overlap only partially. The multiple-boundary/multiple-bridge-type construction in Subsection 5.2 of Kuang's v1 preprint gives a mathematical prototype for such heterogeneous nesting: different boundary sectors can witness different return families, and one high-level trajectory can pass successively through several sectors \citep{kuang2024v1}.

\begin{warning}[Do Not Diagnose Nesting as Multiple Germs]
Nested psychological returns do not imply that the brain or subject ``has multiple boundary germs.'' Even with one fixed germ type, mathematical returns can nest at different scales. This chapter retains only a structural permission: one experiential trajectory may belong to several return intervals organized at different positions and scales.
\end{warning}

\section{The ``Third-Order Consciousness Trap'': When ``the Observer Is Also in the Structure'' Becomes a Tool for Solving the Problem}

Let us consider a subtler return, one that may occur precisely in someone who has read this book. Imagine that anxiety, jealousy, anger after rejection, or persistent rumination comes back. The person immediately remembers the language of Third-Order Consciousness:
\begin{quote}
``The observer is also inside the structure. Good. I will observe the observer correctly now and solve this problem.''
\end{quote}
``The observer is also inside the structure'' is of course not a mathematical theorem; it is a structural insight developed in Part I. The danger begins when it is rewritten as an operational command.

Slowed down, the process may look like
\begin{center}
\resizebox{0.96\linewidth}{!}{$\displaystyle
\begin{array}{ll}
T_1:&\text{an uncomfortable psychological state appears;}\\
T_2:&\text{remember ``the observer is also in the structure'';}\\
T_3:&\text{demand that oneself ``observe the observer correctly'';}\\
T_4:&\text{check whether anxiety, anger, or rumination has decreased;}\\
T_5:&\text{find that the problem remains and judge ``I have not really seen it yet'';}\\
T_6:&\text{invoke the same third-order language again.}
\end{array}
$}
\end{center}
The language becomes more ``meta,'' yet structurally a familiar second-order action may have returned: one subject-position sets the required outcome and monitors whether present experience conforms to it.

What is exempted from further observation is precisely the sentence
\[
\boxed{\text{``I must solve this psychological problem.''}}
\]
It defines success, decides what must disappear, judges whether observation has been good enough, and recruits ``the observer is also in the structure'' as a means to reach a goal. The old observer has been verbally put ``inside,'' while a new solver quietly regains an outside privilege.

This is what the book calls the \emph{Third-Order Consciousness trap}:
\begin{quote}
Learning third-order language does not guarantee third-order organization. When ``Third Order'' is used by an unobserved solver to manage, correct, or eliminate present experience, the vocabulary may remain third-order while the organization has been retaken by a second-order subject.
\end{quote}

\begin{plainbox}[Conceptual distinction: wanting relief and requiring a solution are not the same]
Two easily confused forms of wanting must be separated.

The first is a \emph{directional wish / wellbeing orientation}:
\begin{center}
\resizebox{0.96\linewidth}{!}{$\displaystyle
W:\quad\text{``Of course I hope to suffer less unnecessarily and to be recruited less by old structures.''}
$}
\end{center}
It records direction and valence. Pain motivates leaving pain; sustained anxiety motivates restoration of stability; ordinary agents want relationships, work, and bodily states to become more bearable. Third Order does not require this tendency to disappear.

The second is a \emph{corrective agenda}:
\begin{center}
\resizebox{0.96\linewidth}{!}{$\displaystyle
P:\quad\text{``This state must now be solved by me through some psychological operation.''}
$}
\end{center}
$P$ adds something to $W$: it prescribes what present experience must become, defines successful observation, and recruits observation as an instrument for producing the result. The problem is not that a wish exists. The problem begins when the wish acquires structural privilege, evaluating other experience while excluding its own urgency, fear, and success criterion from the field of experience.
\end{plainbox}

This distinction resonates with Krishnamurti's discussions of sorrow. His pressure test is not ``find a better method to remove suffering quickly,'' but whether one can remain with the fact without religious, secular, or philosophical escape, even temporarily holding the question without rushing to solve it, so that suffering, loneliness, shock, and the movement away from them become visible together; in his own thought, such non-escape is linked to release \citep{krishnamurtiSorrow2026}. Part V preserves this intellectual direction without promoting ``staying with sorrow necessarily produces release'' into an experimentally established psychological mechanism.

The minimum practical boundary can be written
\[
\boxed{\text{I may hope it ends, but I do not require observation to make it end.}}
\]
This cancels no action. One can leave danger, seek treatment, rest, refuse harm, and solve real-world tasks. It cancels only a hidden performance contract: ``if I observe correctly, I should feel better now.''

Immediate comfort therefore cannot be the sole KPI of third-order observation. A more relevant structural question is whether the old organization still monopolizes continuation automatically. When the same discomfort returns, must the system still run along ``name--ruminate--check--self-evaluate--ruminate again,'' or can discomfort, the wish for relief, the urgency to solve it, invocation of the theory, and checking of its effectiveness all appear together without one of them owning the next step?
\[
\boxed{\text{freedom from structural capture}\neq\text{immediate comfort}.}
\]
Comfort can be a long-term orientation or an outcome; it does not issue certificates of success to present observation.

Third Order therefore does not require ``giving up problem solving.'' It requires that the problem, the wish for relief, the urgency to make it vanish now, the act of invoking theory, the effectiveness check, and the self-evaluation ``I should know how to do this by now'' all be allowed into the same record. No layer acquires final authority merely because its name is observer, theory, or solver.

Nor is a ``fourth-order observer'' needed to monitor the third order. That would only copy the same problem one level higher. The requirement is simpler: the wish, purpose, urgency, and evaluation that were previously left outside the structure are now counted as part of it. Third Order removes structural privilege from the observer; it does not remove human agency.

In the return language of this chapter, the trap exposes an important meta-return:
\begin{center}
\resizebox{0.96\linewidth}{!}{$\displaystyle
\begin{gathered}
\widehat R_{\mathrm{meta\text{-}fix}}:\quad
\text{discomfort}\to\text{invoke third-order concept}\to\text{check whether it worked}\\
\to\text{judge ``not solved''}\to\text{invoke the concept again}.
\end{gathered}
$}
\end{center}
A person can ruminate with a theory of ``not ruminating'' and check with a method of ``not checking.'' More sophisticated concepts do not automatically remove return; concepts themselves can enter the return grammar.

This brings the Part-I pressure test of choiceless awareness back to an ordinary case. \emph{Third-Order Consciousness} and Krishnamurti's related discussions emphasize that when observation serves from its first second the project ``I must become a certain result,'' the old managerial center can survive under new language \citep{toc2026,krishnamurti1954,krishnamurti1969}. Part V does not treat that claim directly as a therapeutic mechanism. It turns it into an observable question: when a person invokes third-order language, does the invocation itself form a new monitoring, evaluation, and return?

\section{A Multi-Boundary-Style Audit: Some ``Agains'' Become Visible Only after Recording Sector Passage}

The Third-Order Consciousness trap also shows why a case may require organizing-position passage rather than one ``fundamental cycle.'' Using the weak case notation from Chapter 30, write
\[
\Sigma=\{\sigma_{\mathrm{check}},\sigma_{\mathrm{self}},\sigma_{\mathrm{rule}},\sigma_{\mathrm{meta}}\},
\]
where
\begin{itemize}
\item $\sigma_{\mathrm{check}}$ organizes the search for new information;
\item $\sigma_{\mathrm{self}}$ organizes ``what does this say about who I am?'';
\item $\sigma_{\mathrm{rule}}$ organizes ``what should I do now?'';
\item $\sigma_{\mathrm{meta}}$ organizes ``am I observing correctly, using the theory correctly, and succeeding in changing the state?''
\end{itemize}
An episode may pass through
\[
\sigma_{\mathrm{check}}
\to
\sigma_{\mathrm{self}}
\to
\sigma_{\mathrm{rule}}
\to
\sigma_{\mathrm{meta}}
\to
\sigma_{\mathrm{check}}.
\]
If return-completions and later changes across comparable episodes reliably depend on this sector passage, recording the passage may offer more resolution than a one-loop model.

The relation borrowed from the bridge-type word in Kuang v1, Subsection 5.2, is that a high-level trajectory can pass successively through different boundary types and finite coding need not be concentrated at one point \citep{kuang2024v1}. The psychological side contains no literal bridges and does not rename $\sigma_{\mathrm{meta}}$ a singularity. It tests only a weaker question: does recording organizing-position passage predict later checking, self-blame, simulation, theory invocation, or exit better than a single dominant-loop model?

The boundary-detection discipline of Chapter 30 still applies. Merely possessing the label $\sigma_{\mathrm{meta}}$ does not authorize the researcher to declare every ``I should observe myself'' a boundary-like interface. Stronger evidence requires repeated operation: similar position passages cluster across episodes, old continuation repeatedly loses sufficiency near comparable positions, and later options change stably across the transition. The fissure is still forced out by the system.

If sector passage adds no prediction, delete the sector analysis. If ordinary rumination, metacognitive monitoring, or decentering variables explain the differences already, the additional AF-by-discrete translation is unnecessary.

\section{Why ``I Cannot Stop'' Still Does Not Mean Incompressible}

Ordinary language says ``I cannot stop thinking about this.'' Mathematical incompressibility has a precise meaning: relative to a declared formal return language, no return factor occurs. Psychological persistence, distress, number of repetitions, and formal incompressibility therefore answer different questions.

A highly repetitive checking sequence can be very compressible under a clear coding. The meta-return $\widehat R_{\mathrm{meta\text{-}fix}}$ may have an especially obvious repetitive skeleton. Conversely, the fact that the content changes slightly each time does not grant real thought the mathematical property of incompressibility.

The separation remains
\[
\boxed{\text{psychological persistence}\neq\text{formal incompressibility}.}
\]
Only after psychological variables have been restricted to a predeclared formal model capable of failing does it make sense to ask whether they avoid the return factors declared by that model. ``Incompressible'' is not a new label for ``this person is too fixated'' or ``this problem is too deep.''

\section{Empirical Bridge: How Could the New Meta-Return Be Tested?}

The chapter now yields two levels of empirical question.

The first concerns ordinary repetitive thought. Do participants spontaneously identify some thought episodes as ``the same thing starting again''? Does that judgment depend on action similarity, semantic similarity, an unresolved question, or a shared organizing position? If return-based segmentation adds information, it should improve prediction of later sequence structure, not merely rename existing rumination.

The second comes directly from the Third-Order Consciousness trap. Compare two instruction framings. Both groups learn the same proposition that an observer is organized by past material, rules, and evaluative positions. One group is instructed to use this principle to reduce present discomfort as quickly as possible. The other is instructed only to record discomfort, the impulse to make it disappear, invocation of the theory, and checking of its effects together, without defining ``discomfort decreases'' in advance as the criterion of successful observation. Then record:
\begin{itemize}
\item number of theory re-invocations;
\item self-monitoring of ``am I observing correctly?'';
\item duration and reopening rate of checking/rumination episodes;
\item organizing-position passage;
\item whether the available action set narrows or reopens.
\end{itemize}
The goal is not to assume in advance that the second instruction is ``higher'' or more effective. The actual question is whether explicitly corrective goal framing makes a repeatable meta-monitoring loop more likely, and whether that difference exceeds what general negative affect, rumination, metacognition, and decentering already explain \citep{watkins2008,fleming2012,bernstein2015}.

If no additional difference appears, the Third-Order Consciousness trap remains an evocative philosophical description and should not be upgraded into a psychological mechanism. If a stable difference does appear, only then is there reason to investigate whether it comes from goal framing, self-evaluation, attention allocation, or the observer-position organization emphasized here.

The final question is therefore narrower than ``how do we stop rumination?''
\begin{quote}
\textbf{When ``I am solving my rumination'' begins to repeat, can that new ``again'' be measured separately?}
\end{quote}

\chapter{Narrative Overcompression: When Does ``I Am Always Like This'' Merge Histories That Were Structurally Different?}
\label{chap:narrative-overcompression-en}

\section{Why Narrative Is Necessary: Compression Is Not Itself an Error}

Research on autobiographical memory and narrative identity emphasizes that the past is not a complete recording. It is continually selected, reconstructed, and organized into present self-understanding. Conway and Pleydell-Pearce's self-memory system locates autobiographical recollection in the interaction between self-goals and knowledge structures \citep{conway2000}; Bruner treats narrative not merely as a record but as a major form through which life becomes meaningful \citep{bruner1987}; McAdams and McLean discuss how life stories provide unity and purpose \citep{mcadams2013}.

We therefore never treat narrative compression itself as error. Without compression, a life cannot be remembered or made into an actionable history. The real question is what the compression deletes and whether the deleted differences would change future judgment, action, or continuation.

The minimum psychological motivation is simple: people need generalization, but generalization also needs an explicit domain of validity. A long-term narrative can be true while remaining an organization of some material, some relations, and some scale. Truth does not automatically grant it the right to represent the whole.

\section{``The Projection Is Real, but It Is Not the Whole'': How a Local Judgment Overreaches}

In a discussion of the public cycle of idolization and destruction, \emph{Snowfall on the Clock} offers a formulation useful here: the projection is real, but the projection is not the whole. Excellent performance in one concrete setting can be real; a body of material demanding accountability can also be real. The structural error comes in the next step, when one local manifestation takes possession of the entire person: ``she did very well here'' becomes ``this is what she is,'' or ``there are serious problems here'' becomes ``everything about her was false'' \citep{projection2026}.

The point is not the public figure. It is the more general psychological relation
\[
\boxed{\text{local projection}\not\Rightarrow\text{global identity}.}
\]
A local judgment can be entirely valid. Overreach begins when it erases its own domain of validity.

The same organization appears in self-narrative. Being unanswered at one meeting can be a real disappointment; silence in one relationship can be a real injury; one work mistake can be a real failure. But to continue from these finite facts to
\begin{quote}
``I am always ignored.''\qquad ``I am the kind of person who is never taken seriously.''
\end{quote}
performs an additional global completion. Third Order first requires that this act of completion itself remain visible: from which position, using which episodes, did someone fill how much unobserved space into a whole person?

\begin{principle}[Projection Locality Principle: A Local Projection Cannot Substitute for the Global]
A root, observation position, or local narrative may organize current material truthfully, yet cannot acquire an explanatory monopoly over the global architecture merely through its own salience. Part V therefore records separately whether a judgment is true and how large its valid domain is:
\[
\boxed{\text{local truth}\neq\text{global occupancy right}.}
\]
If new material requires different positions, paths, relations, or time scales to remain visible, the old projection must permit itself to be unfolded. It may not reinterpret every difference as still deeper evidence of the same essence.
\end{principle}

\section{``I Am Always Ignored'': One Sentence Can Compress Four Different Histories}

Let us consider four episodes:
\begin{enumerate}
\item a private message receives no reply;
\item a collaborator replies to a work email three days later;
\item a contribution in a meeting is not taken up by the chair;
\item a family member, while busy, fails to notice a question.
\end{enumerate}
One narrative may compress all four into ``I am always ignored.'' The sentence may accurately express a long-term feeling, yet it does not establish that the four episodes share one return structure. Their entrances, participants, norms, time scales, and available exits may differ.

The sentence also comes from a current observation position. Suppose $\sigma_{\mathrm{self}}$ is asking ``what do these events jointly say about who I am?'' Before being projected into this position, the four episodes may have belonged to work coordination, intimacy, group attention, and family resource allocation. Their projection into $\sigma_{\mathrm{self}}$ may make a meaningful summary. But if the summary is then treated as all the facts, a local root begins to impersonate the global.

The chapter does not require the sentence ``I am always ignored'' to be forbidden. It requires three audit questions to remain beside it:
\[
\begin{gathered}
\text{Which episodes were included?}\\
\text{Which differences were deleted?}\\
\text{If those differences are restored, does the next continuation change?}
\end{gathered}
\]

\section{Good Telescoping and Bad Quotient: Good Compression Must Permit Re-Expansion}

Part III provisionally mapped narrative to a telescoped presentation. In Part V the analogy becomes a practical audit:
\begin{quote}
If four episodes are compressed into one macro-pattern, which intermediate differences are removed? Would those differences change a relevant next prediction or action?
\end{quote}
If the removed variables do not change relevant continuation, the compression may be good. If it merges future-wise very different situations into one identity judgment, it may be overcompression.

A case-level heuristic is
\begin{center}
\resizebox{0.96\linewidth}{!}{$\displaystyle
\boxed{\text{good compression permits re-expansion that restores finite differences governing continuation}.}
$}
\end{center}
Re-expansion matters. A good narrative allows itself to be revised as new material enters. A bad quotient often behaves differently: once an identity label is formed, counterexamples are reinterpreted as even deeper proof of the label. Compression has then changed from a working representation into a final court with no exit.

\section{Single- and Multi-Singular-Boundary Candidates: One Salient Interface Cannot Own Every Fissure}

The projection/global distinction supplies a natural formal motivation, but root and boundary must remain separate. Observation position is closer to root/rooted projection; mathematical boundary governs gluing across pieces; and the singular-boundary-like interface activated on the psychological side of this chapter additionally requires germ-sensitive distinction. Therefore
\[
\text{one observation position}\neq\text{one boundary point}.
\]
They perform different structural tasks.

What can be compared is a shared discipline: a local object cannot monopolize global explanation merely because it is salient. A single-singular-boundary prototype may fit current psychological material: many different episodes and positions may genuinely undergo key reorganizations near one reidentifiable germ-sensitive locus. Then one singular-like locus may organize the continuation differences relevant to the present translation. That does not mean the full mathematical system has only one boundary point, and still less that the subject has one state or one viewpoint.

But suppose repeated cases expose several irreducible transition loci, for example
\begin{center}
\resizebox{0.96\linewidth}{!}{$\displaystyle
\sigma_{\mathrm{scholar}}\leftrightarrow\sigma_{\mathrm{self}}
\text{ repeatedly near }b_1,
\qquad
\sigma_{\mathrm{partner}}\leftrightarrow\sigma_{\mathrm{abandonment}}
\text{ repeatedly near }b_2,
$}
\end{center}
with different return families, sector passages, or stable continuation consequences near $b_1$ and $b_2$. Forcing them back into one ``real core fissure'' would repeat the global-completion error under critique.

Two opposite overreaches must therefore be blocked:
\[
\begin{gathered}
\text{one salient boundary does not automatically become}\\[-1mm]
\text{``the root of everything'';}\\
\text{many local projections do not automatically imply}\\[-1mm]
\text{many boundaries.}
\end{gathered}
\]
In particular,
\[
\boxed{\begin{gathered}\text{many projections}\\[-1mm]\not\Rightarrow\\[-1mm]\text{many boundaries}\end{gathered}}.
\]
What decides between single- and multi-singular-boundary candidates is not how many ``sides'' a person has. It is the endogenous boundary-detection rule from Chapter 30: once analysis rules and locus-identity rules are frozen, how many irreducible, reidentifiable germ-sensitive loci does the system repeatedly force into view?

The structural compatibility can be summarized as
\begin{center}
\resizebox{0.96\linewidth}{!}{$\displaystyle
\boxed{
\begin{gathered}
\text{the projection principle forbids one position from owning the whole;}\\
\text{the multi-boundary principle forbids one interface from monopolizing all fissures.}
\end{gathered}}
$}
\end{center}
This does not make them one mathematical theorem. It identifies a shared refusal to let one local center finalize the global structure.

\section{Spurious Return and Missed Return: Narrative Can Over-Recognize and Under-Recognize}

Narrative overcompression can manufacture a spurious return: structurally different episodes are forced into one label. It can also manufacture a missed return: two episodes with very different surface content are not recognized as sharing an organizing sequence because they lack a common affect or character label.

For example, ``the other person does not reply--I imagine educating them'' in a romantic context and ``a collaborator does not adopt my suggestion--I repeatedly compose a corrective email in my head'' in academic work differ in content domain but may share the finite organization
\begin{center}
\resizebox{0.96\linewidth}{!}{$\displaystyle
\text{expectation not completed}
\to
\text{normative evaluation}
\to
\text{corrective simulation}
\to
\text{simulation reactivates affect}.
$}
\end{center}
This remains only a psychological return candidate. Stronger support must return to concrete finite-episode evidence; the formal witness of Part IV is not pasted directly onto lived material.

A researcher convinced that ``it is fundamentally all about being ignored'' can over-recognize structurally different events as returns. A researcher convinced that every relational setting is independent can miss returns that cross content and people while preserving the same organizing sequence. Narrative can overmerge and over-separate. Return audit restores the requirement of finite evidence rather than obedience to an already completed personality story.

\section{Third Order Is Not a Better Story and Not a Higher Total View}

The third-order move is not to replace ``I am always ignored'' with a more positive sentence such as ``actually everyone values me.'' That is narrative competition. Nor is it to replace the old story with a larger one such as ``all of this is my core trauma'' and allow the new story a higher final authority.

Actual untelescoping reopens the material. What happened each time? Where did the episode begin? From which position does the present judgment speak? What makes this episode enter the old pattern? Which differences change action? Which paths remain uncovered by the present story?

This is consistent with the account of anti-narrative capacity in the Geometry of Existence: not refusing stories, but being able to decompose a complete story back into facts, positions, paths, and remainders \citep{geometry2026}. The further discussion of projection adds another discipline: mature judgment does not require evaluation to stop; it requires each local judgment to remain within its evidential range rather than immediately compress several true local facts into a final answer to ``what is this person really?'' \citep{projection2026}.

The chapter therefore ends not with ``do not generalize,'' but with a stricter question:
\begin{quote}
\textbf{Does a narrative still allow new positions, paths, and fissures to enter, or has it begun to finalize the global structure in advance?}
\end{quote}

\chapter{The Same Return, Different Afters: What Residue Means at the Case Level}
\label{chap:case-residue-en}

\section{First Correct a Psychological Misunderstanding: Change Does Not Mean Never Being Triggered Again}

People often judge change by an extremely severe standard:
\begin{quote}
``If the same problem happens again, I have not changed at all.''
\end{quote}
An opposite narrative says:
\begin{quote}
``If I handled it differently this time, it cannot be the same old pattern.''
\end{quote}
The sentences look opposed, but share one assumption: return and change are mutually exclusive.

That is the psychological problem of this chapter. Familiar bodily tension may recur; ``am I being rejected again?'' may recur; even a similar defensive interaction may complete again. At the same time, the available actions, predictions about others, retrieved past material, and self-evaluation after the episode may have changed.

Part V therefore preserves an ordinary-language statement:
\begin{center}
\resizebox{0.96\linewidth}{!}{$\displaystyle
\boxed{\text{an old organization can complete a return again without monopolizing the future again}.}
$}
\end{center}
The return-residue calculus of Part IV gives this two-ledger distinction a strict mathematical prototype. At the case level, however, we use only the weaker $C^+(E)$ and never write real experience directly as $\rho_\xi(R)$. The formal definition and well-definedness of residue remain in Chapter 27.

\section{A Controlled Case: The Same Conflict Completion, Then What Happens after the Pause?}

Let us consider two collaboration episodes separated by several months. We record only the raw events first.

\textbf{Episode A}
\begin{center}
\resizebox{0.96\linewidth}{!}{$\displaystyle
\begin{gathered}
\text{make suggestion}\to\text{suggestion not adopted}\to\text{defend immediately}\to\text{voice rises}\to\text{discussion pauses}\to\\
\text{repeatedly think ``they do not respect me''}\to\text{reduce communication and rehearse the next counterattack}.
\end{gathered}
$}
\end{center}

\textbf{Episode B}
\begin{center}
\resizebox{0.96\linewidth}{!}{$\displaystyle
\begin{gathered}
\text{make suggestion}\to\text{suggestion not adopted}\to\text{defend immediately}\to\text{voice rises}\to\text{discussion pauses}\to\\
\text{remain uncomfortable but separate ``my idea was not accepted'' from ``I was rejected''}\to\\
\text{ask about the concrete disagreement the next day and continue collaborating}.
\end{gathered}
$}
\end{center}
The first five steps are indeed very similar. But Chapter 30 prohibits declaring return merely because two sequences look similar. The researcher must first disclose the current question $Q$ and the completion coding.

Suppose the case predeclares
\begin{center}
\resizebox{0.96\linewidth}{!}{$\displaystyle
Q=\text{``Do both segments complete the finite sequence `suggestion blocked--defense escalates--pause'?''}
$}
\end{center}
and codes completion from the first five steps only. If independent coding supports the same finite completion in both episodes, only then do we propose the case-level psychological return hypothesis
\[
\operatorname{CompletionCode}(E_A)=\operatorname{CompletionCode}(E_B).
\]
Now inspect what follows the pause. The important result is not the value judgment ``B is more mature than A,'' but the structural fact
\[
C^+(E_A)\neq C^+(E_B).
\]
Both equations must remain. The second does not revoke the first, and the first does not erase the second.

\begin{plainbox}[Why this is not ``telling a more positive story'']
Episode B does not require ``they did not respect me'' to be forcibly rewritten as ``actually they respect me very much.'' The material may still support the judgment that the other person's communication was poor. What changes is that the same finite conflict completion no longer automatically determines later prediction, self-evaluation, and action set. Here we study freedom in continuation, not positive thinking.
\end{plainbox}

\section{What Does \texorpdfstring{$C^+$}{C+} Record? It Is Not a Mysterious ``Psychological Residue''}

The word residue can easily suggest ``how much emotion remains after the event.'' Part V blocks that reading on purpose.

At the case level, $C^+(E)$ is only a continuation profile. For recording purposes one may provisionally write
\[
C^+(E)=(A_E,P_E,M_E,S_E),
\]
where
\begin{itemize}
\item $A_E$: actions still experienced as available;
\item $P_E$: predictions about the next event, the other person's response, or the outcome;
\item $M_E$: past material or comparison sets more readily retrieved afterward;
\item $S_E$: the self/sector position that subsequently acquires organizing authority, such as continuing the task, withdrawing, correcting, shame, self-judgment, or repair.
\end{itemize}
These are a case-coding bundle, not four hypothesized brain modules.

Episode A and B can then be described plainly. After A, the action set narrows, hostile interpretation becomes more likely, ``not being respected'' memories are preferentially retrieved, and a defensive self-position continues to dominate. After B, discomfort remains, but asking, delaying judgment, and continuing the collaboration remain available.

Thus
\[
\boxed{\text{affect still present}\neq\text{continuation still filled by the old structure}.}
\]
A person can feel very bad while the future has reopened. Or the person can feel calm quickly and then follow exactly the same avoidance, checking, or control path next time. ``How comfortable do you feel?'' alone cannot support the distinction this chapter needs.

\section{Put the ``Third-Order Consciousness Trap'' Back into the Two-Coordinate Picture}

Chapter 33 distinguished the directional wish
\[
W:\quad\text{``I hope to suffer less unnecessarily.''}
\]
from the corrective agenda
\[
\begin{aligned}
P:\quad &\text{``This state must now be solved}\\[-1mm]
&\text{through some psychological operation.''}
\end{aligned}
\]
We now show why this distinction does not require old discomfort never to reappear.

Suppose two episodes both contain the same familiar discomfort: chest tension, repeated thoughts about the event, and ``I am doing this again.'' Under one continuation, the system immediately enters
\begin{center}
\resizebox{0.96\linewidth}{!}{$\displaystyle
\text{invoke third-order concept}
\to
\text{check whether it works}
\to
\text{still uncomfortable}
\to
\text{invoke it again}.
$}
\end{center}
Under another, discomfort and the wish to feel better both remain, but they do not automatically trigger effectiveness checking. The person can eat, work, leave a harmful setting, seek help, or temporarily solve nothing.

At a coarse level, both episodes may support the same psychological return: familiar discomfort has been recognized again. Their continuations differ. The first is still captured by $\widehat R_{\mathrm{meta\text{-}fix}}$; the second does not make immediate comfort a gatekeeper for successful observation.

The empirical meaning is therefore
\[
\boxed{\text{freedom from structural capture}\neq\text{absence of recurrence}.}
\]
The return of an old affect, thought, or first action does not decide change by itself. The more informative question is whether, once it returns, it can still write the rest of the path automatically.

Third Order therefore does not need the fiction ``I am no longer triggered.'' If real recurrence is denied in order to protect a new identity---``I have already entered Third Order''---third-order language begins serving a self-image rather than the facts.

\section{Different Afters Do Not Automatically Mean an Additional Singular Boundary}

The previous chapter discussed single- and multi-singular-boundary candidates. Add the reverse discipline here:
\[
C^+(E_1)\neq C^+(E_2)
\quad\not\Rightarrow\quad
\text{there are two distinct boundary points}.
\]
Boundary asks where local organization repeatedly loses continuation sufficiency and requires an additional interface. $C^+$ asks which differences still affect later continuation after return-completion. They are different coordinates.

Different continuations can occur near one stable singular-like interface. Two episodes may both undergo germ-sensitive pressure at the transition ``suggestion rejected $\to$ self-defense,'' yet differ in $C^+$ because of relationship history, present task, acquired information, or subsequent sector passage. That remains compatible with a single-singular-boundary prototype.

A finite multi-singular-boundary candidate gains extra support only if boundary detection across multiple comparable episodes shows that different continuation families depend stably on several irreducible, reidentifiable germ-sensitive transition loci, and merging those loci loses return, sector-passage, or predictive differences.

The three anti-overreach rules are therefore
\[
\begin{gathered}
\text{one projection does not substitute for the global;}\\
\text{one salient boundary does not explain every fissure;}\\
\text{one continuation difference does not manufacture}\\[-1mm]
\text{a boundary out of nothing.}
\end{gathered}
\]
They prevent the theory from adding a new interface every time a new difference is observed until AF-by-discrete becomes unfalsifiable.

\section{The ``Spiral'' Can Be Retained, but It Is Not a Meter of Progress}

The spiral remains a useful image. Viewed from above, the trajectory circles through similar directions; with another coordinate added, each return need not land at exactly the same place. Mathematically, a graph-of-germs lift allows a trajectory to be closed downstairs and displaced upstairs; Chapter 20 explained that relational skeleton.

The psychological side retains only
\begin{center}
\resizebox{0.96\linewidth}{!}{$\displaystyle
\boxed{\text{same recognized return-completion} + \text{different future-relevant continuation}.}
$}
\end{center}
But a spiral cannot automatically mean ``I rise to a higher level each time.'' It can narrow downward: after the same return, available actions become fewer, predictions more rigid, and repair more difficult. It may have no stable direction at all; different contexts may simply yield different continuations.

We therefore refuse to turn spiral height into psychological maturity. Without independent future variables, ``this time I am on a higher turn'' is only another self-narrative.

\section{Case-Level Pressure Test: How to Show That Two Ledgers Really Add Resolution}

If separating return-completion from continuation merely makes the story sound more refined, the empirical bridge has not been completed. At minimum, the two coordinates should be supported by different information sources.

One executable design would train participants to recognize one finite episode structure while arranging that the same completion reliably leads to different futures in different contexts. On test trials, first ask
\begin{quote}
``Is this the same kind of pattern occurring again?''
\end{quote}
Then, independently, ask participants to predict the next event, choose an action, report available options, or perform a later-memory-retrieval task. This obtains the return judgment first and a $C^+$ proxy second, preventing the same sentence from defining both variables.

Stronger support would require all of the following:
\begin{enumerate}
\item after semantic similarity is controlled, two episodes still receive the same stable completion/again judgment;
\item independent next-event predictions or action sets nevertheless show repeatable context-sensitive differences;
\item the continuation differences are not replaceable by one scalar ``stronger/weaker affect'' variable;
\item adding continuation coding predicts later behavior better than a model that records only whether ``it happened again.''
\end{enumerate}

The exit conditions are equally clear. If participants cease to call two episodes the same pattern whenever they predict different futures, return-completion and continuation may not be separable in that task. If $C^+$ differences depend entirely on post-hoc researcher stories and fail across trials, coders, or held-out episodes; or if a simpler learning/context model explains the phenomenon completely, the residue bridge must contract.

\section{Alternative Models and Translation Remainder: Different Afters Belong First to the Problems of Psychology Itself}

Episode B may have a different continuation for many established psychological reasons: learning, context effects, relational security, memory updating, reappraisal, a changed task goal, fatigue differences, or simply more information provided by the collaborator. The book has no reason to bypass those variables and attribute the difference directly to ``residue.''

AF-by-discrete contributes only one structural question: even after two episodes have been merged by the same completion code, which differences should not be quotiented out because they change the future?

The translation remainder is therefore large. $C^+$ does not explain why affect occurs, how learning happens, how body or neural mechanisms implement anything, or which future is healthier, more ethical, or more worth choosing. It only prevents one ``again'' from finalizing those questions in advance.

\section{Case-Level Conclusion: Return Is a Loop, Not a Seal on Fate}

Chapter 34 showed that a true local projection does not thereby own the whole. We can now state the temporal version: a real return does not thereby own the future.

A person can therefore say without contradiction:
\begin{quote}
``Yes, the old pattern happened again. It really did trigger me. But this time it did not automatically write the path that followed.''
\end{quote}
This neither denies the past nor makes change depend on a new identity requirement of ``never being triggered again.'' It simply separates return-completion from continuation.

The chapter retains three claims:
\[
\begin{gathered}
\text{the same ``again'' can lead to different afterwards;}\\
\text{return says that a finite path forms a loop;}\\
\text{it does not write the global structure or the future in advance.}
\end{gathered}
\]
That is what the residue interface is most worth keeping at the case level: not adding a ``deeper remnant'' to experience, but continuing to protect differences that genuinely alter the future after return-completion has already been established.

\chapter{Case-Level Counterexamples: When Should the Present Structural Analysis Exit AF-by-Discrete?}
\label{chap:case-counterexamples-en}

\section{First Specify What Exits and What Does Not}

The Translation Protocol requires every cross-domain model to disclose its boundaries, counterexamples, and evidence level \citep{translation2026}. This matters especially here because the book uses several layers of claims with different strengths. They cannot be bundled into one all-or-nothing package.

Layering is not a back door for saving the theory. It keeps counterexamples positioned correctly. Suppose a case initially assumes one organizing position and later data stably require at least three. The correct conclusion is ``the single-observation-position assumption failed,'' not ``Third-Order Consciousness failed.'' Likewise, if frozen analysis rules require three irreducible candidate singular-like loci that can be reidentified across material, the first object to exit is the single-singular-boundary prototype. A finite multi-singular-boundary candidate has not yet been excluded, still less AF-by-discrete itself. Conversely, ``three loci are enough in the present finite sample'' is not positive evidence for global finiteness. Counterexamples can themselves overreach when they defeat one local assumption and are rewritten as final judgments on the entire theory.

Layering also cannot become a survival tactic. If every failed hypothesis is quietly downgraded to a weaker statement while the abstract continues to advertise the original strong claim, the ``exit'' is nominal only. Part V needs a traceable retreat: retreat exactly as far as the data force. If even the weakest finite singular-interface/persistent-locus correspondence fails, the cross-domain model must actually be left behind.

\begingroup\small
\begin{longtable}{>{\raggedright\arraybackslash}p{0.40\textwidth}>{\raggedright\arraybackslash}p{0.50\textwidth}}
\toprule
If we observe & What should exit first\\
\midrule
The same material requires several irreducible observation positions & The claim of one globally unique root/observation position exits; AF-by-discrete does not thereby exit.\\
Several irreducible, reidentifiable candidate singular-like loci & The single-singular-boundary prototype exits; test a finite multi-singular-boundary candidate before declaring the model failed.\\
As resolution becomes finer, the singular-like connector/interface burden per predeclared comparable unit keeps growing without uniform control & The singular-sector per-unit uniform-control proxy exits first. Because regular boundary occurrences are not measured, this is not a direct empirical refutation of mathematical bounded type.\\
After coding rules are fixed, candidate singular-interface loci proliferate without bound as material grows and cannot be repeatedly reidentified with one finite set & The singular-boundary-germ translation used in this book comes under direct falsification pressure. This is not an empirical refutation of the abstract mathematical AF-by-discrete definition.\\
Root, return, residue, and related distinctions add no independent behavioral or predictive differences & Any Level-III/IV empirical or mechanism upgrade exits; at most the Level-II structural analogy remains.\\
\bottomrule
\end{longtable}
\endgroup

The central discipline is that a counterexample must strike the layer it actually contradicts. Failure of a single-singular-boundary prototype cannot be written as ``AF-by-discrete as a whole is wrong.'' Failure of a singular-sector per-unit uniform-control proxy cannot be written as empirical disproof of mathematical bounded type. In the other direction, once a stronger claim fails, one cannot quietly downgrade to a weaker layer and continue to say that the original proposition was supported.

\section{Three Things That Look Like Failures but Are Not}

Before listing genuine counterexamples, remove three common misreadings. Otherwise the legitimate openness of the model will be mistaken for instability.

\subsection{The Return of Discomfort Is Not the Failure of Third Order}

Chapters 33--35 repeatedly emphasized that the immediate success criterion for third-order observation is neither ``do I feel comfortable now?'' nor ``will the old feeling never return again?'' An old pattern can return; anxiety, shame, anger, or the feeling of waiting can return. What matters is whether they still automatically recruit later interpretation, checking, and action.
\[
\boxed{\text{recurrence of discomfort}\not\Rightarrow\text{failure of third-order organization}.}
\]
If the researcher preloads ``symptoms must immediately decrease'' as the only success criterion, the study first tests a corrective agenda, not whether an observation position has lost final authority.

\subsection{Several Observation Positions Do Not Mean the Structure Has Failed}

The projection-locality principle of Chapter 34 says that a local position can be real without thereby owning the whole \citep{projection2026}. Finding that one person uses several organizing positions across tasks, relationships, self-evaluation, and normative judgment does not refute the model. What it refutes is the narrower assumption that one observation position can speak for the entire system.

Moreover,
\[
\boxed{\text{many projections / roots}\not\Rightarrow\text{many boundaries}.}
\]
Root asks from where organization occurs; boundary asks where existing local organization becomes insufficient to determine continuation. The ledgers remain separate.

\subsection{Several Candidate Singular-Like Loci Are Not a Direct Counterexample to AF-by-Discrete}

This point must remain explicit. Part II already supplied a mathematical prototype with multiple boundary points and bridge types: a return grammar may be organized jointly by the passage order across several boundary sectors. The current psychological operationalization covers only the singular subset. If empirical material stably forces two, three, or more irreducible, reidentifiable singular-like loci, the first hypothesis to exit is therefore the single-singular-boundary hypothesis, not AF-by-discrete itself.

The next question is whether those interfaces are repeated visits to one finite persistent locus set or whether new distinct loci continue to be forced out as material grows. Grouping many loci into finitely many types/families does not answer that question.

\section{Where the Present Cross-Domain Model Is Really Hit: Singular-Interface Finiteness Fails}

Now we reach the central pressure test. In the mathematical native domain, relative to a fixed finite generating set and tile presentation, the infinite-tile picture of AF-by-discrete has only finitely many persistent boundary points/non-AF generator interfaces. ``Finite'' does not mean visually sparse and is not a sparsity intuition the researcher may redefine at will. It depends on a fixed generating language.

A finite psychological experiment cannot prove an infinite-tile statement. Part V can instead build a corresponding falsification pressure test. First freeze a finite event code, the question $Q$, episode segmentation rules, and a transition vocabulary. Then let additional comparable material enter.

Suppose the expanding material shows all of the following:
\begin{enumerate}
\item every new episode requires inventing a new ``special interface'';
\item old interfaces cannot be reidentified in new material;
\item candidate boundaries can be located only after the outcome is known;
\item repeated return, continuation failure, and recruitment of additional context do not cluster near any finite set of positions;
\item every attempted coarse-graining of interfaces immediately loses a crucial continuation difference and there is no stable finite alternative.
\end{enumerate}
Then ``a small number of persistent interfaces organizes a large background'' is no longer a good description of the task. One cannot preserve the model by saying ``the real boundary is hidden deeper'' unless an independent localization rule exists.

\begin{principle}[Exit Condition for the Finite Singular-Interface Submodel]
Suppose material is added in a predeclared order and a fixed cross-batch rule reidentifies candidate singular-boundary-like interface loci. Let $B_m$ be the number of distinct persistent loci that must still be retained after the first $m$ batches. The empirical trend required by the singular-germ correspondence proposed in this book is not merely ``there are finitely many types.'' The loci themselves should repeatedly point back to one finite set. If $B_m$ keeps acquiring new, nonmergeable, future-relevant persistent loci as material grows, and repeated return/continuation evidence gives no sign that $B_m$ is stabilizing, the present singular-boundary-germ translation should exit. This does not amount to a direct empirical refutation of source-side global AF-by-discrete boundary support because regular boundary points are not included in the present psychological operationalization.
\end{principle}

This is the reverse pressure test of the principle that boundaries must be forced out by the system. A good boundary hypothesis becomes more concentrated as the system repeatedly runs. A bad one requires the researcher to draw more fissures in order to save the theory.

\section{Bounded Type Can Fail while AF-by-Discrete Remains Alive}

This distinction is easy to lose and must be stated separately. The mathematical source side now has two different finite-level strengths on the ledger. General groupoid bounded type controls the generator-boundary burden of \emph{each individual finite AF fibre / tile} uniformly across levels. The tile-inflation form used in the author's work additionally controls the number of Bratteli vertices and edges per level, so that an aggregate boundary burden across tile types is uniformly controlled as a consequence.

The empirical side cannot literally instantiate either statement. It therefore uses only a deliberately narrower pressure test. Let $B_m$ count distinct persistent singular-like loci reidentifiable across accumulated material. To construct a bounded-type-inspired proxy, predeclare a family of episode resolutions from coarse to fine and, at every resolution $r$, predeclare a family of comparable episode/block units. For each such unit, count the singular-boundary-like / connector occurrences retained by the coding, and let
\[
N_r
=
\max_{\text{predeclared comparable units at resolution }r}
\{\text{retained singular-like interface occurrences in that unit}\}.
\]
Then ask whether $N_r$ remains uniformly controlled as $r$ becomes finer.

This is deliberately a \emph{per-unit} maximum rather than the total number of occurrences in the whole dataset. A dataset-total count can grow merely because more trials, blocks, people, or contexts are included. The per-unit maximum is the closer structural analogue of the general per-tile bounded-type direction.

Even this analogue is weaker than mathematical bounded type. $N_r$ records only singular-like empirical occurrences, whereas the mathematical finite-tile boundary includes regular as well as singular boundary data. Moreover, the present empirical program does not operationalize analogues of uniformly bounded $|V_n|$ or $|E_n|$. It therefore does not claim to test the full tile-inflation bounded-type form as a whole.

If $N_r$ keeps growing with resolution with no sign of uniform control, the first conclusion is
\[
\boxed{\text{the singular-sector per-unit uniform-control proxy exits}.}
\]
This neither empirically refutes mathematical groupoid bounded type nor automatically eliminates \AFBD{}.

\begin{plainbox}[How a cognitive-science reader should understand the distinction]
At present the empirical side asks two narrower questions. As material accumulates, do candidate singular-like loci keep returning to finitely many reidentifiable positions? As observational resolution becomes finer, is the singular-like interface burden \emph{within a comparable unit} uniformly controlled? Failure of the first damages the finite singular-germ correspondence. Failure of the second damages the per-unit uniform-control proxy. A future study that wishes to mimic the full tile-inflation bounded-type form would have to define, independently and in advance, what counts as the analogue of a finite tile type and then test the complexity of that type inventory separately. The current book does not pretend that this additional operationalization already exists.
\end{plainbox}

\section{Counterexample: Episode Segmentation Has No Cross-Scale Compatibility}

Even with a controlled interface count, the model may fail because episode geometry itself is inappropriate. Suppose segmentations of the same experience at different scales are unrelated: a coarse episode cannot be traced as a composition of finer episodes; the same concrete occurrence has no identifiable embedding after refinement; every change of resolution generates an independent data structure. Then the traverse-decomposition and inflation analogies lose their basis.

Psychological event segmentation can of course be subjective. Event Segmentation Theory and related work show that people organize ongoing activity into discrete events and that event boundaries relate to prediction, attention, and memory updating \citep{zacks2007,kurby2008}. Those literatures do not guarantee the Bratteli-like nested coherence required by this book.

If data persistently show no repeatable occurrence correspondence across scales, the corresponding filtration/persistence module should exit. It cannot be saved by declaring that ``the real nesting exists at a deeper level.''

\section{Counterexample: A Putatively Genuine Return Must Be Revoked by Later Refinement}

The formal theory in Part IV has a hard property: a finite closed witness already present does not disappear under supported inflation merely because ambient context grows. Real psychological judgments can of course be corrected by new facts.

Suppose empirical work shows that a judgment we wanted to call a ``genuine cognitive return'' holds only at a coarse scale and must be withdrawn once finer facts are added: the two episodes did not in fact instantiate the same return-completion. This cannot simply be described as ``the return was refined.'' Two possibilities must be separated:
\begin{enumerate}
\item the original record was only a subjective return hypothesis; new evidence shows it never reached the strength required for formal genuine-return correspondence;
\item the target domain genuinely requires a non-monotone theory in which return-completion itself can be revoked by later evidence.
\end{enumerate}
If the second pattern recurs, it directly strikes the persistence mapping. The mathematical theorem remains true. The cross-domain translation has failed. Mathematics cannot require psychological facts to persist merely because mathematical witnesses do.

\section{Counterexample: Return Judgment Is Driven Mainly by Labels Rather than Finite Organization}

Suppose participants reliably report ``the same thing again'' whenever the same word is present---``rejection,'' ``failure,'' ``anger,'' ``control''---while finite event order, entrance and exit, current question $Q$, organizing position, and later continuation hardly affect the judgment. Then the closed-traverse model adds no explanatory value to return recognition.

The relevant objects may instead be semantic categorization, schema activation, or narrative identity. Those are important without being renamed structural return. A particularly strong counterexample would be that changing the label sharply reduces return judgments while large changes in event structure have little effect as long as the label remains the same. In that case, a semantic baseline should be preferred.

\section{Counterexample: Third Order Leaves No Difference beyond Ordinary Metacognition}

``The observer also enters observation'' does not acquire the status of an independent psychological construct through philosophical language alone.

If future measurement finds that every observable difference attributed to Third Order is completely absorbed by existing measures of metacognition, confidence monitoring, error monitoring, decentering, and related constructs, and that additionally recording how the current observation position is generated from the past, current question, and evaluative rules adds no prediction, behavior, or relational difference, this book has no reason to insist on a stronger empirical third-order construct. Research on metacognition and decentering supplies mature comparison baselines \citep{fleming2012,fleming2012a,bernstein2015}.

This exit condition is especially important because the novelty of ``Third Order'' can easily grant it unnecessary explanatory privilege. If established constructs already do the same work, the correct result is not a renaming but an acknowledgment that the translation added no resolution.

\section{Counterexample: The Phenomenon Appears Only after Participants Learn the Book}

A more subtle risk is that the theory may manufacture the object it claims to discover. Participants may originally never describe a conflict as a ``return'' and may not stably distinguish several organizing positions. After learning the vocabulary of root, boundary, return, and residue, they begin reporting experience in those terms.

That result does not by itself show the model is wrong; conceptual training can improve discrimination. But it cannot directly establish that the same structures existed before training as the same psychological mechanisms.

Research designs must therefore compare:
\begin{itemize}
\item whether participants unfamiliar with the theory show corresponding behavioral and temporal structure;
\item whether neutral language yields the same segmentation, return judgments, and continuation predictions;
\item whether theory training increases sensitivity to pre-existing differences or merely increases use of the vocabulary of the book.
\end{itemize}
If all ``evidence'' appears only after participants have learned the theory and no corresponding difference exists in independent tasks, we have shown at most that a new self-description vocabulary can be learned. That is not evidence for a pre-existing mechanism.

\section{Counterexample: Alternative Resolution Can Be Invented Only after the Fact}

The pNH benchmark in Chapter 21 imposes a strict condition on ``looking from another angle'': an alternative resolution must have a system-internal source and should ideally be declared before the current result is known.

If every counterexample is followed by an ad hoc change in the completion criterion so that any outcome can be made to look completed ``from another angle,'' or every failed root prediction is answered with ``the real observation position was actually somewhere else,'' the third-order observation site has become unfalsifiable narrative freedom.

If the admissible question/frame family cannot be declared in advance, or if arbitrary post-hoc reframing is predictively indistinguishable from the proposed structural resolution shift, the module must exit.

The projection discipline of Chapter 34 applies as well. ``One projection is not the whole'' cannot mean ``there is always another projection that can explain whatever happened.'' Removing local monopoly does not remove explanatory responsibility.

\section{The Strongest Case-Level Counterexample: No Independent Observable Consequence, or a Simpler Model Already Explains Everything}

Every elegant structural distinction must eventually pass this test. At the case and experiment level, if distinguishing root, return, residue, multi-boundary passage, or telescoping changes no repeatable behavioral prediction, memory organization, attention allocation, next action, cross-scale segmentation, or counterfactual judgment, those distinctions remain explanatory language at most.

They must also compete with simple baselines. If semantic similarity, event frequency, general arousal, working-memory load, schema activation, or ordinary metacognitive confidence explains the same variance and AF-by-discrete-inspired coding adds no stable increment, a Level-III empirical upgrade does not hold and a mechanism claim is still less warranted.

The discipline is simple:
\[
\boxed{\text{more structural complexity}\not\Rightarrow\text{more explanation}.}
\]

\section{Chapter Conclusion: A Real Theory Must Let Reality Dismantle It Layer by Layer}

The purpose of this chapter is not to supply a polite checklist titled ``how to refute us.'' It establishes actual deletion rules.

When one observation position cannot represent the whole, delete its global privilege. When one candidate singular boundary is insufficient, allow the data to test a finite multi-singular-boundary picture. When singular-like occurrence burden lacks uniform control, exit the corresponding bounded-type-inspired proxy. When fixed rules cannot force out even finite persistent singular-interface organization, exit the present singular-germ translation. When structural distinctions add no independent empirical value, stop upgrading Level-II correspondence into psychological mechanism. The complete source-side conditions AF-by-discrete and bounded type remain on their own mathematical ledgers.

Layered exit prevents two opposite errors: a model that collapses globally at the first sign of complex reality, and a model that refuses ever to fail by continually adding boundaries, observers, and hidden levels.

\begin{stopbox}[Global exit condition: when should AF-by-discrete really be left behind?]
Under observation, coding, and comparison rules fixed in advance, if the main organizational complexity of the target domain cannot stably be represented as
\[
\boxed{\begin{gathered}
\text{traceable local background}\\[-1mm]
+\ \text{finite persistent continuation interfaces}
\end{gathered}},
\]
and if this decomposition adds no repeatable resolution across scales, return persistence, independent observables, and baseline-model comparisons, AF-by-discrete should no longer serve as the primary structural model for that problem.

The correct response is not to keep looking for a ``deeper boundary.'' It is to change models.
\end{stopbox}

One apparently negative but important result remains. A failed correspondence does not invalidate the psychological facts. The experience is still there. What changes is only our admission that, this time, the singular-germ translation extracted from AF-by-discrete no longer has the right to keep speaking for it.

\chapter{Empirical Bridge I: Turning Root, Episode, and Return into Measurable Questions}
\label{chap:empirical-bridge-I-en}

\par\medskip
\noindent\textbf{Why this chapter matters}\par
The preceding cases have shown that a mathematical distinction has not really returned to psychology if it merely gives us another way to speak. A genuine empirical bridge must answer more ordinary questions: what exactly do we record? How do we know that two analysts are talking about the same thing? What result would make the correspondence fail?

We therefore do not search for a magical scale that directly ``measures'' root, boundary, or return. Instead, we decompose those structural distinctions into ordinary behavioral, temporal, memory, judgment, and prediction variables, and then ask whether they form stable relations. Mathematical objects remain in the mathematical native domain. Psychological experiments test only whether the structural resolution they motivate has empirical content.
\par\medskip

\section{First Discipline: Do Not ``Measure a Mathematical Object''; Measure the Psychological Differences It Requires Us to Separate}

Moving from Level-II structural compatibility toward empirical research usually begins with operationalization. Here one shortcut is especially dangerous: writing mathematical nouns directly into a questionnaire and treating participants' answers as if the mathematical object had been measured.

Asking ``did you feel a boundary appear?'' does not measure boundary. Asking ``did you return to the same pattern?'' does not by itself prove return. It simply asks participants to repeat language already accepted by the researcher.

The required order is slower:
\begin{center}
\resizebox{0.96\linewidth}{!}{$\displaystyle
\text{ordinary psychological/behavioral data}
\longrightarrow
\text{case-level coding}
\longrightarrow
\text{structural correspondence hypothesis},
$}
\end{center}
not
\[
\text{mathematical term}=\text{questionnaire item}=\text{psychological entity}.
\]

The ordinary data can be mundane: where someone presses an event-segmentation key, when a prediction changes, which action is chosen next, which past episode is retrieved first, whether the person says ``this is happening again,'' and whether a judgment changes after new evidence. The case layer then uses $D,\Pi,Q,\Sigma,\widehat{\mathcal R}_{\mathrm{psych}},C^+$ to organize those observations. Only at the final step does the researcher ask whether the relations are stably structurally compatible with formal distinctions such as root, return, boundary, or resolution.

\begin{plainbox}[One observable is never enough]
The empirical work envisioned by this book follows a triangulation discipline: every major structural hypothesis should, wherever possible, be supported by at least two independent indicators and must compete with a simpler psychological explanation. A sudden attention shift alone does not establish root shift; one surprise does not establish a boundary-like interface; one ``here we go again'' does not establish return. Structural correspondence begins to acquire empirical content only when several independent indicators covary under predeclared rules.
\end{plainbox}

\section{Episode: Measure ``Where the Cut Is'' before Asking ``Why the Cut Is There''}

Experience does not arrive with episode boundaries already attached. We usually live through ongoing activity first and later say ``this part,'' ``that time,'' or ``from here it became different.'' Event-segmentation research supplies mature paradigms: participants can mark when they think one event ends and another begins while watching ongoing activity, and these segmentations relate to prediction updating, attention, and memory \citep{zacks2007,kurby2008}.

This is an important empirical starting point, but it does not complete the translation. Keep
\[
\boxed{\text{EST event boundary}\neq\text{the boundary-like interface in this book}.}
\]
The former first asks where participants segment ongoing activity. The latter additionally requires that some positions play a special role in continuation, context recruitment, or organizing-position change.

At minimum, episode measurement should therefore contain three layers:
\begin{enumerate}
\item \textbf{Explicit segmentation:} where participants or independent coders place episode cuts;
\item \textbf{Implicit segmentation:} whether those positions coincide with independent signs such as memory chunking, prediction updating, or reaction-time change;
\item \textbf{Question dependence:} whether $\Pi$ changes reproducibly when $Q$ changes, rather than moving arbitrarily with a retrospective researcher story.
\end{enumerate}
The third matters especially. Part V has repeatedly emphasized that segmentation is not a uniquely true answer read from a God's-eye position. The same material can legitimately preserve different distinctions when the question is ``is the other person's relational intention known?'' versus ``can my next action stop now?'' What should be measured is not a unique correct cut, but whether changes in segmentation are stable and reproducible after the question is declared and whether they affect independent memory, prediction, or action.

\section{Root: Do Not Measure ``Where Am I?''; Measure Which Position Organizes Comparison and Continuation}

Root is among the easiest mathematical objects to misread psychologically. It is not a brain region, is not attention itself, and is not ``the true self.'' The weaker case-level question is: from which distinguished reference position is current experience being organized?

Candidate root/organizing positions can be recorded by $\Sigma$. A position deserves to be retained only if it changes some organization relation, for example:
\begin{itemize}
\item what the same stimulus is compared against;
\item which evaluative rule becomes primary;
\item which past material receives priority in retrieval;
\item which actions are experienced as currently available;
\item how the next event is predicted.
\end{itemize}
A gaze shift, attention shift, or increase in affect is therefore insufficient by itself. Stronger evidence would be that, while the stimulus is held as constant as possible, a change of reference frame simultaneously reorganizes at least two dimensions among comparison, retrieval, and continuation.

Research on sense of ownership and sense of agency distinguishes experiences such as ``this body/experience is mine'' from ``I produced this action'' \citep{braun2018}. These variables can be measured alongside candidate root shifts because both concern self-reference, but they are not identical. If ownership/agency fully explains the observed changes, the root correspondence should contract rather than duplicate an established construct.

The projection discipline of Chapter 34 remains:
\begin{center}
\resizebox{0.96\linewidth}{!}{$\displaystyle
\boxed{\text{a measurable root-like position}\neq\text{the person's global psychological structure}.}
$}
\end{center}
One participant may show several stable organizing positions across tasks, relationships, and times. Discovering multiple positions is not measurement failure; it is a reason not to let the most salient observation position take possession of the whole.

\section{Return: Measure Content Similarity Separately from Finite Return-Completion}

The experience of ``again'' is easily absorbed by semantic similarity. If two episodes both contain ``no reply,'' ``argument,'' or ``checking the phone,'' participants may say ``here it is again.'' The return translation asks a narrower question: did the participant merely recognize similar content, or a finite episode with an entrance, internal passage, and return-completion relation?

Empirical measurement therefore needs at least three information sources:
\begin{enumerate}
\item \textbf{Structural coding:} independent coders use predeclared endpoint, order, and episode rules to decide whether two segments satisfy the same finite completion pattern;
\item \textbf{Subjective recognition:} participants independently judge whether ``the same kind of thing happened again'';
\item \textbf{Later consequence:} the judgment changes memory chunking, next-event prediction, action choice, or information search.
\end{enumerate}
Semantic similarity, lexical labels, frequency, and temporal proximity must be controlled. Otherwise the study can show only that similar things are more likely to be judged similar.

\begin{plainbox}[Minimum return evidence]
The persuasive result is not perfect agreement among all three indicators. It is that, after content similarity is controlled, predeclared finite sequence/endpoint structure still adds prediction for ``again'' judgments or later behavior. Conversely, if return judgment is almost fully explained by semantic similarity, affect intensity, or frequency, the psychological return correspondence should weaken. The mathematical definition of return in Part IV does not fail thereby.
\end{plainbox}

\section{Boundary: Do Not Ask Participants to Point out the Fissure for the Theory; Test Whether the Fissure Reappears by Itself}

Boundary requires the greatest care in Part V. Chapter 30 already fixed the order: do not first ask ``is this a boundary?'' Record first where local organization repeatedly becomes insufficient for continuation.

Stronger empirical signs of a candidate boundary-like interface may include:
\begin{itemize}
\item a current local policy or prediction repeatedly failing at comparable positions;
\item repeated recruitment of more distant context at those positions;
\item clustering of return endpoints or passages between organizing positions near those transitions;
\item stable changes in predictions, action sets, or information needs in $C^+$ after crossing those positions.
\end{itemize}
The crucial words are ``repeatedly'' and ``comparable positions.'' One surprise, affect peak, or attentional interruption is not enough. A stricter design should discover transition clusters in one subset of episodes and then ask in held-out episodes whether the same rule locates them without modification and whether they predict continuation change.

This also turns the distinction between single- and multi-singular-boundary candidates into an empirical question. If several irreducible, reidentifiable transition clusters stably predict different passage/continuation patterns, they should not be merged into one ``core fissure'' merely for theoretical neatness. But
\[
\boxed{\text{many organizing positions}\not\Rightarrow\text{many boundaries}.}
\]
Several root-like positions may repeatedly transform through the same interface family. Only irreducible transition families forced out repeatedly by the system provide empirical reason for a multi-boundary prototype.

\section{Finite Singular-Interface Proxy and Bounded Type: What Empirical Research Can Pressure-Test and What It Cannot Prove}

No finite psychological experiment can prove an infinite mathematical bounded-type statement. The source audit in Part II gives an even more specific reason for caution: ``bounded type'' has a general groupoid formulation and a more explicit finite-rank tile-inflation form with additional presentation-level controls. Part V therefore tests only explicitly declared fragments of those structures.

\begin{plainbox}[Separate locus, occurrence, family, per-unit burden, and type inventory]
From Part V onward five bookkeeping objects must remain distinct.
\begin{itemize}
\item \textbf{Interface occurrence:} a candidate transition occurrence detected in one concrete episode/trial. One persistent locus may be traversed many times.
\item \textbf{Persistent interface locus:} a candidate fissure judged, under fixed coding and reidentification rules, to be ``the same structural position'' across material. This bears the finiteness pressure of the finite singular-interface correspondence, not the complete \AFBD{} boundary count.
\item \textbf{Interface family/type:} a coarser descriptive coding that groups loci by behavioral profile, return grammar, or continuation profile. A finite family count cannot substitute for a finite locus count.
\item \textbf{Per-unit interface burden:} at a predeclared resolution $r$, count how many singular-like interface/connector occurrences must be retained inside each predeclared comparable episode/block unit, and let $N_r$ be the maximum of those counts.
\item \textbf{Structural unit-type inventory:} the number and identity of the comparable unit types used at resolution $r$. The present book does not claim that this quantity is the psychological counterpart of $|V_n|$ or $|E_n|$; any future attempt to model the full tile-inflation bounded-type form would have to justify and test that correspondence independently.
\end{itemize}
\end{plainbox}

For the singular-boundary-germ translation, ask whether the cumulative number $B_m$ of distinct persistent singular-like loci stabilizes after observation, coding, and cross-material reidentification rules are frozen.

For the bounded-type-inspired per-unit proxy, ask whether
\[
N_r
=
\max_{\text{predeclared comparable units at resolution }r}
\{\text{retained singular-like interface occurrences in that unit}\}
\]
remains uniformly controlled as resolution is refined. The total number of occurrences in the entire dataset is not used for this purpose.

The empirical ledger is therefore
\begin{center}
\resizebox{0.96\linewidth}{!}{$\displaystyle
\begin{array}{rcl}
B_m\text{ stabilizes}&\Rightarrow&\text{support for the finite singular-germ correspondence},\\
N_r\text{ uniformly controlled}&\Rightarrow&\text{support for a singular-sector per-unit proxy},\\
B_m\text{ keeps adding new loci}&\Rightarrow&\text{weakens the singular-germ correspondence},\\
N_r\uparrow,\ B_m\text{ stable}&\Rightarrow&\text{weakens only the per-unit proxy first}.
\end{array}
$}
\end{center}

No row is an empirical proof or refutation of the complete mathematical \AFBD{}, the general groupoid bounded-type condition, or the full tile-inflation bounded-type form. Their purpose is narrower: to make the cross-domain correspondence genuinely capable of failure without pretending that psychological samples are Bratteli levels.

\section{How Third-Order Observation Enters Measurement: Not by Adding ``Did You Become Aware?''}

The empirical difficulty of Third Order is that if a researcher asks ``did you realize that you were observing?'' participants can easily give the theoretically correct answer without any organizational change. Third Order acquires empirical content only if the generation conditions of the observation position itself enter the record.

More informative observables include:
\begin{itemize}
\item whether participants can identify which evaluative rule, past reference, and comparison target they just used;
\item whether changing those conditions changes segmentation, return judgment, or self-evaluation accordingly;
\item whether the corrective agenda ``I must make the discomfort go away now'' enters the record as part of the current process rather than remaining an invisible criterion of success;
\item whether theory invocation, effect checking, and re-invocation form a repeatable meta-monitoring loop.
\end{itemize}

This is why Chapter 33 distinguished the directional wish for relief from the corrective agenda of eliminating the state by a method. The empirical comparison is whether observation placed under an explicit symptom-reduction goal produces different self-monitoring, evidence updating, and action choices than observation in which the urge to solve is itself recorded. Those differences must compete with mature constructs in metacognition, confidence monitoring, and decentering \citep{fleming2012,fleming2012a,bernstein2015}.

If those established variables already explain the changes fully, Third Order should not acquire an additional psychological mechanism identity.

\section{An Executable Measurement Ledger}

\begingroup\scriptsize
\begin{longtable}{>{\raggedright\arraybackslash\bfseries}p{0.16\textwidth}>{\raggedright\arraybackslash}p{0.25\textwidth}>{\raggedright\arraybackslash}p{0.19\textwidth}>{\raggedright\arraybackslash}p{0.20\textwidth}}
\toprule
Structural question & Minimum candidate evidence & Must control & Result that weakens correspondence\\
\midrule
episode & explicit segmentation + independent segmentation signs in memory/prediction & $Q$, time scale, task demand & cuts move arbitrarily with retrospective instruction\\
root-like position & a reference-frame change jointly reorganizes at least two dimensions among comparison, retrieval, and continuation & attention, affect intensity, ownership/agency & alleged root shift is only general attention shift\\
return & finite sequence/endpoint coding + independent ``again'' judgment + at least one later behavioral measure & semantic similarity, frequency, lexical label & ``again'' is fully explained by content similarity\\
singular-boundary-like interface & local-sufficiency failure/context recruitment clusters stably in held-out episodes and predicts $C^+$ change & surprise, arousal, researcher-preselected positions & every batch requires newly invented, nonreidentifiable singular-like loci\\
finite multi-singular-boundary & several irreducible, reidentifiable transition loci each have stable passage/continuation consequences & many roots do not imply many boundaries; finite-sample sufficiency does not prove global finiteness & merging several loci loses no information\\
third-order audit & evaluative rule, past reference, and corrective agenda entering the record produce independent updating differences & metacognition, decentering, reappraisal & mature constructs already explain the entire effect\\
\bottomrule
\end{longtable}
\endgroup

\section{When Is It Legitimate to Say That We Have Taken a Step toward Level III?}

A structural correspondence only begins to qualify for movement from Level II toward Level III if all of the following hold:
\begin{enumerate}
\item \textbf{Rules are declared in advance.} $Q$, $\Pi$, candidate interface detection rules, and return coding cannot be changed after the main result is known.
\item \textbf{At least one independent indicator.} Material used to define the structure cannot simultaneously serve as the sole evidence of success.
\item \textbf{Held-out generalization.} The same rule locates the relevant structure in new episodes, new participants, or new contexts.
\item \textbf{Competition with simpler models.} Semantic similarity, ordinary event segmentation, attention, arousal, metacognition, and related baselines must be compared explicitly.
\item \textbf{Negative results have consequences.} If a structural variable adds no prediction, the correspondence must contract, be deleted, or exit rather than being saved by adding another hidden boundary.
\end{enumerate}

\chapter{Empirical Bridge II: Eight Executable Research Designs}
\label{chap:empirical-bridge-II-en}

\par\medskip
\noindent\textbf{Why this chapter matters}\par
The question is no longer whether a mathematical word resembles a psychological phenomenon. We decompose the most important structural distinctions into eight research questions that can fail independently. On a first reading, one can ignore statistical details and ask whether each design does four things: fixes questions and coding rules in advance, tests consequences with independent indicators, lets simpler psychological accounts compete, and specifies what result would genuinely weaken the correspondence.
\par\medskip

All designs in this chapter are research proposals, not completed experiments and not ready-made scales. Their only task is to rewrite Level-II structural compatibility as constrained models that could in the future enter Level III. No experiment will ``measure a mathematical root,'' ``discover a mathematical boundary point,'' or ``validate a groupoid.'' It can only test whether psychological variables exhibit the relations required by the book.

\section{A Common Execution Protocol for All Eight Studies}

To prevent eight studies from silently changing the rules in eight different ways, fix one common protocol.
\begin{enumerate}
\item \textbf{Declare $Q$ and segmentation rules first.} Research questions, episode segmentation, candidate interface-detection rules, and primary outcomes must be fixed before the main result is seen.
\item \textbf{Separate definition material from test material.} If some trials are used to learn thresholds, candidate boundary scores, or organizing-frame classifiers, the main conclusion must be tested on held-out trials, held-out episodes, or new participants.
\item \textbf{Retain at least one ordinary-psychology baseline.} Semantic similarity, frequency, ordinary prediction error, arousal, attention, event segmentation, metacognition, decentering, and related variables get the first opportunity to explain what they can.
\item \textbf{Success criteria cannot equal definitions.} Material used to declare ``a return occurs here'' cannot also be the sole evidence that return has behavioral consequences. Prediction failure used to locate an interface cannot itself prove that the interface changes future continuation.
\item \textbf{Failures are settled locally.} Failure of one design retracts the correspondence it tests. If the failure directly hits the finite singular-interface/persistent-locus hypothesis, the model cannot survive indefinitely by adding hidden boundaries.
\end{enumerate}
The eight studies fall into four groups. A--C ask how repetition becomes recognized return. D and F ask why return-completion is not all future information. E asks whether Third Order adds anything beyond ordinary metacognition and whether a corrective agenda can re-instrumentalize it. G--H ask whether multiple organizing positions and interfaces forced out by the system add prediction.

\section{Study A: Can Raw Repetition Be Separated from an ``Again'' Judgment?}

\textbf{Core distinction.} Content similarity or action repetition is not automatically psychological return. Test whether, beyond similarity, an episode's entrance, internal order, and completion relation still affect ``again.''

\textbf{Design.} Construct paired episodes using short animations, interactive scripts, or textual events. Independently manipulate in advance:
\begin{enumerate}
\item surface similarity of characters, actions, objects, and affect words;
\item a predeclared completion grammar: whether the second segment forms the same finite completion as the first under fixed endpoint and order rules.
\end{enumerate}
This yields four conditions: similar/completed, similar/not completed, dissimilar/completed, dissimilar/not completed. Participants answer only ordinary-language questions such as ``is the same kind of thing happening again?'' They are not taught return, boundary, or witness.

\textbf{Measures and baseline.} Trial-level ``again'' judgment is primary. Semantic/visual similarity ratings are the main baseline; reaction time and confidence are auxiliary. A similarity-only model predicts first, and the analysis then asks whether completion adds held-out prediction.

\textbf{Support.} After content similarity is controlled, predeclared completion structure still reliably improves prediction of ``again'' judgments and generalizes to new material.

\textbf{Negative result.} If judgments are almost entirely explained by semantic/action similarity, or the completion effect exists only in the same material used to define it, the psychological return--``again'' correspondence must contract sharply. Mathematical return is unaffected.

\section{Study B: Does Multi-Scale Episode Segmentation Show Stable Compatibility?}

\textbf{Core distinction.} Mathematical decomposition in Part IV is strict. Psychological experience need not copy strict nesting, but if coarse, medium, and fine segmentations are unrelated, the cross-scale translation loses its empirical support.

\textbf{Design.} Use natural activity videos or ongoing tasks from event-segmentation paradigms \citep{zacks2007,kurby2008}. Ask participants separately to mark the ``finest natural units,'' ``medium meaningful units,'' and ``largest complete activities.'' Randomize order to reduce anchoring from earlier segmentations. A separate participant group performs only one scale to estimate instruction-demand effects.

\textbf{Measures and baseline.} Predefine a cross-scale nesting/alignment score: whether fine-scale boundaries cluster consistently near coarse boundaries across participants and whether coarse episodes can usually be approximately covered by consecutive finer episodes. General boundary salience, action-change magnitude, and visual cuts serve as baselines.

\textbf{Support.} No exact mathematical nesting is required. What matters is repeatable hierarchical compatibility beyond low-level perceptual change, preserved in new videos.

\textbf{Negative result.} If the relation among scales is almost fully explained by visual/action salience, or boundaries at different scales interleave arbitrarily with no repeatable nesting tendency, the psychological translation of decomposition coherence should exit. Mathematical strict decomposition cannot fill gaps in empirical data.

\section{Study C: Does Recognized Return Predict Chunking and the Next Step Better than Bare Repetition?}

\textbf{Core distinction.} A pattern repeating and a participant organizing it as ``again'' may be different events. The study tests incremental prediction without assuming return recognition is a causal memory-chunking mechanism.

\textbf{Design.} Participants learn sequences containing repeated patterns matched in length and frequency. For every repetition, record separately whether the participant subjectively recognizes ``the same pattern again.'' Later, without further ``again'' prompts, test free/serial recall, chunk boundaries, next-event prediction, and reaction time. Training material and primary test material are separate.

\textbf{Measures and baseline.} First predict memory and next responses using repetition frequency, semantic similarity, surprisal, and general familiarity. Then add the trial-level return judgment and test out-of-sample increment.

\textbf{Support.} After bare repetition and familiarity are controlled, episodes recognized as returns form more stable memory chunks or constrain next-event prediction more effectively.

\textbf{Negative result.} If return judgment adds nothing to chunking or prediction, ``return promotes compression/organization'' remains a formal inspiration only. Correlation, if found, still is not causation; a causal upgrade requires independent manipulation of return recognition.

\section{Study D: After the Same Return Judgment, Can Later Continuation Remain Different?}

\textbf{Core distinction.} This is the minimal empirical test of the residue correspondence: sameness of return-completion and sameness of future must be separable.

\textbf{Design.} Train two backgrounds $A$ and $B$. Insert the same finite episode $E$ in both, so that participants stably recognize $E$ as the same pattern, but train the later relations
\[
A+E\to Y,
\qquad
B+E\to Z,
\qquad
Y\neq Z.
\]
At test, first ask whether the same kind of event happened again. Then, without revealing the correct answer, ask participants to predict the next event, choose an action, or indicate what background information should be retrieved.

\textbf{Measures and baseline.} First test whether completion judgment remains stable across $A,B$. Second test whether future prediction still depends on background. A model using only the current episode competes directly with a context-augmented model.

\textbf{Support.} Participants simultaneously maintain
\[
\boxed{\text{same return judgment} + \text{different continuation prediction}.}
\]
This supports the psychological resolution of two coordinates.

\textbf{Negative result.} If recognizing the same pattern necessarily collapses later predictions, or different continuations are fully explained by obvious surface cues, the residue-style translation adds no content. Even support measures only a $C^+$ proxy, not mathematical germ residue.

\section{Study E: Does Third-Order Structural Audit Exceed Ordinary Metacognition and Avoid Re-Instrumentalization by a Corrective Agenda?}

\textbf{Core distinction.} The observer-belongs-to-the-system idea has earlier ancestors in second-order cybernetics and is not claimed as a unique discovery of Third-Order Consciousness. The narrower increment is whether open structural audit exceeds confidence monitoring and decentering, and whether an identified observation position becomes an unaudited manager again under a corrective agenda.

\textbf{Design.} Use a low-risk self-evaluation task: after difficult judgments, participants receive ambiguous, updateable performance feedback. Randomly assign four instruction conditions, without using the phrase Third-Order Consciousness:
\begin{enumerate}
\item \textbf{Metacognitive monitoring:} report confidence and where judgments may be wrong.
\item \textbf{Decentering:} describe current thoughts as mental events without immediately taking them as facts.
\item \textbf{Open structural audit:} record the rules, past references, and comparison targets used in current self-evaluation; also record the wish to feel better and urge to solve the problem, while explicitly not treating immediate discomfort reduction as the criterion of successful observation.
\item \textbf{Corrective structural audit:} record the same rules/references but explicitly set the goal ``use this method to reduce discomfort as quickly as possible,'' with permission to check whether the method is working.
\end{enumerate}
Then add new evidence partially conflicting with the initial self-evaluation and allow participants to gather more information, update judgment, reopen the strategy instruction, or end the task.

\textbf{Primary measures and baselines.} Immediate comfort is secondary. Primary measures include whether evaluative rules can be reported and revised, whether updating after new evidence becomes more evidence-sensitive, how often strategy instructions are reopened or ``do I feel better?'' is checked, and whether next action remains locked by one self-evaluation.

\textbf{Second stage: relativizability of meta-identification.} After Stage 1, participants explicitly write the organizing rule $O$ they just identified. Keep the core event material fixed while changing predeclared question or role context, for example: ``is this feedback sufficient to judge ability?''; ``should more information be collected?''; ``does this rule still apply in a collaboration?'' Compare whether $O$ becomes a new total label across contexts or whether participants can retain its local reality while allowing it to lose current authority elsewhere. Outcomes include evidence updating, memory retrieval, action choice, meta-checking, and forced cross-context generalization of $O$.

\textbf{Support.} Open structural audit would need to add prediction beyond confidence monitoring, decentering, and ordinary context sensitivity; reduce meta-checking/strategy re-invocation relative to corrective framing; increase updateability to new evidence; and allow $O$ to remain locally real without becoming a new total center. Immediate discomfort need not be lower. Such a result would be compatible with
\[
\text{freedom from structural capture}\neq\text{immediate comfort}
\]
and with the relativizability of meta-identification. At most it would provide a psychological increment compatible with the non-Hausdorff-inspired distinction; it would not measure a literal non-Hausdorff germ.

\textbf{Negative result.} If every effect is absorbed by ordinary metacognition/decentering, or open and corrective framings do not differ stably in meta-checking, updating, or continuation, Third Order should not be written as an already independent psychological capacity \citep{fleming2012,bernstein2015}. The study also does not test a mechanism claim that ``fully staying with pain necessarily produces release.''

\section{Study F: Can the Same Material Be ``Unresolved under One Question and Completed under Another''?}

\textbf{Core distinction.} The observation-site benchmark borrowed from pNH in Chapter 21 is not a priority claim for Third-Order Consciousness. It formalizes a shared reflexivity problem: completion can depend on an explicit question $Q$, but the change must be a constrained question-relative resolution rather than arbitrary retrospective storytelling.

\textbf{Design.} Present the same interaction to all participants, for example ``send invitation--wait--no reply--stop contacting.'' Before the stimulus, fix two questions:
\begin{enumerate}
\item $Q_A$: is the other person's relational intention determined?
\item $Q_B$: under the given action rule, must I continue waiting or contacting?
\end{enumerate}
Randomize question order. Add a post-hoc-reframing control that introduces a new interpretive frame only after the event ends.

\textbf{Measures and baseline.} Record completion judgments for both questions, what information is still needed, whether participants continue waiting/contacting, and later information search. Test whether predeclared $Q$ predicts behavioral differences more stably than post-hoc story.

\textbf{Support.} The same material stably supports
\[
Q_A:\text{ unresolved},
\qquad
Q_B:\text{ completed},
\]
and these judgments lead to different information search and action choices.

\textbf{Negative result.} If predeclared questions differ no more from arbitrary retrospective reappraisal, or the judgments have no independent behavioral consequences, the pNH-style observation-site correspondence remains only a methodological benchmark rather than an empirical mechanism.

\section{Study G: Does the Order of Passage through Several Organizing Positions Predict Better than ``One Fundamental Loop''?}

\textbf{Core distinction.} Multiple positions and multiple boundaries are not interchangeable. This study asks a weaker measurable question: once checking, self-evaluation, rule monitoring, and related frames are independently identifiable, does their passage order predict ``again'' and the next step better than a dominant-frame model?

\textbf{Design.} Use calibration trials independent of the main task to establish distinguishable organizing frames $\sigma_A,\sigma_B,\sigma_C$, such as task checking, self-evaluation, and rule monitoring; freeze the classifier/coding rule. In the main task, match surface content while independently manipulating:
\begin{enumerate}
\item whether action/semantic content is similar;
\item whether the unordered set of frames is the same;
\item whether passage order, such as $A\to B\to C\to A$, repeats.
\end{enumerate}
Collect same-pattern/again judgments, next-event predictions, memory chunking, and action choices.

\textbf{Model competition.} Compare at least three held-out models: one dominant frame, unordered set of frames, and ordered passage. Only if ordered passage adds prediction on new trials should a multi-sector grammar be retained.

\textbf{Support.} After semantic similarity and single-frame cues are controlled, passage order still predicts return judgment or continuation. This provides a weaker psychological candidate corresponding to the multi-boundary type-word structure of Kuang v1, Subsection 5.2 \citep{kuang2024v1}.

\textbf{Negative result.} If ordered passage does not outperform single-root/unordered baselines, or frames can be named only after outcomes are known, the multi-sector passage translation should exit. Even if ordered passage succeeds, it establishes only that organizing-position sequence carries information; many organizing positions still do not imply many mathematical boundary points.

\section{Study H: Two-Level Pressure Tests for the Finite Singular-Interface Hypothesis}

\textbf{Core distinction.} The study tests only a psychological finite singular-interface hypothesis, call it $H^{\mathrm{psych}}_{\mathrm{fin}}$: candidate persistent singular-like loci stabilize in a finite-set-like way under declared tasks, scales, and identity rules. It does not directly measure the source-side AF-by-discrete finiteness statement because the regular boundary part lacks a psychological operationalization.

\subsection*{H1: With Known Finite Interfaces, Can the Method Recover Them?}

H1 is a measurement-recovery benchmark, not evidence for natural psychological finiteness. Construct a long sequence-learning or interactive task whose generative mechanism is known to contain only a few hidden transitions requiring remote background. Participants do not know their locations. Continuously record prediction error, context recruitment, episode segmentation, strategy shifts, return-endpoint judgments, and next choices.

Use a discovery set to construct a detector
\begin{center}
\resizebox{0.96\linewidth}{!}{$\displaystyle
I_{\mathrm{detect}}(t)
=
w_1\,\text{local-policy failure}
+w_2\,\text{context recruitment}
+w_3\,\text{return/passage concentration}.
$}
\end{center}
Weights are preassigned or learned once on the discovery set and then frozen. Continuation shift does not enter the detector. Held-out episodes independently test segmentation, information search, and
\[
\Delta C^+(t)=\text{continuation-profile change across the candidate position}.
\]
H1 success shows only that the detector and locus-identity rule can recover finite-interface ground truth. It does not show that open-world psychology naturally has finitely many loci.

\subsection*{H2: Without Finite Ground Truth, Do Loci Continue to Proliferate?}

Use longitudinal experience sampling, long-running task-switch records, repeated interactions, or cross-context behavioral sequences. Before data collection, freeze candidate features, episode rules, cross-material locus-identity rules, thresholds for a new locus, complexity penalties, and held-out outcomes. Then expand time windows, tasks, and social contexts batch by batch without loosening the ``same locus'' criterion because the results are inconvenient.

Let
\begin{center}
\resizebox{0.96\linewidth}{!}{$\displaystyle
B_m=\text{number of distinct persistent singular-like loci still required after the first }m\text{ batches}.
$}
\end{center}
Compare at least a finite-saturation model, in which the new-locus rate eventually decays, with a continuing-growth model allowing new future-relevant loci to continue appearing. A family/type count $K_m$ describes only coarse classification and cannot substitute for $B_m$.

\textbf{Support.} If the new-locus rate declines within the predeclared range, old loci are repeatedly reidentified in held-out contexts and continue predicting $\Delta C^+$, and continuing-growth models show no stable advantage, $H^{\mathrm{psych}}_{\mathrm{fin}}$ receives qualified support.

\textbf{Negative result.} If new loci continue appearing as material, tasks, and contexts expand; the old finite set loses predictive value; and the result is robust across alternative identity rules and complexity penalties, the present finite singular-boundary-germ translation exits. This does not refute a source-side AF-by-discrete groupoid. It says only that this psychological architecture lacks support.

A bounded-type-inspired proxy remains on a separate ledger. For every predeclared resolution $r$, predeclare comparable episode/block units and let $N_r$ be the maximum number of singular-like connector/interface occurrences retained within any one such unit. Lack of uniform control over this per-unit maximum weakens only the singular-sector per-unit uniform-control proxy. It is neither the general mathematical groupoid condition nor the full tile-inflation package: it omits regular boundary data and does not operationalize the $|V_n|,|E_n|$ clauses.

\section{Eight Designs Are Not Eight Passes into Mechanism Theory}

Support from these experiments is local. Success in A shows only that ``again'' may contain structural information beyond similarity. Success in B supplies one empirical footing for cross-scale episode organization. Success in D shows only that return-completion and continuation can separate. Success in G and H1/H2 begins to support organizing passage and a qualified finite singular-interface subarchitecture. None turns a mathematical object into a cognitive mechanism by itself.

Failure, however, must matter. A reasonable minimal order is: begin with A, D, and F to test the three central separations; then B, G, and H to test cross-scale and multi-interface geometry; E separately carries the pressure on the Third-Order claim and the corrective-agenda trap; C links return judgment to memory/prediction but remains correlational until intervention evidence exists.

\begin{stopbox}[Stop line for this chapter]
These designs merely rewrite correspondences as recordable, comparable, falsifiable studies. Executable does not mean validated. If a simpler model explains the result sufficiently, use the simpler model. The cross-domain value of AF-by-discrete can come only from additional resolution.
\end{stopbox}

\chapter{From Structural Mapping to Mechanism: How Evidence Must Be Kept on Separate Ledgers}
\label{chap:evidence-ledgers-en}

\par\medskip
\noindent\textbf{Why this chapter matters}\par
Chapter 38 rewrote many structural correspondences as executable studies. At this point, it is easy to make the next mistake: once an experiment begins to support a structural distinction, we write, ``we found the psychological mechanism.'' That is too fast.

A psychology reader can hold onto four verbs: \emph{describes}, \emph{predicts}, \emph{causes}, \emph{implements}. A model can describe data well and predict new data, yet still not show what causes what or how the relation is implemented in the brain, body, or social world. Here we keep the four verbs from standing in for one another.
\par\medskip

\section{Begin with the Psychological Question: The Same Data Can Support Sentences of Different Strength}

Suppose Study H produces an elegant result. Without telling participants where candidate singular-like interfaces are, repeated continuation failure, distant-context recruitment, concentration of return endpoints, and strategy shifts cluster at a set of positions that can be reidentified across material; those loci continue predicting later segmentation and action change in held-out episodes.

That is important evidence, but it supports at least four different strengths of claim:
\begin{enumerate}
\item \textbf{Descriptive:} repeatable positions co-occur with several behavioral changes.
\item \textbf{Structural/predictive:} a finite singular-interface/persistent-locus model organizes and predicts those changes better than ordinary surprise, frequency, or semantic-similarity baselines.
\item \textbf{Causal/mechanistic:} independently manipulating an organizing process actually changes return judgment or continuation.
\item \textbf{Implementational:} the causal chain is realized by specified neural, bodily, or social processes.
\end{enumerate}
The first three are already different claims; the fourth is another category again. Held-out prediction can refute ``this structure has no empirical content,'' but cannot automatically answer ``is there a boundary detector in the brain?'' Likewise, one neural-imaging correlation cannot prove in reverse that the formal relation is the correct psychological structure.

The discipline is
\[
\boxed{\text{describes}\neq\text{predicts}\neq\text{causes}\neq\text{implements}.}
\]

\section{Five Claim Levels Remain Separate, but Each Becomes More Concrete after Part V}

The Translation Protocol requires formal proof, experimental data, phenomenological description, and historical material not to impersonate one another \citep{translation2026}. Combined with the case protocol, the project keeps five claim levels:

\begingroup\small
\begin{longtable}{>{\raggedright\arraybackslash\bfseries}p{0.14\textwidth}>{\raggedright\arraybackslash}p{0.24\textwidth}>{\raggedright\arraybackslash}p{0.29\textwidth}>{\raggedright\arraybackslash}p{0.13\textwidth}}
\toprule
Level & What can be said & Required evidence & Current status\\
\midrule
Level I & two phenomena resemble one another in expression or image & analogy, explicit failure limits, and statements of what cannot be inferred & allowed but minimized\\
Level II & independently defined relational skeletons show constrained structural compatibility & native-domain fidelity, explicit mapping, counterexamples, exit conditions & overall status of this book\\
Level III & independently operationalized psychological variables satisfy predeclared formal relations and add prediction on new data & preregistered coding, independent observables, held-out tests, baseline comparison, model selection & next goal for local modules; not yet established\\
Level IV & some organizing variables belong to a real cognitive causal chain, such that changing them systematically changes return, segmentation, or continuation & manipulation, mediation, alternative-mechanism controls, replication & not claimed here\\
Level V & AF-by-discrete is the ontology of consciousness or its uniquely real architecture & cross-level ontological argument, competing ontologies, strong counterexamples & explicitly not pursued\\
\bottomrule
\end{longtable}
\endgroup

Level III is the easiest to misunderstand. It does not mean that a mathematical word has been measured. Ordinary psychological variables are defined first; then the data actually exhibit the declared relational constraints. If return judgment depends on an episode relation after semantic similarity is controlled, or a finite singular-locus model continues to outperform simpler models on new data, we obtain an empirically constrained structural model, not a mathematical machine hidden in the brain.

\section{From Level II to Level III: What Succeeds Is the Constrained Relation, Not the Vocabulary}

The minimum threshold can be compressed into six requirements:
\begin{enumerate}
\item \textbf{Ordinary psychological variables first.} Episode boundary, self-evaluation, confidence, memory retrieval, prediction, and action choice must be definable and measurable without AF-by-discrete vocabulary.
\item \textbf{Freeze the mapping first.} Which variables form case proxies, how return hypotheses are coded, how persistent interface loci are reidentified across material, and whether loci are grouped into families must be fixed before the main results.
\item \textbf{Formal relations must genuinely be able to fail.} If every dataset can be explained by adding a hidden sector, hidden germ, or new episode cut, the work remains Level-II narrative.
\item \textbf{Compete with simpler models.} Semantic similarity, frequency, prediction error, event segmentation, metacognition, and decentering are real baselines \citep{fleming2012,fleming2012a,bernstein2015}.
\item \textbf{Predict new data.} Beauty on a discovery set is insufficient. Held-out episodes, other participants, or new tasks should preserve some of the relation.
\item \textbf{Negative results rewrite the theory.} If new data keep forcing distinct persistent loci, do not save the model only by raising the permitted locus count or sorting them into a few families. If ordinary attention absorbs the root shift, stop treating root as an independent psychological construct.
\end{enumerate}

Level III therefore tests
\[
\begin{gathered}
\text{ordinary observables}\longrightarrow\text{frozen coding}\\
\Downarrow\\
\text{formal relation}\longrightarrow\text{held-out consequence}.
\end{gathered}
\]
This is why the book keeps insisting that boundaries must be forced out by the system. If only someone who knows the theoretical answer can identify the boundary, it is not yet an empirical variable; it is a redescription by the researcher.

\section{Accurate Prediction May Still Be a Good Marker Rather than a Mechanism}

Suppose an interface score $I(t)$ reliably predicts the next strategy shift. At least three explanations remain possible:
\begin{enumerate}
\item $I(t)$ captures a process that genuinely participates in organizational transition;
\item $I(t)$ is a reliable marker of another unmeasured variable, such as surprise, task difficulty, or arousal;
\item $I(t)$ is a useful compression statistic organizing outcomes produced by many causes, without itself ``doing'' anything.
\end{enumerate}
The third possibility matters. Formal usefulness does not require every variable to be a causal agent. Contour lines on a map can be accurate without causing a mountain to rise. Likewise, a boundary-like interface can be a position where organization change repeatedly appears without becoming a psychological organ ``responsible for switching.''

We therefore reject
\[
\text{predictive importance}\Longrightarrow\text{causal importance}\Longrightarrow\text{mechanistic entity}.
\]
Each arrow requires additional evidence.

\section{What Would Qualify for Level IV? A Manipulable Causal Chain Must Become Visible}

If future work claims that an organizing boundary causes return judgment, at minimum one needs a manipulation that changes the candidate organizing process while preserving other major explanations. For example, after controlling content similarity, event frequency, and surprise, manipulate transition structure so that distant background must be recruited; then test whether return judgment, episode segmentation, and continuation change in the predeclared direction.

More complete mechanism evidence must answer three questions:
\begin{enumerate}
\item \textbf{What exactly did the intervention change?} The organizing relation itself, or general attention, load, arousal, and demand characteristics?
\item \textbf{Does the intermediate chain appear?} If the hypothesis says ``background recruitment changes $\to$ completion judgment changes $\to$ continuation changes,'' measure the stages separately rather than only comparing final affect ratings.
\item \textbf{Can alternative mechanisms explain the result?} If a simpler prediction-error, memory-retrieval, or metacognitive account gives the same result, AF-by-discrete has no mechanism priority.
\end{enumerate}
Likewise, one cannot claim that third-order observation changes a psychological cycle merely because participants given a Third-Order instruction later feel calmer. The result could come from attention shift, expectancy, linguistic reappraisal, demand, or ordinary decentering. What would matter is a predeclared structural change in the relation among observation position, corrective agenda, meta-checking, and continuation that simpler variables cannot absorb.

\begin{warning}[Hard Stop Line for Level IV]
Without independent variables, an intervention chain, alternative-mechanism controls, and repeatable support, every use of ``mechanism'' in this book remains a future hypothesis. In particular, boundary germ, return residue, AF core, and sector passage cannot be translated directly into brain regions, neural circuits, trauma traces, personality components, or hidden psychological organs.
\end{warning}

\section{``Third-Order Consciousness Works'' Requires Special Protection against Theory-Induced Phenomena}

Third-Order Consciousness has a special evidence risk. Once participants learn phrases such as ``the observer is also in the structure'' or ``do not confuse a local projection with the whole,'' they may rewrite their reports to meet the expectations of the theory. The experiment may then be observing theory compliance.

Future research therefore needs three special forms of evidence.

First, \textbf{spontaneous appearance in untrained conditions}. Without teaching this vocabulary, do participants already record ``how am I evaluating?'' together with the urge to solve the discomfort quickly and the act of checking whether the method worked?

Second, \textbf{transfer}. If Third Order appears only on tasks that repeat the original sentences of the theory and disappears in new relational, memory, or decision tasks, it looks more like learned discourse.

Third, \textbf{counterintuitive outcomes}. A constrained model must allow trained participants to produce results that violate the expectation of the theory, such as open observation failing to reduce meta-checking or corrective agendas failing to narrow continuation. Only if such outcomes are permitted to remain has the theory not trained participants into becoming its own proof.

This turns the Third-Order Consciousness trap into an evidential discipline: theory invocation itself must be a variable, not an unconditional researcher intervention.

\section{How Neural Evidence Enters: Seek Constraints on Implementation, Not Brain Regions Named after Mathematical Terms}

If future work adds EEG, fMRI, physiology, or lesion data, the formulations to avoid are ``we found the boundary region'' or ``this network is the AF core.'' Neither follows from this book and neither is granted automatically by Level III or IV.

A better question is: once a structural distinction has been independently established behaviorally, which neural or bodily processes constrain its implementation? For example, do trials behaviorally classified as the same root shift share a repeatable network reconfiguration; when background retrieval is experimentally manipulated, which physiological measures align with its temporal order; do implementation variables explain individual differences that behavior alone does not?

Two multiplicities remain:
\[
\boxed{
\begin{gathered}
\text{one structural relation can have multiple implementations;}\\
\text{one neural process can participate in multiple structural relations.}
\end{gathered}}
\]
Neural data can constrain a mechanism theory; they cannot stamp a mathematical noun as biologically real. Searching the brain for a ``boundary'' before the behavioral structure is independently established simply copies an unvalidated translation into another domain.

\section{Social and Historical Material Has Its Own Evidential Authority: It Is Not Noise}

The treatment of second-order subjecthood in this book has always included role, language, others, and history. When mechanism is discussed, those materials cannot suddenly be demoted to noise to be controlled away.

If an organizing position exists only within particular relational norms, professional evaluations, or social roles, interviews, historical records, interaction logs, and institutional context can constrain how that position acquired its evaluative rules and action permissions. Such evidence can restrict the interpretation of a psychological task, while still not proving a cognitive causal chain on its own.

Conversely, a stable laboratory return effect cannot declare social generation irrelevant. A more appropriate ledger is: social-historical material explains where some positions and rules come from; behavioral experiments test how those rules organize current judgment; mechanism studies then ask what happens when one process is changed. The sources can constrain one another, but none owns the whole.

This is the projection principle again: a window of evidence can be true without speaking for the global system.

\section{An Upgrade Ledger: What Future Results Would Allow Us to Say, at Most}

\begingroup\scriptsize
\begin{longtable}{>{\raggedright\arraybackslash}p{0.31\textwidth}>{\raggedright\arraybackslash}p{0.25\textwidth}>{\raggedright\arraybackslash}p{0.24\textwidth}}
\toprule
Future result & Claim that can be upgraded & What still cannot be said\\
\midrule
Participants stably and repeatably report ``again'' & a researchable subjective return judgment exists & a formal return witness exists\\
After similarity is controlled, a predeclared relation predicts return/continuation & a particular Level-III structural relation receives support & the relation is a causal mechanism\\
Distinct persistent singular-like loci tend toward one stable finite set in held-out material, while singular-sector occurrence counts are tested separately across resolutions & finite singular-boundary-germ correspondence gains empirical footing & full AF-by-discrete boundary support or mathematical bounded type has been empirically validated\\
Manipulating a candidate organizing process changes the intermediate chain in the predeclared direction & a bounded mechanism hypothesis receives support & the result applies to all consciousness or all contexts\\
Neural/bodily measures show temporal and interventional consistency with an established causal chain & implementation level becomes constrained & a brain region can be called a germ, boundary, or AF core\\
Third-Order instruction changes meta-checking beyond metacognition/decentering baselines & a third-order process remains worth studying as an independent candidate & ``Third-Order Consciousness has been proved'' or is a higher personality rank\\
\bottomrule
\end{longtable}
\endgroup

The table does not diminish the theory. It tells the theory what each step has genuinely earned. Stronger evidence rules out more alternatives, but no attractive result automatically licenses the language of the next level.

\section{Chapter Conclusion: Mechanism Is Not a Prize Awarded to a Beautiful Mathematical Structure}

Formalization matters because it forces more difficult questions: what must be measured independently? Which relation should remain on new data? Which failure deletes the correspondence? If those questions receive support, the empirical status of the model rises.

But ``mechanism'' is not a medal awarded when a formal model becomes sufficiently elegant. It carries a different evidential burden: manipulable variables, traceable intermediate chains, competition among mechanisms, cross-task replication, and independent constraints on implementation.

The evidence discipline can therefore be frozen as
\begin{center}
\resizebox{0.96\linewidth}{!}{$\displaystyle
\boxed{\text{structural mapping supported}\not\Rightarrow\text{psychological mechanism established}\not\Rightarrow\text{neural implementation found}.}
$}
\end{center}

\begin{stopbox}[Stop line]
If future experiments support only structural predictions, the book may move from Level II toward local Level III. That does not grant Level-IV mechanism language. If intervention and mediation evidence later appears, mechanism claims must still remain limited to the variables, tasks, and populations actually studied. Success at one level never signs in advance for the next.
\end{stopbox}

\chapter{Closing Part V: What Cases Can Do and What They Cannot Do}
\label{chap:part-v-close-en}

Part V has returned the structural translation to experience. It does not ``validate AF-by-discrete.'' It requires each important correspondence to say what it records, which baselines it competes with, what it predicts in new material, and when it exits.

\section{What Cases and Experiments Actually Accomplished}

Five empirical disciplines can now be frozen:
\begin{enumerate}
\item record raw material, segmentation rules, structural coding, and outcome variables on separate layers;
\item measure repetition, recognized return, and continuation separately;
\item use different information for interface detection and continuation validation, preventing self-confirmation;
\item compare every complex structural account with simpler models based on similarity, frequency, event segmentation, arousal, metacognition, decentering, and related variables;
\item let negative results strike only the level they actually test, while forbidding the model from surviving indefinitely by adding hidden boundaries after failure.
\end{enumerate}

\section{The Finiteness Actually Activated on the Empirical Side}

The psychological side uses a hypothesis of the form
\begin{center}
\resizebox{0.96\linewidth}{!}{$\displaystyle
H^{\mathrm{psych}}_{\mathrm{fin}}:
\quad
\text{candidate persistent singular-like loci show finite-set-like stabilization in expanding material}.
$}
\end{center}
It is not the same proposition as the source-side AF-by-discrete condition, which also counts regular boundary points and is formulated over all infinite tiles relative to a fixed generator/presentation. Study H1 tests recovery when finite interfaces are part of the known generating structure. H2, by contrast, compares finite-saturation with continuing-growth models in longitudinal/open-world material without assuming finite ground truth.

\section{The Minimum Case-Level Meaning of Third Order}

Cases allow a person to hope that discomfort decreases, take action, leave danger, and repair real-world problems. What Third Order requires to enter the record is which current rule, identity, evaluative goal, or corrective agenda organizes observation and whether it exempts itself from audit.

This is why Study E separates metacognition, decentering, open structural audit, and corrective structural audit. Observer endogeneity already has earlier theoretical ancestors in second-order cybernetics \citep{vonfoerster2003}. Any empirical increment should therefore not be written as ``the first discovery that the observer is in the system.'' It would need to appear as a narrower difference: does an already identified organizing position regain unaudited governing authority, and can that meta-identification itself still be relativized in a new context?

\section{Boundary of Evidence Upgrade}

A case shows that material \emph{can be read} through a structure. Held-out prediction begins to support a constrained empirical model. Stable manipulation and mediation begin to enter mechanism. No step automatically upgrades into an ontology of consciousness.

Part V ends positive construction here. Part VI will only settle boundaries of explanation, failure, and remainder.

\part{Boundaries, Translation Remainders, and Future Problems}

\begin{plainbox}[Let the theory stop where it should]
This part inherits the methodological discipline of the Translation Protocol and the Geometry of Existence: every judgment must disclose its native domain, observation position, range of material, scale, evidence level, and exit conditions \citep{translation2026,geometry2026}. Part V gave empirical material the power to veto the model from below. Part VI goes one step further and removes structural privilege from the theory itself. What the model selects, what it ignores, which geometry it prefers, and before which results it must retreat all enter the record.
\end{plainbox}

\begin{plainbox}[The word ``boundary'' has two meanings in the title of this part]
The boundary points of Part II are strict mathematical objects. Under a fixed finite generating set and presentation, AF-by-discrete uses global finiteness of boundary points over all infinite tiles. In the title of Part VI, ``boundary'' mainly means the boundary of applicability, evidence, and translation. Unless the text explicitly says \emph{mathematical boundary point} or refers back to the Part-II definition, the two meanings are not interchangeable. A methodological statement that ``a theory has boundaries'' cannot prove a mathematical boundary; global finite boundary support cannot automatically tell psychology where consciousness ``has a fissure.''
\end{plainbox}

Part VI is therefore not a disclaimer attached to the end. It performs the final structural audit of the book. Every attractive correspondence developed earlier must now answer one last question: how far does it speak, and where does it begin to fall silent?

\chapter{The Task of Part VI: Why a Theory Must Actively Preserve Blank Space}
\label{chap:part-vi-task-en}

The first five parts established a structural language that can be defined, compared, and allowed to fail. Part VI does not enlarge the model. It asks a different class of questions: what never entered the interface, what received only a Level-II representation, and which mathematical facts cannot be replaced by empirical proxies. Knowing where to stop is itself part of a theory.

\section{Complete Source Geometry and the Present Psychological Interface}

For a fixed finite generating family $S$ and tile presentation, the source-side ambient structure used in this book is
\[
\partial_S^\infty
=
\bigcup_{T\in\mathcal T_\infty}\partial_S T,
\qquad
|\partial_S^\infty|<\infty.
\]
This counts all persistent boundary points, not boundary types. At least three kinds of positions must be separated inside it:

\begingroup\small
\begin{longtable}{>{\raggedright\arraybackslash\bfseries}p{0.22\textwidth}>{\raggedright\arraybackslash}p{0.39\textwidth}>{\raggedright\arraybackslash}p{0.29\textwidth}}
\toprule
Position & Source-side behavior & Present cross-domain status\\
\midrule
boundary-free & the AF infinite-tile piece itself supplies orbital/germ-local background & no psychological boundary analogue is built\\
regular boundary & $H_b=1$; ordinary reversible gluing is possible and the graph of germs carries no extra fibre structure over the orbital graph & no current psychological counterpart\\
singular boundary & $H_\xi\neq1$; nontrivial germ fibres/sheets may appear & source prototype for the present singular-boundary-like translation\\
\bottomrule
\end{longtable}
\endgroup

The psychological quantities $B_m$ and $N_r$ therefore pressure-test only the singular-like sector. They cannot directly prove or refute the complete AF-by-discrete or bounded-type conditions.

\section{After a Singularity Becomes Visible, One More Cut Is Needed}

Hausdorff, non-Hausdorff, and purely non-Hausdorff are not levels of consciousness. They classify the local separability structure of germ differences. A Hausdorff singularity allows a nontrivial difference to remain separable in a sufficiently small neighborhood. Non-Hausdorffness says that at least one real difference is approached arbitrarily closely by identity regions. pNH extends this property across all germ directions.

The intellectual genealogy continues to matter. ``The observer enters the system'' belongs to a reflexivity problem shared by Third-Order Consciousness, second-order cybernetics, and other traditions. A possible contribution of non-Hausdorffness is to refine that shared problem into two stages: becoming visible, and whether what has become visible can then regain a separable final territory. pNH is thus a mathematically generated benchmark, not a proof that Third-Order Consciousness had historical priority \citep{vonfoerster2003}.

\section{Three Kinds of Stopping}

Part VI uses only three reasons for stopping.
\begin{enumerate}
\item \textbf{Principled stop:} the current formal language does not define the variable at all.
\item \textbf{Evidential stop:} a structural question has been defined, but independent data are insufficient.
\item \textbf{Scope stop:} the mathematical direction is clear, but this book has not carried out the complete construction.
\end{enumerate}
Qualia, mineness, neural realization, and heterogeneous subject integration occupy different positions among these stops. They cannot all be renamed ``things the model will explain later.''

\chapter{What Does the Model Actually Explain? A Positive Boundary List}
\label{chap:positive-boundary-list-en}

Here ``explain'' means increasing structural resolution, not supplying a mechanism of consciousness. The positive results actually earned by the book can be compressed into seven items.

\begingroup\small
\begin{longtable}{>{\raggedright\arraybackslash\bfseries}p{0.25\textwidth}>{\raggedright\arraybackslash}p{0.65\textwidth}}
\toprule
Achievement & Statement currently licensed\\
\midrule
local background / interface & a large system can run mainly through an AF-like background while additional continuation information is concentrated on globally finite persistent boundary support; the psychological side tests only a singular-like analogue\\
repetition / return & recurrence and being organized by a subject as ``coming back'' are different judgments; return requires finite completion evidence\\
episode segmentation & how ongoing events become ``this one thing'' must retain entrance, exit, scale, and coding rule\\
root / position & a rooted projection can be locally real and effective without representing the whole system or complete subject\\
persistence / revision & an identified finite witness can remain traceable as context enlarges; this does not mean human memory cannot be reconstructed\\
return-completion / continuation & recognized return can remain while future prediction, action, and memory retrieval differ\\
reflexive organization & the positions and rules used in observation, segmentation, and evaluation can themselves enter audit; this general reflexive requirement overlaps with traditions such as second-order cybernetics, and the independent empirical increment of Third-Order Consciousness remains to be tested \citep{vonfoerster2003}\\
\bottomrule
\end{longtable}
\endgroup

Mathematics therefore supplies not an answer to ``what is consciousness?'' but a set of non-implications that ordinary language often collapses:
\[
\text{repetition}\not\Rightarrow\text{return},
\qquad
\text{return-completion}\not\Rightarrow\text{same continuation},
\]
\[
\text{root}\not\Rightarrow\text{whole system},
\qquad
\text{boundary}\not\Rightarrow\text{singular}.
\]
If these distinctions add no empirical resolution, they should retreat to the mathematical native domain. Only if they reliably improve held-out prediction can they begin to enter Level III locally. The book has not earned stronger psychological mechanism conclusions.

\chapter{Translation Remainder I: Mineness, Qualia, and Embodiment}
\label{chap:remainder-mineness-en}

\section{At Least Three Different Questions Are Hidden inside the Same Word ``Mine''}

The book has repeatedly used phrases such as ``my experience,'' ``subject-related position,'' and ``self-indexed history.'' The apparently simple ``mine'' must now be decomposed. At least three kinds of question cannot replace one another:
\begin{enumerate}
\item \textbf{Structural self-indexing:} which history, memory, commitment, evaluation, or continuation organization receives an episode? Why is today's anger connected to ``I did this last time too''?
\item \textbf{Phenomenological mineness / for-me-ness:} why does this pain, color, sound, or thought appear in a first-person manner for me? Why is experience not merely an item registered in a system?
\item \textbf{Concrete agency/ownership judgments:} was this action ``done by me,'' is this body ``my body,'' is this thought experienced as initiated by me? These questions neighbor mineness but have their own psychological and experimental structures \citep{gallagher2000,braun2018}.
\end{enumerate}
The first is precisely where this book can provide limited formalization. The second is the central remainder of the chapter. Parts of the third can be investigated through existing paradigms for body ownership, agency judgment, and error monitoring, but measurability does not make them identical with phenomenological mineness.

The crucial non-implication is therefore
\[
\boxed{\text{structural self-indexing}\not\Rightarrow\text{phenomenological mineness}.}
\]
An episode stably placed in ``my history'' has acquired an indexing position in the current organization; this can be studied indirectly through memory links, event segmentation, reports, and prediction. Phenomenological mineness asks a different question: why is experience first given as ``for me''? That question lies at the center of long-running debates about subjectivity and consciousness \citep{zahavi2005,gallagher2000}. The present AF-by-discrete model has not defined a variable for first-person givenness and cannot infer it from structural indexing.

\begin{plainbox}[If a variable has never been defined, increasing resolution will not create it]
If a formal language contains no representation of phenomenological mineness, making tiles finer, increasing graph-of-germs resolution, or recording more return witnesses cannot automatically produce a theorem about first-person givenness.
\[
\boxed{\text{finer structural resolution}\neq\text{new phenomenological ontology}.}
\]
Finer resolution can recover structural differences that were previously compressed. It cannot conjure a phenomenal dimension absent from the language.
\end{plainbox}

\section{Root Supplies a Coordinate; It Does Not Generate the Feeling of ``I''}

Choose a root $x$ in the mathematical model. From $x$ one can discuss rooted graphs, paths, traverses, returns, and how the position unfolds under finer germ resolution. What has been obtained is a distinguished reference position. It tells us where organization is rooted, not why the location is experienced as ``me.''

Root therefore cannot support three overreaches:
\begin{itemize}
\item being distinguished does not make it a metaphysical self;
\item supporting self-indexing does not make it a generator of phenomenological mineness;
\item different roots producing different rooted projections does not mean different conscious subjects have been mathematically identified.
\end{itemize}
A simple counterfactual makes the limit clear. Suppose two systems are isomorphic at the rooted labeled structure relevant to this book. At this level mathematics can say that the relevant path, return, boundary/germ, and continuation structures agree. It cannot infer from that isomorphism that the two systems possess the same subjective presence, the same feeling of ``I,'' or any phenomenal consciousness at all. Those conclusions require additional theory and evidence.

This is again the projection principle: root can be a real and useful local coordinate without acquiring ownership of the entire subject.

\section{Qualia: Return-Completion Can Be Described while Pain Still Hurts}

The return language supplies a strong structural distinction: a finite episode can complete a genuine return while continuation afterward differs. That distinction may prove empirically useful for rumination, relational conflict, habit, and learning.

But however precise return-completion becomes, it has not answered:
\begin{quote}
Why does pain appear with this painful quality? Why does shame burn and contract in this way? Why does red have this visual quality?
\end{quote}
Nagel's question of ``what it is like'' and Chalmers's discussion of the explanatory problem of phenomenal consciousness both warn against identifying functional, behavioral, or relational structure with subjective quality without further argument \citep{nagel1974,chalmers1995}. This book does not settle those philosophical disputes. It makes a smaller and stricter claim: the current formal language does not explain qualitative character.

Write, only as bookkeeping,
\[
\text{qualitative character of experience}\in U_{\mathrm{remainder}}.
\]
$U_{\mathrm{remainder}}$ denotes phenomena not yet admitted to the formal interface. It is not a secretly introduced third mathematical native domain and carries no additional mathematical structure.

This stop line also blocks a tempting mistranslation. Return residue $\rho_\xi(R)$ cannot mean ``the feeling that remains,'' and a more complicated germ fibre cannot mean ``a stronger experience.'' More sheets, different continuation, or different residue indicates structural resolution differences, not a natural unit of qualia intensity.

\section{Embodiment: The Body Is Not Background Waiting to Be Coded as an AF Core}

The body was present from the beginning of Part I: sensation, action, posture, interoception, pain, fatigue, and environmental coupling are not decorations later added to a subject. Psychology and cognitive science have their own research programs and measurement paradigms for interoception, body ownership, agency, and embodied selfhood \citep{craig2009,braun2018,gallagher2000}.

Part III extracted only a very limited relational skeleton for controlled correspondence, such as ``local activity can precede stable narrative labeling'' and ``a large background can be organized through finite-scale approximations.'' Nothing authorized
\[
\boxed{\text{body}=\text{AF core}.}
\]
The AF core is an approximately finite mathematical structure. A body contains metabolism, interoception, motor control, posture, pain, physiological regulation, and ongoing environmental coupling. Their object types, evidence sources, and failure modes differ. Compressing the body into an AF core would not make the mathematics more explanatory; it would delete the biological content Part I needs to preserve.

The book does not even decide what status embodiment ultimately has relative to these organizational relations. Bodily processes may be implementation conditions for some structures, may strongly constrain which episodes, root shifts, action options, and continuations are possible, or may play a more constitutive role in stronger embodied/enactive accounts. The present model remains neutral among these options. Any upgrade must introduce bodily variables and independent evidence; it cannot be accomplished by renaming an existing discrete object.

\section{An Untranslated Remainder Cannot Be Redrawn as a ``Deeper Boundary''}

Chapter 41 required the theory to distinguish model failure, insufficient evidence, and an undefined variable. Mineness, qualia, and full embodiment belong mainly to the third category. One move is therefore prohibited.

\begin{warning}[A Remainder Is Not a Hidden Singularity]
If the present model does not explain first-person givenness, qualitative character, or bodily mechanism, it cannot preserve itself by saying ``the real boundary is deeper'' and adding a hidden root, hidden germ, or new singular sector.

An unmodeled phenomenal dimension is not an undiscovered mathematical fissure.
\end{warning}

This discipline is especially compatible with the global finite boundary support of AF-by-discrete and finite singular subset. The book chose a formal framework on purpose, one in which fissures cannot spread everywhere. If qualia, embodiment, and mineness are all renamed boundary whenever explanation becomes difficult, the final part would rhetorically destroy the finite-support discipline established earlier.

A mature response is therefore not to force every remainder back into the diagram, but to retain an edge not filled by mathematics. It is not a wastebasket for failure. It is where the current theory stops possessing the phenomenon before the next body of evidence begins.

\section{First Class of Remainders: What Is Specifically Left Open}

Three long-term remainders are frozen:
\begin{center}
\resizebox{0.96\linewidth}{!}{$\displaystyle
\boxed{\text{phenomenological mineness},\quad\text{qualitative character / qualia},\quad\text{full embodied realization}.}
$}
\end{center}
``Frozen'' does not mean forever unresearchable. It means that the present structural model has not paid the evidential cost of answering them.

\begingroup\small
\begin{longtable}{>{\raggedright\arraybackslash\bfseries}p{0.25\textwidth}>{\raggedright\arraybackslash}p{0.34\textwidth}>{\raggedright\arraybackslash}p{0.31\textwidth}}
\toprule
Remainder & What the current structural model can do at most & What a future upgrade would still require\\
\midrule
phenomenological mineness & describe relations among self-indexing, rooted reference, and subjective report & an independent first-person/phenomenological theory and an empirical bridge capable of constraining it\\
qualia & distinguish episode, return-completion, and continuation while preserving reports of experience & theory and evidence linking structural variables to qualitative character while excluding simpler functional substitutes\\
embodiment & leave structural interfaces for event organization and action/continuation & interoceptive, sensorimotor, physiological, and environmental variables, independently measured in relation to structural processes\\
\bottomrule
\end{longtable}
\endgroup

Remainder does not mean that these questions are unscientific or that structural models are useless. It means only that structural resolution already obtained cannot answer in place of first-person givenness, qualia, and embodiment. The structural world remains available for research; so does the experiential world with its color, pain, and bodily weight.

\chapter{Translation Remainder II: Affect, Neural Implementation, Social Production, and Normativity}
\label{chap:remainder-affect-en}

\section{Affective Valence Is Not Return Residue}

In psychological cases, returns often accompany anger, shame, anxiety, relief, closeness, or aversion. Because these feelings are so salient in life, it is easy to read mathematical $\rho_\xi(R)$ as ``emotional residue.'' That move must stop.

Mathematical residue records a later structural difference left by a witnessed return under germ-resolved continuation. It carries no pleasant/unpleasant value, arousal, motivation, or bodily feeling by itself. The strongest licensed statement remains
\[
\boxed{\text{return-completion and continuation can be structurally separated},}
\]
not
\[
\rho_\xi(R)=\text{emotional residue}.
\]

The psychological side must still separate ordinary variables: pleasant versus unpleasant feeling, arousal, approach versus avoidance, bodily tension, and intended next action. These can correlate with one another and jointly alter segmentation, memory retrieval, or action choice, but the relations require independent measurement.

\begin{principle}[Negative Valence Is Not Structural Failure]
``Feeling bad'' is first a real affective fact. It can signal danger, loss, conflict, or unmet need and can contribute to reasons for action. It cannot by itself prove that a return is wrong, that a boundary has disappeared, or that third-order observation ``failed.'' Feeling comfortable cannot automatically prove structural change either.
\end{principle}

This connects to the distinction made in Part V between wanting relief and a corrective agenda. A person can reasonably want suffering to decrease. What requires audit is whether that wish acquires structural privilege and becomes ``I must use an observation technique to eliminate this state immediately.'' Hence
\[
\boxed{\text{wellbeing orientation}\neq\text{instant structural success criterion}.}
\]
If future research finds that affective variables systematically predict episode cutting, return recognition, or continuation change, that would establish a new mechanism bridge. It would not rename valence as residue.

\section{Affect Can Constrain Structure without Being the Structure Itself}

Leaving affect in the remainder does not expel emotion from the theory. A more mature future model may find that high arousal shortens temporal integration windows, shame changes which memories are recruited, fear renders some continuation options almost inaccessible, or closeness changes which episode a behavior enters. Affect would then impose real constraints on structure.

The direction must remain clear:
\[
\text{affective state}\longrightarrow\text{possible constraint on organization}
\]
does not mean
\[
\text{affective state}=\text{boundary / return / residue}.
\]
A structural model can help ask which organizing relation affect changes. Answering that question is not an explanation of why the affect has the qualitative feel it has; that returns to qualitative character from Chapter 43.

\section{Neural Implementation: Two Bridges Still Separate a Formal Relation from a Brain Process}

Parts II and IV use discrete mathematical objects. Real neural activity includes continuous-time dynamics, distributed networks, neuromodulation, learning plasticity, and multiscale physiological constraints. Assigning a vertex to a brain region, a boundary germ to a synaptic event, or the AF core to the ``default network'' has no support from the definitions or evidence in this book.

A real upgrade requires at least two independent bridges. The first goes from formal relations to repeatable psychological/behavioral processes: can predeclared segmentation, return judgment, root shift, or continuation profiles be measured reliably across new tasks and participants? The second bridge then goes from those established processes to neural realization: what temporal dynamics, network interactions, or physiological manipulations explain them and change them under intervention?

The evidential order is
\[
\begin{gathered}
\text{formal relation}\\
\Downarrow\\
\text{independent psychological process}\\
\Downarrow\\
\text{neural/body realization}.
\end{gathered}
\]
A beautiful correlation at a later level cannot substitute for an earlier bridge.

Implementation is also rarely one-to-one:
\begin{center}
\resizebox{0.96\linewidth}{!}{$\displaystyle
\boxed{\text{structural role}\neq\text{unique neural location},\qquad\text{neural correlate}\neq\text{formal object}.}
$}
\end{center}
Existing metacognition research supplies independent paradigms for confidence, error monitoring, and neural correlates \citep{fleming2012,fleming2012a}; decentering has its own measures and process models \citep{bernstein2015}. Any future mechanism-level work must compete with them rather than define Third Order as a capacity visible only to this theory.

\section{Social Production: The Categories Used by the Present Subject May Come from a Larger History}

Part I used Mead, Goffman, Berger--Luckmann, and related traditions to show how subject positions are shaped by language, roles, expectations of others, and institutional classifications \citep{mead1934,goffman1959,berger1966}. Part V mostly followed one participant's trajectory to control scope. That design choice does not mean social structure is merely external background.

Real-world ``again'' judgments often already carry social categories. Consider ``you are disrespecting me again.'' What is reidentified is not just an action. ``Respect'' can depend on family role, professional hierarchy, gender norms, cultural expectation, or platform etiquette. The category used to complete the return judgment may have been supplied by the social world long before the current episode.

Two questions must therefore be separated:
\[
\begin{gathered}
\text{How does the current system use a category to organize experience?}\\
\text{How was that category produced, authorized,}\\[-1mm]
\text{and sustained in social history?}
\end{gathered}
\]
The first can enter the present structural description. The second requires social theory, historical material, and institutional evidence. AF-by-discrete can record how a classification participates in return, root shift, or continuation; it does not thereby explain who possesses naming authority, which differences institutions amplify, which experiences enter public memory, or which remain unspeakable.

Social origin also does not automatically equal a new mathematical boundary. A category coming from an institution shows that a structural variable may have sources outside the individual. A singular-boundary-like claim must still be supported independently by repeated system pressure, persistent locus, and continuation evidence.

\section{Developmental Origin: A Mature Structural Map Cannot Impersonate Its Own Developmental History}

The book mainly begins from adult material in which language, episode segmentation, autobiographical memory, and subject history are already available. It studies how, given an organizational grammar, events are segmented, returns are recognized, parts of the past become relevant again, and continuation is organized. This is closer to a structural snapshot.

An adult's ability to use a rule today does not mean the rule was present from the beginning. How children develop autobiographical memory, first-person language, role understanding, and long-term self-persistence is a different problem. Developmental work on family narrative, language, and social interaction already shows clear histories of acquisition \citep{nelson2004,conway2000}.

Hence
\[
\boxed{\text{mature organization}\neq\text{developmental origin}.}
\]
A real developmental model would have to track which distinctions appear first, which return categories are learned later, whether persistent loci change through development, and how a previously absent organizing position stabilizes. This book does not derive childhood stages from Bratteli levels and never treats a mathematical level as an age.

\section{Normativity and Ethics: A Structural Diagram Cannot Sign ``Ought'' on Someone's Behalf}

Seeing structure can change action, but it cannot by itself determine obligation. Third-order observation may reveal how anger brings particular parts of the past back into view, how a corrective loop restarts, or how one relationship narrative compresses different episodes into the same ``always.'' Those distinctions can matter enormously. The mathematics cannot conclude
\begin{quote}
``Therefore I should forgive the other person.''
\end{quote}
or
\begin{quote}
``Therefore I must leave this relationship.''
\end{quote}
Action also involves risk, responsibility, commitments, bodily safety, the interests of others, resource constraints, and which consequences a person is willing to bear. Existential traditions especially remind us that theory does not abolish choice and responsibility in concrete situations \citep{sartre2007,reynolds2024}.

The distinction from Part V matters here too. ``I hope to suffer less'' can be a real and reasonable directional wish; Third Order does not cancel it. But the wish does not specify one psychological technique and cannot make ``if I am still distressed, observation failed'' into a law. Likewise, structural freedom from an old loop does not guarantee that every resulting choice is ethically correct.

The book retains only the weaker contribution
\[
\boxed{\begin{gathered}
\text{structural clarification can inform action}\\[-1mm]
\text{but does not determine the ought}
\end{gathered}}.
\]
If Third-Order Consciousness is used to exempt responsibility---``it is only structure, so I do not need to apologize'' or ``the observer is structure, so nobody is responsible''---reflexive observation has become theoretical self-justification.

\section{Four Independent Research Layers: They Intersect but Do Not Possess One Another}

\begingroup\small
\begin{longtable}{>{\raggedright\arraybackslash\bfseries}p{0.23\textwidth}>{\raggedright\arraybackslash}p{0.36\textwidth}>{\raggedright\arraybackslash}p{0.31\textwidth}}
\toprule
Layer & What the current model can ask & What it cannot infer automatically\\
\midrule
affect & whether affective state changes segmentation, return judgment, or continuation & valence/arousal is residue or boundary\\
neural/body realization & whether established psychological structural processes have stable physiological implementation constraints & a brain region, network, or physiological event is a germ, root, or AF core\\
social and developmental production & where categories, roles, and organizing rules come from and how they change & social origin is itself a boundary; mature structure is a developmental stage\\
normativity & how structural clarification changes visible action options and consequences & a mathematical relation directly determines what should be done\\
\bottomrule
\end{longtable}
\endgroup

These layers can constrain one another. Affect can change the integration window; social categories can shape return language; development can determine which root-like positions are available; neural/body implementation can limit which continuations are physically reachable; ethical judgment can determine whether a person chooses a reachable path. Interaction does not imply ontological identity.

\begin{warning}[Second Class of Stop Lines]
This manuscript does not reduce affective valence, neural/body realization, social production, developmental acquisition, or normative judgment to existing AF-by-discrete objects. These are not blanks waiting to be filled with mathematical vocabulary. They are independent research layers with their own variables, evidence, and failure modes. Future bridges require new measurements and new mechanisms at each step; an upgrade cannot be accomplished by giving root, boundary, germ, return, or residue a bolder name.
\end{warning}

The conclusion is not ``we know nothing outside structure.'' We now know which questions structural language can sharpen and which questions must return with their own material. A map can show how roads intersect; it does not thereby live the pain, body, institutions, and responsibilities encountered on those roads.

\chapter{Failure Conditions: What Observations Would Force the Model to Exit?}
\label{chap:failure-matrix-en}

A cross-domain model must permit separate failure at separate levels. The matrix below replaces all-or-nothing judgment.

\begingroup\small
\begin{longtable}{>{\raggedright\arraybackslash\bfseries}p{0.10\textwidth}>{\raggedright\arraybackslash}p{0.36\textwidth}>{\raggedright\arraybackslash}p{0.34\textwidth}}
\toprule
Failure & Observation & What must be withdrawn\\
\midrule
A & episode cuts show no stable relation across coders/scales over time & episode--traverse correspondence\\
B & ``again'' is fully explained by similarity/frequency & return--recognition correspondence\\
C & one center is insufficient, but finitely many organizing positions add stable prediction & only the single-position hypothesis\\
D & $N_r$ keeps growing with resolution & singular-sector per-unit uniform-control proxy; no verdict on mathematical bounded type\\
E & expanding material continually forces new persistent singular-like loci and $B_m$ does not stabilize & $H^{\mathrm{psych}}_{\mathrm{fin}}$ / finite singular-interface translation\\
F & multi-scale segmentations are mutually unconstrained stories & decomposition/persistence translation\\
G & AF-by-discrete-inspired model does not outperform a simple baseline & the corresponding Level-III candidate\\
H & the target phenomenon appears only after learning the theoretical vocabulary & theory-dependent interpretation; demand effects must be recontrolled\\
I & every counterexample is absorbed by adding another hidden boundary & the entire present cross-domain model loses falsifiability\\
\bottomrule
\end{longtable}
\endgroup

Failures D and E must remain separate. $N_r$ is only the maximum singular-like connector/interface burden within a predeclared comparable empirical unit at resolution $r$; $B_m$ tracks distinct persistent singular-like loci across expanding material. Neither includes source-side regular boundary data, and neither operationalizes the Bratteli level-size controls of the tile-inflation package. Positive and negative results therefore cannot directly declare the mathematical \AFBD{} or the full bounded-type tile-inflation package true or false.

Finite data always contain only finitely many observed positions. What can pressure-test a finite-set-like hypothesis is not ``this batch found many,'' but continued emergence of nonmergeable new loci as material expands under frozen identity rules and analysis scale, together with loss of held-out predictive value by the old finite set.

Exit follows one principle: stop first at the level struck by the counterexample. Failure of a local correspondence does not alter a mathematical theorem. Failure of a strong claim cannot be hidden by retreating to a weaker claim while continuing to advertise the original. If even the weakest testable structure fails to outperform simple explanations, the present model should genuinely yield.

\chapter{Reverse Translation and Evidence Ledgers: Preventing Mathematics from Rewriting Psychological Facts}
\label{chap:reverse-translation-en}

Formalization must still be able to return to the raw material. Part III never established a one-to-one map between psychological and mathematical objects, so ``reverse translation'' is not inversion of a function. It audits which segmentation, coding, and structural choices were applied to the original material and what mathematics added beyond them.

\section{Four Object Layers and Six Kinds of Evidence}

Compress the empirical path of the book as
\[
D\longrightarrow C\longrightarrow M\longrightarrow H,
\]
where $D$ is raw material, $C$ is disclosed case coding/segmentation/locus identity, $M$ is formal correspondence, and $H$ is a prediction or mechanism claim. There is no equality sign among them.

Correspondingly, mathematical proof, intellectual/phenomenological description, existing psychological evidence, case coding, structural translation, and mechanism/implementation evidence can speak only for their own windows. A simple restoration test is to temporarily remove the words root, boundary, germ, return, and residue. Can the researcher still explain in ordinary psychological language what was observed, how it was segmented, what would change future prediction, and what result would revoke the judgment? If not, formalization has begun to loop back onto itself.

\section{Finiteness Especially Cannot Be Upgraded in Reverse}

The psychological $H^{\mathrm{psych}}_{\mathrm{fin}}$ asks only whether candidate singular-like loci stabilize in expanding material. The source-side AF-by-discrete finiteness counts all persistent regular and singular boundary support. Success or failure of the former cannot rewrite the latter without a new bridge theorem.

Likewise, finitely many families/types do not imply finitely many distinct loci, and singular-sector per-unit uniform control of $N_r$ is not mathematical bounded type. It omits regular boundary data and leaves the tile-inflation package's $|V_n|,|E_n|$ controls untranslated. ``We have not yet found more'' remains an open empirical conclusion.

\section{Weakening at the Source Must Propagate Downstream}

The result of critical reconstruction cannot be recorded only in Part I. If upstream investigation shows that an intellectual claim has earlier ancestors, every downstream attribution must change accordingly. In this book, observer endogeneity and the anti-final-observer problem are at least partly shared among Third-Order Consciousness, Krishnamurti, second-order cybernetics, and related traditions \citep{vonfoerster2003}. Later pNH discussions, Studies E/F, and the conclusion can therefore present non-Hausdorffness only as a mathematical refinement of a shared problem, not as mathematical ``proof'' of a proposition uniquely owned by Third-Order Consciousness.

\section{Theory Training Can Enter the Loop Too}

A participant saying ``the observer is also in the structure'' may first mean only that the vocabulary was learned. If theory training produces repeated checking of ``am I observing correctly?'' or ``have I really let go of the goal?'' the theory itself has become a new organizing variable. Future experiments must therefore record theory exposure, theory invocation, and meta-checking and test primary structural predictions on tasks that do not depend on repeating the vocabulary.

Revision still follows the minimum-change principle. If $H$ fails, first withdraw the mechanism claim. If $M$ fails, withdraw the correspondence. If $C$ is unstable, repair the coding. Modify $D$ only when the raw facts themselves have been reverified. The final test of formalization is simple: does it help us see the original material more clearly, rather than merely helping us see our own theory?

\chapter{Future Problem I: From the Multi-Boundary Prototype to Heterogeneous Subject Integration}
\label{chap:heterogeneous-integration-en}

\section{Draw the Stop Line Correctly First: Multi-Boundary Already Exists; the Future Problem Is Heterogeneous Integration}

Parts III and IV repeatedly used the single-site case because one distinguished singularity makes return, germ resolution, and witness persistence easiest to see. But the single-site case is a benchmark, not the definition of AF-by-discrete.

In the fixed source setup of this book, if
\[
\partial_S^\infty
=
\bigsqcup_{T\in\mathcal T_\infty}\partial_S T,
\qquad
|\partial_S^\infty|<\infty,
\]
then the positions that genuinely carry boundary data across all infinite tiles already form a finite set. Hence
\[
\boxed{\text{multi-boundary}\neq\text{boundary proliferation}.}
\]
Here one may have $1<|\partial_{S,\mathrm{sing}}^\infty|<\infty$; the standing AF-by-discrete condition remains $|\partial_S^\infty|<\infty$. The singular subset is only the part that participates in the heterogeneous germ integration studied in this chapter.

Kuang's v1, Subsection~5.2, already uses multiple infinite boundary points / bridge types to construct
\[
\Phi(\tau)=(\tau_1,\ldots,\tau_k),
\qquad
\operatorname{typ}\Phi(\tau)=x_1\cdots x_k,
\]
thereby providing a type-word return-coding prototype that passes across multiple boundary sectors \citep{kuang2024v1}. So the question ``can one return pass through several sectors?'' is not something this chapter needs to invent; in that source prototype, at least, the answer is already yes.

What is genuinely missing for the future is this: if those sectors differ in local germ structure, orbit membership, residue behavior, and large-scale continuation, what data are sufficient to describe how they coexist?

\begin{principle}[The Future Problem Is Not to Reimagine Finitely Many Fissures as Infinitely Many Fragments]
The AF-by-discrete global finite-boundary-support condition remains frozen throughout this chapter; it counts the total of regular and singular boundary points. Future heterogeneous integration studies only the finite subset $\partial_{S,\mathrm{sing}}^\infty$, adding relations and decorations at these persistent singular loci. Mathematically, one leaves AF-by-discrete only when total boundary support itself can no longer remain finite. Psychologically, if singular-like loci fail to stabilize, one withdraws only the singular-germ translation used in this book.
\end{principle}

\section{Why ``Several Positions'' Still Does Not Equal ``One Subject''}

Part V repeatedly distinguished root, boundary, and global organization. A root gives only the current rooted projection; a boundary marks only a place where continuation needs additional gluing. Even listing finitely many roots or boundaries does not yet produce a ``whole subject.''

This is exactly consistent with the projection principle of Chapter 34:
\[
\boxed{
\begin{gathered}
\text{one position cannot occupy the whole;}\\
\text{a list of positions is still not the whole.}
\end{gathered}}
\]

The goal of heterogeneous subject integration therefore cannot be to locate ``the truly central root,'' nor can it be merely to write
\[
\xi_1,\ldots,\xi_k
\]
side by side. What must be explained is which local structures can reach one another; which are related only indirectly through finite-scale approximants; which continuation information is preserved in passage; which differences must remain permanently visible; and how those relations maintain coherence across scales.

In other words, the future object to be formalized is not a new super-observer. It is how a finite heterogeneous relational network can form a traceable global organization.

\section{Problem A: A Finite Heterogeneous Singular Profile}

Let
\[
\Xi=\{\xi_1,\ldots,\xi_k\}\subseteq\partial_S^\infty
\]
denote the boundary/singular roots that must currently be retained. Since $\Xi\subseteq\partial_{S,\mathrm{sing}}^\infty\subseteq\partial_S^\infty$ and global boundary support is finite, $k<\infty$ is a known source-side condition, not a future conjecture. Each $\xi_i$ may carry its own local group of germs
\[
H_i=\operatorname{Germ}_{\xi_i}(\mathcal G).
\]

Orbit information must be separated first. If $\xi_i$ and $\xi_j$ lie in the same groupoid orbit and an arrow $\eta:\xi_i\to\xi_j$ is chosen, the corresponding isotropy groups---the local groups of germs---are related by conjugation and hence are isomorphic. Truly non-isomorphic local groups can occur between sectors in different orbits. A future singular profile may therefore need to record at least
\[
\boxed{\text{finite singular set}+\text{orbit partition}+\text{local germ data}.}
\]

But this remains only a local inventory. It does not tell us how different sectors are jointly organized through finite tiles, traverses, and changes of resolution.

\begin{warning}[Finite Traverse Does Not Mean Limiting Germ Arrow]
A finite-tile trajectory
\[
\xi_{i,n}\rightsquigarrow\xi_{j,n}
\]
shows only that two boundary occurrences are linked by a path in that finite approximation. It cannot automatically be upgraded to an arrow $\xi_i\to\xi_j$ in the limiting groupoid of germs. In particular, when two limiting sectors lie in different orbits, finite-level proximity or passage cannot be reverse-translated into an orbit relation that does not exist.
\end{warning}

\section{Problem B: How Regular Background and Exceptional Germ Resolutions Enter One Architecture}

A further fact is easily hidden by the language of ``multiple centers.'' Because global boundary support is finite, only finitely many infinite tiles carry boundary at all; among those, regular-boundary gluing must still be separated from singular-boundary germ resolution. For every boundary-free infinite tile and any root $x$ in it,
\[
\widetilde\Gamma_x\cong(T,x).
\]

Thus a future heterogeneous architecture does not split the entire infinite background into multiple sheets everywhere. Most of the boundary-free background is already self-resolved. The finitely supported boundary-bearing part then divides again into regular boundary gluing and singular germ-sheet resolution. The new work in heterogeneous germ integration is concentrated in the latter, but a global architecture cannot therefore forget the former.

This imposes a strong constraint on future work. An integration framework must simultaneously accommodate
\[
\boxed{
\begin{gathered}
\text{large boundary-free background}\\
+\ \text{finitely supported regular boundary gluing}\\
+\ \text{finitely supported singular germ resolutions},
\end{gathered}}
\]
rather than drawing only a small diagram of ``several singular nodes connected to one another'' and forgetting the infinite background in which they sit.

A graph of germs generated by boundary data can be assembled by gluing finitely or infinitely many infinite-tile copies. In the regular-boundary case there is only reversible labeled-arrow gluing; singular boundaries can additionally create nontrivial germ sheets. What is finite is total boundary support, and within it the singular-producing subset used for nontrivial resolution---not the number of copies inside a glued/resolved graph. Any future global object must continue to keep these layers on separate ledgers.

\section{Problem C: Decorated Transition Structure}

One natural candidate is to treat boundary sectors as vertices, traceable elementary finite passages as directed data, and local groups of germs, residue information, and scale embeddings as decorations. We do not name this object after any existing category or groupoid on purpose, because the required axioms and universal property have not been proved.

Future work must answer at least four questions:
\begin{itemize}
\item how finite-traverse composition is preserved;
\item how local germ/residue data are recorded along a passage without being swallowed by a coarse type label;
\item how conjugacy information for same-orbit sectors and\linebreak non-transportability across different orbits are preserved at once;
\item how higher-level decomposition is compatible with lower-level passage composition.
\end{itemize}

Until these questions are answered, calling different sectors ``several parts of one subject'' remains psychological narration rather than formal integration.

\section{Problem D: Integration without Homogenization}

The principle most worth preserving in this chapter is
\[
\boxed{\text{integration}\neq\text{homogenization}.}
\]

Integration should not require
\[
H_1\cong H_2\cong\cdots\cong H_k,
\]
nor should it require every sector to share one return grammar, one residue, or one continuation geometry. A whole may contain genuinely different local organizations. Integration asks for relations among those differences that are sufficiently stable, composable, and traceable; it does not first erase the differences and then announce that a whole has been obtained.

This carries forward a structural discipline left by the projection picture: a local projection may be real, yet no local projection thereby acquires global occupancy rights. Transferred to multi-boundary geometry, the point is not that ``every sector is a small subject,'' but rather
\[
\boxed{
\begin{gathered}
\text{no single sector may speak for the whole,}\\
\text{and the whole cannot exist only by flattening its sectors.}
\end{gathered}}
\]

From the psychological side, this supplies a future source of questions about heterogeneous organizing positions. Competence-related, relational, and public-evaluative positions, for example, may recruit different memories and have different continuation signatures; the worthwhile question is whether there is a stable passage structure among them. There is currently no evidence identifying these positions literally with the $\xi_i$, and still less evidence identifying their differences with non-isomorphic groups of germs.

\section{Problem E: From Bridge-Type Words to a Heterogeneous Multi-Sector Return Calculus}

The existing source prototype tells us that finite coding of a multi-boundary return should not simply be written as
\[
\mathcal R=\bigcup_i\mathcal R_i.
\]
A more faithful record is a sector/bridge-type word
\[
x_1x_2\cdots x_m,
\]
because one return factor may pass through several sectors and may depend on passage order \citep{kuang2024v1}.

Once local heterogeneity must also be preserved, the future object must contain at least something of the form
\[
\boxed{\begin{gathered}
\text{sector passage}+\text{return-interval incidence}\\[-1mm]
{}+\text{scale}+\text{local germ decoration}
\end{gathered}}.
\]

One would then have to study anew how nested or overlapping return intervals are compatible with sector passages, whether witness persistence admits a uniform formulation across different decorations, and whether global incompressibility can still be described by a controlled forbidden-factor language. No conclusion here can simply be inherited from the single-site benchmark.

\section{What Can Actually Be Asked on the Psychological Side: From ``Many Sides'' to Heterogeneous Passage}

The easiest psychological misreading of this future problem is the vague sentence ``people naturally have many sides.'' That is too weak, because
\[
\text{many projections}\not\Rightarrow\text{many boundaries},
\]
and still less does it imply heterogeneous germ structure.

A future empirical bridge would have to demand something more specific. Under a predeclared coding scheme, finitely many persistent organizing loci would need repeatedly distinguishable local continuation signatures; ordered passages among loci would need to predict later behavior beyond a single-root, unordered-state-set, or ordinary context-switching baseline; and as material expands, those loci would need mainly to be reidentified rather than proliferate without bound.

Only at that level does heterogeneous integration begin to become a testable psychological question rather than another way of saying that ``the subject is complex.''

\begin{stopbox}[Future boundary of this chapter]
The mathematics already available here consists of global finite boundary support, its finite singular-boundary subset, the multi-boundary/type-word return prototype, local groups of germs, finite traverses, and graph-of-germs resolution. What is not yet available is a single integration object that simultaneously preserves orbit partition, heterogeneous local germ data, passage composition, return incidence, and scale coherence; still less is there evidence identifying such a future object with a ``whole subject.''

The question left by this chapter is therefore not ``how can we add more boundaries?'' but a narrower one: when the genuine fissures are already finite, how can their differences jointly form structure without being erased? This has a potential interface with the rumination/checking and multi-sector passage cases in Part V, but at present only as a future research program. Complex psychological loops cannot be reverse-translated into literal singularities, and an elegant decorated object would not mean that psychological integration had been solved.
\end{stopbox}

\chapter{Future Problem II: Minimal States, Cofinal Equivalence, and Subject Invariants}
\label{chap:minimal-states-en}

\section{``Minimal State'' Is, First of All, Not ``Minimal Self''}

This book can already discuss rooted organizing positions, birooted episodes, return witnesses, continuation residue, and the embedding of finitely many persistent singular loci in a large regular background. It still has not constructed a mathematical object called the ``whole subject.'' That stop line must remain.

A future minimal state can therefore, at most, begin with a relative meaning: relative to a declared observation language and a family of continuation tests, it is a state that cannot be merged further without losing those observable distinctions. It is not the smallest constituent of a person, nor the true self hidden behind narrative.

The discipline can be compressed as
\[
\boxed{\text{minimal state relative to a probe language}\neq\text{minimal self}.}
\]

This distinction is consistent with the projection principle of the preceding chapters. If two positions look the same under the current projection, that means only that the present observational interface has not separated them; it does not prove that the underlying objects are identical. Conversely, different labels do not guarantee different structure, because labels may preserve narrative differences that are irrelevant to future organization.

A finiteness discipline must also be fixed in advance. The standing AF-by-discrete hypothesis is
\[
\left|\bigsqcup_{T\in\mathcal T_\infty}\partial_S T\right|<\infty,
\]
that is, global boundary support is finite relative to the fixed generating set and fixed presentation, and hence the singular-boundary subset is also finite. This does \emph{not} imply that the total number of candidate structural states is finite: boundary-free infinite tiles can still contain infinitely many ordinary rooted positions, episodes, and continuation contexts. Hence
\[
\boxed{\text{finite singular-boundary subset}\not\Rightarrow\text{finite subject-state space}.}
\]
Future minimalization must not confuse these two kinds of ``finite.''

\section{Problem E: Relative to What Should Continuation Equivalence Be Defined?}

The Myhill--Nerode-like intuition is attractive because it gives a very clean question: if two current positions are indistinguishable under every relevant future, is there still any need to keep them as two states? In the setting of this book, however, ``every relevant future'' has not yet been defined.

Future work might first fix a family $\mathcal C$ of admissible continuation probes. For a current rooted/sector position $p$, one could record the decorated future profile visible under those probes, schematically
\[
F_{\mathcal C}(p)=
\left(
\begin{array}{c}
\text{admissible continuations, return outcomes,}\\
\text{persistence data, residue/germ data when relevant}
\end{array}
\right).
\]
Here $F_{\mathcal C}$ is only future-work notation, not an object already defined by the present theory.

A possible research direction is to define some form of continuation equivalence only after one has proved that ``same future profile'' gives a stable equivalence relation compatible with composition, witness persistence, and decomposition coherence. At least six obstacles must be addressed:
\begin{itemize}
\item the raw AF core has no intrinsic labels, so literal word equality cannot substitute for structural equivalence;
\item ordinary rooted positions and singular boundary loci must be kept separate; not every state may be assumed to be a boundary state;
\item if a position lies in a singular sector, orbit partition and local germ data cannot silently disappear under the quotient;
\item elementary traverses, composite paths, and a complete continuation language belong to different levels;
\item the equivalence must be compatible with cross-scale witness embedding and decomposition coherence;
\item whether the quotient is a congruence and whether it preserves the required return/residue structure must both be proved.
\end{itemize}

The family of probes must also be stated. If $\mathcal C_1$ tests only coarse futures while $\mathcal C_2$ adds more history- or germ-sensitive tests, two positions indistinguishable under $\mathcal C_1$ may separate again under $\mathcal C_2$. Thus minimality is first a resolution-relative / question-relative candidate concept, not the ontological ``smallest real state.''

\section{Problem F: A Good Quotient and Narrative Overcompression}

If continuation equivalence can eventually be established rigorously, it may give the narrative-compression discussion of Chapter 34 a more intrinsic structural standard. Yet two very different senses of ``same'' must still be separated.

One is observational equivalence under the current probe language: no difference affecting continuation has yet been found. The other is the stronger possibility of structural quotientability: after proof, the relevant differences no longer affect composition, return, persistence, germ decoration, or scale refinement.

Therefore one cannot write
\[
\text{same measured future}\Longrightarrow\text{same state}.
\]
A safer order of research is
\[
\begin{gathered}
\text{same measured future}
\Longrightarrow \text{candidate observational quotient}\\
\Longrightarrow \text{richer probes/coherence tests}
\Longrightarrow \text{possible structural quotient}.
\end{gathered}
\]

This provides one future direction for testing narrative overcompression. If ``I am always like this'' merges two positions that continue to produce different continuation grammars under predeclared future tests, the compression may have deleted a genuinely relevant difference. If two labels remain different for a long time but cannot be separated under progressively richer probes, one may instead ask whether a compression has been missed.

But this language must never become a psychological tribunal. Two experiences being observationally equivalent in the current experiment does not authorize the researcher to announce that they ``are really the same experience,'' still less that they belong to ``the same true self-state.''

\section{Problem G: Cofinal Atlas Equivalence Must First Be Separated from Standard Cofinality}

A Bratteli presentation can be telescoped, so literal birth level, one particular intermediate level, or one finite spelling should not casually be treated as intrinsic data. When comparing two candidate architectures in the future, a natural idea is to allow finite levels to be skipped, or to compare multi-scale decomposition, return persistence, and singular decoration along cofinal subsequences.

This manuscript provisionally calls that idea \emph{cofinal atlas equivalence}, but the phrase is only a working label. It must not be silently identified with the established tail/cofinal equivalence on Bratteli path spaces, nor may it inherit theorems from that notion merely because both contain the word ``cofinal.'' Future work must specify:
\begin{itemize}
\item exactly what objects are being compared;
\item which telescoping/refinement moves are allowed;
\item which witness, return, germ, and orbit data must be transported under those moves;
\item in what sense two presentations carry ``the same structural information.''
\end{itemize}

The generator dependence frozen earlier in this book is especially important. The global finite-boundary-support condition is relative to a fixed finite generating set $S$ and its associated presentation. In the absence of an invariance theorem, raw boundary-point counts, concrete boundary locations, or birth levels cannot be declared subject invariants across presentations. A future equivalence that forgets presentation must first prove which of these data transport naturally and which should be forgotten.

\section{Problem H: Morita-Like Equivalence Cannot Be a Shortcut to ``Subject Identity''}

Morita equivalence is an important language for comparing groupoids in étale groupoid theory \citep{renault1980,nek2022}. It suggests a future question: might two different presentations carry the same global structural information in some categorical or dynamical sense?

At least three gates have not yet been passed. First, the cognitive-side architecture has not been constructed as a natural étale groupoid or another category capable of carrying Morita theory. Second, even if such an object is eventually constructed, Morita equivalence would be an equivalence notion between the mathematical objects so constructed; it would not automatically mean ``the same subject.'' Third, orbit equivalence, tail/cofinal equivalence, continuation equivalence, and Morita equivalence answer different questions. This manuscript does not arrange them into an unproved hierarchy of strength.

The only licensed order at present is therefore
\begin{center}
\resizebox{0.96\linewidth}{!}{$\displaystyle
\boxed{
\text{construct the objects}
\longrightarrow
\text{define an equivalence}
\longrightarrow
\text{ask whether it is suitable for subject comparison}.}
$}
\end{center}
One must not first decide what ``the same person'' means and then search mathematics for an equivalence relation whose name sounds strong enough.

\section{Problem I: Equivalence Must Come Before Subject Invariant}

The word \emph{invariant} is especially prone to acquiring too much authority. Mathematically, a quantity or structure is invariant only after the transformations/equivalences under which it remains unchanged have been specified. Until Problems E--H are solved, a strict ``subject invariant'' has therefore not been defined.

Candidate data worth investigating in the future include
\[
\begin{gathered}
\text{transition/passage architecture},\\
\text{persistent return structure and filtration},\\
\text{orbit-partitioned local germ profile},\\
\text{residue/continuation data},\\
\text{decomposition grammar across scales}.
\end{gathered}
\]
But these are presently only candidate descriptors. Some may be presentation-dependent; some may make sense only for a fixed generating set; some may become genuine invariants only after an appropriate quotient. In particular, global finite boundary support is a standing source-domain hypothesis in this book. Its importance does not license the raw boundary/singular count to become a cross-presentation ``subject constant.''

The psychological side needs another stop line:
\[
\boxed{\text{structural invariant}\neq\text{personality trait}\neq\text{identity essence}.}
\]
Even if a future transition grammar remains stable across tasks and time points, it is first data invariant under a previously defined structural equivalence. Interpreting it as a trait, developmental feature, or personality dimension would require another psychological theory and longitudinal evidence.

\section{A Future-Work Order That Must Not Be Reversed}

Taken together, the chapter gives a concrete research order:
\begin{enumerate}
\item decide the candidate state object and the permitted continuation language;
\item propose continuation equivalence and prove that it is a composable, cross-scale-compatible equivalence/congruence;
\item study whether the quotient preserves return, germ, orbit, and residue data without overcompression;
\item define presentation change / cofinal atlas equivalence and distinguish intrinsic data from presentation artifacts;
\item if a natural groupoid/category genuinely appears, test whether Morita-like equivalence is appropriate;
\item only after an equivalence notion has been frozen is it legitimate to define invariants relative to it;
\item only then return to psychology and ask whether those mathematical states, quotients, or invariants have independently measurable subject-level consequences.
\end{enumerate}

The point is simple: because ``the same subject'' is a question we want to answer, we cannot pre-emptively name a mathematical equivalence ``subject identity.'' A future theory must first prove what it preserves, then discuss what such preservation means for psychology.

\begin{stopbox}[Future boundary and stop line]
Chapter 47 asked how finitely many heterogeneous singular sectors can jointly form structure without being homogenized. We now ask, once a structure is given, which differences may safely be forgotten and what equivalence could allow two presentations or states to count as ``structurally the same.'' These remain open questions. Problems E--I are not part of the theory completed here. One may not infer a finite subject-state space from a finite singular-boundary subset; observational/cofinal/Morita-like equivalence may not be written as subject identity; and candidate descriptors may not be written as personality dimensions. Without new definitions, theorems, and empirical bridges, this manuscript stops here.
\end{stopbox}

\chapter{Whole-Book Concept Passport and Claim Audit}
\label{chap:whole-book-passport-en}

This chapter retains only the final authority table; detailed definitions remain where each concept first appears.

\begingroup\small
\begin{longtable}{>{\raggedright\arraybackslash\bfseries}p{0.18\textwidth}>{\raggedright\arraybackslash}p{0.32\textwidth}>{\raggedright\arraybackslash}p{0.30\textwidth}}
\toprule
Object / claim & Current authority & Must not be upgraded into\\
\midrule
AF-by-discrete & open AF core + discrete complement; under the current compactly generated fixed-generator tile correspondence, global finite persistent boundary support & the empirical fact that ``psychological fissures are naturally finite''\\
boundary point & persistent generator interface; may be regular or singular & a mathematical name for trauma point, event boundary, or affective peak\\
singular germ & $H_\xi\neq 1$; nontrivial local germ fibre & ordinary reflection, contradiction, or role switching\\
Hausdorff / non-Hausdorff / pNH & hierarchy of local separability of germ differences & second-/third-order personality levels\\
organizing position & psychological coordinate currently comparing, evaluating, and recruiting history & the whole ego or phenomenological mineness\\
recognized return & a return-completion judgment beyond similarity, requiring independent evidence & all repetition, relapse, or rumination\\
continuation profile $C^+$ & prediction, action, and memory recruitment that can still differ after return-completion & formal residue itself\\
$H_{\mathrm{fin}}^{\mathrm{psych}}$ & finite-set-like stabilization hypothesis for singular-like loci & empirical verification of the full source-side $H_{\mathrm{fin}}^{\mathrm{AFBD}}$\\
\bottomrule
\end{longtable}
\endgroup

\section{Conceptual Genealogy Also Belongs in the Passport}

Future research must record the source ancestry of a concept as part of its passport. This matters especially for P5/P6$^*$: Third-Order Consciousness supplies an important natural-language organization, but observer endogeneity and the anti-final-observer problem are not its exclusive property; second-order cybernetics and related traditions are explicit competing sources \citep{vonfoerster2003}. The theoretical significance of non-Hausdorff / pNH should therefore be written as ``mathematics adds new resolution to a shared reflexivity problem,'' not as a historical priority claim.

\section{Permanent Separation of Ledgers}

The following relations must never be merged again:
\begin{center}
\resizebox{0.96\linewidth}{!}{$\displaystyle
\begin{gathered}
\text{boundary}\not\Rightarrow\text{singular},\qquad
\text{singular}\not\Rightarrow\text{non-Hausdorff},\qquad
\text{non-Hausdorff}\not\Rightarrow\text{pNH},\\
\text{repetition}\not\Rightarrow\text{return},\qquad
\text{return}\not\Rightarrow\text{same continuation},\qquad
\text{root}\not\Rightarrow\text{whole subject}.
\end{gathered}
$}
\end{center}
Likewise, finitely many types do not imply finitely many loci, and finitely many singular-like loci do not imply a finite subject-state space.

\section{Claim Strength}

Level II permits only a controlled compatibility between relation skeletons independently defined on the two sides. Level III requires independent operationalization, a baseline, and held-out prediction. Level IV requires manipulation, mediation, and control of alternative mechanisms. Level V---``consciousness is essentially AF-by-discrete''---does not belong to this book. When the verb changes from ``is compatible with'' to ``predicts,'' ``causes,'' ``implements,'' or ``is,'' the evidential burden changes with it.

\chapter{Conclusion: A Theory That Knows It Has Not Completed Consciousness}
\label{chap:conclusion-en}

This book began with several ordinary questions. When does experience become organized as ``mine''? Why does similarity become ``again''? Who is cutting, comparing, and judging? Once that observational position enters the picture, can it simply regain supreme authority in a new vocabulary?

\section{What Remains Is Not an Equality Sign}

What remains at the end is not
\[
\text{Third-Order Consciousness}=\text{AF-by-discrete}.
\]
The two sides remain independent native domains throughout. Third-Order Consciousness supplies a family of candidate questions about subject formation, narrative, the observer, and final authority; metacognition, decentering, self-model theory, predictive processing, Buddhist no-self traditions, second-order cybernetics, and other traditions already explain or anticipate a substantial part of that terrain. Observer endogeneity in particular is not an exclusively new discovery of Third-Order Consciousness \citep{vonfoerster2003}. Critical reconstruction permits such reduction and requires it to change downstream attribution.

Mathematics contributes at another point. General germ geometry separates ``arriving at the same point'' from ``which local action brings one there.'' AF-by-discrete adds an extra, falsifiable commitment of finitely supported fissures. The distinctions regular/singular, Hausdorff/non-Hausdorff/pNH, and return-completion/residue then force previously broad natural-language questions to split further.

We therefore leave a family of relations that can no longer be merged casually:
\[
\begin{gathered}
\text{repetition}\not\Rightarrow\text{return},\\
\text{return-completion}\not\Rightarrow\text{same continuation},\\
\text{root}\not\Rightarrow\text{whole system},\\
\text{boundary}\not\Rightarrow\text{singular},\\
\text{observer enters the system}\not\Rightarrow\text{observer has lost final authority}.
\end{gathered}
\]

The last line tightens the reflexivity problem one step further. Placing the observer inside the system resolves only the first problem. The candidate mathematical gain from non-Hausdorffness is to distinguish between ``an organizing difference has become visible'' and ``that difference has reacquired a final territory isolated from other local possibilities.'' This remains only a Level-II benchmark. It becomes genuine cross-domain work only if it generates operational questions in return and exceeds constructs already available in psychology.

\section{Finite Boundaries Organize an Infinite Background}

On the source side, the master picture retained by the book is still
\[
\boxed{\text{infinite AF-like background}+\text{globally finite persistent boundary support}.}
\]
This boundary support contains both regular and singular points. Regular boundary supports ordinary reversible gluing; singular boundary may additionally produce a nontrivial germ fibre. The psychological side currently operationalizes only singular-like interfaces, so $H_{\mathrm{fin}}^{\mathrm{psych}}$ and the full $H_{\mathrm{fin}}^{\mathrm{AFBD}}$ must remain permanently separate.

This finiteness is not a necessary truth derived from Third-Order Consciousness. It is a strong hypothesis that this project has chosen to bear. If open-world material continues to force new, nonmergeable persistent singular-like loci, the finite singular-interface translation should yield. If future work instead requires polynomial-activity-like proliferation or a new heterogeneous architecture, it should say plainly that the model has changed.

\section{Third Order Is Most Dangerous When It Successfully Becomes an Identity}

A person can learn that ``the observer is also inside the structure'' and immediately use the sentence to manage the self: ``I must observe correctly. I must let go quickly. I must prove that this method works.'' Reflexive vocabulary has appeared, yet a new manager may already have reacquired exemption.

The practical sentence finally retained by this book is therefore not ``have no goal,'' but:
\begin{quote}
\textbf{I may hope that it ends, but I do not require observation to be responsible for making it end.}
\end{quote}
Action, responsibility, boundaries, and value remain. What disappears is the right of any observer, theory, or goal to acquire final explanatory authority merely because it is currently organizing experience.

\section{What Has Not Been Completed}

This book has not explained phenomenological mineness, qualia, full embodiment, affective valence, neural realization, social and developmental mechanisms; nor has it defined a complete mathematical subject. A finite singular-boundary subset still does not imply a finite subject-state space. The blanks left by Part VI are not holes waiting for the next term to fill them. They are where the current evidence genuinely stops.

If the work continues, it should first test the distinctions already proposed. Can they be independently operationalized in ordinary psychological language? Do they outperform simple baselines on held-out material? In Study H2, does finite saturation outperform continuing growth? Does meta-identification relativizability add anything beyond decentering, cognitive flexibility, and context sensitivity? If the answer is no, the theory should become smaller.

\section{Keep the Diagram Open}

Mathematics has not spoken the last word for consciousness. It has done something more limited and more specific: it has separated several relations that are easy to complete too early, making definitions, evidence, and failure conditions traceable again.

The final discipline is therefore still this:
\begin{quote}
\textbf{Make definitions strict enough that error has nowhere to hide;}\\
\textbf{make translation restrained enough that both native domains can still breathe;}\\
\textbf{make observation open enough that the observer, too, cannot escape outside the diagram.}
\end{quote}

Formalize without overreach; explain without claiming the final ceiling. See the structure, and leave reality somewhere to continue happening.

\chapter*{Primary Source Index: Full WeChat Permanent Links}
\addcontentsline{toc}{chapter}{Primary Source Index: Full WeChat Permanent Links}
This index preserves the stable source IDs TOC-001--TOC-083, archive page spans, dates, and original Chinese titles. The 83 source records map one-to-one to 83 verified permanent WeChat links. To keep paper copies, screenshots, and plain-text exports independently traceable, every full URL is printed here and remains clickable.

The opening/general preface of \CN{《三阶意识：从生物我、社会我到空位意识》}, which was not separately included in the 619-page compilation, is available at:
\begin{quote}\url{https://mp.weixin.qq.com/s/SDsz7NttA9AX9nk53tZPhw}\end{quote}
This index publishes metadata and permanent links only; it does not reproduce the full text of the WeChat posts.

\begingroup\scriptsize
\begin{longtable}{>{\raggedright\arraybackslash\bfseries}p{0.10\textwidth}>{\raggedright\arraybackslash}p{0.14\textwidth}>{\raggedright\arraybackslash}p{0.56\textwidth}}
\toprule
ID & Date / archive pages & Archived title / full permanent WeChat link\\
\midrule
\endfirsthead
\toprule
ID & Date / archive pages & Archived title / full permanent WeChat link\\
\midrule
\endhead
TOC-001 & 2026-03-15 / 1–6 & \CN{1.1 “我在思考”这句话为什么可疑？ （第一章：意识不是我）——《三阶意识：从 生物我、社会我到空位意识》}\newline \url{https://mp.weixin.qq.com/s/U-q3jrAaWoNG9vnzGISwAQ}\\[1.2mm]
TOC-002 & 2026-03-15 / 7–12 & \CN{1.2 不是先有我，再有意识（第一章：意识不是我）——《三阶意识：从生物我、 社会我到空位意识》}\newline \url{https://mp.weixin.qq.com/s/bO-6MbJAyZOU02sRfP4riQ}\\[1.2mm]
TOC-003 & 2026-03-16 / 13–19 & \CN{1.3 “我”作为界面结果，而非本体起点（第一章：意识不是我）——《三阶意识： 从生物我、社会我到空位意识》}\newline \url{https://mp.weixin.qq.com/s/6K8mZGAizHjVtGIZQFLXNw}\\[1.2mm]
TOC-004 & 2026-03-17 / 20–25 & \CN{1.4 自我是意识的组织效应（第一章：意识不是我）——《三阶意识：从生物我、 社会我到空位意识》}\newline \url{https://mp.weixin.qq.com/s/ToJGG1ff3dWUqOHrcBdHuQ}\\[1.2mm]
TOC-005 & 2026-03-17 / 26–30 & \CN{1.5 本章总结：自我是如何从界面中升起的（第一章：意识不是我）——《三阶意 识：从生物我、社会我到空位意识》}\newline \url{https://mp.weixin.qq.com/s/G-Th6dNuGPwrfDlVXuwozA}\\[1.2mm]
TOC-006 & 2026-03-17 / 31–37 & \CN{2.1 物，不如说物性（第二章：意识作为耦合显现）——《三阶意识：从生物我、 社会我到空位意识》}\newline \url{https://mp.weixin.qq.com/s/Eafhi8ye3Efm6dzVeybJrg}\\[1.2mm]
TOC-007 & 2026-03-18 / 38–43 & \CN{2.2 基底物性、组织物性与显现物性（第二章：意识作为耦合显现）——《三阶意 识：从生物我、社会我到空位意识》}\newline \url{https://mp.weixin.qq.com/s/5ZvhX2IgrkMt12M25f7j8g}\\[1.2mm]
TOC-008 & 2026-03-18 / 44–49 & \CN{2.3 为什么意识既非纯灵，也非副产品（第二章：意识作为耦合显现）——《三阶 意识：从生物我、社会我到空位意识》}\newline \url{https://mp.weixin.qq.com/s/eu4cZXtsUmlIEempNXrHXg}\\[1.2mm]
TOC-009 & 2026-03-19 / 50–54 & \CN{2.4 地球人的意识：一种被条件化的界面显现（第二章：意识作为耦合显现）—— 《三阶意识：从生物我、社会我到空位意识》}\newline \url{https://mp.weixin.qq.com/s/0vwZK0BCVog-F0hTfDd1Sg}\\[1.2mm]
TOC-010 & 2026-03-19 / 55–59 & \CN{2.5 意识的存在论定位：显现，不是实体（第二章：意识作为耦合显现）——《三 阶意识：从生物我、社会我到空位意识》}\newline \url{https://mp.weixin.qq.com/s/tlJ9Qi_ADunVDFt246xjZA}\\[1.2mm]
TOC-011 & 2026-03-20 / 60–65 & \CN{3.1 感官不是窗口，而是滤网（第三章：地球人的意识为何天然带着偏差）—— 《三阶意识：从生物我、社会我到空位意识》}\newline \url{https://mp.weixin.qq.com/s/qCGiKq0OlEN5xPhaEUBGlA}\\[1.2mm]
TOC-012 & 2026-03-20 / 66–70 & \CN{3.2 神经系统不是记录器，而是压缩器（第三章：地球人的意识为何天然带着偏 差）——《三阶意识：从生物我、社会我到空位意识》}\newline \url{https://mp.weixin.qq.com/s/RnRNti6MDYNJbb5PykPwCg}\\[1.2mm]
TOC-013 & 2026-03-20 / 71–75 & \CN{3.3 时间感、空间感与对象感如何被生成（第三章：地球人的意识为何天然带着 偏差）——《三阶意识：从生物我、社会我到空位意识》}\newline \url{https://mp.weixin.qq.com/s/XLDWSExXQJZBjJ1dO4-QgQ}\\[1.2mm]
TOC-014 & 2026-03-21 / 76–81 & \CN{3.4 人类意识的五重界面偏差（第三章：地球人的意识为何天然带着偏差）—— 《三阶意识：从生物我、社会我到空位意识》}\newline \url{https://mp.weixin.qq.com/s/MyGTfpEZjJOBLJGBetf_Nw}\\[1.2mm]
TOC-015 & 2026-03-21 / 82–89 & \CN{3.5 所谓现实感，不过是高稳定性的读取结果（第三章：地球人的意识为何天然 带着偏差）——《三阶意识：从生物我、社会我到空位意识》}\newline \url{https://mp.weixin.qq.com/s/U9gyqPw5Fo-OX4gt0XwniQ}\\[1.2mm]
TOC-016 & 2026-03-22 / 90–94 & \CN{4.1 内感觉：意识最早不是看世界，而是感到自己快要失衡（第四章身体是第一 重界面）——《三阶意识：从生物我、社会我到空位意识》第二部第一阶意识：生 物我}\newline \url{https://mp.weixin.qq.com/s/Vx4CpQ7PgW4B0oudSTz_Hg}\\[1.2mm]
TOC-017 & 2026-03-23 / 95–100 & \CN{4.2 饥饿、疼痛、恐惧与快感：意识的底层语法（第四章身体是第一重界面）—— 《三阶意识：从生物我、社会我到空位意识》第二部第一阶意识：生物我}\newline \url{https://mp.weixin.qq.com/s/s00aKHRhY2MnDbCLynPlDg}\\[1.2mm]
TOC-018 & 2026-03-23 / 101–105 & \CN{4.3 身体边界如何组织出最初的“我在这里”（第四章身体是第一重界面）—— 《三阶意识：从生物我、社会我到空位意识》第二部第一阶意识：生物我}\newline \url{https://mp.weixin.qq.com/s/eRJ3Gy63qcJhr5cITe-QnQ}\\[1.2mm]
TOC-019 & 2026-03-24 / 106–111 & \CN{4.4 第一阶意识不是主体，而是生存聚焦（第四章身体是第一重界面）——《三阶 意识：从生物我、社会我到空位意识》第二部第一阶意识：生物我}\newline \url{https://mp.weixin.qq.com/s/I-Ymxth7I0QoRBeyc8y5tw}\\[1.2mm]
TOC-020 & 2026-03-24 / 112–116 & \CN{4.5 身体不是客体，而是最早的在场格式（第四章身体是第一重界面）——《三阶 意识：从生物我、社会我到空位意识》第二部第一阶意识：生物我}\newline \url{https://mp.weixin.qq.com/s/ulvMWVkn4MoiLhg07YLOYQ}\\[1.2mm]
TOC-021 & 2026-03-25 / 117–122 & \CN{5.1 中心不等于实体（第五章生物我：一个被迫出现的中心）——《三阶意识：从 生物我、社会我到空位意识》第二部第一阶意识：生物我}\newline \url{https://mp.weixin.qq.com/s/-HdnXTz3tL19UFyeDdOpEA}\\[1.2mm]
TOC-022 & 2026-03-25 / 123–129 & \CN{5.2 身体我——意识在场的第一个拓扑结（第五章生物我：一个被迫出现的中心） ——《三阶意识：从生物我、社会我到空位意识》第二部第一阶意识：生物我}\newline \url{https://mp.weixin.qq.com/s/jcon9YSb5opgyCuhthZZ_A}\\[1.2mm]
TOC-023 & 2026-03-26 / 130–136 & \CN{5.3 前反思路径性：行动如何先于自我解释（第五章生物我：一个被迫出现的中 心）——《三阶意识：从生物我、社会我到空位意识》第二部第一阶意识：生物我}\newline \url{https://mp.weixin.qq.com/s/uxduAEgSyNa924kU0Ugbdw}\\[1.2mm]
TOC-024 & 2026-03-26 / 137–143 & \CN{5.4 动物式流动与人类式占有的差别（第五章生物我：一个被迫出现的中心）—— 《三阶意识：从生物我、社会我到空位意识》第二部第一阶意识：生物我}\newline \url{https://mp.weixin.qq.com/s/hjE1EL4tpzjs6mFpUykiNQ}\\[1.2mm]
TOC-025 & 2026-03-27 / 144–149 & \CN{5.5 生物我：为什么这一刀必须切在这里（第五章生物我：一个被迫出现的中心） ——《三阶意识：从生物我、社会我到空位意识》第二部第一阶意识：生物我}\newline \url{https://mp.weixin.qq.com/s/nUshMaTj2roylhJo2JH8cw}\\[1.2mm]
TOC-026 & 2026-03-27 / 150–156 & \CN{6.1 第一阶时间不是历史，而是节律（第六章：时间如何展开意识）——《三阶意 识：从生物我、社会我到空位意识》}\newline \url{https://mp.weixin.qq.com/s/kF4Zo5tC_eKneJaoxvU8RQ}\\[1.2mm]
TOC-027 & 2026-03-28 / 157–162 & \CN{6.2 为什么第一阶意识只能活在“下一秒”（第六章：时间如何把意识推向升阶） ——《三阶意识：从生物我、社会我到空位意识》}\newline \url{https://mp.weixin.qq.com/s/Elk1MRSUKoUfMV7FW9ofZg}\\[1.2mm]
TOC-028 & 2026-03-28 / 163–168 & \CN{6.3 投影为何必须沿时间轴展开（第六章：时间如何展开意识）——《三阶意识： 从生物我、社会我到空位意识》}\newline \url{https://mp.weixin.qq.com/s/4i7QfhNDA_1c79OWWdoe3Q}\\[1.2mm]
TOC-029 & 2026-03-28 / 169–176 & \CN{6.4 实体化的必要幻觉：为何投影必须“缠绕”出一个主体（第六章：时间如何展 开意识）——《三阶意识：从生物我、社会我到空位意识》}\newline \url{https://mp.weixin.qq.com/s/KvbSRC7LwQuY7JD616ZfQA}\\[1.2mm]
TOC-030 & 2026-03-28 / 177–182 & \CN{7.1 进化解释的边界：为什么智人大脑显得“过量” （第七章为显现而准备的大脑） ——《三阶意识：从生物我、社会我到空位意识》第三部第二阶意识：社会我}\newline \url{https://mp.weixin.qq.com/s/FC3v6I9XrBPmlUfChk-tHw}\\[1.2mm]
TOC-031 & 2026-03-29 / 183–188 & \CN{7.2 人类基底物性的异常：一台没有被写死的脑（第七章为显现而准备的大脑） ——《三阶意识：从生物我、社会我到空位意识》第三部第二阶意识：社会我}\newline \url{https://mp.weixin.qq.com/s/h1RGPt-ERwmzW3jWkrJh4w}\\[1.2mm]
TOC-032 & 2026-03-29 / 189–195 & \CN{7.3 组织物性的跃迁：从响应世界到内部重演世界（第七章为显现而准备的大脑） ——《三阶意识：从生物我、社会我到空位意识》第三部第二阶意识：社会我}\newline \url{https://mp.weixin.qq.com/s/OT5SAHrpLKijDFTYz4TuwA}\\[1.2mm]
TOC-033 & 2026-03-30 / 196–202 & \CN{7.4 他者如何进入系统内部（第七章为显现而准备的大脑）——《三阶意识：从生 物我、社会我到空位意识》第三部第二阶意识：社会我}\newline \url{https://mp.weixin.qq.com/s/5Y10d46a511aL3ftfnMEog}\\[1.2mm]
TOC-034 & 2026-03-30 / 203–211 & \CN{7.5 人类大脑过量不是为了活下去，而是为了让显现升级（第七章为显现而准备 的大脑）——《三阶意识：从生物我、社会我到空位意识》第三部第二阶意识：社 会我}\newline \url{https://mp.weixin.qq.com/s/jDOyjg2iwYEvVzE3X7Us3g}\\[1.2mm]
TOC-035 & 2026-03-30 / 212–217 & \CN{8.1 为什么符号会被召唤（第八章语言不是表达意识，而是设定意识）——《三阶 意识：从生物我、社会我到空位意识》第三部第二阶意识：社会我}\newline \url{https://mp.weixin.qq.com/s/7WSn0Dv0WFzniYqeh95Ibw}\\[1.2mm]
TOC-036 & 2026-03-31 / 218–224 & \CN{8.2 命名如何切割经验（第八章语言不是表达意识，而是设定意识）——《三阶意 识：从生物我、社会我到空位意识》第三部第二阶意识：社会我}\newline \url{https://mp.weixin.qq.com/s/0Cy0ydwrEOFjnpkIo52IrA}\\[1.2mm]
TOC-037 & 2026-03-31 / 225–230 & \CN{8.3 主谓宾结构如何制造主体幻觉（第八章语言不是表达意识，而是设定意识） ——《三阶意识：从生物我、社会我到空位意识》第三部第二阶意识：社会我}\newline \url{https://mp.weixin.qq.com/s/wxqkZafjiL8ucvObbAQepA}\\[1.2mm]
TOC-038 & 2026-03-31 / 231–238 & \CN{8.4 语言不是中性工具，而是组织机制（第八章语言不是表达意识，而是设定意 识）——《三阶意识：从生物我、社会我到空位意识》第三部第二阶意识：社会我}\newline \url{https://mp.weixin.qq.com/s/wN9kiiopIlG5SiTxNxBrkQ}\\[1.2mm]
TOC-039 & 2026-04-01 / 239–246 & \CN{8.5 概念如何替代直接经验（第八章语言不是表达意识，而是设定意识）——《三 阶意识：从生物我、社会我到空位意识》第三部第二阶意识：社会我}\newline \url{https://mp.weixin.qq.com/s/p9aBrc3B8MH9eoXScv5Cgw}\\[1.2mm]
TOC-040 & 2026-04-01 / 247–255 & \CN{8.6 当语言接管时间与身体：第二阶意识的正式降临（第八章语言不是表达意识， 而是设定意识）第三部第二阶意识：社会我}\newline \url{https://mp.weixin.qq.com/s/8XhxWIB-AAOEAip6A4uxFQ}\\[1.2mm]
TOC-041 & 2026-04-02 / 256–264 & \CN{9.1 名字、角色与归属：社会缝合的第一层针脚（第九章社会我：身份如何缝合身 体我）——《三阶意识：从生物我、社会我到空位意识》第三部第二阶意识：社会我}\newline \url{https://mp.weixin.qq.com/s/CR_lC-2cDz-wJuJ9ag6qdQ}\\[1.2mm]
TOC-042 & 2026-04-02 / 265–272 & \CN{9.2 镜像机制：我如何在他人眼中成立（第九章社会我：身份如何缝合身体我） ——《三阶意识：从生物我、社会我到空位意识》第三部第二阶意识：社会我}\newline \url{https://mp.weixin.qq.com/s/qloCV04ARkPf5n8-p6r7-w}\\[1.2mm]
TOC-043 & 2026-04-02 / 273–281 & \CN{9.3 客我、主我与叙事性自我：分裂的内部剧场（第九章社会我：身份如何缝合身 体我）——《三阶意识：从生物我、社会我到空位意识》第三部第二阶意识：社会我}\newline \url{https://mp.weixin.qq.com/s/rIFzw_qNZ6NanCF8nIZmMA}\\[1.2mm]
TOC-044 & 2026-04-03 / 282–290 & \CN{9.4 社会不是约束一个现成的我，而是“生产”了我（第九章社会我：身份如何缝 合身体我）——《三阶意识：从生物我、社会我到空位意识》第三部第二阶意识： 社会我}\newline \url{https://mp.weixin.qq.com/s/V10D9z-QWyVpDJPi6Gphiw}\\[1.2mm]
TOC-045 & 2026-04-04 / 291–299 & \CN{9.5 身份感为什么如此坚硬：比肉身更沉重的幻觉（第九章社会我：身份如何缝 合身体我）——《三阶意识：从生物我、社会我到空位意识》第三部第二阶意识： 社会我}\newline \url{https://mp.weixin.qq.com/s/pleGRZLpF6DEoluJXzBXFw}\\[1.2mm]
TOC-046 & 2026-04-04 / 300–308 & \CN{10.1 从匮乏到欲望的符号化（第十章欲望剧场：第二阶意识的真正燃料）—— 《三阶意识：从生物我、社会我到空位意识》第三部第二阶意识：社会我}\newline \url{https://mp.weixin.qq.com/s/sLuVIFQIIocfd63eTVRtMg}\\[1.2mm]
TOC-047 & 2026-04-05 / 309–314 & \CN{10.2 叔本华：欲望为何永不满足（第十章欲望剧场：第二阶意识的真正燃料） ——《三阶意识：从生物我、社会我到空位意识》第三部第二阶意识：社会我}\newline \url{https://mp.weixin.qq.com/s/3juvujoPsqJlqFmuDYxb9Q}\\[1.2mm]
TOC-048 & 2026-04-05 / 315–323 & \CN{10.3 尼采：欲望如何被强行“等级化”（第十章欲望剧场：第二阶意识的真正燃 料）——《三阶意识：从生物我、社会我到空位意识》第三部第二阶意识：社会我}\newline \url{https://mp.weixin.qq.com/s/3oK7Tp4Xj5NA2AOZ0w9GqA}\\[1.2mm]
TOC-049 & 2026-04-06 / 324–333 & \CN{10.4 拉康：我们欲望的不是对象，而是缺口（第十章欲望剧场：第二阶意识的真 正燃料）——《三阶意识：从生物我、社会我到空位意识》第三部第二阶意识：社 会我}\newline \url{https://mp.weixin.qq.com/s/9LEjqoAF_bQRotPKuCawQQ}\\[1.2mm]
TOC-050 & 2026-04-07 / 334–341 & \CN{10.5 福柯与吉拉尔：权力如何设定欲望，模仿如何加速竞争（第十章欲望剧场： 第二阶意识的真正燃料）第三部第二阶意识：社会我}\newline \url{https://mp.weixin.qq.com/s/RxQ4h4SHeAzQJ5iHnjhdRA}\\[1.2mm]
TOC-051 & 2026-04-08 / 342–350 & \CN{10.6 欲望不是自由的发动机，而是二阶剧场的火源（第八章语言不是表达意识， 而是设定意识）第三部第二阶意识：社会我}\newline \url{https://mp.weixin.qq.com/s/zXujp2eruQxucrpVgmi76w}\\[1.2mm]
TOC-052 & 2026-04-08 / 351–359 & \CN{11.1 主体为何要通过编织记忆争夺“叙事权” （第十一章叙事路径我：我们如何 活成一个故事）——《三阶意识》第三部第二阶意识：社会我}\newline \url{https://mp.weixin.qq.com/s/-pzX3DP4D9epZ1gk91PFVQ}\\[1.2mm]
TOC-053 & 2026-04-09 / 360–368 & \CN{11.2 创伤为何容易变成人生的解释中轴（第十一章叙事路径我：我们如何活成一 个故事）——《三阶意识：从生物我、社会我到空位意识》第三部第二阶意识：社 会我}\newline \url{https://mp.weixin.qq.com/s/QZPq6H1qz2hieG0MJqhJCQ}\\[1.2mm]
TOC-054 & 2026-04-09 / 369–376 & \CN{11.3 时间如何被叙事压成命运：从流动到解释权（第十一章叙事路径我：我们如 何活成一个故事）——《三阶意识》第三部第二阶意识：社会我}\newline \url{https://mp.weixin.qq.com/s/6Cpn3pMcNCnGHMkq-uSy2w}\\[1.2mm]
TOC-055 & 2026-04-10 / 377–383 & \CN{11.4 “主体工程”的最终形态：从生存的连续，到故事的连续（第十一章叙事路 径我：我们如何活成一个故事）——《三阶意识》第三部第二阶意识：社会我}\newline \url{https://mp.weixin.qq.com/s/XhhKJnp6sgeUABXCBtU3Ig}\\[1.2mm]
TOC-056 & 2026-04-11 / 384–391 & \CN{11.5 第二阶意识的终点： “本体病”末期与“不朽”的诱惑（第十一章叙事路径 我：我们如何活成一个故事）——《三阶意识》第三部第二阶意识：社会我}\newline \url{https://mp.weixin.qq.com/s/tCkfkpwNwaFu6eZB_3ByQw}\\[1.2mm]
TOC-057 & 2026-04-10 / 392–397 & \CN{从“裂口”到“投影”——关于拉康的理解，齐泽克与存在几何学的根本分野}\newline \url{https://mp.weixin.qq.com/s/Qfo8PpiOCOCdIlLlukvd6A}\\[1.2mm]
TOC-058 & 2026-04-10 / 398–407 & \CN{三阶意识不是层级升级，而是意识投影过程的“光学调焦”——《三阶意识》全书 总序及阅读指南}\newline \url{https://mp.weixin.qq.com/s/Mep8D-w-Ew3pW_Up-UwIDA}\\[1.2mm]
TOC-059 & 2026-04-12 / 408–417 & \CN{12.1 大叙事为何先崩塌：历史、文明与危机出口的失效（第十二章主体自燃与意 义熵爆：当叙事路径突然断裂）——《三阶意识》}\newline \url{https://mp.weixin.qq.com/s/d-jLjiCyQehVswzJj8CVLg}\\[1.2mm]
TOC-060 & 2026-04-13 / 418–429 & \CN{12.2 叙事下沉、主体下沉：从英雄、小人物到破碎形象（第十二章主体自燃与意 义熵爆：当叙事路径突然断裂）——《三阶意识》}\newline \url{https://mp.weixin.qq.com/s/Wd0sYFbsHiYRlfeG0TfjHQ}\\[1.2mm]
TOC-061 & 2026-04-14 / 430–440 & \CN{12.3 “叙事路径我”的过热与“我”的废墟显现（第十二章主体自燃与意义熵爆： 当叙事路径突然断裂）——《三阶意识》}\newline \url{https://mp.weixin.qq.com/s/59sQEibimg7LcUhogexKKQ}\\[1.2mm]
TOC-062 & 2026-04-16 / 441–449 & \CN{12.4 浪漫的谎言与小说的真实：小说家如何率先走出主体幻觉（第十二章主体自 燃与意义熵爆：当叙事路径突然断裂）——《三阶意识》}\newline \url{https://mp.weixin.qq.com/s/5GiH5XnEz0HupWCNcLz2fA}\\[1.2mm]
TOC-063 & 2026-04-17 / 450–459 & \CN{12.5 “主体叙事”崩塌并不意味导向第三阶：犬儒、虚无以及再缝合（第十二章 主体自燃与意义熵爆：当叙事路径突然断裂）——《三阶意识》}\newline \url{https://mp.weixin.qq.com/s/JMPKbj-KExkO8Y7oQuqC4A}\\[1.2mm]
TOC-064 & 2026-04-18 / 460–468 & \CN{12.6 悬置、凝视、延异：焦点中让出的观察空位（第十二章主体自燃与意义熵爆： 当叙事路径突然断裂）——《三阶意识》}\newline \url{https://mp.weixin.qq.com/s/RiQ49txE7JJOyMQr2w3Azg}\\[1.2mm]
TOC-065 & 2026-04-14 / 469–477 & \CN{为什么必须先解剖“我” ，才能谈“无我”——从西方文学的主体搭建史，重返意 识工程的内部}\newline \url{https://mp.weixin.qq.com/s/SHe3y4JgHkjrt-omTAheGw}\\[1.2mm]
TOC-066 & 2026-04-19 / 478–484 & \CN{13.1 观察为何是革命，而不是修养（第十三章观察革命：从二阶意识到三阶意识 的革命性转向）——《三阶意识》第四部：第三阶意识}\newline \url{https://mp.weixin.qq.com/s/ONtZ9T4oKWHmWQVUgQfoEw}\\[1.2mm]
TOC-067 & 2026-04-19 / 485–490 & \CN{13.2 观察者不是中立者，而是过去本身（第十三章观察革命：从二阶意识到三阶 意识的革命性转向）——《三阶意识》第四部：第三阶意识}\newline \url{https://mp.weixin.qq.com/s/Gn0oFPOonaq-PzUGN65NxQ}\\[1.2mm]
TOC-068 & 2026-04-20 / 491–495 & \CN{13.3 当观察者即被观察者：冲突为何第一次失效（第十三章观察革命：从二阶意 识到三阶意识的革命性转向）——《三阶意识》第四部：第三阶意识}\newline \url{https://mp.weixin.qq.com/s/cXZA5492zHTTeomv5SdWQg}\\[1.2mm]
TOC-069 & 2026-04-20 / 496–500 & \CN{13.4 无选择的觉察：观察不是选择更好的自己（第十三章观察革命：从二阶意识 到三阶意识的革命性转向）——《三阶意识》第四部：第三阶意识}\newline \url{https://mp.weixin.qq.com/s/9fCctNc3SgvfF5E_xud_tw}\\[1.2mm]
TOC-070 & 2026-04-21 / 501–506 & \CN{13.5 从观察到关系：革命为何发生在人际之中（第十三章观察革命：从二阶意识 到三阶意识的革命性转向）——《三阶意识》第四部：第三阶意识}\newline \url{https://mp.weixin.qq.com/s/4fqAJjf8mIVXC6SXhf74Zw}\\[1.2mm]
TOC-071 & 2026-04-21 / 507–512 & \CN{13.6 无中心的在场：三阶意识不是答案，而是清醒在场（第十三章观察革命：从 二阶意识到三阶意识的革命性转向）——《三阶意识》第四部：第三阶意识}\newline \url{https://mp.weixin.qq.com/s/wX_XO733-KNXFIENdGN6LQ}\\[1.2mm]
TOC-072 & 2026-04-29 / 513–520 & \CN{14.1 本体论维度转向：从“实体”到“交织” （第十四章场我如何开始：一场哥白 尼式的意识革命）——《三阶意识》第四部第三阶意识}\newline \url{https://mp.weixin.qq.com/s/qKnbbEAqQokvukQmlP-huw}\\[1.2mm]
TOC-073 & 2026-04-29 / 521–528 & \CN{14.2 认识论维度转向：从“我的意义”到“意义如何发生” （第十四章场我如何 开始：一场哥白尼式的意识革命）——《三阶意识》第四部第三阶意识}\newline \url{https://mp.weixin.qq.com/s/ia73syn189xQI4nLNjozsg}\\[1.2mm]
TOC-074 & 2026-04-30 / 529–536 & \CN{14.3 伦理学维度转向：从“使用他者”到“让他者显现” （第十四章场我如何开 始：一场哥白尼式的意识革命）——《三阶意识》第四部第三阶意识}\newline \url{https://mp.weixin.qq.com/s/S1mSypWZlCkBWftFLpzHUg}\\[1.2mm]
TOC-075 & 2026-05-01 / 537–543 & \CN{14.4 时间现象学维度：从“我将成为”到“显现的现在” （第十四章场我如何开 始：一场哥白尼式的意识革命）——《三阶意识》第四部第三阶意识}\newline \url{https://mp.weixin.qq.com/s/-MHbVUzjPaPGHWfrIxPNDg}\\[1.2mm]
TOC-076 & 2026-05-02 / 544–552 & \CN{14.5 “场我”不是结果，而是持续的在场（第十四章场我如何开始：一场哥白尼 式的意识革命）——《三阶意识》第四部第三阶意识}\newline \url{https://mp.weixin.qq.com/s/p_3hX8U8Gg1dgu15u9ByvQ}\\[1.2mm]
TOC-077 & 2026-05-05 / 553–560 & \CN{15.1 AI 不是智能神话，而是意识结构的放大器（第十五章 AI 革命与三阶意识： 后主体时代的认知主权）——《三阶意识》}\newline \url{https://mp.weixin.qq.com/s/PCkV7VYhZeFKgSBU7v1kkg}\\[1.2mm]
TOC-078 & 2026-05-06 / 561–570 & \CN{15.2 二阶意识的外包时代：当“我最像我”的部分开始失效（第十五章 AI 革命 与三阶意识：后主体时代的认知主权）——《三阶意识》}\newline \url{https://mp.weixin.qq.com/s/CbbvGCR2HCa-ztZSFMRuOQ}\\[1.2mm]
TOC-079 & 2026-05-07 / 571–579 & \CN{15.3 替身我与空位之争：AI 会帮助让位，还是制造更强的新我（第十五章 AI 革 命与三阶意识：后主体时代的认知主权）——《三阶意识》}\newline \url{https://mp.weixin.qq.com/s/8YoWBS878odDr7PgucwoKQ}\\[1.2mm]
TOC-080 & 2026-05-08 / 580–589 & \CN{15.4 不可外包之物：三阶意识的现实门槛（第十五章 AI 革命与三阶意识：后主 体时代的认知主权）——《三阶意识》}\newline \url{https://mp.weixin.qq.com/s/24_ygT0lF20HBpehIychZw}\\[1.2mm]
TOC-081 & 2026-05-08 / 590–598 & \CN{15.5 认知主权的重写：不是坚持自我，而是识别自我如何被制造（第十五章 AI 革命与三阶意识：后主体时代的认知主权）——《三阶意识》}\newline \url{https://mp.weixin.qq.com/s/BJwsey0TzdfhiwUDTs1_bQ}\\[1.2mm]
TOC-082 & 2026-05-09 / 599–607 & \CN{15.6 自由不是主体的奖赏：AI 时代对自由的重新定义（第十五章 AI 革命与三阶 意识：后主体时代的认知主权）——《三阶意识》}\newline \url{https://mp.weixin.qq.com/s/lVX_5Tf0Bo5cfjQGMwPsYA}\\[1.2mm]
TOC-083 & 2026-05-10 / 608–619 & \CN{后记：三阶意识不是三种人，而是意识的三种显现方式——《三阶意识》}\newline \url{https://mp.weixin.qq.com/s/8ZkwCzxbMWXPqecABzJLMg}\\[1.2mm]
\bottomrule
\end{longtable}
\endgroup

\begin{plainbox}[Scope discipline]
The 83 entries above correspond to TOC-001--TOC-083 in the 619-page archive. Other \emph{Snowfall on the Clock} articles cited directly in the main text have their full permanent WeChat links printed in the reference entries below.
\end{plainbox}


\renewcommand{\bibname}{References}
\begin{thebibliography}{99}

\bibitem[Baddeley(1992)]{baddeley1992}
A. Baddeley.
\newblock Working memory.
\newblock \emph{Science}, 255(5044):556--559, 1992.
\newblock doi:10.1126/science.1736359.

\bibitem[Bartholdi et~al.(2022)]{bnz2022}
L. Bartholdi, V. Nekrashevych, and T. Zheng.
\newblock Growth of groups with linear Schreier graphs.
\newblock 2022, arXiv:2205.01792.

\bibitem[Berger and Luckmann(1966)]{berger1966}
P. L. Berger and T. Luckmann.
\newblock \emph{The Social Construction of Reality: A Treatise in the Sociology of Knowledge}.
\newblock Doubleday, Garden City, NY, 1966.

\bibitem[Bernstein et~al.(2015)]{bernstein2015}
A. Bernstein, Y. Hadash, Y. Lichtash, G. Tanay, K. Shepherd, and D. M. Fresco.
\newblock Decentering and related constructs: A critical review and metacognitive processes model.
\newblock \emph{Perspectives on Psychological Science}, 10(5):599--617, 2015.
\newblock doi:10.1177/1745691615594577.

\bibitem[Bratteli(1972)]{bratteli1972}
O. Bratteli.
\newblock Inductive limits of finite dimensional $C^*$-algebras.
\newblock \emph{Transactions of the American Mathematical Society}, 171:195--234, 1972.
\newblock doi:10.1090/S0002-9947-1972-0312282-2.

\bibitem[Braun et~al.(2018)]{braun2018}
N. Braun, S. Debener, N. Spychala, E. Bongartz, P. Sörös, H. H. O. Müller, and A. Philipsen.
\newblock The senses of agency and ownership: A review.
\newblock \emph{Frontiers in Psychology}, 9:535, 2018.
\newblock doi:10.3389/fpsyg.2018.00535.

\bibitem[Bruner(1987)]{bruner1987}
J. Bruner.
\newblock Life as narrative.
\newblock \emph{Social Research}, 54(1):11--32, 1987.

\bibitem[Cantu(2019)]{cantu2019}
J. C. Cantu.
\newblock \emph{Periodic Groups via Orbital Graphs}.
\newblock Doctoral dissertation, Texas A\&M University, 2019.

\bibitem[Chalmers(1995)]{chalmers1995}
D. J. Chalmers.
\newblock Facing up to the problem of consciousness.
\newblock \emph{Journal of Consciousness Studies}, 2(3):200--219, 1995.

\bibitem[Conway and Pleydell-Pearce(2000)]{conway2000}
M. A. Conway and C. W. Pleydell-Pearce.
\newblock The construction of autobiographical memories in the self-memory system.
\newblock \emph{Psychological Review}, 107(2):261--288, 2000.
\newblock doi:10.1037/0033-295X.107.2.261.

\bibitem[Cooley(1902)]{cooley1902}
C. H. Cooley.
\newblock \emph{Human Nature and the Social Order}.
\newblock Charles Scribner's Sons, New York, 1902.

\bibitem[Craig(2009)]{craig2009}
A. D. Craig.
\newblock How do you feel---now? The anterior insula and human awareness.
\newblock \emph{Nature Reviews Neuroscience}, 10:59--70, 2009.
\newblock doi:10.1038/nrn2555.

\bibitem[Crowell(2012)]{crowell2012}
S. Crowell, editor.
\newblock \emph{The Cambridge Companion to Existentialism}.
\newblock Cambridge University Press, Cambridge, 2012.

\bibitem[Exel(1998)]{exel1998}
R. Exel.
\newblock Partial actions of groups and actions of inverse semigroups.
\newblock \emph{Proceedings of the American Mathematical Society}, 126(12):3481--3494, 1998.
\newblock doi:10.1090/S0002-9939-98-04575-3.

\bibitem[Exel(2008)]{exel2008}
R. Exel.
\newblock Inverse semigroups and combinatorial $C^*$-algebras.
\newblock \emph{Bulletin of the Brazilian Mathematical Society}, 39:191--313, 2008.
\newblock doi:10.1007/s00574-008-0080-7.

\bibitem[Fleming and Dolan(2012)]{fleming2012}
S. M. Fleming and R. J. Dolan.
\newblock The neural basis of metacognitive ability.
\newblock \emph{Philosophical Transactions of the Royal Society B}, 367(1594):1338--1349, 2012.
\newblock doi:10.1098/rstb.2011.0417.

\bibitem[Fleming et~al.(2012)]{fleming2012a}
S. M. Fleming, R. J. Dolan, and C. D. Frith.
\newblock Metacognition: computation, biology and function.
\newblock \emph{Philosophical Transactions of the Royal Society B}, 367(1594):1280--1286, 2012.
\newblock doi:10.1098/rstb.2012.0021.

\bibitem[Flynn(2006)]{flynn2006}
T. R. Flynn.
\newblock \emph{Existentialism: A Very Short Introduction}.
\newblock Oxford University Press, Oxford, 2006.

\bibitem[von Foerster(2003)]{vonfoerster2003}
H. von Foerster.
\newblock \emph{Understanding Understanding: Essays on Cybernetics and Cognition}.
\newblock Springer, New York, 2003.
\newblock doi:10.1007/b97451.

\bibitem[Foucault(1978)]{foucault1978}
M. Foucault.
\newblock \emph{The History of Sexuality, Volume 1: An Introduction}.
\newblock Translated by R. Hurley. Pantheon, New York, 1978.

\bibitem[Gallagher(2000)]{gallagher2000}
S. Gallagher.
\newblock Philosophical conceptions of the self: implications for cognitive science.
\newblock \emph{Trends in Cognitive Sciences}, 4(1):14--21, 2000.
\newblock doi:10.1016/S1364-6613(99)01417-5.

\bibitem[Gillespie(1995)]{gillespie1995}
M. A. Gillespie.
\newblock \emph{Nihilism Before Nietzsche}.
\newblock University of Chicago Press, Chicago, 1995.

\bibitem[Giordano et~al.(1995)]{gps1995}
T. Giordano, I. F. Putnam, and C. F. Skau.
\newblock Topological orbit equivalence and $C^*$-crossed products.
\newblock \emph{Journal für die reine und angewandte Mathematik}, 469:51--111, 1995.

\bibitem[Girard(1965)]{girard1965}
R. Girard.
\newblock \emph{Deceit, Desire, and the Novel: Self and Other in Literary Structure}.
\newblock Johns Hopkins University Press, Baltimore, 1965.

\bibitem[Goffman(1959)]{goffman1959}
E. Goffman.
\newblock \emph{The Presentation of Self in Everyday Life}.
\newblock Anchor Books, New York, 1959.

\bibitem[Kross and Ayduk(2011)]{kross2011}
E. Kross and O. Ayduk.
\newblock Making meaning out of negative experiences by self-distancing.
\newblock \emph{Current Directions in Psychological Science}, 20(3):187--191, 2011.
\newblock doi:10.1177/0963721411408883.

\bibitem[Krishnamurti(1954)]{krishnamurti1954}
J. Krishnamurti.
\newblock \emph{The First and Last Freedom}.
\newblock Harper, New York, 1954.

\bibitem[Krishnamurti(1969)]{krishnamurti1969}
J. Krishnamurti.
\newblock \emph{Freedom from the Known}.
\newblock Harper \& Row, New York, 1969.

\bibitem[Krishnamurti and Bohm(1985)]{krishnamurtiBohm1985}
J. Krishnamurti and D. Bohm.
\newblock \emph{The Ending of Time}.
\newblock Gollancz, London, 1985.

\bibitem[Krishnamurti Foundation Trust(2026a)]{kft2026a}
Krishnamurti Foundation Trust.
\newblock \emph{The Observer and the Observed}.
\newblock Curated archive of Krishnamurti talks, accessed August 27, 2026.
\newblock \url{https://kfoundation.org/the-observer-and-the-observed/}.

\bibitem[Krishnamurti Foundation Trust(2026b)]{kft2026b}
Krishnamurti Foundation Trust.
\newblock Public Talk 5, New Delhi, 24 December 1970.
\newblock Transcript, accessed August 27, 2026.
\newblock \url{https://kfoundation.org/transcript/public-talk-5-new-delhi-24-december-1970/}.

\bibitem[Kuang(2024, v1)]{kuang2024v1}
Z. Kuang.
\newblock Growth of groups with incompressible elements, I.
\newblock arXiv:2402.16238v1, 26 February 2024. See especially Subsection~5.2 (multiple types of bridges) and Example~5.8.

\bibitem[Kuang(2025)]{kuang2025tiling}
Z. Kuang.
\newblock Growth of finitely generated subgroups of the topological full groups of inverse semigroups of bounded type.
\newblock arXiv:2409.19491v2, revised 29 May 2025. See especially Section~3 (expanding tile inflations).
\newblock doi:10.48550/arXiv.2409.19491.

\bibitem[Kuang(2026a)]{kuang2026a}
Z. Kuang.
\newblock Growth of groups with incompressible elements, I.
\newblock \emph{Groups, Geometry, and Dynamics}, published online first 21 April 2026.
\newblock doi:10.4171/GGD/954.

\bibitem[Kuang(2026b)]{kuang2026b}
Z. Kuang.
\newblock Near full groups of bounded type: full group completion.
\newblock arXiv:2607.25729v2, 21 August 2026.
\newblock doi:10.48550/arXiv.2607.25729.

\bibitem[Kurby and Zacks(2008)]{kurby2008}
C. A. Kurby and J. M. Zacks.
\newblock Segmentation in the perception and memory of events.
\newblock \emph{Trends in Cognitive Sciences}, 12(2):72--79, 2008.
\newblock doi:10.1016/j.tics.2007.11.004.

\bibitem[Lacan(2006)]{lacan2006}
J. Lacan.
\newblock \emph{Écrits: The First Complete Edition in English}.
\newblock Translated by B. Fink. W. W. Norton, New York, 2006.

\bibitem[Matui(2012)]{matui2012}
H. Matui.
\newblock Homology and topological full groups of étale groupoids on totally disconnected spaces.
\newblock \emph{Proceedings of the London Mathematical Society}, 104(1):27--56, 2012.
\newblock doi:10.1112/plms/pdr029.

\bibitem[McAdams and McLean(2013)]{mcadams2013}
D. P. McAdams and K. C. McLean.
\newblock Narrative identity.
\newblock \emph{Current Directions in Psychological Science}, 22(3):233--238, 2013.
\newblock doi:10.1177/0963721413475622.

\bibitem[Mead(1934)]{mead1934}
G. H. Mead.
\newblock \emph{Mind, Self, and Society}.
\newblock University of Chicago Press, Chicago, 1934.

\bibitem[Metzinger(2003)]{metzinger2003}
T. Metzinger.
\newblock \emph{Being No One: The Self-Model Theory of Subjectivity}.
\newblock MIT Press, Cambridge, MA, 2003.

\bibitem[Nagel(1974)]{nagel1974}
T. Nagel.
\newblock What is it like to be a bat?
\newblock \emph{The Philosophical Review}, 83(4):435--450, 1974.
\newblock doi:10.2307/2183914.

\bibitem[Nekrashevych(2018)]{nek2018}
V. Nekrashevych.
\newblock Palindromic subshifts and simple periodic groups of intermediate growth.
\newblock \emph{Annals of Mathematics}, 187(3):667--719, 2018.
\newblock doi:10.4007/annals.2018.187.3.2.

\bibitem[Nekrashevych(2022)]{nek2022}
V. Nekrashevych.
\newblock \emph{Groups and Topological Dynamics}.
\newblock Graduate Studies in Mathematics, vol.~223. American Mathematical Society, Providence, RI, 2022.

\bibitem[Nelson and Fivush(2004)]{nelson2004}
K. Nelson and R. Fivush.
\newblock The emergence of autobiographical memory: A social cultural developmental theory.
\newblock \emph{Psychological Review}, 111(2):486--511, 2004.
\newblock doi:10.1037/0033-295X.111.2.486.

\bibitem[Nietzsche(1974)]{nietzsche1974}
F. Nietzsche.
\newblock \emph{The Gay Science}.
\newblock Translated by W. Kaufmann. Vintage Books, New York, 1974. Original German editions 1882 and 1887.

\bibitem[Nietzsche(1997)]{nietzsche1997}
F. Nietzsche.
\newblock \emph{On the Genealogy of Morality}.
\newblock Translated by C. Diethe. Cambridge University Press, Cambridge, 1997.

\bibitem[Nolen-Hoeksema et~al.(2008)]{nolen2008}
S. Nolen-Hoeksema, B. E. Wisco, and S. Lyubomirsky.
\newblock Rethinking rumination.
\newblock \emph{Perspectives on Psychological Science}, 3(5):400--424, 2008.
\newblock doi:10.1111/j.1745-6924.2008.00088.x.

\bibitem[Putnam(2010)]{putnam2010}
I. F. Putnam.
\newblock Orbit equivalence of Cantor minimal systems: A survey and a new proof.
\newblock \emph{Expositiones Mathematicae}, 28(2):101--131, 2010.
\newblock doi:10.1016/j.exmath.2009.06.002.

\bibitem[Renault(1980)]{renault1980}
J. Renault.
\newblock \emph{A Groupoid Approach to $C^*$-Algebras}.
\newblock Lecture Notes in Mathematics 793. Springer, Berlin, 1980.

\bibitem[Reynolds and Renaudie(2024)]{reynolds2024}
J. Reynolds and P.-J. Renaudie.
\newblock Jean-Paul Sartre.
\newblock In E. N. Zalta and U. Nodelman, editors, \emph{The Stanford Encyclopedia of Philosophy}, Fall 2024 edition. Metaphysics Research Lab, Stanford University, 2024.
\newblock \url{https://plato.stanford.edu/archives/fall2024/entries/sartre/}.

\bibitem[Ricoeur(1992)]{ricoeur1992}
P. Ricoeur.
\newblock \emph{Oneself as Another}.
\newblock Translated by K. Blamey. University of Chicago Press, Chicago, 1992.

\bibitem[Rosenthal(2005)]{rosenthal2005}
D. M. Rosenthal.
\newblock \emph{Consciousness and Mind}.
\newblock Clarendon Press, Oxford, 2005.

\bibitem[Sartre(2007)]{sartre2007}
J.-P. Sartre.
\newblock \emph{Existentialism Is a Humanism}.
\newblock Translated by C. Macomber. Yale University Press, New Haven, 2007. Original public lecture delivered in Paris in late 1945; French publication 1946.

\bibitem[Schopenhauer(1969)]{schopenhauer1969}
A. Schopenhauer.
\newblock \emph{The World as Will and Representation}, volume~1.
\newblock Translated by E. F. J. Payne. Dover, New York, 1969.

\bibitem[Seth(2015)]{seth2015}
A. K. Seth.
\newblock The cybernetic Bayesian brain: From interoceptive inference to sensorimotor contingencies.
\newblock In T. Metzinger and J. M. Windt, editors, \emph{Open MIND}, 35(T). MIND Group, Frankfurt am Main, 2015.
\newblock doi:10.15502/9783958570108.

\bibitem[Siderits(2007)]{siderits2007}
M. Siderits.
\newblock \emph{Buddhism as Philosophy: An Introduction}.
\newblock Ashgate, Aldershot, 2007.

\bibitem[Sidki(2000)]{sidki2000}
S. N. Sidki.
\newblock Automorphisms of one-rooted trees: growth, circuit structure, and acyclicity.
\newblock \emph{Journal of Mathematical Sciences (New York)}, 100(1):1925--1943, 2000.
\newblock doi:10.1007/BF02677504.

\bibitem[Tulving and Pearlstone(1966)]{tulving1966}
E. Tulving and Z. Pearlstone.
\newblock Availability versus accessibility of information in memory for words.
\newblock \emph{Journal of Verbal Learning and Verbal Behavior}, 5(4):381--391, 1966.
\newblock doi:10.1016/S0022-5371(66)80048-8.

\bibitem[Vygotsky(1986)]{vygotsky1986}
L. S. Vygotsky.
\newblock \emph{Thought and Language}.
\newblock Revised edition, edited by A. Kozulin. MIT Press, Cambridge, MA, 1986. Original Russian publication 1934.

\bibitem[Watkins(2008)]{watkins2008}
E. R. Watkins.
\newblock Constructive and unconstructive repetitive thought.
\newblock \emph{Psychological Bulletin}, 134(2):163--206, 2008.
\newblock doi:10.1037/0033-2909.134.2.163.

\bibitem[Zacks et~al.(2007)]{zacks2007}
J. M. Zacks, N. K. Speer, K. M. Swallow, T. S. Braver, and J. R. Reynolds.
\newblock Event perception: a mind-brain perspective.
\newblock \emph{Psychological Bulletin}, 133(2):273--293, 2007.
\newblock doi:10.1037/0033-2909.133.2.273.

\bibitem[Zahavi(2005)]{zahavi2005}
D. Zahavi.
\newblock \emph{Subjectivity and Selfhood: Investigating the First-Person Perspective}.
\newblock MIT Press, Cambridge, MA, 2005.

\bibitem[Krishnamurti Meditation Workshop(2026)]{krishnamurtiSorrow2026}
\CN{克里希那穆提冥思坊}.
\newblock \CN{《【言谈录】与痛苦共处，就像怀抱一件奇珍异宝》} [\emph{Living with Sorrow as If Holding a Rare Treasure}].
\newblock WeChat Official Account \CN{“克里希那穆提冥思坊”}, August 29, 2026; excerpted from \emph{Meeting Life}, ``The Ending of Sorrow is Love.''
\newblock \url{https://mp.weixin.qq.com/s/JMp86TVrguTFgyEJgAsTPA}.

\bibitem[Peng(2026)]{toc2026}
\href{https://www.amazon.ca/Guoxin-Peng/e/B0GLHTBNVC/ref=dp_byline_cont_ebooks_1}{Guoxin Peng}.
\newblock \emph{Third-Order Consciousness: From the Biological Self and Social Self to Vacant Consciousness (The Geometry of Existence Book 4)}.
\newblock Edited by \href{https://www.amazon.ca/s/ref=dp_byline_sr_ebooks_2?ie=UTF8&field-author=Pinchuan+Xie&text=Pinchuan+Xie&sort=relevancerank&search-alias=digital-text}{Pinchuan Xie}. Independently published, 2026. Amazon ASIN B0H7SF3FDN.
\newblock \url{https://www.amazon.ca/dp/B0H7SF3FDN}.

\bibitem[Kuang(2026)]{zenodo22202015}
Kuang, Z.
\newblock \emph{Third-Order Consciousness and AF-by-Discrete: Bilingual Research Edition}.
\newblock Zenodo, 2026. DOI: \href{https://doi.org/10.5281/zenodo.22202015}{10.5281/zenodo.22202015}.

\bibitem[Snowfall on the Clock(2026b)]{translation2026}
\CN{钟上的降雪} (Peng Guoxin \& AI).
\newblock \CN{《存在几何学的翻译协议——跨学科、跨语言、跨文理与跨文明思维路径的界面方法》} [\emph{The Translation Protocol of the Geometry of Existence---An Interface Method for Cross-Disciplinary, Cross-Linguistic, Cross-Arts-and-Sciences, and Cross-Civilizational Thought}].
\newblock WeChat Official Account \CN{“钟上的降雪”}, July 22, 2026.
\newblock \url{https://mp.weixin.qq.com/s/o5s5PQ4A-sewowwQODx4WQ}.

\bibitem[Snowfall on the Clock(2026c)]{geometry2026}
\CN{钟上的降雪} (Peng Guoxin \& AI).
\newblock \CN{《在结构堵塞与结构越位之间——符号堆积、意义通胀时代的关系疏通与维度切换》} [\emph{Between Structural Blockage and Structural Overreach---Relational Unblocking and Dimensional Switching in an Age of Symbolic Accumulation and Meaning Inflation}].
\newblock WeChat Official Account \CN{“钟上的降雪”}, August 27, 2026.
\newblock \url{https://mp.weixin.qq.com/s/Vc-fYyX4s1ZpwxNYEZmXWQ}.

\bibitem[Snowfall on the Clock(2026d)]{projection2026}
\CN{钟上的降雪} (Peng Guoxin \& AI).
\newblock \CN{《当“神像”倒下以后，我们看见了什么？——从仲树“塌房”谈投影、路径与毁神机制》} [\emph{What Do We See after the ``Idol'' Falls?---Projection, Path, and the Mechanism of De-Idolization}].
\newblock WeChat Official Account \CN{“钟上的降雪”}, August 26, 2026.
\newblock \url{https://mp.weixin.qq.com/s/nr_ZrhTEnkW6b76lExS6qQ}.

\bibitem[Snowfall on the Clock(2026e)]{snowfall2026e}
\CN{钟上的降雪} (Peng Guoxin \& AI).
\newblock \CN{《主体是生命投影结构逼出的奇点——从哥德尔的不完备、佩雷尔曼的手术到三阶意识的维度切换》} [\emph{The Subject as a Singularity Forced by the Projective Structure of Life}].
\newblock WeChat Official Account \CN{“钟上的降雪”}, June 13, 2026.
\newblock \url{https://mp.weixin.qq.com/s/T4Meetx--gL0ExvP0GF2UA}.

\bibitem[Snowfall on the Clock(2026f)]{snowfall2026f}
\CN{钟上的降雪} (Peng Guoxin \& AI).
\newblock \CN{《（六）奇点不是缺陷，而是结构无法继续原样运作之处》} [\emph{(VI) A Singularity Is Not a Defect, but Where Structure Can No Longer Continue Unchanged}].
\newblock WeChat Official Account \CN{“钟上的降雪”}, June 24, 2026.
\newblock \url{https://mp.weixin.qq.com/s/XEP3uvPGV0lFkSQMYBBIag}.

\bibitem[Snowfall on the Clock(2026g)]{snowfall2026g}
\CN{钟上的降雪} (Peng Guoxin \& AI).
\newblock \CN{《从黑洞边缘态到意识拓扑：真正重要的，也许不是内部，而是边界》} [\emph{From Black-Hole Edge States to the Topology of Consciousness: Perhaps What Matters Is Not the Interior but the Boundary}].
\newblock WeChat Official Account \CN{“钟上的降雪”}, June 25, 2026.
\newblock \url{https://mp.weixin.qq.com/s/JS_WNdCRBMDdAxJe-AGxeA}.

\bibitem[Snowfall on the Clock(2026h)]{snowfall2026h}
\CN{钟上的降雪} (Peng Guoxin \& AI).
\newblock \CN{《从主体到场：节点自觉如何生成去中心化秩序——关于场意识、空位差异与自组织的存在几何学》} [\emph{From Subject to Field: How Node Awareness Generates Decentered Order}].
\newblock WeChat Official Account \CN{“钟上的降雪”}, July 24, 2026.
\newblock \url{https://mp.weixin.qq.com/s/uarbLlu0WI7zVBps_luqhA}.

\end{thebibliography}

\end{document}
